% Options for packages loaded elsewhere
\PassOptionsToPackage{unicode}{hyperref}
\PassOptionsToPackage{hyphens}{url}
\documentclass[
]{article}
\usepackage{xcolor}
\usepackage{amsmath,amssymb}
\setcounter{secnumdepth}{-\maxdimen} % remove section numbering
\usepackage{iftex}
\ifPDFTeX
  \usepackage[T1]{fontenc}
  \usepackage[utf8]{inputenc}
  \usepackage{textcomp} % provide euro and other symbols
\else % if luatex or xetex
  \usepackage{unicode-math} % this also loads fontspec
  \defaultfontfeatures{Scale=MatchLowercase}
  \defaultfontfeatures[\rmfamily]{Ligatures=TeX,Scale=1}
\fi
\usepackage{lmodern}
\ifPDFTeX\else
  % xetex/luatex font selection
\fi
% Use upquote if available, for straight quotes in verbatim environments
\IfFileExists{upquote.sty}{\usepackage{upquote}}{}
\IfFileExists{microtype.sty}{% use microtype if available
  \usepackage[]{microtype}
  \UseMicrotypeSet[protrusion]{basicmath} % disable protrusion for tt fonts
}{}
\makeatletter
\@ifundefined{KOMAClassName}{% if non-KOMA class
  \IfFileExists{parskip.sty}{%
    \usepackage{parskip}
  }{% else
    \setlength{\parindent}{0pt}
    \setlength{\parskip}{6pt plus 2pt minus 1pt}}
}{% if KOMA class
  \KOMAoptions{parskip=half}}
\makeatother
\usepackage{longtable,booktabs,array}
\newcounter{none} % for unnumbered tables
\usepackage{calc} % for calculating minipage widths
% Correct order of tables after \paragraph or \subparagraph
\usepackage{etoolbox}
\makeatletter
\patchcmd\longtable{\par}{\if@noskipsec\mbox{}\fi\par}{}{}
\makeatother
% Allow footnotes in longtable head/foot
\IfFileExists{footnotehyper.sty}{\usepackage{footnotehyper}}{\usepackage{footnote}}
\makesavenoteenv{longtable}
\usepackage{graphicx}
\makeatletter
\newsavebox\pandoc@box
\newcommand*\pandocbounded[1]{% scales image to fit in text height/width
  \sbox\pandoc@box{#1}%
  \Gscale@div\@tempa{\textheight}{\dimexpr\ht\pandoc@box+\dp\pandoc@box\relax}%
  \Gscale@div\@tempb{\linewidth}{\wd\pandoc@box}%
  \ifdim\@tempb\p@<\@tempa\p@\let\@tempa\@tempb\fi% select the smaller of both
  \ifdim\@tempa\p@<\p@\scalebox{\@tempa}{\usebox\pandoc@box}%
  \else\usebox{\pandoc@box}%
  \fi%
}
% Set default figure placement to htbp
\def\fps@figure{htbp}
\makeatother
\setlength{\emergencystretch}{3em} % prevent overfull lines
\providecommand{\tightlist}{%
  \setlength{\itemsep}{0pt}\setlength{\parskip}{0pt}}
\usepackage{bookmark}
\IfFileExists{xurl.sty}{\usepackage{xurl}}{} % add URL line breaks if available
\urlstyle{same}
\hypersetup{
  hidelinks,
  pdfcreator={LaTeX via pandoc}}

\author{}
\date{}

\begin{document}

ЧАСТКОВО ДИФЕРЕНЦІАЛ РІВНЯННЯ для Вчені та Інженери

Стенлі Дж. Фарлоу, професор математики, Університет Мену

\section{С. Фарлоу}\label{ux441.-ux444ux430ux440ux43bux43eux443}

\section{РІВНЯННЯ З ПОХІДНИМИ ПОХІДНИМИ ДЛЯ НАУКОВЦІВ І
ІНЖЕНЕРІВ}\label{ux440ux456ux432ux43dux44fux43dux43dux44f-ux437-ux43fux43eux445ux456ux434ux43dux438ux43cux438-ux43fux43eux445ux456ux434ux43dux438ux43cux438-ux434ux43bux44f-ux43dux430ux443ux43aux43eux432ux446ux456ux432-ux456-ux456ux43dux436ux435ux43dux435ux440ux456ux432}

Перекладено з англійської \textbf{A.I.} Плісом

Редактор

С. І. Похожаєва

5\textbar{} 66K 22.161.6 φ24 υ ΥΠΚ 517.2

\section{Фарлоу С.}\label{ux444ux430ux440ux43bux43eux443-ux441.}

F 24 Рівняння з частими похідними для науковців та інженерів. Москва,
видавництво Мир, 1985, 384 с.

King's American Mathematician, який є підручником з теорії
диференціальних рівнянь у похідних похідних. Він характеризується
компактністю, ясністю та ясністю подачі, а також неформальним підходом
до подачі матеріалу. Вона містить багато ілюстрацій, графіків і діаграм,
але замість суворих доказів часто даються роздуми, засновані на інтуїції
або аналогіях.

Для інженерів і нематематиків --- біологів, диміків, а також студентів
університету.
\(\Phi = \frac{1702050000-380}{041(01)-85} 26-85, \text{ q. 1}\) BBK
22.161.6

Редакційна колегія журналу Literature on Mathematical Sciences

© переклад російською, Мир, 1985

© John Wiley \& Sons, Inc.~1982. Всі права захищені. Авторизований
переклад з англомовного видання, опублікований John Wiley \& Sons, Inc.

\section{ВІД РЕДАКТОРА
ПЕРЕКЛАДУ}\label{ux432ux456ux434-ux440ux435ux434ux430ux43aux442ux43eux440ux430-ux43fux435ux440ux435ux43aux43bux430ux434ux443}

Книга, запропонована для читачів, слід розглядати як підручник для
курсів: «Диференціальні рівняння в похідних похідних» та «Рівняння
математичної фізики» для студентів технічних університетів. Вона також
буде корисною для інженерів різних спеціальностей. У книзі представлено
традиційні методи розв'язання диференціальних рівнянь у похідних
похідних, включаючи деякі чисельні методи. Водночас особлива увага
приділяється методу розділення змінних і методу інтегральних
перетворень, і цю частину книги можна використовувати на практичних
заняттях в університетах

Орієнтація книги на широку «нематематичну» аудиторію визначила стиль
викладання розглянутих питань: інтуїтивний підхід і велика увага до
фізичного значення не лише самих рівнянь, а й меж і початкових умов для
різних задач. Вважається, що читачі знайомі лише з основами
диференціального та інтегрального числення та елементами теорії
звичайних диференціальних рівнянь.

Фахівці можуть не задовольнитися інтуїтивним описом певних понять і
методів автором, але більш суворий виклад неминуче призведе до
ускладнення книги і звузить коло її читачів, що, очевидно, не було
наміром автора.

Більшим запереченням може бути відсутність загальних результатів щодо
задач для диференціальних рівнянь у похідних з змінними коефіцієнтами,
включно з рівняннями вищого порядку. Однак читач має чітко розуміти, що
ця книга є вступним курсом з теорії диференціальних рівнянь у часткових
похідних. Питання лінійної теорії диференціальних рівнянь, які автор не
включив, можна знайти в літературі, наведеній наприкінці книги.

Матеріал книги подається автором з великою методологічною майстерністю.
Однак не все було однаково успішним, і винний читач знайде це,
наприклад, у вивченні принципу Дюамеля, сферичних заповітів та деяких
інших уривків у книзі.

Без сумніву, запропонована книга С. Фарлоу знайде читача не лише серед
студентів і інженерів, а й серед викладачів, які можуть багато чого
навчитися з конкретного способу презентації та розробки спеціальних
математичних курсів.

\section{ПЕРЕДМОВА}\label{ux43fux435ux440ux435ux434ux43cux43eux432ux430}

За останні кілька років спостерігається значне зростання кількості
учнів, які починають вивчати диференціальні рівняння з похідними у своїх
одинадцятих класах. Крім того, багато з них спеціалізуються на сферах,
де інтуїція важливіша за математичну строгість. Тому, пишучи цю книгу, я
прагнув створити текст, який стимулює інтуїтивне мислення читачів, не
втрачаючи при цьому значної математичної строгості. Можливо, на одному
полюсі, подати предмет на високому рівні «епсилон-дельта», але тоді, як
це зазвичай буває, студенти не знатимуть, що робити з отриманою
інформацією. Інша крайність --- повне ігнорування математичних тонкощів,
але тоді ні учні, ні вчителі не знатимуть, як діяти. Я намагався знайти
розумний компроміс між цими крайнощами.

Ця книга виникла з курсу лекцій, які я читав протягом останніх п'яти
років. Розташування матеріалу не є загальноприйнятим --- вони майже
незалежні один від одного

Друг лекцій.

Основна увага приділяється двом найважливішим аналітичним методам:
розділенню змінних та інтегральним перетворенням. Також представлені
деякі нестандартні розділи, такі як метод Монте-Карло, варіаційне
числення, теорія керування, теорія потенціалів, інтегральні рівняння. Це
пов'язано з тим, що більшості студентів доведеться стикатися з цими
питаннями під час навчання, і якщо вони не ознайомляться з ними зараз,
навряд чи коли-небудь систематично їх вивчатимуть.

Цю книгу можна використовувати як навчальний посібник для одно- або
двосеместрового курсу як для початківців, так і для старших студентів.
Читачі повинні знати лише диференціальне та інтегральне числення та
звичайні диференціальні рівняння. Матеріал більшості лекцій подається в
одному-двох уроках. Типовий односеместровий курс може включати лекції
1-13, 15-17, 19-20, 22-23, 25-27, 30-32, 37-39. Усі 47 лекцій легко
читати протягом двох семестрів. Водночас часу буде достатньо для
вирішення проблем.

Автор вдячний Вайлі за пропозицію написати книгу, а також рецензентам
--- професорам Крісу Рорресу та М. Куршиду Алі, які зробили низку
корисних коментарів до тексту книги. Будь-які пропозиції щодо подальшого
покращення цього підручника будуть вдячні як від учнів, так і від
викладачів. Я також вдячний Дороті, Сюзан, Александер і Дезі Фарлоу.

Стенлі Дж. Фарлоу

\section{ВСТУП ДО ТЕОРІЇ ДИФЕРЕНЦІАЛЬНИХ РІВНЯНЬ У ПОХІДНИХ
ПОХІДНИХ}\label{ux432ux441ux442ux443ux43f-ux434ux43e-ux442ux435ux43eux440ux456ux457-ux434ux438ux444ux435ux440ux435ux43dux446ux456ux430ux43bux44cux43dux438ux445-ux440ux456ux432ux43dux44fux43dux44c-ux443-ux43fux43eux445ux456ux434ux43dux438ux445-ux43fux43eux445ux456ux434ux43dux438ux445}

МЕТА ЛЕКЦІЇ: Показати, що таке диференціальні рівняння в похідних
похідних, де і як вони виникають і як їх розв'язують. Коротко розгляньте
методи класифікації рівнянь і надайте список базових понять, які будуть
детально вивчені на наступних лекціях.

Більшість фізичних явищ у таких сферах, як гідродинаміка, електрика та
магнетизм, механіка, оптика та теплопередача, загалом можна описати за
допомогою рівнянь у похідних похідних (PNP). Більшість рівнянь у
математичній фізиці є диференціальними рівняннями з похідними похідними.
Правда, з деякими спрощеними припущеннями ці рівняння зводяться до
звичайних диференціальних рівнянь, але повний опис таких систем неминуче
призводить до використання диференціальних рівнянь у похідних похідних.

\section{Що таке диференціальні рівняння у похідних
похідних?}\label{ux449ux43e-ux442ux430ux43aux435-ux434ux438ux444ux435ux440ux435ux43dux446ux456ux430ux43bux44cux43dux456-ux440ux456ux432ux43dux44fux43dux43dux44f-ux443-ux43fux43eux445ux456ux434ux43dux438ux445-ux43fux43eux445ux456ux434ux43dux438ux445}

Рівняння з частковими похідними --- це рівняння, яке містить
диференціали у похідних похідних. На відміну від звичайних
диференціальних рівнянь (TAC), де невідома функція залежить лише від
однієї змінної, у рівняннях часткових похідних вона залежить від кількох
змінних (наприклад, температура \(u(x, t)\) залежить від координати x і
часу t).

Тепер наведемо деякі з найважливіших диференціальних рівнянь у похідних
похідних. Для спрощення запису ми використаємо наступні нотації

{[} u\_t \textbackslash equiv
\textbackslash frac\{\textbackslash partial u\}\{\textbackslash partial
t\} {]}

, \(u_x \equiv \frac{\partial u}{\partial x}\) ,
\(u_{xx} \equiv \frac{\partial^2 u}{\partial x^2}\) , \ldots{}

\section{Деякі диференціальні рівняння в похідних
похідних}\label{ux434ux435ux44fux43aux456-ux434ux438ux444ux435ux440ux435ux43dux446ux456ux430ux43bux44cux43dux456-ux440ux456ux432ux43dux44fux43dux43dux44f-ux432-ux43fux43eux445ux456ux434ux43dux438ux445-ux43fux43eux445ux456ux434ux43dux438ux445}

{[} u\_t = u\_\{xx\} {]}

(одновимірне рівняння теплопровідності), \(u_t = u_{xx} + u_{yy}\)
(двовимірне рівняння теплопровідності),
\(u_{rr} + \frac{1}{r}u_r + \frac{1}{r^2}u_{00} = 0\) (рівняння Лапласа
за полярними координатами), \(u_{tt} = u_{xx} + u_{yy} + u_{zz}\)
(тривимірне хвильове рівняння), \(u_{tt} = u_{xx} + zu_t + \beta u\)
(телеграфне рівняння).

\section{Примітки щодо наведених
прикладів}\label{ux43fux440ux438ux43cux456ux442ux43aux438-ux449ux43eux434ux43e-ux43dux430ux432ux435ux434ux435ux43dux438ux445-ux43fux440ux438ux43aux43bux430ux434ux456ux432}

У всіх наведених вище прикладах функція невідома і залежить від більше
ніж однієї змінної. Змінна і (яку ми диференціюємо) називається залежною
змінною. Змінні, за якими відбувається диференціювання, називаються
незалежними змінними. Наприклад, у рівнянні.

{[} u\_t = u\_\{xx\} {]}

Залежна змінна \(u(x, t)\) є функцією двох незалежних змінних x і t; та
в рівнянні

{[} u\_t = u\_\{rr\} + \textbackslash frac\{1\}\{r\} u\_r +
\textbackslash frac\{1\}\{r\^{}2\} u\_\{00\} {]}

змінна \(u(r, \theta, t)\) залежить від \(r, \theta, t\) .

\section{Чому необхідно вивчати диференціальні рівняння у похідних
похідних?}\label{ux447ux43eux43cux443-ux43dux435ux43eux431ux445ux456ux434ux43dux43e-ux432ux438ux432ux447ux430ux442ux438-ux434ux438ux444ux435ux440ux435ux43dux446ux456ux430ux43bux44cux43dux456-ux440ux456ux432ux43dux44fux43dux43dux44f-ux443-ux43fux43eux445ux456ux434ux43dux438ux445-ux43fux43eux445ux456ux434ux43dux438ux445}

Більшість фізичних законів природи можна сформулювати через рівняння з
похідними похідними. Прикладами є рівняння Максвелла, закон
теплопередачі Ньютона, рівняння Нав'є-Стокса, рівняння руху Ньютона та
рівняння Шредінгера в кватичній механіці. У всіх цих рівняннях фізичні
явища описуються мовою просторових і часових похідних. Похідні
з'являються в рівняннях, оскільки вони описують найважливіші фізичні
величини (такі як швидкість, прискорення, сила, тертя, потік, струм
тощо). Таким чином, виникають диференціальні рівняння в частинних
похідних, які містять невідому функцію, яку потрібно визначити.

Мета цієї книги --- показати читачеві:

\begin{enumerate}
\def\labelenumi{\arabic{enumi}.}
\item
  Як сформулювати фізичну задачу у формі диференціального рівняння у
  похідних похідних (побудова математичної моделі).
\item
  Як розв'язати диференціальне рівняння в похідних (з урахуванням
  початкових і граничних умов).
\end{enumerate}

Перш ніж перейти до побудови математичних моделей, давайте коротко
розглянемо методи розв'язання рівнянь у похідних похідних.

\section{Як розв'язати диференціальні рівняння у похідних
похідних?}\label{ux44fux43a-ux440ux43eux437ux432ux44fux437ux430ux442ux438-ux434ux438ux444ux435ux440ux435ux43dux446ux456ux430ux43bux44cux43dux456-ux440ux456ux432ux43dux44fux43dux43dux44f-ux443-ux43fux43eux445ux456ux434ux43dux438ux445-ux43fux43eux445ux456ux434ux43dux438ux445}

Гарне питання. Виявляється, існує цілий арсенал методів, придатних для
практичного використання. Найважливішими є ті, у яких диференціальні
рівняння з похідними похідними зводяться до звичайних диференціальних
рівнянь. Перелічимо десять методів розв'язання диференціальних рівнянь у
похідних похідних:

\begin{enumerate}
\def\labelenumi{\arabic{enumi}.}
\tightlist
\item
  Метод розділення змінних. Рівняння з факторами, отриманими з n
  незалежними змінними, зводиться до n звичайних-
\end{enumerate}

диференціальнірівняння 1).

\begin{enumerate}
\def\labelenumi{\arabic{enumi}.}
\setcounter{enumi}{1}
\item
  Метод інтегральних перетворень. Рівняння в похідних з n незалежними
  змінними зводиться до рівняння з похідними у (n-1) незалежних змінних;
  Відповідно, рівняння з частковими похідними з двома незалежними
  змінними можна звести до звичайного диференціального рівняння \(^2\)
  ).
\item
  Метод перетворення координат. Початкове диференціальне рівняння в
  похідних похідних зводиться до звичайного рівняння з частковими
  похідними або до іншого, простішого рівняння у похідних похідних
  шляхом відповідного перетворення координат (наприклад, обертання осей
  координат тощо).
\item
  Перетворення залежної змінної. Початкове диференціальне рівняння в
  похідних похідних перетворюється у таке рівняння з похідними похідними
  для іншої невідомої функції, яка
\end{enumerate}

легше розв'язати, ніж початкові 3).

\begin{enumerate}
\def\labelenumi{\arabic{enumi}.}
\setcounter{enumi}{4}
\tightlist
\item
  Чисельні методи. Початкове диференціальне рівняння в похідних похідних
  зводиться до системи різницей, яку розв'язують методом ітерацій на
  комп'ютері. У багатьох випадках це єдиний спосіб розв'язати
  диференціальне рівняння з похідними похідними. Окрім різницьких
  методів для розв'язання рівнянь з частковими про-
\end{enumerate}

\begin{enumerate}
\def\labelenumi{\alph{enumi})}
\setcounter{enumi}{4}
\tightlist
\item
  Ефективніше використовувати загальні перетворення незалежних і
  залежних змінних, включаючи похідні. Ед.
\end{enumerate}

4) Це стосується повного розділення змінних. При частковому розділенні
змінних одне диференціальне рівняння в похідних похідних зводиться до
кількох рівнянь у часткових похідних з меншою кількістю незалежних
змінних. --- Ед. Ед.

2) Це відбувається у випадку одномірного інтегрального перетворення. У
випадку k-вимірного інтегрального перетворення LCP з n незалежними
змінними зводиться до рівняння часткових з n---k незалежними змінними.
--- Ед. Ед.

Існують також інші чисельні методи, зокрема ті, що базуються на
наближенні розв'язків за поліноміальними поверхнями (наближення за
сплайнами).

\begin{enumerate}
\def\labelenumi{\arabic{enumi}.}
\setcounter{enumi}{5}
\tightlist
\item
  Метод теорії збурень. Початкова нелінійна задача зводиться до
  послідовності лінійних задач, апрокенн-
\end{enumerate}

нелінійна задача 1).

\begin{itemize}
\tightlist
\item
  \begin{enumerate}
  \def\labelenumi{\arabic{enumi}.}
  \setcounter{enumi}{6}
  \tightlist
  \item
    Метод функцій Гріна. Початкові та граничні умови замінюються
    системою найпростіших джерел, і задача розв'язується для кожного
    найпростішого джерела. Повне розв'язання початкової задачі
    отримується шляхом підсумовування розв'язків елементарних джерел.
  \end{enumerate}
\item
  \begin{enumerate}
  \def\labelenumi{\arabic{enumi}.}
  \setcounter{enumi}{7}
  \tightlist
  \item
    Метод інтегральних рівнянь. Рівняння з частковими похідними
    зводяться до інтегрального рівняння (рівняння, в якому невідома
    функція знаходиться під знаком інтеграла). Існує багато різних
    методів розв'язання інтегральних рівнянь.
  \end{enumerate}
\item
  \begin{enumerate}
  \def\labelenumi{\arabic{enumi}.}
  \setcounter{enumi}{8}
  \tightlist
  \item
    Варіаційні методи. Замість рівняння з частковими похідними
    диференціальними розв'язують задачу мінімізації. Виявляється, що
    функція, яка дає мінімум певному виразу (наприклад, загальній
    енергії системи), одночасно є розв'язком початкового рівняння
    похідниху похідних 2.
  \end{enumerate}
\item
  \begin{enumerate}
  \def\labelenumi{\arabic{enumi}.}
  \setcounter{enumi}{9}
  \tightlist
  \item
    Метод розкладу за власними функціями. Розв'язок рівняння часткових
    похідних шукається як послідовність власних функцій. Ці власні
    функції знаходяться як розв'язки так званої задачі власних значень,
    що відповідає початковій задачі для рівняння часткових похідних.
  \end{enumerate}
\end{itemize}

\section{Типи диференціальних рівнянь у похідних
похідних}\label{ux442ux438ux43fux438-ux434ux438ux444ux435ux440ux435ux43dux446ux456ux430ux43bux44cux43dux438ux445-ux440ux456ux432ux43dux44fux43dux44c-ux443-ux43fux43eux445ux456ux434ux43dux438ux445-ux43fux43eux445ux456ux434ux43dux438ux445}

Рівняння в похідних похідних можна класифікувати за багатьма
характеристиками. Класифікація рівнянь важлива, оскільки, як виявилося,
кожен клас має свою загальну теорію та методи розв'язання рівнянь.

Тут ми наведемо шість основних методів класифікації

рівняння.

\begin{enumerate}
\def\labelenumi{\arabic{enumi}.}
\tightlist
\item
  Порядок рівняння. Порядок рівняння --- це найвищий порядок часткових
  похідних, включених до рівняння. Наприклад,
\end{enumerate}

Примітка. Ед.

1) Точніше, це метод лінеаризації нелінійних задач. Загальна теорія
збурень містить лінійні та нелінійні теорії збурень, кожна з яких, у
свою чергу, є регулярною та особливою. --- Ед. Ед.

Ці методи застосовні до спеціальних задач, що зумовлено існуванням
відповідних виразів --- функціоналів для цих задач.

Застосовувався раніше

{[} u\_t = u\_\{xx\} {]}

(рівняння другого порядку), \(u_t = u_x\) (рівняння першого порядку),
\(u_1 = uu_{xxx} + \sin x\) (рівняння трьох порядків).

\begin{enumerate}
\def\labelenumi{\arabic{enumi}.}
\setcounter{enumi}{1}
\tightlist
\item
  Кількість змінних. Кількість змінних --- це кількість незалежних
  змінних. Наприклад,
\end{enumerate}

{[} u\_t = u\_\{xx\} {]}

(рівняння з двома змінними \(x\) і \(t\) ),
\(u_t = u_{rr} + \frac{1}{r} u_r + \frac{1}{r^2} u_{\theta\theta}\)
(рівняння з трьома змінними \(r\) ,

\begin{enumerate}
\def\labelenumi{\arabic{enumi}.}
\setcounter{enumi}{2}
\tightlist
\item
  Лінійність. Рівняння з частковими похідними можуть бути лінійними та
  нелінійними. У лінійних рівняннях залежна змінна та всі її часткові
  похідні включаються лінійно, зокрема, вони не множаться одна на одну,
  не множаться в квадраті тощо). Точніше, лінійне рівняння другого
  порядку з двома незалежними змінними є рівнянням вигляду
\end{enumerate}

{[} (1.1) Au\_\{xx\} + Bu\_\{xy\} + Cu\_\{yy\} + Du\_x + Eu\_y + Fu = G
{]}

де A, B, C, D, E, F і G --- це константи або задані функції
пояснювальних змінних x і y.

{[} u\_\{tt\} = e\^{}\{-t\}u\_\{xx\} + \textbackslash sin t {]}

(лінійне рівняння), \(uu_{xx} + u_t = 0\) (нелінійне рівняння),
\(u_{xx} + yu_{yy} = 0\) (лінійне рівняння), \(xu_x + yu_y + u^2 = 0\)
(нелінійне рівняння).

\begin{itemize}
\tightlist
\item
  \begin{enumerate}
  \def\labelenumi{\arabic{enumi}.}
  \setcounter{enumi}{3}
  \tightlist
  \item
    Однорідність. Рівняння (1.1) називають однорідним, якщо права
    частина \(G(x, y)\) ідентична нулю для всіх x і y. Якщо \(G(x, y)\)
    не ідентичне нулю, то рівняння називають неоднорідним.
  \end{enumerate}
\item
  \begin{enumerate}
  \def\labelenumi{\arabic{enumi}.}
  \setcounter{enumi}{4}
  \tightlist
  \item
    Типи шансів. Якщо коефіцієнти A, B, C, D, E і F рівняння (1.1) є
    сталими, то це рівняння називається рівнянням зі сталими
    коефіцієнтами (інакше рівнянням зі змінними коефіцієнтами).
  \end{enumerate}
\item
  \begin{enumerate}
  \def\labelenumi{\arabic{enumi}.}
  \setcounter{enumi}{5}
  \tightlist
  \item
    Три основні типи лінійних рівнянь. Усі лінійні рівняння з частковими
    похідними другого порядку виду (1.1) належать до одного з трьох
    типів: а) параболічних, б) гіперболічних, в) еліптичних.
  \end{enumerate}
\end{itemize}

Параболічний тип. Рівняння параболічного типу описують процеси
теплопровідності та дифузії і визначаються умовою \(B^2 - 4AC = 0\) .

Гіперболічний тип. Рівняння гіперболічного типу описів

осцилюючі системи та хвильові рухи визначаються умовою \(B^2-4AC>0\) .

Еліпічний тип. Рівняння еліптичного типу описують усталені процеси і
визначаються умовним \(B^2-4AC<0\) . Приклади

\begin{enumerate}
\def\labelenumi{\alph{enumi})}
\tightlist
\item
\end{enumerate}

{[} u\_i = u\_\{xx\} {]}

, \(B^2 - 4AC = 0\) (параболічна), (гіперболічна), (гіперболічна),
(гіперболічна \(y > 0\) ), (гіперболічна), (еліптична), (еліптична),
параболічна на \(y = 0\) , гіперболічна на \(y < 0\) .

(У випадку змінних коефіцієнтів тип рівняння може змінюватися від точки
до точки.)

\section{ЗАУВАЖЕННЯ}\label{ux437ux430ux443ux432ux430ux436ux435ux43dux43dux44f}

\begin{itemize}
\tightlist
\item
  \begin{enumerate}
  \def\labelenumi{\arabic{enumi}.}
  \tightlist
  \item
    Загалом, значення \(B^2-4AC\) є функцією незалежних змінних.
    Відповідно, базовий тип рівняння може змінюватися в межах визначення
    рівняння, хоча це не є обов'язковим.
  \end{enumerate}
\item
  \begin{enumerate}
  \def\labelenumi{\arabic{enumi}.}
  \setcounter{enumi}{1}
  \tightlist
  \item
    У рівнянні (1.1) незалежними змінними є x і y. У багатьох задачах
    однією з двох незалежних змінних є час, і рівняння (1.1) можна
    записати через x і t.
  \end{enumerate}
\end{itemize}

\pandocbounded{\includegraphics[keepaspectratio]{images/_page_10_Figure_9.jpeg}}

PHC 1.1. Діаграма класифікації диференціальних рівнянь у похідних
похідних.

\begin{enumerate}
\def\labelenumi{\arabic{enumi}.}
\setcounter{enumi}{2}
\tightlist
\item
  Рис. 1.1. Представлено діаграму класифікації для диференціальних
  рівнянь у похідних похідних.
\end{enumerate}

\section{ЗАВДАННЯ}\label{ux437ux430ux432ux434ux430ux43dux43dux44f}

\begin{itemize}
\tightlist
\item
  \begin{enumerate}
  \def\labelenumi{\arabic{enumi}.}
  \tightlist
  \item
    Класифікуємо наступні рівняння відповідно до всіх характеристик,
    зазначених у днаграмі на рис. 1.1:
  \end{enumerate}
\item
  \begin{enumerate}
  \def\labelenumi{\alph{enumi})}
  \tightlist
  \item
    \(u_t = u_{xx} + 2u_x + u\) ,
  \end{enumerate}
\end{itemize}

\begin{enumerate}
\def\labelenumi{\arabic{enumi})}
\setcounter{enumi}{5}
\tightlist
\item
  \(u_t = u_{xx} + e^{-t}\) , B) \(u_{xx} + 3u_{xy} + u_{yy} = \sin x\)
  , \(u_{tt} = uu_{xxx} + e^{-t}\) .
\end{enumerate}

\begin{itemize}
\tightlist
\item
  \begin{enumerate}
  \def\labelenumi{\arabic{enumi}.}
  \setcounter{enumi}{1}
  \tightlist
  \item
    Скільки розв'язків \(u_1 = u_{xx}\) рівняння існує? Спробуйте знайти
    рішення, як \(u(x, t) = e^{a\lambda + bt}\) .
  \end{enumerate}
\item
  \begin{enumerate}
  \def\labelenumi{\arabic{enumi}.}
  \setcounter{enumi}{2}
  \tightlist
  \item
    Якщо функції \(u_1(x, t)\) і \(u_2(x, t)\) задовольняють рівняння
    (1.1), чи задовольняє сума цих функцій його? Доведи це.
  \end{enumerate}
\item
  \begin{enumerate}
  \def\labelenumi{\arabic{enumi}.}
  \setcounter{enumi}{3}
  \tightlist
  \item
    Ймовірно, найпростіше диференціальне рівняння в похідних ---
  \end{enumerate}
\end{itemize}

{[} \textbackslash frac\{\textbackslash partial u(x,
y)\}\{\textbackslash partial x\} = 0. {]}

Чи можете ви розв'язати це рівняння? (Знайдіть усі функції, що
задовольняють це рівняння.)

\begin{enumerate}
\def\labelenumi{\arabic{enumi}.}
\setcounter{enumi}{4}
\tightlist
\item
  Що можна сказати про рівняння
\end{enumerate}

{[} \textbackslash frac\{\textbackslash partial\^{}2 u(x,
y)\}\{\textbackslash partial x \textbackslash partial y\} = 0 {]}

?

Чи можете ви знайти всі розв'язки цього рівняння? Скільки їх? Порівняйте
з кількістю розв'язків звичайного диференціального рівняння

{[} \textbackslash frac\{d\textsuperscript{2y\}\{dx}2\} = 0. {]}

\section{ЗАДАЧІ ТИПУ ДИФУЗІЇ (ПАРАБОЛІЧНІ
РІВНЯННЯ)}\label{ux437ux430ux434ux430ux447ux456-ux442ux438ux43fux443-ux434ux438ux444ux443ux437ux456ux457-ux43fux430ux440ux430ux431ux43eux43bux456ux447ux43dux456-ux440ux456ux432ux43dux44fux43dux43dux44f}

МЕТА ЛЕКЦІЇ: Показати, як параболічні рівняння використовуються для
розв'язання задач теплопровідності та дифузії. Наведено кілька прикладів
параболічних рівнянь, а також пояснюється фізичне значення різних членів
\((u_1, u_2, u_3, u_4, u_6)\) , включених у ці рівняння.

На конкретному прикладі розглядається основна ідея зміщеної задачі для
параболічного рівняння. Одна з головних цілей цієї лекції --- надати
читачеві інтуїтивне уявлення

Розуміння параболічних проблем.

Почнемо з розгляньмо просту фізичну задачу і те, як її можна описати
математичною моделлю з рівнянням у похідних похідних.

Потім ми ускладнимо задачу і покажемо, як нові фізичні умови призводять
до появи нових рівнянь у похідних похідних. У цій лекції диференціальні
рівняння не виводяться і не розв'язуються. Про це ми поговоримо пізніше.

\section{Простий експеримент з
теплопровідністю}\label{ux43fux440ux43eux441ux442ux438ux439-ux435ux43aux441ux43fux435ux440ux438ux43cux435ux43dux442-ux437-ux442ux435ux43fux43bux43eux43fux440ux43eux432ux456ux434ux43dux456ux441ux442ux44e}

Припустимо, що ми провели простий експеримент, який розбили на такі
кроки:

КРОК 1. Візьміть достатньо довгий (скажімо, \(L=2\,\mathrm{m}\) )
стрижень (наприклад, мідний) діаметром \(2\,\mathrm{cm}\) , бічна
поверхня (але не кінці) якої термічно ізольована. Замість стрижня можна
використовувати мідну трубку, покриту теплоізоляцією зовні та зсередини.

Іншими словами, тепло може проходити лише через кінці стрижня і не може
проходити через його бічну поверхню.

КРОК 2. Тепер помістимо цей стрижень у пристрій із фіксованою
температурою \(T_{\mathfrak{o}}\) (градуси Цельсія) на достатньо довгий
час, щоб температура всередині стрижня стала

Так само, як і температура пристрою. Для простоти припустимо, що
температура пристрою \(T_0 = 10^{\circ}\) C.

КРОК 3. Давайте витягнемо стрижень з нашого пристрою у певний момент
часу, що зручно взяти як нуль t = 0, і приєднаємо до нього дві термопари
з кінців. Завдання цих елементів --- підтримувати фіксовані температури
\(T_1\) і \(T_2\) на кінцях (наприклад, \(T_1=0^{\circ}\) C і
\(T_2=50^{\circ}\) C). Іншими словами, два термостати постійно
контролюють температуру на кінцях стрижня, і якщо температура
відхиляється від встановлених значень \(T_1\) і \(T_2\) , активуються
потужні нагрівальні або охолоджувальні елементи (пристрої) для
відповідного регулювання температури. Наш експеримент ілюструє RPS, 2 1,

\pandocbounded{\includegraphics[keepaspectratio]{images/_page_13_Figure_4.jpeg}}

Рис. 2.1. Схема експерименту: a --- це термопара, яка підтримує
встановлену температуру \(T_1\) на лівому кінці стрижня; \(\delta\) ---
це термопара, яка підтримує встановлену температуру \(T_2\) на правому
кінці стрижня.

КРОК 4. Тепер ми будемо слідкувати за температурним профілем у стрижні
за допомогою пристрою візуалізації. Чому ми хочемо провести такий
експеримент? Це інше питання, про яке ми поговоримо пізніше.

Це кінець опису нашого експерименту.

Головна мета цієї лекції --- показати, як цю фізичну проблему та її
модифікації можна змоделювати мовою параболічних рівнянь.

\section{Математична модель
теплопровідності}\label{ux43cux430ux442ux435ux43cux430ux442ux438ux447ux43dux430-ux43cux43eux434ux435ux43bux44c-ux442ux435ux43fux43bux43eux43fux440ux43eux432ux456ux434ux43dux43eux441ux442ux456}

Опис нашої фізичної проблеми вимагає трьох типів відносин.

\begin{itemize}
\tightlist
\item
  \begin{enumerate}
  \def\labelenumi{\arabic{enumi}.}
  \tightlist
  \item
    Рівняння з частковими похідними, які описують фізичне явище
    теплопровідності.
  \end{enumerate}
\item
  \begin{enumerate}
  \def\labelenumi{\arabic{enumi}.}
  \setcounter{enumi}{1}
  \tightlist
  \item
    Граничні умови, що описують процеси теплопередачі на межах.
  \end{enumerate}
\item
  \begin{enumerate}
  \def\labelenumi{\arabic{enumi}.}
  \setcounter{enumi}{2}
  \tightlist
  \item
    Початкові твердження, які описують стан системи на початку процесу.
  \end{enumerate}
\end{itemize}

\section{Рівняння
теплопровідності}\label{ux440ux456ux432ux43dux44fux43dux43dux44f-ux442ux435ux43fux43bux43eux43fux440ux43eux432ux456ux434ux43dux43eux441ux442ux456}

Базове одновимірне рівняння теплопровідності записується у вилці

(2.1)

{[} u\_t = \textbackslash alpha\^{}2 u\_\{xx\}, \textbackslash quad 0
\textless{} x \textless{} L, \textbackslash quad 0 \textless{} t
\textless{} \textbackslash infty. {]}

Це рівняння пов'язує значення \(u_t\) --- швидкість зміни температури з
часом (виміряна в градусах/с)

та величину \(u_{xx}\) --- це увігнутий профілю температури \(u(x, t)\)
(який є мірою різниці між температурою в заданій точці та температурою в
сусідніх точках).

\pandocbounded{\includegraphics[keepaspectratio]{images/_page_14_Figure_14.jpeg}}

Рис. 2.2. Стрілки показують зміну температури відповідно до рівняння
\(u_t = \alpha^2 u_{xx}\) .

Це рівняння буде виведено з базового рівняння збереження тепла в
наступних лекціях, але наразі ми розглянемо його заданим. Це рівняння
каже, що температура \(u\left(x,\ t\right)\) (у певній точці t у певній
точці x стрижня) зростає \((u_t>0)\) або зменшується \((u_t<0)\) залежно
від того, чи є друга похідна \(u_{xx}\) позитивною чи негативною. Фіг.
Рисунок 2.2 ілюструє зміну температури в різних точках на елементі.

Давайте подивимось, як значення \(u_{xx}\) можна інтерпретувати мовою
теплопровідності. Припустимо, ми наблизимо значення \(u_{xx}\) фінальних
постів

{[} u\_\{xx\}(x, t) \textbackslash cong
\textbackslash frac\{1\}\{\textbackslash Delta x\^{}2\} {[}u(x +
\textbackslash Delta x, t) - 2u(x, t) + u(x - \textbackslash Delta x,
t){]}. {]}

Це співвідношення можна переписати у вигляді

{[} u\_\{\textbackslash lambda x\}(x, t) \textbackslash simeq
\textbackslash frac\{2\}\{\textbackslash Delta x\^{}2\}
\textbackslash left{[} \textbackslash frac\{u(x + \textbackslash Delta
x, t) + u(x - \textbackslash Delta x, t)\}\{2\} - u(x, t)
\textbackslash right{]}. {]}

Тепер ми можемо дати наступну інтерпретацію значення \(u_{xx}\) .

\begin{enumerate}
\def\labelenumi{\arabic{enumi}.}
\item
  Якщо температура \(u(x, t)\) менша за середнє значення температури у
  двох сусідніх точках, то \(u_{xx} > 0\) (тут загальний тепловий потік
  по осі x додатний).
\item
  Якщо температура \(u(x, t)\) дорівнює середньому значенню температур у
  двох сусідніх точках, то \(u_{xx} = 0\) (тут загальний тепловий потік
  по осі x дорівнює нулю).
\item
  Якщо температура \(u\left(x,\ t\right)\) вища за середню температуру у
  двох сусідніх точках, то \(u_{xx}<0\) (тут --- повний потік
\end{enumerate}

тепла вздовж осі x є від'ємною).

Усе це ілюструється на рис. 2.2. Іншими словами, якщо температура в
точці x більша за середню температуру у двох сусідніх точках
\(x - \Delta x\) і \(x + \Delta x\) , то температура в точці x
зменшиться. Відповідно, точна швидкість зниження температури, а саме
величина \(u_t\) , пропорційна цій різниці. Коефіцієнт пропорційності
\(\alpha^2\) визначається властивостями матеріалу. Ми розглянемо цю
константу детальніше у майбутніх лекціях.

\section{Граничні умови
(GI)}\label{ux433ux440ux430ux43dux438ux447ux43dux456-ux443ux43cux43eux432ux438-gi}

Усі фізичні проблеми характеризуються наявністю певних меж. Отже, щоб
рухатися вперед, ми повинні включити межі в математичну модель, щоб
адекватно описати фізичну проблему. У нашому експерименті граничні умови
(ГІ) дуже прості. Оскільки температура на кінцях x=0 і x=L завжди
залишається сталою і дорівнює \(T_1\) та \(T_2\) , можливо записати

(2.2)

\begin{verbatim}
\left\{ \begin{aligned} u\left(0,\ t\right) = T_{t}, \\ u\left(L,\ t\right) = T_{2}, \end{aligned} \right. \quad 0 < t < \infty,
\end{verbatim}

\section{Початкові умови
(NU)}\label{ux43fux43eux447ux430ux442ux43aux43eux432ux456-ux443ux43cux43eux432ux438-nu}

Усі фізичні процеси мають починатися у певний момент часу (зазвичай
приймається як нуль t=0), тому ми повинні встановити фізичні умови саме
в цю точку. Відтоді, як ми почали контролювати температуру з моменту,
коли паличка мала постійну температуру \(T_{\mathrm{ }o}\) , ми

(2.3)

{[} u(x, 0) = T\_0, \textbackslash quad 0 \textbackslash le x
\textbackslash le L {]}

Тепер ми створили математичну модель експерименту. Система співвідношень
(2.1), (2.2) та (2.3) називається задачею зміщення для рівняння
теплопровідності і зазвичай записується як

(2.4)

{[} \textbackslash begin\{aligned\} \& (\textbackslash text\{YUII\})
\textbackslash quad u\_t = \textbackslash alpha\^{}2 u\_\{xx\}, \& 0
\textless{} x \textless{} L, ~0 \textless{} t \textless{}
\textbackslash infty, \textbackslash{} \& \textbackslash left(
\textbackslash Gamma \textbackslash text\{Y\} \textbackslash right)
\textbackslash quad \textbackslash begin\{cases\} u ~(0, ~t) = T\_1,
\textbackslash{} u ~(L, ~t) = T\_2, \textbackslash end\{cases\} \& 0
\textless{} t \textless{} \textbackslash infty, \textbackslash{} \&
(\textbackslash text\{HY\}) \textbackslash quad u ~(x, ~0) = T\_0,
\textbackslash quad 0 \textbackslash leqslant x \textbackslash leqslant
L. \textbackslash end\{aligned\} {]}

Цікаво, що це зовсім не очевидно: чи існує одна функція \(u(x, t)\), яка
задовольняє умови (2.4), і чи описує ця функція температуру в стрижні?
Отже, наше безпосереднє завдання --- знайти цю функцію \(u(x, t)\) з
(2.4).

Перш ніж завершити цю лекцію, давайте розглянемо деякі варіації основної
проблеми. Почнемо з невеликих змін рівняння теплопровідності.

\section{Деякі рівняння типу
дифузії}\label{ux434ux435ux44fux43aux456-ux440ux456ux432ux43dux44fux43dux43dux44f-ux442ux438ux43fux443-ux434ux438ux444ux443ux437ux456ux457}

Теплопередача через бічну поверхню пропорційна різниці температур

Рівняння

{[} u\_t = \textbackslash alpha\^{}2 u\_\{xx\} - \textbackslash beta (u
- u\_0) {]}

Описує теплопровідність елемента, враховуючи не лише дифузію теплових
\(\alpha^2 u_{xx}\) вздовж елемента, а й теплопередачу через бічну
поверхню елемента. Відтік тепла ( \(\beta > 0\) ) або його приплив (
\(\beta < 0\) ) пропорційний різниці між температурою u стрижня та
температурою навколишнього середовища \(u_0\) ( \(\beta\) --- це сталий
коефіцієнт пропорційності). Якщо коефіцієнт \(\beta\) великий порівняно
з \(\alpha^2\) , то теплопотік уздовж стрижня буде малим порівняно з

Протікання через бічну поверхню, і тому тепло буде проходити через неї
відповідно до закону \(u_t = -\beta (u - u_a)\) .

У хімії, де величина відіграє роль концентрації, рівняння

{[} u\_1 = \textbackslash alpha\^{}2 u\_\{xx\} - \textbackslash beta (u
- u\_0) {]}

означає, що швидкість зміни \(u_i\) кількості речовини залежить як від
дифузії \(\alpha^2 u_{xx}\) (у напрямку осі x), так і від появи (
\(\beta < 0\) ) або розпаду ( \(\beta > 0\) ) речовини в хімічній
реакції, і пропорційна різниці між двома концентраціями u і
\(u_{\theta}\) .

\section{Внутрішнє джерело
тепла}\label{ux432ux43dux443ux442ux440ux456ux448ux43dux454-ux434ux436ux435ux440ux435ux43bux43e-ux442ux435ux43fux43bux430}

Неоднородне рівняння

{[} u\_t = \textbackslash alpha\^{}2 u\_\{xx\} + f(x, t) {]}

відповідає випадку, коли всередині елемента є джерело тепла (розташоване
вздовж усього елемента і діє в будь-який момент t). Наприклад, може бути
так, що дріт з електричним струмом проходить через стрижень і опір є
постійним джерелом тепла \(f(x,t)\) = K.

\section{Рівняння конвективної
дифузії}\label{ux440ux456ux432ux43dux44fux43dux43dux44f-ux43aux43eux43dux432ux435ux43aux442ux438ux432ux43dux43eux457-ux434ux438ux444ux443ux437ux456ux457}

Припустимо, що домішка поширюється вздовж потоку, що рухається зі
швидкістю v. Очевидно, що концентрація \(u\left(x,\,t\right)\) домішок
змінюється залежно від x (позитивний напрямок осі x обирається вздовж
потоку) і часу t. Швидкість зміни концентрації \(u_t\) визначається
рівнянням конвективної дифузії

{[} u\_t = \textbackslash alpha\^{}2 u\_\{xx\} - v u\_x. {]}

Термін \(\alpha^2 u_{xx}\) описує внесок дифузії, а \(vu_x\) ---
конвективний компонент. Що важливіше, дифузія або конвекція залежить від
співвідношення \(\alpha^2\) до v. Димові частинки піднімаються
конвективно через рух гарячого повітря і одночасно розсіюють через
вихорові рухи повітря. Окрім цих модифікацій рівняння теплопровідності,
можливо змінювати граничні умови на кінцях стрижня відповідно до
фізичної ситуації. Деякі з цих змін ми обговоримо у третій лекції.

\section{ЗАУВАЖЕННЯ}\label{ux437ux430ux443ux432ux430ux436ux435ux43dux43dux44f-1}

Рівняння теплопровідності \(u_t = \alpha^2(x) u_{xx}\) з змінним
коефіцієнтом \(\alpha(x)\) відповідає задачі, в якій дифузія залежить
від x (тобто матеріал є гетерогенним). Наприклад, мідь і сталь---

\pandocbounded{\includegraphics[keepaspectratio]{images/_page_18_Picture_4.jpeg}}

PHC. 2.3.

Пластини контактують, як показано на рис. 2.3. Припустимо, що
температура \(u(0,t)=0^{\circ}\) C підтримується на лівому боці міді, а
температура \(u(L,t)=20^{\circ}\) C --- з правого боку сталі, тоді
рівняння, що описує теплопровідність у цій системі, виглядає так:

{[} u\_t = \textbackslash alpha\^{}2(x) , u\_\{\_\{! A , X\}\},
\textbackslash quad 0 \textless{} x \textless{} L, {]}

де

{[} \textbackslash alpha(x) = \textbackslash begin\{cases\}
\textbackslash alpha\_1 \& \textbackslash text\{(коэффициент диффузии в
меди)\}, \textbackslash quad 0 \textless{} x \textless{} L/2,
\textbackslash{} \textbackslash alpha\_2 \&
\textbackslash text\{(коэффициент диффузии в стали)\},
\textbackslash quad L/2 \textless{} x \textless{} L.
\textbackslash end\{cases\} {]}

\#\#\# ЗАВДАННЯ

\begin{enumerate}
\def\labelenumi{\arabic{enumi}.}
\tightlist
\item
  Якщо початкова температура в стрижні
\end{enumerate}

{[} u(x, 0) = \textbackslash sin \textbackslash pi x,
\textbackslash quad 0 \textbackslash leqslant x \textbackslash leqslant
1, {]}

а граничні умови такі:

{[} u(0, t) = 0, {]}

\(u(1, t) = 0,\) Тоді як поводитиметься температура в стрижні
\(U(x, t)\) при T \textgreater{} 0? 2. Припустимо, що в стрижні є
постійне внутрішнє джерело тепла, так що рівняння теплопровідності в
стрижні записується як

{[} u\_t = \textbackslash alpha\^{}2 u\_\{xx\} + 1, \textbackslash quad
0 \textless{} x \textless{} 1. {]}

Припустимо, що ми зафіксували температуру на межі з значеннями
\(u(0, t)\) = 0 і \(u(1, t)\) = 1. Якою буде стаціонарна температура у
стрижні? Іншими словами, який c-незалежний розподіл \(U(x, t)\)
перешкоджує \(u(x, t)\)?

\begin{enumerate}
\def\labelenumi{\arabic{enumi}.}
\setcounter{enumi}{2}
\tightlist
\item
  Припустимо, що металевий стрижень втрачає тепло через бічну поверхню
  згідно з рівнянням
\end{enumerate}

{[} u\_1 = \textbackslash alpha\^{}2 u\_\{xx\} - \textbackslash beta u
{]}

, \(0 < x < 1\) .

І припустимо, що ми підтримуємо кінці стрижня при температурах
\(U(0, T)\) = 1, \(U(1, T)\) = 1. Знайдіть стаціонарний розподіл
температури по стрижню. Побудуйте графік цього розподілу. У яких
моментах відбувається тепловий потік?

\begin{enumerate}
\def\labelenumi{\arabic{enumi}.}
\setcounter{enumi}{3}
\tightlist
\item
  Припустимо, що термоізольований металевий стрижень довжини L = 1 має
  початкову температуру \(\sin(3\pi x)\) , а фіксовану температуру --- 0
  і \(10\,^{\circ}\) C на лівому та правому кінцях.
\end{enumerate}

\section{Лекція 3}\label{ux43bux435ux43aux446ux456ux44f-3}

\section{ГРАНИЧНІ УМОВИ В ЗАДАЧАХ ДИФУЗІЙНОГО
ТИПУ}\label{ux433ux440ux430ux43dux438ux447ux43dux456-ux443ux43cux43eux432ux438-ux432-ux437ux430ux434ux430ux447ux430ux445-ux434ux438ux444ux443ux437ux456ux439ux43dux43eux433ux43e-ux442ux438ux43fux443}

МЕТА ЛЕКЦІЇ: Показати, як виникають різні граничні умови у задачах
теплопровідності та дифузії, а також ввести важливе поняття теплового
потоку.

Розглядаються три найважливіші типи граничних умов:

\begin{itemize}
\tightlist
\item
  \begin{enumerate}
  \def\labelenumi{\arabic{enumi}.}
  \tightlist
  \item
    u = \(g(t)\) (температура задається на межі),
  \end{enumerate}
\item
  \texttt{2.\ \textbackslash{}\ \textbackslash{}frac\{\textbackslash{}partial\ u\}\{\textbackslash{}partial\ n\}\ +\ \textbackslash{}lambda\ u\ =\ g\textbackslash{}left(t\textbackslash{}right)}
  (встановлення температури навколишнього середовища);
\end{itemize}

n --- вектор зовнішньої нормалі до межі),
\texttt{3.\ \textbackslash{}frac\{\textbackslash{}partial\ u\}\{\textbackslash{}partial\ n\}\ =\ g(t)}
(визначений тепловий потік через межу).

У задачах теплопровідності зазвичай використовують три основні типи
граничних умов. Ця лекція наводить приклади того, як різні фізичні умови
експерименту призводять до різних типів граничних умов.

\section{Граничні умови першого виду (температура задається на
межі)}\label{ux433ux440ux430ux43dux438ux447ux43dux456-ux443ux43cux43eux432ux438-ux43fux435ux440ux448ux43eux433ux43e-ux432ux438ux434ux443-ux442ux435ux43cux43fux435ux440ux430ux442ux443ux440ux430-ux437ux430ux434ux430ux454ux442ux44cux441ux44f-ux43dux430-ux43cux435ux436ux456}

Розглянемо тепловий потік в одномірному стрижні, показаний на рис. 3 1,
і припустимо, що температура кінців змінюється відповідно до законів
\(g_1(t)\) і \(g_2(t)\) .

\pandocbounded{\includegraphics[keepaspectratio]{images/_page_20_Picture_4.jpeg}}

Рис. 31. Температура встановлюється на межі.

Як згадувалося в попередній лекції, кожен пристрій, що підтримує
встановлену температуру на кінцях, складається з термостата та
нагрівального елемента для підтримки відповідних теплових умов. Проблеми
з граничними умовами першого виду дуже поширені. У деяких випадках
складність полягає у пошуку контролерів граничних умов, тобто температур
\(g_1(t)\) та \(g_2(t)\) меж, які призводять до певної зміни температури
всередині стрижня. Наприклад, у металургії часто необхідно обирати
граничні умови контролю, щоб температура металу всередині печі
змінювалася з часом, а температурний градієнт був невеликим.

\pandocbounded{\includegraphics[keepaspectratio]{images/_page_20_Picture_7.jpeg}}

Рис. 3.2. Коливання температури на межі.

Граничні умови першого виду також виникають у багатовимірних задачах
теплопровідності. Як приклад можна навести цікаву задачу визначення
температурного поля всередині кругового диска з радіусом \emph{R} (див.
рис. 3.2), коли гранична температура задається у полярних координатах
формулою \(u(R, \theta, t) = \cos t \sin \theta\) .

Звісно, щоб розв'язати задачу, потрібно знати початкову температуру, але
в цьому випадку початкові умови перестають впливати на розв'язок через
дуже короткий час, і отримана температура всередині диска визначається
лише значеннями температури на межі.

\section{Граничні умови другого типу (з урахуванням температури
навколишнього
середовища)}\label{ux433ux440ux430ux43dux438ux447ux43dux456-ux443ux43cux43eux432ux438-ux434ux440ux443ux433ux43eux433ux43e-ux442ux438ux43fux443-ux437-ux443ux440ux430ux445ux443ux432ux430ux43dux43dux44fux43c-ux442ux435ux43cux43fux435ux440ux430ux442ux443ux440ux438-ux43dux430ux432ux43aux43eux43bux438ux448ux43dux44cux43eux433ux43e-ux441ux435ux440ux435ux434ux43eux432ux438ux449ux430}

Припустимо, що ми знову розглядаємо термічно ізольований мідний
стержень, але тепер відмовляємося від заданого температурного режиму на
кінцях стрижня. Натомість давайте наблизимо кінці стрижня до контакту з
двома середовищами. Нехай температура одного з них \(g_1(t)\) , а іншого
\(g_2(t)\) . Іншими словами, припустимо, що лівий кінець стрижня
поміщений у контейнер із рідиною, температура якої змінюється відповідно
до закону \(g_1(t)\) . А правий кінець поміщається в іншу рідину з
температурою \(g_2(t)\) (рис. 3, 3)

\pandocbounded{\includegraphics[keepaspectratio]{images/_page_21_Picture_6.jpeg}}

Рис. 3.3. Конвективна теплопередача через межі, a --- це рідина при
температурі \(g_1(t)\) , \(\delta\) --- рідина при температурі
\(g_2(t)\) .

Задаючи граничні умови такого типу, ми не можемо вважати, що граничні
температури стрижня збігаються з температурами рідини. Але ми знаємо
закон Ньютона: якщо температура одного з кінців стрижня менша за
температуру відповідної рідини, тепло потече в стрижень зі швидкістю,
пропорційною різницям температур. Іншими словами, для одномірного
стрижня з межами x=0 і x=1 закон теплопередачі Ньютона формулюється так:

(3.1)

\begin{verbatim}
\{
\end{verbatim}

вихідний тепловий потік на \(x=0\) дорівнює
\(h\left[u\left(0,\,t\right)-g_{1}\left(t\right)\right]\) , вихідний
тепловий потік при \(x=L\) дорівнює
\(h\left[u\left(L,\,t\right)-g_{2}\left(t\right)\right]\) ,

де h --- коефіцієнт теплопередачі, який показує, скільки калорій
проходить через межу за одну секунду при різниці температур у один
градус. Вихідний потік дорівнює кількості кало-

Пройшов крізь палочний спис за одну секунду. Зверніть увагу, що вихідний
тепловий потік буде позитивним на кінці стрижня, де температура стрижня
вища за температуру навколишнього середовища. Рівняння (3.1) разом із
відомим законом теплопровідності Фур'є тепер можна використовувати для
отримання Р и с

\pandocbounded{\includegraphics[keepaspectratio]{images/_page_22_Figure_3.jpeg}}

. 3.4. Ілюстрація закону Фур'є: \(\frac{\partial u}{\partial n} > 0\)
означає, що тепло надходить у цю область;
\(\frac{\partial u}{\partial n} < 0\) означає, що тепло виходить із цієї
зони; n --- напрямок зовнішньої нормалі;
\(\frac{\partial u}{\partial n} -\) Зміна температури в напрямку N;
\(\frac{\partial u}{\partial n} = -\frac{\partial u}{\partial (-n)}\) є
похідною у напрямку внутрішньої нормальної --- це похідна, взята
\(\mathbf{e}\) протилежному знаку протилежного знака пормалі.

граничні умови. Закон Фур'є дає друге представлення (перше міститься у
(3.1)) для вихідного теплового потоку. Прирівнюючи ці два вирази до
потоку, ми отримуємо граничні умови, які шукаємо. Сформулюємо закон
Фур'є (встановлений експериментально):

(3.2) Тепловий потік, що проходить через межу площі, пропорційний
нормальній похідній температури у напрямку внутрішньої норми.

Цей закон стверджує, що якщо температура швидко підвищується у напрямку
зовнішньої температури до межі D (рисунок 3.4), то потік буде текти з
навколишнього середовища до області D.

У нашій одновимірній задачі закон Фур'є має вигляд:

(3.3)

{[} \textbackslash begin\{cases\} \textbackslash text\{вытекающий поток
тепла при \} x=0 \textbackslash text\{ равен \} k
\textbackslash frac\{\textbackslash partial u\}\{\textbackslash partial
x\}, \textbackslash{} \textbackslash text\{вытекающий поток тепла при \}
x=' \textbackslash text\{ равен \} -k
\textbackslash frac\{\textbackslash partial u\}\{\textbackslash partial
x\}, \textbackslash end\{cases\} {]}

де k --- теплопровідність матеріалу, яка є мірою того, наскільки добре
матеріал проводить тепло. Погано провідні матеріали мають коефіцієнт
близький до нуля в одиницях GHS, тоді як мідь та алюміній мають значення
значно вище нуля.

\pandocbounded{\includegraphics[keepaspectratio]{images/_page_23_Figure_5.jpeg}}

Рис. 3.5. Ще одна ілюстрація закону Фур'є.

Закон Фур'є (3.3) фактично описує теплопровідність всюди всередині
стрижня, а не лише на його межі; Наприклад, (див. рис. 3.5).

(3.4) Потік, що проходить через точку

x

зліва направо \(=-k\frac{\partial n}{\partial x}\) Закон Фур'є (3.4)
стверджує: якщо \(u_x < 0\) , то в точці \(x_0\) тепло тече зліва
направо, якщо \(u_x > 0\) , то в точці \(x_0\) тепло тече справа наліво.
(Тепло завжди переходить від вищих температур до нижчих.)

Нарешті, якщо використовувати вирази (3.1) і (3.3) для теплового потоку,
то отримуємо бажані граничні умови для задачі, описаної на рис. 3.3, у
чисто математичній формі

{[} (\textbackslash Gamma Y) \textbackslash quad
\textbackslash begin\{cases\}
\textbackslash frac\{\textbackslash partial
u\textbackslash left(0,,t\textbackslash right)\}\{\textbackslash partial
x\} =
\textbackslash frac\{h\}\{k\}\textbackslash left{[}u\textbackslash left(0,,t\textbackslash right)
-
g\_\{1\}\textbackslash left(t\textbackslash right)\textbackslash right{]},
\textbackslash{} \textbackslash frac\{\textbackslash partial
u\textbackslash left(L,,t\textbackslash right)\}\{\textbackslash partial
x\} =
-\textbackslash frac\{h\}\{k\}\textbackslash left{[}u\textbackslash left(L,,t\textbackslash right)
-
g\_\{2\}\textbackslash left(t\textbackslash right)\textbackslash right{]}.
\textbackslash end\{cases\} \textbackslash quad 0 \textless{} t
\textless{} \textbackslash infty, {]}

Дуже часто константа h/k позначається як \(\lambda\) , а граничні умови
теплового потоку на межі записуються у вигляді

(3.5)

{[} u\_\{x\}(0, t) = \textbackslash lambda {[}u(0, t) - g\_\{1\}(t){]},
u\_\{x\}(L, t) = -\textbackslash lambda {[}u(L, t) - g\_\{2\}(t){]}. {]}

Для випадку вищих розмірностей отримують подібні граничні умови.
Наприклад, якщо межа кругового диска осяяна рухомою рідиною з
температурою \(g(\theta,t)\) , то граничні умови записуються як

{[} \textbackslash frac\{\textbackslash partial
u\}\{\textbackslash partial r\}(R, \textbackslash theta, t) =
-\textbackslash frac\{h\}\{k\} {[}u(R, \textbackslash theta, t) -
g(\textbackslash theta, t){]}. {]}

Тут \(\frac{\partial u}{\partial r}(R, \theta, t)\) позначає зовнішню
нормальну похідну (у позитивному напрямку r-осі) значення u, обчислену в
точці \((R, \theta)\) межі. Ми назвемо такі граничні умови лінійними
(оскільки вони лінійні в U і \(u_r\) ), але неоднорідними, оскільки
функція \(g(\theta, t)\) знаходиться праворуч.

\section{Граничні умови третього виду (враховується потік; зокрема,
враховується випадок термічно ізольованих
меж)}\label{ux433ux440ux430ux43dux438ux447ux43dux456-ux443ux43cux43eux432ux438-ux442ux440ux435ux442ux44cux43eux433ux43e-ux432ux438ux434ux443-ux432ux440ux430ux445ux43eux432ux443ux454ux442ux44cux441ux44f-ux43fux43eux442ux456ux43a-ux437ux43eux43aux440ux435ux43cux430-ux432ux440ux430ux445ux43eux432ux443ux454ux442ux44cux441ux44f-ux432ux438ux43fux430ux434ux43eux43a-ux442ux435ux440ux43cux456ux447ux43dux43e-ux456ux437ux43eux43bux44cux43eux432ux430ux43dux438ux445-ux43cux435ux436}

Жоден потік не проходить через термоізольовані межі, тому нормальна
похідна (зовнішня чи внутрішня) має зрівнятися з нулем на межі, оскільки
потік пропорційний нормальній похідній. У випадку одномірного елемента з
ізольованими кінцями x=0 і x=L, граничні умови мають вигляд

{[} u\_\{\textbackslash lambda\}(0, t) = 0, {]}

\(u_{\lambda}(L, t) = 0.\) \(0 < t < \infty,\) У двовимірній області
ізоляція межі означає, що нормальна похідна температури на межі
перетворюється на нуль. Припустимо, наприклад, що круглий диск термічно
ізольований уздовж межі. Тоді граничну умову слід записати так:
\(u_r(R, \theta, t) = 0\) для всіх \(0 \le \theta < 2\pi\) і всіх
\(0 < t < \infty\) .

З іншого боку, якщо вказати кількість тепла, що проходить через межу
нашого диска, то граничні умови набувають вигляду

{[} u\_r(R, \textbackslash theta, t) = f(\textbackslash theta, t), {]}

де \(f(0, t)\) --- це кількість тепла, що переходить у круглий диск із
зовнішнього джерела.

Тепер проілюструємо різні типи граничних умов.

\section{Типові граничні умови для одномірної задачі
теплопровідності}\label{ux442ux438ux43fux43eux432ux456-ux433ux440ux430ux43dux438ux447ux43dux456-ux443ux43cux43eux432ux438-ux434ux43bux44f-ux43eux434ux43dux43eux43cux456ux440ux43dux43eux457-ux437ux430ux434ux430ux447ux456-ux442ux435ux43fux43bux43eux43fux440ux43eux432ux456ux434ux43dux43eux441ux442ux456}

Припустимо, що у нас є мідний стрижень довжиною 200 см, бікова поверхня
якого термічно ізольована, а початкова температура становить 0 °C.
Припустимо також, що верхній кінець стрижня x=0 термічно ізольований, а
нижній кінець x=200 змивається рухомою водою, яка має постійну
температуру \(g_2(t)=20\) °C (див. рис. 3.6).

\pandocbounded{\includegraphics[keepaspectratio]{images/_page_25_Picture_6.jpeg}}

Рис. 3.6. Змішане завдання.

Математична модель цієї задачі описується наступними чотирма
відношеннями:

{[} (3.6) \textbackslash qquad \textbackslash begin\{aligned\}
(\textbackslash mathrm\{Y\}\textbackslash mathrm\{H\}\textbackslash Pi)
\& u\_t = \textbackslash alpha\^{}2 u\_\{xx\}, \& 0 \textless{} x
\textless{} 200, \& 0 \textless{} t \textless{} \textbackslash infty,
\textbackslash{} u\_x (0, , t) = 0, \& \textbackslash{} u\_x (200, , t)
= -\textbackslash frac\{h\}\{k\} \textbackslash left{[} u , (200, , t) -
20 \textbackslash right{]}, \& 0 \textless{} t \textless{}
\textbackslash infty, \textbackslash{} (\textbackslash mathrm\{HY\}) \&
u , (x, , 0) = 0 ,
\^{}\{\textbackslash mathrm\{o\}\}\textbackslash mathrm\{C\}, \& 0
\textbackslash leqslant x \textbackslash leqslant 200, \&
\textbackslash end\{aligned\} {]}

де \(\alpha^2 = 1,16\) см 2/с --- теплова дифузія міді;

k = 0,93 кал/см·с·°C --- теплопровідність міді;

h --- коефіцієнт теплопередачі. Визначення цього співвідношення ---
велика проблема. Він вимірюється швидкістю виходу тепла від другого
кінця стрижня до води. Це залежить від швидкості потоку води, стану
поверхні, яку слід змивати тощо.

Читач може спробувати провести експеримент, щоб визначити

ділення значення цього коефіцієнта.

\section{ЗАУВАЖЕННЯ}\label{ux437ux430ux443ux432ux430ux436ux435ux43dux43dux44f-2}

\begin{enumerate}
\def\labelenumi{\arabic{enumi}.}
\tightlist
\item
  Типова задача внутрішньої теплопровідності для квадрата показана на
  рис. 3.7.
\end{enumerate}

У цій задачі початкова температура \(u\left(x,\,y,\,t\right)\) при t = 0
задається всередині квадрата, і рівняння з похідними з частковими має
вигляд Р и с

\pandocbounded{\includegraphics[keepaspectratio]{images/_page_26_Figure_9.jpeg}}

. 3.7. Типові граничні умови для задачі дифузії в квадраті: a ---
кількість тепла, \(f_1(t)\) проходить через межу; b --- термоізольована
межа; a --- це навколишня температура Равії \(10^{\circ}\) C; z ---
температура підтримується рівною \(10^{\circ}\) C.

а граничні умови будуть виконані на \(0 < t < \infty\) і визначатимуть
температуру в будь-який момент. Якою б не була ця температура, вона
повинна відповідати граничним умовам, зазначеним на рис. 3.7.

\begin{enumerate}
\def\labelenumi{\arabic{enumi}.}
\setcounter{enumi}{1}
\tightlist
\item
  Зверніть увагу, що гранична умова
\end{enumerate}

{[} u\_r(R, \textbackslash theta, t) = -\textbackslash frac\{h\}\{k\}
{[}u(R, \textbackslash theta, t) - g(\textbackslash theta, t){]} {]}

На колі не означає, що гранична температура дорівнює \(g(\theta,t)\) ,
але якщо коефіцієнт теплопередачі h великий, то ці умови фактично
еквівалентні вимогі \(g(\theta,t)\) температури на межі.

\section{ЗАВДАННЯ}\label{ux437ux430ux432ux434ux430ux43dux43dux44f-1}

\begin{enumerate}
\def\labelenumi{\arabic{enumi}.}
\item
  Намалюйте контури графіків для розв'язання змішаної задачі (3.6) для
  різних моментів часу. Чи відповідають ваші графи визначеним граничним
  умовам? Яка стаціонарна температура стрижня? Очевидно, ваша відповідь
  базується на інтуїції?
\item
  Як би ви інтерпретували наступне змішане завдання?
\end{enumerate}

{[} \textbackslash begin\{aligned\} (\textbackslash text\{YHII\}) \&
u\_t = \textbackslash alpha\^{}2 u\_\{xx\}, \& 0 \textless{} x
\textless{} 1, \& 0 \textless{} t \textless{} \textbackslash infty,
\textbackslash{} (\textbackslash text\{FY\}) \&
\textbackslash begin\{cases\}
u\textbackslash left(0,,t\textbackslash right) = 0, \textbackslash{}
u\_x\textbackslash left(1,,t\textbackslash right) = 1, \textbackslash{}
u\textbackslash left(x,,0\textbackslash right) =
\textbackslash sin\textbackslash left(\textbackslash pi
x\textbackslash right), \& 0 \textbackslash leqslant x
\textbackslash leqslant 1. \textbackslash{} \textbackslash end\{cases\}
\textbackslash end\{aligned\} {]}

Чи можете ви графічно показати рішення цієї проблеми в різні моменти
часу? Чи буде це рішення, як правило, стаціонарним? Це очевидно?

\begin{enumerate}
\def\labelenumi{\arabic{enumi}.}
\setcounter{enumi}{2}
\tightlist
\item
  Яку фізичну інтерпретацію ви можете дати проблемі
\end{enumerate}

{[} \textbackslash begin\{aligned\} (\textbackslash text\{УЧП\}) \& u\_t
= \textbackslash alpha\^{}2 u\_\{xx\}, \& 0 \textless{} x \textless{} 1,
\& 0 \textless{} t \textless{} \textbackslash infty, \textbackslash{}
(\textbackslash text\{ГУ\}) \& \textbackslash begin\{cases\} u\_x , (0,
, t) = 0, \& 0 \textless{} t \textless{} \textbackslash infty,
\textbackslash{} u\_x , (1, , t) = 0, \& 0 \textless{} t \textless{}
\textbackslash infty, \textbackslash end\{cases\} \textbackslash{}
(\textbackslash text\{НУ\}) \& u , (x, , 0) = \textbackslash sin
(\textbackslash pi x), \& 0 \textbackslash leqslant x
\textbackslash leqslant 1? \textbackslash{}
\textbackslash end\{aligned\} {]}

Чи можете ви графічно показати рішення цієї проблеми в різні моменти
часу?

Що можна сказати про стабільну температуру?

\begin{enumerate}
\def\labelenumi{\arabic{enumi}.}
\setcounter{enumi}{3}
\tightlist
\item
  Припустимо, що бічна поверхня металевого стрижня не ізольована і має
  початкову температуру 20°C, але температура одного кінця стає 50°C.
  Потім стрижень промивають рідиною з температурою 30°C. Як виглядає
  змішана задача, що відповідає цьому випадку?
\end{enumerate}

\section{ВИВЕДЕННЯ РІВНЯННЯ
ТЕПЛОПРОВІДНОСТІ}\label{ux432ux438ux432ux435ux434ux435ux43dux43dux44f-ux440ux456ux432ux43dux44fux43dux43dux44f-ux442ux435ux43fux43bux43eux43fux440ux43eux432ux456ux434ux43dux43eux441ux442ux456}

МЕТА ЛЕКЦІЇ: Показати, як можна отримати одномірне рівняння
теплопровідності

{[} u\_t = \textbackslash alpha\^{}2 u\_\{xx\} + f(x, t), {]}

спираючись лише на фундаментальний закон збереження кількості тепла.
Розглядається залежність швидкості теплопередачі від основних
термофізичних параметрів: теплопровідності, теплоємності та густини
матеріалу. Розглядаються різні модифікації базового рівняння
теплопровідності.

У кожній галузі знань існує набір базових принципів, правильність яких
очевидна. Усі інші твердження мають бути виведені з цих базових
принципів. Звісно, те, що очевидно одній людині, може здатися сумнівним
для іншої. Історія будь-якої науки полягає у постійному пошуку більш
фундаментальних принципів, що лежать в основі цієї науки. Основні
принципи мають бути настільки універсальними, щоб ніхто не сумнівався.

? --- Припущення A --- Припущення B --- Припущення C --- Припущення D
--- ? Рис. 4.1. Аксіоматичний метод.

Наприклад, одна людина може вважати, що всі факти певної науки можна
вивести з одного базового припущення, яке ми позначимо як B. На основі
припущення B вона або вона може довести теорему C, що призведе до
теореми D, а яка, у свою чергу, призведе до доведення багатьох інших
теорем (див. рисунок 4.1). Ось як виглядає прогрес у науці --- дослідник
може отримувати все більше нових результатів. Найвидатніші фізики,
хіміки та біологи або саме на цьому шляху.

З іншого боку, замість доведення нових теорем можна спробувати, якщо
можливо, знайти новий принцип (назвемо його A), який є більш
фундаментальним за B. Припущення A природно вважається більш
фундаментальним, ніж B.

Якщо B можна вивести з A. Отже, якщо можна знайти новий принцип A, тоді
межа наших знань зміщується до сфери більш фундаментальних ідей про
певну науку. У теорії теплопровідності основним принципом є закон
збереження енергії (теплова енергія). Усі інші твердження випливають із
цього базового принципу (рис. 4.2).

Закон збереження енергії

{[} u\_t = \textbackslash alpha\^{}2 u\_\{xx\} + f(x,t) {]}

Інші властивості теплопровідності Рис. 4.2. Енергозбереження: наріжний
камінь теплопровідності.

Звісно, ми могли б забути цю лекцію і використати рівняння
теплопровідності як фундаментальне поняття (дехто може так подумати,
якщо це очевидно), але тоді ми обдуримо більш серйозних студентів, адже
закон збереження вважається найфундаментальнішим принципом у науці.
Науковці, моделюючи свої задачі, часто записують відношення, що
виражають закони збереження, а потім переписують їх у вигляді рівнянь у
похідних похідних.

Повернемося до головної мети нашої лекції --- вивести рівняння
теплопровідності з рівняння збереження енергії.

\section{Виведення рівняння
теплопровідності}\label{ux432ux438ux432ux435ux434ux435ux43dux43dux44f-ux440ux456ux432ux43dux44fux43dux43dux44f-ux442ux435ux43fux43bux43eux43fux440ux43eux432ux456ux434ux43dux43eux441ux442ux456-1}

Розглянемо однорідний стрижень довжини L, щодо якого зробимо такі
припущення:

\begin{enumerate}
\def\labelenumi{\arabic{enumi}.}
\tightlist
\item
  Стрижень виготовлений з одного однорідного провідного матеріалу.
\end{enumerate}

\pandocbounded{\includegraphics[keepaspectratio]{images/_page_29_Figure_10.jpeg}}

Рис. 4.3. Тонкий теплопровідний стержень.

\begin{enumerate}
\def\labelenumi{\arabic{enumi}.}
\setcounter{enumi}{1}
\item
  Бікова поверхня стрижня термічно ізольована (тепло може поширюватися
  лише вздовж осі x).
\item
  Стрижень тонка, тобто температура всіх точок у кожному поперечному
  перерізі стрижня є сталою. Якщо розглядати частину вудки на сегменті
  \([x, x+\Delta x]\) і використовувати
\end{enumerate}

закон збереження кількості тепла, тоді можна написати:

(4.1) Загальна зміна кількості тепла на сегменті \([x, x + \Delta x] =\)
дорівнює загальній кількості тепла, що пройшло через межі + Нульова
кількість тепла, утвореного всередині сегмента \([x, x + \Delta x]\) .

Оскільки загальна кількість тепла (у калоріях) всередині сегмента
\(\{x, x + \Delta x\}\) в будь-який момент часу t обчислюється формулою:

Загальна кількість тепла всередині сегмента
\([x, x+\Delta x] = \int\limits_{-\infty}^{x+\Delta x} coAu\) (s, t) ds,

де c --- питома теплоємність матеріалу (показує здатність матеріалу
накопичувати тепло),

Щільність матеріалу,

A --- це площа поперечного перерізу стрижня.

\textbf{va}con збереження енергії (4.1) можна надати у такій
математичній формі:

(4.2)

{[} \textbackslash frac\{d\}\{dt\}
\textbackslash int\_\{x\}\^{}\{x+\textbackslash Delta x\}
c\textbackslash rho Au(s, t) ds = c\textbackslash rho A
\textbackslash int\_\{x\}\^{}\{x+\textbackslash Delta x\} u\_\{t\}(s, t)
ds = eA \textbackslash left{[} u\_\{x\}(x+\textbackslash Delta x, t) -
u\_\{x\}(x, t) \textbackslash right{]} + A
\textbackslash int\_\{x+\textbackslash Delta
x\}\^{}\{x+\textbackslash Delta x\} f(s, t) ds. {]}

\textbf{де} k --- теплопровідність матеріалу (вказує на здатність
матеріалу проводити тепло), \(f(x, t)\) --- об'ємна потужність
зовнішнього джерела тепла (калорії/см·с).

Задача полягає в тому, щоб записати рівняння (4.2) у формі, яка не
містить інтегралів. Щоб розв'язати цю задачу, нагадаємо читачеві теорему
про середнє значення з ходу інтегральної нечисельності.

\section{Теорема про середнє
значення}\label{ux442ux435ux43eux440ux435ux43cux430-ux43fux440ux43e-ux441ux435ux440ux435ux434ux43dux454-ux437ux43dux430ux447ux435ux43dux43dux44f}

Якщо функція \(f(x)\) є неперервною на сегменті {[}a, b{]}, то існує
принаймні одна точка \(\xi \in [a, b]\) така, що

{[} \textbackslash int\_\{a\}\^{}\{b\} f(x) dx = f(\textbackslash xi)
(b-a). {]}

Застосовуючи цей результат до рівняння (4.2), ми отримуємо наступний
результат

{[} coAu\_\{t\}(\textbackslash xi, t) \textbackslash Delta x =
kA\{u\_\{x\}(x + \textbackslash Delta x, t) - u\_\{x\}(x, t)\} +
Af(\textbackslash xi, t) \textbackslash Delta x, {]}

{[} x \textless{} \textbackslash xi \textless{} x + \textbackslash Delta
x, {]}

або

{[} u\_t(\textbackslash xi, t) =
\textbackslash frac\{k\}\{c\textbackslash rho\} \textbackslash left\{
\textbackslash frac\{u\_x(x + \textbackslash Delta x, t) - u\_x(x,
t)\}\{\textbackslash Delta x\} \textbackslash right\} +
\textbackslash frac\{1\}\{c\textbackslash rho\} f(\textbackslash xi, t).
{]}

Спрямовуючи \textbackslash Delta x до 0, отримуємо потрібне рівняння

(4.3)

{[} u\_t(x, t) = \textbackslash alpha\^{}2 u\_\{xx\}(x, t) + F(x, t),
{]}

де

{[}
\textbackslash alpha\^{}\{2\}=\textbackslash frac\{k\}\{c\textbackslash rho\}
{]}

--- коефіцієнт теплової дифузії,
F\textbackslash left(x,~t\textbackslash right)\(=\frac{1}{c\rho}f\left(x,\ t\right)\)---
густина джерел тепла.

Ми могли б зупинитися на цьому. Однак ми також розглянемо випадок, коли
бічна поверхня стрижня не термічно ізольована. Припустимо, що кількість
теплового потоку через бічну поверхню стрижня в цьому випадку
пропорційна різниці між температурою стрижня \(u(x,t)\) і температурою
навколишнього середовища, яка залишається сталою і дорівнює нулю. У
цьому випадку закон збереження кількості тепла веде до рівняння

{[} (4.4) u\_t = \textbackslash alpha\^{}2 u\_\{xx\} -
\textbackslash beta u + F(x, t), {]}

де \textbackslash beta --- коефіцієнт пропорційності потоку через бічну
поверхню.

\section{ПРИМІТКИ}\label{ux43fux440ux438ux43cux456ux442ux43aux438}

\begin{enumerate}
\def\labelenumi{\arabic{enumi}.}
\tightlist
\item
  Як зазначалося вище, константа k називається теплопровідністю речовини
  і чисельно дорівнює кількості тепла (у калоріях), що проходить за
  секунду через пластину товщиною 1 см і площею поперечного перерізу 1
  см³ при різниці температур на протилежних гранях пластини рівною 1°C.
  Числові значення теплопровідності для різних матеріалів можна знайти в
  довідниках з фізики або хімії. Типові значення теплопровідності
  варіюються від 1 (для міді) до значень, близьких до нуля для хороших
  теплоізоляторів. Якщо стрижень зроблений із однорідного матеріалу, то
  значення k не залежить від x. Для деяких матеріалів значення k
  залежить від температури, і тому рівняння теплопровідності в цьому
  випадку стає нелінійним:
\end{enumerate}

{[} u\_\{t\} = \textbackslash frac\{1\}\{c\textbackslash rho\}
\textbackslash frac\{\textbackslash partial\}\{\textbackslash partial
x\} \{ k (u) \textbar{} u\_\{x\} \}. {]}

Однак найчастіше теплопровідність k дуже слабо впливає на температуру u,
і цю нелінійність можна ігнорувати.

\begin{enumerate}
\def\labelenumi{\arabic{enumi}.}
\setcounter{enumi}{1}
\item
  Значення c зазвичай називають питомою теплоємністю речовини. Він
  характеризує здатність речовини зберігати теплову енергію. Наприклад,
  запечена картопля має високу теплоємність, оскільки може зберігати
  велику кількість тепла на одиницю маси (саме тому вона зберігає високу
  температуру так довго). Чисельна теплоємність дорівнює кількості
  калорій тепла, які потрібно передати 1 граму речовини, щоб підвищити
  її температуру на 1°C. Для більшості розглянутих задач значення c
  можна вважати сталим, не записуючи ні x, ні i. Значення питомої
  теплоємності для конкретних матеріалів можна знайти у фізичних
  довідниках.
\item
  Одиниці вимірювання основних величин, включених у рівняння
\end{enumerate}

теплопровідність (у системі GHS) наступна:

i-температура (градуси Цельсія), \(u_t\) --- швидкість зміни температури
(°C/s), \(u_x\) --- нахил температурної кривої (°C/см), \(u_{xx}\) ---
опуклість температурної кривої (°C/см²),

\begin{verbatim}
c = щільність тепла (калор/г \cdot °C)

k — теплопровідність (кал/см \cdot при \cdot °C), \rho — щільність (г/<sup>см3</sup>), \alpha^2 — це теплова дифузія (см<sup>²</sup>/с).

4. Слід зазначити, що теплова дифузія матеріалу \alpha^2 = \frac{k}{\epsilon \rho} прямо пропорційна теплопровідності матеріалу і обернено пропорційна густині \rho та питому теплоємності \epsilon ; Ми сподіваємося, що цей зв'язок відповідає інтуїтивним ідеям читача.
\end{verbatim}

\section{ЗАВДАННЯ}\label{ux437ux430ux432ux434ux430ux43dux43dux44f-2}

\begin{verbatim}
1. Підставити розмірність кожної з u,\ u_t,\ \dots величин у рівняння
\end{verbatim}

{[} u\_t = \textbackslash alpha\^{}2 u\_\{xx\} - \textbackslash beta u
{]}

і переконайтеся, що всі терміни мають однаковий розмір
\{\}\^{}\{\textbackslash circ\}\textbackslash mathrm\{C\}/c .

\begin{enumerate}
\def\labelenumi{\arabic{enumi}.}
\setcounter{enumi}{1}
\tightlist
\item
  Підставити розмірність кожної з величин у рівняння
\end{enumerate}

{[} u\_t = \textbackslash alpha\^{}2 u\_\{xx\} - v u\_x, {]}

де \textbackslash upsilon вимірюється в одиницях швидкості, і
переконайтеся, що всі терми мають однакову розмірність.

3 Отримаємо рівняння теплопровідності

{[} u\_\{t\} = \textbackslash frac\{1\}\{c\textbackslash rho\}
\textbackslash frac\{\textbackslash partial\}\{\textbackslash partial
x\} {[}h(x) u\_\{x\}{]} + f(x, t) {]}

Для випадку, коли теплопровідність \(k\left( x\right)\) залежить від
координати X.

\begin{enumerate}
\def\labelenumi{\arabic{enumi}.}
\setcounter{enumi}{3}
\tightlist
\item
  Припустимо, що \(u(x, t)\) --- це значення концентрації речовини в
  потоку, що рухається зі швидкістю v. Припустимо, що концентрація
  змінюється внаслідок процесів дифузії та конвекції. Виведемо рівняння
\end{enumerate}

{[} u\_1 = \textbackslash alpha\^{}2 u\_\{xx\} - v u\_x, {]}

Виходячи з того, що в будь-який момент часу в певному {[}x, x +
\textbackslash Delta x{]} сегменті матерії не з'являється і не зникає.

ПРИМІТКА. Закон збереження матерії виглядає так:

\begin{verbatim}
Зміни маси речовини на сегменті [x, x + \Delta x] = = Зміна внаслідок дифузії через межі + Зміна внаслідок переносу через межі.
\end{verbatim}

Лекція 5

\section{РОЗДІЛЕННЯ
ЗМІННИХ}\label{ux440ux43eux437ux434ux456ux43bux435ux43dux43dux44f-ux437ux43cux456ux43dux43dux438ux445}

МЕТА ЛЕКЦІЇ: Познайомити читача з потужним методом розділення змінних і
показати, як цей метод можна використати для розв'язання добре відомої
проблеми дифузії. Оскільки цей метод створює труднощі для деяких учнів
(головним чином через громіздкі алгебранічні перетворення), ми
супроводжуємо презентацію інтуїтивними міркуваннями.

Основна ідея методу полягає в тому, щоб розкласти початкову умову на
найпростіші компоненти, знайти відгук системи на кожен найпростіший
компонент і підсумувати всі відповіді. Таким чином можна знайти
відповідь на довільну початкову умову.

Існуючий метод поділу змінних певною мірою зменшує можливість такої
інтерпретації методу, але все ж не виключає його.

Метод розділення змінних є одним із найвідоміших методів розв'язання
зсувних задач і використовується, якщо:

\begin{itemize}
\tightlist
\item
  \begin{enumerate}
  \def\labelenumi{\arabic{enumi}.}
  \tightlist
  \item
    Рівняння є лінійним і однорідним (не обов'язково з постійними
    коефіцієнтами).
  \end{enumerate}
\item
  \begin{enumerate}
  \def\labelenumi{\arabic{enumi}.}
  \setcounter{enumi}{1}
  \tightlist
  \item
    Граничні умови визначаються як
  \end{enumerate}
\end{itemize}

{[} \textbackslash alpha u\_x(0, t) + \textbackslash beta u(0, t) = 0,
{]}

\textbackslash gamma \(u_x\)(1, t) + \textbackslash delta \(u(1, t)\) =
0, де \textbackslash alpha , \textbackslash beta , \textbackslash gamma
і \textbackslash delta --- константи (граничні умови, задані в такій
формі, називаються лінійними однорідними граничними умовами).

Метод був створений за часів Фур'є (зазвичай відомий як метод Фур'є) і,
здається, є найпопулярнішим (де це застосовно).

\pandocbounded{\includegraphics[keepaspectratio]{images/_page_34_Figure_8.jpeg}}

Рис. 5.1. Схематичне представлення задачі дифузії.

Замість того, щоб вивчати метод у загальному випадку, давайте спочатку
розглянемо конкретну задачу (загальний випадок обговоримо пізніше).
Розглянемо змішану задачу типу дифузії: щоб знайти розв'язок

(PPP)

{[} u\_t = \textbackslash alpha\^{}2 u\_\{xx\} {]}

{[} , 0 \textless{} x \textless{} 1 , 0 \textless{} t \textless{}
\textbackslash infty , {]}

що задовольняє граничні умови

(GI)

{[} \textbackslash begin\{cases\} u(0, t) = 0, \& 0 \textless{} t
\textless{} \textbackslash infty, \textbackslash{} u(1, t) = 0, \& 0
\textless{} t \textless{} \textbackslash infty,
\textbackslash end\{cases\} {]}

та початкова умова

(ГУ)

{[} u(x, 0) = \textbackslash varphi(x), 0 \textbackslash le x
\textbackslash le 1. {]}

Перш ніж почати розділяти змінні, давайте дамо фізичну інтерпретацію
нашої проблеми. Отже, існує стрижень скінченної довжини, кінці якого
підтримуються при сталій температурі, рівній нулю (насправді кінці можна
підтримувати при значно вищій температурі, значення якої вважається
початковою точкою). Додаткові дані про завдання подаються як початкова
умова. Наша мета --- знайти розподіл температури \(u\left(x,t\right)\) в
наступний час.

\section{Загальні принципи методу розділення
змінних}\label{ux437ux430ux433ux430ux43bux44cux43dux456-ux43fux440ux438ux43dux446ux438ux43fux438-ux43cux435ux442ux43eux434ux443-ux440ux43eux437ux434ux456ux43bux435ux43dux43dux44f-ux437ux43cux456ux43dux43dux438ux445}

Для найпростішого диференціального рівняння в похідних похідних
розділення змінних є пошуком розв'язків у вигляді

{[} u(x, t) = X(x)T(t), {]}

де \(X(x)\) --- це функція, що залежить лише від змінної x, а \(T(t)\)
--- функція, що залежить лише від t. Це розв'язок у певному сенсі є
найпростішим, оскільки температура \(u(x, t)\), представлена у такій
формі, зберігатиме «форму» профілю у різний час часу t (див. рисунок
5.2).

\pandocbounded{\includegraphics[keepaspectratio]{images/_page_35_Figure_6.jpeg}}

РИСУНОК 5.2 Графік функції X(x)\$\$T(t) у різних точках часу t.

\begin{verbatim}
Загальна ідея полягає в тому, щоб знайти нескінченну кількість таких часткових диференціальних розв'язків (які задовольняють граничні умови). Ці найпростіші функції u_n\left(x,\ t\right) = X_n\left(x\right)T_n\left(t\right) (які називаються фундаментальними розв'язками) є, так би мовити, елементарними будівельними блоками, з яких будується розв'язок нашої проблеми. Розв'язок нашої задачі u\left(x,\ t\right) у вигляді такої лінійної комбінації фундаментальних розв'язків X_n\left(x\right)T_n\left(t\right) , що отримана сума
\end{verbatim}

{[} \textbackslash sum\_\{n=1\}\^{}\{\textbackslash infty\} A\_n X\_n(x)
T\_n(t) {]}

задовольняє початкові умови. Оскільки ця сума задовольняє рівняння та
граничні умови, вона є розв'язком початкової задачі. Ми маємо робити всі
ці розрахунки достатньо детально.

\section{Розділення
змінних}\label{ux440ux43eux437ux434ux456ux43bux435ux43dux43dux44f-ux437ux43cux456ux43dux43dux438ux445-1}

КРОК 1. (Пошук елементарних розв'язків рівняння з частковими похідними.)

Ми хочемо знайти функцію \(u(x, t)\), яка є розв'язком

Наступного завдання:

{[} \textbackslash begin\{array\}\{ll\}
(\textbackslash mathrm\{Y\}\^{}\{\textbackslash mathbf\{U\}\}\textbackslash Pi)
\& u\_t = \textbackslash alpha\^{}2 u\_\{xx\}, \& 0 \textless{} x
\textless{} 1, \& 0 \textless{} t \textless{} \textbackslash infty,
\textbackslash{} (\textbackslash Gamma \textbackslash mathrm\{Y\}) \&
\textbackslash begin\{cases\}
u\textbackslash left(0,,t\textbackslash right) = 0, \& 0 \textless{} t
\textless{} \textbackslash infty, \textbackslash{}
u\textbackslash left(1,,t\textbackslash right) = 0, \& 0 \textless{} t
\textless{} \textbackslash infty, \textbackslash end\{cases\}
\textbackslash{} (\textbackslash mathrm\{H\}\textbackslash mathrm\{Y\})
\& u\textbackslash left(x,,0\textbackslash right) =
\textbackslash varphi\textbackslash left(x\textbackslash right), \& 0
\textless{} x \textbackslash leq 1. \textbackslash end\{array\} {]}

Пошукаємо розв'язки, які представлені у вигляді \(u(x, t)\) = \(X(x)\)
\(T(t)\). Для цього підставляємо вираз \(X(x)\) \(T(t)\) у рівняння.
Внаслідок підстановки отримуємо

{[} X(x) T'(t) = q\^{}2 X'\,'(x) T(t). {]}

\textbf{Тепер} виконаємо операцію, властиву цьому методу:
\emph{поділимо} обидві сторони останнього рівняння на
\textbackslash alpha\^{}2 \(X(x)\) \(T(t)\) , у результаті чого
отримуємо

{[} \textbackslash frac\{T'(t)\}\{\textbackslash alpha\^{}2 T(t)\} =
\textbackslash frac\{X'\,'(x)\}\{X(x)\}. {]}

з нульовими граничними умовами. Цей новий виклик можна подолати

\begin{itemize}
\item
  \begin{enumerate}
  \def\labelenumi{\arabic{enumi}.}
  \tightlist
  \item
    Шляхом ділення змінних (якщо нова задача виявиться однорідною, тобто
    \(F(x, t)\) = 0).
  \end{enumerate}

  \begin{itemize}
  \tightlist
  \item
    \begin{enumerate}
    \def\labelenumi{\arabic{enumi}.}
    \setcounter{enumi}{1}
    \tightlist
    \item
      Методом інтегральних перетворень і розкладу власними функціями,
      якщо F(x, t) \textbackslash neq 0 .
    \end{enumerate}
  \end{itemize}
\end{itemize}

Метод розділення змінних, про який йшлося в попередній лекції, є досить
потужним, а розв'язки, отримані за його допомогою, подаються у зручній
формі. Однак читач повинен розуміти, що цей метод не застосовний до всіх
задач. Щоб метод розділення змінних був застосовним, граничні умови
мають бути лінійними, однорідними, тобто представленими як

(6.1)

{[} \textbackslash begin\{aligned\} \textbackslash alpha\_1 u\_x(0, t) +
\textbackslash beta\_1 u(0, t) \&= 0, \textbackslash{}
\textbackslash alpha\_2 u\_x(L, t) + \textbackslash beta\_2 u(L, t) \&=
0. \textbackslash end\{aligned\} {]}

Мета цієї лекції --- показати, як виникає задача з неоднорідними
граничними умовами форми

(6.2)

{[} \textbackslash begin\{aligned\} \&
(\textbackslash mathsf\{У\}\textbackslash mathsf\{Ч\}\textbackslash Pi)
\& u\_t = \textbackslash alpha\^{}2 u\_\{xx\}, \textbackslash{} \&
(\textbackslash mathsf\{\textbackslash Gamma\}\textbackslash mathsf\{Y\})
\& \textbackslash begin\{cases\} \textbackslash alpha\_1 u\_x , (0, , t)
+ \textbackslash beta\_1 u , (0, , t) = g\_1 , (t) \textbackslash{}
\textbackslash alpha\_2 u\_x , (L, , t) + \textbackslash beta\_2 u , (L,
, t) = g\_2 , (t) \textbackslash end\{cases\} \textbackslash text\{
(неоднородные \}
\textbackslash mathsf\{\textbackslash Gamma\}\textbackslash mathsf\{Y\}),
\textbackslash{} \&
(\textbackslash mathsf\{H\}\textbackslash mathsf\{Y\}) \& u , (x, , 0) =
\textbackslash varphi , (x) \textbackslash end\{aligned\} {]}

можна розв'язати, зведивши її до задачі з однорідними граничними
умовами.

Нове завдання можна розв'язати іншими методами

(наприклад, шляхом розкладу за власними функціями).

Почнемо з надзвичайно простого прикладу перетворення гетерогенних
граничних умов на однорідні.

\section{Перетворення неоднорідних граничних умов у однорідні граничні
умови}\label{ux43fux435ux440ux435ux442ux432ux43eux440ux435ux43dux43dux44f-ux43dux435ux43eux434ux43dux43eux440ux456ux434ux43dux438ux445-ux433ux440ux430ux43dux438ux447ux43dux438ux445-ux443ux43cux43eux432-ux443-ux43eux434ux43dux43eux440ux456ux434ux43dux456-ux433ux440ux430ux43dux438ux447ux43dux456-ux443ux43cux43eux432ux438}

Розглянемо проблему поширення тепла в термоізольованому стрижні, кінці
якого підтримуються при сталих температурах \(k_1\) і \(k_2\) , тобто

{[} \textbackslash begin\{array\}\{lll\}
(\textbackslash mathrm\{Y\}\textbackslash mathrm\{H\}\textbackslash Pi)
\& u\_t = \textbackslash alpha\^{}2 u\_\{xx\}, \& 0 \textless{} x
\textless{} L, \& 0 \textless{} t \textless{} \textbackslash infty,
\textbackslash{} (6.3) \&
(\textbackslash Gamma\textbackslash mathrm\{Y\}) \&
\textbackslash begin\{cases\} u ~(0, ~t) = k\_t, \textbackslash{} u ~(L,
~t) = k\_2, \textbackslash end\{cases\} \& 0 \textless{} t \textless{}
\textbackslash infty, \textbackslash{}
(\textbackslash mathrm\{H\}\textbackslash mathrm\{Y\}) \& u ~(x, ~0) =
\textbackslash varphi ~(x), \& 0 \textbackslash leqslant x
\textbackslash leqslant L. \textbackslash end\{array\} {]}

Складність цієї задачі полягає в тому, що, оскільки граничні умови в ній
є гетерогенними, ми не можемо розв'язати її шляхом розділення змінних.
Однак очевидно, що в t \textbackslash to \textbackslash infty розв'язок
нашої задачі схиляється до стаціонарного розв'язку, який лінійно (вздовж
x) змінюється від температури \(k_1\) до температури \(k_2\) (рис. 6.1).
Іншими словами, розумно припустити, що температуру в нашій задачі можна
представити як суму

\section{двох
термінів}\label{ux434ux432ux43eux445-ux442ux435ux440ux43cux456ux43dux456ux432}

Стаціонарна + тимчасова. Частина розв'язку, що залежить від початкової
умови, прямує до нуля при \(t \\to \\infty\). Тому покладемо

{[} u(x, t) = \textbackslash left{[}k\_1 +
\textbackslash frac\{x\}\{L\}(k\_2 - k\_1)\textbackslash right{]} + U(x,
t). {]}

У цьому випадку наше завдання --- знайти перехідну температуру (U(x,
t)). Заміна

{[} u(x, t) = \textbackslash left{[}k\_1 +
\textbackslash frac\{x\}\{L\}(k\_2 - k\_1)\textbackslash right{]} + U(x,
t) {]}

У початковій задачі (6.3) ми отримуємо нову задачу відносно невідомої
функції U\textbackslash left(x,~t\textbackslash right) . Ви можете
вирішити проблему Р и с

\pandocbounded{\includegraphics[keepaspectratio]{images/_page_45_Figure_7.jpeg}}

\begin{verbatim}
. 6.1. Розв'язання задачі (6.3) у різний час часу a- початкова температура u\left(x,\ 0\right)=\phi\left(x\right); b - стаціонарна температура u\left(x,\ \infty\right)=k_{1}+\frac{x}{L}\left(k_{2}-k_{1}\right); b - перехідна температура.

Відносно нова невідома функція U\left(x,\,t\right) і додає її до стаціонарного розв'язку, що дає бажану функцію u\left(x,\,t\right) . Виконавши ці прості перетворення з (6.3),
\end{verbatim}

Ми отримаємо

(6.4)

{[} U\_\{t\} = \textbackslash alpha\^{}\{2\}U\_\{xx\}
\textbackslash quad 0 \textless{} x \textless{} L, {]}

{[} (\textbackslash Gamma Y) \textbackslash quad
\textbackslash begin\{cases\} U(0, t) = 0, \textbackslash{} U(L, t) = 0,
\textbackslash end\{cases\} \textbackslash quad 0 \textless{} t
\textless{} \textbackslash infty, {]}

{[} (HY) \textbackslash quad U(x, 0) = \textbackslash varphi(x) -
\textbackslash left{[}k\_\{1\} + \textbackslash frac\{x\}\{L\}(k\_\{2\}
- k\_\{1\})\textbackslash right{]} =
\textbackslash overline\{\textbackslash varphi\}(x), {]}

де \textbackslash overline\{\textbackslash phi\}(x) --- нова, але відома
початкова умова.

Ця задача стосується не лише однорідного рівняння, а, на щастя,
однорідних граничних умов, що дозволяє розв'язувати її методом
розділення змінних. Насправді, читач, ймовірно, пам'ятає рішення:

(6.5)

{[} U(x, t) = \textbackslash sum\_\{n=1\}\^{}\{\textbackslash infty\}
a\_n e\textsuperscript{\{-(n\textbackslash pi\textbackslash alpha)}4 t\}
\textbackslash sin(n\textbackslash pi x/L), {]}

рей

{[} a\_n = \textbackslash frac\{2\}\{L\}
\textbackslash int\_\{-\textbackslash infty\}\^{}\{L\}
\textbackslash overline\{\textbackslash phi\}(\textbackslash xi)
\textbackslash sin(\textbackslash pi n\textbackslash xi/L)
d\textbackslash xi. {]}

Це все для сердечників із фіксованою температурою на кінцях. А як щодо
більш реалістичних GU з правою стороною, залежною від часу? Основні ідеї
тут ті ж, що й у попередній задачі, але вони набагато складніші.

\section{Перетворення граничних умов, залежних від часу, у
нульові}\label{ux43fux435ux440ux435ux442ux432ux43eux440ux435ux43dux43dux44f-ux433ux440ux430ux43dux438ux447ux43dux438ux445-ux443ux43cux43eux432-ux437ux430ux43bux435ux436ux43dux438ux445-ux432ux456ux434-ux447ux430ux441ux443-ux443-ux43dux443ux43bux44cux43eux432ux456}

Розглянемо типову задачу

(PPP)

{[} u\_t = \textbackslash alpha\^{}2 u\_\{xx\} {]}

, 0 \textless{} x \textless{} L , 0 \textless{} t \textless{}
\textbackslash infty , 0 \textless{} t \textless{} \textbackslash infty
, 0 \textless{} t \textless{} \textbackslash infty , 0 \textless{} t
\textless{} \textbackslash infty , 0 \textless{} t \textless{}
\textbackslash infty , 0 \textless{} t \textless{} \textbackslash infty
, 0 \textless{} t \textless{} \textbackslash infty , 0 \textless{} t
\textless{} \textbackslash infty , 0 \textless{} t \textless{}
\textbackslash infty , 0 \textless{} t \textless{} \textbackslash infty
, 0 \textless{} t \textless{} \textbackslash infty , 0 \textless{} t
\textless{} \textbackslash infty , 0 \textless{} t \textless{}
\textbackslash infty , 0 \textless{} t \textless{} \textbackslash infty
, 0 \textless{} t \textless{} \textbackslash infty , 0 \textless{} t
\textless{} \textbackslash infty , 0 \textless{} t \textless{}
\textbackslash infty , 0 \textless{} t \textless{} \textbackslash infty
, 0 \textless{} t \textless{} \textbackslash infty , 0 \textless{} t
\textless{} \textbackslash infty , 0 \textless{} t \textless{}
\textbackslash infty , 0 \textless{} t \textless{} \textbackslash infty
, 0 \textless{} t \textless{} \textbackslash infty , 0 \textless{} t
\textless{} \textbackslash infty , 0 \textless{} t \textless{}
\textbackslash infty , 0 \textless{} t \textless{} \textbackslash infty
, , 0 \textless{} t \textless{} \textbackslash infty , 0 \textless{} t
\textless{} \textbackslash infty , 0 \textless{} t \textless{}
\textbackslash infty , 0 \textless{} t \textless{} \textbackslash infty
, 0 \textless{} t \textless{} \textbackslash infty , 0 \textless{} t
\textless{} \textbackslash infty , 0 \textless{} t \textless{}
\textbackslash infty , 0 \textless{} t \textless{} \textbackslash infty
, 0 \textless{} t \textless{} \textbackslash infty , 0 \textless{} t
\textless{} \textbackslash infty , 0 \textless{} t \textless{}
\textbackslash infty , 0 \textless{} t \textless{} \textbackslash infty
, 0 \textless{} t \textless{} \textbackslash infty , 0 \textless{} t
\textless{} \textbackslash infty , 0 \textless{} t \textless{}
\textbackslash infty , 0 \textless{} t \textless{} \textbackslash infty
, 0 \textless{} t \textless{} \textbackslash infty , 0 \textless{} t
\textless{} \textbackslash infty , 0 \textless{} t \textless{}
\textbackslash infty , 0 \textless{} t \textless{} \textbackslash infty
, 0 \textless{} t \textless{} \textbackslash infty , 0 \textless{} t
\textless{} \textbackslash infty , 0 \textless{} t \textless{}
\textbackslash infty , 0 \textless{} t \textless{} \textbackslash infty
, 0 \textless{} t \textless{} \textbackslash infty , 0 \textless{} t
\textless{} \textbackslash infty , 0 \textless{} t \textless{}
\textbackslash infty , 0 \textless{} t \textless{} \textbackslash infty
, 0 \textless{} t \textless{} \textbackslash infty , 0 \textless{} t
\textless{} \textbackslash infty , 0 \textless{} t \textless{}
\textbackslash infty , 0 \textless{} t \textless{} \textbackslash infty
, 0 \textless{} t \textless{} \textbackslash infty , 0 \textless{} t
\textless{} \textbackslash infty , 0 \textless{} t \textless{}
\textbackslash infty , 0 \textless{} t \textless{} \textbackslash infty
, 0 \textless{} t \textless{} \textbackslash infty , 0 \textless{} t
\textless{} \textbackslash infty , 0 \textless{} t \textless{}
\textbackslash infty , 0 \textless{} t \textless{} \textbackslash infty
, 0 \textless{} t \textless{} \textbackslash infty , 0 \textless{} t
\textless{} \textbackslash infty , 0 \textless{} t \textless{}
\textbackslash infty , 0 \textless{} t \textless{} \textbackslash infty
, 0 \textless{} t \textless{} \textbackslash infty , 0 \textless{} t
\textless{} \textbackslash infty , 0 \textless{} t \textless{}
\textbackslash infty , 0 \textless{} t \textless{} \textbackslash infty
, 0 \textless{} t \textless{} \textbackslash infty , 0 \textless{} t
\textless{} \textbackslash infty , 0 \textless{} t \textless{}
\textbackslash infty , 0 \textless{} t \textless{} \textbackslash infty
, 0 \textless{} t \textless{} \textbackslash infty , 0 \textless{} t
\textless{} \textbackslash infty , 0 \textless{} t \textless{}
\textbackslash infty , 0 \textless{} t \textless{} \textbackslash infty
, 0 \textless{} t \textless{} \textbackslash infty , 0 \textless{} t
\textless{} \textbackslash infty , 0 \textless{} t \textless{}
\textbackslash infty , 0 \textless{} t \textless{} \textbackslash infty
, 0 \textless{} t \textless{} \textbackslash infty , 0 \textless{} t
\textless{} \textbackslash infty , 0 \textless{} t \textless{}
\textbackslash infty , 0 \textless{} t \textless{} \textbackslash infty
, 0 \textless{} t \textless{} \textbackslash infty , 0 \textless{} t
\textless{} \textbackslash infty , 0 \textless{} t \textless{}
\textbackslash infty , 0 \textless{} t \textless{} \textbackslash infty
, 0 \textless{} t \textless{} \textbackslash infty , 0 \textless{} t
\textless{} \textbackslash infty , 0 \textless{} t \textless{}
\textbackslash infty , 0 \textless{} t \textless{} \textbackslash infty
, 0 \textless{} t \textless{} \textbackslash infty , 0 \textless{} t
\textless{} \textbackslash infty , 0 \textless{} t \textless{}
\textbackslash infty , 0 \textless{} t \textless{} \textbackslash infty
, 0 \textless{} t \textless{} \textbackslash infty , 0 \textless{} t
\textless{} \textbackslash infty , 0 \textless{} t \textless{}
\textbackslash infty , 0 \textless{} t \textless{} \textbackslash infty
, 0 \textless{} t \textless{} \textbackslash infty , 0 \textless{} t
\textless{} \textbackslash infty , 0 \textless{} t \textless{}
\textbackslash infty , 0 \textless{} t \textless{} \textbackslash infty
, 0 \textless{} t \textless{} \textbackslash infty , 0 \textless{} t
\textless{} \textbackslash infty , 0 \textless{} t \textless{}
\textbackslash infty , 0 \textless{} t \textless{} \textbackslash infty
, 0 \textless{} t \textless{} \textbackslash infty , 0 \textless{} t
\textless{} \textbackslash infty , 0 \textless{} t \textless{}
\textbackslash infty , 0 \textless{} t \textless{} \textbackslash infty
, 0 \textless{} t \textless{} \textbackslash infty , 0 \textless{} t
\textless{} \textbackslash infty , 0 \textless{} t \textless{}
\textbackslash infty , 0 \textless{} t \textless{} \textbackslash infty
, 0 \textless{} t \textless{} \textbackslash infty , 0 \textless{} t
\textless{} \textbackslash infty , 0 \textless{} t \textless{}
\textbackslash infty , 0 \textless{} t \textless{} \textbackslash infty
, 0 \textless{} t \textless{} \textbackslash infty , 0 \textless{} t
\textless{} \textbackslash infty , 0 \textless{} t \textless{}
\textbackslash infty , 0 \textless{} t \textless{} \textbackslash infty
, 0 \textless{} t \textless{} \textbackslash infty , 0 \textless{} t
\textless{} \textbackslash infty , 0 \textless{} t \textless{}
\textbackslash infty , 0 \textless{} t \textless{} \textbackslash infty
, 0 \textless{} t \textless{} \textbackslash infty , 0 \textless{} t
\textless{} \textbackslash infty , , 0 \textless{} t \textless{}
\textbackslash infty , 0 \textless{} t \textless{} \textbackslash infty
, 0 \textless{} t \textless{} \textbackslash infty , 0 \textless{} t
\textless{} \textbackslash infty , 0 \textless{} t \textless{}
\textbackslash infty , 0 \textless{} t \textless{} \textbackslash infty
, 0 \textless{} t \textless{} \textbackslash infty , 0 \textless{} t
\textless{} \textbackslash infty , 0 \textless{} t \textless{}
\textbackslash infty , 0 \textless{} t \textless{} \textbackslash infty
, 0 \textless{} t \textless{} \textbackslash infty , 0 \textless{} t
\textless{} \textbackslash infty , 0 \textless{} t \textless{}
\textbackslash infty , 0 \textless{} t \textless{} \textbackslash infty
, 0 \textless{} t \textless{} \textbackslash infty , 0 \textless{} t
\textless{} \textbackslash infty , 0 \textless{} t \textless{}
\textbackslash infty , 0 \textless{} t \textless{} \textbackslash infty
, 0 \textless{} Щоб перетворити ці граничні умови на нульові умови, ми
(після деяких спроб і помилок) зупинилися на такій формі розв'язку:

(6.7)

{[} u(x, t) = A(t){[}1-x/L{]} + B(t){[}x/L{]} + U(x, t), {]}

де функції \(A\left(t\right)\) і \(B\left(t\right)\) обираються так, що
«квазі-стаціонарна» частина розв'язку (6.7)

(6.8)

{[} S(x, t) = A(t){[}1 - x/L{]} + B(t){[}x/L{]} {]}

задовольняло граничні умови початкової задачі. У цьому випадку функція
U\textbackslash left(x,~t\textbackslash right) задовольнятиме однорідні
граничні умови.

Підставляючи функцію \(S(x, t)\) у граничні умови

{[} S(0, t) = g\_1(t), {]}

\(S_x\)(L, t) + hS(L, t) = \(g_2\)(t) призводить до двох рівнянь, з яких
можна визначити \(A(t)\) і \(B(t)\). У результаті ми отримуємо

(6.9)

{[} A(t) = g\_1(t), B(t) = \textbackslash frac\{g\_1(t) + Lg\_2(t)\}\{1
+ Lh\}. {]}

Отже.

{[} u(x, t) = g\_1(t) {[}1 - x/L{]} + \textbackslash frac\{g\_1(t) +
Lg\_2(t)\}\{1 + Lh\} {[}x/L{]} + U(x, t). {]}

Якщо ми підставимо ці вирази для you(x, t) у початкову задачу (це
просять читача), отримаємо нову задачу для невідомої функції
\(U(x, t)\):

\begin{verbatim}
\begin{array}{llllllllllllllllllllllllllllllllllll
\end{verbatim}

Тепер ми стикаємося з новою проблемою однорідних граничних умов (на
жаль, рівняння стало гетерогенним). Ми не можемо розв'язати цю задачу
методом розділення змінних, але якщо читач погоджується чекати, то за
кілька лекцій він ознайомиться з розв'язком цієї задачі методом цілісних
перетворень і розкладання за власними функціями.

\section{ЗАУВАЖЕННЯ.}\label{ux437ux430ux443ux432ux430ux436ux435ux43dux43dux44f.}

\begin{enumerate}
\def\labelenumi{\arabic{enumi}.}
\item
  Як зазначалося раніше, мета цієї лекції --- навчити читача, як
  перетворювати задачі з неунірольними граничними умовами на задачі з
  нульовими граничними умовами. Після цього може статися, що нове
  диференціальне рівняння в похідних похідних стане однорідним. Якщо це
  станеться, ми можемо вважати себе щасливцями, оскільки в цьому випадку
  проблему можна розв'язати методом розділення змінних (див. перший
  приклад з цієї лекції). Якщо перетворене рівняння виявиться
  гетерогенним, тоді нову задачу потрібно розв'язати іншим методом.
\item
  Лінійні неоднорідні граничні умови найзагальної форми
\end{enumerate}

{[} \textbackslash alpha\_1 u\_x(0, t) + \textbackslash beta\_1 u(0, t)
= g\_1(t), {]}

\textbackslash alpha\_2 \(u_x\)(L, t) + \textbackslash beta\_2
\(u(L, t)\) = \(g_2\)(t) Ви також можете перетворити умови на нульові
граничні умови, використовуючи метод, описаний у другому прикладі.
Звісно, нове рівняння, ймовірно, буде гетерогенним.

\begin{itemize}
\tightlist
\item
  \begin{enumerate}
  \def\labelenumi{\arabic{enumi}.}
  \setcounter{enumi}{2}
  \tightlist
  \item
    Деякі методи розв'язку не містять вимог щодо унікальності граничних
    умов і, отже, не потребують попередніх перетворень. Пізніше, коли ми
    ознайомимося з перетворенням Лапласа, побачимо, що нульові граничні
    умови є необов'язковими (у випадку нульових граничних умов розв'язок
    дещо простіший, ніж в інших випадках).
  \end{enumerate}
\item
  \begin{enumerate}
  \def\labelenumi{\arabic{enumi}.}
  \setcounter{enumi}{3}
  \tightlist
  \item
    Для граничних умов типу
  \end{enumerate}
\end{itemize}

{[} u(0, t) = g\_1(t), {]}

\(u(L, t)\) = \(g_2\)(t) Метод, описаний у другому прикладі, призводить
до наступної форми розв'язання задач:

{[} u(x, t) = \textbackslash left\{ g\_1(t) +
\textbackslash frac\{x\}\{L\} \textbackslash left{[} g\_2(t) - g\_1(t)
\textbackslash right{]} \textbackslash right\} + U(x, t). {]}

\section{ЗАВДАННЯ}\label{ux437ux430ux432ux434ux430ux43dux43dux44f-3}

\begin{enumerate}
\def\labelenumi{\arabic{enumi}.}
\tightlist
\item
  Знайти розв'язок наступної змішаної задачі:
\end{enumerate}

{[} \textbackslash begin\{array\}\{lll\} (\textbackslash text\{УЧП\}) \&
u\_t = \textbackslash alpha\^{}2 u\_\{xx\} + x \textbackslash cos t, \&
0 \textless{} x \textless{} 1, \textbackslash{} (\textbackslash Gamma
\textbackslash text\{У\}) \& \textbackslash begin\{cases\} u ~(0, ~t) =
1, \textbackslash{} u\_x ~(1, ~t) + hu ~(1, ~t) = 1,
\textbackslash end\{cases\} \& 0 \textless{} t \textless{}
\textbackslash infty, \textbackslash{} (\textbackslash text\{HY\}) \& u
~(x, ~0) = \textbackslash sin (\textbackslash pi x), \& 0
\textbackslash leqslant x \textbackslash leqslant 1,
\textbackslash end\{array\} {]}

попередньо перетворюючи його на задачу з нульовими граничними умовами.
Чи відповідає рішення, яке ви отримуєте, вашим задумам щодо цієї
проблеми?

Завдання

(PPP)

{[} u\_t = u\_\{xx\}, \textbackslash quad 0 \textless{} x \textless{} 1,
{]}

(GI)

{[} \textbackslash begin\{cases\} u(0, t) = 0, \& 0 \textless{} t
\textless{} \textbackslash infty, \textbackslash{} u(1, t) = 1, \& 0
\textless{} t \textless{} \textbackslash infty,
\textbackslash end\{cases\} {]}

(НУ) \(u(x, 0)\) = x\^{}2, \textbackslash quad 0 \textbackslash le x
\textbackslash le 1, Перетворіть у задачу з нульовими граничними умовами
і розв'яжіть її. Як виглядатиме рішення у різних

Миті часу? Чи відповідає це вашим інтуїтивним ідеям? Що таке стаціонарне
рішення? Як виглядає перехідне рішення?

\begin{enumerate}
\def\labelenumi{\arabic{enumi}.}
\setcounter{enumi}{2}
\tightlist
\item
  Завдання
\end{enumerate}

{[} \textbackslash begin\{array\}\{ll\} (\textbackslash text\{УЧП\}) \&
u\_t = u\_\{xx\}, \& 0 \textless{} x \textless{} 1, \textbackslash{}
(\textbackslash Gamma \textbackslash text\{У\}) \&
\textbackslash begin\{cases\} u\_x(0, t) = 0, \textbackslash{} u\_x(1,
t) + hu(1, t) = 1, \textbackslash{} (\textbackslash text\{HY\}) \& u(x,
0) = \textbackslash sin(\textbackslash pi x),
\textbackslash end\{cases\} \& 0 \textless{} x \textless{} 1,
\textbackslash{} 0 \textless{} t \textless{} \textbackslash infty,
\textbackslash{} 0 \textbackslash le x \textbackslash le 1,
\textbackslash end\{array\} {]}

Перетворимо на задачу з нульовими граничними умовами. Чи буде нове
рівняння однорідним?

Лекція 7

\section{РОЗВ'ЯЗАННЯ СКЛАДНІШИХ ЗАДАЧ ШЛЯХОМ ПОДІЛУ
ЗМІННИХ}\label{ux440ux43eux437ux432ux44fux437ux430ux43dux43dux44f-ux441ux43aux43bux430ux434ux43dux456ux448ux438ux445-ux437ux430ux434ux430ux447-ux448ux43bux44fux445ux43eux43c-ux43fux43eux434ux456ux43bux443-ux437ux43cux456ux43dux43dux438ux445}

МЕТА ЛЕКЦІЇ: Показати, як більш складні проблеми теплопровідності можна
розв'язати методом розділення змінних. У цій лекції ми головним чином
розглядаємо такі задачі, розв'язання яких дозволить читачеві краще
ознайомитися з методом ділення змінних. Сподіваємося, що читач зможе
екстраполювати викладені тут ідеї для власних завдань.

У тій же лекції читач ознайомиться з задачею за її власними функціями,
або, як її ще називають, проблемою Штурма-Ліувіля. Обговорюються деякі
властивості проблеми Штурма-Ліувіля.

Як згадувалося вище, мета цієї лекції --- розв'язати змішану задачу
шляхом розділення змінних. Можливо, саме таким методом читач вирішить
свою проблему. Водночас ми сподіваємося, що він зможе перенести метод,
розглянутий у цій лекції, для власних завдань.

Почнемо з одномірної задачі теплопровідності, в якій одна з граничних
умов містить похідну.

\section{Задача теплопровідності з похідною в граничній
умові}\label{ux437ux430ux434ux430ux447ux430-ux442ux435ux43fux43bux43eux43fux440ux43eux432ux456ux434ux43dux43eux441ux442ux456-ux437-ux43fux43eux445ux456ux434ux43dux43eux44e-ux432-ux433ux440ux430ux43dux438ux447ux43dux456ux439-ux443ux43cux43eux432ux456}

Розглянемо пристрій, показаний на рис. 7.1. Припустимо, що верхній
кінець стрижня підтримується при нульовій температурі, тобто \(u(0, t)\)
= 0, а нижній кінець занурений у воду, температура якої також стала і
дорівнює нулю (у цьому випадку --- нуль Р и с

\pandocbounded{\includegraphics[keepaspectratio]{images/_page_50_Picture_4.jpeg}}

. 7.1. Схематичне представлення зсувної задачі.

\pandocbounded{\includegraphics[keepaspectratio]{images/_page_50_Figure_6.jpeg}}

Рис. 7.2. Природа кривих, які задовольняють

{[} (\textbackslash Gamma Y) = {]}

{[} \textbackslash begin\{cases\} u(0, t) = 0, \textbackslash{} u\_x(1,
t) = -hu(1, t). \textbackslash end\{cases\} {]}

відповідає певній еталонній температурі). Теплообмін між стрижнем і
водою відбувається згідно із законом Ньютона, який записується у x=1 як
наступна гранична умова:

{[} u\_x(1, t) = -hu(1, t). {]}

Тепер припустимо, що початкова температура стрижня дорівнює \(u(x, 0)\)
= x, але одразу при t \textgreater{} 0 включаємо граничні умови. Щоб
визначити температуру в стрижні, необхідно розв'язати проблему зміщення

\begin{verbatim}
\begin{aligned} (\mathrm{УЧ}\Pi) & u_t = \alpha^2 u_{xx}, & 0 < x < 1, & 0 < t < \infty, \\ (7.1) & (\Gamma\mathrm{У}) & \begin{cases} u_t(0, t) = 0 \\ u_x(1, t) + hu_t(1, t) = 0 \end{cases} & \text{(однородные } \Gamma\mathrm{У}), \\ (\mathrm{H}\mathrm{У}) & u_t(x, 0) = x, & 0 \leq x \leq 1. \end{aligned}
\end{verbatim}

Щоб використати метод розділення змінних, розбиємо процес розв'язання на
кілька етапів.

\textbf{КРОК} 1. (Перетворення диференціального рівняння в похідних у
два звичайних диференціальних рівняння)

Підставляючи \(u(x, t)\) = X(x)\$\(T(t)\)у рівняння (7.1),
отримуємо\(XT' = \alpha^2 X''T\)Поділивши інші частини цього рівняння
на\$\textbackslash alpha\^{}2 XT отримуємо

{[} \textbackslash frac\{T'\}\{\textbackslash alpha\^{}2 T\} =
\textbackslash frac\{X'\,'\}\{X\} . {]}

Оскільки ліва частина останнього рівняння залежить лише від часу, а
права --- лише від x (і оскільки x і t незалежні), обидві сторони
рівняння мають бути константами. Прирівнюючи обидві частини до однієї
сталої \(\mu\) , отримуємо

(7.2)

{[} T' - \textbackslash mu \textbackslash alpha\^{}\{2\} T = 0, X'\,' -
\textbackslash mu X = 0. {]}

Процес розщеплення змінних завершено.

КРОК 2. (Визначення константи відокремлення)

По-перше, константа і не повинна бути додатною, інакше \(^{1}\) )
функція \(T\left( t\right)\) буде зростати експоненційно, що призведе до
нескінченного підвищення температури u=XT. Очевидно, що таке рішення не
має фізичного значення.

По-друге, припустимо, що \(\mu = 0\) . Тоді

{[} X'\,' = 0 {]}

і, отже.

{[} X(x) = A + \textbackslash hat\{B\}x. {]}

Однак, оскільки граничні умови вимагають, що

{[} u(0, t) = X(0) T(t) = 0, {]}

\(u(1, t) = (hX(1) + X'(1)) T(t) = 0,\) Ми одразу дійшли висновку, що

{[} X(0) = 0 {]}

\(\Rightarrow A = 0\) \(X'(1) + hX(1) = 0\) \(\Rightarrow B = 0\) .

Це означає, що \(u(x, t) \equiv 0\) . Іншими словами, лише нульове
розв'язання відповідає нульовій константі розділення. Ми відкинемо його,
оскільки нас цікавлять ненульові розв'язки.

Якщо, нарешті, \(\mu < 0\) , можна поставити \(\mu = -\lambda^2\) і
переписати рівняння (7.2) у вигляді

{[} T' + \textbackslash lambda\^{}2 \textbackslash alpha\^{}2 T = 0, {]}

{[} X' + \textbackslash lambda\^{}2 X = 0. {]}

1) Розв'язок \(X(x)\) відповідної задачі з граничним значенням буде лише
тривіальним, \(X(x)\) = 0. --- Ед. Ед.

Загальні розв'язки цих рівнянь легко знайти:

{[} T(t) = Ae\^{}\{-(\textbackslash lambda \textbackslash alpha)\^{}2
t\}, {]}

{[} X(x) = B\textbackslash sin(\textbackslash lambda x) +
C\textbackslash cos(\textbackslash lambda x). {]}

Отже, будь-яка функція провідника

(7.3)

{[} u(x, t) = e\^{}\{-(\textbackslash lambda \textbackslash alpha)\^{}2
t\} {[}A \textbackslash sin(\textbackslash lambda x) + B
\textbackslash cos(\textbackslash lambda x){]} {]}

при будь-яких значеннях \(\lambda\) A і B задовольняє рівняння (читач
може перевірити це самостійно простою заміною).

Тепер нам потрібно з'ясувати, які з цих функцій задовольняють граничним
умовам

(7.4)

{[} u(0, t) = 0, u\_x(1, t) + hu(1, t) = 0. {]}

Підстановка розв'язків виду (7.3) у граничні умови (7.4) дозволяє знайти
відношення, яким мають задовольняти значення \(\lambda\) , \(\Lambda\)
та B:

{[} Be\textsuperscript{\{-(\textbackslash lambda\textbackslash alpha)}2
t\} = 0 \textbackslash Rightarrow B = 0, {]}

\(A\lambda e^{-(\lambda\alpha)^2 t} \cos \lambda + hAe^{-(\lambda\alpha)^2 t} \sin \lambda = 0.\)
Після простих алгебранхічних перетворень останнє рівняння зводиться до
вигляду

{[} tg \textbackslash lambda =
-\textbackslash lambda/\textbackslash hbar {]}

.

Іншими словами, щоб визначити \(\lambda\) , потрібно знайти точки
перетину графів функцій tg \(\lambda\) і --- \(\lambda/h\) (рисунок
7.3).

Числові значення значень \(\lambda_1\) , \(\lambda_2\) , \ldots{} можна
знайти за допомогою комп'ютера, якщо задано значення h Р и с

\pandocbounded{\includegraphics[keepaspectratio]{images/_page_52_Figure_15.jpeg}}

. 7.3. Перетин графів функцій \(tg(\lambda)\) і \(-\lambda/h\) .

викликати власні значення задачі граничного значення

(7.5)

{[} X'\,' + \textbackslash lambda\^{}\{2\}X = 0, {]}

{[} \textbackslash begin\{cases\} X(0) = 0 \textbackslash{} X'(1) +
hX(1) = 0. \textbackslash end\{cases\} {]}

Іншими словами, з цими значеннями параметра \textbackslash lambda
існують нетривіальні розв'язки задачі (7.5). Власні значення задачі
(7.5) визначаються як корені рівняння \textbackslash log
\textbackslash lambda = -\textbackslash lambda/h . Їхні числові значення
знаходилися на комп'ютері для випадку h = 1. Таблиця 7.1 показує
результати розрахунків перших кількох власних значень.

\begin{verbatim}
| n | \lambda_B | rt . | \lambda_n |
\end{verbatim}

\textbar-------------\textbar----------------------\textbar--------\textbar----------------\textbar{}

123 \textbar{} 2,024,917,98 \textbar{} 45 \textbar{} 11,0814,20
\textbar{}

\begin{verbatim}
Розв'язки задачі (7.5), що відповідають власним значенням \lambda_n , називаються власними функціями X_n(x) . Для розглянутого завдання власні функції виглядають так:
\end{verbatim}

{[} X\_n(x) = \textbackslash sin(\textbackslash lambda\_n x) {]}

і зображені на рис. 7.4.

\pandocbounded{\includegraphics[keepaspectratio]{images/_page_53_Figure_10.jpeg}}

Рис. 7.4. Власні функції \(X_n\)(x) завдання (7.5) при h=1; Усі власні
функції задовольняють граничні умови.

КРОК 3, (Пошук фундаментальних рішень)

Тепер у нас є нескінченний набір функцій (фундаментальних розв'язків) у
нашому розпорядженні

{[} u\_n(x, t) = X\_n(x) T\_n(t) = e\^{}\{-(\textbackslash lambda\_n
x)\^{}2 t\} \textbackslash sin(\textbackslash lambda\_n x), {]}

кожна з яких задовольняє рівняння та граничні умови. Залишився ще один
крок --- загалом

{[} u(x, t) = \textbackslash sum\_\{n=1\}\^{}\{\textbackslash infty\}
a\_n X\_n(x) T\_n(t) = {]}

{[} = \textbackslash sum\_\{n=1\}\^{}\{\textbackslash infty\} a\_n
e\^{}\{-(\textbackslash lambda\_n c t\^{}2 t)\}
\textbackslash sin(\textbackslash lambda\_n x) {]}

Вибрати коефіцієнти \(a_n\) так, щоб функція \(u(x, t)\) задовольняла
початкову умову \(u(x, 0)\) = x (у цьому випадку функція \(u(x, t)\)
задовольнятиме рівняння та граничні умови, оскільки, по-перше, кожен
член суми їх задовольняє, а по-друге, і рівняння, і граничні умови є
лінійними та однорідними).

Отже, ймовірність \(a_n\) слід визначати з цієї умови

(7 6)

{[} u(x, 0) = x =
\textbackslash sum\_\{n=1\}\^{}\{\textbackslash infty\} a\_n
\textbackslash sin(\textbackslash lambda\_n x). {]}

Ми маємо зробити останній крок.

КРОК 4. (Розклад початкової умови на ряд за власними функціями)

Щоб знайти коефіцієнти \(a_n\) у розкладі (7,6) на власні функції,
помножте обидві частини цього співвідношення на \(\sin(\lambda_m x)\) і
інтегруйте від 0 до 1. У результаті ми отримуємо

{[} \textbackslash int\_\{0\}\^{}\{2\} \textbackslash xi
\textbackslash sin(\textbackslash lambda\_\{m\} \textbackslash xi)
d\textbackslash xi =
\textbackslash sum\_\{n=1\}\^{}\{\textbackslash infty\} a\_\{n\}
\textbackslash int\_\{0\}\^{}\{1\}
\textbackslash sin(\textbackslash lambda\_\{n\} \textbackslash xi)
\textbackslash sin(\textbackslash lambda\_\{m\} \textbackslash xi)
d\textbackslash xi = {]}

{[} = a\_\{m\} \textbackslash int\_\{0\}\^{}\{1\}
\textbackslash sin\^{}\{2\}(\textbackslash lambda\_\{m\}
\textbackslash xi) d\textbackslash xi = {]}

{[} = a\_\{m\} \textbackslash left(
\textbackslash frac\{\textbackslash lambda\_\{m\} - \textbackslash sin
\textbackslash lambda\_\{m\} \textbackslash cos
\textbackslash lambda\_\{m\}\}\{2\textbackslash lambda\_\{m\}\}
\textbackslash right). {]}

Розв'язавши останнє відношення відносно \(a_m\) і замінивши m на n,
отримуємо бажаний результат:

(7.7)

{[} a\_n =
\textbackslash frac\{2\textbackslash lambda\_n\}\{(\textbackslash lambda\_n
- \textbackslash sin \textbackslash lambda\_n \textbackslash cos
\textbackslash lambda\_n)\}
\textbackslash int\_0\^{}\{\textbackslash infty\} \textbackslash xi
\textbackslash sin (\textbackslash lambda\_n \textbackslash xi)
d\textbackslash xi. {]}

Таблиця 7.2.

\section{\texorpdfstring{Шанси \(a_n\) у формулі
(7.8)}{Шанси a\_n у формулі (7.8)}}\label{ux448ux430ux43dux441ux438-a_n-ux443-ux444ux43eux440ux43cux443ux43bux456-7.8}

{\def\LTcaptype{none} % do not increment counter
\begin{longtable}[]{@{}llll@{}}
\toprule\noalign{}
n & an & n & \(a_n\) \\
\midrule\noalign{}
\endhead
\bottomrule\noalign{}
\endlastfoot
123 & 0,240,22-0,03 & 45 & -0.11-0.09  \\
\end{longtable}
}

Отже, розв'язок задачі (7.1) записується у вигляді

(7.8)

{[} u(x, t) = \textbackslash sum\_\{n=1\}\^{}\{\textbackslash infty\}
a\_n e\^{}\{-i\textbackslash lambda\_n \textbackslash alpha\_n\^{}\{2\}
t\} \textbackslash sin(\textbackslash lambda\_n x), {]}

де коефіцієнти \(a_n\) задаються формулами (7.7). Перші п'ять
коефіцієнтів \(a_n\) нашої задачі були знайдені чисельно. Їхні значення
наведені в Таблиці 1. 7.2.

\pandocbounded{\includegraphics[keepaspectratio]{images/_page_55_Figure_8.jpeg}}

Рис. 7.5. Графічне представлення функції (7.8); межа умовно
задовольняється у всі моменти часу.

Відповідно, перші три члени розв'язання змішаної задачі

(7.9)

{[} \textbackslash begin\{aligned\} u\_t \&= u\_\{xx\}, \& 0 \textless{}
x \textless{} 1, \& 0 \textless{} t \textless{} \textbackslash infty,
\textbackslash{} u\_t \&= 0, \& 0 \textless{} t \textless{}
\textbackslash infty, \textbackslash{} u\_x \&= 0, \& 0 \textless{} t
\textless{} \textbackslash infty, \textbackslash{} u\_x \&= 0, \& 0
\textless{} t \textless{} \textbackslash infty, \textbackslash{} u\_x
\&= 0, \& 0 \textless{} t \textless{} \textbackslash infty,
\textbackslash{} u\_x \&= 0, \& 0 \textless{} t \textless{}
\textbackslash infty, \textbackslash{} u\_x \&= 0, \& 0 \textless{} t
\textless{} \textbackslash infty, \textbackslash{} u\_x \&= 0, \& 0
\textless{} t \textless{} \textbackslash infty, \textbackslash{} u\_x
\&= 0, \& 0 \textless{} t \textless{} \textbackslash infty,
\textbackslash{} u\_x \&= 0, \& 0 \textless{} t \textless{}
\textbackslash infty, \textbackslash{} u\_x \&= 0, \& 0 \textless{} t
\textless{} \textbackslash infty, \textbackslash{} u\_x \&= 0, \& 0
\textless{} t \textless{} \textbackslash infty, \textbackslash{} u\_x
\&= 0, \& 0 \textless{} t \textless{} \textbackslash infty,
\textbackslash{} u\_x \&= 0, \& 0 \textless{} t \textless{}
\textbackslash infty, \textbackslash{} u\_x \&= 0, \& 0 \textless{} t
\textless{} \textbackslash infty, \textbackslash{} u\_x \&= 0, \& 0
\textless{} t \textless{} \textbackslash infty, \textbackslash{} u\_x
\&= 0, \& 0 \textless{} t \textless{} \textbackslash infty,
\textbackslash{} u\_x \&= 0, \& 0 \textless{} t \textless{}
\textbackslash infty, \textbackslash{} u\_x \&= 0, \& 0 \textless{} t
\textless{} \textbackslash infty, \textbackslash{} u\_x \&= 0, \& 0
\textless{} t \textless{} \textbackslash infty, \textbackslash{} u\_x
\&= 0, \& 0 \textless{} t \textless{} \textbackslash infty,
\textbackslash{} u\_x \&= 0, \& 0 \textless{} t \textless{}
\textbackslash infty, \textbackslash{} u\_x \&= 0, \& 0 \textless{} t
\textless{} \textbackslash infty, \textbackslash{} u\_x \&= 0, \& 0
\textless{} t \textless{} \textbackslash infty, \textbackslash{} u\_x
\&= 0, \& 0 \textless{} t \textless{} \textbackslash infty,
\textbackslash{} u\_x \&= 0, \& 0 \textless{} t \textless{}
\textbackslash infty, \textbackslash{} u\_x \&= 0, \& 0 \textless{} t
\textless{} \textbackslash infty, \textbackslash{} u\_x \&= 0, \& 0
\textless{} t \textless{} \textbackslash infty, \textbackslash{} u\_x
\&= 0, \& 0 \textless{} t \textless{} \textbackslash infty,
\textbackslash{} u\_x \&= 0, \& 0 \textless{} t \textless{}
\textbackslash infty, \textbackslash{} u\_x \&= 0, \& 0 \textless{} t
\textless{} \textbackslash infty, \textbackslash{} u\_x \&= 0, \& 0
\textless{} t \textless{} \textbackslash infty, \textbackslash{} u\_x
\&= 0, \& 0 \textless{} t \textless{} \textbackslash infty,
\textbackslash{} u\_x \&= 0, \& 0 \textless{} t \textless{}
\textbackslash infty, \textbackslash{} u\_x \&= 0, \& 0 \textless{} t
\textless{} \textbackslash infty, \textbackslash{} u\_x \&= 0, \& 0
\textless{} t \textless{} \textbackslash infty, \textbackslash{} u\_x
\&= 0, \& 0 \textless{} t \textless{} \textbackslash infty,
\textbackslash{} u\_x \&= 0, \& 0 \textless{} t \textless{}
\textbackslash infty, \textbackslash{} u\_x \&= 0, \& 0 \textless{} t
\textless{} \textbackslash infty, \textbackslash{} u\_x \&= 0, \& 0
\textless{} t \textless{} \textbackslash infty, \textbackslash{} u\_x
\&= 0, \& 0 \textless{} t \textless{} \textbackslash infty,
\textbackslash{} u\_x \&= 0, \& 0 \textless{} t \textless{}
\textbackslash infty, \textbackslash{} u\_x \&= 0, \& 0 \textless{} t
\textless{} \textbackslash infty, \textbackslash{} u\_x \&= 0, \& 0
\textless{} t \textless{} \textbackslash infty, \textbackslash{} u\_x
\&= 0, \& 0 \textless{} t \textless{} \textbackslash infty,
\textbackslash{} u\_x \&= 0, \& 0 \textless{} t \textless{}
\textbackslash infty, \textbackslash{} u\_x \&= 0, \& 0 \textless{} t
\textless{} \textbackslash infty, \textbackslash{} u\_x \&= 0, \& 0
\textless{} t \textless{} \textbackslash infty, \textbackslash{} u\_x
\&= 0, \& 0 \textless{} t \textless{} \textbackslash infty,
\textbackslash{} u\_x \&= 0, \& 0 \textless{} t \textless{}
\textbackslash infty, \textbackslash{} u\_x \&= 0, \& 0 \textless{} t
\textless{} \textbackslash infty, \textbackslash{} u\_x \&= 0, \& 0
\textless{} t \textless{} \textbackslash infty, \textbackslash{} u\_x
\&= 0, \& 0 \textless{} t \textless{} \textbackslash infty,
\textbackslash{} u\_x \&= 0, \& 0 \textless{} t \textless{}
\textbackslash infty, \textbackslash{} u\_x \&= 0, \& 0 \textless{} t
\textless{} \textbackslash infty, \textbackslash{} u\_x \&= 0, \& 0
\textless{} t \textless{} \textbackslash infty, \textbackslash{} u\_x
\&= 0, \& 0 \textless{} t \textless{} \textbackslash infty,
\textbackslash{} u\_x \&= 0, \& 0 \textless{} t \textless{}
\textbackslash infty, \textbackslash{} u\_x \&= 0, \& 0 \textless{} t
\textless{} \textbackslash infty, \textbackslash{} u\_x \&= 0, \& 0
\textless{} t \textless{} \textbackslash infty, \textbackslash{} u\_x
\&= 0, \& 0 \textless{} t \textless{} \textbackslash infty,
\textbackslash{} u\_x \&= 0, \& 0 \textless{} t \textless{}
\textbackslash infty, \textbackslash{} u\_x \&= 0, \& 0 \textless{} t
\textless{} \textbackslash infty, \textbackslash{} u\_x \&= 0, \& 0
\textless{} t \textless{} \textbackslash infty, \textbackslash{} u\_x
\&= 0, \& 0 \textless{} t \textless{} \textbackslash infty,
\textbackslash{} u\_x \&= 0, \& 0 \textless{} t \textless{}
\textbackslash infty, \textbackslash{} u\_x \&= 0, \& 0 \textless{} t
\textless{} \textbackslash infty, \textbackslash{} u\_x \&= 0, \& 0
\textless{} t \textless{} \textbackslash infty, \textbackslash{} u\_x
\&= 0, \& 0 \textless{} t \textless{} \textbackslash infty,
\textbackslash{} u\_x \&= 0, \& 0 \textless{} t \textless{}
\textbackslash infty, \textbackslash{} u\_x \&= 0, \& 0 \textless{} t
\textless{} \textbackslash infty, \textbackslash{} u\_x \&= 0, \& 0
\textless{} t \textless{} \textbackslash infty, \textbackslash{} u\_x
\&= 0, \& 0 \textless{} t \textless{} \textbackslash infty,
\textbackslash{} u\_x \&= 0, \& 0 \textless{} t \textless{}
\textbackslash infty, \textbackslash{} u\_x \&= 0, \& 0 \textless{} t
\textless{} \textbackslash infty, \textbackslash{} u\_x \&= 0, \& 0
\textless{} t \textless{} \textbackslash infty, \textbackslash{} u\_x
\&= 0, \& 0 \textless{} t \textless{} \textbackslash infty,
\textbackslash{} u\_x \&= 0, \& 0 \textless{} t \textless{}
\textbackslash infty, \textbackslash{} u\_x \&= 0, \& 0 \textless{} t
\textless{} \textbackslash infty, \textbackslash{} u\_x \&= 0, \& 0
\textless{} t \textless{} \textbackslash infty, \textbackslash{} u\_x
\&= 0, \& 0 \textless{} t \textbackslash end\{aligned\} {]}

мають вигляд

{[} u(x, t) = = 0.24e\^{}\{-4t\} \textbackslash sin(2.02x) +
0.22e\^{}\{-24t\} \textbackslash sin(4.9x) - 0.03e\^{}\{-48.3t\}
\textbackslash sin(7.98x) + \textbackslash dots {]}

Фіг. На рисунку 7-5 показані графіки цього розв'язку для кількох точок
часу. Рекомендуємо читачеві відповісти на такі питання: Чи відповідає
тип рішення, яке ви отримуєте, з вашими інтуїтивними ідеями? Чи
виконуються граничні умови?

\section{ЗАУВАЖЕННЯ}\label{ux437ux430ux443ux432ux430ux436ux435ux43dux43dux44f-3}

Задача власних значень (7.5) є особливим випадком більш загальної
задачі, яку зазвичай називають проблемою Штурма-Ліувіля:

(7.10)

{[} \textbackslash begin\{aligned\} \& (\textbackslash text\{ОДУ\})
\textbackslash quad {[}p(x)y'{]}' - q(x)y + \textbackslash lambda r(x)y
= 0, \textbackslash quad 0 \textless{} x \textless{} 1, \textbackslash{}
\& (\textbackslash alpha\_1 y(0) + \textbackslash beta\_1 y'(0) = 0,
\textbackslash{} \& (\textbackslash alpha\_2 y(1) +
\textbackslash beta\_1 y'(1) = 0. \textbackslash end\{aligned\} {]}

Якщо розв'язати задачу з однорідними граничними умовами методом
розділення змінних, то функції \(X(x)\) завжди знаходяться як розв'язки
деякої задачі Штурма--Ліувіля. Очевидно, що задача власних значень (7.5)
є особливим випадком задачі (7.10).

Перелічимо основні властивості задачі (7.10), які були встановлені
Штурмом і Ліувілем за певних природних обмежень.

функції \(p(x)\), \(q(x)\) та \(r(x)\).

\begin{enumerate}
\def\labelenumi{\arabic{enumi}.}
\tightlist
\item
  Існує нескінченна послідовність власних значень, що задовольняють
  нерівності
\end{enumerate}

{[} \textbackslash lambda\_1 \textless{} \textbackslash lambda\_2
\textless{} \textbackslash lambda\_3 \textless{} \textbackslash ldots
\textless{} \textbackslash lambda\_n \textless{} \textbackslash ldots
\textbackslash to \textbackslash infty {]}

.

\begin{enumerate}
\def\labelenumi{\arabic{enumi}.}
\setcounter{enumi}{1}
\item
  Кожне власне значення \(\lambda_n\) відповідає одній (точній до
  константного множника) ненульовій власній функції \(y_n(x)\) .
\item
  Якщо \(y_n(x)\) і \(y_m(x)\) дві різні власні функції, що відповідають
  різним власним значенням \(\lambda_n\) і \(\lambda_m\) , то вони є
  ортогональними з вагою \(r(x)\) на сегменті {[}0, 1{]}, тобто вони
  задовольняють це відношення
\end{enumerate}

{[} \textbackslash int\_\{0\}\^{}\{1\} r(x) y\_\{n\}(x) y\_\{m\}(x) dx =
0. {]}

Більше деталей про проблему Штурма-Ліувіля можна знайти в літературі,
наведеній наприкінці книги.

1) Див., наприклад, Тіхоног А. Н., Самарський А. А. Рівняння
математичної фізики. --- Москва: Наука, 1972. --- ред. Ед.

\section{ЗАВДАННЯ}\label{ux437ux430ux432ux434ux430ux43dux43dux44f-4}

\begin{enumerate}
\def\labelenumi{\arabic{enumi}.}
\tightlist
\item
  Розв'язати наступну задачу для рівняння теплопровідності методом
  розкладу
\end{enumerate}

(PPP)

{[} u\_1 = u\_\{xx\}, \textbackslash quad 0 \textless{} x \textless{} 1,
\textbackslash quad 0 \textless{} t \textless{} \textbackslash infty,
{]}

(GI)

{[} \textbackslash begin\{cases\} u(0, t) = 0, \textbackslash{} u\_x(1,
t) = 0, \textbackslash{} u(x, 0) = x, \textbackslash quad 0
\textbackslash le x \textbackslash le 1. \textbackslash end\{cases\} {]}

Чи відповідає рішення, яке ви знаходите, з інтонативними
представленнями? Що таке стаціонарне рішення?

\begin{enumerate}
\def\labelenumi{\arabic{enumi}.}
\setcounter{enumi}{1}
\tightlist
\item
  Знайдіть власні цінності та функції проблеми Ассолта-Ліувіля
\end{enumerate}

(ODY)

{[} X' + \textbackslash lambda X = 0, \textbackslash quad 0 \textless{}
x \textless{} 1, {]}

(FY)

{[} \textbackslash begin\{cases\} X(0) = 0, \textbackslash{} X'(1) = 0.
\textbackslash end\{cases\} {]}

Задайте специфічні вирази для функцій \(p(x)\), \(q(x)\) та \(r(x)\) із
загальної задачі Штурма--Ліувіля для цього випадку.

\begin{enumerate}
\def\labelenumi{\arabic{enumi}.}
\setcounter{enumi}{2}
\tightlist
\item
  Розв'язати задачу теплопровідності для планки з термоізольованими
  межами
\end{enumerate}

(PPP)

{[} u\_t = u\_\{xx\}, \textbackslash quad 0 \textless{} x \textless{} 1,
\textbackslash quad 0 \textless{} t \textless{} \textbackslash infty,
{]}

(GI)

{[} \textbackslash begin\{cases\} u\_x (0, t) = 0, \textbackslash{} u\_x
(1, t) = 0, \textbackslash end\{cases\} \textbackslash quad 0
\textless{} t \textless{} \textbackslash infty, {]}

(ГУ) \(u(x, 0) = x, \quad 0 \le x \le 1.\) Чи відповідає рішення, яке ви
отримуєте, з вашою інтерпретацією цієї проблеми? Що таке стаціонарне
рішення? Яке його фізичне значення?

\begin{enumerate}
\def\labelenumi{\arabic{enumi}.}
\setcounter{enumi}{3}
\tightlist
\item
  Знайдіть власні цінності та функції завдань
\end{enumerate}

(ОДА)

{[} X'\,' + \textbackslash lambda X = 0, \textbackslash quad 0
\textless{} x \textless{} 1, {]}

(GI)

{[} \textbackslash begin\{cases\} X'(0) = 0, \textbackslash{} X'(1) = 0.
\textbackslash end\{cases\} {]}

\#\#\# ПЕРЕТВОРЕННЯ СКЛАДНИХ РІВНЯНЬ У ПРОСТУ ФОРМУ

МЕТА ЛЕКЦІЇ Показати, як PPP для функції \(u(x, t)\) можна перетворити
на простіший PPP для нової функції \(w(x, t)\). Вибір правильної
трансформації зазвичай базується на ненавмисному намірі. У нашій лекції

{[} u\_t = \textbackslash alpha\^{}2 u\_\{xx\} - \textbackslash beta u,
{]}

{[} u\_t = \textbackslash alpha\^{}2 u\_\{xx\} - v u\_x,
\textbackslash quad (v = \textbackslash text\{const\}) {]}

переходимо до найпростішого рівняння теплопровідності

{[} w\_1 = \textbackslash alpha\^{}2 w\_\{xx\} {]}

внаслідок відповідних перетворень

{[} u(x, t) = e\^{}\{-\textbackslash beta t\} w(x, t) {]}

{[} u(x, t) = e\^{}\{\textbackslash epsilon t (1 - \textbackslash delta
t - 2) 2\textbackslash alpha\^{}2\} w(x, t). {]}

Розв'язавши просте рівняння теплопровідності за функцією \(\omega(x,t)\)
, ми можемо знайти розв'язки початкових рівнянь, перетворивши функцію
\(\omega(x,t)\) відповідно до цих формул (звісно, ці перетворення слід
застосовувати і до початкових і граничних умов).

Після прочитання двох останніх лекцій читач може створити враження, що
метод розділення змінних застосовується \emph{лише} до рівняння

{[} u\_1 = \textbackslash alpha\^{}2 u\_\{xx\} {]}

Насправді рівняння теплопотоку є найпростішим, але не єдиним
параболічним рівнянням, яке можна розв'язати методом розділення змінних.
Як уже згадувалося, якщо рівняння є лінійним і однорідним, його можна
розв'язати, поділивши змінні. Наприклад, двовимірна задача
теплопровідності в колі описується рівнянням

{[} u\_t = \textbackslash alpha\^{}2 \textbackslash left{[} u\_\{rr\} +
\textbackslash frac\{1\}\{r\} u\_r + \textbackslash frac\{1\}\{r\^{}2\}
u\_\{\textbackslash theta\textbackslash theta\} \textbackslash right{]},
{]}

І хоча це рівняння з \emph{змінними коефіцієнтами}, його можна звести до
трьох звичайних диференціальних рівнянь.

У цій лекції ми покажемо читачеві, що іноді, замість того, щоб
безпосередньо розв'язувати диференціальне рівняння в похідних, краще
перетворити його на простіше рівняння, яке можна розв'язати методом
розділення змінних або іншим методом. Наступний приклад ілюструє цей
підхід.

\section{Перетворення проблеми з бічним поверхневим обміном у проблему з
термічно ізольованою бічною
поверхнею}\label{ux43fux435ux440ux435ux442ux432ux43eux440ux435ux43dux43dux44f-ux43fux440ux43eux431ux43bux435ux43cux438-ux437-ux431ux456ux447ux43dux438ux43c-ux43fux43eux432ux435ux440ux445ux43dux435ux432ux438ux43c-ux43eux431ux43cux456ux43dux43eux43c-ux443-ux43fux440ux43eux431ux43bux435ux43cux443-ux437-ux442ux435ux440ux43cux456ux447ux43dux43e-ux456ux437ux43eux43bux44cux43eux432ux430ux43dux43eux44e-ux431ux456ux447ux43dux43eux44e-ux43fux43eux432ux435ux440ux445ux43dux435ux44e}

Розглянемо наступне завдання:

(8.1)

{[} u\_\{t\} = \textbackslash alpha\^{}\{2\}u\_\{xx\} -
\textbackslash beta u, \textbackslash quad 0 \textless{} x \textless{}
1, \textbackslash quad 0 \textless{} t \textless{} \textbackslash infty,
{]}

{[} \textbackslash begin\{cases\} u(0, t) = 0, \textbackslash{} u(1, t)
= 0, \textbackslash{} u(x, 0) = \textbackslash varphi(x),
\textbackslash quad 0 \textbackslash le x \textbackslash le 1,
\textbackslash end\{cases\} {]}

(8.1)

{[} (FY) \textbackslash quad {]}

{[} \textbackslash begin\{cases\} u(0, t) = 0, \textbackslash{} u(x, 0)
= \textbackslash varphi(x), \textbackslash quad 0 \textbackslash le x
\textbackslash le 1, \textbackslash end\{cases\} {]}

де термін \(-\beta u\) описує потік тепла через бічну поверхню стрижня
(рис. 8.1).

\pandocbounded{\includegraphics[keepaspectratio]{images/_page_59_Figure_7.jpeg}}

Рис. 81. Задача теплопровідності, описана рівнянням
\(u_t = \alpha^t u_{\lambda\lambda} - \beta u\) : \(\alpha\) --- це
втрата гепле через поверхню Боха; \(\delta\) -- дифузія вздовж стрижня.

Наша мета --- ввести нову температурну \(w\left(x,\ t\right)\) замість
\(u\left(x,\ t\right)\) , щоб рівняння для \(w\left(x,\ t\right)\) стало
простішим за початкове рівняння

{[} u\_t = \textbackslash alpha\^{}2 u\_\{xx\} - \textbackslash beta u
{]}

.

Це досить загальний метод розв'язання диференціальних рівнянь у
частинних похідних, але вибір відповідного перетворення базується на
інтуїтивних уявленнях про поведінку розв'язків початкового рівняння.
Наприклад, у задачі (8.1) температура \(u(x, t)\) у кожній точці \(x_0\)
змінюється внаслідок двох факторів:

\begin{itemize}
\tightlist
\item
  \begin{enumerate}
  \def\labelenumi{\arabic{enumi})}
  \tightlist
  \item
    дифузія тепла вздовж стрижня (через компонентний \(\alpha^2 u_{xx}\)
    ),
  \end{enumerate}
\item
  \begin{enumerate}
  \def\labelenumi{\arabic{enumi})}
  \setcounter{enumi}{1}
  \tightlist
  \item
    теплопередача через бічну поверхню стрижня (через член RI).
  \end{enumerate}
\end{itemize}

Найважливішим елементом нашого міркування є наступне: якщо не було
дифузії вздовж стрижня ( \(\alpha = 0\) ),

тоді температура в кожній точці стрижня експоненційно знижується до нуля
згідно із законом

{[} u(x\_0 t) = u(x\_0, 0) e\^{}\{-\textbackslash beta t\} {]}

.

Виходячи з цього спостереження, запитаємо, чи можна розв'язок проблеми
(8.1) представити як утворення двох факторів

{[} (8.2) u(x, t) = e\^{}\{-\textbackslash beta t\}w(x, t) {]}

або:

\begin{verbatim}
\begin{pmatrix}
\end{verbatim}

Температура елемента з \textbackslash{} Неізольовані боком Поверхня

де \(w(x, t)\) описує розподіл температури, що виникає лише через
дифузію. Давайте подивимось, що станеться, якщо підставити цей вираз для
температури у задачу (8.1). Після виконання відповідних розрахунків
(читач може зробити їх сам), ми отримаємо

\begin{verbatim}
Нова задача майже точно така ж, як і оригінальна, але більше не має додатку 6u. Тому замість розв'язання задачі (8.1) ми розв'язуємо перетворену задачу (8.3), а потім множимо отриманий розв'язок w(x, t) на e^{-\beta t} . У нашому випадку розв'язок задачі (8.3) вже було отримано раніше методом розділення змінних, тому ми можемо одразу записати
\end{verbatim}

(8.4)

{[} w(x, t) = \textbackslash sum\_\{n=1\}\^{}\{\textbackslash infty\}
a\_n
e\textsuperscript{\{-(n\textbackslash pi\textbackslash alpha)}\{\textbackslash frac\{n\}\{2\}\}t\}
\textbackslash sin(n\textbackslash pi x), {]}

{[} a\_n = 2 \textbackslash int\_0\^{}\{\textbackslash infty\}
\textbackslash phi(\textbackslash xi)
\textbackslash sin(n\textbackslash pi \textbackslash xi)
d\textbackslash xi, {]}

і, отже, розв'язок початкової задачі (8.1) має вигляд

{[} u(x, t) = e\^{}\{-\textbackslash beta t\}w(x, t). {]}

Нижче, у розділі «Нотатки», наведено ще один приклад проблеми,
розв'язаної цим методом.

\section{ПРИМІТКИ}\label{ux43fux440ux438ux43cux456ux442ux43aux438-1}

\begin{enumerate}
\def\labelenumi{\arabic{enumi}.}
\tightlist
\item
  Вирішити проблему
\end{enumerate}

{[} u\_t = u\_\{xx\} - u, \textbackslash quad 0 \textless{} x
\textless{} 1, \textbackslash quad 0 \textless{} t \textless{}
\textbackslash infty, {]}

{[} \textbackslash begin\{cases\} u(0, t) = 0, \textbackslash{} u(1, t)
= 0, \textbackslash{} u(x, 0) = \textbackslash sin(\textbackslash pi x)
+ 0.5\textbackslash sin(3\textbackslash pi x)
\textbackslash end\{cases\} {]}

За методом, описаним вище, ми

\begin{itemize}
\tightlist
\item
  \begin{enumerate}
  \def\labelenumi{\alph{enumi})}
  \tightlist
  \item
    тимчасово відкинути конвективний член u у початковому рівнянні.
  \end{enumerate}
\item
  \begin{enumerate}
  \def\labelenumi{\alph{enumi})}
  \setcounter{enumi}{1}
  \tightlist
  \item
    записати розв'язок змішаної задачі без конвективного члена
  \end{enumerate}
\end{itemize}

{[} u(x, t) = e\textsuperscript{\{-\textbackslash pi}2 t\}
\textbackslash sin(\textbackslash pi x) +
0.5e\textsuperscript{\{-(3\textbackslash pi)}2 t\}
\textbackslash sin(3\textbackslash pi x), {]}

\begin{enumerate}
\def\labelenumi{\alph{enumi})}
\setcounter{enumi}{2}
\tightlist
\item
  помножимо це розв'язання на конвективний фактор
  \(e^{-\beta t}=e^{-t}\)отримаємо остаточне розв'язання початкової
  задачі
\end{enumerate}

{[} u(x, t) = e\^{}\{-t\} \textbackslash left{[}
e\textsuperscript{\{-\textbackslash pi}2 t\}
\textbackslash sin(\textbackslash pi x) +
0.5e\textsuperscript{\{-(3\textbackslash pi)}2 t\}
\textbackslash sin(3\textbackslash pi x) \textbackslash right{]}. {]}

\begin{enumerate}
\def\labelenumi{\arabic{enumi}.}
\setcounter{enumi}{1}
\tightlist
\item
  Рівняння конвективної дифузії
\end{enumerate}

{[} u\_t = \textbackslash alpha\^{}2 u\_\{xx\} - v u\_x {]}

(коефіцієнт і є константою) також можна перетворити у вигляд

{[} w\_t = \textbackslash alpha\^{}2 w\_\{xx\}. {]}

Відповідне перетворення має вигляд

{[} u(x, t) =
e\textsuperscript{\{v(x-vt/2)/2\textbackslash alpha}2\}w(x, t). {]}

Це перетворення підкреслює експоненціальний фактор, який повністю
зумовлений рухом середовища. Зверніть увагу, що цей множник описує
хвилю, що поширюється зліва направо зі швидкістю
\textbackslash upsilon/2 . Читач матиме можливість скористатися цією
трансформацією.

\section{ЗАВДАННЯ}\label{ux437ux430ux432ux434ux430ux43dux43dux44f-5}

\begin{enumerate}
\def\labelenumi{\arabic{enumi}.}
\item
  Розв'язати проблему конвективної дифузії

  \textbackslash begin\{array\}\{lll\}
  (\textbackslash mathrm\{Y\}\textbackslash mathrm\{H\}\textbackslash Pi)
  \& u\_t = u\_\{xx\} - u\_x, \& 0 \textless{} x \textless{} 1, \& 0
  \textless{} t \textless{} \textbackslash infty, \textbackslash{}
  (\textbackslash mathrm\{\textbackslash Gamma\}\textbackslash mathrm\{Y\})
  \& \textbackslash left\{ \textbackslash begin\{array\}\{ll\} u ~(0,
  ~t) = 0, \& 0 \textless{} t \textless{} \textbackslash infty,
  \textbackslash{} u ~(1, ~t) = 0, \& 0 \textless{} t \textless{}
  \textbackslash infty, \textbackslash{} u ~(x, ~0) = e\^{}\{x/2\}, \& 0
  \textbackslash leqslant x \textbackslash leqslant 1,
  \textbackslash end\{array\}

  \textbackslash Rightarrow. \textbackslash{}
  \textbackslash end\{aligned\}
\end{enumerate}

методом преобразования к более простому уравнению. На що похоже
розв'язок? Можна дать следующую интерпретацию цієї задачи:
\(u(x, t)\)--- концентрация некоторого вещества в движущейся среде
(среда движется слева направо со скоростью v=1), начальное распределение
копцентрации задастся функцией \(e^{x/2}\) , на границах концентрация
поддерживается нулевой (можна считать, що там стоят фильтры).
Соответствует ли ваще розв'язок такой интерпретации?

\begin{enumerate}
\def\labelenumi{\arabic{enumi}.}
\setcounter{enumi}{1}
\item
  Решите залачу

  \textbackslash begin\{aligned\} \textbackslash text\{(YUII)\} \& u\_t
  = u\_\{xx\} - u + x, \& 0 \textless{} x \textless{} 1, \& 0
  \textless{} t \textless{} \textbackslash infty, \textbackslash{}
  \textbackslash text\{(FY)\} \& \textbackslash left\{
  \textbackslash begin\{aligned\} u ~(0, ~t) = 0, \textbackslash{} u
  ~(1, ~t) = 1, \textbackslash{} u ~(x, ~0) = 0, \& 0 \textbackslash le
  x \textbackslash le 1, \textbackslash end\{aligned\}
  \textbackslash end\{aligned\}

  \textbackslash Rightarrow. \textbackslash{}
  \textbackslash end\{aligned\}
\end{enumerate}

\begin{itemize}
\item
  а) преобразовав исоднородные граничные условия к однородным;
\item
  б) преобразовав исходное уравнение к повому уравнению, не содержащему
  члена -u;
\item
  в) решив получившуюся в результате эт іх преобразований задачу.
\item
  \begin{enumerate}
  \def\labelenumi{\arabic{enumi}.}
  \setcounter{enumi}{2}
  \tightlist
  \item
    Решите задачу
  \end{enumerate}

  \textbackslash begin\{array\}\{lll\} (\textbackslash text\{Y 4II\}) \&
  u\_t = u\_\{\textbackslash text\{vx\}\} - u, \& 0 \textless{} x
  \textless{} 1, \& 0 \textless{} t \textless{} \textbackslash infty,
  \textbackslash{} (\textbackslash text\{\textbackslash Gamma Y\}) \&
  \textbackslash begin\{cases\} u(0,,t) = 0, \textbackslash{} u(1,,t) =
  0, \textbackslash{} u(x,,0) = \textbackslash sin\{(\textbackslash pi
  x)\}, \& 0 \textbackslash leqslant x \textbackslash leqslant 1,
  \textbackslash end\{cases\}
\end{itemize}

применяя метод поділу змінних непосредственно к исходному уравнению и не
проводя никаких предварительных преобразований. Совпадает ли полученное
вами розв'язок с решением, полученным ранее після преобразования

{[} u(x, t) = e\^{}\{-t\}w(x, t) {]}

?

\section{РОЗВ'ЯЗОК НЕОДНОРОДНЫХ УЧП МЕТОДОМ РАЗЛОЖЕНИЯ ПО СОБСТВЕННЫМ
ФУНКЦИЯМ.}\label{ux440ux43eux437ux432ux44fux437ux43eux43a-ux43dux435ux43eux434ux43dux43eux440ux43eux434ux43dux44bux445-ux443ux447ux43f-ux43cux435ux442ux43eux434ux43eux43c-ux440ux430ux437ux43bux43eux436ux435ux43dux438ux44f-ux43fux43e-ux441ux43eux431ux441ux442ux432ux435ux43dux43dux44bux43c-ux444ux443ux43dux43aux446ux438ux44fux43c.}

ЦЕЛЬ ЛЕКЦИИ: Научиться решать смешанные задачи вида

\begin{verbatim}
\begin{array}{lll} \text{(YUII)} & u_t = \alpha^2 u_{xx} + f(x,\ t), & 0 < x < 1, & 0 <t < \infty, \\ \text{(FY)} & \left\{ \begin{array}{ll} \alpha_1 u_x \left( 0,\ t \right) + \beta_1 u \left( 0,\ t \right) = 0, \\ \alpha_2 u_x \left( 1,\ t \right) + \beta_2 u \left( 1,\ t \right) = 0, \\ u \left( x,\ 0 \right) = q \left( x \right), & 0 \leqslant x \leqslant 1, \end{array}

\right. \\ \end{aligned}
\end{verbatim}

Розв'язок гетерогенного UBP можна знайти у вигляді серії

{[} u\textbackslash left( x,~t\textbackslash right)
=\textbackslash sum\_\{n=1\}\^{}\{\textbackslash infty\}T\_\{n\}\textbackslash left(
t\textbackslash right) X\_\{n\}\textbackslash left(
x\textbackslash right) , {]}

де \(X_n\left(x\right)\) --- власні функції відповідної однорідної
задачі

\begin{verbatim}
\begin{aligned} (\text{YH}\Pi) & u_t = \alpha^2 u_{x_1}, & 0 < x < 1, & 0 < t < \infty, \\ (\text{\GammaV}) & \left\{ \begin{aligned} \alpha_1 u_x (0, \ t) + \beta_1 u \ (0, \ t) = 0, \\ \alpha_2 u_x \ (1, \ t) + \beta_2 u \ (1, \ t) = 0, \\ u \ (x, \ 0) = \phi \ (x), & 0 \leqslant x \leqslant 1, \end{aligned} \end{aligned}

\right. \\ \end{aligned}
\end{verbatim}

А функції \(T_n\)(t) визначаються як розв'язки деяких звичайних
диференціальних рівнянь.

На лекції 6 ми ознайомилися з тим, як перетворювати неоднорідні граничні
умови на однорідні. На жаль, рівняння залишається гетерогенним, і ми
стоїмо перед завданням

{[} (9.1) \textbackslash quad u\_\{t\} = \textbackslash alpha\^{}\{2\}
u\_\{xx\} + f(x, t), \textbackslash quad 0 \textless{} x \textless{} 1,
\textbackslash quad 0 \textless{} t \textless{} \textbackslash infty,
{]}

{[} (9.1) \textbackslash quad (\textbackslash Gamma Y)
\textbackslash quad {]}

{[} \textbackslash begin\{cases\} \textbackslash alpha\_\{1\}
u\_\{x\}(0, t) + \textbackslash beta\_\{1\} u(0, t) = 0,
\textbackslash{} \textbackslash alpha\_\{2\} u\_\{x\}(1, t) +
\textbackslash beta\_\{2\} u(1, t) = 0, \textbackslash{} u(x, 0) =
\textbackslash varphi(x), \textbackslash quad 0 \textbackslash le x
\textbackslash le 1. \textbackslash end\{cases\} {]}

Мета нашої лекції --- опанувати метод розв'язання таких задач. Цей метод
аналогічний методу варіації довільних констант.

теорія звичайних диференціальних рівнянь називається методом розкладу за
власними функціями.

Ідея методу надзвичайно проста. Відомо, що розв'язок задачі (9.1) при
\(f(x, t)\) = 0 (так звана відповідна однорідна задача) записується як
ряд

{[}
u\textbackslash left(x,~t\textbackslash right)=\textbackslash sum\_\{n=1\}\textsuperscript{\{\textbackslash infty\}a\_\{n\}e}\{-\textbackslash left(\textbackslash lambda\_\{n\}\textbackslash alpha\textbackslash right)\^{}\{2\}\}/X\_\{n\}\textbackslash left(x\textbackslash right),
{]}

де \(X_n\)(x) --- власні функції задачі Штурма-Ліувіля,

(9.2)

{[} X'\,' + \textbackslash lambda\^{}2 X = 0, \textbackslash{}
\textbackslash alpha\_1 X'(0) + \textbackslash beta\_1 X(0) = 0,
\textbackslash{} \textbackslash alpha\_2 X'(1) + \textbackslash beta\_2
X(1) = 0. {]}

Запитаймо себе, чи можна також шукати розв'язок гетерогенної проблеми
(9.1) у вигляді серії власних функцій, але більш загальної форми

{[} u(x, t) = \textbackslash sum\_\{n=1\}\^{}\{\textbackslash infty\}
T\_n(t) X\_n(x). {]}

Існують певні фізичні підстави для такого припущення. Насправді терміни
\(f(x,t)\) описують джерело тепла, розташоване всередині стрижня. Це
означає, що залежність рішення від часу не буде виражена показниками

{[}
e\textsuperscript{\{-(\textbackslash lambda\_n\textbackslash alpha)}2
t\} {]}

,

Це стосувалося однорідного завдання, але з деякими іншими функціями.

Щоб показати цей метод у дії і не ускладнювати подання деталями, ми
застосуємо його до розв'язання однієї дуже простої задачі.

\section{Після розв'язання методом демаркації за власними
функціями,}\label{ux43fux456ux441ux43bux44f-ux440ux43eux437ux432ux44fux437ux430ux43dux43dux44f-ux43cux435ux442ux43eux434ux43eux43c-ux434ux435ux43cux430ux440ux43aux430ux446ux456ux457-ux437ux430-ux432ux43bux430ux441ux43dux438ux43cux438-ux444ux443ux43dux43aux446ux456ux44fux43cux438}

Розглянемо гетерогенну проблему

(9.3) (PY)

{[} \textbackslash begin\{aligned\} u\_t \&= \textbackslash alpha\^{}2
u\_\{xx\} + f(x, t), \& 0 \textless{} x \textless{} 1, \& 0 \textless{}
t \textless{} \textbackslash infty, \textbackslash{} u(0, t) \&= 0, \& 0
\textless{} t \textless{} \textbackslash infty, \textbackslash{} u(1, t)
\&= 0, \& 0 \textless{} t \textless{} \textbackslash infty,
\textbackslash{} u(x, 0) \&= \textbackslash varphi(x), \& 0
\textbackslash leqslant x \textbackslash leqslant 1.
\textbackslash end\{aligned\} {]}

Процес розв'язання задачі подається у вигляді послідовності наступних
кроків.

КРОК 1. Основна ідея методу полягає в розкладанні густини джерела
\(f(x, t)\) на ряд власними функціями

{[} f(x, t) = f\_1(t) X\_1(x) + f\_2(t) X\_2(x) + \textbackslash dots +
f\_n(t) X\_n(x) + \textbackslash dots {]}

а також визначити реакції системи
\(u_n\)\textbackslash left(x,,t\textbackslash right) на вплив кожного
компонента f\_\(n\left(t\right)\)X\_\(n\left(x\right)\) . Якщо знайти
всі відповіді, рішення проблеми виглядатиме так:

{[} u(x, t) = \textbackslash sum\_\{n=1\}\^{}\{\textbackslash infty\}
u\_n(x, t). {]}

Головна складність цього методу --- розклад густини джерела на
компоненти \(f_n\)(t),\(X_n\)(x) . Виявляється, що \(X_n\)(x) факторами
в цій задачі є власними векторами системи Штурма--Ліувіля, які виникають
при розв'язанні відповідної однорідної задачі методом розділення
змінних, а саме

{[} \textbackslash begin\{array\}\{c\} u\_t = \textbackslash alpha\^{}2
u\_\{xx\} \textbackslash quad \textbackslash text\{(отметим, что \} f(x,
~t) = 0), \textbackslash{} u(0, ~t) = 0, \textbackslash{} u(1, ~t) = 0,
\textbackslash{} u(x, ~0) = \textbackslash varphi(x).
\textbackslash end\{array\} {]}

Завдання Штурма-Ліувіля в цьому випадку виглядає так:

{[} X'\,' + \textbackslash lambda\^{}2 X = 0, {]}

\(X(0)\) = 0, \(X(1)\) = 0, а його розв'язки є функціями

{[} X\_n(x) = \textbackslash sin(n\textbackslash pi x),
\textbackslash quad n = 1, 2, 3 \textbackslash dots {]}

Отже, розклад на густину джерела представлений як ряд

{[} (9.5) ~f(x,t) = f\_1(t) \textbackslash sin(\textbackslash pi x) +
f\_2(t) \textbackslash sin(2\textbackslash pi x) + \textbackslash dots +
f\_n(t) \textbackslash sin(n\textbackslash pi x) + \textbackslash dots
{]}

\textbackslash mathcal\{U\} нарешті, щоб знайти функції \(f_n\)(t) ,
помножте обидві сторони останнього відношення на
\textbackslash sin(m\textbackslash pi x) і інтегруйте від 0 до 1 за
координатою x.

У результаті ми отримуємо

{[} \textbackslash int\_\{0\}\^{}\{1\} f(x, t)
\textbackslash sin(m\textbackslash pi x) dx =
\textbackslash sum\_\{n=1\}\^{}\{\textbackslash infty\} f\_\{n\}(t)
\textbackslash int\_\{0\}\^{}\{1\} \textbackslash sin(m\textbackslash pi
x) \textbackslash sin(n\textbackslash pi x) dx = {]}

{[} = \textbackslash frac\{1\}\{2\} f\_\{m\}(t) {]}

або (після заміни t на p)

(9.6)

{[} f\_\{n\}(t) = 2 \textbackslash int\_\{0\}\^{}\{1\} f(x, t)
\textbackslash sin(n\textbackslash pi x) dx. {]}

Останнє співвідношення встановлює співвідношення між коефіцієнтами
\(f_n(t)\) і часткою джерела тепла \(f(x, t)\).

КРОК 2. (Пошук \(u_n(x, t) = X_n(x) T_n(t)\) реакції на вхідний вплив
\(f_n(t) X_n(x)\) .)

Замінимо густину джерела на його розклад у рядок

{[} f(x, t) = \textbackslash sum\_\{n=1\}\^{}\{\textbackslash infty\}
f\_n(t) \textbackslash sin(n\textbackslash pi x) {]}

і намагатися знайти індивідуальні відповіді, тобто функції,
\(T_{\mathbf{n}}(t)\) у

{[} u(x, t) = \textbackslash sum\_\{n=1\}\^{}\{\textbackslash infty\}
T\_n(t) \textbackslash sin(n\textbackslash pi x). {]}

Оскільки розв'язок задачі має вигляд

{[} u(x, t) = \textbackslash sum\_\{n=1\}\^{}\{\textbackslash infty\}
T\_n(t) \textbackslash sin(n\textbackslash pi x), {]}

потім підставимо цей вираз у початкову задачу

{[} u\_\{t\} = \textbackslash alpha\^{}\{2\} u\_\{xx\} +
\textbackslash sum\_\{n=1\}\^{}\{\textbackslash infty\} f\_\{n\}(t)
\textbackslash sin(n\textbackslash pi x), {]}

{[} u(0, t) = 0, {]}

{[} u(1, t) = 0, {]}

{[} u(x, 0) = \textbackslash varphi(x), {]}

У результаті ми отримаємо

{[} \textbackslash sum\_\{n=1\}\^{}\{\textbackslash infty\} T'\emph{n(t)
\textbackslash sin(n\textbackslash pi x) = -\textbackslash alpha\^{}2
\textbackslash sum}\{n=1\}\^{}\{\textbackslash infty\}
(n\textbackslash pi)\^{}2 T\_n(t) \textbackslash sin(n\textbackslash pi
x) + + \textbackslash sum\_\{n=1\}\^{}\{\textbackslash infty\} f\_n(t)
\textbackslash sin(\textbackslash pi n\textbackslash lambda), {]}

(9.7)

{[} \textbackslash sum\_\{n=1\}\^{}\{\textbackslash infty\} T\_n(t)
\textbackslash sin 0 = 0 {]}

(задовольняється ідентично),
\textbackslash sum\_\{n=1\}\^{}\{\textbackslash infty\} \(T_n\)(t)
\textbackslash sin (n\textbackslash pi) = 0 (задовольняється ідентично),
\textbackslash sum\_\{n=1\}\^{}\{\textbackslash infty\} \(T_n\)(0)
\textbackslash sin (n\textbackslash pi x) = \textbackslash phi(x) .

Переписування рівняння та початкової умови

(PPP)

{[} \textbackslash sum\_\{n=1\}\^{}\{\textbackslash infty\} {[}T'\_n +
(n\textbackslash pi\textbackslash alpha)\^{}2 T\_n - f\_n(t){]}
\textbackslash sin(n\textbackslash pi x) = 0, {]}

(НУ)

{[} \textbackslash sum\_\{n=1\}\^{}\{\textbackslash infty\} T\_n(0)
\textbackslash sin(n\textbackslash pi x) = \textbackslash varphi(x), {]}

можна побачити, що функції \(T_n\)(\textbackslash ell) є розв'язками
задачі Коші

(9.8)

{[} T\_\{n\} + (n\textbackslash pi\textbackslash alpha)\^{}\{2\}
T\_\{n\} = f\_\{n\}(t), {]}

{[} T\_\{n\}(0) = 2 \textbackslash int\_\{0\}\^{}\{1\}
\textbackslash varphi(\textbackslash xi)
\textbackslash sin(n\textbackslash pi\textbackslash xi)
d\textbackslash xi = a\_\{n\}. {]}

Усі ці TAC легко розв'язати, а їхні рішення записуються у вигляді

\begin{enumerate}
\def\labelenumi{(\arabic{enumi})}
\setcounter{enumi}{98}
\tightlist
\item
\end{enumerate}

{[} T\_n(t) = a\_n
e\textsuperscript{\{-(n\textbackslash pi\textbackslash alpha)}2 t\} +
\textbackslash int\_0\^{}t
e\textsuperscript{\{-(n\textbackslash pi\textbackslash alpha)}2
(t-\textbackslash tau)\} f\_n(\textbackslash tau) d\textbackslash tau,
{]}

З того місця, де легко отримати повне розв'язання проблеми (9.3):

(9.10)

{[} u(x, t) = \textbackslash sum\_\{n=1\}\^{}\{\textbackslash infty\}
T\_n(t) \textbackslash sin(n\textbackslash pi x) = {]}

{[} = \textbackslash sum\_\{n=1\}\^{}\{\textbackslash infty\}
\textbackslash left{[} a\_n
e\textsuperscript{\{-(n\textbackslash pi\textbackslash alpha)}n t\}
\textbackslash sin(n\textbackslash pi x) \textbackslash right{]} +
\textbackslash sum\_\{n=1\}\^{}\{\textbackslash infty\}
\textbackslash left{[} \textbackslash sin(n\textbackslash pi x)
\textbackslash int\_0\^{}t
e\textsuperscript{\{-(n\textbackslash pi\textbackslash alpha)}n
(t-\textbackslash tau)\} f\_n(\textbackslash tau) d\textbackslash tau
\textbackslash right{]} {]}

Перехідна частина (обумовлена стаціонарною частиною (умовлена початковою
умовою) правою частиною \(f(x, t)\)).

Розв'язок (9.10) показує, що температурна реакція складається з двох
частин: перша частина зумовлена початковими умовами, а друга ---
густиною джерела \(f(x,t)\). Зверніть увагу, що термін «стаціонарна
частина» не є найкращим описом другої частини температурної реакції,
оскільки він не означає, що ця частина не залежить від часу
(стаціонарний стан може бути періодичним, якщо \(f(x, t)\) є періодичною
функцією часу t).

Проблема повністю вирішена. Перш ніж перейти до наступної теми, давайте
покажемо, як використовувати описаний вище метод для розв'язання

Конкретне завдання.

70

\section{Розв'язання конкретної задачі методом розкладу за власними
функціями}\label{ux440ux43eux437ux432ux44fux437ux430ux43dux43dux44f-ux43aux43eux43dux43aux440ux435ux442ux43dux43eux457-ux437ux430ux434ux430ux447ux456-ux43cux435ux442ux43eux434ux43eux43c-ux440ux43eux437ux43aux43bux430ux434ux443-ux437ux430-ux432ux43bux430ux441ux43dux438ux43cux438-ux444ux443ux43dux43aux446ux456ux44fux43cux438}

Розглянемо просту задачу

(9.11)

{[} \textbackslash begin\{aligned\} u\_t \&= \textbackslash alpha\^{}2
u\_\{xx\} + \textbackslash sin\{(3\textbackslash pi x)\}, \& 0
\textless{} x \textless{} 1, \textbackslash{}
(\textbackslash mathbf\{FY\}) \& \textbackslash begin\{cases\} u(0, t)
\&= 0, \textbackslash{} u(1, t) \&= 0, \textbackslash end\{cases\} \& 0
\textless{} t \textless{} \textbackslash infty, \textbackslash{}
(\textbackslash mathbf\{HY\}) \& u(x, 0) \&=
\textbackslash sin\{(\textbackslash pi x)\}, \& 0 \textbackslash leq x
\textbackslash leq 1. \textbackslash end\{aligned\} {]}

Наша мета --- обчислити коефіцієнти \(T_n\)(t) при розкладі

{[}
u\textbackslash left(x,,t\textbackslash right)=\textbackslash sum\_\{n=1\}\^{}\{\textbackslash infty\},T\_\{n\}\textbackslash left(t\textbackslash right)\textbackslash sin\textbackslash left(n\textbackslash pi
x\textbackslash right) {]}

(Власні функції цієї задачі \(X_n\)(x) збігатися з власними функціями
попереднього). Якщо підставити цю декомпозицію на початкову задачу,
отримаємо наступну задачу Коші для визначення функцій \(T_n\)(t) :

{[} T\_n + (n\textbackslash pi\textbackslash alpha)\^{}2 T\_n =
\textbackslash int\_n (t) = 2 \textbackslash int\_0\^{}1
\textbackslash sin(3\textbackslash pi x)
\textbackslash sin(n\textbackslash pi x) dx = {]}

{[} \textbackslash begin\{cases\} 1 \&n = 3, \textbackslash{} 0 \& n
\textbackslash neq 3, \textbackslash end\{cases\} {]}

{[} T\_n(0) = 2 \textbackslash int\_0\^{}1
\textbackslash sin(n\textbackslash xi)
\textbackslash sin(n\textbackslash pi\textbackslash xi)
d\textbackslash xi = {]}

{[} \textbackslash begin\{cases\} 1 \&n = 1, \textbackslash{} 0 \& n
\textbackslash neq 1. \textbackslash end\{cases\} {]}

Перепишем эти уравнения последовательно для n=1, 2\ldots{}

{[} \textbackslash begin\{aligned\} (n=1) \& T\_1' +
(\textbackslash pi\textbackslash alpha)\^{}2 T\_1 = 0 \textbackslash{}
T\_1(0) = 1 \& \textbackslash Rightarrow T\_1(t) =
e\textsuperscript{\{-(\textbackslash pi\textbackslash alpha)}2 t\},
\textbackslash{} (n=2) \& T\_2' +
(2\textbackslash pi\textbackslash alpha)\^{}2 T\_2 = 0
\textbackslash T\_2~ (0) = 0 \& \textbackslash Rightarrow T\_2(t) = 0,
\textbackslash{} (n=3) \& T\_3' +
(3\textbackslash pi\textbackslash alpha)\^{}2 T\_3 = 4 \textbackslash{}
T\_3(0) = 0 \& \textbackslash Rightarrow T\_3(t) =
\textbackslash frac\{1\}\{(3\textbackslash pi\textbackslash alpha)\^{}2\}
\textbackslash left{[}1 -
e\textsuperscript{\{-(3~pi\textbackslash alpha)}2
t\}\textbackslash right{]}, \textbackslash{} (n \textbackslash geqslant
4) \& T\_n' + (n\textbackslash pi\textbackslash alpha)\^{}2 T\_n = 0
\textbackslash{} T\_n(0) = 0 \& \textbackslash Rightarrow T\_n(0) = 0.
\textbackslash end\{aligned\} {]}

Отже, розв'язок задачі можна записати у вигляді

{[} u(x, t) =
e\textsuperscript{\{-(\textbackslash pi\textbackslash alpha)}2 t\}
\textbackslash sin(\textbackslash pi x) +
\textbackslash frac\{1\}\{(3\textbackslash pi\textbackslash alpha)\^{}2\}
{[}1 - e\textsuperscript{\{-(3\textbackslash pi\textbackslash alpha)}2
t\}{]} \textbackslash sin(3\textbackslash pi x). {]}

(9.12)

\section{ЗАУВАЖЕННЯ}\label{ux437ux430ux443ux432ux430ux436ux435ux43dux43dux44f-4}

\begin{enumerate}
\def\labelenumi{\arabic{enumi}.}
\item
  Метод розкладу власних функцій є одним із найпотужніших методів
  розв'язання неоднорідних диференціальних рівнянь у похідних похідних.
  Пізніше, коли ми ознайомимося з інтегральними перетвореннями,
  побачимо, що існують й інші методи розв'язання задач такого типу.
\item
  Власні функції \(X_n(x)\) , за допомогою якої виконується розклад,
  залежать від задачі, яку потрібно розв'язати, точніше від рівнянь і
  граничних умов. Читача просять розв'язати задачу 4 наприкінці лекції,
  щоб перевірити, чи дійсно він навчився знаходити власні функції
  \(X_n(x)\) .
\item
  Якщо читач пам'ятає основні факти з теорії звичайних диференціальних
  рівнянь, то він знає, що розв'язки, що відповідають неоднорідним
  членам вигляду
\end{enumerate}

{[} P\_n(x) e\^{}\{\textbackslash alpha x\} {]}

{[} \textbackslash begin\{cases\}
\textbackslash sin\{(\textbackslash beta x)\}, \textbackslash{}
\textbackslash cos\{(\textbackslash beta x)\},
\textbackslash end\{cases\} {]}

можно найти методом неопределенных коэффициентов. То же самое и здесь.
Задача (9.11) может быть решена таким же способом.

\section{ЗАДАЧИ}\label{ux437ux430ux434ux430ux447ux438}

\section{1. Решение
задачи}\label{ux440ux435ux448ux435ux43dux438ux435-ux437ux430ux434ux430ux447ux438}

{[} \textbackslash begin\{array\}\{ll\} (\textbackslash text\{УЧП\}) \&
u\_t = u\_\{xx\} + \textbackslash sin\{(3\textbackslash pi x)\}, \& 0
\textless{} x \textless{} 1, \textbackslash{} (\textbackslash Gamma
\textbackslash text\{У\}) \& \textbackslash begin\{cases\} u ~(0, ~t) =
0, \textbackslash{} u ~(1, ~t) = 0, \textbackslash end\{cases\} \& 0
\textless{} t \textless{} \textbackslash infty \textbackslash{}
(\textbackslash text\{HV\}) \& u ~(x, ~0) =
\textbackslash sin\{(\textbackslash pi x)\}, ~0 \textbackslash leqslant
x \textbackslash leqslant 1 \textbackslash end\{array\} {]}

дается формулой (9.12). Соответствует ли это решение вашим интуитивным
представлениям? Что представляет собой полученное решение?

\begin{enumerate}
\def\labelenumi{\arabic{enumi}.}
\setcounter{enumi}{1}
\tightlist
\item
  Решите задачу
\end{enumerate}

{[} \textbackslash begin\{array\}\{ll\} (\textbackslash text\{YUII\}) \&
u\_t = u\_\{xx\} + \textbackslash sin\{(\textbackslash pi x)\} +
\textbackslash sin\{(2\textbackslash pi x)\}, \& 0 \textless{} x
\textless{} 1, \& 0 \textless t \textless{} \textbackslash infty,
\textbackslash{} (\textbackslash text\{\textbackslash GammaY\}) \&
\textbackslash begin\{cases\}
u\textbackslash left(0,~t\textbackslash right) = 0, \textbackslash{}
u\textbackslash left(1,~t\textbackslash right) = 0,
\textbackslash end\{cases\} \& 0 \textless{} t \textless{}
\textbackslash infty, \textbackslash{} (\textbackslash text\{HY\}) \&
u\textbackslash left(x,~0\textbackslash right) = 0, \& 0
\textbackslash leqslant x \textbackslash leqslant 1.
\textbackslash end\{array\} {]}

\section{3. Розв'яжіть
задачу}\label{ux440ux43eux437ux432ux44fux436ux456ux442ux44c-ux437ux430ux434ux430ux447ux443}

(PPP)

{[} u\_t = u\_\{xx\} + \textbackslash sin(\textbackslash pi x),
\textbackslash quad 0 \textless{} x \textless{} 1, \textbackslash quad 0
\textless{} t \textless{} \textbackslash infty, {]}

(GI)

{[} \textbackslash begin\{cases\} u(0, t) = 0, \textbackslash{} u(1, t)
= 0, \textbackslash{} u(x, 0) = 1, \textbackslash quad 0 \textless{} x
\textless{} 1, \textbackslash end\{cases\} {]}

методом разложения по собственным функциям.

\begin{enumerate}
\def\labelenumi{\arabic{enumi}.}
\setcounter{enumi}{3}
\tightlist
\item
  Найдите решение задачи
\end{enumerate}

{[} \textbackslash begin\{array\}\{ll\} (\textbackslash text\{YHII\}) \&
u\_t = u\_\{xx\} + \textbackslash sin\{(\textbackslash lambda\_1 x)\},
\& 0 \textless{} x \textless{} 1, \& 0 \textless t \textless{}
\textbackslash infty, \textbackslash{}
(\textbackslash text\{\textbackslash GammaY\}) \&
\textbackslash begin\{cases\} u\textbackslash left(0,
~t\textbackslash right) = 0, \textbackslash{} u\_x\textbackslash left(1,
~t\textbackslash right) + u\textbackslash left(1,
~t\textbackslash right) = 0, \textbackslash{} 0 \textless{} t
\textless{} \textbackslash infty, \textbackslash end\{cases\}
\textbackslash{} (\textbackslash text\{HY\}) \& u\textbackslash left(x,
~0\textbackslash right) = 0, \& 0 \textbackslash leqslant x
\textbackslash leqslant 1 \textbackslash end\{array\} {]}

методом разложения по собственным функциям, если
\textbackslash lambda\_1 ---первый корень уравнения
\textbackslash operatorname\{tg\} \textbackslash lambda =
-\textbackslash lambda . Каковы собственные функция этой задачи?

\section{5. Решите
задачу}\label{ux440ux435ux448ux438ux442ux435-ux437ux430ux434ux430ux447ux443}

{[} \textbackslash begin\{array\}\{ll\} (\textbackslash text\{УЧП\}) \&
u\_t = u\_\{xx\} \& 0 \textless{} x \textless{} 1, \textbackslash quad 0
\textless{} t \textless{} \textbackslash infty, \textbackslash{}
(\textbackslash Gamma \textbackslash text\{У\}) \&
\textbackslash begin\{cases\} u(0, t) = 0, \textbackslash{} u(1, t) =
\textbackslash cos t, \textbackslash end\{cases\} \& 0 \textless{} t
\textless{} \textbackslash infty, \textbackslash{}
(\textbackslash text\{HY\}) \& u(x, 0), \textbackslash quad 0
\textbackslash leqslant x \textbackslash leqslant 1,
\textbackslash end\{array\} {]}

Спочатку перетворивши граничні умови в нуль. Отримана задача
розв'язується методом розкладання за власними функціями.

\section{ІНТЕГРАЛЬНІ ПЕРЕТВОРЕННЯ (СИНУСОЇДАЛЬНІ ТА КОСИНУСНІ
ПЕРЕТВОРЕННЯ)}\label{ux456ux43dux442ux435ux433ux440ux430ux43bux44cux43dux456-ux43fux435ux440ux435ux442ux432ux43eux440ux435ux43dux43dux44f-ux441ux438ux43dux443ux441ux43eux457ux434ux430ux43bux44cux43dux456-ux442ux430-ux43aux43eux441ux438ux43dux443ux441ux43dux456-ux43fux435ux440ux435ux442ux432ux43eux440ux435ux43dux43dux44f}

МЕТА ЛЕКЦІЇ: Ввести концепцію інтегрального перетворення та показати, як
використання інтегральних перетворень дозволяє перейти від CPP з n
змінними до нового рівняння з n---1 змінними.

Лекція показує, що інтегральні перетворення можна інтерпретувати як
розклад вхідної дії на елементарні (частотний розклад), визначення
реакції системи на елементарні впливи та подальше підсумовування всіх
отриманих результатів.

У інтегральних перетвореннях операція диференціювання переходить у
операцію множення, і, відповідно, деякі часткові похідні переходять у
алгебраїчні вирази.

Для розв'язання задач дифузії на необмежених інтервалах вводять
синусоїдальні та косинусні перетворення. Розв'язок, отриманий у цьому
випадку, цікавий тим, що містить функцію, яка є додатковою до
ймовірнісного інтеграла.

Інтегральне перетворення зазвичай називається перетворенням, яке кожній
функції \(f(t)\) призначає нову функцію \(F(s)\) відповідно до формули

{[} F(s) = \textbackslash int\_\{A\}\^{}\{B\} K(s, t) f(t) dt. {]}

Зверніть увагу, що початкова функція залежить від t, а отримана функція
--- від s. Функція \(K(s,t)\) називається ядром перетворення. Це
головний елемент, який відрізняє одну трансформацію від іншої. Зазвичай
ядро обирається так, щоб перетворення мало певні задані властивості.
Межі інтегрування також залежать від типу трансформації.

Голосною, що лежить в основі інтегральних перетворень, є здатність
позбутися часткових похідних однієї з незалежних змінних. Це означає, що
перетворене рівняння матиме на одну змінну менше, ніж початкове.
Наприклад, якщо застосувати інтегральне перетворення до рівняння

Щоб позбутися похідної часу, отримуємо звичайне диференціальне рівняння
з незалежною змінною X.

{[} u\_\{xx\} + u\_\{cy\} + u\_\{zz\} = 0 {]}

застосувати перетворення Фур'є до змінної x, тоді похідна \(u_{xx}\)
буде усунена, і в новому рівнянні залишаться лише дві незалежні змінні y
та z. Тоді отримуємо звичайне диференціальне рівняння з однією
незалежною змінною. Іншими словами, інтегральні перетворення спрощують
початкове завдання. Потім розв'язується перетворення задачі, і,
застосовуючи обернене перетворення до її розв'язку, отримується
розв'язок початкового перетворення. Схема описаного алгоритму показана
на рис. 10.1

\pandocbounded{\includegraphics[keepaspectratio]{images/_page_72_Figure_5.jpeg}}

РИС 10.1, Загальна схема застосування інтегрованих перетворень.

Фіг. На рисунку 10.1 видно, що кожне перетворення прямого інтеграла
пов'язане з оберненим перетворенням, яке відновлює початкову функцію з
трансформованої.

Прямі та обернені перетворення зазвичай називають парою перетворень.
Таблиця 1. Рисунок 10.1 показує кілька пар перетворень, які ми будемо
використовувати при розв'язуванні рівнянь із факторами.

\begin{verbatim}
Слід зазначити, що в цих перетвореннях використовуються два типи символів. Наприклад, у випадку перетворення Лапласа позначення \mathcal{L}[f] вказує, що функція f перетворена, тоді як позначення F(s) означає, що перетворення породжує функцію, яка залежить від s.
\end{verbatim}

У цій лекції ми не прагнемо вивчати всі пари перетворень. Ми
зосередимося лише на перетвореннях Фур'є синуса та косинуса (пари 1 і
2). Трохи пізніше ми вивчаємо іншу з цих пар. Ознайомляючись із
трансформаціями, ми дізнаємося про їхні переваги та недоліки,
дізнаємося, як

Таблиця 10.1 Деякі пари інтегральних перетворень

\begin{enumerate}
\def\labelenumi{\arabic{enumi}.}
\tightlist
\item
  \textbackslash begin\{cases\} \textbackslash mathcal\{F\}\_s{[}f{]} =
  F(\textbackslash omega) =
  \textbackslash frac\{2\}\{\textbackslash pi\}
  \textbackslash int\_0\^{}\textbackslash infty f(t)
  \textbackslash sin(\textbackslash omega t) dt \&
  \textbackslash text\{Синус-преобразование\} \textbackslash{}
  \textbackslash mathcal\{F\}\_s\^{}\{-1\}{[}F{]} = f(t) =
  \textbackslash int\_0\^{}\textbackslash infty F(\textbackslash omega)
  \textbackslash sin(\textbackslash omega t) d\textbackslash omega \&
  \textbackslash text\{Обратное синус-преобразование\} \textbackslash{}
  \textbackslash mathcal\{F\}\_s\^{}\{-1\}{[}F{]} = f(t) =
  \textbackslash int\_0\^{}\textbackslash infty f(t)
  \textbackslash cos(\textbackslash omega t) dt \&
  \textbackslash text\{Косниус-преобразование\} \textbackslash{}
  \textbackslash mathcal\{F\}\_c{[}f{]} = F(\textbackslash omega) =
  \textbackslash frac\{2\}\{\textbackslash pi\}
  \textbackslash int\_0\^{}\textbackslash infty f(t)
  \textbackslash cos(\textbackslash omega t) d\textbackslash omega \&
  \textbackslash text\{Обратное косинус-преобразование\}
  \textbackslash{} \textbackslash mathcal\{F\}\_c\^{}\{-1\}{[}F{]} =
  f(t) = \textbackslash int\_0\^{}\textbackslash infty
  F(\textbackslash omega) \textbackslash cos(\textbackslash omega t)
  d\textbackslash omega \& \textbackslash text\{Обратное
  преобразование\} \textbackslash{}
  \textbackslash mathcal\{F\}\emph{s\^{}\{-1\}{[}F{]} = f(x) =
  \textbackslash frac\{1\}\{\textbackslash sqrt\{2\textbackslash pi\}\}
  \textbackslash int}\{-\textbackslash infty\}\^{}\textbackslash infty
  f(x) e\^{}\{-i\textbackslash omega x\} dx \&
  \textbackslash text\{Преобразование\} \textbackslash Phi
  \textbackslash text\{урье\} \textbackslash{}
  \textbackslash mathcal\{F\}\_s\^{}\{-1\}{[}F{]} = f(x) =
  \textbackslash frac\{1\}\{\textbackslash sqrt\{2\textbackslash pi\}\}
  \textbackslash int\_0\^{}\textbackslash infty F(\textbackslash omega)
  e\^{}\{i\textbackslash omega x\} d\textbackslash omega \&
  \textbackslash text\{Обратное преобразование\} \textbackslash{}
  \textbackslash mathcal\{F\}\emph{s\^{}\{-1\}{[}F{]} = f(x) =
  \textbackslash int}\{n=1\}\^{}\textbackslash infty
  \textbackslash int\_0\^{}\textbackslash infty f(x)
  \textbackslash sin(n\textbackslash pi x/L) dx \&
  \textbackslash text\{Конечное синус-преобразование\} \textbackslash{}
  \textbackslash mathcal\{F\}\emph{s\^{}\{-1\}{[}F\_n{]} = f(x) =
  \textbackslash sum}\{n=1\}\^{}\textbackslash infty S\_n
  \textbackslash sin(n\textbackslash pi x/L) \&
  \textbackslash text\{Обратное конечное синус-преобразование\}
  \textbackslash{} \textbackslash mathcal\{F\}\_s\^{}\{-1\}{[}F\_n{]} =
  f(x) = \textbackslash frac\{2\}\{L\} \textbackslash int\_0\^{}L f(x)
  \textbackslash cos(n\textbackslash pi x/L) dx \&
  \textbackslash text\{Конечное косинус-преобразование\}
  \textbackslash{}
  \textbackslash mathcal\{F\}\emph{s\^{}\{-1\}{[}F\_n{]} = f(x) =
  \textbackslash frac\{C\_0\}\{2\} +
  \textbackslash sum}\{n=1\}\^{}\textbackslash infty C\_n
  \textbackslash cos(n\textbackslash pi x/L) \&
  \textbackslash text\{Обратное конечное косинус-преобразование\}
  \textbackslash{} \textbackslash mathcal\{F\}\_s\^{}\{-1\}{[}F{]} =
  f(t) = \textbackslash frac\{1\}\{2\textbackslash pi i\}
  \textbackslash int\_0\^{}\textbackslash infty f(x) e\^{}\{-it\} dx \&
  \textbackslash text\{Преобразование\} \&
  \textbackslash text\{Лапласа\} \textbackslash{}
  \textbackslash mathcal\{L\}\^{}\{-1\}{[}F{]} = f(t) =
  \textbackslash frac\{1\}\{2\textbackslash pi i\}
  \textbackslash int\_0\^{}\textbackslash infty f(x) e\^{}\{-it\} dx \&
  \textbackslash text\{Обратное преобразование\} \textbackslash{}
  \textbackslash mathcal\{L\}\^{}\{-1\}{[}F{]} = f(t) =
  \textbackslash frac\{1\}\{2\textbackslash pi i\}
  \textbackslash int\_0\^{}\textbackslash infty f(x) e\^{}\{-it\} dx \&
  \textbackslash text\{Обратное преобразование\} \textbackslash{}
  \textbackslash mathcal\{L\}\^{}\{-1\}{[}F{]} = f(t) =
  \textbackslash frac\{1\}\{2\textbackslash pi i\}
  \textbackslash int\_0\^{}\textbackslash infty f(x) e\^{}\{-it\} dx \&
  \textbackslash text\{Обратное преобразование\} \textbackslash{}
  \textbackslash mathcal\{L\}\^{}\{-1\}{[}F{]} = f(t) =
  \textbackslash frac\{1\}\{2\textbackslash pi i\}
  \textbackslash int\_0\^{}\textbackslash infty f(x) e\^{}\{-it\} dx \&
  \textbackslash text\{Обратное преобразование\} \textbackslash{}
  \textbackslash mathcal\{L\}\^{}\{-1\}{[}F{]} = f(t) =
  \textbackslash frac\{1\}\{2\textbackslash pi i\}
  \textbackslash int\_0\^{}\textbackslash infty f(x) e\^{}\{-it\} dx \&
  \textbackslash text\{Обратное преобразование\} \textbackslash{}
  \textbackslash mathcal\{L\}\^{}\{-1\}{[}F{]} = f(t) =
  \textbackslash frac\{1\}\{2\textbackslash pi i\}
  \textbackslash int\_0\^{}\textbackslash infty f(x) e\^{}\{-it\} dx \&
  \textbackslash text\{Обратное преобразование\} \textbackslash{}
  \textbackslash mathcal\{L\}\^{}\{-1\}{[}F{]} = f(t) =
  \textbackslash frac\{1\}\{2\textbackslash pi i\}
  \textbackslash int\_0\^{}\textbackslash infty f(x) e\^{}\{-it\} dx \&
  \textbackslash text\{Обратное преобразование\} \textbackslash{}
  \textbackslash mathcal\{L\}\^{}\{-1\}{[}F{]} = f(t) =
  \textbackslash frac\{1\}\{2\textbackslash pi i\}
  \textbackslash int\_0\^{}\textbackslash infty f(x) e\^{}\{-it\} dx \&
  \textbackslash text\{Обратное преобразование\} \textbackslash{}
  \textbackslash mathcal\{L\}\^{}\{-1\}{[}F{]} = f(t) =
  \textbackslash frac\{1\}\{2\textbackslash pi i\}
  \textbackslash int\_0\^{}\textbackslash infty f(x) e\^{}\{-it\} dx \&
  \textbackslash text\{Обратное преобразование\} \textbackslash{}
  \textbackslash mathcal\{L\}\^{}\{-1\}{[}F{]} = f(t) =
  \textbackslash frac\{1\}\{2\textbackslash pi i\}
  \textbackslash int\_0\^{}\textbackslash infty f(x) e\^{}\{-it\} dx \&
  \textbackslash text\{Обратное преобразование\} \textbackslash{}
  \textbackslash mathcal\{L\}\^{}\{-1\}{[}F{]} = f(t) =
  \textbackslash frac\{1\}\{2\textbackslash pi i\}
  \textbackslash int\_0\^{}\textbackslash infty f(x) e\^{}\{-it\} dx \&
  \textbackslash text\{Обратное преобразование\} \textbackslash{}
  \textbackslash mathcal\{L\}\^{}\{-1\}{[}F{]} = f(t) =
\end{enumerate}

Продовження

\begin{enumerate}
\def\labelenumi{\arabic{enumi}.}
\setcounter{enumi}{6}
\tightlist
\item
\end{enumerate}

{[} \textbackslash begin\{cases\} H\{f\} = F\_n(\textbackslash xi) =
\textbackslash int\textbackslash limits\_0\^{}\textbackslash infty r
J\_n(\textbackslash xi r) f(r) , dr \&
\textbackslash text\{Преобразование Ханкеля\} \textbackslash{}
H\^{}\{-1\}{[}F\_n{]} = f(r) =
\textbackslash int\textbackslash limits\_0\^{}\textbackslash infty
\textbackslash xi J\_n(\textbackslash xi r) F\_n(\textbackslash xi) ,
d\textbackslash xi \& \textbackslash text\{Обратное преобразование
Ханкеля\} \textbackslash end\{cases\} {]}

вони пов'язані між собою при їх застосуванні тощо. Однак перед тим, як
перейти до вивчення інтегральних перетворень, буде корисно зрозуміти, що
зазвичай називають спектром функції (або спектральним розкладом
функції).

\section{Спектр
функцій}\label{ux441ux43fux435ux43aux442ux440-ux444ux443ux43dux43aux446ux456ux439}

Інтегральні перетворення та спектр функції тісно пов'язані. Насправді
інтегральне перетворення можна інтерпретувати як розклад функції на
певний спектр компонент. Цей спектр залежить від обраного перетворення,
але кожна функція має певний спектр відносно цього перетворення Р и с

\pandocbounded{\includegraphics[keepaspectratio]{images/_page_74_Picture_7.jpeg}}

. 10 2 Розклад періодичної функції на ряд синусами та косинусами;
Послідовність квадратних імпульсів апроксимується синусами та косинусами

Наприклад, розглянемо розклад періодичної функції на ряд синусами та
косинусами (ряд Фур'є) 1:

{[} f(x) = \textbackslash sum\_\{n=0\}\^{}\{\textbackslash infty\}
\textbackslash left{[} A\_n \textbackslash cos(nx) + B\_n
\textbackslash sin(nx) \textbackslash right{]} {]}

(див. рис. 10.2).

\begin{verbatim}
Тут коефіцієнти A_n і B_n позначають внесок \cos(nx) і \sin(nx) у функцію f(x), тоді як величина
\end{verbatim}

{[} V\textbackslash overline\{A\_n\textsuperscript{2+B\_n}3\} {]}

1) Фур'є-ряди детальніше розглядаються у лекції 11,

(називається спектром функції) показує частку амплітуди функції
\(f(x)\), що припадає на частоту n.

Наприклад, якщо функція \(f(x)\) --- це просто сума

синуси та косинуси

{[} f(x) = 1 + \textbackslash sin x +
\textbackslash frac\{1\}\{5\}\textbackslash sin(3x) + \textbackslash cos
x + \textbackslash frac\{1\}\{2\}\textbackslash cos(2x) +
\textbackslash frac\{1\}\{4\}\textbackslash cos(4x), {]}

тоді його спектр (дискретний) виглядатиме так, як показано на рис. 10.3.

\begin{verbatim}
Зчитуючи значення значення \sqrt{A_n^2 + B_n^2} з графіка, можна визначити амплітудне значення компоненти функції f(x), яка має частоту n.
\end{verbatim}

\pandocbounded{\includegraphics[keepaspectratio]{images/_page_75_Figure_8.jpeg}}

Рис. 10 3. Дискретний спектр функції \(f(r)\).

Періодичні функції розкладаються на нескінченні ряди (вони мають
дискретний спектр), тоді як неперіодичні функції повинні мати
неперервний спектр. Якщо функція визначена лише на скінченному
інтервалі, її можна періодично розширювати на всю вісь і розкладати на
ряд Фур'є, який представлятиме нашу функцію в межах визначення області.

І хоча неперіодичну функцію не можна розкласти на нескінченний ряд
синусами та косинусами, можна спробувати записати неперервний аналог
ряду Фур'є:

{[} f(x) =
\textbackslash int\_\{-\textbackslash infty\}\^{}\{+\textbackslash infty\}
{[}C(\textbackslash omega)\textbackslash cos(\textbackslash omega x) +
S(\textbackslash omega)\textbackslash sin(\textbackslash omega x){]}
d\textbackslash omega, {]}

де функції S(\textbackslash omega) і C(\textbackslash omega) слугують
мірою спектральної густини синуса і косинуса у функції \(f(x)\), а також
значення

{[} V \textbackslash overline\{S\^{}2(\textbackslash omega) +
C\^{}2(\textbackslash omega)\} {]}

Він слугує мірою загального внеску частотного \textbackslash omega у
функцію \(f(x)\) і називається спектром (неперервним) функції \(f(x)\).
Виходячи з наведеного вище інтуїтивного визначення спектра функції,
давайте проаналізуємо інтегральні перетворення «гвинт за гвинтом».
Спершу запишемо кілька властивостей цих перетворень, які ми будемо
широко використовувати.

\section{Синусоїдальні та косинусні перетворення
похідних}\label{ux441ux438ux43dux443ux441ux43eux457ux434ux430ux43bux44cux43dux456-ux442ux430-ux43aux43eux441ux438ux43dux443ux441ux43dux456-ux43fux435ux440ux435ux442ux432ux43eux440ux435ux43dux43dux44f-ux43fux43eux445ux456ux434ux43dux438ux445}

\begin{enumerate}
\def\labelenumi{\arabic{enumi}.}
\tightlist
\item
\end{enumerate}

{[} \textbackslash mathcal\{F\}\_s{[}f'{]} = -\textbackslash omega
\textbackslash mathcal\{F\}\_c{[}f{]} {]}

(доведено інтегруванням по частинах).

\begin{enumerate}
\def\labelenumi{\arabic{enumi}.}
\tightlist
\item
\end{enumerate}

{[} \textbackslash mathcal\{F\}\_s{[}f'{]} = -\textbackslash omega
\textbackslash mathcal\{F\}\_c{[}f{]} {]}

\begin{verbatim}
2. \mathcal{F}_s[f''] = \frac{2}{\pi} \omega f(0) - \omega^2 \mathcal{F}_s[f] (10.1) 3.
\end{verbatim}

{[} \textbackslash mathcal\{F\}\emph{\{\textbackslash rho\}{[}f'{]} =
\textbackslash frac\{2\}\{\textbackslash pi\} f(0) +
\textbackslash omega \textbackslash mathcal\{F\}}\{s\}{[}f{]} {]}

.

\begin{enumerate}
\def\labelenumi{\arabic{enumi}.}
\setcounter{enumi}{3}
\tightlist
\item
\end{enumerate}

{[} \textbackslash mathcal\{F\}\emph{\{c\}{[}f'\,') =
-\textbackslash frac\{2\}\{\textbackslash pi\} f}\{c\}(0) -
\textbackslash omega\^{}\{2\} \textbackslash mathcal\{F\}\_\{c\}{[}f{]}.
{]}

Деякі інші властивості прямих і обернених синусоїдальних і спеційних
перетворень наведені в таблиці наприкінці книги. Тепер покажемо, як
синусоїдальне перетворення можна використати для розв'язання важливої
змішаної задачі.

\section{Розв'язок задачі дифузії на напівпрямому синусоїдальному
метрі}\label{ux440ux43eux437ux432ux44fux437ux43eux43a-ux437ux430ux434ux430ux447ux456-ux434ux438ux444ux443ux437ux456ux457-ux43dux430-ux43dux430ux43fux456ux432ux43fux440ux44fux43cux43eux43cux443-ux441ux438ux43dux443ux441ux43eux457ux434ux430ux43bux44cux43dux43eux43cux443-ux43cux435ux442ux440ux456}

Задача, яку ми хочемо розв'язати, формулюється так:

(PPP)

{[} u\_t = \textbackslash alpha\^{}2 u\_\{xx\} {]}

{[} , 0 \textless{} x \textless{} \textbackslash infty , 0 \textless{} t
\textless{} \textbackslash infty , (GU) u(0, t) = A , 0 \textless{} t
\textless{} \textbackslash infty , (NU) u(x, 0) = 0 , 0
\textbackslash le x \textless{} \textbackslash infty . {]}

\pandocbounded{\includegraphics[keepaspectratio]{images/_page_76_Figure_14.jpeg}}

\begin{verbatim}
Рис. 10.4. Задача дифузії в напівнескінченній області; u\left(0,\,t\right)= = A, u(x, 0) = 0.
\end{verbatim}

КРОК 1. Розберемо розв'язання проблеми на три досить прості кроки.
Діятимемо відповідно до наступного плану. На першому кроці виконуємо
перетворення синуса Фур'є на змінній x.

тоді отримуємо звичайне диференціальне рівняння для змінної t. Для
реалізації цього плану застосовуємо синусоїдальне перетворення до обох
частин рівняння з частковими похідними. У результаті ми отримуємо

{[} \textbackslash mathcal\{F\}\_s{[}u\_t{]} = \textbackslash alpha\^{}2
\textbackslash mathcal\{F\}\_s{[}u\_\{xx\}{]}. {]}

Розглянемо кожну частину цього співвідношення.
\textbackslash mathcal\{F\}\_s{[}\(u_t\){]} : часткову похідну \(u_t\) в
цій задачі можна назвати зовнішньою похідною, оскільки ми використовуємо
перетворення на x. У цьому випадку можна записати

{[}
\textbackslash mathcal\{F\}\emph{\{s\}\textbackslash left{[}u\_\{t\}\textbackslash right{]}
= \textbackslash frac\{2\}\{\textbackslash pi\}
\textbackslash int}\{0\}\^{}\{\textbackslash infty\}
u\_\{t\}\textbackslash left(x, t\textbackslash right)
\textbackslash sin\textbackslash left(\textbackslash omega
x\textbackslash right) dx = {]}

{[} =
\textbackslash frac\{\textbackslash partial\}\{\textbackslash partial
t\} \textbackslash left{[}\textbackslash frac\{2\}\{\textbackslash pi\}
\textbackslash int\_\{0\}\^{}\{\textbackslash infty\}
u\textbackslash left(x, t\textbackslash right)
\textbackslash sin\textbackslash left(\textbackslash omega
x\textbackslash right) dx\textbackslash right{]} = {]}

{[} = \textbackslash frac\{d\}\{dt\}
\textbackslash mathcal\{F\}\_\{s\}\textbackslash left{[}u\textbackslash right{]}
= {]}

{[} = \textbackslash frac\{d\}\{dt\}
U\textbackslash left(t\textbackslash right). {]}

Тут ми взяли операцію диференціації поза знаком інтеграла, застосовуючи
відповідну теорему інтегрального числення. Зверніть увагу, що початкова
функція u залежить від змінних x і t, а також від її образу

{[} \textbackslash mathcal\{F\}\_\{s\}{[}u{]} = U(\textbackslash omega,
t) {]}

залежить від змінних \(\omega\) і t. Нова змінна \(\omega\) відіграє
роль параметра в новій задачі, і тому можна припустити, що образи
функцій залежать лише від однієї змінної t:

{[} \textbackslash mathcal\{F\}\_s{[}u{]} = U(t) {]}

. \({\mathscr F}_s[u_{xx}]\) : Для цього елемента дійсні наступні
рівняння:

{[} \textbackslash begin\{split\}
\&\textbackslash mathcal\{F\}\_s\textbackslash left{[}u\_\{xx\}\textbackslash right{]}
= \textbackslash frac\{2\}\{\textbackslash pi\} \textbackslash omega u
~(0, \textbackslash t) - \textbackslash omega\^{}2
\textbackslash mathcal\{F\}\_s\textbackslash left{[}u\textbackslash right{]}
= \textbackslash{} \&= \textbackslash frac\{2\}\{\textbackslash pi\}
\textbackslash omega u ~(0, \textbackslash t) -
\textbackslash omega\^{}2 U ~(t) = \textbackslash{} \&=
\textbackslash frac\{2A\textbackslash omega\}\{\textbackslash pi\} -
\textbackslash omega\^{}2 U ~(t). \textbackslash end\{split\} {]}

Зверніть увагу, що при виведенні відношень (10.1) ми розглядали функцію
f як функцію однієї змінної x. Тут випадок дещо інший, оскільки
\(u\left(x,\ t\right)\) залежить і від x, і від t. Однак ті ж формули
можна використовувати, якщо перетворення виконується однією змінною, а
решта вважаються константами.

У нашому випадку перетворення виконується змінною x, тобто змінну t
можна вважати сталою. Звернімо також увагу на той факт, що вже на цьому
етапі розв'язання задачі використовується гранична умова
\(u\left(0,\ t\right)=A\) . Підставляючи знайдені вирази у рівняння
\(u_{t}=\alpha^{2}u_{xx}\) , отримуємо звичайне диференціальне рівняння

{[} \textbackslash frac\{dU\}\{dt\} = \textbackslash alpha\^{}2
\textbackslash left{[} -\textbackslash omega\^{}2 U(t) +
\textbackslash frac\{2A\textbackslash omega\}\{\textbackslash pi\}
\textbackslash right{]}. {]}

Єдине, що існує --- це початкова умова функції \(U(t)\). Отримуємо його,
застосовуючи перетворення до початкової умови \(u(x, 0)\) = 0:

{[} \textbackslash mathcal\{F\}\_\{s\}{[}u(x, 0){]} = U(0) = 0. {]}

Це завершує перший крок у побудові розв'язку методом інтегрального
перетворення: замість початкової змішаної задачі ми отримали задачу Коші
для звичайного диференціального рівняння:

(10.2)

{[} \textbackslash begin\{aligned\} (O \textbackslash Pi V) \&
\textbackslash frac\{dU\}\{di\} + \textbackslash omega\^{}2
\textbackslash alpha\^{}2 U =
\textbackslash frac\{2A\textbackslash omega
\textbackslash alpha\^{}2\}\{\textbackslash pi\}, \textbackslash{} (H V)
\& U(0) = 0. \textbackslash end\{aligned\} {]}

КРОК 2. Щоб розв'язати отриману задачу Коші, можна використати будь-який
елементарний метод розв'язання звичайних диференціальних рівнянь. У
результаті ми отримуємо

{[} U(t) =
\textbackslash frac\{2A\}\{\textbackslash pi\textbackslash omega\} (1 -
e\textsuperscript{\{-\textbackslash omega}2 \textbackslash alpha\^{}2
t\}). {]}

Отже, перед нами маємо результат застосування синусоїдального
перетворення до бажаної функції \(u\left(x,\ t\right)\) . Останній крок
--- застосувати обернене перетворення до функції \(U\left(t\right)\) і
отримати потрібне рішення \(u\left(x,\ t\right)\) :

{[} u(x, t) = \textbackslash mathcal\{F\}\_\{x\}\^{}\{-1\}{[}U{]}. {]}

КРОК 3. Щоб отримати розв'язок, можна або безпосередньо обчислити
інтеграл, що входить до оберненого перетворення, або використати таблиці
перетворень. Використовуючи таблиці, ми знаходимо

{[} u(x, t) = A
\textbackslash operatorname\{erfc\}(x/2\textbackslash alpha
\textbackslash sqrt\{t\}), {]}

де erfc(x) --- функція, доповнююча до ймовірнісного інтегралу, а
визначається відношенням

{[} \textbackslash operatorname\{erfc\}(x) =
\textbackslash frac\{2\}\{\textbackslash sqrt\{2\textbackslash pi\}\}
\textbackslash int\_\{x\}\^{}\{\textbackslash infty\}
e\textsuperscript{\{-t}\{2\}\} dt. {]}

Графік цієї функції показано на рис. 10.5.

Ймовірнісний інтеграл і додаткова функція до нього не виражаються через
елементарні функції. Вони добре популярні, і відповідні таблиці можна
знайти в довідниках з фізики та хімії.

\pandocbounded{\includegraphics[keepaspectratio]{images/_page_79_Figure_4.jpeg}}

Рис. 10.5. Графи функцій erf (x) та eric (x).

\section{Інтерпретація
рішення}\label{ux456ux43dux442ux435ux440ux43fux440ux435ux442ux430ux446ux456ux44f-ux440ux456ux448ux435ux43dux43dux44f}

Розв'язок

{[} u(x, t) = A
\textbackslash operatorname\{erfc\}{[}x/2\textbackslash alpha V
\textbackslash overline\{t\}{]} {]}

дає багато їжі для роздумів. На рисунку 10.6 показані графіки
ймовірнісної комплементарної функції на різних значеннях часу t.

\pandocbounded{\includegraphics[keepaspectratio]{images/_page_79_Figure_10.jpeg}}

Рис. 10.6. Розподіл температури в напівнескінченному стрижні, на кінці
якого підтримується постійна температура.

Вони відрізняються один від одного лише характеристичною шкалою вздовж
осі x.

\section{ЗАВДАННЯ}\label{ux437ux430ux432ux434ux430ux43dux43dux44f-6}

\begin{itemize}
\tightlist
\item
  \begin{enumerate}
  \def\labelenumi{\arabic{enumi}.}
  \tightlist
  \item
    Довести ідентичності (10.1).
  \end{enumerate}
\item
  \begin{enumerate}
  \def\labelenumi{\arabic{enumi}.}
  \setcounter{enumi}{1}
  \tightlist
  \item
    Розв'яжіть задачу Коші (10.2).
  \end{enumerate}
\item
  \begin{enumerate}
  \def\labelenumi{\arabic{enumi}.}
  \setcounter{enumi}{2}
  \tightlist
  \item
    Використайте синусоїдальне або коспусне перетворення для розв'язання
  \end{enumerate}
\end{itemize}

\textbf{z}adachu

{[} \textbackslash begin\{aligned\} (\textbackslash text\{УЧП\}) \& u\_t
= \textbackslash alpha\^{}2 u\_\{xx\}, \& 0 \textless{} x \textless{}
\textbackslash infty, \textbackslash{} (\textbackslash text\{ГУ\}) \&
u\_x (0, ~t) = 0, \& 0 \textless{} t \textless{} \textbackslash infty,
\textbackslash{} (\textbackslash text\{НУ\}) \& u (x, ~0) = H (1-x), \&
0 \textbackslash leqslant x \textless{} \textbackslash infty,
\textbackslash end\{aligned\} {]}

де \(H(x)\) --- функція Хевісайду

{[} H(x) = {]}

{[} \textbackslash begin\{cases\} 0, \& x \textless{} 0,
\textbackslash{} 1, \& 0 \textbackslash le x.
\textbackslash end\{cases\} {]}

Входящая в початкова умова функція изображена на рисунке.

\pandocbounded{\includegraphics[keepaspectratio]{images/_page_80_Picture_7.jpeg}}

Как буде выглядеть график розв'язки в различные моменты времени?

\begin{verbatim}
5. Найдите прообраз константы при синус-преобразовании, т. е. найдите f(x), якщо F(\omega) = c = \mathrm{const.} # # РЯДЫ И ПРЕОБРАЗОВАНИЯ ФУРЬЕ
\end{verbatim}

ЦЕЛЬ ЛЕКЦИИ: Ввести понятие ряда Фурье и показать, каким образом
периодическую функцию \(f(x)\) можна представить в гиде ряда по синусам
и косинусам:

{[} f(x) = \textbackslash frac\{a\_0\}\{2\} +
\textbackslash sum\_\{n=1\}\^{}\{\textbackslash infty\} {[}a\_n
\textbackslash cos(n\textbackslash pi x) + b\_n
\textbackslash sin(n\textbackslash pi x){]}. {]}

Якщо же функція \(f(x)\) непериодическая и определена на
(-\textbackslash infty, +\textbackslash infty) , то показать, как ряд
Фурье переходит в \emph{преобразование Фурье} и как функцию \(f(x)\)
можна разложить в непрерывную совокупность простейших функцій. Эго
разложение (интеграл Фурье) можна записать в комплексной форме:

{[} f(x) =
\textbackslash frac\{1\}\{\textbackslash sqrt\{2\textbackslash pi\}\}
\textbackslash int\_\{-\textbackslash infty\}\^{}\{+\textbackslash infty\}
\textbackslash left{[}
\textbackslash frac\{1\}\{\textbackslash sqrt\{2\textbackslash pi\}\}
\textbackslash int\_\{-\textbackslash infty\}\^{}\{+\textbackslash infty\}
f(x) e\^{}\{-i\textbackslash xi x\} dx \textbackslash right{]}
e\^{}\{i\textbackslash xi x\} d\textbackslash xi, {]}

откуда получаются прямое и обратное преобразования Фурье

{[} \textbackslash mathcal\{F\}{[}f{]} = F(\textbackslash xi) =
\textbackslash frac\{1\}\{\textbackslash sqrt\{2\textbackslash pi\}\}
\textbackslash int\_\{-\textbackslash infty\}\^{}\{+\textbackslash infty\}
f(x) e\^{}\{-i\textbackslash xi x\} dx \textbackslash qquad
\textbackslash text\{(пряме перетворення Фур'є),\} {]}

{[} \textbackslash mathcal\{F\}\^{}\{-1\}{[}F{]} = f(x) =
\textbackslash frac\{1\}\{\textbackslash sqrt\{2\textbackslash pi\}\}
\textbackslash int\_\{-x\}\^{}\{+x\} F(x) e\^{}\{i\textbackslash xi x\}
d\textbackslash xi \textbackslash qquad \textbackslash text\{(обернене
перетворення Фур'є)\}. {]}

Важная роль рядов Фурье в теории уравнений з частинними похідними
обусловлена тем, що периодические функції, определенные на
\((-\infty, +\infty)\) , или функції, определенные на конечном
интервате, можна разложить в ряд по сипусам и косинусам. Це позволяет
упростить исходную задачу. Наприклад, пилообразная функція

{[} f(x) = x, -L \textless{} x \textless{} L, {]}

\(f(x + 2L) = f(x)\) (условие периодичности),

изображенная на рис. 11.1, может быть представлена в виде ряда Фурье

(11.1)

{[} f(x) = \textbackslash frac\{a\_0\}\{2\} +
\textbackslash sum\_\{n=1\}\^{}\{\textbackslash infty\}
\textbackslash left{[} a\_n \textbackslash cos(n\textbackslash pi x/L) +
b\_n \textbackslash sin(n\textbackslash pi x/L) \textbackslash right{]},
{]}

в котором коефіцієнти Фурье \(a_n\) и \(b_n\) находятся по формулам
Эйлера

{[} a\_n = \textbackslash frac\{1\}\{L\}
\textbackslash int\_\{-L\}\^{}\{L\} f(x)
\textbackslash cos(n\textbackslash pi x/L) dx = 0, \textbackslash quad n
= 0, 1, 2, \textbackslash dots, {]}

{[} b\_n = \textbackslash frac\{1\}\{L\}
\textbackslash int\_\{-L\}\^{}\{L\} f(x)
\textbackslash sin(n\textbackslash pi x/L) dx = -(-1)\^{}n
(2L/n\textbackslash pi), \textbackslash quad n = 1, 2,
\textbackslash dots. {]}

Эти ингегралы легко вычисляются. Чтобы вывести формулы Эйлера для
коефіцієнтів \(a_n\) и \(b_n\) , умножим обе части разложения (11.1) на
\textbackslash sin\{(n\textbackslash pi x/L)\} или
\textbackslash cos\{(n\textbackslash pi x/L)\} и проинтегрируем получив-
Р И С

\pandocbounded{\includegraphics[keepaspectratio]{images/_page_82_Figure_7.jpeg}}

. 11.1. Разложение пилообразной функції в ряд Фурье.

\begin{verbatim}
шееся соотношение от -L до L. Ортогональность функцій \{\sin(n\pi x/L)\} и \{\cos(n\pi x/L)\} позволяет срязу получить формулы для коефіцієнтів a_n и b_n (см. задачу 6 в конце лекции). Для пилообразной функції разложение в ряд Фурье принимает вид
\end{verbatim}

(11.3)

{[} f(x) = \textbackslash frac\{2L\}\{\textbackslash pi\}
\textbackslash left{[} \textbackslash sin(\textbackslash pi x/L) -
\textbackslash frac\{1\}\{2\} \textbackslash sin(2\textbackslash pi x/L)
+ \textbackslash frac\{1\}\{3\} \textbackslash sin(3\textbackslash pi
x/L) - \textbackslash dots \textbackslash right{]}. {]}

Каждое слагаемое ряда принято называть гармоникой, частота каждой
следующей гармоники выше предыдущей. Частоты всех гармоник кратны
основной частоте, период которой совпадает

\begin{verbatim}
с периодом функції f(x). Один из недостатков рядов Фурье состоит в том, що в ряд Фурье можна разложить лише периодическую функцию. Конечно, якщо функція определена на конечном интервале (наприклад, f(x) = x при 0 \le x \le 1 ), ее можна периодически продолжить на всю числовую ось и затем использовать разложение (11.1).
\end{verbatim}

Тот факт, що сумма ряда Фурье определена и периодична вне интервала
{[}0, 1{]}, ничему не противоречит, поскольку мы интересуемся поведением
функції лише внутри интервала. На самом деле для функції, определенной
внутри интервала, можна указать много различных разложений в ряд Фурье,
поскольку функцию можна продолжить за пределы интервала различными
способами. Сразу же отметим, що полученные ряды будут сходиться к
исходной функції с различными скоростями.

Читатель не повинен думать, що всякую периодическую функцию можна
разложить в ряд Фурье. Пока мы знаем лише, що якщо функцию \(f(x)\)
можна разложить в ряд Фурье (11.1), то коефіцієнти \(a_n\) и \(b_n\)
находятся по формулам Эйлера (11.2). Более того, даже якщо функція
разложима в ряд Фурье, мы не знаем, можна ли ряд почленно
дифференцировать. И в самом деле, легко видеть, що производную функції
\(f(x)\) = x (пилообразной функції) нельзя найти почленным
дифференцированием ряда Фурье (11.3).

В литературе, перечисленной в конце книги, можна найти точные условия
разложимости функції \(f(x)\) в ряд Фурье, а також условия почленной
дифференцируемости получающегося ряда. Нам буде достаточно ознакомиться
с известной теоремой Дирихле.

\emph{Теорема Дирихле} (достаточные условия разложимости функції в ряд
Фурье):

Якщо \(f(x)\)---ограниченная периодическая функція, имеющая на каждом
периоде конечное число максимумов, минимумов и точек разрыва, то ряд
Фурье функції \(f(x)\) сходится к \(f(x)\) в каждой точке непрерывности
функції и к среднему значению пределов функції слева и справа в каждой
точке разрыва функції \(f(x)\).

\begin{verbatim}
Наприклад, на рис. 11.1 ряд Фурье сходится к функції f(x) во всех точках, за исключением точк x=\pm L, \pm 3L, \ldots (точки разрыва). В каждой точке разрыва ряд сходится к нулю (среднее арифметическое значение функції в точках \pm L и \pm L ).
\end{verbatim}

Теперь мы педголовлены для знакомства с преобразованием Фурье. Однако
целесообразно буде познакомиться сначала с понятием частотного спектра
периодической функції.

\section{Дискретный частотный спектр периодической
функції.}\label{ux434ux438ux441ux43aux440ux435ux442ux43dux44bux439-ux447ux430ux441ux442ux43eux442ux43dux44bux439-ux441ux43fux435ux43aux442ux440-ux43fux435ux440ux438ux43eux434ux438ux447ux435ux441ux43aux43eux439-ux444ux443ux43dux43aux446ux456ux457.}

Якщо \(f(x)\) --- периодическая функція, то разложение в ряд Фурье можна
интерпретировать как представление функції \(f(x)\) последовательностью
чисел \{\(c_n\)\} вида

{[} c\_n = \textbackslash sqrt\{a\_n\^{}2 + b\_n\^{}2\},
\textbackslash quad n = 0, 1, 2, \textbackslash dots, {]}

где коефіцієнти \(c_n\) служат мерой вклада соответствующих частотных
компонент в функцию \(f(x)\). Наприклад, пилообразная функція \(f(x)\)
разлагается в ряд Фурье

{[} f(x) = \textbackslash frac\{2L\}\{\textbackslash pi\}
\textbackslash left{[} \textbackslash sin(\textbackslash pi x/L) -
\textbackslash frac\{1\}\{2\} \textbackslash sin(2\textbackslash pi x/L)
+ \textbackslash frac\{1\}\{3\} \textbackslash sin(3\textbackslash pi
x/L) - \textbackslash dots \textbackslash right{]} {]}

и, отже, частотный спектр \{\(c_n\)\} записывается в виде
\(c_n=2L/n\pi\)для n=1,~2,~\textbackslash dots и
\(c_0\)=\textbar{}\(a_0\)\textbar=0 (см. рис. 11.2).

\pandocbounded{\includegraphics[keepaspectratio]{images/_page_84_Figure_8.jpeg}}

PNC. 11 2. Дискретный частотный спектр пилообразной функції a---
основная частота (период 2L); b---n-ая гармоника (период 2L/n).

Последовательность \{\(c_n\)\} напоминает спектральное разложение белого
света по цветам в спектроскове.

Введем теперь понятие преобразования Фурье.

\section{Преобразование
Фурье}\label{ux43fux440ux435ux43eux431ux440ux430ux437ux43eux432ux430ux43dux438ux435-ux444ux443ux440ux44cux435}

\begin{verbatim}
Как уже отмечалось рансе, непериодические функції, заданные на (-\infty, +\infty) , не разлагаются в ряды Фурье Однако для некоторых таких функцій можна построить аналог ряда Фурье. Не вдаваясь в детали доказательства, можна утверждать, що для непериодических функцій, определенных на (-\infty, \infty) , ряд Фурье
\end{verbatim}

{[} f(x\_\{l\} = \textbackslash frac\{a\_\{0\}\}\{2\} +
\textbackslash sum\_\{n=1\}\^{}\{\textbackslash infty\}
\textbackslash left{[} a\_\{n\} \textbackslash cos(n\textbackslash pi
x/L) + b\_\{n\} \textbackslash sin(n\textbackslash pi x/L)
\textbackslash right{]} {]}

переходит в интеграл Фирье (разложение по непрерывному спектру частот)

(11.4)

{[} f(x) = \textbackslash int\_\{0\}\^{}\{\textbackslash infty\}
a(\textbackslash xi) \textbackslash cos(\textbackslash xi x)
d\textbackslash xi +
\textbackslash int\_\{0\}\^{}\{\textbackslash infty\}
b(\textbackslash xi) \textbackslash sin(\textbackslash xi x)
d\textbackslash xi, {]}

гле

(11.5)

{[} a(\textbackslash xi) = \textbackslash frac\{1\}\{\textbackslash pi\}
\textbackslash int\_\{0\}\^{}\{\textbackslash infty\} f(x)
\textbackslash cos(\textbackslash xi x) dx {]}

{[} , b(\textbackslash xi) =
\textbackslash frac\{1\}\{\textbackslash pi\}
\textbackslash int\_\{0\}\^{}\{\textbackslash infty\} f(x)
\textbackslash sin(\textbackslash xi x) dx . {]}

\begin{verbatim}
Из (11.4) видно, що интеграл Фурье осуществляет разложение функції f(x) в спектр по всем частотам 0 < \xi < \infty (а не Р И С
\end{verbatim}

\pandocbounded{\includegraphics[keepaspectratio]{images/_page_85_Figure_7.jpeg}}

11.3. Частотные спектры некоторых функцій.

лише по частотам, кратным основной частоте, как це було у периодических
функцій). По аналогии с рядами Фурье определим частотный спектр по
формуле

{[} C(\textbackslash xi) =
\textbackslash sqrt\{a\^{}2(\textbackslash xi) +
b\^{}2(\textbackslash xi)\}. {]}

Частотный спектр служит мерой вклада частоты 🕇 в функ-

иню \(f(x)\).

\begin{verbatim}
На рис. 11.3 изображены некоторые функції f(x) и их спектры. Из рис. 11.3 видно, що функції, графики которых имеют угловые точки, обладают широким частотным спектром. Объясняется це тем, що для передачи излома графкка требуются высокочастотные спектральные компоненты. С другой стороны, якщо f(x) = \sin \xi_0 x —простая периодическая функція, то ее частотный спектр обращается в нуль всюду, за исключением точки \xi = \xi_0 .
\end{verbatim}

Теперь мы готовы к тому, чтобы познакомиться с экспоненциальной формой
преобразования Фурье (формулы 11.5 известны как синус- и
косинус-преобразования Фурье). Воспользуемся фор-

мулой Эйлера

{[} e\^{}\{i\textbackslash theta\} = \textbackslash cos
\textbackslash theta + i \textbackslash sin \textbackslash theta {]}

и носле несложных преобразований приведем (11.4) к виду

(11.6)

{[} f(x) =
\textbackslash frac\{1\}\{\textbackslash sqrt\{2\textbackslash pi\}\}
\textbackslash int\_\{-\textbackslash infty\}\^{}\{+\textbackslash infty\}
\textbackslash left{[}
\textbackslash frac\{1\}\{\textbackslash sqrt\{2\textbackslash pi\}\}
\textbackslash int\_\{-\textbackslash infty\}\^{}\{+\textbackslash infty\}
f(x) e\^{}\{-i\textbackslash xi x\} dx \textbackslash right{]}
e\^{}\{i\textbackslash xi x\} d\textbackslash xi, {]}

яке принято называть интегральным представлением Фурье. Отсюда получаем
пару інтегральних перетворень

(11.7)

{[} \textbackslash mathcal\{F\}{[}f{]} = F(\textbackslash xi) =
\textbackslash frac\{1\}\{\textbackslash sqrt\{2\textbackslash pi\}\}
\textbackslash int\_\{-\textbackslash infty\}\^{}\{+\textbackslash infty\}
f(x) e\^{}\{-i\textbackslash xi x\} dx {]}

(прямое преобразование Фурье),

{[} \textbackslash mathcal\{F\}\^{}\{-1\}{[}F{]} = f(x) =
\textbackslash frac\{1\}\{\textbackslash sqrt\{2\textbackslash pi\}\}
\textbackslash int\_\{-\textbackslash infty\}\^{}\{+\textbackslash infty\}
F(\textbackslash xi) e\^{}\{i\textbackslash xi x\} d\textbackslash xi
{]}

(обратное преобразование Фурье).

В следующей лекции мы остановимся на изучении свойств цих
преобразований.

\section{ЗАМЕЧАНИЯ}\label{ux437ux430ux43cux435ux447ux430ux43dux438ux44f}

\begin{verbatim}
1. Фурье-образ F(\xi) функції f(x) может оказаться комплекснозначной функцией. Наприклад, для функції
\end{verbatim}

{[} f(x) = {]}

{[} \textbackslash begin\{cases\} 0, \& x \textbackslash le 0,
\textbackslash{} e\^{}\{-x\}, \& x \textgreater{} 0,
\textbackslash end\{cases\} {]}

Фур'є-образ має вигляд

{[} F(\textbackslash xi) =
\textbackslash frac\{1\}\{\textbackslash sqrt\{2\textbackslash pi\}\}
\textbackslash frac\{1 - i\textbackslash xi\}\{1 +
\textbackslash xi\^{}2\}. {]}

\begin{enumerate}
\def\labelenumi{\arabic{enumi}.}
\setcounter{enumi}{1}
\tightlist
\item
  Абсолютне значення зображення Фур'є \(|F(\xi)|\) називається частотним
  спектром функції \(f(x)\). Наприклад, частотний спектр функції
  \(f(x)\) з попереднього зауваження визначається формулою
\end{enumerate}

\textbar F(\textbackslash xi)\textbar{} =
\textbackslash sqrt\{\textbackslash frac\{1\}\{2\textbackslash pi(1+\textbackslash xi\^{}2)\}\}

(зчитувач повинен мати змогу знайти модуль комплексного числа).

\begin{enumerate}
\def\labelenumi{\arabic{enumi}.}
\setcounter{enumi}{2}
\tightlist
\item
  Не всі функції піддаються перетворенню Фур'є (інтеграл (11.7) може не
  існувати). Наприклад, функції \(f(x)\) = c, \(\sin x\) , \(e^x\) , не
  мають \(x^2\) образу Фур'є. Лише ті функції, які мають тенденцію
  швидко нульовою на \(|x| \to \infty\) 1). У прикладних задачах ми
  застосовуємо перетворення Фур'є до температурних, хвильових та інших
  фізичних полів, які мають тенденцію до нуля на \(|x| \to \infty\) .
\end{enumerate}

\section{ЗАВДАННЯ}\label{ux437ux430ux432ux434ux430ux43dux43dux44f-7}

\begin{enumerate}
\def\labelenumi{\arabic{enumi}.}
\tightlist
\item
  Розташуйте функцію у рядку Фур'є над синусами
\end{enumerate}

{[} f(x) = {]}

{[} \textbackslash begin\{cases\} -1, \& -1 \textless{} x \textless{} 0,
\textbackslash{} 1, \& 0 \textbackslash le x \textless{} 1,
\textbackslash end\{cases\} {]}

{[} f(x+2) = f(x) {]}

(умова періодичності).

Побудуйте суму перших 2, 3 і 4 членів, щоб слідувати збіжності ряду до
\(f(x)\). Побудуйте графіки частотного спектра функції \(f(x)\).

\begin{itemize}
\tightlist
\item
  \begin{enumerate}
  \def\labelenumi{\arabic{enumi}.}
  \setcounter{enumi}{1}
  \tightlist
  \item
    Покажіть, що якщо розрізнити ряд Фур'є (11.3) функції пилки за
    членом, отримаємо послідовність, яка не збігається до похідної
    пилоподібної функції.
  \end{enumerate}
\item
  \begin{enumerate}
  \def\labelenumi{\arabic{enumi}.}
  \setcounter{enumi}{2}
  \tightlist
  \item
    Побудуйте частотний спектр наступних періодичних функцій:
  \end{enumerate}
\item
  \begin{enumerate}
  \def\labelenumi{\alph{enumi})}
  \tightlist
  \item
    \(f(x) = \sin x\) ,
  \end{enumerate}
\item
  \begin{enumerate}
  \def\labelenumi{\arabic{enumi})}
  \setcounter{enumi}{5}
  \tightlist
  \item
    \(f(x) = \sin x + \cos 2x,\) c)
    \(f(x) = \sin x + \cos x + 0.5 \sin 3x\) .
  \end{enumerate}
\end{itemize}

\begin{enumerate}
\def\labelenumi{\arabic{enumi}.}
\setcounter{enumi}{3}
\tightlist
\item
  Знайти \(F(\xi)\) зображення Фур'є та частотний спектр
  \(C(\xi) = |F(\xi)|\) функцію
\end{enumerate}

{[} f(x) = {]}

{[} \textbackslash begin\{cases\} 1, \& -1 \textless{} x \textless{} 1,
\textbackslash{} 0 \& \textbackslash text\{в остальных точках.\}
\textbackslash end\{cases\} {]}

Тут ми говоримо лише про перетворення Власіка-Фур'є. Ед.

\begin{verbatim}
5. Показати, що абсолютне значення функції F(\xi) = \frac{1}{1+i\xi} дорівнює
\end{verbatim}

\textbar F(\textbackslash xi)\textbar{} =
\textbackslash sqrt\{\textbackslash frac\{1\}\{1+\textbackslash xi\^{}2\}\}.

ПРИМІТКА. Спочатку помножте чисельник і знаменник на
1-i\textbackslash xi , а потім позбавтеся комплексної одиниці i у
знаменнику. \^{}1 )

\begin{verbatim}
6. Перевірити ортогональність функцій \{\sin(n\pi x/L)\} та \{\cos(n\pi x/L)\} на сегменті [-L, L]:
\end{verbatim}

{[} \textbackslash int\_\{-L\}\^{}\{L\}
\textbackslash sin(m\textbackslash pi x/L)
\textbackslash sin(n\textbackslash pi x/L) dx = {]}

{[} \textbackslash begin\{cases\} 0, \& m \textbackslash neq n,
\textbackslash{} L, \&m = n, \textbackslash end\{cases\} {]}

{[} \textbackslash int\_\{-L\}\^{}\{L\}
\textbackslash cos(m\textbackslash pi x/L)
\textbackslash cos(n\textbackslash pi x/L) dx = {]}

{[} \textbackslash begin\{cases\} 0, \& m \textbackslash neq n,
\textbackslash{} L, \&m = n, \textbackslash end\{cases\} {]}

{[} \textbackslash int\_\{-L\}\^{}\{L\}
\textbackslash sin(m\textbackslash pi x/L)
\textbackslash cos(n\textbackslash pi x/L) dx = 0 \textbackslash text\{
npn Bcex \} m, n = 1, 2, \textbackslash dots {]}

\section{Лекція 12}\label{ux43bux435ux43aux446ux456ux44f-12}

\section{ПЕРЕТВОРЕННЯ ФУР'Є ТА ЙОГО ЗАСТОСУВАННЯ ДО РОЗВ'ЯЗАННЯ
ДИФЕРЕНЦІАЛЬНИХ РІВНЯНЬ У ПОХІДНИХ
ПОХІДНИХ}\label{ux43fux435ux440ux435ux442ux432ux43eux440ux435ux43dux43dux44f-ux444ux443ux440ux454-ux442ux430-ux439ux43eux433ux43e-ux437ux430ux441ux442ux43eux441ux443ux432ux430ux43dux43dux44f-ux434ux43e-ux440ux43eux437ux432ux44fux437ux430ux43dux43dux44f-ux434ux438ux444ux435ux440ux435ux43dux446ux456ux430ux43bux44cux43dux438ux445-ux440ux456ux432ux43dux44fux43dux44c-ux443-ux43fux43eux445ux456ux434ux43dux438ux445-ux43fux43eux445ux456ux434ux43dux438ux445}

МЕТА ЛЕКЦІЇ: Проілюструвати важливі властивості перетворення Фур'є та
показати, як ці властивості можна використовувати для розв'язання PPP.
Зокрема, щоб показати, як під впливом перетворення Фур'є операція
диференціації замінюється множенням, а диференціальні рівняння
перетворюються на алгебранхічні рівняння. У цій лекції вводиться поняття
згортки.

\begin{verbatim}
<sup>2)</sup> Легше |F(\xi)| = \frac{|1|}{|1+i\xi|} = \frac{1}{\sqrt{1+\xi^2}} - \Pi_{PVM}, pcd. Якщо функція f(x) визначена в точці -\infty < x < \infty , то її чотиризображення визначається формулою
\end{verbatim}

(12.1)

{[} \textbackslash mathcal\{F\}{[}f{]} = F(\textbackslash xi) =
\textbackslash frac\{1\}\{\textbackslash sqrt\{2\textbackslash pi\}\}
\textbackslash int\_\{-\textbackslash infty\}\^{}\{+\textbackslash infty\}
f(x) e\^{}\{-i\textbackslash xi x\} dx. {]}

Отже, якщо у нас є функція \(f(x)\), визначена на дійсній осі, ми
підставляємо її у формулу (12.1) і отримуємо нову функцію,
F(\textbackslash xi) визначену для -\textbackslash infty \textless{}
\textbackslash xi \textless{} \textbackslash infty . Водночас не слід
забувати про конвенції існування раннього перетворення Фур'є (у цьому
випадку класичного).

\begin{verbatim}
Таблиця 1. 12.1 показує деякі функції та їхні фур'є-образи. Більш розгорнута таблиця функцій та їхніх фур'є-образів наведена в додатку. З наведеної вище таблиці видно, що образ Фур'є функції може виявитися комплекснозначною функцією аргументу \xi . У першому прикладі цієї таблиці функція F(\xi) містить як множник уявну одиницю i, і таким чином є комплекснозначною функцією екзистенційного аргументу ( \xi змінюється від -\infty до +\infty ). Іншими словами, аргумент \xi є реальним, а значення функції є складними.
\end{verbatim}

Перетворення Фур'є (як і більшість інших інтегральних перетворень)
виявилося дуже корисним, оскільки воно замінює операцію диференціації на
операцію множення і таким чином перетворює диференціальні рівняння на
алгебранхічні рівняння. Існує багато інших властивостей, які роблять
перетворення Фур'є зручним інструментом для розв'язання задач. Тут ми
перелічимо лише найважливіші з них.

Таблиця 12.1. Перетворення Фур'є деяких функцій

\begin{verbatim}
| Функція / (a) | \Gamma Перетворення Фур'є F ( \xi ) |
\end{verbatim}

\textbar----------------------------------------------------------------------------------\textbar-------------------------------------------------------------------------------------------\textbar{}

1.

{[} f(x) = \textbackslash begin\{cases\} e\^{}\{-x\}, \& x
\textbackslash ge 0, \textbackslash{} -e\^{}\{-x\}, \& x \textless{} 0.
\textbackslash end\{cases\} {]}

\begin{verbatim}
| F(\xi) = -i \sqrt{\frac{2}{\pi}} \frac{\xi}{1 + \xi^2} (Xomnekches funktsia). |
\end{verbatim}

2.

{[} f(x) = \textbackslash begin\{cases\} 1, \& -1 \textless{} x
\textless{} 1, \textbackslash{} 0, \& x \textbackslash le -1, x
\textbackslash ge 1 \textbackslash end\{cases\} {]}

\begin{verbatim}
| F(\xi) = \sqrt{\frac{\frac{2}{\pi}}{\frac{\sin \xi}{\xi}}} (Beg.ectronnal (yul.ung). |

| 3. f(x) = e^{-x^2} | F\left(\xi\right)=\frac{1}{\sqrt{2}}e^{-\sqrt{\xi}/2\right)^2} (основна синя функція). |
\end{verbatim}

\section{Властивості перетворення Фур'є Властивість 1 (пара перетворень
Фур'є)}\label{ux432ux43bux430ux441ux442ux438ux432ux43eux441ux442ux456-ux43fux435ux440ux435ux442ux432ux43eux440ux435ux43dux43dux44f-ux444ux443ux440ux454-ux432ux43bux430ux441ux442ux438ux432ux456ux441ux442ux44c-1-ux43fux430ux440ux430-ux43fux435ux440ux435ux442ux432ux43eux440ux435ux43dux44c-ux444ux443ux440ux454}

\begin{verbatim}
Після застосування перетворення Фур'є до функції f(x), визначеної для -\infty < x < \infty , отримуємо нову функцію F(\xi) за формулою
\end{verbatim}

{[} \textbackslash mathscr\{F\}{[}f{]} = F(\textbackslash xi) =
\textbackslash frac\{1\}\{\textbackslash sqrt\{2\textbackslash pi\}\}
\textbackslash int\_\{-\textbackslash infty\}\^{}\{+\textbackslash infty\}
f(x) e\^{}\{-i\textbackslash xi x\} dx. {]}

Якщо застосувати до функції, F(\textbackslash xi) визначеної для
-\textbackslash infty \textless{} \textbackslash xi \textless{}
\textbackslash infty , обернене перетворення Фур'є

{[} \textbackslash mathscr\{F\}\^{}\{-1\}{[}F{]} = f(x) =
\textbackslash frac\{1\}\{\textbackslash sqrt\{2\textbackslash pi\}\}
\textbackslash int\_\{-\textbackslash infty\}\^{}\{+\textbackslash infty\}
F(\textbackslash xi) e\^{}\{i\textbackslash xi x\} d\textbackslash xi,
{]}

\textbf{Тоді} отримуємо початкову функцію \(f(x)\).

\pandocbounded{\includegraphics[keepaspectratio]{images/_page_90_Figure_7.jpeg}}

Рис. 12.1. Графи функції та її фур'є-образу.

Наприклад, (див. рис. 12.1)

{[} e\^{}\{-\textbar x\textbar\}
\textbackslash xrightarrow\{\textbackslash mathcal\{F\}\}
\textbackslash sqrt\{\textbackslash frac\{2\}\{\textbackslash pi\}\}
\textbackslash xrightarrow\{1 + \textbackslash xi\^{}2\}
\textbackslash xrightarrow\{\textbackslash mathcal\{F\}\^{}\{-1\}\}
e\^{}\{-\textbar x\textbar\}. {]}

\section{Властивість 2 (Перетворення
лінійності)}\label{ux432ux43bux430ux441ux442ux438ux432ux456ux441ux442ux44c-2-ux43fux435ux440ux435ux442ux432ux43eux440ux435ux43dux43dux44f-ux43bux456ux43dux456ux439ux43dux43eux441ux442ux456}

Перетворення Фур'є --- це лінійне перетворення, тобто

{[}
\textbackslash mathscr\{F\}{[}af+bg{]}=a\textbackslash mathscr\{F\}{[}f{]}+b\textbackslash mathscr\{F\}{[}g{]}.
{]}

Це легко побачити. Читач може витратити кілька хвилин, щоб ознайомитися
з цією нерухомістю, яка використовується «на кожному кроці». Наприклад,
перетворення Фур'є виразу

{[} e\^{}\{-x\} + 3e\textsuperscript{\{-x}2\} {]}

має вигляд

{[} \textbackslash mathscr\{F\}{[}e\^{}\{-c\}{]} +
3\textbackslash mathscr\{F\}{[}e\textsuperscript{\{-c}*\}{]}. {]}

\section{Властивість 3 (Часткова конвертація
похідних)}\label{ux432ux43bux430ux441ux442ux438ux432ux456ux441ux442ux44c-3-ux447ux430ux441ux442ux43aux43eux432ux430-ux43aux43eux43dux432ux435ux440ux442ux430ux446ux456ux44f-ux43fux43eux445ux456ux434ux43dux438ux445}

\begin{verbatim}
Якщо нас цікавить, як перетворення Фур'є діє на похідні, то перш за все слід звернути увагу на змінні, за якими диференційована функція. Наприклад, якщо перетворення Фур'є виконується на змінній \boldsymbol{x} (тобто змінна \boldsymbol{x} інтегрована у перетворення), і якщо часткова похідна функції \boldsymbol{u}(\boldsymbol{x},t) на тій самій змінній \boldsymbol{x} перетворена, то наступні формули перетворення похідних є дійсними
\end{verbatim}

{[} \textbackslash begin\{split\}
\textbackslash mathcal\{F\}\textbackslash left\{u\_\{x\}\textbackslash right\}
\&=
\textbackslash frac\{1\}\{\textbackslash sqrt\{2\textbackslash pi\}\}
\textbackslash int\_\{-\textbackslash infty\}\^{}\{+\textbackslash infty\}
u\_\{x\}\textbackslash left(x,,t\textbackslash right)
e\^{}\{-i\textbackslash xi x\} dx = i\textbackslash xi
\textbackslash mathcal\{F\}\textbackslash left{[}u\textbackslash right{]},
\textbackslash{}
\textbackslash mathcal\{F\}\textbackslash left\{u\_\{xx\}\textbackslash right\}
\&=
\textbackslash frac\{1\}\{\textbackslash sqrt\{2\textbackslash pi\}\}
\textbackslash int\_\{-\textbackslash infty\}\^{}\{+\textbackslash infty\}
u\_\{xx\}\textbackslash left(x,,t\textbackslash right) ,
\textbackslash delta\^{}\{-i\textbackslash xi x\} dx =
-\textbackslash xi\^{}\{2\}
\textbackslash mathcal\{F\}\textbackslash left{[}u\textbackslash right{]}.
\textbackslash end\{split\} {]}

С другой стороны, если подвергнуть преобразованию частную производную
\(u_t\)(x, t) (а переменная интегрирования в преобразовании, как и
прежде, будет x), получаем следующие формулы:

{[} \textbackslash begin\{aligned\}
\textbackslash mathcal\{F\}\textbackslash left{[}u\_\{t\}\textbackslash right{]}
\&=
\textbackslash frac\{1\}\{\textbackslash sqrt\{2\textbackslash pi\}\}
\textbackslash int\_\{-\textbackslash infty\}\^{}\{+\textbackslash infty\}
u\_\{t\}\textbackslash left(x, ~t\textbackslash right)
e\^{}\{-i\textbackslash frac\{\textbackslash pi\}\{2\}x\} dx =
\textbackslash frac\{\textbackslash partial\}\{\textbackslash partial
t\}
\textbackslash mathcal\{F\}\textbackslash left{[}u\textbackslash right{]},
\textbackslash{}
\textbackslash mathcal\{F\}\textbackslash left{[}u\_\{tt\}\textbackslash right{]}
\&=
\textbackslash frac\{1\}\{\textbackslash sqrt\{2\textbackslash pi\}\}
\textbackslash int\_\{-\textbackslash infty\}\^{}\{+\textbackslash infty\}
u\_\{tt\}\textbackslash left(x, ~t\textbackslash right)
e\^{}\{-i\textbackslash frac\{\textbackslash pi\}\{2\}x\} dx =
\textbackslash frac\{\textbackslash partial\^{}\{2\}\}\{\textbackslash partial
t\^{}\{2\}\}
\textbackslash mathcal\{F\}\textbackslash left{[}u\textbackslash right{]}.
\textbackslash end\{aligned\} {]}

\section{Властивість 4 (Властивість
згортки)}\label{ux432ux43bux430ux441ux442ux438ux432ux456ux441ux442ux44c-4-ux432ux43bux430ux441ux442ux438ux432ux456ux441ux442ux44c-ux437ux433ux43eux440ux442ux43aux438}

Кожне інтегральне перетворення має так звану властивість згортки. Факт у
тому, що образ Фур'є добутку двох функцій загалом не відрізняється від
добутку образів Фур'є цих функцій, тобто

{[} \textbackslash mathcal\{F\}{[}f(x)g(x){]} \textbackslash neq
\textbackslash mathcal\{F\}{[}f{]}\textbackslash mathcal\{F\}{[}g{]}.
{]}

Однак у теорії перетворень існує операція, яка називається згорткою
функцій f і g (позначена як f*g), яка певною мірою виконує роль добутку,
а саме відношення

(12.2)

{[} \textbackslash mathcal\{F\}{[}/*g{]} =
\textbackslash mathcal\{F\}{[}/{]}\textbackslash mathcal\{F\}{[}g{]}.
{]}

Що це за загадкова операція під назвою згортка? Вона визначається
формулою

(12.3)

{[} (f*g)(x) =
\textbackslash frac\{1\}\{\textbackslash sqrt\{2\textbackslash pi\}\}
\textbackslash int\_\{-\textbackslash infty\}\^{}\{+\textbackslash infty\}
f(x-\textbackslash xi) g(\textbackslash xi) d\textbackslash xi. {]}

Неважко показати, що співвідношення (12,2) дійсно виконується. З
визначення видно, що згортка двох функцій \(f(x)\) і \(g(x)\) є новою
функцією (f*g)(x) аргументу x.

\section{Приклад обчислення згортки двох
функцій}\label{ux43fux440ux438ux43aux43bux430ux434-ux43eux431ux447ux438ux441ux43bux435ux43dux43dux44f-ux437ux433ux43eux440ux442ux43aux438-ux434ux432ux43eux445-ux444ux443ux43dux43aux446ux456ux439}

Нехай задано дві функції

{[} f(x) = x, {]}

{[} g(x) = e\textsuperscript{\{-x}2\}. {]}

Тоді їхня згортка дорівнює

{[} (f \textbackslash circ g)(x) =
\textbackslash frac\{1\}\{\textbackslash sqrt\{2\textbackslash pi\}\}
\textbackslash int\_\{-\textbackslash infty\}\^{}\{\textbackslash infty\}
(x - \textbackslash xi) e\textsuperscript{\{-\textbackslash xi}2\}
d\textbackslash xi = x/\textbackslash sqrt\{2\}. {]}

Тут ми використали відому формулу

{[}
\textbackslash int\_\{-\textbackslash infty\}\^{}\{+\textbackslash infty\}
e\textsuperscript{\{-\textbackslash xi}2\} d\textbackslash xi =
\textbackslash sqrt\{\textbackslash pi\}. {]}

Важливість концепції згортки (12.3) полягає в тому, що дуже часто,
розв'язуючи диференціальні рівняння, ми стикаємося з необхідністю
застосовувати обернене перетворення Фур'є до виразу, який є добутком
двох образів Фур'є, тобто потрібно знайти

(12.4)

{[}
\textbackslash mathcal\{F\}\^{}\{-1\}\textbackslash left\{\textbackslash mathcal\{F\}\textbackslash left{[}f\textbackslash right{]}\textbackslash mathcal\{F\}\textbackslash left{[}g\textbackslash right{]}\textbackslash right\}.
{]}

Застосовуючи обернене перетворення до обох частин (12.2), отримуємо

{[} f * g = \textbackslash mathcal\{J\}\^{}\{-1\} \{
\textbackslash mathcal\{F\} {[}f{]} \textbackslash mathcal\{F\} {[}g{]}
\}. {]}

Отже, щоб визначити (12.4), потрібно знайти прототипи кожного з
факторів, тобто функцій f і g, а потім обчислити їхню згортку. Тепер ми
готові розв'язати одну важливу проблему: теорію диференціальних рівнянь
у похідних похідних.

\section{Розв'язання проблеми
Коші}\label{ux440ux43eux437ux432ux44fux437ux430ux43dux43dux44f-ux43fux440ux43eux431ux43bux435ux43cux438-ux43aux43eux448ux456}

Розглянемо проблему поширення тепла в \emph{нескінченному} стрижні, якщо
задано початкову температуру \(u(x, 0) = \varphi(x)\) .

Цю задачу зазвичай називають задачою Коші для рівняння теплопровідності
або задачою початкових умов

{[} \textbackslash begin\{aligned\} (\textbackslash text\{УЧП\}) \& u\_t
= \textbackslash alpha\^{}2 u\_\{xx\}, \& -\textbackslash infty
\textless{} x \textless{} \textbackslash infty, \& 0 \textless{} t
\textless{} \textbackslash infty, \textbackslash{}
(\textbackslash text\{12.6\}) \& (\textbackslash text\{HV\}) \& u(x,0) =
q(x), \& -\textbackslash infty \textless{} x \textless{}
\textbackslash infty. \textbackslash end\{aligned\} {]}

Розв'язок задачи можна разбить на три основных шага.

ШАГ 1 (Преобразование задачи).

Поскольку пространственная переменная x изменяется в пределах от
\(-\infty\) до \(+\infty\) , мы подвергнем уравнение и початкова умова
(12.6) преобразованию Фурье по переменной \(x^{(1)}\) (x---переменная
интегрирования в преобразовании). Сделав це, получим

{[} \textbackslash mathcal\{F\}{[}u\_t{]} = \textbackslash alpha\^{}2
\textbackslash mathcal\{F\}{[}u\_\{xx\}{]}, {]}

{[} \textbackslash mathcal\{F\}{[}u(x, 0){]} =
\textbackslash mathcal\{F\}{[}\textbackslash varphi(x){]}. {]}

Воспользуемся свойствами преобразования Фурье. Тогда

{[} \textbackslash frac\{\textbackslash frac\{dU(t)\}\{dt\} =
-\textbackslash alpha\^{}2 \textbackslash xi\^{}2 U(t), \}\{U(0) =
\textbackslash Phi(\textbackslash xi)\} = (\textbackslash text\{here \}
\textbackslash Phi - \textbackslash phi) \textbackslash text\{rje-image
funcin \} \textbackslash phi), {]}

где \(U(t)\) = \textbackslash mathcal\{F\}{[}\(u(x, t)\){]} . Читате ть
повинен обратить внимание на то, що функція \(U(t)\) на салот доле
зависит не лише от t, но и от \textbackslash xi , но, поскольку в
дифференциальном уравнении (12.7) величина \textbackslash xi играет роль
константы, мы не будем записывать ее в число аргументов, т. е. будем
считать, що U = \(U(t)\).

ШАГ 2 (Розв'язок преобразованной задачи).

Поскольку новая переменная играет роль гонстанты в задаче (12.7), ее
розв'язок легко находится:

(12.8)

{[} U(t) = \textbackslash Phi(\textbackslash xi)
e\textsuperscript{\{-x}2 e\^{}\{tt\}\}. {]}

ШАГ 3 (Нахождение обратного преобразования). Искомое розв'язок
\(u(x, t)\) находится по формуле

{[} u(x, t) = \textbackslash mathcal\{F\}\^{}\{-1\}
{[}U(\textbackslash xi, t){]} = {]}

{[} = \textbackslash mathcal\{F\}\^{}\{-1\}
{[}\textbackslash Phi(\textbackslash xi) e\^{}\{-\textbackslash alpha
\textbackslash xi\^{}2 t\}{]}. {]}

Предполагается, що преобразование Фурье цих функцій существует. Прим.
ред.

Вот здесь нас выручит теорема о свертке (12.5). Используя эту теорему,
получаем

(12.9)

{[} u(x, t) = \textbackslash mathcal\{F\}\^{}\{-1\}
\textbackslash left{[} \textbackslash Phi(\textbackslash xi)
e\^{}\{-\textbackslash alpha r \textbackslash xi\^{}\{2\} t\}
\textbackslash right{]} = {]}

{[} = \textbackslash mathcal\{F\}\^{}\{-1\} \textbackslash left{[}
\textbackslash Phi(\textbackslash xi) \textbackslash right{]} *
\textbackslash mathcal\{F\}\^{}\{-1\} \textbackslash left{[}
e\textsuperscript{\{-\textbackslash alpha}\{2\}
\textbackslash xi\^{}\{2\} t\} \textbackslash right{]} = {]}

{[} = \textbackslash Phi(x) * \textbackslash left{[}
\textbackslash frac\{1\}\{2\textbackslash alpha
\textbackslash sqrt\{\textbackslash pi t\}\}
e\textsuperscript{\{-x}\{2\}/4\textbackslash alpha\^{}\{2\} t\}
\textbackslash right{]} \textbackslash text\{ (знайдено в таблиці)\} =
{]}

{[} = \textbackslash frac\{1\}\{2\textbackslash alpha
\textbackslash sqrt\{\textbackslash pi t\}\}
\textbackslash int\_\{-\textbackslash infty\}\^{}\{+\textbackslash infty\}
\textbackslash Phi(\textbackslash xi)
e\textsuperscript{\{-(x-\textbackslash xi)}\{2\}/4\textbackslash alpha\^{}\{2\}
t\} d\textbackslash xi. {]}

Дело сделано! Формула (12.9) дает розв'язок нашей задачи. Перш ніж чем
поставить точку, давайте проанализируем полученный результат. Отметим,
що почынтегральная функція состои с

из двух сомножителей:

\begin{enumerate}
\def\labelenumi{\arabic{enumi}.}
\item
  Начальной температури \(q(x)\),

  \begin{enumerate}
  \def\labelenumii{\arabic{enumii}.}
  \setcounter{enumii}{1}
  \tightlist
  \item
    Функції G(x, t) = \textbackslash frac\{1\}\{2\textbackslash alpha
    \textbackslash sqrt\{\textbackslash pi t\}\}
    e\textsuperscript{\{-(x-\textbackslash xi)}2/4c\^{}2t\} , яка
    називається функцией Грина или функцией источника.
  \end{enumerate}
\end{enumerate}

Можна показать, що функція источника \(G(x, t)\) описывает отклик
системи на начальный температурный единичный импульс Р И С

\pandocbounded{\includegraphics[keepaspectratio]{images/_page_94_Figure_9.jpeg}}

\begin{verbatim}
. 12.2. Отклик G(x,t) на температурный импульс в точке x=\xi , графчк G(x,t) похож на график нормального распределения, при цьому время игриег роль среднеквадратического уклонения: a—функція G(x,t) при пебольших значениях t, \delta —функція G(x,t) сущестьенно уменьши...ась.

в ючке x=\xi . Другими словами, функція G(\kappa,t) описывает распределение температури в стержие в мочент времени t, якщо на точку x=\xi подействовал единичный тепловой импульс (см. рис. 12.2).

Теперь формуле (12.9) можна дать следующую интерпретацию: начальную температуру u\left(x,\,0\right)=\phi\left(x\right) можна представлять как конининуальное множество точечных импульсов каждой величины q(\xi) (в точке x = \xi ). Каждый точечный импульс дает распределение температури \phi(\xi) G(x, t). Результирующее распределение находится суммированием (интегрированием) температур точечных источников по формуле (12.9). Позже мы увидим, що це одно из проявлений общего принципа суперпозиции.
\end{verbatim}

\section{ЗАМЕЧАНИЕ}\label{ux437ux430ux43cux435ux447ux430ux43dux438ux435}

Главный недостаток преобразования Фурье в том, що его можна применять не
ко всем функциям. Наприклад, функції вида

{[} f(x) = \textbackslash text\{const\}, {]}

{[} f(x) = e\^{}x, {]}

{[} f(x) = \textbackslash sin x {]}

не имеют фурье-образа, так как интеграт

{[} \textbackslash mathcal\{F\}{[}f{]} =
\textbackslash frac\{1\}\{\textbackslash sqrt\{2\textbackslash pi\}\}
\textbackslash int\_\{-\textbackslash pi\}\^{}\{\textbackslash pi\} f(x)
e\^{}\{-i\textbackslash xi x\} dx {]}

расходится. Фурье-образ имеют лише те функції, які достаточно быстро
сгремятся к нулю при \textbar x\textbar{} \textbackslash to
\textbackslash infty . Кроме того, нельзя применять преобразование Фурье
по времени для розв'язки задач с начальными условиями, так как 0
\textless{} t \textless{} \textbackslash infty (якщо не рассматривать
продолжения функцій на полуось в соответствующем классе функцій).

\section{ЗАДАЧИ}\label{ux437ux430ux434ux430ux447ux438-1}

\begin{enumerate}
\def\labelenumi{\arabic{enumi}.}
\tightlist
\item
  Найдите фурье-образ функції
\end{enumerate}

{[} f(x) = {]}

{[} \textbackslash begin\{cases\} 0, \& x \textbackslash leq 0,
\textbackslash{} e\^{}\{-x\}, \& 0 \textless{} x,
\textbackslash end\{cases\} {]}

используя таблицы, приведенные в приложении.

\begin{itemize}
\tightlist
\item
  \begin{enumerate}
  \def\labelenumi{\arabic{enumi}.}
  \setcounter{enumi}{1}
  \tightlist
  \item
    Убедитесь в 10м, що прямое и обратное преобразования Фурье являются
    линейными преобразованиями.
  \end{enumerate}
\item
  \begin{enumerate}
  \def\labelenumi{\arabic{enumi}.}
  \setcounter{enumi}{2}
  \tightlist
  \item
    Используя преобразование Фурье, решите задачу Коши
  \end{enumerate}
\end{itemize}

{[} \textbackslash begin\{aligned\} (\textbackslash text\{УЧП\}) \& u\_t
= \textbackslash alpha\^{}2 u\_\{xx\}, \& -\textbackslash infty
\textless{} x \textless{} \textbackslash infty, \textbackslash{}
(\textbackslash text\{HV\}) \& u(x, 0) = e\textsuperscript{\{-x}2\}, \&
-\textbackslash infty \textless{} x \textless{} \textbackslash infty.
\textbackslash end\{aligned\} {]}

\begin{enumerate}
\def\labelenumi{\arabic{enumi}.}
\setcounter{enumi}{3}
\tightlist
\item
  Проверьте справедливость соотношений
\end{enumerate}

{[} \textbackslash begin\{aligned\}
\textbackslash mathcal\{F\}\textbackslash left{[}u\_\{x\}\textbackslash right{]}
\&= i\textbackslash xi
\textbackslash mathcal\{F\}\textbackslash left{[}u\textbackslash right{]},
\textbackslash{}
\textbackslash mathcal\{F\}\textbackslash left{[}u\_\{xx\}\textbackslash right{]}
\&= -\textbackslash xi\^{}\{2\}
\textbackslash mathcal\{F\}\textbackslash left{[}u\textbackslash right{]}.
\textbackslash end\{aligned\} {]}

УКАЗАНИЕ, Воспользуйтесь интегрированием по частям.

\begin{enumerate}
\def\labelenumi{\arabic{enumi}.}
\setcounter{enumi}{4}
\tightlist
\item
  Проверьте, що свертку двух функцій f и g можна записать в двух
  эквивалентных формах
\end{enumerate}

(чорт)(x) =
\textbackslash frac\{1\}\{\textbackslash sqrt\{2\textbackslash pi\}\}
\textbackslash int\_\{-\textbackslash infty\}\^{}\{+\textbackslash infty\}
f(\textbackslash xi) g(x - \textbackslash xi) d\textbackslash xi

или

(чорт)(x) =
\textbackslash frac\{1\}\{\textbackslash sqrt\{2\textbackslash pi\}\}
\textbackslash int\_\{-\textbackslash infty\}\^{}\{+\textbackslash infty\}
f(x-\textbackslash xi) g(\textbackslash xi) d\textbackslash xi.

Лекція 13 \# ПРЕОБРАЗОВАНИЕ ЛАПЛАСА

ЦЕЛЬ ЛЕКЦИИ: Ввести понятие преобразования Лапласа

{[} \textbackslash mathscr\{L\}{[}f{]} = F(s) =
\textbackslash int\textbackslash limits\_0\^{}\textbackslash infty
f(t),e\^{}\{-st\}dt {]}

(преобразование Лапласа),

познакомиться с формулой обращения преобразования Лапласа

{[} \textbackslash mathscr\{L\}\^{}\{-1\}{[}F{]} = f(t) =
\textbackslash frac\{1\}\{2\textbackslash pi t\}
\textbackslash int\_\{c-t\textbackslash infty\}\^{}\{c+t\textbackslash infty\}
F(s) e\^{}\{st\} ds {]}

(обратное преобразование Лапласа)

и проиллюстрировать их свойства.

Преобразование Лапласа имеет важное преимущество перед преобразованием
Фурье, так как в подынтегральное вираз входит быстро затухающий
множитель \(e^{-st}\) , и поэтому преобразование Лапласа применимо к
более широкому классу функцій по сравнению с преобразованием Фурье.
Кроме того, преобразование Лапласа имеет дело с функциями, определенными
на промежутке \([0, \infty)\) , а значит, особенно удобно проводить его
по переменной t.

Після знакомства с основными свойствами преобразования Лапласа мы решим
несколько важных задач для уравнений з частинними похідними 113 всех
інтегральних перетворень, рассматриваемых в цієї книге, преобразование
Лапласа

(13 1)

{[} \textbackslash mathscr\{Z\}{[}f{]} =
\textbackslash int\_\{0\}\^{}\{\textbackslash infty\} f(t) e\^{}\{-st\}
dt, {]}

по-видимому, является единственным, с которым читатель уже встречался
ранее, поскольку це один из мощных методов преобразования обыкновенных
дифференциальных уравнений в алгебранческие. Более того, преобразование
Лапласа позволяет преобразовывать уравнение з частинними похідними в
обыкновенное дифференциальное уравнение и его можна проводить по любой
из тех змінних \(x, y, z, t \ldots\) , які изменяются от 0 до \(\infty\)
(чаще всего преобразование Лапласа проводят по переменной t).

Якщо при применении преобразования Лапласа к обыкновенному
дифференциальному уравшению получаются алгебранческие уравнения, то в
случае уравнений з частинними похідними преобразование Лапласа приводит
либо к новому уравнению з частинними похідними с меньшим числом змінних,
либо к обыкновенному дифференциальному уравнению Після цього возникает
вопрос о выборе метода розв'язки преобразование, или решать ее методом
разделения змінних и т д.). Перш ніж чем приступить к решению очень
интересных задач, перечислим некоторые свойства преобразования Лапласа

\section{Свойства преобразования
Лапласа}\label{ux441ux432ux43eux439ux441ux442ux432ux430-ux43fux440ux435ux43eux431ux440ux430ux437ux43eux432ux430ux43dux438ux44f-ux43bux430ux43fux43bux430ux441ux430}

Свойство 1 (Прямое и обратное преобразования)

Преобразование Лапласа и формула его обращения имеют вид

{[} \textbackslash mathscr\{L\}{[}f{]} = F(s) =
\textbackslash int\_\{0\}\^{}\{\textbackslash infty\} f(t) ,
e\^{}\{-st\} , dt \textbackslash quad \textbackslash text\{(Перетворення
Лапласа),\} {]}

(13.2)

{[} \textbackslash mathscr\{L\}\^{}\{-1\}{[}F{]} = f(t) =
\textbackslash frac\{1\}\{2\textbackslash pi i\}
\textbackslash int\_\{c-i\textbackslash infty\}\^{}\{c+i\textbackslash infty\}
F(s) , e\^{}\{st\} , ds \textbackslash quad
\textbackslash text\{(формула інверсії).\} {]}

Преобразование Лапласа обладает важным преимуществом по сравнению с
преобразованием Фурье. Це преимущество обусловлено наличием затухающего
множителя \(e^{-st}\) в подынтегральном выражении, що обеспечивает
возможность применения преобразования Лапласа к более широкому классу
функцій (множитель \(e^{t\xi t}\) в преобразовании Фурье не является
затухающим, поскольку модуль цього множителя дорівнює единице). Более
точное представление об цьому классе функцій дает следующая теорема:

\section{Достаточные условия существования преобразования
Лапласа}\label{ux434ux43eux441ux442ux430ux442ux43eux447ux43dux44bux435-ux443ux441ux43bux43eux432ux438ux44f-ux441ux443ux449ux435ux441ux442ux432ux43eux432ux430ux43dux438ux44f-ux43fux440ux435ux43eux431ux440ux430ux437ux43eux432ux430ux43dux438ux44f-ux43bux430ux43fux43bux430ux441ux430}

Ести

\begin{enumerate}
\def\labelenumi{\arabic{enumi})}
\item
  для любого положительного числа A функція \(f(t)\) кусочнонепрерывна
  на интервале \(0 \le t \le A\) ;
\item
  существуют такие константы M, a и T, що \(|f(t)| \leq Me^{at}\) для
  всех t \textgreater{} T, то преобразование Лапласа
\end{enumerate}

{[} \textbackslash mathscr\{L\}{[}f(\textbackslash ell){]} = F(s) =
\textbackslash int\_\{0\}\^{}\{\textbackslash infty\}
f(\textbackslash ell) e\^{}\{-st\} dt {]}

существует при всех s \textgreater{} a.

Приведем примеры некоторых функцій и их преобразований Лапласа

В определении преобразования Лапласа подразумевается, що
s---действительная переменная, принадлежащая промежутку \((0, \infty)\)
. Однако можна (а зачастую желательно) распространить це определение на
комплексные значения s. Це позволяет написать формулу обращения
преобразований Лапласа

{[} \textbackslash mathscr\{L\}\^{}\{-1\}{[}F{]} =
\textbackslash frac\{1\}\{2\textbackslash pi i\}
\textbackslash int\_\{\textbackslash epsilon -
i\textbackslash infty\}\^{}\{\textbackslash epsilon +
i\textbackslash infty\} F(s) e\^{}\{st\} ds. {]}

Использование цієї формулы подразумевает знакомство с техникой
интегрирования по контурам в комплексной плоскости и теорией вычетов Мы
не станем отвлекаться на изучение здесь цих вопросов, а при обращении
преобразования Лапласа будем пользоваться таблицей, приведенной в
приложении

\section{Свойство 2 (Преобразования частинних
похідних)}\label{ux441ux432ux43eux439ux441ux442ux432ux43e-2-ux43fux440ux435ux43eux431ux440ux430ux437ux43eux432ux430ux43dux438ux44f-ux447ux430ux441ux442ux438ux43dux43dux438ux445-ux43fux43eux445ux456ux434ux43dux438ux445}

Пусть \(u(x, t)\) --- функція змінних, и мы хотим найти результаты
применения преобразования Лапласа к различным частным производным
\(u_t\) , \(u_{tt}\) , \(u_x\) , \(u_{\lambda x}\) . . . . Поскольку
преобразование Лапласа проводится по переменной t (t --- переменная
интегрирования), правила преобразования имеют вид

{[}
\textbackslash mathcal\{L\}\textbackslash left{[}u\_\{t\}\textbackslash right{]}
= \textbackslash int\_\{0\}\^{}\{\textbackslash infty\} u\_\{t\}(x, t)
e\^{}\{-st\} dt = sU(x, s) - u(x, 0), {]}

{[}
\textbackslash mathcal\{L\}\textbackslash left{[}u\_\{tt\}\textbackslash right{]}
= \textbackslash int\_\{0\}\^{}\{\textbackslash infty\} u\_\{tt\}(x, t)
e\^{}\{-st\} dt = s\^{}\{2\}U(x, s) - su(x, 0) - u\_\{t\}(x, 0), {]}

{[}
\textbackslash mathcal\{L\}\textbackslash left{[}u\_\{x\}\textbackslash right{]}
= \textbackslash int\_\{0\}\^{}\{\textbackslash infty\} u\_\{x\}(x, t)
e\^{}\{-st\} dt = \textbackslash frac\{\textbackslash partial
U\}\{\textbackslash partial x\}(x, s), {]}

{[}
\textbackslash mathcal\{L\}\textbackslash left{[}u\_\{xx\}\textbackslash right{]}
= \textbackslash int\_\{0\}\^{}\{\textbackslash infty\} u\_\{xx\}(x, s)
e\^{}\{-st\} dt =
\textbackslash frac\{\textbackslash partial\^{}\{2\}U\}\{\textbackslash partial
x\^{}\{2\}\}(x, s), {]}

\begin{enumerate}
\def\labelenumi{\arabic{enumi}.}
\tightlist
\item
\end{enumerate}

{[} f(t) = 1 {]}

{[} , 0 \textless{} t \textless{} \textbackslash infty (M = 1, a = 0)
{]}

\pandocbounded{\includegraphics[keepaspectratio]{images/_page_99_Figure_3.jpeg}}

{[} F(s) = \textbackslash int\_\{0\}\^{}\{\textbackslash infty\}
e\^{}\{-st\} dt = \textbackslash frac\{1\}\{s\} {]}

\pandocbounded{\includegraphics[keepaspectratio]{images/_page_99_Figure_5.jpeg}}

2

{[} f(t) = e\^{}\{2t\} {]}

{[} , 0 \textless{} t \textless{} \textbackslash infty (M = 1, n = 2 )
{]}

\pandocbounded{\includegraphics[keepaspectratio]{images/_page_99_Figure_7.jpeg}}

{[} F(s) = \textbackslash frac\{1\}\{s-2\} {]}

{[} , s \textgreater{} 2 a {]}

\pandocbounded{\includegraphics[keepaspectratio]{images/_page_99_Figure_9.jpeg}}

\begin{enumerate}
\def\labelenumi{\arabic{enumi}.}
\setcounter{enumi}{2}
\tightlist
\item
\end{enumerate}

{[} f(t) = \textbackslash sin(\textbackslash omega t) {]}

{[} (M=1, \textbackslash alpha=0) {]}

\pandocbounded{\includegraphics[keepaspectratio]{images/_page_99_Figure_11.jpeg}}

\pandocbounded{\includegraphics[keepaspectratio]{images/_page_99_Figure_12.jpeg}}

\pandocbounded{\includegraphics[keepaspectratio]{images/_page_99_Figure_13.jpeg}}

\begin{enumerate}
\def\labelenumi{\arabic{enumi}.}
\setcounter{enumi}{3}
\tightlist
\item
  \(f(t)\) = \(e^{t^2}\) (не имеет преобразования Лапласа)
\end{enumerate}

Рис. 13.1. Графики некоторых функцій и их изображений по Лапласу.

\begin{verbatim}
где U\left(x,\,s\right)=\mathcal{L}\left[u\left(x,\,t\right)\right] . Правила преобразования похідних u_{x} и u_{xy} вытекают из теоремы о дифференцировании интеграла по параметру
\end{verbatim}

{[}
\textbackslash frac\{d\}\{dx\}\textbackslash int\_\{a\}\textsuperscript{\{b\}f\textbackslash left(x,~y\textbackslash right)dy=\textbackslash int\_\{a\}}\{b\}\textbackslash frac\{df\}\{dx\}\textbackslash left(x,~y\textbackslash right)dy,
{]}

а формулы для преобразования похідних \(u_t\) и \(u_{tt}\) можна
получить интегрированием по частям.

\section{Свойство 3 (Теорема о
свертке)}\label{ux441ux432ux43eux439ux441ux442ux432ux43e-3-ux442ux435ux43eux440ux435ux43cux430-ux43e-ux441ux432ux435ux440ux442ux43aux435}

Свертка играет здесь ту же роль, що и в теории преобразования Фурье, но
вводится несколько иначе.

\section{Определение конечной
свертки}\label{ux43eux43fux440ux435ux434ux435ux43bux435ux43dux438ux435-ux43aux43eux43dux435ux447ux43dux43eux439-ux441ux432ux435ux440ux442ux43aux438}

Конечная свертка двух функцій f и g определяется формулами

{[} (f * g) (t) = {]}

{[} \textbackslash begin\{cases\} \textbackslash int\_0\^{}t
f(\textbackslash tau) g(t-\textbackslash tau) d\textbackslash tau,
\textbackslash{} \textbackslash int\_0\^{}t f(t-\textbackslash tau)
g(\textbackslash tau) d\textbackslash tau \textbackslash end\{cases\}
{]}

(Можна довести, що ці інтеграли рівні). Тут інтегрування здійснюється
вздовж скінченного інтервалу від 0 до t, а не вздовж нескінченної (
\(\infty\) , \(+\infty\) ), як це було раніше з визначенням нескінченної
згортки.

Ось приклад скінченної згортки двох функцій. Нехай

{[} f(t) = t, \textbackslash quad g(t) = t; {]}

тоді

{[} (f * g)(t) = \textbackslash int\_\{0\}\^{}\{t\} \textbackslash tau(t
- \textbackslash tau) d\textbackslash tau = t\^{}\{3\}/6. {]}

Як і у випадку нескінченної згортки, її найважливіша властивість
визначається формулою

(13.4)

{[} \textbackslash mathscr\{L\}{[}f * g{]} =
\textbackslash mathscr\{L\}{[}f{]} \textbackslash mathscr\{L\}{[}g{]}
{]}

або еквівалентна формула

(13.5)

{[} \textbackslash mathscr\{L\}\^{}\{-1\}\textbackslash left\{
\textbackslash left{[}\textbackslash mathscr\{L\}\textbackslash left{[}f\textbackslash right{]}\textbackslash mathscr\{L\}\textbackslash left{[}g\textbackslash right{]}\textbackslash right\}
= f*g. \textbackslash right. {]}

Ці формули дозволяють змінити перетворення Лапласа у випадках, коли
образ можна представити як добуток \(\mathscr{L}[f]\) і
\(\mathscr{L}[g]\) таких множників, для яких

Функції F і G легко знайти. Після пошуку функцій f і g можна знайти їхню
згортку. Наприклад,

{[}
\textbackslash mathcal\{Z\}\^{}\{-1\}\textbackslash left{[}\textbackslash frac\{1\}\{s\}
\textbackslash cdot
\textbackslash frac\{1\}\{s\^{}2+1\}\textbackslash right{]} =
\textbackslash int\_0\^{}t \textbackslash sin \textbackslash tau ,
d\textbackslash tau = 1 - \textbackslash cos t, {]}

{[} F(s) = \textbackslash frac\{1\}\{s\}
\textbackslash xrightarrow\{\textbackslash mathcal\{Z\}\^{}\{-1\}\} f(t)
= 1 \textbackslash quad G(s)
\textbackslash xrightarrow\{\textbackslash mathcal\{Z\}\^{}\{-1\}\} g(t)
= \textbackslash sin t. {]}

Тепер ми готові до важливого змішаного виклику.

\section{Теплопровідність у напівнескінченному
середовищі}\label{ux442ux435ux43fux43bux43eux43fux440ux43eux432ux456ux434ux43dux456ux441ux442ux44c-ux443-ux43dux430ux43fux456ux432ux43dux435ux441ux43aux456ux43dux447ux435ux43dux43dux43eux43cux443-ux441ux435ux440ux435ux434ux43eux432ux438ux449ux456}

Розглянемо глибокий резервуар з рідиною, і нехай бічна поверхня
резервуара буде теплоізольованою. Нехай \(u_0\) буде початковою
температурою рідини, а температура повітря над рідиною ---

\pandocbounded{\includegraphics[keepaspectratio]{images/_page_101_Picture_7.jpeg}}

Рис. 13.2. Схематичне представлення проблеми теплопровідності: a ---
тепло тече всередину, якщо ви(0, t) \textless{} 0, і назовні, якщо ви(0,
t) \textgreater{} 0; b --- резервуар настільки глибокий, що гранична
умова внизу не впливає на розв'язок при значеннях x.

кістка дорівнює нулю (орієнтир температур) (див. рис. 13.2). Наше
завдання --- знайти температурне поле в рідині на різних глибинах і в
різний час. Іншими словами, нам потрібно розв'язати проблему

(13.6)

{[} \textbackslash begin\{array\}\{lll\} (\textbackslash text\{УЧП\}) \&
u\_t = u\_\{xx\}, \& 0 \textless{} x \textless{} \textbackslash infty,
\& 0 \textless{} t \textless{} \textbackslash infty, \textbackslash{}
(\textbackslash Gamma \textbackslash text\{У\}) \& u\_x (0, t) - u (0,
t) = 0, \& 0 \textless{} t \textless{} \textbackslash infty,
\textbackslash{} (\textbackslash text\{НУ\}) \& u (x, 0) = u\_0, 0
\textless{} x \textless{} \textbackslash infty.
\textbackslash end\{array\} {]}

Щоб розв'язати цю задачу, ми застосуємо перетворення Лапласа до змінної
t. Зверніть увагу, що можна також застосувати перетворення Лапласа до
змінної x, оскільки воно змінюється від 0 до \textbackslash infty .
Після застосування перетворення

Ми отримуємо звичайне диференціальне рівняння за змінною X.

(13.7)

{[} (ODY) \textbackslash quad sU(x) - u\_0 =
\textbackslash frac\{d\textsuperscript{2U\}\{dx}2\}, \textbackslash quad
0 \textless{} x \textless{} \textbackslash infty, {]}

{[} (\textbackslash Gamma Y) \textbackslash quad
\textbackslash frac\{dU\}\{dx\}(0) = U(0) {]}

(Ми перетворили лише рівняння та граничну умову, а не початкову умову.)
Маємо звичайне диференціальне рівняння другого порядку, але з однією
граничною умовою при x=0. Насправді існує ще одна гранична умова. Він
визначається фізичними міркуваннями і фіксується як \(U(x) \to 0\) ,
коли \(x \to +\infty\) . Зверніть увагу, що для спрощення позначення ми
опустили параметр s із аргументів функції \(U(x)\) всюди, оскільки
рівняння (13.7) є диференціальним рівнянням лише відносно x.

Щоб розв'язати задачу (13.7), запишіть загальне рішення (загальний
розв'язок однорідного + часткового розв'язку гетерогенного) TAC:

{[} U(x) = c\_1 e\^{}\{V\_\{5x\}\} + c\_2 e\^{}\{-V\_\{5x\}\} + u\_0/s.
{]}

Підставляючи цей вираз у граничну умову (13.7), можна визначити
константи \(c_1\) та \(c_2\) (одразу зрозуміло, що \(c_3 = 0\) ,

\pandocbounded{\includegraphics[keepaspectratio]{images/_page_102_Figure_7.jpeg}}

Рис. 13.3. Граф ймовірнісного інтеграла (eri) та його комплементарної
функції (eric).

Інакше температура буде зростати нескінченно зі збільшенням координати
X). Визначаючи сталу \(c_2\) з граничної умови при x=0, отримуємо
фінальний вираз для \(U\left(x\right)\) (13.8)

{[} U(x) = u\_0 \textbackslash left\{ \textbackslash frac\{e\^{}\{-V
\textbackslash widehat\{s\} x\}\}\{s(V \textbackslash widehat\{s\} +
1)\} \textbackslash right\} + \textbackslash frac\{1\}\{s\}. {]}

Залишається зробити останній крок. Щоб визначити температурне поле,
потрібно обчислити \(u\left(x,\,t\right)\)

{[} u(x, t) = \textbackslash mathcal\{L\}\^{}\{-1\}{[}U(x, s){]} {]}

(ми повернулися до написання всіх аргументів функції \(U(x, s)\)). Щоб
повернути перетворення Лапласа, ми використаємо таблиці, наведені в
додатку. У результаті ми отримуємо

(13.9)

{[} u(x, t) = u\_0 - u\_0 \textbackslash left{[}
\textbackslash text\{eric\} \textbackslash left(
\textbackslash frac\{x\}\{2\} \textbackslash sqrt\{t\}
\textbackslash right) + \textbackslash text\{eric\} \textbackslash left(
\textbackslash sqrt\{t\} + \textbackslash frac\{x\}\{2\}
\textbackslash sqrt\{t\} \textbackslash right) e\^{}\{x+t\}
\textbackslash right{]}, {]}

де

{[} \textbackslash operatorname\{eric\}(x) =
\textbackslash frac\{2\}\{\textbackslash sqrt\{\textbackslash pi\}\}
\textbackslash int\_\{x\}\^{}\{\textbackslash infty\}
e\textsuperscript{\{-\textbackslash xi}\{2\}\} d\textbackslash xi {]}

є функцією, доповнюючою до ймовірнісного інтеграла. Графіки цієї функції
показані на рис. 13.3. Аналіз цього рішення P H C

\pandocbounded{\includegraphics[keepaspectratio]{images/_page_103_Figure_7.jpeg}}

13.4. Температура всередині напівнескінченного середовища в різний час.

Потребує дуже малих зусиль. Фіг. (13.4) представлено результати
чисельного аналізу, отриманого на комп'ютері, оснащеному плоттером.

\section{ПРИМІТКИ}\label{ux43fux440ux438ux43cux456ux442ux43aux438-2}

\begin{itemize}
\tightlist
\item
  \begin{enumerate}
  \def\labelenumi{\arabic{enumi}.}
  \tightlist
  \item
    Перетворення Лапласа також можна застосувати до неоднорідного
    диференціального рівняння з частковими похідними (для застосування
    методу розділення змінних рівняння має бути однорідним), згідно з
    ним, загалом, його можна застосовувати до рівнянь із змінними
    коефіцієнтами (метод розділення змінних також може застосовуватися
    до змінних коефіцієнтів). Таблиця 1. 13.2 вказує сфери застосування
    кожного з двох методів.
  \end{enumerate}
\item
  \begin{enumerate}
  \def\labelenumi{\arabic{enumi}.}
  \setcounter{enumi}{1}
  \tightlist
  \item
    Для розв'язання змішаних і граничних задач також використовуються
    перетворення Ганкеля та Мелліна, які, однак, відрізняються від
    перетворення Лапласа в одному аспекті. Перетворення Лапласа замінює
    похідну операцією множення формули
  \end{enumerate}
\end{itemize}

{[} \textbackslash mathscr\{L\}{[}y'{]} =
s\textbackslash mathscr\{L\}{[}y{]} - y(0), {]}

Таблиця 13.2 Порівняння методів розділення перетворення Лапласа та
змінних

{\def\LTcaptype{none} % do not increment counter
\begin{longtable}[]{@{}
  >{\raggedright\arraybackslash}p{(\linewidth - 4\tabcolsep) * \real{0.2907}}
  >{\raggedright\arraybackslash}p{(\linewidth - 4\tabcolsep) * \real{0.3140}}
  >{\raggedright\arraybackslash}p{(\linewidth - 4\tabcolsep) * \real{0.3953}}@{}}
\toprule\noalign{}
\begin{minipage}[b]{\linewidth}\raggedright
\end{minipage} & \begin{minipage}[b]{\linewidth}\raggedright
Метод
\end{minipage} & \begin{minipage}[b]{\linewidth}\raggedright
\end{minipage} \\
\midrule\noalign{}
\endhead
\bottomrule\noalign{}
\endlastfoot
& Перетворення Лапласа & Розділення \\
Гетерогенні PPP & Так & Так \\
Гетерогенні ГІ & Так & Так \\
Змінні ймовірності & Так & de' \\
Нелінійні рівняння & Так & Так \\
\end{longtable}
}

а перетворення Ганкеля і Мелліна замінюють диференціальний оператор на
множник. Наприклад, для перетворення Ганкеля, визначеного формулою

{[} H{[}y{]} = \textbackslash int\_\{0\}\^{}\{\textbackslash infty\} r
J\_\{o\}(\textbackslash xi r) y(r) dr, {]}

справедливе співвідношення

{[}
H\textbackslash left{[}y'\,`\textbackslash left(r\textbackslash right)+\textbackslash frac\{1\}\{r\}~y'\textbackslash left(r\textbackslash right)\textbackslash right{]}=-\textbackslash xi\^{}\{2\}H\textbackslash left{[}y\textbackslash right{]}.
{]}

Таким чином можна розв'язати спеціальне диференціальне рівняння з
змінними коефіцієнтами (рівняння Бесселя). 3. Перетворення Лапласа на
змінній t можна інтерпретувати як проекцію площини xt на осі x, що
призводить до У ц ь

\pandocbounded{\includegraphics[keepaspectratio]{images/_page_104_Figure_9.jpeg}}

ому випадку початкові рівняння, граничні та початкові умови
перетворюються на нове диференціальне рівняння та нові граничні умови
(див. діаграму вище).

\section{ЗАВДАННЯ}\label{ux437ux430ux432ux434ux430ux43dux43dux44f-8}

\begin{enumerate}
\def\labelenumi{\arabic{enumi}.}
\tightlist
\item
  Переконайтеся, що наступна формула для перетворення часткової похідної
  \(u_t\) є істинною:
\end{enumerate}

{[} \textbackslash mathscr\{L\}\textbackslash left{[}u\_\{t\}(x,
t)\textbackslash right{]} = sU(x, s) - u(x, 0). {]}

\begin{enumerate}
\def\labelenumi{\arabic{enumi}.}
\setcounter{enumi}{1}
\tightlist
\item
  Використати перетворення Лаціаса для розв'язання задачі Коші
\end{enumerate}

{[} \textbackslash begin\{aligned\} (\textbackslash text\{УЧП\}) \& u\_t
= \textbackslash alpha\^{}2 u\_\{xx\}, \& -\textbackslash infty
\textless{} x \textless{} \textbackslash infty, \& 0 \textless{} t
\textless{} \textbackslash infty, \textbackslash{}
(\textbackslash text\{HV\}) \& u(x, 0) = \textbackslash sin x, \&
-\textbackslash infty \textless{} x \textless{} \textbackslash infty.
\textbackslash end\{aligned\} {]}

\begin{enumerate}
\def\labelenumi{\arabic{enumi}.}
\setcounter{enumi}{2}
\tightlist
\item
  С помощью преобразования Лапласа по переменной t решите вадачу
\end{enumerate}

{[} \textbackslash begin\{aligned\}
(\textbackslash text\{УЧ\textbackslash Pi\}) \& u\_t =
u\_\{x,\textbackslash text\{v\}\}, \& 0 \textless{} x \textless{}
\textbackslash infty, \& 0 \textless{} t \textless{}
\textbackslash infty, \textbackslash{} (\textbackslash text\{ГУ\}) \&
u\textbackslash left(0,~t\textbackslash right) = \textbackslash sin t,
\& 0 \textless{} t \textless{} \textbackslash infty, \textbackslash{}
(\textbackslash text\{HY\}) \&
u\textbackslash left(x,~0\textbackslash right) = 0, \& 0
\textbackslash leqslant x \textless{} \textbackslash infty.
\textbackslash end\{aligned\} {]}

Дайте фізичну інтерпретацію цієї зали.

\begin{enumerate}
\def\labelenumi{\arabic{enumi}.}
\setcounter{enumi}{3}
\tightlist
\item
  Розв'язати задачу з граничним значенням
\end{enumerate}

(ОДА)

{[} \textbackslash frac\{d\textsuperscript{2U\}\{dx}2\} - sU = A,
\textbackslash quad 0 \textless{} x \textless{} 1, {]}

{[} (\textbackslash Gamma Y) \textbackslash quad {]}

{[} \textbackslash begin\{cases\} \textbackslash frac\{dU\}\{dx\} (0) =
0 \textbackslash{} U(1) = 0. \textbackslash end\{cases\} {]}

Лекція 14

\section{ПРИНЦИП
ДЮАМЕЛЯ}\label{ux43fux440ux438ux43dux446ux438ux43f-ux434ux44eux430ux43cux435ux43bux44f}

МЕТА ЛЕКЦІЇ: Показати, як перетворення Лапласа можна використати для
виявлення цікавих особливостей розв'язання диференціальних рівнянь.
Зокрема, за допомогою алгебранічних перетворень образів Лапласа деякого
розв'язку рівняння часткових похідних можливо отримати відомий принцип
Дюамеля. Інтерпретація цього принципу в теорії звичайних диференціальних
рівнянь відома, але ми покажемо, як його використовувати для розв'язання
деяких змішаних задач для UPP.

технології.

Окрім того, що перетворення Лапласа є потужним методом розв'язання
часткових диференціальних рівнянь, воно дозволяє виявити фізичну
сутність отриманих розв'язків. За допомогою трансформації Лапласа ми
проілюструємо дуже важливий і цікавий принцип Дюамеля. Перш ніж перейти
до принципу Дюамеля, розглянемо проблему, яка дуже часто виникає у

\section{Теплопровідність у стрижні, кінці якого підтримуються при
заданій
температурі}\label{ux442ux435ux43fux43bux43eux43fux440ux43eux432ux456ux434ux43dux456ux441ux442ux44c-ux443-ux441ux442ux440ux438ux436ux43dux456-ux43aux456ux43dux446ux456-ux44fux43aux43eux433ux43e-ux43fux456ux434ux442ux440ux438ux43cux443ux44eux442ux44cux441ux44f-ux43fux440ux438-ux437ux430ux434ux430ux43dux456ux439-ux442ux435ux43cux43fux435ux440ux430ux442ux443ux440ux456}

Дуже часто необхідно знайти температуру всередині області, якщо відомо,
що на межі області вона змінюється з часом відповідно до певного закону
(граничні умови, залежні від часу).

{[} u(0,t)=0 \textbackslash left( {]}

\begin{verbatim}
\begin{array}{c} \\ \\ \\ \\ \\ \\ \\ \\ \\ \\ \\ \\ \\ \\ \\ \\ \\ \\ \
\end{verbatim}

Рис. 14.1. Часово-залежні граничні умови.

Наприклад, розглянемо термоізольований стержень, правий кінець якого
підтримується у \(f(t)\) (див. рисунок 14.1). Потрібно вирішити цю
проблему

(14.1)

{[} u\_\{t\} = u\_\{xx\}, \textbackslash quad 0 \textless{} x
\textless{} 1, \textbackslash quad 0 \textless{} t \textless{}
\textbackslash infty, {]}

{[} (\textbackslash Gamma Y) \textbackslash quad
\textbackslash begin\{cases\} u(0, t) = 0, \textbackslash{} u(1, t) =
f(t), \textbackslash end\{cases\} \textbackslash quad 0 \textless{} t
\textless{} \textbackslash infty, {]}

{[} (HY) \textbackslash quad u(x, 0) = 0, \textbackslash quad 0
\textbackslash le x \textbackslash le 1. {]}

Сподіваємося, що розв'язання цієї задачі (14.1) буде легким, оскільки ми
знаємо розв'язок простішої задачі, де температура в кінці є сталою:

Насправді, якщо обидві задачі ((14.1) і (14.2)) розв'язуються одночасно
за допомогою перетворення Лапласа, то отримаємо визначний результат:
розв'язок задачі (14.1) виражається через розв'язання задачі (14.2).

Отже, розв'язуючи обидві задачі одночасно, ми отримуємо

\section{Легке завдання (14.2) (Постійні
солдати)}\label{ux43bux435ux433ux43aux435-ux437ux430ux432ux434ux430ux43dux43dux44f-14.2-ux43fux43eux441ux442ux456ux439ux43dux456-ux441ux43eux43bux434ux430ux442ux438}

Застосування перетворення -Лапласа до задачі (14.2)

{[} \textbackslash frac\{d\textsuperscript{2W\}\{dx}2\} - sW(x) = 0, {]}

{[} W(0) = 0, {]}

{[} W(1) = 1/s {]}

.

Пемеріль ODU

{[} W(x, s) = \textbackslash frac\{1\}\{s\} \textbackslash left{[}
\textbackslash frac\{\textbackslash sinh(x
\textbackslash sqrt\{s\})\}\{\textbackslash sinh(\textbackslash sqrt\{s\})\}
\textbackslash right{]}. {]}

Пошук оберненого перетворення

{[} w(x, t) = {]}

{[} = x + \textbackslash frac\{2\}\{\textbackslash pi\}
\textbackslash sum\_\{n=1\}\^{}\{\textbackslash infty\}
\textbackslash frac\{(-1)\^{}n\}\{n\}
e\textsuperscript{\{-(n\textbackslash pi)}2 t\} \textbackslash times {]}

\(\times \sin(n\pi x)\) (розв'язання задачі з постійним ГІ).

\section{Важке завдання (14.1) (Солдати, що залежать від
часу)}\label{ux432ux430ux436ux43aux435-ux437ux430ux432ux434ux430ux43dux43dux44f-14.1-ux441ux43eux43bux434ux430ux442ux438-ux449ux43e-ux437ux430ux43bux435ux436ux430ux442ux44c-ux432ux456ux434-ux447ux430ux441ux443}

Застосування перетворення Лапласа до задачі (14.1)

{[} \textbackslash frac\{d\textsuperscript{2U\}\{dx}2\} - sU(x) = 0, {]}

{[} U\textbackslash left(
\textbackslash mathbf\{0\}\textbackslash right) =0, {]}

{[} U(1) = F(s) {]}

.

{[} U(x, s) = F(s) \textbackslash left{[}
\textbackslash frac\{\textbackslash sinh(x
\textbackslash sqrt\{s\})\}\{\textbackslash sinh(\textbackslash sqrt\{s\})\}
\textbackslash right{]}. {]}

Розділіть і помножте на z

{[} U(x, s) = F(s) \textbackslash left\{ s \textbackslash left{[}
\textbackslash frac\{\textbackslash sinh(x
\textbackslash sqrt\{s\})\}\{s \textbackslash sinh(s)\}
\textbackslash right{]} \textbackslash right\}. {]}

Використовуючи співвідношення \(\mathscr{L}[w_t] = sW - w(x, 0),\)
отримуємо

{[} U(x, s) = F(s) \textbackslash mathcal\{L\}{[}w\_t{]}. {]}

Відповідно,

{[} u(x, t) = \textbackslash mathcal\{L\}\^{}\{-1\} \{ F(s)
\textbackslash mathcal\{L\}{[}w\_t{]} \} = {]}

{[} = \textbackslash mathcal\{L\}\^{}\{-1\} {[} F(s) {]} *
\textbackslash mathcal\{L\}\^{}\{-1\}{[}w\_t{]} = {]}

{[} = f(t) * w\_t(t) = {]}

{[} = \textbackslash int\_0\^{}t w\_t(x, t-\textbackslash tau)
f(\textbackslash tau) d\textbackslash tau = {]}

(після інтегрування частинами)

{[} = \textbackslash int\_\{0\}\^{}\{t\} w(x, t-\textbackslash tau)
f'(\textbackslash tau) d\textbackslash tau + f(0) w(x, t). {]}

(Розв'язок задачі з часово-залежними ГІ виражається через розв'язок
задачі з постійними ГІ

розв'язано у задачах з константою

Отже, ми висловили розв'язок задач із граничними умовами, що змінюються
в часі, через розв'язання задачі з постійними граничними умовами.
Відповідні формули такі:

(14.3)

{[} u(x, t) = \textbackslash int\_\{0\}\^{}\{t\} w\_\{t\}(x, t -
\textbackslash tau) f(\textbackslash tau) d\textbackslash tau = {]}

{[} = \textbackslash int\_\{0\}\^{}\{t\} w(x, t - \textbackslash tau)
f'(\textbackslash tau) d\textbackslash tau + f(0) w(x, t) {]}

та є математичною формою принципу Дюамеля. Тепер, щоб отримати розв'язок
задачі (14.1), залишається замінити розв'язок задачі з сталими
граничними умовами у формулу (14.3). Слід зазначити, що в цьому випадку
ми маємо використовувати другу формулу з (14.3), оскільки ряд для
\(w(x, t)\) не можна диференціювати за змінною t (диференціація термів
веде до дивергентного ряду).

Принцип Дюамеля має й інші корисні властивості.

\section{Важливість принципу
Дюамеля}\label{ux432ux430ux436ux43bux438ux432ux456ux441ux442ux44c-ux43fux440ux438ux43dux446ux438ux43fux443-ux434ux44eux430ux43cux435ux43bux44f}

У цьому прикладі ми спочатку розв'язали просту задачу з постійними
граничними умовами, а потім використали принцип Дюамеля (14.3), щоб
отримати розв'язок задачі з змінними граничними умовами. Однак дуже
часто навіть просту задачу з постійними граничними умовами неможливо
розв'язати аналітично. У цьому випадку рішення проблеми з постійними
граничними умовами можна знайти експериментально, тобто шляхом оснащення
установки відповідними приладами та вимірювання відгуку системи. Тоді,
використовуючи принцип Дюамеля, можна знайти розв'язок задачі за умов
довільних граничних умов, що змінюються за часом. Насправді, якщо ми
знаємо з спостережень функцію \(w(x, t)\), яка відповідає сталим
граничним умовам, то розв'язок задачі з довільними граничними умовами,
що містить функцію \(f(t)\), знаходиться згідно з формулою (14.3).

\section{ПРИМІТКИ}\label{ux43fux440ux438ux43cux456ux442ux43aux438-3}

Існує ще одна форма принципу Дюамеля. Він дозволяє висловити рішення
проблеми

(14.4)

{[} \textbackslash begin\{aligned\} \&
(\textbackslash text\{УЧ\}\textbackslash Pi) \& u\_t = u\_\{xx\}, \& 0
\textless{} x \textless{} 1, \& 0 \textless{} t \textless{}
\textbackslash infty, \textbackslash{} \& \textbackslash left(
\textbackslash Gamma \textbackslash text\{У\} \textbackslash right) \&
\textbackslash begin\{cases\} u \textbackslash left( 0, ~t
\textbackslash right) = 0, \textbackslash{} u \textbackslash left( 1, ~t
\textbackslash right) = f \textbackslash left( t \textbackslash right),
\textbackslash end\{cases\} \& 0 \textless{} \textbackslash textbf\{I\}
\textless{} \textbackslash infty, \textbackslash{} \&
(\textbackslash text\{HY\}) \& u \textbackslash left( x, ~0
\textbackslash right) = 0, \& 0 \textbackslash leqslant x
\textbackslash leqslant 1, \textbackslash end\{aligned\} {]}

розв'язавши \(w\left(x,\,t\right)\) інше просте завдання

(

PU \()\) \(w_t=w_{xx}, \quad 0< x<1, \quad 0< t<\infty,\) \((\) 14,5
\()\) \((\) GU \()\)

{[} \textbackslash begin\{cases\} w~(0,~t)=0 \textbackslash{}
w~(1,~t)=\textbackslash delta~(t) \textbackslash end\{cases\} {]}

температурний імпульс на \(t=0\) , \((\) HU \()\)
\(w\ (x,\ 0)=0, \quad 0\leqslant x\leqslant 1.\) У цьому випадку
відповідна формула має вигляд

{[} u(x, t) = \textbackslash int\_\{0\}\^{}\{t\} w(x,
t-\textbackslash tau) f(\textbackslash tau) d\textbackslash tau {]}

і дозволяє знайти температуру \(u\left(x,\ t\right)\) при довільній
граничній температурі \(f\left(t\right)\) , якщо відома відповідь
системи \(w\left(x,\ t\right)\) на температурний імпульс.

\section{ЗАЛАЧІ}\label{ux437ux430ux43bux430ux447ux456}

\begin{itemize}
\tightlist
\item
  \begin{enumerate}
  \def\labelenumi{\arabic{enumi}.}
  \tightlist
  \item
    Доведіть принцип Дюамеля (14.6), застосувавши перетворення Лапласа
    до задач (14.4) і (14.5) та застосовуючи ті ж міркування, що й у
    виведенні (14.3). ПРИМІТКА. Образ Лапласа функції імпульсу
    \(\delta(t)\) визначається формулою \(\mathscr{L}[\delta(t)] = 1\) .
  \end{enumerate}
\item
  \begin{enumerate}
  \def\labelenumi{\arabic{enumi}.}
  \setcounter{enumi}{1}
  \tightlist
  \item
    Показати, що часткова похідна функції \(w(t)\)
  \end{enumerate}
\end{itemize}

{[} w(x, t) = x + \textbackslash frac\{2\}\{\textbackslash pi\}
\textbackslash sum\_\{n=1\}\^{}\{\textbackslash infty\}
\textbackslash frac\{(-1)\^{}n\}\{n\}
e\textsuperscript{\{-(n\textbackslash pi)}2 t\}
\textbackslash sin(n\textbackslash pi x) {]}

Її не можна обчислити за диференціацією за членами ряду, оскільки
диференційована послідовність розходиться в будь-якому X.

\begin{enumerate}
\def\labelenumi{\arabic{enumi}.}
\setcounter{enumi}{2}
\tightlist
\item
  Припустимо, що лівий кінець термічно ізольованого стрижня
  підтримується при нульовій температурі, а правий кінець піддається
  імпульсному нагріванню. Припустимо, що початкова температура стрижня
  дорівнює нулю (деяка початкова опорна температура), а після дії
  теплового
\end{enumerate}

Імпульс --- температура вимірюється в точці з координатою \(x\!=\!0.5\)
у послідовні моменти часу. Внаслідок цього таблиця

{\def\LTcaptype{none} % do not increment counter
\begin{longtable}[]{@{}lll@{}}
\toprule\noalign{}
Час & Температура & \\
\midrule\noalign{}
\endhead
\bottomrule\noalign{}
\endlastfoot
\(t_1\) & \(w_t\) & \\
\(t_2 = 2t_1\) & \(w_*\) & \\
\(t_3 = 3t_1\) & w 2 & \\
• & • & \\
• & • & \\
. • & • & \\
\(t_n = nt_i\) & \(w_{\pi}\) & \\
\end{longtable}
}

Використовуючи ці дані, приблизно знайдіть температурну відповідь u (0,
5; t) у формі

\begin{itemize}
\item
  \begin{enumerate}
  \def\labelenumi{\alph{enumi})}
  \tightlist
  \item
    \(u(1, t) = \sin t\) ,
  \end{enumerate}
\item
  \begin{enumerate}
  \def\labelenumi{\alph{enumi})}
  \setcounter{enumi}{1}
  \tightlist
  \item
    \(u(1, t)\) = \(f(t)\) (\(f(t)\) --- довільна функція).
  \end{enumerate}
\item
  \begin{enumerate}
  \def\labelenumi{\arabic{enumi}.}
  \setcounter{enumi}{3}
  \tightlist
  \item
    Використовуючи принцип Дюамеля, розв'яжіть змішану задачу
  \end{enumerate}

  \textbackslash begin\{aligned\} (\textbackslash mathrm\{Y\}
  \textbackslash mathrm\{H\} \textbackslash Pi) \& u\_t =
  \textbackslash alpha\^{}2 u\_\{xx\}, \textbackslash quad 0 \textless{}
  x \textless{} 1, \textbackslash quad 0 \textless{} t \textless{}
  \textbackslash infty, \textbackslash{} (\textbackslash mathrm\{F\}
  \textbackslash mathrm\{Y\}) \& \textbackslash left\{
  \textbackslash begin\{aligned\} u ~(0, ~t) = 0, \textbackslash{} u
  ~(1, ~t) = \textbackslash sin t, \textbackslash end\{aligned\}
  \textbackslash end\{aligned\}

  \textbackslash right. \& 0 \textless{} t \textless{}
  \textbackslash infty, \textbackslash{} (\textbackslash mathrm\{H\}
  \textbackslash mathrm\{Y\}) \& u ~(x, ~0) = 0, \textbackslash quad 0
  \textbackslash leqslant r \textbackslash leqslant 1.
  \textbackslash end\{aligned\}
\end{itemize}

\section{\texorpdfstring{КОНВЕКТИВНИЙ ЧЛЕН \(U_x\) У ЗАДАЧІ
ДИФУЗІЇ}{КОНВЕКТИВНИЙ ЧЛЕН U\_x У ЗАДАЧІ ДИФУЗІЇ}}\label{ux43aux43eux43dux432ux435ux43aux442ux438ux432ux43dux438ux439-ux447ux43bux435ux43d-u_x-ux443-ux437ux430ux434ux430ux447ux456-ux434ux438ux444ux443ux437ux456ux457}

МЕТА ЛЕКЦІЇ: Показати, як термін \(u_x\) у задачі дифузії

{[} u\_t = Du\_\{xx\} - Vu\_x {]}

Дифузійний термін Конвективний член

описує явище конвекції. Наведене вище рівняння конвективної дифузії
зустрічається у багатьох задачах і описує як дифузію, так і конвекцію.
Роль кожного з цих механізмів визначається співвідношенням коефіцієнтів
D і V.

Оскільки конвекція речовини спричинена її рухом разом із навколишнім
середовищем, природно перейти у систему координат, яка рухається разом
із середовищем. Внаслідок такого перетворення конвективний член буде
виключено, і отримане рівняння можна розв'язати одним із відомих
методів. Перетворивши отриманий розв'язок у фіксовану систему відліку,
отримуємо розв'язок початкової задачі.

Поки що ми розглядали теплопровідність (або певний вид дифузії) в
одномірних системах. Тепер розглянемо проблему розподілу концентрації
певної речовини, що виходить із землі і входить у висхідні повітряні
потоки. Ми врахуємо як дифузію матерії, так і конвективний транспорт
повітрям, який рухається зі швидкістю V. Очевидно, що в цьому випадку
конвективний механізм передачі речовини відіграє важливішу роль, ніж
механізм дифузії (роль кожного механізму визначається співвідношенням
коефіцієнтів дифузії та швидкості повітря). Наше завдання --- вивести
рівняння конвективної дифузії

{[} u\_t = Du\_\{xx\} - Vu\_x {]}

Щоб переконатися, що рівняння конвективної дифузії правильне, ми
використаємо два базові факти.

\begin{enumerate}
\def\labelenumi{\arabic{enumi}.}
\tightlist
\item
  Потік через конвекцію.
\end{enumerate}

Конвективний потік речовини зліва направо через поперечний переріз у цій
точці дорівнює Vu(x, t), де V --- швидкість середовища (см/с), а
\(u(x, t)\) --- концентрація речовини (г/см) (рис. 15.1).

\pandocbounded{\includegraphics[keepaspectratio]{images/_page_112_Picture_4.jpeg}}

{[} \textbackslash begin\{aligned\} \textbackslash vdots
\textbackslash{} \textbackslash vdots \textbackslash{}
\textbackslash vdots \textbackslash{} \textbackslash vdots
\textbackslash{} \textbackslash vdots \textbackslash{}
\textbackslash vdots \textbackslash{} \textbackslash vdots
\textbackslash{} \textbackslash vdots \textbackslash{}
\textbackslash vdots \textbackslash{} u\_i + Vu\_x = 0 \textbackslash{}
(\textbackslash textit\{SpaBhenue nepenoca nep8ozo nops@ka\})
\textbackslash{} \textbackslash vdots \textbackslash{} u(x,0) = 0
\textbackslash end\{aligned\} {]}

\begin{verbatim}
РИС. 15.1. Кількість речовини, що транспортується за одну секунду через поперечний розріз у точці x шляхом конвекції, дорівнює Vu\left( v,\ t \right). РИС. 15.2. Проблема суто конвективного переносу.

2. Потік через дифузію. {\it Диффузионный поток} речовини зліва направо через переріз у заданій точці дорівнює - Du_x(x, t) , де D — коефіцієнт дифузії.
\end{verbatim}

Якщо підставити ці два вирази у закон збереження з лекції 3, отримаємо
базове диференціальне рівняння в частинних похідних

{[} u\_t = Du\_\{xx\} - Vu\_x. {]}

Щоб зрозуміти, як конвективний член впливає на природу рішення, спочатку
розглянемо лише конвективний механізм переносу. Типовим прикладом є
проблема скидання речовини в річку (швидкість течії річки V). Нехай x
--- це відстань нижче за течією від місця скидання. Якщо речовина не
дифундує у воді, то концентрація \(u(x, t)\) речовини є розв'язком
наступної проблеми (рис. 15.2):

(PPP)

{[} u\_t = -Vu\_x {]}

, 0 \textless{} x \textless{} \textbackslash infty , 0 \textless{} t
\textless{} \textbackslash infty , (\textbackslash Gamma Y) \(u(0, t)\)
= P \textbackslash leftarrow Постійна концентрація на місці розряду,

\begin{verbatim}
(НУ) u(x, 0) = 0 \leftarrow На початку річка була чистою.
\end{verbatim}

Перед розв'язуванням проблеми слід подумати, яким має бути її рішення.
Очевидно, забруднення річки

(кількість речовини у воді) спочатку буде нуль, а потім, підтримувана
константою x=0, рухатиметься вниз за течією зі швидкістю V. Щоб
переконатися в цьому, розв'яжемо задачу (15.1). Оскільки рівняння (15.1)
є лінійним, а гранична умова також лінійна, можна сподіватися розв'язати
цю задачу методом розділення змінних або методом інтегральних
перетворень. Однак, оскільки координата x тут не обмежена, метод
розділення змінних не застосовується. Використайте трансформацію Лапласа
у часі t.

\section{Розв'язання рівняння транспорту за допомогою перетворення
Лапласа}\label{ux440ux43eux437ux432ux44fux437ux430ux43dux43dux44f-ux440ux456ux432ux43dux44fux43dux43dux44f-ux442ux440ux430ux43dux441ux43fux43eux440ux442ux443-ux437ux430-ux434ux43eux43fux43eux43cux43eux433ux43eux44e-ux43fux435ux440ux435ux442ux432ux43eux440ux435ux43dux43dux44f-ux43bux430ux43fux43bux430ux441ux430}

Використання перетворення Лапласа

{[} U(x) = \textbackslash int\_\{0\}\^{}\{\textbackslash infty\} u(x, t)
e\^{}\{-st\} dt {]}

Задача конвективного транспорту (15.1) може бути зведена до вигляду

{[} sU(x) = -V \textbackslash frac\{dU\}\{dx\}, \textbackslash quad 0
\textless{} x \textless{} \textbackslash infty, {]}

{[} U(0) = P/s. {]}

Розв'язуючи цю дуже просту задачу Коші, ми отримуємо
\(U(x) = Pe^{-(sx/V)}/s\) .

Використовуючи таблиці обернених перетворень, ми знаходимо розв'язок
задачі \(u(x, t) = \mathcal{L}^{-1}[U] = PH(t-x/V),\) де H (ξ) ---
покрокова функція Хевісайда

{[} H(\textbackslash xi) = {]}

{[} \textbackslash begin\{cases\} 0, \& \textbackslash xi \textless{} 0,
\textbackslash{} 1, \& \textbackslash xi \textbackslash geqslant 0.
\textbackslash end\{cases\} {]}

Отже, розв'язок нашої проблеми можна надати такій формі:

{[} u(x, t) = {]}

{[} \textbackslash begin\{cases\} 0, \& t \textless{} x/V,
\textbackslash{} P, \& t \textbackslash ge x/V.
\textbackslash end\{cases\} {]}

Конечно, все це очень просто! Не сложнее, чем сбрасывать що-нибудь на
ленту транспортера, а затем наблюдать за его движением. Однако все
становится очень интересным, якщо примесь диффундирует в среде. Чтобы
разобраться в том, що про-исходит с движущейся волной при наличии
диффузии, решим

следующую задачу

(15.2)

{[} \textbackslash begin\{aligned\} (3 \textbackslash text\{ HI\})
\textbackslash quad u\_t \&= Du\_\{xx\} - Vu\_x, \textbackslash quad
-\textbackslash infty \textless{} x \textless{} \textbackslash infty,
\textbackslash{} (HY) \textbackslash quad u(x, 0) \textbackslash quad 1
- H(x), \textbackslash quad -\textbackslash infty \textless{} x
\textless{} \textbackslash infty. \textbackslash end\{aligned\} {]}

где, как обычно, \(H(x)\) --- функція Хевисайда Начальное распределение
концентрации изображено на рис 15 4.

\pandocbounded{\includegraphics[keepaspectratio]{images/_page_114_Figure_4.jpeg}}

Рис. 15.3 Конвективная волна a---передиції фронт примеси, \(\delta\)
---на одну единицу вперед волнового фронта

Обратим внимание на то, що в задаче (15.2) мы отодвинули границу в ---
\(\infty\) (теперь решается задача Коши), так що она не оказывает теперь
влияние на оценку розв'язок. Чтобы решить задачу (15.2), можна
воспользоваться преобразованием Лапласа Р И С

\pandocbounded{\includegraphics[keepaspectratio]{images/_page_114_Figure_7.jpeg}}

. 15 4 початкові умови в задаче конвективной диффузии.

ло переменной t или преобразованием Фурье по переменной x. Однако в
цьому случае гораздо интереснее поступить совсем по-другому Введем новую
систему координат, яка движется вдоль старой со скоростью V. Другими
словами, вместо системи координат, привязанной к берегу реки, мы
рассмотрим систему координат, яка движется со скоростью фронта примеси
(конечно, при налични диффузии фронт буде «смазан»). С точки зрения
математики це означает, що мы заменяем пространственную координату x на
новую \(\xi = x - Vt\) .

Теперь ясно, що

когда \(\xi = 0\) , мы находимся па фронте распространяющейся примеси.

когда \(\xi = 1\) , мы находимся на одну единицу длины впереди фронта.

когда \(\xi = -1\) , мы находимся на одну единицу длины позади фронта.

Наша задача --- преобразовать исходную задачу с НУ

{[} \textbackslash begin\{aligned\} (\textbackslash mathrm\{Y\} +
\textbackslash Pi) \& u\_1 = Du\_\{xx\} - Vu\_x, \&
-\textbackslash infty \textless{} x \textless{} \textbackslash infty,
\textbackslash{} (\textbackslash mathrm\{H\} \textbackslash mathrm\{Y\})
\& u(x, 0) = 1 - H(x), \& -\textbackslash infty \textless{} x
\textless{} \textbackslash infty, \textbackslash end\{aligned\} {]}

у нову задачу в рухомій системі координат, розв'язати нову задачу, а
потім перетворити розв'язок на старі координати (x, t). Спочатку
замінимо змінні (замінимо незалежні змінні). Замість старих координат
(x, t) вводимо нові \((\xi, \tau)\) відповідно до формул

{[} \textbackslash xi = x - Vt, {]}

{[} \textbackslash tau = t. {]}

Слід зазначити, що літери \textbackslash tau і t позначають одну й ту ж
Ведичіну. Такі позначення вводяться для уникнення

плутанина у наступних формулах. Записати диференціальне рівняння в
похідних у нових координатах (~, т), мы воспользуемся следующими
формулами перехода;

MATHBLK0547X

MATHBLK0548X

MATHBLK0549X

Функциональная диаграмма на рис. 15.5 делает эти формулы совсем
очевидными. Эта диаграмма особенно полезна при вычислении частных
производных функций и и и по переменным Р И С

\pandocbounded{\includegraphics[keepaspectratio]{images/_page_115_Picture_14.jpeg}}

15.5 Диаграмма, иллюстрирующая функциональную зависимость переменных

\begin{verbatim}
\xi   и   \tau  , поскольку   \xi   зависит и от x, и от t. Переменная   \tau   зависит только от t. Это все, что касается преобразований. Теперь подставим найденные выражения   u_t  ,   u_x   и   u_{vx}   в уравнение (15 2) и получим
\end{verbatim}

MATHBLK0550X

откуда

MATHBLK0551X.

\begin{verbatim}
Следовательно, новая задача с начальным условием в персменных   \xi   и   \tau   имеет вид
\end{verbatim}

(УЧП)

MATHBLK0552X

\begin{verbatim}
(НУ)   u(\xi, 0) = 1 - H(\xi), \quad -\infty < \xi < \infty.

Так как   \xi = x   при t = 0, то новые граничные условия совпадают со старыми. Эта задача была решена в лекции 12 с помощью интегрального преобразования Фурье, и ее решение записывается в виде
\end{verbatim}

MATHBLK0553X

\begin{verbatim}
где   \phi(\beta)   — начальное условие. Следовательно, в нашем случае
\end{verbatim}

MATHBLK0554X

После подстановки

MATHBLK0555X

получаем окончательный результат

(15.3)

MATHBLK0556X

MATHBLK0557X

{[} \textbackslash begin\{cases\} \textbackslash frac\{1\}\{2\}
\textbackslash left{[} 1 + \textbackslash operatorname\{erf\}
\textbackslash left(
\textbackslash frac\{-\textbackslash xi\}\{2\textbackslash sqrt\{D\textbackslash tau\}\}
\textbackslash right) \textbackslash right{]}, \& \textbackslash xi
\textless{} 0, \textbackslash{} \textbackslash frac\{1\}\{2\}
\textbackslash operatorname\{erfc\} \textbackslash left(
\textbackslash frac\{\textbackslash xi\}\{2\textbackslash sqrt\{D\textbackslash tau\}\}
\textbackslash right), \& 0 \textbackslash leq \textbackslash xi.
\textbackslash end\{cases\} {]}

MATHBLK0558X

MATHBLK0559X

Графики этого решения для различных моментов времени изображены на рис.
15.6. Осталось сделать последний шаг --- записать решение нашей задачи в
координатах x и t. Подставляя

\begin{verbatim}
MATHBLK0560X,   \tau = t
\end{verbatim}

в формулы (15.3), получаем

(15.4)

MATHBLK0561X

MATHBLK0562X

{[} \textbackslash begin\{cases\} \textbackslash frac\{1\}\{2\}
\textbackslash left{[} 1 +
\textbackslash operatorname\{erf\}\textbackslash left(\textbackslash frac\{Vt
- x\}\{2\textbackslash sqrt\{Dt\}\}\textbackslash right)
\textbackslash right{]}, \& Vt \textgreater{} x, \textbackslash{}
\textbackslash frac\{1\}\{2\}
\textbackslash operatorname\{erfc\}\textbackslash left(\textbackslash frac\{x
- Vt\}\{2\textbackslash sqrt\{Dt\}\}\textbackslash right), \& Vt
\textbackslash leq x. \textbackslash end\{cases\} {]}

MATHBLK0563X

MATHBLK0564X

MATHBLK0565X

{[} \textbackslash begin\{aligned\} u\_t = u\_\{xx\} - 2u\_x, \&
-\textbackslash infty \textless{} x \textless{} \textbackslash infty, \&
0 \textless{} t \textless{} \textbackslash infty, \textbackslash{} u(x,
0) = \textbackslash sin x, \& -\textbackslash infty \textless{} x
\textless{} \textbackslash infty. \textbackslash end\{aligned\} {]}

MATHBLK0566X

MATHBLK0567X

MATHBLK0568X

{[} \textbackslash begin\{aligned\} u\_t = u\_\{xx\} - 2u\_x, \&
-\textbackslash infty \textless{} x \textless{} \textbackslash infty, ~0
\textless{} t \textless{} \textbackslash infty, \textbackslash{} u(x,
~0) = e\^{}x \textbackslash sin x, \& -\textbackslash infty \textless{}
x \textless{} \textbackslash infty. \textbackslash end\{aligned\} {]}

MATHBLK0569X

MATHBLK0570X

MATHBLK0571X

{[} \textbackslash begin\{aligned\} (\textbackslash text\{YHI\}) \& u\_t
= -2u\_x, \textbackslash{} (\textbackslash text\{HY\}) \& u\_t(x, 0) =
e\textsuperscript{\{-x}2\}. \textbackslash end\{aligned\} {]}

MATHBLK0572X

MATHBLK0573X

Проверьте ваш ответ.

\begin{enumerate}
\def\labelenumi{\arabic{enumi}.}
\setcounter{enumi}{3}
\tightlist
\item
  Решите задачу
\end{enumerate}

MATHBLK0574X

MATHBLK0575X

Проверьте ваше решение. Как выглядит решение в различные моменты
времени?

УКАЗАНИЕ. Переход к движущейся системе отсчета позволяет сразу отбросить
слагаемое Vu\_x в уравнении конвективной диффузии. После решения новой
задачи

MATHBLK0576X

MATHBLK0577X

MATHBLK0578X

\begin{verbatim}
сделайте замену переменных   \xi = x - Vt  ,   \tau = t  . В нашем случае можно в явном виде вычислить интеграл
\end{verbatim}

MATHBLK0579X

Такой интеграл уже встречался в лекции 12. Может быть, читателю будет
удобнее обратиться к таблицам преобразования Фурье.

\section{ОДНОМЕРНОЕ ВОЛНОВОЕ УРАВНЕНИЕ (ГИПЕРБОЛИЧЕСКИЕ
УРАВНЕНИЯ)}\label{ux43eux434ux43dux43eux43cux435ux440ux43dux43eux435-ux432ux43eux43bux43dux43eux432ux43eux435-ux443ux440ux430ux432ux43dux435ux43dux438ux435-ux433ux438ux43fux435ux440ux431ux43eux43bux438ux447ux435ux441ux43aux438ux435-ux443ux440ux430ux432ux43dux435ux43dux438ux44f}

ЦЕЛЬ ЛЕКЦИИ: Показать, что колебания сгруны описываются уравнением

MATHBLK0580X

и что это уравнение --- следствие законов Ньютона. Рассматриваются и
некоторые другие виды волнового уравнения:

MATHBLK0581X

MATHBLK0582X

{[} \textbackslash begin\{split\} u\_\{tt\} \&=
\textbackslash alpha\^{}2 u\_\{xx\} + F\textbackslash left(x,
~t\textbackslash right), \textbackslash{} u\_\{tt\} \&=
\textbackslash alpha\^{}2 u\_\{xx\} - \textbackslash beta u\_t -
\textbackslash lambda u + F\textbackslash left(x,
~t\textbackslash right). \textbackslash end\{split\} {]}

MATHBLK0583X

MATHBLK0584X

До сих пор мы занкмались физическими процессами, которые описываются
одномерными параболическими уравнениями (диффузионные задачи). Теперь мы
приступим к изучению еще одного класса уравнений с частными
производными---гиперболических уравнений. Начнем с одномерного волнового
уравнения, которое описывает, в частности, поперечные колебания струны.

\section{Колебания
струны}\label{ux43aux43eux43bux435ux431ux430ux43dux438ux44f-ux441ux442ux440ux443ux43dux44b}

Рассмотрим малые колебания струны с закрепленными концами. Предположим,
что струна туго натянута, сделана из однородного материала и колеблется
в одной плоскости (рис. 16.1).

\pandocbounded{\includegraphics[keepaspectratio]{images/_page_119_Figure_10.jpeg}}

РИС. 16.1. Поперечные колебания струны: a --- смещение струны от
положения равновесия.

Действием сил тяжести на струну будем пренебрегать. Для построения
математической модели колебаний струны рассмотрим все силы, действующие
на небольшой участок струны (рис. 16.2). Оказывается, что волновое
уравнение---это не что иное, как закон движения Ньютона (изменение
импульса равно сумме действующих сил), примененный к струне. Глядя на
рис. 16.1, мы можем представить те силы, которые действуют на струну в
направлений, перпендикулярном оси х.

\pandocbounded{\includegraphics[keepaspectratio]{images/_page_120_Figure_2.jpeg}}

\begin{verbatim}
РИС. 16.2. Малый участок   (x, x + \Delta x)   колеблющейся струны.

1. Сила, обусловленная натяжением струны (  \alpha^2 u_{xx}  ). Величина поперечной компоненты силы натяжения находится по формуле
\end{verbatim}

Сила =

\begin{verbatim}
MATHBLK0585X  \simeq T [u_x(x + \Delta x, t) - u_x(x, t)].  2. Внешняя сила F (x, t).
\end{verbatim}

Внешняя сила \(F(x, t)\), приложенная к струне, может произвольным
образом зависеть от x и t. Некоторые примеры внешних сил приводятся
ниже:

а) гравитационная сила \(F(x, t)\) = -mg;

б) импульсы, распределенные вдоль струны и действующие на нее в
различные моменты времени:

в) силы, возникающие при действии звуковой волны на кожу барабана (позже
мы подробно изучим колебания мембраны, которые описываются волновым
уравнением).

\begin{verbatim}
3. Силы трения, действующие на струну (—  \beta u_t  ). Если струна колеблется в среде, то возникает сила сопротивления, которая пропорциональна скорости   u_t  .
\end{verbatim}

\begin{enumerate}
\def\labelenumi{\arabic{enumi}.}
\setcounter{enumi}{3}
\tightlist
\item
  Возвращающая сила (- үи).
\end{enumerate}

Эта сила направлена противоположно смещению струны. Если смещение u
положительно, то сила отрицательна.

Если применить теперь уравнения движения Ньютона к небольшому участку
струны, то получим

MATHBLK0586X

\begin{verbatim}
где   \rho   — линейная плотность струны. Разделим обе части этого сротношения на   \Delta x   и устремим   \Delta x   — 0. В результате получим хорошо известное телеграфное уравнение
\end{verbatim}

(16.1)

MATHBLK0587X

\begin{verbatim}
В этом уравнении следовало бы разделить   \alpha  ,   \beta  ,   \gamma   и F(x,t) на   \rho  , но мы ради простоты сохраним для этих величин прежние обозначения. Искомое уравнение получено. Приведем интерпретацию простейшего волнового уравнения
\end{verbatim}

MATHBLK0588X

основанную на интуитивных соображениях.

\section{Интуитивная интерпретация волнового
уравнения}\label{ux438ux43dux442ux443ux438ux442ux438ux432ux43dux430ux44f-ux438ux43dux442ux435ux440ux43fux440ux435ux442ux430ux446ux438ux44f-ux432ux43eux43bux43dux43eux432ux43eux433ux43e-ux443ux440ux430ux432ux43dux435ux43dux438ux44f}

У читателя может возникнуть вопрос: почему уравнение вида (16.2) должно
описывать нечто такое, что похоже на колебания скрипичной струны? Чтобы
разобраться в этом, необходимо понять,

\pandocbounded{\includegraphics[keepaspectratio]{images/_page_121_Picture_9.jpeg}}

\begin{verbatim}
РИС. 16 3. Интерпретация волнового уравнения   u_{tt} = \alpha^2 u_{xx}

что величина   u_{tt}   является вертикальным ускорением струны в точке x. Следовательно, уравнение (16.2) можно интерпретировать следующим образом: ускорение струны, обусловленное натяжением, в каждой точке тем больше, чем больше вогнутость струны   u_{xx}   в данной точке (коэффициентом пропорциональности служит константа   \alpha^2 = T/\rho  ) (см. рис. 133).
\end{verbatim}

\section{ЗАМЕЧАНИЯ}\label{ux437ux430ux43cux435ux447ux430ux43dux438ux44f-1}

\begin{verbatim}
1. Волновое уравнение   u_{tt} = \alpha^2 u_{xx}   описывает также продольные и крутильные колебания стержня. В продольных колебаниях смещения параллельны стержню, а через   u\left(x,\,t\right)   обозначается величина продольного смещения относительно равновесного положения. Такие колебания возникают, например, при ударе молотком по торцу стержня. Аналогично описываются крутиль-
\end{verbatim}

ные колебания

MATHBLK0589X

Здесь \emph{k} --- модуль Юнга, характеризующий упругость материала.
Материалы, обладающие большим модулем Юнга, колеблются на более высоких
частотах. Звуковые волны представляют собой продольные волны.

\begin{verbatim}
2. Если линейная плотность струны   \rho(x)   зависит от координаты, то волновое урагнение записывается в виде
\end{verbatim}

MATHBLK0590X

т. е. получается УЧП с переменными коэффициентами.

\begin{verbatim}
3. Поскольку волновое уравнение   u_{tt} = \alpha^2 u_{xx}   содержит производную по времени второго порядка, для получения единственного решения при t>0 необходимо задать   \partial sa   начальных условия:
\end{verbatim}

MATHBLK0591X

(начальное смещение струны), \(u_t\)(x, 0) = \(g(x)\) (начальная
скорость струны).

Этим волновое уравнение отличается от уравнения теплопроводности, для
которого необходимо задать только одно начальное условие.

\begin{enumerate}
\def\labelenumi{\arabic{enumi}.}
\setcounter{enumi}{3}
\tightlist
\item
  С помощью волнового уравнения можно описать распределение
  электрического тока в проводе. Из законов Кирхгофа получается
  следующая система двух уравнений с частными производными первого
  порядка:
\end{enumerate}

(16.3)

MATHBLK0592X

где

ж --- координата вдоль провода,

\emph{t} --- время,

\(i(x, t)\) --- распределение тока вдоль провода,

\begin{verbatim}
v\left(x,\ t\right)   — распределение потенциала вдоль провода,
\end{verbatim}

С -- емкость провода на единицу длины,

G --- утечка на единицу длины,

Р --- сопротивление на единицу длины,

-- индуктивность провода на единицу длины.

Уравнения (16.3) принято называть системой телеграфных уравнений. В
такой форме они обычно не используются. Для преобразования этой системы
продифференцируем первое уравнение по x, второе продифференцируем по t,
умножим обе части на C и вычтем из первого. В результате получаем

MATHBLK0593X

Воспользуемся теперь вторым уравнением из (16.3)

MATHBLK0594X

окончательно получаем

(16.4)

MATHBLK0595X

Это уравнение для тока, известное под названием \emph{телеграфного
уравнения}, является гиперболическим уравнением второго порядка (если
конечно ии C, ни L не равны нулю, в противном случае уравнение
становится параболическим).

Напряжение описывается точно таким же уравнением:

(16.5)

MATHBLK0596X

Если G = R = 0, то уравнения (16.4)---(16.5) упрощаются:

(16.6)

MATHBLK0597X

\section{ЗАДАЧИ}\label{ux437ux430ux434ux430ux447ux438-2}

\begin{itemize}
\tightlist
\item
  \begin{enumerate}
  \def\labelenumi{\arabic{enumi}.}
  \tightlist
  \item
    Получите уравнение (16.5) для \emph{и} из системы двух уравнений
    первого порядка (16.3)
  \end{enumerate}
\item
  \begin{enumerate}
  \def\labelenumi{\arabic{enumi}.}
  \setcounter{enumi}{1}
  \tightlist
  \item
    Зная физический смысл каждого члена волнового уравнения, что вы
    можете сказать о поведении во времени решения следующей задачи:
  \end{enumerate}
\end{itemize}

(YYII)

\begin{verbatim}
MATHBLK0598X,   0 < x < 1  ,   0 < t < \infty  ,
\end{verbatim}

(ГУ)

MATHBLK0599X

MATHBLK0600X

{[} \textbackslash begin\{cases\} u(0, t) = 0, \& 0 \textless{} t
\textless{} \textbackslash infty, \textbackslash{} u(1, t) = 0, \& 0
\textless{} t \textless{} \textbackslash infty,
\textbackslash end\{cases\} {]}

MATHBLK0601X

MATHBLK0602X

MATHBLK0603X

{[} \textbackslash begin\{cases\} u(x, 0) =
\textbackslash sin(\textbackslash pi x), \& 0 \textbackslash le x
\textbackslash le 1, \textbackslash{} u\_t(x, 0) = 0. \& 0
\textbackslash le x \textbackslash le 1, \textbackslash end\{cases\} {]}

MATHBLK0604X

MATHBLK0605X

MATHBLK0606X

Дайте физическую интерпретацию этой задачи.

\begin{enumerate}
\def\labelenumi{\arabic{enumi}.}
\setcounter{enumi}{3}
\tightlist
\item
  Для уравнения \(u_{tt} = u_{xx}\)найдите все решения вида
\end{enumerate}

MATHBLK0607X

Будет ли решением сумма двух решений такого вида?

\section{ФОРМУЛА
ДАЛАМБЕРА}\label{ux444ux43eux440ux43cux443ux43bux430-ux434ux430ux43bux430ux43cux431ux435ux440ux430}

ЦЕЛЬ ЛЕКЦИИ: Построить решение задачи Коши для волнового уравнения

MATHBLK0608X

MATHBLK0609X

Эта задача описывает движение бесконечной струны с заданными начальными
условиями. Она была решена в 1750 году французским математиком
Даламбером. Формула Даламбера

MATHBLK0610X

позволяет легко найти решение, если известны начальные условия. Кроме
того, она позволяет дать интересную физическую интерпретацию решения на
языке бегущих волн.

\begin{verbatim}
Как, наверное, помнит читатель, в параболническом случае мы сначала решали диффузионные задачи на конечном отрезке (методом разделения переменных), а уже затем перешли к неограниченным промежуткам (—   \infty < x < \infty  ), где применяли преобразование Фурье. В гиперболическом случае мы поступим наоборот. Начнем с решения одномерного волнового уравнения на всей действительной оси, т. е. с задачи Коши
\end{verbatim}

(17.1)

MATHBLK0611X

MATHBLK0612X

содержащей только начальные условия. Эту задачу можно было бы решить
либо с помощью преобразования Фурье (по переменной t), либо с помощью
преобразования Лапласа (по переменной t). Однако мы воспользуемся другим
методом (методом канонических координат), который, как мы надеемся,
заинтересует читателя. Этот метод базируется на тех же идеях, что и
метод перехода к движущейся системе отсчета, с которым мы познакомились
в лекции 15. Итак, приступим к решению задачи (17.1).

\section{Решение одномерного волнового уравнения. Формула
Даламбера}\label{ux440ux435ux448ux435ux43dux438ux435-ux43eux434ux43dux43eux43cux435ux440ux43dux43eux433ux43e-ux432ux43eux43bux43dux43eux432ux43eux433ux43e-ux443ux440ux430ux432ux43dux435ux43dux438ux44f.-ux444ux43eux440ux43cux443ux43bux430-ux434ux430ux43bux430ux43cux431ux435ux440ux430}

Решение задачи (17.1) разобьем на несколько шагов:

\begin{verbatim}
ШАГ 1 (Замена координат (x, t) новыми каноническими координатами   (\xi, \eta)  ).
\end{verbatim}

Для решения задачи (17.1) воспользуемся тем, что если заменить две
независимые переменные x и t новыми пространственновременными
координатами

MATHBLK0613X

MATHBLK0614X

то уравнение

MATHBLK0615X

преобразуется в уравнение

MATHBLK0616X.

В этом легко убедиться, если воспользоваться формулами преобразования
координат:

(17.2)

MATHBLK0617X

Подставляя эти выражения для производных в волновое уравнение, получим

MATHBLK0618X.. Первый шаг сделан.

ШАГ 2 (Решение преобразованного уравнения).

Преобразованное уравнение легко решить двумя последовательными
интегрированиями (сначала по переменной ξ, а затем по η). После
интегрирования по ξ получаем

MATHBLK0619X

\begin{verbatim}
— произвольная функция от   \eta  .
\end{verbatim}

Интегрирование последнего соотношения по \textbf{η} приводит к общему
решению

MATHBLK0620X

\begin{verbatim}
где   \phi(\eta)   — первообразная функция   \phi_1(\eta)  , а   \psi(\xi)   — произвольная функция. Итак, общее решение уравнения
\end{verbatim}

MATHBLK0621X

записывается в виде

(17.3)

MATHBLK0622X

\begin{verbatim}
где   \phi(\eta)   и   \psi(\xi)   — произвольные функции своих аргументов. Папример, читатель может проверить, что функции

MATHBLK0623X  u(\xi, \eta) = 1/\eta + tg \xi,    u(\xi, \eta) = \eta^2 + e

являются решениями уравнения   u_{\mathrm{ }nh}=0  . Второй шаг сделан.

\text{IIIA}\Gamma   3 (Возвращение к старым координатам x и t).
\end{verbatim}

Для нахождения общего решения, т. е. всех решений уравнения
\(u_{tt} = c^2 u_{xx}\),подставим

MATHBLK0624X

MATHBLK0625X

в решение

MATHBLK0626X

В результате получаем

(17.4)

MATHBLK0627X

Это общее решение волнового уравнения. С физической точки эрения оно
интересно тем, что представляет сумму двух бегущих волн Р И С

\pandocbounded{\includegraphics[keepaspectratio]{images/_page_126_Figure_14.jpeg}}

. 17.1. Волна \(u(x, t)\) = \(e^{-(x-ct)^2}\) движется слева направо

произвольной формы, движущихся в прогивоположных направлениях со
скоростью c.~Например, функции

MATHBLK0628X

(волна, Лвижущаяся слева направо),

MATHBLK0629X

(волна, движущаяся справа налево),

MATHBLK0630X

(две волны, движущиеся в противоположных направлениях)

являются типичными решениями волиового уравнения. На рисунке 17.1
изображена простейшая бегущая волна.

ШАГ 4 (Учет начальных условий).

Напомним читателю общий метод решения задачи Коши в теории обыкновенных
дифференциальных уравнений. Сначала находят общее решение уравнения, а
затем оно подставляется в начальные условия для того, чтобы найти
конкретные значения произвольных постоянных. У нас ситуация аналогична.
Для решения задачи Коши для волнового уравнения мы подставим общее
решение

MATHBLK0631X

(содержащее две произвольные функции) в начальные условия

MATHBLK0632X \(u_1\)(x, 0) = \(g(x)\), чтобы найти конкретные выражения
для произвольных функций ф и ф. После подстановки получаем

(17.5)

MATHBLK0633X

\begin{verbatim}
Если теперь проинтегрировать второе уравнение (17.5), то получится алгебранческая система двух уравнений относительно неизвестных функций   \varphi(x)   и   \psi(x)  . В самом деле, после интегрирования в пределах от   x_0   до x получаем
\end{verbatim}

(17.6)

MATHBLK0634X

\begin{verbatim}
Решаем теперь уравнение (17.6) совместно с первым уравнением (17.5) и получаем следующие выражения для функций   \varphi(x)   и   \psi(x)  :
\end{verbatim}

(17.7)

MATHBLK0635X

MATHBLK0636X

Следовательно, общее решение задачи (17.1) дается формулой

(17.8)

MATHBLK0637X

Цель достигнута. Мы получили общее решение вадачи (17.1). Решение вида
(17.8) принято называть формулой Даламбера. Читателю предлагается
самостоятельно проследить за расстановкой пределов интегрирования, чтобы
убедиться в том, что они меняются от x-ct до x+ct. Задача полностью
решена.

Прежде чем завершить лекцию, приведем несколько примеров, показывающих,
как применять формулу Даламбера в конкретных валачах.

\section{Примеры применения формулы
Даламбера}\label{ux43fux440ux438ux43cux435ux440ux44b-ux43fux440ux438ux43cux435ux43dux435ux43dux438ux44f-ux444ux43eux440ux43cux443ux43bux44b-ux434ux430ux43bux430ux43cux431ux435ux440ux430}

\section{1. Движение синусондальной
волны}\label{ux434ux432ux438ux436ux435ux43dux438ux435-ux441ux438ux43dux443ux441ux43eux43dux434ux430ux43bux44cux43dux43eux439-ux432ux43eux43bux43dux44b}

Рассмотрим начальные условия вида

MATHBLK0638X \(u_1\)(x, 0) = 0. По формуле Даламбера получаем решение

MATHBLK0639X

\begin{verbatim}
Его можно интерпретировать следующим образом. Начальное смещение струны   u(x, 0) = \sin x   делится на две одинаковые части
\end{verbatim}

MATHBLK0640X

{[} H \textbackslash frac\{\textbackslash sin x\}\{2\} . {]}

Каждая из частей распространяется со скоростью с в виде бегущих волн.
Одна из волн перемещается слева направо, а вторая---в противоположном
направлении. Читателю предлагается подумать над тем, как выглядит в этом
случае результирующая волна.

\section{2. Распространение простейшего прямоугольного
импульса}\label{ux440ux430ux441ux43fux440ux43eux441ux442ux440ux430ux43dux435ux43dux438ux435-ux43fux440ux43eux441ux442ux435ux439ux448ux435ux433ux43e-ux43fux440ux44fux43cux43eux443ux433ux43eux43bux44cux43dux43eux433ux43e-ux438ux43cux43fux443ux43bux44cux441ux430}

В этом случае начальные условия задаются следующим образом:

MATHBLK0641X

\(u_t\)(x, 0) = 0 (в остальных точках).

Поскольку начальное смещение расшепляется на две полуволны, движущиеся в
противоположных направлениях, результирующее движение будет таким, как
показано на рис. 17.2.

\pandocbounded{\includegraphics[keepaspectratio]{images/_page_129_Figure_2.jpeg}}

РИС, 17.2. Начальное возмущение распадается на две бегущие волны.

\section{3. Задана начальная
скорость}\label{ux437ux430ux434ux430ux43dux430-ux43dux430ux447ux430ux43bux44cux43dux430ux44f-ux441ux43aux43eux440ux43eux441ux442ux44c}

\begin{verbatim}
Предположим теперь, что в начальный момент струна находится в положении равновесия. Придадим струне начальную скорость (как в фортепиано) вида   \sin x  , т. е.
\end{verbatim}

MATHBLK0642X

MATHBLK0643X

Решение находим по формуле Даламбера

MATHBLK0644X

Оно представляет сумму двух бегущих косинусоидальных воли. Читатель
должен спросить себя, кажется ли ему это решение разумным.

На этом лекция 17 заканчивается. В следующей лекции мы покажем, как
формула Даламбера помогает дать интерпретацию в плоскости переменных x и
t.

\section{ЗАМЕЧАНИЯ}\label{ux437ux430ux43cux435ux447ux430ux43dux438ux44f-2}

\begin{itemize}
\item
  \begin{enumerate}
  \def\labelenumi{\arabic{enumi}.}
  \tightlist
  \item
    Общее решение уравнения с частными производными второго порядка
    содержит две произвольные функции, а общее решение обыкновенного
    дифференциального уравнения содержит две произвольные постоянные.
    Значит, у уравнения с частными производными больше решений, чем у
    обыкновенного дифференциального уравнения.
  \end{enumerate}
\item
  2 Замена переменных для приведения уравнения с частными производными к
  более простому виду является одним из общих методов теории. Новые
  координаты (ξ, η) принято называть каноническими координатами. Мы
  будем обращаться к ним и в дальнейшем, особенно при изучении
  гиперболических уравнений.
\item
  3 Метод нахождения общего решения уравнения с частными производными и
  подстановки этого решения в начальные условия не является общим
  методом решения уравнений с частными производными. Решение, которое
  рассматривается в данной лекции,--- единственный пример использования
  этого метода. Обычно невозможно найти общее решение уравнения с
  частными производными, но даже если удается это сделать, оказывается
  очень сложным подставить это решение в краевые условия.
\end{itemize}

\section{ЗАДАЧН}\label{ux437ux430ux434ux430ux447ux43d}

\begin{enumerate}
\def\labelenumi{\arabic{enumi}.}
\item
  Проверьте, что решение Даламбера (17-8) удовлетворяет начальным
  условиям задачи (17.1).

  \begin{enumerate}
  \def\labelenumii{\arabic{enumii}.}
  \setcounter{enumii}{1}
  \tightlist
  \item
    Подставьте (17.7) в общее решение u(x, t) = \textbackslash psi(x -
    ct) + \textbackslash psi(x + \textbackslash epsilon t) и получите
    формулу Даламбера.
  \end{enumerate}
\item
  Найдите решение задачи Коши:
\end{enumerate}

MATHBLK0645X

Как выглядит это решение в различные моменты времени? 4. Найдите решение
задачи Коши:

(УЧП)

MATHBLK0646X

(НУ)

MATHBLK0647X

Постройте график решения в различные моменты времени. 5. Решая совместно
уравнения (17.6) и первое уравнение из (17.5),

получите формулы (17.7).

\section{ФОРМУЛА ДАЛАМБЕРА
(ПРОДОЛЖЕНИЕ)}\label{ux444ux43eux440ux43cux443ux43bux430-ux434ux430ux43bux430ux43cux431ux435ux440ux430-ux43fux440ux43eux434ux43eux43bux436ux435ux43dux438ux435}

ЦЕЛЬ ЛЕКЦИИ! Показать, как с помощью формулы Даламбера можно найти
решение волнового уравнения на полуограниченной прямой

(УЧП)

MATHBLK0648X

\begin{verbatim}
(ГУ)   u(0, t) = 0, \quad 0 < t < \infty,
\end{verbatim}

(НУ)

MATHBLK0649X

Дополнительно дается интерпретация формулы Даламбера в координатной
плоскости x, t.

На прошлой лекции было показано, что если заданы начальное смещение
струны \(u(x, 0)\) = \(f(x)\) и начальная скорость \(u_t\)(x, 0) =
\(g(x)\) , то выражение

(18.i)

MATHBLK0650X

описывает смещение струны в произвольный момент времени. В этой лекции
мы познакомим читателя с одной интересной интерпретацией этой формулы в
плоскости переменных x и t (пространственновременная или фазовая
плоскость) и покажем, как изменить формулу Даламбера, чтобы получить
решение задачи о колебаниях полубесконечной струны.

Начнем с интерпретации формулы (18.1) в плоскости переменных x и t.

\section{Пространственно-временная интерпретация формулы
Даламбера}\label{ux43fux440ux43eux441ux442ux440ux430ux43dux441ux442ux432ux435ux43dux43dux43e-ux432ux440ux435ux43cux435ux43dux43dux430ux44f-ux438ux43dux442ux435ux440ux43fux440ux435ux442ux430ux446ux438ux44f-ux444ux43eux440ux43cux443ux43bux44b-ux434ux430ux43bux430ux43cux431ux435ux440ux430}

На предыдущей лекции мы показали, что решение задачи Коши

(18.2)

MATHBLK0651X

MATHBLK0652X

{[} u(x, t) = \textbackslash frac\{1\}\{2\} {[}f(x-ct) + f(x+ct){]} +
\textbackslash frac\{1\}\{2c\} \textbackslash int\_\{x-ct\}\^{}\{x+ct\}
g(\textbackslash xi) d\textbackslash xi. {]}

MATHBLK0653X

{[} u(x, 0) = f(x), {]}

MATHBLK0654X

{[} u(x, t) = \textbackslash frac\{1\}\{2\} {[}f(x-ct) + f(x+ct){]}. {]}

MATHBLK0655X

{[} x-ct=x\_0-ct\_0 {]}

MATHBLK0656X

{[} (Y \textbackslash text\{ ЧП\}) \textbackslash qquad u\_\{tt\} =
c\^{}2 u\_\{xx\}, \textbackslash quad -\textbackslash infty \textless{}
x \textless{} \textbackslash infty, \textbackslash quad 0 \textless{} t
\textless{} \textbackslash infty, {]}

MATHBLK0657X

{[} (HY) \textbackslash qquad {]}

MATHBLK0658X

MATHBLK0659X

MATHBLK0660X

На рис. 18.2 решение изображено в плоскости переменных x, t, а не u, x,
как в предыдущей лекции.

\pandocbounded{\includegraphics[keepaspectratio]{images/_page_133_Figure_5.jpeg}}

РИС. 18.2. Решение задачи Коши (18.2) в xt-плоскости a --- передний
фронт волны, движущейся налево; b --- задний фронт волны; b --- переднии
фронт волны, движущейся направо

Теперь займемся интерпретацией формулы Даламбера, когда начальное
смещение равно нулю, а начальная скорость произвольна.

СЛУЧАЙ 2 (Начальное смещение равно нулю, начальная скорость
произвольна).

Рассмотрим начальные условия вида

MATHBLK0661X \(u_t\)(x, 0) = \(g(x)\). Решение в этом случае имеет вид

MATHBLK0662X

\begin{verbatim}
Значение величины u в точке   (x_{\mathrm{ }e},\ t_{\mathrm{ }e})   можно интерпретировать как интеграл от начальной скорости в пределах от   x_{\mathrm{ }e}-ct_{\mathrm{ }e}   до Р И С
\end{verbatim}

\pandocbounded{\includegraphics[keepaspectratio]{images/_page_134_Figure_2.jpeg}}

. 18.3. Интерпретация задачи с начальной скоростью в x1-плоскости.

\begin{verbatim}
x_{\bullet}+ct_{\bullet}   (рис. 18.3). В качестве примера на рис. 18.4 приводится графическая интерпретация решения задачи
\end{verbatim}

(13.4)

MATHBLK0663X

MATHBLK0664X

MATHBLK0665X

{[} u(x, t) = \textbackslash frac\{1\}\{2c\}
\textbackslash int\_\{x-ct\}\^{}\{x+ct\} g(\textbackslash xi)
d\textbackslash xi = {]}

MATHBLK0666X

{[} = \textbackslash frac\{1\}\{2c\}
\textbackslash int\_\{x-ct\}\^{}\{x+ct\} 0 d\textbackslash xi,
\textbackslash quad (x, t) \textbackslash in \textbackslash text\{
области 1,\} {]}

MATHBLK0667X

{[} = \textbackslash frac\{1\}\{2c\}
\textbackslash int\_\{x-ct\}\^{}\{x+ct\} d\textbackslash xi,
\textbackslash quad (x, t) \textbackslash in \textbackslash text\{
области 2,\} {]}

MATHBLK0668X

{[} = \textbackslash frac\{1\}\{2c\}
\textbackslash int\_\{x-ct\}\^{}\{-1\} d\textbackslash xi,
\textbackslash quad (x, t) \textbackslash in \textbackslash text\{
области 3,\} {]}

MATHBLK0669X

{[} = \textbackslash frac\{1\}\{2c\}
\textbackslash int\_\{x-ct\}\^{}\{1\} d\textbackslash xi,
\textbackslash quad (x, t) \textbackslash in \textbackslash text\{
области 4,\} {]}

MATHBLK0670X

{[} = \textbackslash frac\{1\}\{2c\}
\textbackslash int\_\{x-ct\}\^{}\{x+ct\} 0 d\textbackslash xi,
\textbackslash quad (x, t) \textbackslash in \textbackslash text\{
области 5,\} {]}

MATHBLK0671X

{[} = \textbackslash frac\{1\}\{2c\}
\textbackslash int\_\{x-ct\}\^{}\{x+ct\} d\textbackslash xi,
\textbackslash quad (x, t) \textbackslash in \textbackslash text\{
области 6.\} {]}

MATHBLK0672X

MATHBLK0673X

{[} \textbackslash begin\{aligned\}
(\textbackslash text\{Y\}4\textbackslash Pi) \& u\_\{tt\} = c\^{}2
u\_\{xx\}, \& 0 \textless{} x \textless{} \textbackslash infty, \& 0
\textless{} t \textless{} \textbackslash infty, \textbackslash{}
(\textbackslash Gamma \textbackslash text\{Y\}) \& u(0, t) = 0, \& 0
\textless{} t \textless{} \textbackslash infty, \textbackslash{}
(\textbackslash text\{H\}\textbackslash text\{Y\}) \&
\textbackslash begin\{cases\} u(x, 0) = f(x), \textbackslash{} u\_t(x,
0) = g(x). \textbackslash end\{cases\} \& 0 \textless{} x \textless{}
\textbackslash infty, \textbackslash{} \textbackslash end\{aligned\} {]}

MATHBLK0674X

MATHBLK0675X

\section{Решение задачи для полубесконечной струны с помощью формулы
Даламбера}\label{ux440ux435ux448ux435ux43dux438ux435-ux437ux430ux434ux430ux447ux438-ux434ux43bux44f-ux43fux43eux43bux443ux431ux435ux441ux43aux43eux43dux435ux447ux43dux43eux439-ux441ux442ux440ux443ux43dux44b-ux441-ux43fux43eux43cux43eux449ux44cux44e-ux444ux43eux440ux43cux443ux43bux44b-ux434ux430ux43bux430ux43cux431ux435ux440ux430}

Наша задача---найти волновые движения полубесконечной струны, левый
конец которой закреплен, при заданных начальных условиях. Для решения
задачи (18.6) мы воспользуемся тем же приемом, что и при отыскании
общего решения волнового уравнения на всей прямой. Подставляя общее
решение

MATHBLK0676X

в начальные условия, находим (как в предыд, щей лекции)

(18.7)

MATHBLK0677X

MATHBLK0678X

\begin{verbatim}
"Теперь перед нами проблема, которой не было в задаче о бесконечной струне. Решение   u\left(x,\t\right)   должно быть определено всюду только внутри первого квадранта   (x>0,\,t>0)   плоскости переменных x и t. Значит, мы должны уметь вычислить функции
\end{verbatim}

MATHBLK0679X

\begin{verbatim}
при всех   -\infty < x-ct < \infty  ,   \psi(x+ct)   при всех   0 < x+ct < \infty  .

К сожалению, первая из формул (18.7) позволяет вычислять   \varphi(x-ct)   только при   x-ct\geqslant 0  , поскольку в начальных данных функции f(x) и g(x) определены только для положительных значений аргумента.

Если   x-ct \geqslant 0  , то после подстановки (18.7) в общее решение
\end{verbatim}

MATHBLK0680X

получаем

MATHBLK0681X

MATHBLK0682X

{[} \textbackslash begin\{split\}
u\textbackslash left(x,,t\textbackslash right) \&=
\textbackslash phi\textbackslash left(x-ct\textbackslash right) +
\textbackslash psi\textbackslash left(x+ct\textbackslash right) =
\textbackslash{} \&= \textbackslash frac\{1\}\{2\}
\textbackslash left{[} f\textbackslash left(x-ct\textbackslash right) +
f\textbackslash left(x+ct\textbackslash right) \textbackslash right{]} +
\textbackslash frac\{1\}\{2c\}
\textbackslash int\textbackslash limits\_\{x-ct\}\^{}\{x+ct\}
g\textbackslash left(\textbackslash xi\textbackslash right) ,
d\textbackslash xi. \textbackslash end\{split\} {]}

MATHBLK0683X

MATHBLK0684X

\begin{verbatim}
Что же делать, если x < ct? Пришла пора воспользоваться граничным условием u(0,t)=0. С помощью этого граничного условия мы доопределяем функцию   \varphi(x-ct)   при x < ct. Для этого подставим общее решение   u(x,t)=\varphi(x-ct)+\psi(x+ct)   в граничное условие u(0,t)=0. В результате получим   \varphi(-ct)=-\psi(ct)  ,
\end{verbatim}

откуда

MATHBLK0685X

Подстановка этого выражения для ф в общее решение дает

MATHBLK0686X

Комбинируя решения для x \textless{} ct и x \textgreater{} ct,
окончательно получаем (см. рис. 18.6)

MATHBLK0687X

MATHBLK0688X

{[} \textbackslash begin\{cases\} \textbackslash frac\{1\}\{2\}
\textbackslash left{[} f(x - ct) + f(x + ct) \textbackslash right{]} +
\textbackslash frac\{1\}\{2c\} \textbackslash int\_\{x - ct\}\^{}\{x +
ct\} g(\textbackslash xi) d\textbackslash xi, \& x \textbackslash ge ct,
\textbackslash{} \textbackslash frac\{1\}\{2\} \textbackslash left{[}
f(x + ct) - f(ct - x) \textbackslash right{]} +
\textbackslash frac\{1\}\{2c\} \textbackslash int\_\{ct - x\}\^{}\{x +
ct\} g(\textbackslash xi) d\textbackslash xi, \& x \textless{} ct.
\textbackslash end\{cases\} {]}

MATHBLK0689X

MATHBLK0690X

Этим завершается наша лекция. Интерпретация решения (18.8) будет
приведена в замечаниях.

\pandocbounded{\includegraphics[keepaspectratio]{images/_page_137_Figure_9.jpeg}}

\begin{verbatim}
РИС. 18.6. Интерпретация задачи Коши для полуограниченной струны в xt-плоскости a—возмущение отражается от границы;   \delta  —возмущение распространяется вдоль характеристики из точки (  x_0+ct_0  , 0).
\end{verbatim}

\section{ЗАМЕЧАНИЯ}\label{ux437ux430ux43cux435ux447ux430ux43dux438ux44f-3}

\begin{verbatim}
1. Решение (18.8) полностью соогветствует нашим представлениям о полубесконечной струне с граничным условием u(0, t) = 0. При   x \geqslant ct   оно совпадает с формулой Даламбера для бесконечной струны. Если x < ct, то формула модифицируется таким образом, что в нее включается волна, отраженная от границы (при отражении знак волны меняется на противоположный).
\end{verbatim}

\begin{enumerate}
\def\labelenumi{\arabic{enumi}.}
\setcounter{enumi}{1}
\tightlist
\item
  Если граничное условие \(u(0, t)\) = 0 изменить на другое, то и
  решение (18.8), естественно, изменится. Можно получить в явном виде
  решения для граничных условий
\end{enumerate}

MATHBLK0691X

или

MATHBLK0692X

Дополнительную информацию по этим вопросам читатель может найти в
литературе, указанной в конце книги.

\begin{enumerate}
\def\labelenumi{\arabic{enumi}.}
\setcounter{enumi}{2}
\item
  Прямые линии

  MATHBLK0693X x - ct = \textbackslash text\{const\} называются
  характеристиками. Возмущение распространяется вдоль характеристик.
  Понятие характеристики тесно связано с уравнениями гиперболического
  типа.
\end{enumerate}

\section{ЗАДАЧИ}\label{ux437ux430ux434ux430ux447ux438-3}

\begin{enumerate}
\def\labelenumi{\arabic{enumi}.}
\tightlist
\item
  Решите смещанную задачу для полубесконечной струны
\end{enumerate}

MATHBLK0694X

MATHBLK0695X

Постройте графики решения для разл чных моментов времени.

\begin{enumerate}
\def\labelenumi{\arabic{enumi}.}
\setcounter{enumi}{1}
\tightlist
\item
  Для решения предыдущей задачи можно воспользоваться следующим
  алгоритмом:
\end{enumerate}

а) продолжить начальные условия на всю действительную ось по формулам

\begin{verbatim}
MATHBLK0696X  u_t(x, 0) = 0, -\infty < x < 0.  б) усреднить две волны, движущиеся налево и направо (как это было в предыдущей лекции);

**в**) рассмотреть полученное решение при   x \ge 0  .
\end{verbatim}

Используя этот алгоритм, постройте графики решения для различных
моментов времени t, если начальное условие задано графически.

\pandocbounded{\includegraphics[keepaspectratio]{images/_page_139_Figure_3.jpeg}}

\begin{enumerate}
\def\labelenumi{\arabic{enumi}.}
\setcounter{enumi}{2}
\tightlist
\item
  Решите смешанную задачу для полубесконечной струны
\end{enumerate}

MATHBLK0697X

MATHBLK0698X

Воспользуйтесь методом, аналогичным рассмотренному в лекции. Дайте
физическую интерпретацию этой задачи.

\begin{enumerate}
\def\labelenumi{\arabic{enumi}.}
\setcounter{enumi}{3}
\tightlist
\item
  Предположим, что колебания струны описываются уравнением
  \(u_{tt} = u_{xx}\),начальное смещение изображено на рисунке, а
  начальная скорость равна нулю. Изобразите решение этой задачи в
  xt-плоскости. Обратите внимание на то, что начальные условия в этой
  задаче разрывны.
\end{enumerate}

\pandocbounded{\includegraphics[keepaspectratio]{images/_page_139_Figure_8.jpeg}}

\section{ВОЛНОВОЕ УРАВНЕНИЕ И ГРАНИЧНЫЕ
УСЛОВИЯ}\label{ux432ux43eux43bux43dux43eux432ux43eux435-ux443ux440ux430ux432ux43dux435ux43dux438ux435-ux438-ux433ux440ux430ux43dux438ux447ux43dux44bux435-ux443ux441ux43bux43eux432ux438ux44f}

ЦЕЛЬ ЛЕКЦИИ: Познакомыть читателя с основными типами граничных условий
для волнового уравнения

MATHBLK0699X

на конечном отрезке.

В дальнейшем будут рассматриваться три основных типа граничных условий.

\begin{enumerate}
\def\labelenumi{\arabic{enumi}.}
\tightlist
\item
  Заданный режим (граничные условия 1-го рода)
\end{enumerate}

MATHBLK0700X \(u(L, t)\) = \(g_2\)(t). 2. Заданные силы (граничные
условия 2-го рода)

MATHBLK0701X \(u_x\)(L, t) = \(g_2\)(t). 3. Упругое закрепление
(граничные условия 3-го рода)

\begin{verbatim}
MATHBLK0702X  u_x(L, t) - \gamma_2 u(L, t) = g_2(t).  Могут встречаться и произвольные комбинации этих трех типоз граничных условий. В лекции будет показан физический смысл этих граничных условий.
\end{verbatim}

Итак, пока что мы рассмотрели единственный вид волнового
движения---поперечные колебания струны. Читатель должен понимать, что
это только один из видов волновых движений.

Ниже приведен перечень основных типов волн:

\begin{itemize}
\tightlist
\item
  \begin{enumerate}
  \def\labelenumi{\arabic{enumi}.}
  \tightlist
  \item
    Звуковые волны (продольные волны).
  \end{enumerate}
\item
  \begin{enumerate}
  \def\labelenumi{\arabic{enumi}.}
  \setcounter{enumi}{1}
  \tightlist
  \item
    Электромагнитные волны, в том числе световые.
  \end{enumerate}
\item
  \begin{enumerate}
  \def\labelenumi{\arabic{enumi}.}
  \setcounter{enumi}{2}
  \tightlist
  \item
    Колебания твердых тел (продольные, поперечные и крутильные).
  \end{enumerate}
\item
  \begin{enumerate}
  \def\labelenumi{\arabic{enumi}.}
  \setcounter{enumi}{3}
  \tightlist
  \item
    Волны вероятности в квантовой механике.
  \end{enumerate}
\item
  \begin{enumerate}
  \def\labelenumi{\arabic{enumi}.}
  \setcounter{enumi}{4}
  \tightlist
  \item
    Волны на воде (поперечные волны).
  \end{enumerate}
\item
  \begin{enumerate}
  \def\labelenumi{\arabic{enumi}.}
  \setcounter{enumi}{5}
  \tightlist
  \item
    Колебания струны (поперечные волны).
  \end{enumerate}
\end{itemize}

Цель нашей лекции --- обсудить различные типы граничных условий, которые
возникают при решении физических задач волнового характера. Мы
остановимся только на одномерных задачах и линейных граничных условиях.
Обычно различают граничные условия трех типов

\begin{enumerate}
\def\labelenumi{\arabic{enumi}.}
\tightlist
\item
  Заданные режимы в граничных точках (граничные условия 1-го рода)
\end{enumerate}

MATHBLK0703X \(u(L, t)\) = \(g_2\)(t). 2. Заданные силы в граничных
точках (граничные условия 2-го рода)

MATHBLK0704X \(u_x\)(L, t) = \(g_2\)(t). 3. Упругое закрепление в
граничных точках (граничные условия 3-го рода)

\begin{verbatim}
MATHBLK0705X  u_x(L, t) - \gamma_2 u(L, t) = g_2(t).  Мы начнем с граничных условий 1-го рода.
\end{verbatim}

\begin{enumerate}
\def\labelenumi{\arabic{enumi}.}
\tightlist
\item
  Заданные режимы в граничных точках.
\end{enumerate}

Нам предстоит решить задачу вида

(19.1)

MATHBLK0706X

MATHBLK0707X

{[} \textbackslash begin\{aligned\} u\_\{tt\} \&= c\^{}2 u\_\{xx\}, \& 0
\textless{} x \textless{} 1, \& 0 \textless{} t \textless{}
\textbackslash infty, \textbackslash{} u(0, t) \&= g\_1(t), \& 0
\textless{} t \textless{} \textbackslash infty, \textbackslash{} u(1, t)
\&= g\_2(t), \& 0 \textless{} t \textless{} \textbackslash infty,
\textbackslash{} (HY) \&\textbackslash begin\{cases\} u(x, 0) \&= f(x),
\textbackslash{} u\_t(x, 0) \&= g(x), \textbackslash end\{cases\} \& 0
\textbackslash leqslant x \textbackslash leqslant 1,
\textbackslash end\{aligned\} {]}

MATHBLK0708X

MATHBLK0709X

{[} u\_x(0, t) = 0, \textbackslash quad u\_x(L, t) = 0. {]}

MATHBLK0710X

{[} u\_\{x\}(1, t) = \textbackslash frac\{1\}\{k\} v(t) - (k -
\textbackslash text\{модуль Юнга\}). {]}

MATHBLK0711X

{[} u\_\{x\}(0, t) = \textbackslash frac\{h\}\{T\} u(0, t), {]}

MATHBLK0712X

{[} u\_\{x\}(L, t) = -\textbackslash frac\{h\}\{T\} u(L, t). {]}

MATHBLK0713X

{[} u\_\{x\}(0, t) - \textbackslash frac\{h\}\{T\}u(0, t) = 0, {]}

MATHBLK0714X

{[} u\_\{x\}(L, t) + \textbackslash frac\{h\}\{T\}u(L, t) = 0. {]}

MATHBLK0715X

MATHBLK0716X

{[} \textbackslash begin\{aligned\} u\_x(0, T) \&=
\textbackslash frac\{h\}\{T\} {[}u(0, t) -
\textbackslash theta\_1(t){]}, \textbackslash{} u\_x(L, T) \&=
-\textbackslash frac\{h\}\{T\} {[}u(L, t) -
\textbackslash theta\_2(t){]}. \textbackslash end\{aligned\} {]}

MATHBLK0717X

MATHBLK0718X

Этим мы завершаем обсуждение наиболее важных граничных условий,
связанных с гиперболическими задачами. Следующие несколько лекций мы
посвятим решению задач, содержащих полобные граничные условия.

\section{ЗАМЕЧАНИЯ}\label{ux437ux430ux43cux435ux447ux430ux43dux438ux44f-4}

\begin{enumerate}
\def\labelenumi{\arabic{enumi}.}
\tightlist
\item
  В лекции не рассматривалось граничное условие еще одного типа, когда
  струна подвергается действию силы, пропорциональной скорости и
  направленной в противоположном направ-
\end{enumerate}

лении. Граничное условие такого типа (на левом конце) записывается в
виде

MATHBLK0719X

\begin{enumerate}
\def\labelenumi{\arabic{enumi}.}
\setcounter{enumi}{1}
\tightlist
\item
  Нелинейное упругое закрепление левого конца приводит к граничному
  условию вида
\end{enumerate}

MATHBLK0720X

\begin{verbatim}
где   \varphi(u)   — произвольная функция величины u. Например, граничное условие
\end{verbatim}

MATHBLK0721X

говорит о том, что возвращающая сила на левом конце пропорциональна кубу
смещения (а не первой степени и, как это будет согласно закону Гука в
линейном случае).

\begin{enumerate}
\def\labelenumi{\arabic{enumi}.}
\setcounter{enumi}{2}
\tightlist
\item
  Если к нижнему концу пружины, совершающей продольные колебания,
  прикрепить массу \emph{m}, то граничное условие следует задавать в
  виде
\end{enumerate}

MATHBLK0722X

\section{ЗАДАЧИ}\label{ux437ux430ux434ux430ux447ux438-4}

Опираясь на интуитивные представления о граничных условиях различных
типов, изобразить в общих чертах решение задачи

MATHBLK0723X

MATHBLK0724X

\begin{verbatim}
\begin{aligned} (\mathrm{Y} \, \Pi) & u_{tt} = c^2 u_{xx}, \quad 0 < x < 1, \quad 0 < t < \infty, \\ (\mathrm{\Gamma} \, \mathrm{Y}) & \left\{ \begin{aligned} u \, (0, \, t) = 0, \\ u \, (1, \, t) = \sin t, \end{aligned}
\end{verbatim}

MATHBLK0725X

MATHBLK0726X

MATHBLK0727X

{[} \textbackslash begin\{aligned\} u , (x, , 0) = 0, \textbackslash{}
u\_\{t\} , (x, , 0) = 0, \textbackslash end\{aligned\} {]}

MATHBLK0728X

\begin{verbatim}
\right. & 0 \leqslant x \leqslant 1, \end{aligned}
\end{verbatim}

MATHBLK0729X

MATHBLK0730X

MATHBLK0731X

в различные моменты времени. Могли бы вы угадать решение этой задачи?

\begin{enumerate}
\def\labelenumi{\arabic{enumi}.}
\setcounter{enumi}{2}
\tightlist
\item
  Во что перейдет граничное условие
\end{enumerate}

MATHBLK0732X

npit

\begin{verbatim}
- a)   h \rightarrow \infty  ,
    - 6)   h \rightarrow 0  .
\end{verbatim}

Согласуется ли результат с вашей интунцией?

\begin{enumerate}
\def\labelenumi{\arabic{enumi}.}
\setcounter{enumi}{3}
\tightlist
\item
  Изобразить в общих чертах решение задачи
\end{enumerate}

(УЧП)

MATHBLK0733X

(ГУ)

MATHBLK0734X

MATHBLK0735X

{[} \textbackslash begin\{cases\} u(0, t) = 0, \textbackslash{} u\_x(1,
t) = -u(1, t), \textbackslash{} u(x, 0) = x, \textbackslash{} u\_t(x, 0)
= 0. \textbackslash end\{cases\} {]}

MATHBLK0736X

MATHBLK0737X

MATHBLK0738X

{[} \textbackslash begin\{cases\} u(x, 0) = x, \textbackslash{} u\_t(x,
0) = 0. \textbackslash end\{cases\} {]}

MATHBLK0739X

MATHBLK0740X

MATHBLK0741X

\begin{verbatim}
можно найти стандартным методом разделения переменных. При этом оказывается, что решение   u\left(x,\,t\right)  , представленное в виде ряда
\end{verbatim}

MATHBLK0742X

является суперпозицией простейших колебаний. Пространственная часть
\(X_n\)(x) каждого элементарного колебания является решением
(собственной функцией) некоторой краевой задачи Штурма---Лиувилля.

\begin{verbatim}
Итак, мы научились решать волиовое уравнение   u_{II}=c^2u_{xx}   на всей прямой   -\infty < x < \infty  . Формула Даламбера представляет решение в виде суммы двух бегущих волн, распространяющихся в противоположные стороны. Если же рассматривать это же уравнение на ограниченном промежутке 0 < x < L, то бегущих волн уже не станет, так как они будут взаимодействовать с границами. Вместо них возникают другие волны, которые называют стоячими волнами. Посмотрим, например, что случится, если гитарную струну, закрепленную в точках x=0, x=L, привести в движение. Для ответа на этот вопрос нам необходимо решить следующую смешанную задачу гиперболического типа:
\end{verbatim}

MATHBLK0743X

MATHBLK0744X

Оказывается, в этом случае бегущие волны отражаются от границ таким
образом, что результирующие колебания становятся не бегущими, а
сохраняющими форму в одном положении, т. е.

\pandocbounded{\includegraphics[keepaspectratio]{images/_page_147_Picture_5.jpeg}}

РИС. 20.1. Три типичных стоячих волны \(X(x)\) \(T(t)\).

превращаются в стоячие волны Несколько примеров стоячих волн показано на
рис. 20.1. Если известны профили этих стоячих волн \(X_n\)(x) и характер
колебаний каждой волны \(T_n\)(t) , то решение задачи о колебаниях
гитарной струны представляет собой суперпозицию простейших колебаний
\(X_n\)(x),\(T_n\)(t) и запи-

сывается в виде

MATHBLK0745X

\begin{verbatim}
Единственное, что остается сделать нам,— это подобрать коэффициенты   c_n   так, чтобы при t=0 удовлетворялись начальные условия   u\left(x,\ 0\right)=f\left(x\right)   и   u_t\left(x,\ 0\right)=g\left(x\right)  .
\end{verbatim}

Задачу о колебаниях гитарной струны мы решим методом разделения
переменных.

\section{Решение задачи о колебаниях ограниченной струны методом
разделения
переменных}\label{ux440ux435ux448ux435ux43dux438ux435-ux437ux430ux434ux430ux447ux438-ux43e-ux43aux43eux43bux435ux431ux430ux43dux438ux44fux445-ux43eux433ux440ux430ux43dux438ux447ux435ux43dux43dux43eux439-ux441ux442ux440ux443ux43dux44b-ux43cux435ux442ux43eux434ux43eux43c-ux440ux430ux437ux434ux435ux43bux435ux43dux438ux44f-ux43fux435ux440ux435ux43cux435ux43dux43dux44bux445}

Чтобы решить смешанную задачу

(УЧП)

MATHBLK0746X

MATHBLK0747X

MATHBLK0748X

{[} \textbackslash begin\{cases\} u(0, t) = 0, \textbackslash{} u(L, t)
= 0, \textbackslash end\{cases\} {]}

MATHBLK0749X

MATHBLK0750X

MATHBLK0751X

MATHBLK0752X

{[} \textbackslash begin\{cases\} u(x, 0) = f(x), \textbackslash{}
u\_t(x, 0) = g(x), \textbackslash end\{cases\} {]}

MATHBLK0753X

MATHBLK0754X

найдем сначала стоячне волны, т. е. решения вида

MATHBLK0755X

Подставив это выражение в уравнение (20.1) и разделив переменные,
получим два обыкновенных дифференциальных уравнения

MATHBLK0756X

MATHBLK0757X

\begin{verbatim}
где константа   \lambda   — пока что любое действительное число —   \infty < \lambda < \infty   .

Результаты исследования решений этих двух уравнений при различных значениях   \lambda   приведены на рис 20.2.

Теперь мы должны *опбросить* те решения, которые либо неограничены при   t \to \infty  , либо становятся тождественно равными нулю, если потребовать выполнения граничных условий u(0, t) = u(L, t) = 0.

В качестве упражнения читателю предлагается проверить, что нетривиальные и ограниченные решения можно получить только при отрицательных значениях параметра   \lambda  . Таким образом, наша задача— найти такие значения констант A, B, C, D и отрица-

тельной константы   \lambda = -\beta^2  , чтобы выражение
\end{verbatim}

(20.2)

MATHBLK0758X

удовлетворяло граничным условиям. Если это удастся, то мы будем
располагать множеством простейших колебаний струны, каждое из которых
удовлетворяет уравнению с частными произ- Р И С

\pandocbounded{\includegraphics[keepaspectratio]{images/_page_149_Figure_5.jpeg}}

20.2. Построение стоячих воли при раздичных значениях параметра λ.

водными и граничным условиям. Затем останется только найти такую
линейную комбинацию этих колебаний, чтобы сумма удовлетворяла начальным
условиям при t=0.

\begin{verbatim}
Подстановка (20.2) в граничные условия u(0, t) = u(L, t) = 0 дает   u(0, t) = X(0) T(t) = D[A \sin(\alpha\beta t) + B \cos(\alpha\beta t)] = 0 \Rightarrow D = 0  , (20.3)   u(L, t) = X(L) T(t) = C \sin(\beta L) [A \sin(\alpha\beta t) + B \cos(\alpha\beta t)] = 0 \Rightarrow \sin(\beta L) = 0  .

Другими словами, константа разделения   \beta   (мы можем вместо   \lambda   искать константу   \beta  ) должна удовлетворять уравнению   \sin{(\beta L)} = 0  , откуда
\end{verbatim}

MATHBLK0759X

\begin{verbatim}
Заметим, что если во втором уравнении (20.3) положить C=0, то получится тривиальное решение   X\left(x\right)T\left(t\right)\equiv0  . Таким образом, мы нашли последовательность элементарных колебаний струны (будем помечать их индексом n)
\end{verbatim}

MATHBLK0760X

(20.4)

или

MATHBLK0761X

\begin{verbatim}
где   a_n  ,   b_n  ,   R_n   и   \delta_n  —произвольные постоянные. Заметим, что элементарные колебания удовлетворяют граничным условиям. Читателю должно быть ясно, что каждое элементарное колебание представляет собой стоячую волну. В стоячей волне все точки колеблются с одной и той же частотой, а форма волны подобна волнам, изображенным на рис. 20.3.
\end{verbatim}

\pandocbounded{\includegraphics[keepaspectratio]{images/_page_150_Picture_5.jpeg}}

РИС. 20.3. Стоячие волны \(u_n\)(x, t) = \(X_n\)(x) \(T_n\)(t) .

В нашей задаче уравнение и граничные условия линейны и однородны,
поэтому произвольная сумма элементарных стоячих волн также является
решением волнового уравнения и удовлетворяет граничным условиям. Мы
должны теперь найти такую суперпозицию стоячих волн, чтобы она
удовлетворяла также и начальным условиям, тогда получим решение нашей
задачи. Подстановка суммы

MATHBLK0762X

в начальные условия

\begin{verbatim}
MATHBLK0763X  u_t(x, \theta) = g(x)  приводит к двум уравнениям
\end{verbatim}

MATHBLK0764X

MATHBLK0765X

Воспользовавшись соотношениями ортогональности

MATHBLK0766X

MATHBLK0767X

{[} \textbackslash begin\{cases\} 0, \& m \textbackslash neq n,
\textbackslash{} L/2, \& m = n, \textbackslash end\{cases\} {]}

MATHBLK0768X

MATHBLK0769X

находим коэффициенты \(a_n\) и \(b_n\)

(20.5)

MATHBLK0770X

Итак, задача решена. Ее решение записывается в виде

(20.6)

MATHBLK0771X

а коэффициенты \(a_n\) и \(b_n\) находятся по формулам (20.5). Прежде
чем закончить лекцию, сделаем несколько замечаний.

\section{ЗАМЕЧАНИЯ}\label{ux437ux430ux43cux435ux447ux430ux43dux438ux44f-5}

\begin{enumerate}
\def\labelenumi{\arabic{enumi}.}
\tightlist
\item
  Если начальная \emph{скорость} струны равна нулю, то решение (20.6)
  принимает вид
\end{enumerate}

MATHBLK0772X

и его можно интерпретировать следующим образом. Предположим, что мы
разложили начальное смещение струны

MATHBLK0773X

в ряд по синусам

MATHBLK0774X

Тогда каждый член такого разложения должен совершать колебания

MATHBLK0775X

Если теперь просуммировать все эти колебания, то получится решение нашей
задачи. Например, предположим, что

\begin{verbatim}
начальное смещение струны задается функцией   f(x) = \sin(\pi x/L) + 0.5 \sin(3\pi x/ L) + 0,25 \sin(5\pi x/L)  .
\end{verbatim}

Полный отклик системы на начальное условие такого вида равен сумме
откликов на каждое слагаемое, т. е.

MATHBLK0776X

\begin{enumerate}
\def\labelenumi{\arabic{enumi}.}
\setcounter{enumi}{1}
\tightlist
\item
  Слагаемое с номером n в разложении (20.6) называется \textbf{п-й модой
  колебаний} или \textbf{п-й гармоникой}. Используя тригонометрические
  формулы, можно записать эти гармоники в виде
\end{enumerate}

MATHBLK0777X

\begin{verbatim}
где   R_n   и   \delta_n   — новые произвольные постоянные (амплитуда и фазовый угол). Эта новая форма представления n-й моды оказывается более полезной при анализе колебаний. Частота колебаний од (рад/с) и-й моды находится по формуле
\end{verbatim}

MATHBLK0778X

где Т --- натяжение струны, р --- линейная плотность струны. Отметим,
также, что частота n-й гармоники в n раз больше основной частоты (n=1).
Это свойство частоты характерно не для всех типов колебаний. Приятный
звук скривичной или гитарной струны обусловлен кратностью частогы
обертонов основной частоте. В барабане частоты гармоник высокого порядка
не кратны основной частоте колебаний.

\section{ЗАДАЧИ}\label{ux437ux430ux434ux430ux447ux438-5}

\begin{enumerate}
\def\labelenumi{\arabic{enumi}.}
\tightlist
\item
  Найдите решение задачи (20.1), если начальные условия задаются
  формулами
\end{enumerate}

MATHBLK0779X \(u_t\)(x, 0) = 0. Постройте графики решения для различных
моментов времени. Является ли это решение периодическим по времени?
Каков его период?

\begin{enumerate}
\def\labelenumi{\arabic{enumi}.}
\setcounter{enumi}{1}
\item
  Найдите решение задачи о колебаниях струны (20.1) при начальных
  условиях вида

  MATHBLK0780X u\_t(x, 0) = \textbackslash sin(3\textbackslash pi x/L).
  Как выглядит график решения в различные моменты времени?

  \begin{itemize}
  \tightlist
  \item
    \begin{enumerate}
    \def\labelenumii{\arabic{enumii}.}
    \setcounter{enumii}{2}
    \tightlist
    \item
      Покажите, что если на рис. 20.2 величина \textbackslash lambda
      \textbackslash geqslant 0 , то решения X(x) T(t) будут либо не
      ограничены, либо тождественно равны нулю.
    \end{enumerate}
  \end{itemize}
\end{enumerate}

\begin{itemize}
\item
  Найдите решение задачи о колебаниях струны при начальных условиях вида

  MATHBLK0781X u\_1(x, 0) = (3\textbackslash pi
  \textbackslash alpha/L)\textbackslash sin(3\textbackslash pi x/L). 5.
  Гитарная струна длины L=1 оттягивается за среднюю точку на величину h
  (см. рис.) Начальное положение струны можно задать в виде
\end{itemize}

MATHBLK0782X

MATHBLK0783X

MATHBLK0784X

MATHBLK0785X

{[} \textbackslash begin\{aligned\} (\textbackslash mathrm\{Y\}
\textbackslash Pi \textbackslash Pi) \& u\_\{tt\} =
\textbackslash alpha\^{}2 u\_\{xx\} - \textbackslash beta u\_t, \& 0
\textless{} x \textless{} 1, \& 0 \textless{} t \textless{}
\textbackslash infty, \textbackslash{} (\textbackslash Gamma
\textbackslash mathrm\{Y\}) \& \textbackslash begin\{cases\} u ~(0, ~t)
= 0, \textbackslash{} u ~(1, ~t) = 0, \textbackslash end\{cases\} \& 0
\textless{} t \textless{} \textbackslash infty, \textbackslash{}
(\textbackslash mathrm\{H\} \textbackslash mathrm\{Y\}) \&
\textbackslash begin\{cases\} u ~(x, ~0) = f ~(x), \textbackslash{} u\_1
~(x, ~0) = 0. \textbackslash end\{cases\} \& 0 \textbackslash leqslant x
\textbackslash leqslant 1, \textbackslash end\{aligned\} {]}

MATHBLK0786X

MATHBLK0787X

MATHBLK0788X

{[} \textbackslash begin\{aligned\} (\textbackslash mathrm\{Y\}
\textbackslash mathrm\{H\} \textbackslash Pi) \& u\_\{tt\} =
\textbackslash alpha\^{}2 u\_\{xx\} + Kx, \textbackslash quad 0
\textless{} x \textless{} 1, ~0 \textless{} t \textless{}
\textbackslash infty, \textbackslash{}
(\textbackslash mathrm\{\textbackslash Gamma\}
\textbackslash mathrm\{Y\}) \& \textbackslash begin\{cases\} u ~(0, ~t)
= 0, \textbackslash{} u ~(1, ~t) = 0, \textbackslash end\{cases\} \& 0
\textless{} t \textless{} \textbackslash infty, \textbackslash{}
(\textbackslash mathrm\{H\} \textbackslash mathrm\{Y\}) \&
\textbackslash begin\{cases\} u ~(x, ~0) = f ~(x), \textbackslash{} u\_t
~(x, ~0) = 0. \textbackslash end\{cases\} \& 0 \textbackslash leqslant x
\textbackslash leqslant 1, \textbackslash end\{aligned\} {]}

MATHBLK0789X

MATHBLK0790X

\section{КОЛЕБАНИЯ БАЛКИ (УРАВНЕНИЕ С ЧАСТНЫМИ ПРОИЗВОДНЫМИ ЧЕТВЕРТОГО
ПОРЯДКА)}\label{ux43aux43eux43bux435ux431ux430ux43dux438ux44f-ux431ux430ux43bux43aux438-ux443ux440ux430ux432ux43dux435ux43dux438ux435-ux441-ux447ux430ux441ux442ux43dux44bux43cux438-ux43fux440ux43eux438ux437ux432ux43eux434ux43dux44bux43cux438-ux447ux435ux442ux432ux435ux440ux442ux43eux433ux43e-ux43fux43eux440ux44fux434ux43aux430}

ЦЕЛЬ ЛЕКЦИИ: Показать, каким образом задача о колебании балки приводит к
уравнению четвертого порядка и как эта задача решается в случае свободно
опирающихся концов. Проводится сравнение колебаний балки с колебаниями
скрипичной струны.

\pandocbounded{\includegraphics[keepaspectratio]{images/_page_154_Figure_5.jpeg}}

Главное отличие поперечных колебаний тонкой балки от поперечных
колебаний струны в том, что балка оказывает сопротивление изгибу. Не
вдаваясь в механику тонких балок, мы примем, что учет сопротивления
изгиба приводит (вместо волнового уравнения) к уравнению четвертого
порядка

(21.1)

MATHBLK0791X

{[} rge \textbackslash alpha\^{}2 = K/\textbackslash rho , {]}

K---модуль сдвига (чем больше K, тем жестче балка, тем выше частота
колебаний),

р --- линейная плотность балки (масса/ед. длины).

Поскольку читатель впервые встречается в нашей книге с уравнением выше
чем второго порядка, будет полезно решить типовую задачу о колебаниях
балки. Позже мы познакомимся и с другими задачами теории балок.

\section{Свободно опирающаяся
балка}\label{ux441ux432ux43eux431ux43eux434ux43dux43e-ux43eux43fux438ux440ux430ux44eux449ux430ux44fux441ux44f-ux431ux430ux43bux43aux430}

Рассмотрим малые колебания тонкой балки, концы которой свободно
опираются на две опоры. Говоря «свободно опираются», мы подразумеваем,
что концы балки не перемещаются, но на-

клоны балки в концевых точках могут изменяться (концы балки вакреплены с
помощью штифтовых устройств, рис. 21.1).

Совершенно ясно, что на концах балки должны выполняться граничные
условия

MATHBLK0792X \(u(1, t)\) = 0, но совсем не очевидно, что на концах балки
должны удовлетворяться еще два граничных условия:

MATHBLK0793X \(u_{xx}\)(1, t) = 0. Используя теорию тонких балок, можно
показать, что изгибающий момент в балке пропорционален величине
\(u_{xx}\) , а изгибающий момент в свободно опирающемся конце должен
быть ра- Р И С

\pandocbounded{\includegraphics[keepaspectratio]{images/_page_155_Figure_8.jpeg}}

.21.1. Свободно оппрающаяся балка.

вен нулю. Следовательно, изображенная на рис. 21.1 колеблющаяся балка
описывается следующей смешанной задачей (для простоты величина о принята
равной единице):

(21.2)

MATHBLK0794X

MATHBLK0795X

{[} \textbackslash begin\{aligned\} \& (\textbackslash text\{Y\}
\textbackslash cdot \textbackslash Pi) \textbackslash quad u\_\{tt\} =
-u\_\{xxxx\}, \textbackslash quad 0 \textless{} x \textless{} 1,
\textbackslash quad 0 \textless{} t \textless{} \textbackslash infty,
\textbackslash{} \& \textbackslash begin\{cases\} u \textbackslash cdot
(0, ~t) = 0, \textbackslash{} u\_\{xx\} \textbackslash cdot (0, ~t) = 0,
\textbackslash{} u \textbackslash cdot (1, ~t) = 0, \textbackslash{}
u\_\{xx\} \textbackslash cdot (1, ~t) = 0, \textbackslash end\{cases\}
\textbackslash quad 0 \textless{} t \textless{} \textbackslash infty,
\textbackslash{} \& \textbackslash begin\{cases\} u \textbackslash cdot
(x, ~0) = f(x), \textbackslash{} u\_\{1\}(x, ~0) = g(x).
\textbackslash end\{cases\} \textbackslash quad 0
\textbackslash leqslant x \textbackslash leqslant 1,
\textbackslash end\{aligned\} {]}

MATHBLK0796X

MATHBLK0797X

Для решения этой задачи воспользуемся методом разделения переменных.
Будем искать только периодические решения, т. е. колебания вида

(21.3)

MATHBLK0798X

Отметим, что выбор решения вида (21.3) на самом деле означает, что
константу разделения в методе разделения переменных мы взяли
отрицательной.

\begin{verbatim}
Подставим теперь (21.3) в уравнение колебаний балки и получим уравнение для   X\left(x\right)
\end{verbatim}

MATHBLK0799X

общее решение которого записывается в виде

MATHBLK0800X

Для определения коэффициентов C, D, E и F подставим общее решение в
граничные условия

MATHBLK0801X

MATHBLK0802X

{[} \textbackslash begin\{split\}
u\textbackslash left(0,~t\textbackslash right) \&= 0
\textbackslash Rightarrow
X\textbackslash left(0\textbackslash right)T\textbackslash left(t\textbackslash right)
= 0 \textbackslash Rightarrow
X\textbackslash left(0\textbackslash right) = 0
\textbackslash Rightarrow C + E = 0 \textbackslash{}
u\_\{xx\}\textbackslash left(0,~t\textbackslash right) \&= 0
\textbackslash Rightarrow
X'\,`\textbackslash left(0\textbackslash right)T\textbackslash left(t\textbackslash right)
= 0 \textbackslash Rightarrow
X'\,'\textbackslash left(0\textbackslash right) = 0
\textbackslash Rightarrow -C + E = 0 \textbackslash{}
u\textbackslash left(1,~t\textbackslash right) \&= 0
\textbackslash Rightarrow
D\textbackslash sin\textbackslash sqrt\{\textbackslash omega\} +
F\textbackslash sin\textbackslash sqrt\{\textbackslash omega\} = 0,
\textbackslash{} u\_\{xx\}\textbackslash left(1,~t\textbackslash right)
\&= 0 \textbackslash Rightarrow
-D\textbackslash sin\textbackslash sqrt\{\textbackslash omega\} +
F\textbackslash sin\textbackslash sqrt\{\textbackslash omega\} = 0.
\textbackslash end\{split\} {]}

MATHBLK0803X

MATHBLK0804X

{[} F \textbackslash sinh \textbackslash sqrt\{\textbackslash omega\} =
0, {]}

MATHBLK0805X

{[} D \textbackslash sin \textbackslash sqrt\{\textbackslash omega\} =
0. {]}

MATHBLK0806X

{[} F = 0 {]}

MATHBLK0807X

{[} \textbackslash omega\_n = (n\textbackslash pi)\^{}2, {]}

MATHBLK0808X

{[} u\_n(x, t) = X\_n(x) T\_n(t) = {[}a\_n
\textbackslash sin(n\textbackslash pi)\^{}2 t + b\_n
\textbackslash cos(n\textbackslash pi)\^{}2 t{]}
\textbackslash sin(n\textbackslash pi x). {]}

MATHBLK0809X

{[} u(x, t) = \textbackslash sum\_\{n=1\}\^{}\{\textbackslash infty\}
\textbackslash left{[} a\_n \textbackslash sin(n\textbackslash pi)\^{}2
t + b\_n \textbackslash cos(n\textbackslash pi)\^{}2 t
\textbackslash right{]} \textbackslash sin(n\textbackslash pi x) {]}

MATHBLK0810X

{[} u(x, 0) = f(x) =
\textbackslash sum\_\{n=1\}\^{}\{\textbackslash infty\} b\_n
\textbackslash sin(n\textbackslash pi x), {]}

MATHBLK0811X

{[} u\_t(x, 0) = g(x) =
\textbackslash sum\_\{n=1\}\^{}\{\textbackslash infty\}
(n\textbackslash pi)\^{}2 a\_n \textbackslash sin(n\textbackslash pi x).
{]}

MATHBLK0812X

{[} a\_n = \textbackslash frac\{2\}\{(n\textbackslash pi)\^{}2\}
\textbackslash int\_0\^{}1 g(x) \textbackslash sin(n\textbackslash pi x)
dx, {]}

MATHBLK0813X

{[} b\_n = 2 \textbackslash int\_0\^{}1 f(x)
\textbackslash sin(n\textbackslash pi x) dx. {]}

MATHBLK0814X

{[} u(x, 0) = \textbackslash sin(\textbackslash pi x) + 0.5
\textbackslash sin(3\textbackslash pi x), {]}

MATHBLK0815X

{[} a\_n=0 {]}

MATHBLK0816X

{[} u(x, t) = \textbackslash cos(\textbackslash pi\^{}2 t)
\textbackslash sin(\textbackslash pi x) + 0.5
\textbackslash cos(9\textbackslash pi\^{}2 t)
\textbackslash sin(3\textbackslash pi x) {]}

MATHBLK0817X

{[} u(x, t) = \textbackslash cos(\textbackslash pi t)
\textbackslash sin(\textbackslash pi x) + 0.5
\textbackslash cos(3\textbackslash pi t)
\textbackslash sin(3\textbackslash pi x) {]}

MATHBLK0818X

{[} u(0, t) = 0, {]}

MATHBLK0819X

{[} u(x, t) = \textbackslash sum\_\{n=1\}\^{}\{\textbackslash infty\}
X\_n(x) \textbackslash left{[} a\_n
\textbackslash sin(\textbackslash omega\_n t) + b\_n
\textbackslash cos(\textbackslash omega\_n t) \textbackslash right{]},
{]}

MATHBLK0820X

{[} (\textbackslash text\{УЧП\}) \textbackslash quad u\_\{tt\} +
u\_\{xxxx\} = 0, \textbackslash quad 0 \textless{} x \textless{} 1,
\textbackslash quad 0 \textless{} t \textless{} \textbackslash infty,
{]}

MATHBLK0821X

{[} (\textbackslash text\{ГУ\}) \textbackslash quad {]}

MATHBLK0822X

{[} \textbackslash begin\{cases\} u(0, t) \&= 0, \textbackslash{}
u\_x(0, t) \&= 0, \textbackslash{} u\_\{xx\}(1, t) \&= 0,
\textbackslash{} u\_\{xxx\}(1, t) \&= 0, \textbackslash end\{cases\} {]}

MATHBLK0823X

MATHBLK0824X

{[} (\textbackslash text\{HV\}) \textbackslash quad {]}

MATHBLK0825X

{[} \textbackslash begin\{cases\} u(x, 0) \&= f(x), \textbackslash{}
u\_t(x, 0) \&= g(x). \textbackslash end\{cases\} {]}

MATHBLK0826X

MATHBLK0827X

6 M 601

УКАЗАНИЕ. Хотя собственные функции \(X_n\)(x) в этой задаче уже не
обычные синусы, согласно общей теории Штурма --- Ли вилля они будут
ортогональны на отрезке {[}0, 1{]}.

\begin{enumerate}
\def\labelenumi{\arabic{enumi}.}
\setcounter{enumi}{1}
\tightlist
\item
  Найдите решение для свободно опирающейся на оба конца балки при
  начальных условиях
\end{enumerate}

MATHBLK0828X

{[} u\_t(x, 0) = \textbackslash sin(\textbackslash pi x). 0
\textbackslash le x \textbackslash le 1, {]}

\begin{enumerate}
\def\labelenumi{\arabic{enumi}.}
\setcounter{enumi}{2}
\tightlist
\item
  Решите задачу 2 с начальными условиями
\end{enumerate}

MATHBLK0829X

MATHBLK0830X

{[} \textbackslash begin\{aligned\} u(x, 0) = 1 - x\^{}2,
\textbackslash{} u\_t(x, 0) = 0. \textbackslash end\{aligned\} {]}

MATHBLK0831X

{[} \textbackslash quad 0 \textbackslash le x \textbackslash le 1, {]}

MATHBLK0832X

MATHBLK0833X

(см. рис. 22.2).

\pandocbounded{\includegraphics[keepaspectratio]{images/_page_161_Figure_10.jpeg}}

РИС. 22.2. Задача теплопроводности в плоскости переменных (х, т).

Наша цель --- дать новую эквивалентную (22.1) постановку задачи таким
образом, чтобы

в новой задаче не было никаких физических параметров (вроде α);

\begin{enumerate}
\def\labelenumi{\arabic{enumi})}
\setcounter{enumi}{1}
\item
  начальные и граничные условия стали проще.

  Для осуществления этой цели, мы введем три новых безразмерных
  переменных U, \textbackslash xi и \textbackslash tau , чтобы заменить
  прежние u, x и t по схеме
\end{enumerate}

MATHBLK0834X

\begin{verbatim}
(безразмерная температура),   x \to \xi   (безразмерная длина),   t \to \tau   (безразмерное время).
\end{verbatim}

Мы проведем все три преобразования одновременно.

\begin{verbatim}
# Преобразование зависимой переменной  u \rightarrow U
\end{verbatim}

Определим функцию \(U(x, t)\) по формуле

MATHBLK0835X.

\begin{verbatim}
Ясно, что новая температура   U\left(x,\,t\right)   безразмерна, поскольку мы делим градусы Цельсия на градусы Цельсия. Новые граничные условия для   U\left(x,\,t\right)   при x=0 и x=L примут вид   U\left(0,\,t\right)=0   и   U\left(L,\,t\right)=1  . Давайте поставим задачу для новой функции   U\left(x,\,t\right)  . После несложных преобразований задача (22.1) переходит в новую задачу
\end{verbatim}

задачу

MATHBLK0836X

\begin{verbatim}
Если угодно, то здесь можно остановиться, решить задачу для   U\left(x,\, t\right)   и затем найти   u\left(x,\,t\right)   по формуле
\end{verbatim}

MATHBLK0837X

Однако мы пойдем дальше и займемся преобразованием независимых
переменных x и t.

\begin{verbatim}
# Преобразование пространственной координаты  x \rightarrow \xi

Представляется совершенно очевидным, как следует выбрать безразмерную переменную   \xi  . Поскольку   0 \leqslant x \leqslant L  , полагаем
\end{verbatim}

MATHBLK0838X.

После вычисления производных

MATHBLK0839X

MATHBLK0840X

\begin{verbatim}
становится очевидной постановка следующей задачи (в переменных U,   \xi  , t):
\end{verbatim}

MATHBLK0841X

\begin{verbatim}
Мы проделали уже две трети пути. Последний шаг связан с введением безразмерного времени   \tau   так, чтобы исключить коэффициент   [\alpha/L]^2   из дифференциального уравнения.

# Преобразование времени  t \rightarrow \tau

Как выбрать безразмерное время? Это не так очевидно, как в случае первых двух переменных. Однако, поскольку наша цель — исключить константу   [\alpha/L]^2   из уравнения (22.3), мы поступим следующим образом.

- 1. Попытаемся провести замену вида   \tau = ct  , где c неизвестная пока константа.

- 2. Вычислим   u_t = u_{\tau} \tau_t = c u_{\tau}  .
\end{verbatim}

\begin{itemize}
\tightlist
\item
  \begin{enumerate}
  \def\labelenumi{\arabic{enumi}.}
  \setcounter{enumi}{2}
  \tightlist
  \item
    Подставим эту производную в уравнение (22.3) и получим
  \end{enumerate}
\end{itemize}

MATHBLK0842X

\begin{verbatim}
Следовательно, нужно выбрать   c = [\alpha/L]^2  , а значит, новое время должно вводиться по формуле
\end{verbatim}

MATHBLK0843X

\begin{verbatim}
Применяя это преобразование к предыдущей задаче (22.3), можно полностью перейти к безразмерным переменным U,   \xi   и   \tau_i
\end{verbatim}

MATHBLK0844X

Новая безразмерная система обладает следующими свойствами:

\begin{enumerate}
\def\labelenumi{\arabic{enumi}.}
\item
  Уравнение не содержит параметров.
\item
  Граничные условия имеют простой вид.
\item
  Начальное условие изменилось несущественно (т. е. оно содержит только
  известные функции).

  \begin{enumerate}
  \def\labelenumii{\arabic{enumii}.}
  \setcounter{enumii}{3}
  \tightlist
  \item
    Задача стала проще и компактнее по сравнению с исходной. Решение
    этой задачи можно построить раз и навсегда, поэтому если читатель
    преобразует свою исходную задачу (22.1) к безразмерному виду (22.4)
    и найдет ее решение U(\textbackslash xi, \textbackslash tau) в
    учебнике или научном журнале, то он может найти решение исходной
    задачи (22.1) простым пересчетом
  \end{enumerate}
\end{enumerate}

MATHBLK0845X

Этим завершается наше обсуждение проблемы преобразования задач к
безразмерному виду. Не существует универсальных правил, по которым
следует вводить новые переменные, здесь следует опираться на физическую
интуицию и пробовать различные возможности.

Мы закончим эту лекцию простым примером преобразования задачи к
безразмерному виду, решения новой задачи и затем возвращения к исходным
переменным.

\section{Пример преобразования гиперболической задачи к безразмерному
виду}\label{ux43fux440ux438ux43cux435ux440-ux43fux440ux435ux43eux431ux440ux430ux437ux43eux432ux430ux43dux438ux44f-ux433ux438ux43fux435ux440ux431ux43eux43bux438ux447ux435ux441ux43aux43eux439-ux437ux430ux434ux430ux447ux438-ux43a-ux431ux435ux437ux440ux430ux437ux43cux435ux440ux43dux43eux43cux443-ux432ux438ux434ux443}

Рассмотрим колеблющуюся струну

(22.5)

MATHBLK0846X

MATHBLK0847X

С помощью преобразования независимых переменных к безразмерному виду
(преобразовывать u нет необходимости)

\begin{verbatim}
MATHBLK0848X,   \tau = [\alpha/L] t
\end{verbatim}

получаем новую задачу

MATHBLK0849X

MATHBLK0850X

MATHBLK0851X

{[} \textbackslash begin\{cases\} u(0, \textbackslash tau) = 0,
\textbackslash{} u(1, \textbackslash tau) = 0, \textbackslash{}
u(\textbackslash xi, 0) =
\textbackslash sin(\textbackslash pi\textbackslash xi) +
0.5\textbackslash sin(3\textbackslash pi\textbackslash xi),
\textbackslash{} u\_\{\textbackslash xi\}(\textbackslash xi, 0) = 0,
\textbackslash end\{cases\} {]}

MATHBLK0852X

{[} \textbackslash quad 0 \textless{} \textbackslash xi
\textbackslash leqslant 1, {]}

MATHBLK0853X

{[} (\textbackslash text\{HY\}) \textbackslash quad {]}

MATHBLK0854X

{[} \textbackslash begin\{cases\} u(\textbackslash xi, 0) =
\textbackslash sin(\textbackslash pi\textbackslash xi) +
0.5\textbackslash sin(3\textbackslash pi\textbackslash xi),
\textbackslash{} u\_\{\textbackslash xi\}(\textbackslash xi, 0) = 0,
\textbackslash end\{cases\} {]}

MATHBLK0855X

MATHBLK0856X

решением которой является функция

MATHBLK0857X

Если вернуться обратно к координатам x и t, то решение задачи (22.5)
принимает вид

MATHBLK0858X

\section{ЗАМЕЧАНИЯ}\label{ux437ux430ux43cux435ux447ux430ux43dux438ux44f-6}

\begin{itemize}
\tightlist
\item
  \begin{enumerate}
  \def\labelenumi{\arabic{enumi}.}
  \tightlist
  \item
    Анализ размерности особенно важен в численных методах, поскольку
    большинство программ написаны для решения общих математических задач
    и не подразумевают использование большого числа физических
    параметров. Желающие воспользоваться этими программами должны
    преобразовать свою задачу к соответствующему виду, решить задачу на
    ЭВМ, а затем привести численные результаты к нужной размерности.
  \end{enumerate}
\item
  \begin{enumerate}
  \def\labelenumi{\arabic{enumi}.}
  \setcounter{enumi}{1}
  \tightlist
  \item
    Анализ размерностей позволяет математикам работать с уравнением в
    частных производных, исключив из него множество параметров и
    констант, которые не затрагивают математического существа задачи.
  \end{enumerate}
\item
  \begin{enumerate}
  \def\labelenumi{\arabic{enumi}.}
  \setcounter{enumi}{2}
  \tightlist
  \item
    Не всегда обязательно преобразовывать все переменные к безразмерному
    виду, иногда достаточно преобразовать одну или две переменных.
  \end{enumerate}
\end{itemize}

\section{ЗАДАЧИ}\label{ux437ux430ux434ux430ux447ux438-6}

\begin{enumerate}
\def\labelenumi{\arabic{enumi}.}
\item
  С помощью замены переменных

  MATHBLK0859X, \textbackslash tau = {[}\textbackslash alpha/L{]} t
\end{enumerate}

преобразуйте задачу (22.5) к безразмерному виду (22.6).

\begin{enumerate}
\def\labelenumi{\arabic{enumi}.}
\setcounter{enumi}{1}
\tightlist
\item
  Найдите безразмерную форму для задачи
\end{enumerate}

MATHBLK0860X

MATHBLK0861X

\begin{verbatim}
\begin{aligned} (\text{Y} \, \Pi) & u_t = \alpha^2 u_{xx}, & 0 < x < L, \\ (\Gamma \, \text{Y}) & \left\{ \begin{aligned} u \, (0, \, t) = T_1, \\ u \, (L, \, t) = 0, \\ u \, (x, \, 0) = T_2, & 0 \leqslant x \leqslant L. \end{aligned}
\end{verbatim}

MATHBLK0862X

MATHBLK0863X

\begin{enumerate}
\def\labelenumi{\arabic{enumi}.}
\setcounter{enumi}{2}
\tightlist
\item
  С помощью замены
\end{enumerate}

MATHBLK0864X

преобразуйте задачу (22.1) в задачу (22.2).

\begin{verbatim}
4. Почему замена переменной   \tau = \alpha t   приводит к исключению параметра   \alpha   из волнового уравнения   u_{tt} = \alpha^2 u_{xx}  ? В чем физический смысл этой замены? Напоминаем вам, что   \alpha  —скорость
\end{verbatim}

распространения волн. При поиске новых координат в большинстве случаев
главную роль играет интуиция.

Как выбрать новую пространственную координату ξ, чтобы исключить и из
уравнения

MATHBLK0865X

Лекция 23

\section{КЛАССИФИКАЦИЯ УРАВНЕНИЙ С ЧАСТНЫМИ ПРОИЗВОДНЫМИ (КАНОНИЧЕСКАЯ
ФОРМА ГИПЕРБОЛИЧЕСКОГО
УРАВНЕНИЯ)}\label{ux43aux43bux430ux441ux441ux438ux444ux438ux43aux430ux446ux438ux44f-ux443ux440ux430ux432ux43dux435ux43dux438ux439-ux441-ux447ux430ux441ux442ux43dux44bux43cux438-ux43fux440ux43eux438ux437ux432ux43eux434ux43dux44bux43cux438-ux43aux430ux43dux43eux43dux438ux447ux435ux441ux43aux430ux44f-ux444ux43eux440ux43cux430-ux433ux438ux43fux435ux440ux431ux43eux43bux438ux447ux435ux441ux43aux43eux433ux43e-ux443ux440ux430ux432ux43dux435ux43dux438ux44f}

ЦЕЛЬ ЛЕКЦИИ: Показать, что линейное уравнение второго порядка с двумя
независимыми переменными

MATHBLK0866X

(A, B, C, D, E, F и G являются функциями x и y или могут быть
константами) относится к одному из следующих типов:

\begin{itemize}
\item
  \begin{enumerate}
  \def\labelenumi{\arabic{enumi})}
  \tightlist
  \item
    гиперболическому (если B\^{}2 4AC \textgreater{} 0 ),
  \end{enumerate}
\item
  \begin{enumerate}
  \def\labelenumi{\arabic{enumi})}
  \setcounter{enumi}{1}
  \tightlist
  \item
    параболическому (если B\^{}2 4AC = ),
  \end{enumerate}
\item
  \begin{enumerate}
  \def\labelenumi{\arabic{enumi})}
  \setcounter{enumi}{2}
  \tightlist
  \item
    эллиптическому (если 0 B\^{}2 4AC \textless{} 0 ).
  \end{enumerate}

  Показать, какие новые переменные \textbackslash xi =
  \textbackslash xi(x, y) и \textbackslash eta = \textbackslash eta(x,
  y) можно ввести вместо x и y, чтобы упростить исходное уравнение. В
  новых переменных \textbackslash xi и \textbackslash eta уравнение с
  частными производными приводится к одному из следующих видов (в
  вависимости от того, будет ли величина B\^{}2 - 4AC положительна,
  равна нулю или отрицательна):

  \begin{itemize}
  \tightlist
  \item
    \begin{enumerate}
    \def\labelenumi{\arabic{enumi}.}
    \tightlist
    \item
      \textbackslash end\{aligned\}
    \end{enumerate}
  \end{itemize}
\end{itemize}

{[} \textbackslash begin\{cases\}
u\_\{\textbackslash xi\textbackslash xi\}
u\_\{\textbackslash eta\textbackslash eta\} =
\textbackslash Psi(\textbackslash xi, \textbackslash eta, u,
u\_\{\textbackslash xi\}, u\_\{\textbackslash eta\}) і
\textbackslash text\{(дві канонічні форми гіперболічного рівняння),\}
\textbackslash end\{cases\} {]}

\begin{itemize}
\item
  \begin{enumerate}
  \def\labelenumi{\arabic{enumi}.}
  \setcounter{enumi}{1}
  \tightlist
  \item
    \(u_{\eta\eta} = \Phi(\xi, \eta, u, u_{\xi}, u_{\eta}) (\)каноническая
    форма параболического уравнения),
  \end{enumerate}

  \begin{itemize}
  \tightlist
  \item
    \begin{enumerate}
    \def\labelenumi{\arabic{enumi}.}
    \setcounter{enumi}{2}
    \tightlist
    \item
      u\_\{\textbackslash xi\textbackslash xi\} +
      u\_\{\textbackslash eta\textbackslash eta\} =
      \textbackslash Phi(\textbackslash xi, \textbackslash eta, u,
      u\_\{\textbackslash xi\}, u\_\{\textbackslash eta\}) (каноническая
      форма эллиптического уравнения),
    \end{enumerate}
  \end{itemize}

  где \textbackslash Phi и \textbackslash Psi ---функции от частных
  производных первого порядка, зависимой переменной u и новых
  независимых переменных \textbackslash xi и \textbackslash eta .
  Конкретный вид функций \textbackslash Phi и \textbackslash Psi зависит
  от исходного уравнения.
\end{itemize}

У читателя может возникнуть мысль, что глава, посвященная классификации
уравнений, должна размещаться в самом начале книги. Вероятно, это
правильно, и многие книги действительно начинаются с этого материала.
Однако правда и то, что многие студенты не могут с энтузиазмом изучать
нечто такое, о чем у них нет никакого представления. Именно по этой
причине мы отложили изучение темы, посвященной классификации уравнений с
частными производными, до настоящего момента.

Наша задача - классифицировать все уравнения вида

(23 1)

{[} u\_\{xx\} + x\^{}2 u\_\{yy\} = 0 {]}

где A, B, C, D, E, F и G---вообще говоря, функции {[}x и y. Согласно
предлагаемой схеме классификации уравнение называется

\begin{enumerate}
\def\labelenumi{\arabic{enumi})}
\item
  гиперболическим в точке (\(x_0\), \(y_0\)) , если B\^{}2(\(x_0\),
  \(y_0\)) - 4A(\(x_0\), \(y_0\)) \textbackslash помножити на
  \textbackslash помножити на C(\(x_0\), \(y_0\)) \textgreater{} 0 2)
  параболнческим в точке (\(x_0\), \(y_0\)) , если B\^{}2(\(x_0\),
  \(y_0\)) - 4A(\(x_0\), \(y_0\)) \textbackslash помножити на C(\(x_0\),
  \(y_0\)) = 0 ,
\item
  эллиптическим в точке (\(x_0\), \(y_0\)) , если B\^{}2(\(x_0\),
  \(y_0\)) - 4A(\(x_0\), \(y_0\)) \textbackslash помножити на
  \textbackslash помножити на C(\(x_0\), \(y_0\)) \textless{} 0 и в
  зависимости от типа его можно привести к одной из канонических форм.
  Для того чтобы было легче разобраться в классификационной схеме, мы
  приведем четыре примера гиперболических, параболических и
  эллиптических уравнений.
\end{enumerate}

\section{Примеры гиперболических, параболических и эллиптических
уравнений
10}\label{ux43fux440ux438ux43cux435ux440ux44b-ux433ux438ux43fux435ux440ux431ux43eux43bux438ux447ux435ux441ux43aux438ux445-ux43fux430ux440ux430ux431ux43eux43bux438ux447ux435ux441ux43aux438ux445-ux438-ux44dux43bux43bux438ux43fux442ux438ux447ux435ux441ux43aux438ux445-ux443ux440ux430ux432ux43dux435ux43dux438ux439-10}

\begin{enumerate}
\def\labelenumi{\arabic{enumi}.}
\tightlist
\item
  У равнение теплопроводности \(u_t = u_{xx}\)--- линейное уравнение
  второго порядка вида (23.1) с коэффициентами
\end{enumerate}

{[} (\textbackslash text\{UBP\}) \textbackslash qquad {]}

{[} , B=0 , C=0 , D=0 , E=-1 , F=0 , G=0 , {]}

так что B\^{}2-4AC=0 при всех x и t. Следовательно, это уравнение будет
параболическим при всех x и t. В общем уравнении (23 1) переменная y
играет роль времени t. Результат не изменится, если в общем уравнении
поменять местами x и y.

\begin{enumerate}
\def\labelenumi{\arabic{enumi}.}
\setcounter{enumi}{1}
\tightlist
\item
  Волновое уравнение \(u_{1t} = u_{xx}\)также принадлежит классу
  уравнений вида (23.1) с коэффициентами
\end{enumerate}

, \(B=0\) , \(C=-1\) , \(D=E=F=G=0\) .

Так как \(B^2 - 4AC = 4\) для всех x и t, то это уравнение является
еиперболическим при всех x и t.

1) Следует указать, что приведенные в примерах уравнения являются
частными случаями более общих параболических, гиперболических в
эллиптических уравнений, а основные их свойства характерны и для общего
случая.

\begin{enumerate}
\def\labelenumi{\arabic{enumi}.}
\setcounter{enumi}{2}
\item
  У равнение Лапласа \(u_{xx} + u_{yy} = 0\) является эллиптическим при
  всех x и y, так как \(B^2 - 4AC = - 4 < 0\) .
\item
  Линейное уравнение \(xu_{xx} + u_{yy} = \sin x\) с переменными
  коэффициентами также принадлежит классу уравнений вида (23.1), но для
  него \(B^2-4AC=-4x\) и, следовательно, уравнение будет эллиптическим
  при x \textgreater{} 0,
\end{enumerate}

лараболическим при x=0, гиперболическим при x \textless{} 0.

Этот пример иллюстрирует тот факт, что тип уравнения о переменными
коэффициентами может изменяться от точки к точке.

Читатель должен обратить внимание на следующее: тип уравнения (23.1)
определяется только коэффициентами при вторых првизводных и никак не
зависит ни от коэффициентов при первых производных и самой функции, ни
от свободного члена.

Обратимся теперь к главной теме нашей лекции --- приведению
гиперболических уравнений к каноническому виду. Известно, что если в
данной области пространства уравнение является гиперболическим, то
вместо x и y можно ввести координаты \(\xi\) и \(\eta\)
(характеристические переменные) таким образом, что уравнение примет
простейшую форму

(23.2)

Это уравнение содержит только одну производную второго порядка
\(u_{\xi\eta}\) , а функция \textbackslash Phi(\textbackslash xi,)
\textbackslash eta,u,\(u_{\xi}\),\(u_{\eta}\)) зависит только от новых
независимых переменных § и η, зависимой переменной и и первых
производных из и ип. Конкретный вид функции Ф зависит, конечно, от
исходного уравнения и формул перехода к новым координатам § и п.
Научимся отыскивать эти функции.

\section{Каноническая форма гиперболического
уравнения}\label{ux43aux430ux43dux43eux43dux438ux447ux435ux441ux43aux430ux44f-ux444ux43eux440ux43cux430-ux433ux438ux43fux435ux440ux431ux43eux43bux438ux447ux435ux441ux43aux43eux433ux43e-ux443ux440ux430ux432ux43dux435ux43dux438ux44f}

Рассмотрим уравнение общего вида

(23.3)

\begin{enumerate}
\def\labelenumi{\arabic{enumi}.}
\setcounter{enumi}{4}
\tightlist
\item
  Разработайте схему метода Монте-Карло для розв'язки смешанной
  параболической задачи
\end{enumerate}

и пусть в интересующей нас области B\^{}2 - 4AC \textgreater{} 0 . Нам
нужно ввести новые координаты 1)

\begin{verbatim}
\begin{array}{ll} (\text{УЧП}) & u_t = \alpha^2 u_{xx} - \beta u_x - \gamma u, \quad 0 < x < 1, \quad 0 < t < \infty \\ (\text{ГУ}) & \left\{ \begin{array}{ll} u\left(0,\ t\right) = f\left(t\right), & 0 < t < \infty, \\ u\left(1,\ t\right) = g\left(t\right), & 0 < t < \infty, \\ u\left(x,\ 0\right) = \varphi\left(x\right), & 0 \leqslant x \leqslant 1. \end{array}

\right. \end{aligned}
\end{verbatim}

так, чтобы в уравнении (23.3) осталась только одна частная производная
второго порядка \(u_{\xi\eta}\) (не следует пытаться искать такую замену
переменных, с помощью которой гиперболическое уравне-

1) Необходимо, чтобы это преобразование координат было локально обратимо
и дважды дифференцируемо. --- Прим. ред.

ние можно было бы привести и параболической или заливитической
канонической форме, это невозможно сделать).

Вычислим сначала частные производные

{[} J{[}y{]} = \textbackslash int\_\{b\}\^{}\{a\} F(x, y, y') dx, {]}

{[} J{[}y{]} = \textbackslash int\_\{0\}\^{}\{1\} {[}y\^{}\{2\}(x) +
y'\^{}\{2\}(x){]} dx. {]}

{[} \textbackslash begin\{aligned\} u\_x = u\_\textbackslash xi
\textbackslash xi\_x + u\_\textbackslash eta \textbackslash eta\_x,
\textbackslash{} u\_y = u\_\textbackslash xi \textbackslash xi\_y +
u\_\textbackslash eta \textbackslash eta\_y, \textbackslash{} (23.4)
~u\_\{xx\} = u\_\{\textbackslash xi\textbackslash xi\}
\textbackslash xi\_x\^{}2 + 2u\_\{\textbackslash xi\textbackslash eta\}
\textbackslash xi\_x \textbackslash eta\_x +
u\_\{\textbackslash eta\textbackslash eta\} \textbackslash eta\_x\^{}2 +
u\_\{\textbackslash xi\} \textbackslash xi\_\{xx\} +
u\_\{\textbackslash eta\} \textbackslash eta\_\{xx\}, \textbackslash{}
u\_\{xy\} = u\_\{\textbackslash xi\textbackslash xi\}
\textbackslash xi\_x\^{}2 + u\_\{\textbackslash xi\textbackslash eta\}
(\textbackslash xi\_x \textbackslash eta\_y + \textbackslash xi\_y
\textbackslash eta\_x) + u\_\{\textbackslash eta\textbackslash eta\}
\textbackslash eta\_x \textbackslash eta\_y + u\_\{\textbackslash xi\}
\textbackslash xi\_\{xy\} + u\_\{\textbackslash eta\}
\textbackslash eta\_\{xy\}, \textbackslash{} u\_\{yy\} =
u\_\{\textbackslash xi\textbackslash xi\} \textbackslash xi\_y\^{}2 +
2u\_\{\textbackslash xi\textbackslash eta\} \textbackslash xi\_y
\textbackslash eta\_y + u\_\{\textbackslash eta\textbackslash eta\}
\textbackslash eta\_y\^{}2 + u\_\{\textbackslash xi\}
\textbackslash xi\_\{yy\} + u\_\{\textbackslash eta\}
\textbackslash eta\_\{yy\}. \textbackslash end\{aligned\} {]}

{[} \textbackslash frac\{df(x)\}\{dx\} = 0 {]}

{[} T = \textbackslash int\_\{0\}\^{}\{T\} dt =
\textbackslash int\_\{a\}\^{}\{L\} \textbackslash frac\{dt\}\{ds\} ds =
\textbackslash int\_\{0\}\^{}\{L\} \textbackslash frac\{ds\}\{v\} =
\textbackslash frac\{1\}\{\textbackslash sqrt\{2mg\}\}
\textbackslash int\_\{0\}\^{}\{L\}
\textbackslash frac\{ds\}\{\textbackslash sqrt\{y\}\} =
\textbackslash frac\{1\}\{\textbackslash sqrt\{2mg\}\}
\textbackslash int\_\{a\}\^{}\{b\}
\textbackslash sqrt\{\textbackslash frac\{1+y'\^{}\{2\}\}\{y\}\} dy {]}

Подставляя эти соотношения в исходное уравнение (23.3), после простых,
но громоздких вычислений получаем

(23.5)

{[} J{[}y{]} = \textbackslash int\_\{a\}\^{}\{b\} F(x, y, y') dx. {]}

{[} u(r, \textbackslash theta) = u\_0(r, \textbackslash theta) +
\textbackslash varepsilon u\_1(r, \textbackslash theta) +
\textbackslash varepsilon\^{}2 u\_2(r, \textbackslash theta) +
\textbackslash dots {]}

{[} \textbackslash Delta u + \textbackslash varepsilon u\^{}2 = 0 {]}

{[} u\_0 (1, \textbackslash theta) + \textbackslash varepsilon u\_1 (1,
\textbackslash theta) + \textbackslash varepsilon\^{}2 u\_2 (1,
\textbackslash theta) + \textbackslash ldots = \textbackslash cos
\textbackslash theta. {]}

{[} P\_\{0\} {]}

{[} \textbackslash qquad u\_\{\textbackslash theta\}(r,
\textbackslash theta) = r \textbackslash cos \textbackslash theta, {]}

Читателю предоставится возможность проделать эти выкладки при решении
задач.

\begin{verbatim}
Наш следующий шаг—выбрать функции   \xi = \xi(x, y)   и   \eta = \eta(x, y)   таким образом, чтобы коэффициенты   \overline{A}   и   \overline{C}   обратились в нуль. Это позволит нам привести исходное уравиение к канонической форме. Итак, необходимо потребовать, чтобы
\end{verbatim}

\(\bar{C} = A\eta_x^2 + B\eta_x\eta_y + C\eta_y^2 = 0.\) Эти уравнения
можно представить в виде

{[} \textbackslash begin\{split\}
\&A\textbackslash left{[}\textbackslash xi\_\{x\}/\textbackslash xi\_\{y\}\textbackslash right{]}\^{}\{2\}+B\textbackslash left{[}\textbackslash xi\_\{x\}/\textbackslash xi\_\{y\}\textbackslash right{]}+C=0,\textbackslash{}
\&A\textbackslash left{[}\textbackslash eta\_\{x\}/\textbackslash eta\_\{y\}\textbackslash right{]}\textsuperscript{\{2\}+B\textbackslash left{[}\textbackslash eta\_\{x\}/\textbackslash eta\_\{y\}\textbackslash right{]}}\{2\}+C=0.
\textbackslash end\{split\} {]}

{[} P\_\{2\} {]}

{[} \textbackslash begin\{bmatrix\}
\textbackslash xi\_x/\textbackslash xi\_y \textbackslash end\{bmatrix\}
{]}

{[} P\_\{1\} {]}

(характеристические

{[} \textbackslash begin\{cases\} \textbackslash Delta u\_\{1\} =
-r\^{}\{2\} \textbackslash cos\^{}\{2\} \textbackslash theta =
-\textbackslash frac\{r\^{}\{2\}\}\{2\} \textbackslash left{[}1 +
\textbackslash cos(2\textbackslash theta)\textbackslash right{]} =
-\textbackslash frac\{r\^{}\{2\}\}\{2\} -
\textbackslash frac\{r\^{}\{2\}\}\{2\}
\textbackslash cos(2\textbackslash theta), \textbackslash{} u\_\{1\}(1,
\textbackslash theta) = 0. \textbackslash end\{cases\} {]}

уравнения).

\begin{verbatim}
Kaждое из квадратных уравнений для   [\xi_x/\xi_y]   и   [\eta_x/\eta_y]   имеет два кория, но мы оставляем только по одному, но так чтобы они были различны.
\end{verbatim}

\pandocbounded{\includegraphics[keepaspectratio]{images/_page_170_Figure_3.jpeg}}

\begin{verbatim}
РИС. 23.1. Характеристики   \xi(x, y) = c   и   \eta(x, y) = c  .

Задача свелась к нахождению двух функций   \xi(x, y)   и   \eta(x, y)  , таких, чтобы отношения   [\xi_x/\xi_y]   и   [\eta_x/\eta_y]   удовлетворяли уравнению (23.7). Найти такие функции очень просто 1), если обратить внимание на рис. 23.1. Чтобы понять, каким образом можно это сделать, рассмотрим простое уравнение
\end{verbatim}

{[} u\_h(r, \textbackslash theta) =
\textbackslash sum\_\{n=0\}\^{}\{\textbackslash infty\} r\^{}n {[}a\_n
\textbackslash cos(n\textbackslash theta) + b\_n
\textbackslash sin(n\textbackslash theta){]}. {]}

Его характеристики определяются из уравнений

{[} \textbackslash Delta u = -\textbackslash frac\{r\^{}2\}\{2\} -
\textbackslash frac\{r\^{}3\}\{2\}
\textbackslash cos\{(2\textbackslash theta)\} {]}

Решая эти уравнения относительно у, получаем

{[} u\_p(r, \textbackslash theta) = Ar\^{}4 + Br\^{}4
\textbackslash cos(2\textbackslash theta). {]}

{[} u\_p(r, \textbackslash theta) = -\textbackslash frac\{r\^{}4\}\{32\}
- \textbackslash frac\{r\^{}4\}\{24\}
\textbackslash cos(2\textbackslash theta). {]}

Чтобы найти \(\xi\) и \(\eta\) , разрешим полученные уравнения
относительно констант \(c_1\) и \(c_2\) , оставляя их в правой части и
перенося остальные слагаемые в левую.

Выражения, стоящие в левых частях, примем за искомые функции § и пр

{[} \textbackslash sum\_\{n=0\}\^{}\{\textbackslash infty\} {[}a\_n
\textbackslash cos(n\textbackslash theta) + b\_n
\textbackslash sin(n\textbackslash theta){]} -
\textbackslash frac\{1\}\{32\} - \textbackslash frac\{1\}\{24\}
\textbackslash cos(2\textbackslash theta) = 0. {]}

{[} u\_1(r, \textbackslash theta) = \textbackslash frac\{1\}\{32\} +
\textbackslash frac\{1\}\{24\}r\^{}2\textbackslash cos(2\textbackslash theta)
-
\textbackslash frac\{r\^{}4\}\{24\}\textbackslash cos(2\textbackslash theta)
- \textbackslash frac\{r\^{}4\}\{32\} = \textbackslash frac\{(r\^{}4 -
1)\}\{32\} - \textbackslash frac\{(r\^{}4 -
r\^{}2)\}\{24\}\textbackslash cos(2\textbackslash theta). {]}

1) См., например, Тихонов А. Н., Самарский А. А. Уравнения
математической физики. --- М.: Наука, 1972 и др. издания. --- Прим. ред.

\begin{verbatim}
Очевидно, введенные таким образом функции   \xi   и   \eta   удовлетворяют уравнениям характеристик. Эти новые координаты изображены на рис. 23.2. На этом мы заканчиваем обсуждение вопроса о том, как найти новые координаты. Осталось последнее — определить новый вид исходного уравнения.
\end{verbatim}

\pandocbounded{\includegraphics[keepaspectratio]{images/_page_171_Figure_3.jpeg}}

РИС. 23.2. Характеристические координаты для уравнения \(u_{xx}\) -
4u\_\{uu\} + \(u_x = 0\).

\begin{verbatim}
Это совсем просто: чтобы найти каноническую форму уравнения с частными производными, подставим новые координаты   \xi(x, y)   и   \eta(x,) y)   в уравнение
\end{verbatim}

{[} u = u\_0 + u\_1 = r \textbackslash cos \textbackslash theta -
\textbackslash frac\{(r\^{}4 - 1)\}\{32\} - \textbackslash frac\{(r\^{}4
- r\^{}2)\}\{24\} \textbackslash cos (2\textbackslash theta). {]}

,

\begin{verbatim}
где   \overline{A}  ,   \overline{B}  ,   \overline{C}  ,   \overline{D}  ,   \overline{E}  ,   \overline{F}   и   \overline{G}   определены формулами (23.6).
\end{verbatim}

Прежде чем закончить лекцию, давайте на конкретном примере посмотрим,
как «работает» общий метод приведения к каноническому виду.

Приведение гиперболического уравнения
y\textsuperscript{2u\_\{xx\}-x}2u\_\{yy\}=0 к канонической форме

Рассмотрим уравнение

{[} P\_\{2\} {]}

{[} , x\textgreater{} , 0 y\textgreater0 , {]}

которое является гиперболическим в первом квадранте. Найдем новые
координаты таким образом, чтобы привести исходное уравнение к
канонической форме.

ШАГ 1. Решаем уравнения характеристик

(напомним, что этот шаг

{[} u \textbackslash left( 1 + \textbackslash frac\{1\}\{4\}
\textbackslash sin \textbackslash theta, ~\textbackslash theta
\textbackslash right) = \textbackslash cos \textbackslash theta,
\textbackslash quad 0 \textbackslash le \textbackslash theta
\textbackslash le 2\textbackslash pi. {]}

обеспечивает \(\overline{A} = \overline{C} = 0\) ),

\pandocbounded{\includegraphics[keepaspectratio]{images/_page_172_Figure_2.jpeg}}

РИС, 23.3. Новые карактеристические координаты.

Проинтегрировав эти обыкновенные дифференциальные уравнения методом
разделения переменных, получаем два неявных соотнощения (при желании
можно было бы явно выразить у через х)

\pandocbounded{\includegraphics[keepaspectratio]{images/_page_349_Figure_5.jpeg}}

Следовательно, новые координаты следует вводить по формулам

,\\
\(\eta = y^2 + x^2\) .

Эти новые координаты изображены на рис. 23.3. Теперь следует получить
новое уравнение. Вычислим все коэффициенты

\(\widetilde{A} = 0\) (так и должно быть; мы специально выбирали
\(\xi \xi\) и \(\eta\) так, чтобы этот коэффициент обращался в нуль).

{[}
u\textbackslash left(1+\textbackslash frac\{1\}\{4\}\textbackslash sin\textbackslash theta,~\textbackslash theta\textbackslash right)=\textbackslash cos\textbackslash theta
{]}

,

\(\overline{C} = 0\) (по той же причине, что и \(\overline{A} = 0\) ),

{[} u(1+\textbackslash epsilon\textbackslash sin\textbackslash theta,
\textbackslash theta) = \textbackslash cos\textbackslash theta,
\textbackslash quad 0 \textbackslash le \textbackslash epsilon
\textbackslash le 1/4. {]}

{[} f(x+h) = f(x) + f'(x)h + f'\,'(x)\textbackslash frac\{h\^{}2\}\{2!\}
+ \textbackslash dots, {]}

,

\pandocbounded{\includegraphics[keepaspectratio]{images/_page_350_Figure_2.jpeg}}

,

подставим их в уравнение

и получим

{[}
u(1+\textbackslash varepsilon\textbackslash sin\textbackslash theta,~\textbackslash theta)=u(1,~\textbackslash theta)+u\_r(1,~\textbackslash theta)(\textbackslash varepsilon\textbackslash sin\textbackslash theta)+u\_\{rr\}(1,~\textbackslash theta)\textbackslash frac\{(\textbackslash varepsilon\textbackslash sin\textbackslash theta)\^{}2\}\{21\}+\textbackslash ldots
{]}

ШАГ 2. Выражая, наконец, x и y через \(\xi\) и \(\eta\) , окончательно
получаем

(23.8)

{[} \textbackslash Delta u = 0, \textbackslash quad 0 \textless{} r
\textless{} 1 + \textbackslash varepsilon \textbackslash sin
\textbackslash theta, {]}

Задача решена. Теперь мы знаем, что нужно сделать, чтобы найти новые
координаты и привести уравнение к канонической форме.

\section{ЗАМЕЧАНИЯ}\label{ux437ux430ux43cux435ux447ux430ux43dux438ux44f-7}

\begin{enumerate}
\def\labelenumi{\arabic{enumi}.}
\tightlist
\item
  На самом деле для гиперболических уравнений существует две
  канонические формы. Вторую можно получить из первов заменой переменных
  вида
\end{enumerate}

\textbackslash beta = \textbackslash beta (\textbackslash xi,
\textbackslash eta) = \textbackslash xi - \textbackslash eta.
Воспользуемся этой заменой для преобразования уравнения (23.8). Так как

{[} (46.8) u = u\_0 + \textbackslash varepsilon u\_1 +
\textbackslash varepsilon\^{}2 u\_2 + \textbackslash dots. {]}

{[} P\_\{0\} {]}

{[} \textbackslash begin\{array\}\{l\} u\_\{\textbackslash xi\} =
u\_\{\textbackslash alpha\}\textbackslash alpha\_\{\textbackslash xi\} +
u\_\{\textbackslash beta\}\textbackslash beta\_\{\textbackslash xi\} =
u\_\{\textbackslash alpha\} + u\_\{\textbackslash beta\}.
\textbackslash{} u\_\{\textbackslash eta\} =
u\_\{\textbackslash alpha\}\textbackslash alpha\_\{\textbackslash eta\}
+ u\_\{\textbackslash beta\}\textbackslash beta\_\{\textbackslash eta\}
= u\_\{\textbackslash alpha\} - u\_\{\textbackslash beta\},
\textbackslash{} u\_\{\textbackslash xi\textbackslash eta\} =
u\_\{\textbackslash alpha\textbackslash alpha\}\textbackslash alpha\_\{\textbackslash eta\}
+
u\_\{\textbackslash alpha\textbackslash beta\}\textbackslash beta\_\{\textbackslash eta\}
+
u\_\{\textbackslash beta\textbackslash alpha\}\textbackslash alpha\_\{\textbackslash eta\}
+
u\_\{\textbackslash beta\textbackslash beta\}\textbackslash beta\_\{\textbackslash eta\}
= u\_\{\textbackslash alpha\textbackslash alpha\} -
u\_\{\textbackslash beta\textbackslash beta\},
\textbackslash end\{array\} {]}

{[} \textbackslash begin\{cases\} \textbackslash Delta u\_\{0\} = 0, \&
0 \textless{} r \textless{} 1 \textbackslash text\{ (внутри круга),\}
\textbackslash{} u\_\{0\}(1, \textbackslash theta) = \textbackslash cos
\textbackslash theta, \& u\_\{0\}(r, \textbackslash theta) = r
\textbackslash cos \textbackslash theta, \textbackslash end\{cases\} {]}

{[} P\_\{1\} {]}

TO

(23 9)

{[} \textbackslash begin\{cases\} \textbackslash Delta u\_\{1\} = 0, \&
0 \textless{} r \textless{} 1 \textbackslash text\{ (внутри круга),\}
\textbackslash{} u\_\{1\}(1, \textbackslash theta) = -\textbackslash sin
\textbackslash theta \textbackslash frac\{\textbackslash partial
u\_\{0\}(1, \textbackslash theta)\}\{\textbackslash partial r\} =
-\textbackslash sin \textbackslash theta \textbackslash cos
\textbackslash theta \textbackslash end\{cases\} {]}

При желании можно выразить \textbackslash alpha и \textbackslash beta
через исходные координаты x и y. В результате получаем

{[} u = u\_0 + \textbackslash frac\{1\}\{4\}u\_1 +
\textbackslash left(\textbackslash frac\{1\}\{4\}\textbackslash right)\^{}2
u\_1 + \textbackslash cdots {]}

\begin{itemize}
\item
  Читатель может спросить, зачем нужно классифицировать и приводить к
  канонической форме уравнения с частными производными.
\item
  а) Разбиение уравнений на классы гиперболических, параболических и
  эллиптических соответствует разбиению физических процессов на три
  основных класса: волновых, диффузионных и стационарных. Математические
  особенности решений уравнений этих трех типов совершенно различны.
\item
  б) В большинстве работ, посвященных решению гиперболических задач,
  предполагается, что исходное уравнение записано в канонической форме,
  т. е. в виде
\end{itemize}

{[} u\_u + \textbackslash frac\{1\}\{4\} u\_t {]}

Если у нас есть уравнение и мы хотим получить его решение, мы должны
преобразовать его к канонической форме и воспользоваться известными
результатами.

в) Для численного решения канонического гиперболического уравнения
создано много различных программ для ЭВМ.

\begin{verbatim}
Функция   \Phi   (  \xi  ,   \xi \eta  , u,   u_{\xi}  ,   u_{\eta}  ) вводится в   \Theta  ВМ в виде подпрограммы, так что сначала необходимо привести уравнение к канонической форме. После того как решение в новых координатах найдено, его всегда можно преобразовать обратно к исходным координатам.
\end{verbatim}

\section{ЗАДАЧИ}\label{ux437ux430ux434ux430ux447ux438-7}

\begin{itemize}
\item
  \begin{enumerate}
  \def\labelenumi{\arabic{enumi}.}
  \tightlist
  \item
    Выясните, являются ли приведенные ниже уравнения гыперболическими,
    параболическими или эллиптическими:
  \end{enumerate}

  \begin{itemize}
  \tightlist
  \item
    : \textbf{a)} u\_\{xx\} u\_\{xy\} = 0 , - 6) u\_\{tt\} = u\_\{xx\} +
    u\_x + hu , - \textbf{B)} u\_\{xx\} + 3u\_\{yy\} =
    \textbackslash sin x , - r) u\_\{xx\} + u\_\{yy\} = f(x, y) , -
    \textbackslash mathbf\{D\} ) ~u\_\{rr\} +
    \textbackslash frac\{1\}\{r\} u\_r +
    \textbackslash frac\{1\}\{r\^{}2\} u\_\{\textbackslash theta
    \textbackslash theta\} = f(r, ~\textbackslash theta).
  \end{itemize}
\item
  \begin{enumerate}
  \def\labelenumi{\arabic{enumi}.}
  \setcounter{enumi}{1}
  \tightlist
  \item
    Получите соотношения (23.4), (23.5) и (23.6).
  \end{enumerate}
\item
  \begin{enumerate}
  \def\labelenumi{\arabic{enumi}.}
  \setcounter{enumi}{2}
  \tightlist
  \item
    Убедитесь в том, что уравнение
  \end{enumerate}
\end{itemize}

{[} u\_t = (1+x) u\_\{xx\}, \textbackslash{} u(x, 0) = q(x),
\textbackslash quad -\textbackslash infty \textless{} x \textless{}
\textbackslash infty, {]}

гиперболического типа при всех x и y и найдите x a p a
\textbackslash kappa m e p u c - m u + e c \textbackslash kappa u + e c
\textbackslash kappa u + e c \textbackslash kappa u + e c
\textbackslash kappa u + e c \textbackslash kappa u + e c
\textbackslash kappa u + e c \textbackslash kappa u + e c
\textbackslash kappa u + e c \textbackslash kappa u + e c
\textbackslash kappa u + e c \textbackslash kappa u + e c
\textbackslash kappa u + e c \textbackslash kappa u + e c
\textbackslash kappa u + e c \textbackslash kappa u + e c
\textbackslash kappa u + e c \textbackslash kappa u + e c
\textbackslash kappa u + e c \textbackslash kappa u + e e c
\textbackslash kappa u + e c \textbackslash kappa u + e c
\textbackslash kappa u + e c \textbackslash kappa u + e c
\textbackslash kappa u + e c \textbackslash kappa u + e c
\textbackslash kappa u + e c \textbackslash kappa u + e c
\textbackslash kappa u + e c \textbackslash kappa u + e c
\textbackslash kappa u + e c \textbackslash kappa u + e c
\textbackslash kappa u + e c \textbackslash kappa u + e c
\textbackslash kappa u + e c \textbackslash kappa u + e c
\textbackslash kappa u + e c \textbackslash kappa u + e c
\textbackslash kappa u + e c \textbackslash kap u + e c
\textbackslash kap

{[} pa u + e c \textbackslash kappa u + e c \textbackslash kappa u + e c
\textbackslash kappa u + e c \textbackslash kappa u + e c
\textbackslash kappa u + e c \textbackslash kappa u + e c
\textbackslash kappa u + e c \textbackslash kappa u + e c
\textbackslash kappa u + e c \textbackslash kappa u + e c
\textbackslash kappa u + e c \textbackslash kappa u + e c
\textbackslash kappa u + e c \textbackslash kappa u + e c
\textbackslash kappa u + e c \textbackslash kappa u + e c
\textbackslash kappa u + e c \textbackslash kappa u + e c
\textbackslash kappa u + e c \textbackslash kappa u + e c
\textbackslash kappa u + e c \textbackslash kappa u + e e c
\textbackslash kappa u + e c \textbackslash kappa u + e c
\textbackslash kappa u + e c \textbackslash kappa u + e c
\textbackslash kappa u + e c \textbackslash kappa u + e c
\textbackslash kappa u + e c \textbackslash kappa u + e c
\textbackslash kappa u + e c \textbackslash kappa u + e c
\textbackslash kappa u + e c \textbackslash kappa u + e c
\textbackslash kappa u + e c \textbackslash kappa u + e c
\textbackslash kappa u + e c \textbackslash kappa u + e c
\textbackslash kappa u + e c \textbackslash kappa u + e c
\textbackslash kappa u + e c \textbackslash kappa u + e c
\textbackslash kappa u + e c \textbackslash kappa u + e c
\textbackslash kappa u + e c \textbackslash kappa u + e c
\textbackslash kappa u + e c \textbackslash kappa u + e c
\textbackslash kappa u + e c \textbackslash kappa u + e c
\textbackslash kappa u + e c \textbackslash kappa u + e c
\textbackslash kappa u + e c \textbackslash kappa u + e c
\textbackslash kappa u + e c \textbackslash kappa u + e c
\textbackslash kappa u + e c \textbackslash kappa u + e c
\textbackslash kappa u + e c \textbackslash kappa u + e c
\textbackslash kappa u + e c \textbackslash kappa u + e u + e c
\textbackslash kappa u + e c \textbackslash kappa u + e c
\textbackslash kappa u + e c \textbackslash kappa u + e c
\textbackslash kappa u + e c \textbackslash kappa u + e c
\textbackslash kappa u + e c \textbackslash kappa u + e c
\textbackslash kappa u + e c \textbackslash kappa u + e c
\textbackslash kappa u + e c \textbackslash kappa u + e c
\textbackslash kappa u + e c \textbackslash kappa u + e c
\textbackslash kappa u + e c \textbackslash kappa u + e c
\textbackslash kappa u + e c \textbackslash kappa u + e c
\textbackslash kappa u + e c ~kappa u + e c \textbackslash kappa u + e c
\textbackslash kappa u + e c \textbackslash kappa u + e c
\textbackslash kappa u + e c \textbackslash kappa u + e c
\textbackslash kappa u + e c \textbackslash kappa u + e c
\textbackslash kappa u + e c \textbackslash kappa u + e c
\textbackslash kappa u + e c \textbackslash kappa u + e c
\textbackslash kappa u + e c \textbackslash kappa u + e c
\textbackslash kappa u + e c \textbackslash kappa u + e c
\textbackslash kappa u + e c \textbackslash kappa u + e c
\textbackslash kappa u + e e c \textbackslash kappa u + e c
\textbackslash kappa u + e c \textbackslash kappa u + e c
\textbackslash kappa u + e c \textbackslash kappa u + e c
\textbackslash kappa u + e c \textbackslash kappa u + e c
\textbackslash kappa u + e c \textbackslash kappa u + e c
\textbackslash kappa u + e c \textbackslash kappa u + e c
\textbackslash kappa u + e c \textbackslash kappa u + e c
\textbackslash kappa u + e c \textbackslash kappa u + e c
\textbackslash kappa u + e c \textbackslash kappa u + e c
\textbackslash kappa u + e c \textbackslash kappa u + e c
\textbackslash kappa u + e c \textbackslash kappa u + e c
\textbackslash kappa u + e c \textbackslash kappa u + e c
\textbackslash kappa u + e c \textbackslash kappa u + e c
\textbackslash kappa u + e c \textbackslash kappa u + e c
\textbackslash kappa u + e c \textbackslash kappa u + e c
\textbackslash kappa u + e c \textbackslash kappa u + e c
\textbackslash kappa u + e c \textbackslash kappa u + e c
\textbackslash kappa u + e c \textbackslash kappa u + e c
\textbackslash kappa u + e c \textbackslash kappa u + e c
\textbackslash kappa u + e c \textbackslash kappa u + e c
\textbackslash kappa u + e u + e c \textbackslash kappa u + e c
\textbackslash kappa u + e c \textbackslash kappa u + e c
\textbackslash kappa u + e c \textbackslash kappa u + e c
\textbackslash kappa u + e c \textbackslash kappa u + e c
\textbackslash kappa u + e c \textbackslash kappa u + e c
\textbackslash kappa u + e c \textbackslash kappa u + e c
\textbackslash kappa u + e c \textbackslash kappa u + e c
\textbackslash kappa u + e c \textbackslash kappa u + e c
\textbackslash kappa u + e c \textbackslash kappa u + e c
\textbackslash kappa u + e c \textbackslash kappa u + e c ~kappa u + e c
\textbackslash kappa u + e c \textbackslash kappa u + e c
\textbackslash kappa u + e c \textbackslash kappa u + e c
\textbackslash kappa u + e c \textbackslash kappa u + e c
\textbackslash kappa u + e c \textbackslash kappa u + e c
\textbackslash kappa u + e c \textbackslash kappa u + e c
\textbackslash kappa u + e c \textbackslash kappa u + e c
\textbackslash kappa u + e c {]}

\begin{enumerate}
\def\labelenumi{\arabic{enumi}.}
\setcounter{enumi}{3}
\tightlist
\item
  Приведите уравнение из предыдущей задачи к каноническому виду
\end{enumerate}

{[} u\_1 = (1 + \textbackslash varepsilon\_1) u\_\{xx\} {]}

\begin{enumerate}
\def\labelenumi{\arabic{enumi}.}
\setcounter{enumi}{4}
\tightlist
\item
  Продолжите вадачу 4 и приведите уравнение ко второй канонической форме
\end{enumerate}

{[} u = u\_1 + \textbackslash varepsilon u\_1 +
\textbackslash varepsilon\^{}2 u\_2 + \textbackslash cdots {]}

\begin{enumerate}
\def\labelenumi{\arabic{enumi}.}
\setcounter{enumi}{5}
\tightlist
\item
  Найдите характеристики уравнения
\end{enumerate}

{[} P\_\{0\} {]}

Преобразуйте уравнение к новым координатам, решите и, вернувшись к
координатам x и y, получите решение исходной вадачи.

\section{ВОЛНОВОЕ УРАВНЕНИЕ В СВОБОДНОМ ПРОСТРАНСТВЕ (ДВУМЕРНЫЕ И
ТРЕХМЕРНЫЕ
ЗАДАЧИ)}\label{ux432ux43eux43bux43dux43eux432ux43eux435-ux443ux440ux430ux432ux43dux435ux43dux438ux435-ux432-ux441ux432ux43eux431ux43eux434ux43dux43eux43c-ux43fux440ux43eux441ux442ux440ux430ux43dux441ux442ux432ux435-ux434ux432ux443ux43cux435ux440ux43dux44bux435-ux438-ux442ux440ux435ux445ux43cux435ux440ux43dux44bux435-ux437ux430ux434ux430ux447ux438}

ЦЕЛЬ ЛЕКЦИИ: Решить задачу с начальными условиями

(УЧП)

\begin{verbatim}
\begin{cases} \frac{\partial u_{0}}{\partial t} = \frac{\partial^{2} u_{1}}{\partial x^{2}}, \\ u_{0}(x, 0) = \varphi(x), \\ P_{1} \begin{cases} \frac{\partial u_{1}}{\partial t} - \frac{\partial^{2} u_{1}}{\partial x^{2}} = x \frac{\partial^{2} u_{0}}{\partial x^{2}}, \\ u_{1}(x, 0) = 0 \\ \vdots & \vdots \end{cases}
\end{verbatim}

{[} \textbackslash begin\{aligned\} \textbackslash Delta u + u\^{}2 = 0,
\& 0 \textbackslash leqslant r \textless{} 1, \textbackslash quad 0
\textbackslash leqslant \textbackslash theta \textless{}
2\textbackslash pi, \textbackslash{} u(1, \textbackslash theta) =
\textbackslash cos \textbackslash theta, \& 0 \textbackslash leqslant
\textbackslash theta \textless{} 2\textbackslash pi,
\textbackslash end\{aligned\} {]}

{[} \textbackslash begin\{cases\} -\textbackslash infty \textless{} x
\textless{} \textbackslash infty, \textbackslash{} -\textbackslash infty
\textless{} y \textless{} \textbackslash infty, \textbackslash{}
-\textbackslash infty \textless{} z \textless{} \textbackslash infty,
\textbackslash end\{cases\} {]}

{[} u(r, \textbackslash theta) = u\_0(r, \textbackslash theta) +
\textbackslash frac\{1\}\{4\}u\_1(r, \textbackslash theta) {]}

{[} \textbackslash Delta u = 0, {]}

{[} q\_\{xx\} + \textbackslash varphi\_\{uu\} = 0 {]}

{[} \textbackslash begin\{cases\} u(x, y, z, 0) =
\textbackslash varphi(x, y, z), \textbackslash{} u\_\{t\}(x, y, z, 0) =
\textbackslash psi(x, y, z) \textbackslash end\{cases\} {]}

{[} q\_\{uu\} + q\_\{vv\} = 0 {]}

\pandocbounded{\includegraphics[keepaspectratio]{images/_page_353_Figure_9.jpeg}}

в трехмерном пространстве и показать, что полученное решение
удовлетворяет \emph{принципу Гюйгенса}. Для решения соответствующей
двумерной задачи

(УЧП)

{[} w = f(z) {]}

\pandocbounded{\includegraphics[keepaspectratio]{images/_page_355_Figure_5.jpeg}}

Рис. 47.1. Площина комплексної змінної \(\boldsymbol{z}\) та
некольорових формул бігунів

\pandocbounded{\includegraphics[keepaspectratio]{images/_page_355_Figure_7.jpeg}}

{[} \textbackslash begin\{cases\} -\textbackslash infty \textless{} x
\textless{} \textbackslash infty, \textbackslash{} -\textbackslash infty
\textless{} y \textless{} \textbackslash infty,
\textbackslash end\{cases\} {]}

{[} u + iv = (x + iy)\^{}2 = x\^{}2 - y\^{}2 + 2ixy {]}

\begin{verbatim}
\left\{
\end{verbatim}

(НУ)

{[} \textbackslash begin\{cases\} u(x, y, 0) = \textbackslash varphi(x,
y), \textbackslash{} u\_t(x, y, 0) = \textbackslash psi(x, y)
\textbackslash end\{cases\} {]}

используется метод спуска и демонстрируется, что двумерные решения не
удовлетворяют принципу Гюйгенса. И наконец, с помощью метода спуска
показано, что решение одномерного уравнения выражается формулой
Даламбера, с которой мы познакомились ранее.

В предыдущих лекциях мы рассматривали задачу с начальными условиями для
бесконечной колеблющейся струны и показали, что ее решение дается
формулой Даламбера. Читателю должно быть ясно, что в трехмерном
пространстве одномерное волновое уравнение описывает плоские волны.
Например, звуковые или электромагнитные волны на достаточно больших
расстояниях от источников можно считать плоскими волнами и,
следовательно, описывать одномерным уравнением. Существует следующая
терминология:

\begin{itemize}
\tightlist
\item
  \begin{enumerate}
  \def\labelenumi{\arabic{enumi})}
  \tightlist
  \item
    одномерные волны называются плоскими волнами,
  \end{enumerate}
\item
  \begin{enumerate}
  \def\labelenumi{\arabic{enumi})}
  \setcounter{enumi}{1}
  \tightlist
  \item
    двумерные волны называются цилиндрическими волнами.
  \end{enumerate}
\end{itemize}

\begin{enumerate}
\def\labelenumi{\arabic{enumi})}
\setcounter{enumi}{2}
\tightlist
\item
  трехмерные волны называются сферическими волнами.
\end{enumerate}

Другими словами, одномерное волновое уравнение может описывать либо
плоские волны в пространствах большего числа измерений, либо колебания
струны. Задача лекции --- обобщить формулу Даламбера на случай двух и
трех измерений.

\section{Волны в трехмерном
пространстве}\label{ux432ux43eux43bux43dux44b-ux432-ux442ux440ux435ux445ux43cux435ux440ux43dux43eux43c-ux43fux440ux43eux441ux442ux440ux430ux43dux441ux442ux432ux435}

Рассмотрим сферические волны в трехмерном пространстве при ваданных
начальных условиях, т. е. решим задачу с начальными условиями

(24.1)

MATHBLK0677X

MATHBLK0678X

{[} \textbackslash begin\{cases\} -\textbackslash infty \textless{} x
\textless{} \textbackslash infty, \textbackslash{} -\textbackslash infty
\textless{} y \textless{} \textbackslash infty, \textbackslash{}
-\textbackslash infty \textless{} z \textless{} \textbackslash infty,
\textbackslash end\{cases\} {]}

MATHBLK0679X

MATHBLK0680X

(24.1)

MATHBLK0681X

{[} \textbackslash begin\{cases\} u(x, y, z, 0) =
\textbackslash varphi(x, y, z), \textbackslash{} u\_\{t\}(x, y, z, 0) =
\textbackslash psi(x, y, z). \textbackslash end\{cases\} {]}

MATHBLK0682X

MATHBLK0683X

MATHBLK0684X

{[} \textbackslash begin\{aligned\} u\_\{tt\} \&= c\^{}2
\textbackslash Delta u, \textbackslash{} (HY) \&
\textbackslash begin\{cases\} u(x, y, z, 0) = 0, \textbackslash{}
u\_t(x, y, z, 0) = \textbackslash psi(x, y, z),
\textbackslash end\{cases\} \textbackslash end\{aligned\} {]}

MATHBLK0685X

MATHBLK0686X

где экаком \(\Delta\) обозначен дифференциальный оператор

MATHBLK0687X

Эту задачу можно решить методом \(\Phi\) урье \(^{1}\) ), и оказывается,
что решение представимо в виде

MATHBLK0688X

где \(\overline{\psi}\) --- \emph{среднее} значение начального
распределения \(\psi\) \emph{по сфере} радиуса ct с центром в точке (x,
y, z), т. е.

MATHBLK0689X

\(+ ct \cos \varphi\) ) \((ct)^2 \sin \varphi d\varphi d\theta\) .

Аргументы функции \(\psi\) пробегают по поверхности сферы, если величины
\(\theta\) и \(\phi\) изменяются в пределах \([0, 2\pi]\) и \([0, \pi]\)
соответственно (рис. 24.1). Эту формулу можно интерпретировать сле-

См Владимиров В. С. Уравнения математической физики. --- М.: Наука, 1971
и др. издания. --- Прим ред.

дующим образом: каждая точка пространства излучает расходящуюся (со
скоростью c) сферическую волну, через t секунд точка с координатами (x,
y, z) окажется под воздействием начального Р И С

\pandocbounded{\includegraphics[keepaspectratio]{images/_page_177_Picture_3.jpeg}}

. 24 1. Решение является средним значением начатьного распределения на
сфере.

возмущения, сосредоточенного на сфере раднуса ct с центром в заданной
точке (см. рис. 24.2).

Формула (24.3) позволяет рассчитать на ЭВМ решение задачи для
большинства начальных условий. Может быть, читателю будет интересно
найти решение для некоторых простых функций ф.

\pandocbounded{\includegraphics[keepaspectratio]{images/_page_177_Picture_7.jpeg}}

РИС. 24.2. Начальное возмущение \(\psi\) распространяется из каждой
точки во все стороны.

Чтобы получить полное решение, рассмотрим вторую половину задачи

(24.4)

MATHBLK0690X

MATHBLK0691X

MATHBLK0692X

MATHBLK0693X

{[} \textbackslash begin\{aligned\} (\textbackslash text\{YAII\}) \&
u\_\{tt\} = c\^{}2 \textbackslash Delta u, \textbackslash{}
(\textbackslash text\{HY\}) \& \textbackslash begin\{cases\} u = 0,
\textbackslash{} u\_t = \textbackslash varphi,
\textbackslash end\{cases\} \textbackslash end\{aligned\} {]}

MATHBLK0694X

MATHBLK0695X

{[} u =
\textbackslash frac\{\textbackslash partial\}\{\textbackslash partial
t\} {[}t\textbackslash overline\{\textbackslash varphi\}{]}. {]}

MATHBLK0696X

MATHBLK0697X

MATHBLK0698X

MATHBLK0699X

MATHBLK0700X

MATHBLK0701X

\begin{enumerate}
\def\labelenumi{\arabic{enumi}.}
\setcounter{enumi}{4}
\tightlist
\item
  \(\cos(mx), m=1, 2, \dots\)
\end{enumerate}

MATHBLK0702X

MATHBLK0703X

MATHBLK0704X

MATHBLK0705X

MATHBLK0706X

MATHBLK0707X

x

MATHBLK0708X

MATHBLK0709X

\_ \textbar{}

MATHBLK0710X

{[} \textbackslash begin\{cases\} u(x, y, 0) = \textbackslash varphi(x,
y), \textbackslash{} u\_t(x, y, 0) = \textbackslash psi(x, y),
\textbackslash end\{cases\} {]}

MATHBLK0711X

MATHBLK0712X

MATHBLK0713X

мы рассмотрим розв'язок трехмерной задачи при начальных условиях,
зависящих лише от двух змінних x и y.

Сделав це, получим, що трехмерная формула

{[} \textbackslash Delta u = u\_\{xx\} + u\_\{yy\} {]}

для розв'язки и буде описывать цилиндрические хвилі и, отже, даст
розв'язок двумерной задачи. Цей метод називається методом спуска. Після
всех вычислений 1) (отнюдь не тривиальных) получаем

{[} \textbackslash Delta u = u\_\{rr\} +
\textbackslash frac\{1\}\{r\}u\_r +
\textbackslash frac\{1\}\{r\^{}2\}u\_\{\textbackslash theta\textbackslash theta\}
{]}

Перед нами розв'язок двумерного волнового уравнения, и хотя, вероятнее
всего, его практическое применение подразумевает использование ЭВМ,
интересно проанализировать его с точки зрения принципа Гюйгенса. В цьому
решении оба интеграла от начальных возмущений берутся по внутренности
(слово «внутренность» играет здесь решающую роль) круга радиусом сt е
центром в точке (x, y). Другими словами, якщо проанализировать це
розв'язок так же, как мы анализировали розв'язок в трехмерном случае, мы
увидим, що начальное возмущение распространяется с резко очерченным
передним фронтом и размытым задвим фронтом. Таким образом, принцип
Гюйгенса не выполняется в пространстве двух измерений.

И наконец, предположим, що начальные возмущения \(\phi\) и \(\phi\)
зависят лише от одной переменной. В цьому случае возникают плоские
хвилі, и після применения метода спуска получаем

\&lt;sup\textgreater2) См., наприклад, Тихонов А. Н., Самарский А. А.
Уравнения математической физики. --- М.: Наука, 1972 и др. издания. ---
Прим. ред.

хорошо известную формулу Даламбера

{[} \textbackslash Delta u = u\_\{xx\} + u\_\{yy\} + u\_\{zz\} {]}

Отметим, що и в цьому случае применение метода спуска требует проведения
нетривиальных вычислений. В формуле Даламбера начальное смещение ф
приводит к образованию резкого заднего фронта, а начальная скорость ф
его размывает. Другими словами, одномерный случай необычен тем, що
принцип Гюйгенса выполняется для начальных смещений и не выполняется для
начальных скоростей. Можна утверждать, що в целом принцип Гюйгенса в
одномерном случае не выполняется.

\section{ЗАМЕЧАНИЯ}\label{ux437ux430ux43cux435ux447ux430ux43dux438ux44f-8}

Мы не стали подробно излагать метод спуска, поскольку не котели
останавливаться на вопросе о том, как можна понизить размерность
входящих в формулы интегралов. Все эти обчислення можна найти в учебной
литературе, список которой приведен в конце книги. Общая идея
заключается в том, що при решении задач в пространстве невысокой
размерности можна воспользоваться известным решением аналогичной задачи
в пространстве более высокой размерности, а затем упростить его, считая,
що начальные и граничні умови не зависят от части змінних. Читателю
должно быть ясно, що метод спуска можна применять не лише к
рассмотренной задаче.

\section{ЗАДАЧИ}\label{ux437ux430ux434ux430ux447ux438-8}

\begin{enumerate}
\def\labelenumi{\arabic{enumi}.}
\tightlist
\item
  Покажите, що в одномерном случае розв'язок задачи
\end{enumerate}

{[} \textbackslash Delta u = u\_\{rr\} +
\textbackslash frac\{1\}\{r\}u\_r +
\textbackslash frac\{1\}\{r\^{}2\}u\_\{\textbackslash theta\textbackslash theta\}
+ u\_\{zz\} {]}

{[} \textbackslash Delta u = u\_\{rr\} +
\textbackslash frac\{2\}\{r\}u\_r +
\textbackslash frac\{1\}\{r\^{}2\}u\_\{\textbackslash theta\textbackslash theta\}
+
\textbackslash frac\{\textbackslash cot\textbackslash theta\}\{r\^{}2\}
{]}

{[} \textbackslash Delta u = u\_\{rr\} +
\textbackslash frac\{2\}\{r\}u\_r +
\textbackslash frac\{1\}\{r\^{}2\}u\_\{\textbackslash theta\textbackslash theta\}
+
\textbackslash frac\{\textbackslash cot\textbackslash theta\}\{r\^{}2\}u\_0
+
\textbackslash frac\{1\}\{r\^{}2\textbackslash sin\textbackslash theta\}u\_\{\textbackslash phi\textbackslash phi\}
{]}

{[} \textbackslash begin\{cases\} u(x, ~0) = \textbackslash varphi(x),
\textbackslash{} u\_t(x, ~0) = 0 \textbackslash end\{cases\} {]}

\pandocbounded{\includegraphics[keepaspectratio]{images/_page_370_Figure_1.jpeg}}

\pandocbounded{\includegraphics[keepaspectratio]{images/_page_370_Picture_2.jpeg}}

\pandocbounded{\includegraphics[keepaspectratio]{images/_page_370_Figure_3.jpeg}}

{[} \textbackslash begin\{aligned\} \& (\textbackslash forall
\textbackslash mathsf\{Y\}\textbackslash Pi) \textbackslash quad
u\_\{tt\} = c\^{}2 u\_\{xx\}, \textbackslash{} \&
(\textbackslash mathsf\{H\} \textbackslash mathsf\{Y\})
\textbackslash quad \textbackslash begin\{cases\}
u\textbackslash left(x, ~0\textbackslash right) = 0, \textbackslash{}
u\_t\textbackslash left(x, ~0\textbackslash right) =
\textbackslash phi\textbackslash left(x\textbackslash right).
\textbackslash end\{cases\} \textbackslash end\{aligned\} {]}

{[} -\textbackslash infty \textless{} x \textless{}
+\textbackslash infty {]}

{[} -\textbackslash infty \textless{} y \textless{}
+\textbackslash infty {]}

\begin{enumerate}
\def\labelenumi{\arabic{enumi}.}
\setcounter{enumi}{1}
\tightlist
\item
  Результаты задачи 1 используйте для розв'язки следующей задачи:
\end{enumerate}

(УЧП)

{[} -\textbackslash infty \textless{} z \textless{}
+\textbackslash infty {]}

(НУ)

\begin{verbatim}
\begin{array}{lll}
\end{verbatim}

{[} \textbackslash begin\{cases\} u(x, 0) = x, \textbackslash{} u\_t(x,
0) = 0. \textbackslash end\{cases\} {]}

{[} -\textbackslash infty \& \textless{} x \textless{}
+\textbackslash infty \& r \textbackslash geqslant 0, ~0
\textbackslash leqslant \textbackslash theta \textbackslash leqslant
2\textbackslash pi, \textbackslash{} {]}

\begin{enumerate}
\def\labelenumi{\arabic{enumi}.}
\setcounter{enumi}{2}
\tightlist
\item
  Проиллюстрируйте и прокомментируйте розв'язок трехмерной задачи
\end{enumerate}

(УЧП)

{[} -\textbackslash infty \& \textless{} y \textless{}
+\textbackslash infty \& -\textbackslash alpha \textbackslash leqslant z
\textbackslash leqslant \textbackslash alpha, \textbackslash{} {]}

{[} , (x, y, z) \textbackslash in R\^{}3 , u = 0 u\_t = {]}

MATHBLK0738X

{[} \textbackslash begin\{cases\} 1, \& x\^{}2 + y\^{}2 + z\^{}2
\textbackslash leqslant 1, \textbackslash{} 0, \& \textbackslash text\{в
інших точках.\} \textbackslash end\{cases\} {]}

MATHBLK0739X

\begin{enumerate}
\def\labelenumi{\arabic{enumi}.}
\setcounter{enumi}{3}
\tightlist
\item
  Решите аналогичную (см. задачу 3) задачу для двумерного уравнения
\end{enumerate}

(УЧП)

(НУ)

{[} -\textbackslash infty \& \textless{} z \textless{}
+\textbackslash infty \& x = r \textbackslash cos \textbackslash theta,
\textbackslash{} {]}

{[} y = r \textbackslash sin \textbackslash theta, \textbackslash{} {]}

\begin{enumerate}
\def\labelenumi{\arabic{enumi}.}
\setcounter{enumi}{4}
\tightlist
\item
  Решите аналогичную (см. задачу 3) задачу для одномерных плоских волн
\end{enumerate}

(УЧП)

(НУ)

{[} z = z. {]}

\begin{verbatim}
\end{aligned}
\end{verbatim}

\begin{itemize}
\tightlist
\item
  \begin{enumerate}
  \def\labelenumi{\arabic{enumi}.}
  \setcounter{enumi}{5}
  \tightlist
  \item
    Какова физическая природа того факта, що принцип Гюйгенса не
    справедлив в двумерном случае?
  \end{enumerate}
\item
  \begin{enumerate}
  \def\labelenumi{\arabic{enumi}.}
  \setcounter{enumi}{6}
  \tightlist
  \item
    Воспользуйтесь формулой Лейбница
  \end{enumerate}
\end{itemize}

{[} r \textbackslash geqslant 0 {]}

и продифференцируйте интеграл

{[} \textbackslash Delta u = 0 {]}

по времени t.

\section{КОНЕЧНЫЕ ПРЕОБРАЗОВАНИЯ ФУРЬЕ (СИНУС- И
КОСИНУС-ПРЕОБРАЗОВАНИЯ)}\label{ux43aux43eux43dux435ux447ux43dux44bux435-ux43fux440ux435ux43eux431ux440ux430ux437ux43eux432ux430ux43dux438ux44f-ux444ux443ux440ux44cux435-ux441ux438ux43dux443ux441--ux438-ux43aux43eux441ux438ux43dux443ux441-ux43fux440ux435ux43eux431ux440ux430ux437ux43eux432ux430ux43dux438ux44f}

ЦЕЛЬ ЛЕКЦИИ: Ввести два новых интегральных преобразования (конечные
спиус- и косинус-преобразования):

{[} \textbackslash Delta u + \textbackslash lambda\^{}2 u = 0 {]}

(конечное синус-преобразование),

{[} \textbackslash Delta u = k {]}

(конечное косинус-преобразование),

{[} \textbackslash Delta u + k(E - V)u = 0 {]}

(обратное синус-преобразование),

(обратное косинус-преобразование),

и показать, как с помощью этих преобразований решать краевые задачи (в
частности, неоднородные). Ранее мы познакомились с регулярными
преобразованиями Фурье и Лапласа и увидели, как применение этих
преобразований к уравнению с частными производными сводит задачу к
решению обыкновениого дифференциального уравнения. Обычное
преобразование Фурье подразумевает, что преобразуемая переменная
изменяется от \(-\infty\) до \(+\infty\) , и, следовательно, оно
применяется для решения задач в свободном пространстве (без границ). В
этой лекции мы покажем, как решать краевые задачи (с границами),
преобразуя ограниченные переменные (с таким преобразованием мы
встречаемся впервые).

Давайте на время забудем о наших задачах, а займемся самими
преобразованнями. Дадим определение преобразований, приведем формулы
обращения и основные свойства. Для чего нужны преобразования и как ими
пользоваться, мы узнаем несколько позже. Впрочем, если говорить кратко,
то суть интегральных преобразований в разложении входящих в задачу
функций по

частотам, решении семейства задач для каждой частоты и после-

дующем суммировании всех полученных результатов.

Начнем с функции \(f(x)\), определенной на интервале {[}0, L{]}.
Конечные синус- и косинус-преобразования этой функции определяются
формулами

Конечное синус-преобразование
\(S[f] = S_n = \frac{2}{L} \int_{\sigma}^{L} f(x) \sin(n\pi x/L) dx\)
\(n = 1, 2, \dots\)
\(C[f] = C_n = \frac{2}{L} \int_{\sigma}^{L} f(x) \cos(n\pi x/L) dx\)
\(n = 0, 1, \dots\)

Конечное косинус-преобразование

Читатель, наверное, заметил, что эти преобразования не отличаются от
формул для определения коэффициентов Фурье при разложении их в ряд по
синусам или косинусам. Формулы обращения для этих преобразований---это
обычные ряды Фурье по синусам и косинусам.

Обратное синус-преобразование

Обратное косинус-преобразование

Отметим, що в обратном косинус-преобразовании суммирование начинается с
n=0, а в синус-преобразовании --- с n=1.

\section{Примеры
синус-преобразования}\label{ux43fux440ux438ux43cux435ux440ux44b-ux441ux438ux43dux443ux441-ux43fux440ux435ux43eux431ux440ux430ux437ux43eux432ux430ux43dux438ux44f}

Пусть

тогда

{[} u\_\{tt\} = c\^{}2 \textbackslash Delta u {]}

График функції \(f(x)\) и ее преобразования изображены на рис. 25.1.
Обратное преобразование записывается в виде

{[} u\_\{tt\} = c\^{}2 \textbackslash Delta u - hu\_t {]}

Знаете ли вы, как выглядит график цієї функції ене отрезка {[}0, 1{]}?
Подумайте над этим. Отметим також, що синус-преобразование функції
\(f(x)\) определено лише для положительных целых n.~Другими словами,
конечные синус- и косинус-преобразования преобразуют функції в числовые
последовательности.

\pandocbounded{\includegraphics[keepaspectratio]{images/_page_185_Figure_3.jpeg}}

Рис. 25.1. Графики функції \(f(x)\) = 1 и ее преобразования.

\section{Свойства
преобразований}\label{ux441ux432ux43eux439ux441ux442ux432ux430-ux43fux440ux435ux43eux431ux440ux430ux437ux43eux432ux430ux43dux438ux439}

Перш ніж чем приступить к решению задач, мы повинні получить некоторые
свойства цих преобразований.

Якщо \(u(x, t)\) --- функція \(\partial \theta yx\) змінних и мы
проводим преобразование по переменной x, то

{[} u\_t = \textbackslash alpha\^{}2 u\_\{xx\} {]}

{[} u\_t = \textbackslash alpha\^{}2 u\_\{xx\} - hu\_x {]}

(Заметим, що преобразование проводилось по переменной x, и полученная
последовательность зависит лише от времени t.)

А как быть с похідними? Ниже приводятся несколько полезных формул

{[} u\_t = \textbackslash alpha\^{}2 u\_\{xx\} - ku {]}

{[} u\_t = \textbackslash alpha\^{}2 u\_\{xx\} + f(x, t) {]}

{[} \textbackslash begin\{split\}
S\textbackslash left{[}u\_\{t\}\textbackslash right{]} \&=
\textbackslash frac\{dS\textbackslash left{[}u\textbackslash right{]}\}\{dt\},
\textbackslash quad
S\textbackslash left{[}u\_\{tt\}\textbackslash right{]} =
\textbackslash frac\{d\textsuperscript{\{2\}S\textbackslash left{[}u\textbackslash right{]}\}\{dt}\{2\}\},
\textbackslash{} S\textbackslash left{[}u\_\{xx\}\textbackslash right{]}
\&=
-\textbackslash left{[}n\textbackslash pi/L\textbackslash right{]}\^{}\{2\}S\textbackslash left{[}u\textbackslash right{]}
+
\textbackslash frac\{2n\textbackslash pi\}\{L\^{}\{2\}\}\textbackslash left{[}u\textbackslash left(0,~t\textbackslash right)
+
(-1)\^{}\{n+2\}u\textbackslash left(L,~t\textbackslash right)\textbackslash right{]},
\textbackslash{} C\textbackslash left{[}u\_\{xx\}\textbackslash right{]}
\&=
-\textbackslash left{[}n\textbackslash pi/L\textbackslash right{]}\^{}\{2\}C\textbackslash left{[}u\textbackslash right{]}
-
\textbackslash frac\{2\}\{L\}\textbackslash left{[}u\_\{x\}\textbackslash left(0,~t\textbackslash right)
+
(-1)\^{}\{n+1\}u\_\{x\}\textbackslash left(L,~t\textbackslash right)\textbackslash right{]}.
\textbackslash end\{split\} {]}

Уравнение одномерной диффузии Уравнение конвективной диффузии

рівняння теплопровідності с поглощением рівняння теплопровідності с
источником

\section{\texorpdfstring{\textbf{КРОССВОРД}}{КРОССВОРД}}\label{ux43aux440ux43eux441ux441ux432ux43eux440ux434}

\pandocbounded{\includegraphics[keepaspectratio]{images/_page_371_Figure_3.jpeg}}

{[} (\textbackslash text\{УЧП\}) \textbackslash quad u\_\{tt\} =
u\_\{xx\} + \textbackslash sin(\textbackslash pi x), \textbackslash quad
0 \textless{} x \textless{} 1, \textbackslash quad 0 \textless{} t
\textless{} \textbackslash infty, {]}

{[} P\_\{0\} {]}

{[} (\textbackslash text\{ГУ\}) \textbackslash quad {]}

{[} \textbackslash begin\{cases\} u(0, t) = 0, \textbackslash{} u(1, t)
= 0, \textbackslash end\{cases\} {]}

\begin{verbatim}
\begin{cases} \frac{\partial u_{0}}{\partial t} = \frac{\partial^{2} u_{1}}{\partial x^{2}}, \\ u_{0}(x, 0) = \varphi(x), \\ P_{1} \begin{cases} \frac{\partial u_{1}}{\partial t} - \frac{\partial^{2} u_{1}}{\partial x^{2}} = x \frac{\partial^{2} u_{0}}{\partial x^{2}}, \\ u_{1}(x, 0) = 0 \\ \vdots & \vdots \end{cases}
\end{verbatim}

{[} \textbackslash begin\{cases\} u(x, 0) = 1, \textbackslash{} u\_t(x,
0) = 0, \textbackslash end\{cases\} {]}

Решение задачи выполним за три следующих шага:

ШАГ І (Определение подходящего преобразования.)

Поскольку переменная x меняется от 0 до 1, мы воспользуемся конечным
преобразованием, а именно синус-преобразованием; позже станет ясно
почему. Мы могли бы решить эту задачу с помощью преобразования Лапласа
по переменной t (по трудоемкости этот метод не отличается от конечного
синус-преобразования).

ШАГ 2 (Выполнение преобразования).

Для удобства будем использовать обозначение \(S_n(t) = S[u]\) . Применим
синус-преобразование к исходному уравнению

{[} \textbackslash begin\{aligned\} \textbackslash Delta u + u\^{}2 = 0,
\& 0 \textbackslash leqslant r \textless{} 1, \textbackslash quad 0
\textbackslash leqslant \textbackslash theta \textless{}
2\textbackslash pi, \textbackslash{} u(1, \textbackslash theta) =
\textbackslash cos \textbackslash theta, \& 0 \textbackslash leqslant
\textbackslash theta \textless{} 2\textbackslash pi,
\textbackslash end\{aligned\} {]}

и воспользуемся тождеством для синус-преобразования, получим

{[} \textbackslash begin\{split\}
\textbackslash frac\{d\textsuperscript{2S\_n\textbackslash left(t\textbackslash right)\}\{dt}2\}
\&= -\textbackslash left(\textbackslash pi
n\textbackslash right)\^{}2S\_n\textbackslash left(t\textbackslash right)
+
2n\textbackslash pi\textbackslash left{[}u\textbackslash left(0,~t\textbackslash right)
+
\textbackslash left(-1\textbackslash right)\^{}\{n+1\}u\textbackslash left(1,~t\textbackslash right)\textbackslash right{]}
+ D\_n\textbackslash left(t\textbackslash right) = \textbackslash{} \&=
-\textbackslash left(n\textbackslash pi\textbackslash right)\^{}2S\_n\textbackslash left(t\textbackslash right)
+ D\_n\textbackslash left(t\textbackslash right),
\textbackslash end\{split\} {]}

{[} u(r, \textbackslash theta) = u\_0(r, \textbackslash theta) +
\textbackslash frac\{1\}\{4\}u\_1(r, \textbackslash theta) {]}

где

{[} \textbackslash Delta u = 0, {]}

{[} q\_\{xx\} + \textbackslash varphi\_\{uu\} = 0 {]}

{[} q\_\{uu\} + q\_\{vv\} = 0 {]}

Якщо теперь преобразовать початкові умови краевой задачи, то получим
початкові умови для обыкновенного дифференциального уравнения

\pandocbounded{\includegraphics[keepaspectratio]{images/_page_353_Figure_9.jpeg}}

{[} \textbackslash begin\{cases\} 4/n\textbackslash pi, \& n = 1, 3,
\textbackslash dots, \textbackslash{} 0, \& n = 2, 4,
\textbackslash dots, \textbackslash end\{cases\} {]}

{[} w = f(z) {]}

Итак, решим теперь новое семейство задач Кощи:

(ОДУ)

\pandocbounded{\includegraphics[keepaspectratio]{images/_page_355_Figure_5.jpeg}}

Рис. 47.1. Площина комплексної змінної \textbackslash boldsymbol\{z\} та
некольорових формул бігунів

{[} \textbackslash begin\{cases\} 1, \& n = 1, \textbackslash{} 0, \& n
= 2, 3, \textbackslash dots, \textbackslash end\{cases\} {]}

\pandocbounded{\includegraphics[keepaspectratio]{images/_page_355_Figure_7.jpeg}}

{[} u + iv = (x + iy)\^{}2 = x\^{}2 - y\^{}2 + 2ixy {]}

\begin{verbatim}
\begin{cases} S_n(0) = \begin{cases} 4/n\pi, & n = 1, 3, \dots, \\ 0, & n = 2, 4, \dots, \end{cases}

\left\{
\end{verbatim}

{[} \textbackslash begin\{cases\} \textbackslash frac\{dS\_n(0)\}\{dt\}
= 0. \textbackslash end\{cases\} {]}

{[} \textbackslash begin\{aligned\}
u=x\textsuperscript{2-y}2,\textbackslash{} v=2xy,
\textbackslash end\{aligned\} {]}

В результате получаем

{[} \textbackslash Rightarrow. {]}

,

где

Один із способів розв'язання цієї проблеми --- застосувати перетворення
Фур'є до змінної x, але набагато краще конформно відобразити верхню
половину площини у площину (u, v). Читачеві не буде складно побачити, що
конформне відображення

{[} \textbackslash begin\{cases\} 0, \& n = 2, 4, \textbackslash dots,
\textbackslash{}
(4/n\textbackslash pi)\textbackslash cos(n\textbackslash pi t), \& n =
3, 5, 7, \textbackslash dots \textbackslash end\{cases\} {]}

Вона має такі властивості (див. RNS. 47.4).

\begin{verbatim}
1. Вона відображає верхню півплощину площини z у смуги -\infty < u < \infty , 0 < v < \pi , у площині w.
\end{verbatim}

Следовательно, решение нашей задачи можно представить в виде

\section{ЗАМЕЧ
АНИЯ}\label{ux437ux430ux43cux435ux447-ux430ux43dux438ux44f}

\begin{enumerate}
\def\labelenumi{\arabic{enumi}.}
\tightlist
\item
  Для розв'язки задачи методом конечного синус- или
  косинуспреобразования краевые условия при x=0 и x=L повинні иметь вид
\end{enumerate}

{[} \textbackslash begin\{cases\} u(0, t) = f(t) \textbackslash{} u(L,
t) = g(t) \textbackslash end\{cases\} {]}

{[} \textbackslash begin\{cases\} (\textbackslash text\{УЧП\}) \&
\textbackslash varphi\_\{uu\} + \textbackslash varphi\_\{vv\} = 0,
\textbackslash quad -\textbackslash infty \textless{} u \textless{}
\textbackslash infty, \textbackslash quad 0 \textless{} v \textless{}
\textbackslash pi, \textbackslash{} (\textbackslash text\{47.2\}) \&
\textbackslash begin\{cases\} \textbackslash varphi(u, 0) = 0,
\textbackslash{} \textbackslash varphi(u, \textbackslash pi) = 1.
\textbackslash end\{cases\} \textbackslash end\{cases\} {]}

{[} w = u + iv =
\textbackslash ln\textbackslash left{[}\textbackslash frac\{z-1\}\{z+1\}\textbackslash right{]}
= {]}

{[} =
\textbackslash ln\textbackslash left{[}\textbackslash frac\{z-1\}\{z+1\}\textbackslash right{]}
+ i
\textbackslash arg\textbackslash left{[}\textbackslash frac\{z-1\}\{z+1\}\textbackslash right{]}
{]}

\begin{verbatim}
\begin{array}{ll} (\mathrm{Y} \mathrm{U} \Pi) & u_t = u_{xx}, \quad 0 < x < 1, \quad 0 < t < \infty, \\ (\mathrm{\Gamma} \mathrm{Y} & \left\{ \begin{array}{ll} u_x \left( 0, \ t \right) = 0 \\ u_x \left( 1, \ t \right) = 0, \end{array}
\end{verbatim}

{[} v =
\textbackslash arg\textbackslash left{[}\textbackslash frac\{z-1\}\{z+1\}\textbackslash right{]}
=
\textbackslash arg\textbackslash left{[}\textbackslash frac\{x+iy-1\}\{x+1+iy\}\textbackslash right{]}
= {]}

{[} =
\textbackslash arg\textbackslash left{[}\textbackslash frac\{x\textsuperscript{2+y}2-1+i2y\}\{(x+1)\textsuperscript{2+y}2\}\textbackslash right{]}
= {]}

\begin{enumerate}
\def\labelenumi{\arabic{enumi}.}
\setcounter{enumi}{1}
\tightlist
\item
  Решите в общем виде задачу
\end{enumerate}

{[} =
\textbackslash arg\textbackslash left{[}\textbackslash frac\{2y\}\{x\textsuperscript{2+y}2-1\}\textbackslash right{]}
{]}

{[} \textbackslash varphi(x, y) =
\textbackslash frac\{1\}\{\textbackslash pi\}
\textbackslash arctan\textbackslash left{[}\textbackslash frac\{2y\}\{x\^{}2
+ y\^{}2 - 1\}\textbackslash right{]}. {]}

\begin{verbatim}
\right. & 0 < t < \infty, \\ (\text{НУ}) & u(x, 0) = 0, \quad 0 \leqslant x \leqslant 1. \end{aligned}
\end{verbatim}

{[} x\^{}\{2\} + y\^{}\{2\} = 1, {]}

{[} (x-1)\^{}\{2\} + y\^{}\{2\} = 9, {]}

{[} \textbackslash begin\{aligned\} \& (\textbackslash text\{УЧП\}) \&
\textbackslash phi\_\{xx\} + \textbackslash phi\_\{yy\} = 0, \&
\textbackslash text\{внутри \} D, \textbackslash{} \&
(\textbackslash text\{ГУ\}) \& \textbackslash begin\{cases\}
\textbackslash varphi(x, y) = 1 \& \textbackslash text\{на \} x\^{}2 +
y\^{}2 = 1, \textbackslash{} \textbackslash varphi(x, y) = 2 \&
\textbackslash text\{на \} (x - 1)\^{}2 + y\^{}2 = 9.
\textbackslash end\{cases\} \textbackslash end\{aligned\} {]}

\# ПРИНЦИП СУПЕРПОЗИЦИИ--- ОСНОВА ТЕОРИИ ЛИНЕЙНЫХ СИСТЕМ

ЦЕЛЬ ЛЕКЦИИ: Познакомить читателя с принципом суперпозиции и показать,
как с его помощью исходную задачу можно разложить на подзадачи, затем
решить сразу все подзадачи и, сложив полученные результаты, найти
решение исходной задачи. Мы хотим также показать, что два основных
метода решения линейных уравнений---метод разделения переменных и метод
интегральных преобразований---используют принцип суперпозиции.

Предположим, что инженеру нужно найти отклик линейной системы и на
входное воздействие f.~Общий подход к решению задачи таков:

\begin{itemize}
\tightlist
\item
  \begin{enumerate}
  \def\labelenumi{\arabic{enumi})}
  \tightlist
  \item
    разложить f на элементарные воздействия \(f = \sum f_k\) ;
  \end{enumerate}
\item
  \begin{enumerate}
  \def\labelenumi{\arabic{enumi})}
  \setcounter{enumi}{1}
  \tightlist
  \item
    найти отклик системы \(u_k\) на воздействие \(f_k\) ;
  \end{enumerate}
\end{itemize}

\begin{enumerate}
\def\labelenumi{\arabic{enumi})}
\setcounter{enumi}{2}
\tightlist
\item
  сложить отклики \(u = \sum_{i=1}^{n} u_{i}\) .
\end{enumerate}

Если система линейная, то полученная сумма и является откликом системы
на входное воздействие f.~В этом состоит принцип суперпоэиции (см. рнс.
26.1).

\pandocbounded{\includegraphics[keepaspectratio]{images/_page_189_Picture_10.jpeg}}

РИС, 26.1. Основная идея принципа суперпоэнции; L---дифференциальный
оператор; f---вход; u---выход.

Эту основную идею можно использовать при решении сложных смешанных
задач, разбивая их на простые подзадачи, ре-

шая каждую из них и затем суммируя полученные результаты (конечно,
дифференциальное уравнение и граничные условия должны быть линейны).

\section{Разложение смешанной задачи на две более
простые}\label{ux440ux430ux437ux43bux43eux436ux435ux43dux438ux435-ux441ux43cux435ux448ux430ux43dux43dux43eux439-ux437ux430ux434ux430ux447ux438-ux43dux430-ux434ux432ux435-ux431ux43eux43bux435ux435-ux43fux440ux43eux441ux442ux44bux435}

Предположим, нам необходимо решить линейную задачу (обозначим ее P)

(УЧП)

(Р) (ГУ)

Проблема полягає в тому, щоб знайти конформне відображення, яке
перетворює цю область на щось простіше, де нашу проблему легко
розв'язати. У цьому випадку очевидно, що ця проста ділянка має бути
кільцем. Відповідна експозиція

\pandocbounded{\includegraphics[keepaspectratio]{images/_page_359_Figure_10.jpeg}}

{[} \textbackslash begin\{cases\} u(0, t) = 0, \textbackslash{} u(1, t)
= 0, \textbackslash{} u(x, 0) = \textbackslash sin(2\textbackslash pi
x), \textbackslash quad 0 \textbackslash le x \textbackslash le 1.
\textbackslash end\{cases\} {]}

{[} w = 2t \textbackslash left{[} \textbackslash frac\{z-s\}\{z-t\}
\textbackslash right{]}, {]}

Перед нами неоднородное рівняння теплопровідності, так що метод поділу
змінних неприменим. Можпо, конечно, воспользоваться конечным
синус-преобразованием по переменной x или преобразованием Лапласа по
переменной t, но вместо цього мы рассмотрим две подзадачи

{[} u = \textbackslash gamma {[}(x-s)(x-t)-\textbackslash gamma\^{}2{]},
{]}

{[} v = \textbackslash gamma {[}y(x-t){]}, {]}

{[} \textbackslash begin\{cases\} u(0, t) = 0, \textbackslash{} u(1, t)
= 0, \textbackslash{} u(1, t) = 0, \textbackslash end\{cases\} {]}

{[} \textbackslash gamma = 2t/{[}(x-t)\^{}2 + y\^{}2{]}. {]}

\begin{verbatim}
\begin{aligned} \varphi_{uu} + \varphi_{vv} &= 0, & \text{в кольче,} \\ (\Gamma Y) & & & & & & & \\ (\Gamma Y) & & & & & & \\ (\Gamma Y) & & & & & \\ & & & & & \\ (\Gamma Y) & & & & & \\ & & & & \\ & & & & \\ & & & & \\ & & & & \\ & & & \\ & & & \\ & & & \\ & & & \\ & & & \\ & & & \\ & & \\ & & \\ & & \\ & & \\ & & \\ & & \\ & & \\ & & \\ & & \\ & & \\ & & \\ & & \\ & & \\ & & \\ & & \\ & & \\ & & \\ & & \\ & & \\ & & \\ & & \\ & & \\ & & \\ & & \\ & & \\ & & \\ & & \\ & & \\ & & \\ & & \\ & & \\ & & \\ & & \\ & & \\ & & \\ & & \\ & & \\ & & \\ & & \\ & & \\ & & \\ & & \\ & & \\ & & \\ & & \\ & & \\ & & \\ & & \\ & & \\ & & \\ & & \\ & & \\ & & \\ & & \\ & & \\ & & \\ & & \\ & & \\ & & \\ & & \\ & & \\ & & \\ & & \\ & & \\ & & \\ & & \\ & & \\ & & \\ & & \\ & & \\ & & \\ & & \\ & & \\ & & \\ & & \\ & & \\ & & \\ & & \\ & & \\ & & \\ & & \\ & & \\ & & \\ & & \\ & & \\ & & \\ & & \\ & & \\ & & \\ & & \\ & & \\ & & \\ & & \\ & & \\ & & \\ & & \\ & & \\ & & \\ & & \\ & & \\ & & \\ & & \\ & & \\ & & \\ & & \\ & & \\ & & \\ & & \\ & & \\ & & \\ & & \\ & & \\ & & \\ & & \\ & & \\ & & \\ & & \\ & & \\ & & \\ & & \\ & & \\ & & \\ & & \\ & & \\ & & \\ & & \\ & & \\ & & \\ & & \\ & & \\ & & \\ & & \\ & & \\ & & \\ & & \\ & & \\ & & \\ & & \\ & & \\ & & \\ & & \\ & & \\ & & \\ & & \\ & & \\ & & \\ & & \\ & & \\ & & \\ & & \\ & & \\ & & \\ & & \\ & & \\ & & \\ & & \\ & & \\ & & \\ & & \\ & & \\ & & \\ & & \\ & & \\ & & \\ & & \\ & & \\ & & \\ & & \\ & & \\ & & \\ & & \\ & & \\ & & \\ & & \\ & & \\ & & \\ & & \\ & & \\ & & \\ & & \\ & & \\ & & \\ & & \\ & & \\ & & \\ & & \\ & & \\ & & \\ & & \\ & & \\ & & \\ & & \\ & & \\ & & \\ & & \\ & & \\ & & \\ & & \\ & & \\ & & \\ & & \\ & & \\ & & \\ & & \\ & & \\ & & \\ & & \\ & & \\ & & \\ & & \\ & & \\ & & \\ & & \\ & & \\ & & \\ & & \\ & & \\ & & \\ & & \\ & & \\ & & \\ & & \\ & & \\ & & \\ & & \\ & & \\ & & \\ & & \\ & & \\ & & \\ & & \\ & & \\ & & \\ & & \\ & & \\ & & \\ & & \\ & & \\ & & \\ & & \\ & & \\ & & \\ & & \\ & & \\ & & \\ & & \\ & & \\ & & \\ & & \\ & & \\ & & \\ & & \\ & & \\ & & \\ & & \\ & & \\ & & \\ & & \\ & & \\ & & \\ & & \\ & & \\ & & \\ & & \\ & & \\ & & \\ & & \\ & & \\ & & \\ & & \\ & & \\ & & \\ & & \\ & & \\ & & \\ & & \\ & & \\ & & \\ & & \\ & & \\ & & \\ & & \\ & & \\ & & \\ & & \\ & & \\ & & \\ & & \\ & & \\ & & \\ & & \\ & & \\ & & \\ & & \\ & & \\ & & \\ & & \\ & & \\ & & \\ & & \\ & & \\ & & \\ & &
\end{verbatim}

Каждая из цих задач легко решается и в то же время ясно, що
\textbackslash mathit\{suma\} решений задач \(P_1\) и \(P_2\) является
решением исходной вадачи, т. е.

{[} \textbackslash varphi(r) = a \textbackslash ln r + b,
\textbackslash quad r = u\^{}2 + v\^{}2. {]}

Розв'язок задачи \(P_{i}\) Розв'язок задачи \(P_i\) Розв'язок задачи
\(P_i\) Можна показать, що в общем случае розв'язок задачи

{[} \textbackslash varphi(u, v) = 0.57 \textbackslash ln(u\^{}2 +
v\^{}2) + 1 {]}

{[} \textbackslash varphi(x, y) = 0.57 \textbackslash ln(u\^{}2 +
v\^{}2) + 1, {]}

{[} \textbackslash begin\{cases\} u\_t = u\_\{xx\} + f(x, t),
\textbackslash{} u(0, t) = 0, \textbackslash{} u(1, t) = 0,
\textbackslash{} u(x, 0) = \textbackslash varphi(x)
\textbackslash end\{cases\} {]}

{[} w = \textbackslash ln \textbackslash left{[}
\textbackslash frac\{z-1\}\{z+1\} \textbackslash right{]}. {]}

{[} z = re\^{}\{i\textbackslash theta\} {]}

\begin{verbatim}
\begin{aligned} (\text{Y} \forall \Pi) & \qquad \phi_{\text{vx}} + \phi_{yy} = 0, \quad 0 < x < \infty, \quad 0 < y < \infty, \\ \left\{ \begin{aligned} \phi(x, \ 0) = \begin{cases} 1, & 0 < x < 1, \\ 0, & 1 \leqslant x < \infty, \end{cases} \\ \phi(0, \ y) = \begin{cases} 1, & 0 < y < 1, \\ 0, & 1 \leqslant y < \infty. \end{cases} \end{aligned} \end{aligned}
\end{verbatim}

{[} \textbackslash begin\{cases\} u\_t = u\_\{xx\} + f(x, t),
\textbackslash{} u(0, t) = 0, \textbackslash{} u(1, t) = 0,
\textbackslash{} u(x, 0) = 0 \textbackslash end\{cases\} {]}

Зверніть увагу на зображення.

\pandocbounded{
\includegraphics[keepaspectratio]{images/_page_362_Figure_1.jpeg}}

{[} \textbackslash begin\{cases\} u\_t = u\_\{xx\}, \textbackslash{}
u(0, t) = 0, \textbackslash{} u(1, t) = 0, \textbackslash{} u(x, 0) =
\textbackslash varphi(x). \textbackslash end\{cases\} {]}

\pandocbounded{\includegraphics[keepaspectratio]{images/_page_362_Figure_2.jpeg}}

\begin{enumerate}
\def\labelenumi{\arabic{enumi}.}
\setcounter{enumi}{3}
\tightlist
\item
  Розв'язати змішану задачу Діріхле-Неймана з кутовим порошком за
  допомогою розчину під кутом 45°
\end{enumerate}

\pandocbounded{\includegraphics[keepaspectratio]{images/_page_362_Figure_5.jpeg}}

проследим, как оно решается методом конечного синус-преобразования. Вот
что мы должны сделать: разложить еходное воздействие \(f(x, t)\) на
компоненты, найти отклик \(U_n\) на каждую компоненту, а затем сложить
эти отклики. Математически это может быть и не совсем очевидно, но
давайте внимательно проследим за ходом решения. Разложим входящие в
уравнение с частными производными члены в ряды Фурье по синусам:

ПРИМІТКА. Комплексна функція w = \textbackslash ln z = \textbackslash ln
\textbar z\textbar{} + i \textbackslash arg(z) відображає промінь
\textbackslash theta = \(c_1\) у площині z у пряму v = \(c_1\) , а коло
r = \(c_2\) у пряму u = \textbackslash ln \(c_2\) у площині w (див.
рисунки).

\pandocbounded{\includegraphics[keepaspectratio]{images/_page_362_Figure_7.jpeg}}

where

{[} H(x-a) = {]}

{[} \textbackslash begin\{cases\} 0, x \textless{} a, \textbackslash{}
1, x \textbackslash geqslant a, \textbackslash end\{cases\} {]}

Отметим, во-первых, що коефіцієнти \(A_n\) , \(B_n\) , \(F_n\)
действительно являются функциями времени t, поскольку мы применяли
синус-преобразование к функції двух змінних x, t, u, вовторых, що мы
разложили входное воздействие \(f(x,t)\) на простые компоненты
\(F_n\)(t) . Найдем теперь отклики \(U_n\)(t) на каждую момпоненту
\(F_n\)(t) , а затем сложим все отклики \(U_n\)(t) и получим розв'язок
\(u(x,t)\).

Чтобы найти отклики \(U_n\)(t) , необходимо проделать небольшие выкладки
с коэффициентами \(A_n\)(t) и \(B_n\)(t) , чтобы в подынтегральные
вирази входила лише функція u, а не ее производные \(u_i\) и \(u_{xx}\)
. Интегрируя по частям, получаем

Функція Хевісайду,

{[} \textbackslash begin\{split\} A\_n(t) \&=
2\textbackslash int\_0\^{}1 u\_t \textbackslash sin(n\textbackslash pi
x) , dx = \textbackslash frac\{d\}\{dt\} \textbackslash bigg{[}
2\textbackslash int\_0\^{}1 u(x, t)
\textbackslash sin(n\textbackslash pi x) , dx \textbackslash bigg{]} =
\textbackslash frac\{dU\_n(t)\}\{dt\} ,, \textbackslash{} B\_n(t) \&=
2\textbackslash int\_0\^{}1 u\_\{xx\}
\textbackslash sin(n\textbackslash pi x) , dx =
-(n\textbackslash pi)\^{}2 , U\_n(t) ,+ \textbackslash{}
\&\textbackslash quad + 2n\textbackslash pi , \textbackslash big{[} u(0,
t) + (-1)\^{}\{n+1\} , u(1, t) \textbackslash big{]},
\textbackslash end\{split\} {]}

{[} H(x-a) = {]}

{[} \textbackslash begin\{cases\} 0, \& x \textless{} a,
\textbackslash{} 1, \& x \textbackslash geqslant a,
\textbackslash end\{cases\} {]}

{[} \textbackslash begin\{cases\} u(0, t) = 0, \textbackslash{} u(1, t)
= 0 \textbackslash end\{cases\} {]}

в последнее соотношение, получаем

{[} f(x) = \textbackslash mathcal\{F\}\^{}\{-1\}{[}F{]} =
\textbackslash frac\{1\}\{\textbackslash sqrt\{2\textbackslash pi\}\}
\textbackslash int\_\{-\textbackslash infty\}\^{}\{+\textbackslash infty\}
F(\textbackslash omega) e\^{}\{-i\textbackslash omega x\}
d\textbackslash omega \textbackslash qquad F(\textbackslash omega) =
\textbackslash mathcal\{F\}{[}f{]} =
\textbackslash frac\{1\}\{\textbackslash sqrt\{2\textbackslash pi\}\}
\textbackslash int\_\{-\textbackslash infty\}\^{}\{+\textbackslash infty\}
f(x) e\^{}\{-i\textbackslash omega x\} dx {]}

и, отже, (26.1) принимает вид

Поскольку это разложение является тождеством по x, его коэффициенты
должны быть равны нулю, \(\tau\) . е.

.

Таким образом, мы получили связь между входом и выходом системи, между
коэффициентами \(F_n\) и \(U_n\) . Перш ніж чем решать уравнения для
\(U_n\)(t) , обратимся к начальному условию

{[} \textbackslash begin\{cases\} 1, \textbar x\textbar{} \textless{} a
\textbackslash{} 0, \textbar x\textbar{} \textgreater{} a
\textbackslash end\{cases\} {]}

Якщо разложить \(u(x, 0)\) в ряд Фурье по синусам и приравнять его нулю:

то определяются начальные условия для функции \(U_n\left(t\right)\) , т.
е.

, n = 1, 2, \textbackslash dots Мы получили разложение нашей смешанной
задачи на ряд простых задач об отклике системи на элементарные входные
воздействия:

(OДУ)

{[} \textbackslash begin\{cases\} 1, \textbar x\textbar{} \textless{} 1
\textbackslash{} 0, \textbar x\textbar{} \textgreater{} 1
\textbackslash end\{cases\} {]}

(HУ) \(U_n\)(0) = 0, n = 1, 2, \ldots{} Каждое из цих уравнений легко
решается методом интегрирующего множителя (или с помощью преобразования
Лапласа). В любом случае получаем

Таким образом, мы получили отклики \(U_n(t)\) на простейшие входные
воздействия \(F_n(t)\) . Для получения решения исходной задачи осталось
совершить последний шаг---сложить все отклики

Читатель, должно быть, заметил, що каждый отклик \(U_n\)(t) умножается
на весовой множитель \textbackslash sin(n\textbackslash pi x) . Конечно,
в разложении функції \(f(x, t)\) на компоненты \(F_n\)(t) присутствуют
те же весовые множители.

\section{ЗАМЕЧАНИЕ}\label{ux437ux430ux43cux435ux447ux430ux43dux438ux435-1}

Применение конечного синус-преобразования позволяет записывать розв'язок
в виде ряда, по для большинства других інтегральних перетворень
розв'язки представляются в виде интегралов (непрерывное разложение).
Читателю можна предложить дать интерпретацию задачи

{[} \textbackslash delta(x-a) {]}

\(f(x)\) * \(g(x)\)

{[} (1+x\textsuperscript{2)}\{-1\} {]}

{[} xe\textsuperscript{\{-a+x}4\}, a \textgreater{} 0 {]}

\begin{verbatim}
\begin{array}{ll} (\mathrm{Y} + \Pi) & u_{tt} = u_{xx} + \sin{(3\pi x)}, & 0 < x < 1, & 0 < t < \infty, \\ (\mathrm{\Gamma} \mathrm{Y}) & \left\{ \begin{array}{ll} u\left(0, \ t\right) = 0, \\ u\left(1, \ t\right) = 0, \\ \end{array}
\end{verbatim}

\(H(x+a)\) - \(H(x-a)\)

{[} \textbackslash frac\{a\}\{x\textsuperscript{2+a}2\} {]}

{[} \textbackslash frac\{2ax\}\{(x\textsuperscript{2+a}2)\^{}2\} {]}

{[} \textbackslash begin\{array\}\{ll\} u\textbackslash left(x,
~0\textbackslash right) = \textbackslash sin\{(\textbackslash pi x)\},
\textbackslash{} u\_\{t\}\textbackslash left(x, ~0\textbackslash right)
= 0, \textbackslash end\{array\} {]}

Для каждой из подзадач выберите метод розв'язки по своему BKYCY,

\begin{enumerate}
\def\labelenumi{\arabic{enumi}.}
\setcounter{enumi}{2}
\tightlist
\item
  Предположим, що \(u_1\) и \(u_2\) --- розв'язки следующих уравнений:
\end{enumerate}

{[} \textbackslash begin\{cases\} \textbackslash cos(ax),
\textbar x\textbar{} \textgreater{} \textbackslash pi/2a
\textbackslash{} 0, \textbar x\textbar{} \textgreater{}
\textbackslash pi/2a \textbackslash end\{cases\} {]}

\begin{enumerate}
\def\labelenumi{\alph{enumi})}
\tightlist
\item
\end{enumerate}

;\\
b) \(u_t = u_{xx} + e^t\) ;

\begin{enumerate}
\def\labelenumi{\Alph{enumi})}
\setcounter{enumi}{1}
\tightlist
\item
\end{enumerate}

;

{[} F) u\_t = u\_\{xx\} + u\^{}2 . {]}

{[} \textbackslash begin\{cases\} 1-\textbar x\textbar,
\textbar x\textbar{} \textless{} 1 \textbackslash{} 0,
\textbar x\textbar{} \textgreater{} 1 \textbackslash end\{cases\} {]}

\textbackslash mathbf\{\textbackslash Pi\} ля каких уравнений функція
\(u_1\) + \(u_2\) також буде решением? Какие выводы вы сможете сделать,
исходя из ваших ответов?

\begin{enumerate}
\def\labelenumi{\arabic{enumi}.}
\setcounter{enumi}{3}
\tightlist
\item
  Найдите четыре смещанные задачи, сумма решений которых дает розв'язок
  следующей задачи:
\end{enumerate}

{[} \textbackslash begin\{aligned\} u\_t \&= u\_\{xx\} + f(x, t), \& 0
\textless{} x \textless{} 1, \& 0 \textless{} t \textless{}
\textbackslash infty, \textbackslash{} \& \textbackslash begin\{cases\}
u(0, t) = g\_1(t), \textbackslash{} u(1, t) = g\_2(t),
\textbackslash end\{cases\} \& 0 \textless{} t \textless{}
\textbackslash infty, \textbackslash{} \& u(x, 0) =
\textbackslash varphi(x), \& 0 \textbackslash le x \textbackslash le 1.
\textbackslash end\{aligned\} {]}

{[} \textbackslash cos(ax) {]}

{[} \textbackslash frac\{2\}\{\textbackslash pi\}
\textbackslash frac\{\textbackslash sin
a\textbackslash omega\}\{\textbackslash omega\} {]}

\begin{enumerate}
\def\labelenumi{\arabic{enumi}.}
\setcounter{enumi}{4}
\tightlist
\item
  Решите задачу Коши
\end{enumerate}

(OAY)

{[} \textbackslash frac\{2\}\{\textbackslash pi\}
\textbackslash frac\{\textbackslash sin
a\textbackslash omega\}\{(\textbackslash omega\textsuperscript{2+a}2)\^{}2\}
{]}

(HY) \(U_n\)(0) = 0. Можете ли вы проверить полученное розв'язок?

\begin{enumerate}
\def\labelenumi{\arabic{enumi}.}
\setcounter{enumi}{5}
\tightlist
\item
  Предположим, що функції \(u_1\) и \(u_2\) удовлетворяют линейным
  однородным граничным условиям
\end{enumerate}

{[} \textbackslash frac\{\textbackslash pi\}\{2\}
e\^{}\{-a+\textbackslash omega\} {]}

\(u_x\)(1, t) + \(h_2\) \(u(1, t)\) = 0. Буде ли функція \(u_i\) +
\(u_i\) удовлетворять этим условиям? , .

Лекція 27

\section{УРАВНЕНИЯ ПЕРВОГО ПОРЯДКА (МЕТОД
ХАРАКТЕРИСТИК)}\label{ux443ux440ux430ux432ux43dux435ux43dux438ux44f-ux43fux435ux440ux432ux43eux433ux43e-ux43fux43eux440ux44fux434ux43aux430-ux43cux435ux442ux43eux434-ux445ux430ux440ux430ux43aux442ux435ux440ux438ux441ux442ux438ux43a}

ЦЕЛЬ ЛЕКЦИИ: Ввести понятие уравнения з частинними похідними первого
порядка (до сих пор мы рассматривали уравнения второго порядка) и
познакомиться с важным методом розв'язки задач с начальными
условиями---методом характеристик. Задача, которой мы теперь займемся,
имеет вид

(УЧП)

{[} \textbackslash frac\{\textbackslash pi\}\{2\}
e\^{}\{-a+\textbackslash omega\} {]}

{[} , -\textbackslash infty \textless{} x \textless{}
\textbackslash infty , 0 \textless{} t \textless{} \textbackslash infty
, {]}

\begin{verbatim}
(HV)  u(x, 0) = \varphi(x), \quad -\infty < x < \infty.  Сразу же отметим, що впервые мы будем решать задачу с змінними коэффициентами. Оказывается, якщо от координат (x, t) перейти к новым, характеристическим координатам  (s, \tau)  , то наше УЧП превратится в обыкновенное дифференциальное уравнение. Затем можна решить обыкновенное дифференциальное уравнение, т. е. найти  u(s, \tau)  , а на последнем шаге выразить s и t через x и t и получить u(x, t).
\end{verbatim}

Напомним читателю, що, когда мы решали уравнение дифф

{[} \textbackslash frac\{\textbackslash pi\}\{2\}
e\^{}\{-a+\textbackslash omega\} {]}

, -\textbackslash infty \textless{} x \textless{} \textbackslash infty ,
0 \textless{} t \textless{} \textbackslash infty константа
\textbackslash alpha\^{}2 играла роль коэффициента диффузии, а v---была
скоростью среды. Очевидно, що якщо \textbackslash alpha=0 (т. е.
диффузии нет), то розв'язок буде смещаться как целое вдоль оси x со
скоростью v (у нас осталась лише конвекция). Другими словами, якщо
початкова умова \(u(x,0)\)=\textbackslash varphi(x) , то соответствующее
розв'язок уравнения

{[} \textbackslash frac\{\textbackslash pi\}\{2\}
e\^{}\{-a+\textbackslash omega\} {]}

буде иметь вид \(u(x, t)\) = \textbackslash varphi(x-vt) .

Це дает нам основания думать, що розв'язок уравнения первого порядка

{[} \textbackslash frac\{\textbackslash pi\}\{2\}
e\^{}\{-a+\textbackslash omega\} {]}

переносится вдоль потока со скоростью

{[} \textbackslash frac\{2ax\}\{(x\textsuperscript{2+a}2)\^{}2\} {]}

.

Конечно, якщо a и b---константы, то розв'язок представляет собой бегущую
с постоянной скоростью волну. Якщо же \(a(x, t)\) и \(b(x, t)\)

\pandocbounded{\includegraphics[keepaspectratio]{images/_page_196_Figure_10.jpeg}}

\begin{verbatim}
Рис. 27.1. Начальное значение в точке x оказывает влияние на розв'язок лише вдоль характеристики в плоскости змінних (x,t). Параметр s изменяется вдоль характеристик от нуля до бесконечности, начиная с линии начального условия;  \tau  постоянна вдоль каждой характеристики;  \tau  изменяется вдоль линии начального условия.
\end{verbatim}

зависят от x и t, то скорость потока изменяется как вдоль потока, так и
во времени (читатель сможет увидеть, що начальная кривая сильно
искажается). Здесь так много аналогий с конвекцией!

Вернемся к нашей основной задаче:

(УЧП)

{[} \textbackslash frac\{\textbackslash cos(ax), \textbar x\textbar{}
\textgreater{} \textbackslash pi/2a\}\{0, \textbar x\textbar{}
\textgreater{} \textbackslash pi/2a\} {]}

(НУ) \(u(x, 0)\) = \textbackslash varphi(x), -\textbackslash infty
\textless{} x \textless{} \textbackslash infty. Розв'язок цього
линейного уравнения первого порядка основывается на следующем физическом
факте: початкова умова в некоторой точке x переносится в tx-плоскости
вдоль линии, яка

називається характеристи-кой (см. рис. 27.1)

Этим наше уразнение отличается от других уравнений (таких, как рівняння
теплопровідності \(u_t = u_{xx} )\),для которых начальное значение в
точке x оказывает влияние на розв'язок во всех точках пространства и во
все моменты времени. Якщо читатель помнит, то начальное сме- Р И С

\pandocbounded{\includegraphics[keepaspectratio]{images/_page_197_Figure_5.jpeg}}

\begin{verbatim}
. 27.2. Начальное возмущение  u\left(x,0\right)  в точке порождает две хвилі. Начальное возмущение струны согласно волновому уравнению распространяется вдоль двух характеристик.

щение скрипичной струны в точке x оказывает влияние на розв'язок вдоль  \partial syx  линий в tx-плоскости (соответствующих двум бегущим волнам) (см. рис. 27.2).

Естественно, напрашивается идея ввести две новые координаты s и  \tau  (вместо x и t) так, чтобы
\end{verbatim}

\begin{itemize}
\item
  s изменялась вдоль характеристических кривых,
\item
  т изменялась вдоль начальной кривой (желательно вдоль линии t=0).
\end{itemize}

Рассмотрим сначала новую координату s. Якщо s выбрать так, що она
удовлетворяет приведенному выше условию, то уравнение

{[} \textbackslash frac\{1-\textbar x\textbar, \textbar x\textbar{}
\textless{} 1\}\{0, \textbar x\textbar{} \textgreater{} 1\} {]}

превратится в обыкновенное дифференциальное уравнение

{[} \textbackslash cos(ax) {]}

Конечно, остается открытым вопрос, как найти эти характеристики. Ответ
прост: мы выберем характеристики \{{[}\(x(s)\), \(t(s)\){]}: 0
\textless{} s \textless{} \textbackslash infty\} так, чтобы они
удовлетворяли системе

{[} \textbackslash cos(ax) {]}

{[} \textbackslash frac\{\textbackslash pi\}\{2\}
{[}\textbackslash delta(\textbackslash omega+a) +
\textbackslash delta(\textbackslash omega-a){]} {]}

{[} \textbackslash begin\{cases\} \textbackslash frac\{dx\}\{ds\} = a(x,
t), \textbackslash{} \textbackslash frac\{dt\}\{ds\} = b(x, t).
\textbackslash end\{cases\} {]}

{[} f(x) = \textbackslash int\_\{0\}\^{}\{\textbackslash infty\}
F(\textbackslash omega) \textbackslash sin(\textbackslash omega x)
d\textbackslash omega {]}

{[} F(\textbackslash omega) =
\textbackslash frac\{?\}\{\textbackslash pi\}
\textbackslash int\_\{0\}\^{}\{\textbackslash infty\} f(x)
\textbackslash sin(\textbackslash omega x) dx {]}

Ясно, що в цьому случае

{[} 0 \textless{} x \textless{} \textbackslash infty {]}

Другими словами, на характеристиках \{{[}\(x(s)\), \(t(s)\){]}: 0
\textless{} s \textless{} \textbackslash infty\} уравнение з частинними
похідними становится обыкновенным дифференциальным уравнением. Поясним
эти идеи на конкретном примере.

\section{ПРИМЕР}\label{ux43fux440ux438ux43cux435ux440}

Предположим, що нам потрібно решить следующую задачу Коши с постоянными
коэффициентами:

{[} 0 \textless{} \textbackslash omega \textless{} x {]}

f'\,'(0)

{[} \textbackslash begin\{array\}\{lll\} (\textbackslash forall
\textbackslash Pi \textbackslash Pi) \& u\_x + u\_t + 2u = 0, \&
-\textbackslash infty \textless{} x \textless{} \textbackslash infty, \&
0 \textless{} t \textless{} \textbackslash infty \textbackslash{}
(\textbackslash forall \textbackslash Psi) \& u(x, 0) =
\textbackslash sin x, \& -\textbackslash infty \textless{} x \textless{}
\textbackslash infty. \textbackslash end\{array\} {]}

{[} -\textbackslash omega\^{}2 F(\textbackslash omega) +
\textbackslash frac\{2\}\{\textbackslash pi\} \textbackslash omega f(0)
{]}

{[} \textbackslash frac\{1\}\{\textbackslash alpha\}
F\textbackslash left(\textbackslash frac\{\textbackslash omega\}\{\textbackslash alpha\}\textbackslash right)
{]}

{[} \textbackslash frac\{dx\}\{ds\} = 1, \textbackslash quad
\textbackslash frac\{dt\}\{ds\} = 1, \textbackslash quad 0 \textless{} s
\textless{} \textbackslash infty. {]}

{[} \textbackslash frac\{2\textbackslash omega\}\{\textbackslash pi
(\textbackslash alpha\^{}2 + \textbackslash omega\^{}2)\} {]}

{[} x(s) = s + c\_1, \textbackslash quad t(s) = s + c\_2. {]}

{[} x\^{}\{-1/2\} {]}

{[} x(0) = \textbackslash tau, {]}

{[} {[}2/\textbackslash pi\textbackslash omega{]}\^{}\{1/2\} {]}

{[} x = s + \textbackslash tau {]}

{[} \textbackslash frac\{2\}\{\textbackslash pi\textbackslash omega\}
{[}1-\textbackslash cos(\textbackslash omega a){]} {]}

{[} x-t=\textbackslash tau {]}

1

{[} \textbackslash frac\{du\}\{ds\} + 2u = 0, \textbackslash quad 0
\textless{} s \textless{} \textbackslash infty, {]}

{[} e\^{}\{-1\}a + \textbackslash omega {]}

{[} u(0) = \textbackslash sin \textbackslash tau. {]}

{[}
\textbackslash frac\{1\}\{2\}e\^{}\{-\textbackslash alpha\}\textbackslash sin\textbackslash omega
{]}

{[}
u\textbackslash left(s,,\textbackslash tau\textbackslash right)=\textbackslash sin\textbackslash tau
e\^{}\{-2s\} {]}

{[}
\textbackslash frac\{1-e\^{}\{-a\textbackslash omega\}\}\{\textbackslash omega\}
{]}

{[} x = s + \textbackslash tau {]}

{[}
\textbackslash frac\{2\}\{\textbackslash pi\}F'\,'(\textbackslash omega)
{]}

{[} s=t {]}

{[} f(x) = \textbackslash int\_\{0\}\^{}\{\textbackslash infty\}
F(\textbackslash omega) \textbackslash cos(\textbackslash omega x)
d\textbackslash omega \textbackslash qquad \textbackslash qquad F(x) =
\textbackslash frac\{2\}\{\textbackslash pi\}
\textbackslash int\_\{0\}\^{}\{\textbackslash infty\} f(x)
\textbackslash cos(\textbackslash omega x) dx {]}

{[} u(x, t) = \textbackslash sin(x-t)e\^{}\{-xt\} {]}

{[} 0 \textless{} x \textless{} \textbackslash infty
\textbackslash qquad \textbackslash qquad 0 \textless{}
\textbackslash omega \textless{} \textbackslash infty {]}

{[} (\textbackslash forall \textbackslash Pi) \textbackslash quad a(x,
t) u\_x + b(x, t) u\_t + c(x, t) u = 0, \textbackslash quad
-\textbackslash infty \textless{} x \textless{} \textbackslash infty,
\textbackslash{} 0 \textless{} t \textless{} \textbackslash infty, {]}

j'\,'(x)

{[} u(x, 0) = f(x), \textbackslash quad -\textbackslash infty
\textless{} x \textless{} \textbackslash infty. {]}

{[} S\_n = \textbackslash frac\{2\}\{\textbackslash pi\}
\textbackslash int\_0\^{}\{\textbackslash pi\} f(x)
\textbackslash sin(nx) dx. {]}

{[} \textbackslash frac\{dx\}\{ds\} = a(x, t), \textbackslash qquad
\textbackslash frac\{dt\}\{ds\} = b(x, t). {]}

{[} y = \textbackslash frac\{\textbackslash pi\}\{b - a\} (x - a). {]}

{[} x = x (s, \textbackslash tau), {]}

{[} f(x) = \textbackslash sum\_\{n=1\}\^{}\{\textbackslash infty\} S\_n
\textbackslash sin(nx). {]}

{[} t = t (s, \textbackslash tau). {]}

{[} f(x) = \textbackslash sum\_\{n=1\}\^{}\{\textbackslash infty\} S\_n
\textbackslash sin(nx) {]}

{[} \textbackslash frac\{du\}\{ds\} + c{[}x(s, \textbackslash tau), t(s,
\textbackslash tau){]}u = 0, 0 \textless{} s \textless{}
\textbackslash infty, {]}

{[} S\_n = \textbackslash frac\{2\}\{\textbackslash pi\}
\textbackslash int\_\{0\}\^{}\{\textbackslash pi\} f(x)
\textbackslash sin(nx) dx {]}

{[} (\textbackslash text\{УЧП\}) {]}

{[} 0 \textbackslash le x \textbackslash le \textbackslash pi {]}

{[} \textbackslash frac\{dx\}\{ds\} = x {]}

{[} n = 1, 2, \ldots{} {]}

{[} \textbackslash frac\{dt\}\{ds\} = 1 {]}

f'\,'(x)

{[} x = \textbackslash tau e\^{}s, {]}

{[} \textbackslash frac\{\textbackslash pi\}\{2\} e\^{}\{ax\} {]}

{[} t = s. {]}

{[} C\_n = \textbackslash frac\{2\}\{\textbackslash pi\}
\textbackslash int\_0\^{}\{\textbackslash pi\} f(x)
\textbackslash cos(nx) dx. {]}

{[} \textbackslash frac\{du\}\{ds\} + su = 0, \textbackslash quad 0
\textless{} s \textless{} \textbackslash infty, {]}

{[} y = \textbackslash frac\{\textbackslash pi\}\{b-a\}(x-a). {]}

{[} u(0) = F(\textbackslash tau). {]}

{[} f(x) = \textbackslash frac\{C\_0\}\{2\} +
\textbackslash sum\_\{n=1\}\^{}\{\textbackslash infty\} C\_n
\textbackslash cos(nx). {]}

{[} u(s, \textbackslash tau) = F(\textbackslash tau) e\^{}\{-s\}
\textbackslash ell\^{}2 {]}

{[} f(x) = \textbackslash frac\{C\_0\}\{2\} +
\textbackslash sum\_\{n=1\}\^{}\{\textbackslash infty\} C\_n
\textbackslash cos(nx) \textbackslash qquad C\_n =
\textbackslash frac\{2\}\{\textbackslash pi\}
\textbackslash int\_\{0\}\^{}\{\textbackslash pi\} f(x)
\textbackslash cos nx , dx {]}

{[} s=t {]}

{[} 0 \textbackslash le x \textbackslash le \textbackslash pi
\textbackslash qquad n = 0, 1, 2, \textbackslash dots {]}

{[} u(x, t) = F(xe\^{}\{-t\}) e\textsuperscript{\{-t}2/2\} {]}

f'\,'(x)

{[} u(x, t) = \textbackslash sin(xe\^{}\{-t\})
e\textsuperscript{\{-t}2/2\} {]}

{[} -n\^{}\{2\}C\_\{n\} - \textbackslash frac\{2\}\{\textbackslash pi\}
\{f'(0) - (-1)\^{}\{n\} f'(n)\} {]}

{[} u\_x + u\_t = 0 {]}

{[} \textbackslash begin\{array\}\{ll\} ( \textbackslash forall
\textbackslash forall \textbackslash Pi ) \& xu\_x + tu\_t + 2u = 0,
\textbackslash quad -\textbackslash infty \textless{} x \textless{}
\textbackslash infty, \textbackslash quad 1 \textless{} t \textless{}
\textbackslash infty, \textbackslash{} ( H \textbackslash forall ) \& u
\textbackslash left( x, ~1 \textbackslash right) = \textbackslash sin x,
\textbackslash quad -\textbackslash infty \textless{} x \textless{}
\textbackslash infty. \textbackslash end\{array\} {]}

{[} \textbackslash begin\{cases\} 2, \&n=0 \textbackslash{} 0, \&n=1, 2,
\textbackslash dots \textbackslash end\{cases\} {]}

{[} \textbackslash begin\{array\}\{ll\}
(\textbackslash mathrm\{Y\}\textbackslash mathrm\{H\}\textbackslash Pi)
\& au\_x + bu\_y + cu\_t + du = 0, \& -\textbackslash infty \textless{}
x \textless{} \textbackslash infty, \textbackslash{} \&
-\textbackslash infty \textless{} y \textless{} \textbackslash infty,
\textbackslash{} \& 0 \textless{} t \textless{} \textbackslash infty,
\textbackslash end\{array\} {]}

где a, b, c и d-заданные константы. Ответ проверьте.

\begin{enumerate}
\def\labelenumi{\arabic{enumi}.}
\setcounter{enumi}{3}
\tightlist
\item
  Решите задачу
\end{enumerate}

{[} \textbackslash begin\{cases\} 1, \& n=m \textbackslash{} 0, \&
n\textbackslash neq m \textbackslash end\{cases\} {]}

{[} \textbackslash begin\{aligned\} u\_x + u\_t + tu = 0, \&
-\textbackslash infty \textless{} x \textless{} \textbackslash infty, \&
0 \textless{} t \textless{} \textbackslash infty \textbackslash{} u(x,
0) = F(x), \& -\textbackslash infty \textless{} x \textless{}
\textbackslash infty, \& \textbackslash end\{aligned\} {]}

\begin{enumerate}
\def\labelenumi{\arabic{enumi}.}
\setcounter{enumi}{5}
\tightlist
\item
\end{enumerate}

{[} (\textbackslash text\{УЧП\}) \textbackslash quad u\_x + 2u\_y + 2u =
0, \textbackslash quad -\textbackslash infty \textless{} x \textless{}
\textbackslash infty, \textbackslash quad -\textbackslash infty
\textless{} y \textless{} \textbackslash infty. {]}

{[} u\_i + g(u)u\_x = 0 {]}

\begin{verbatim}
\begin{cases}
\end{verbatim}

ТАБЛИЦЯ F. Перетворення Лапласа

{[} u(x, 0) = \textbackslash varphi(x) {]}

\textbar----------------------------------\textbar{}

1/с, с \textgreater{} 0 \textbar{}

{[} u\_t + f\_x = 0. {]}

{[} \textbackslash int\_\{a\}\^{}\{b\}
\textbackslash frac\{\textbackslash partial u\}\{\textbackslash partial
t\}(x, t) dx = -\textbackslash int\_\{a\}\^{}\{b\}
\textbackslash frac\{\textbackslash partial f\}\{\textbackslash partial
x\}(x, t) dx. {]}

{[} u\_t + f\_x = 0. {]}

{[} f(u) = Au (1-u). {]}

{[} u\_t + \textbackslash frac\{df\}\{du\} u\_x = u\_t + g(u) u\_x = 0.
{]}

{[} u\_t + 2uu\_x = 0. {]}

{[} \textbackslash begin\{aligned\}
(\textbackslash mathsf\{Y\}\textbackslash mathsf\{P\}\textbackslash Pi)
\& u\_t + g\textbackslash left(u\textbackslash right)u\_x = 0, \&
-\textbackslash infty \textless{} x \textless{} \textbackslash infty, \&
0 \textless{} t \textless{} \textbackslash infty, \textbackslash{}
(\textbackslash mathsf\{H\}\textbackslash mathsf\{Y\}) \&
u\textbackslash left(x, ~0\textbackslash right) =
\textbackslash varphi\textbackslash left(x\textbackslash right), \&
-\textbackslash infty \textless{} x \textless{} \textbackslash infty.
\textbackslash end\{aligned\} {]}

F(s-a) \textbar{}

F(s) G(s) \textbar{}

Вспомним предыдущую лекцию, где мы изучали уравнение конвективного
переноса:

\begin{verbatim}
| 15. f^{(n)}(t) , (n-похідна) | s^{n}F(s) - s^{n-1}f(0) - \dots - f^{(n-1)}(0) |
\end{verbatim}

\textbar------------------------------------------------------------------------\textbar--------------------------------------------------\textbar{}

\begin{verbatim}
| 16. ж. (at) | \frac{1}{a}F\left(\frac{s}{a}\right),  a>0 |.
\end{verbatim}

В этом уравнении \(u(x, t)\)---это концентрация вещества в потоке,
который движется со скоростью v. Можно провести аналогию между
нелинейным уравнением

и уравнением переноса. Будем считать, что частина воды в гочке \(x_0\)
движется со скоростью \(g(u)\) (вдоль по течению или против него). Тогда
через t секунд координата частины будет

\begin{verbatim}
| 19. Ерік (A/2 <b>V</b> <del>T</del> ) | \frac{1}{s}e^{-a} |
\end{verbatim}

(уравнение характеристик).

Вспомним, что концентрация \(u(x, t)\) не меняется вдоль характеристики.
Если нам известна начальная концентрация u(\(x_a\), 0) ,

\pandocbounded{\includegraphics[keepaspectratio]{images/_page_206_Figure_8.jpeg}}

Рис. 28.2. Характеристики уравнения \(u_t\)+\(g (u)\) \(u_x=0\);
a---общий вид характеристики, выходящей из т. ( \(x_0\) , 0), вдоль
характеристики решение остается постоянным и равным величине u ( \(x_0\)
, 0); b---пример характеристики, начинающейся из т. (2, 0), уравнения
\(u_t\)+3uu\_x=0 с \(u (2, 0)\)=1.

то уравнение характеристик принимает вид

x = \(x_0\) + g{[}u(\(x_0\), 0){]} t (уравнение характеристики,
начинающейся в точке (\(x_0\), 0) ).

Пусть, например, мы решаем задачу

(УЧП),

(НУ) \(u(x, 0)\) =

\begin{verbatim}
\begin{cases} 0, & x < 0, \\ 1, & x \ge 0. \end{cases} \end{cases}
\end{verbatim}

Чтобы найти характеристику, выходящую из точки, например (2, 0), мы
должны написать

MATHBLK0741X

MATHBLK0742X

Можно утверждать, что решение нашей задачи \(u(x, t)\) равно единице на
линин x=2+3t, т. е. \(u(2+3t, t)\)=1 при всех t\textgreater0. На рис.
28.2 изображены характеристики для общего случая и для этой конкретной
задачи.

Теперь ясно, что, зная характеристики в каждой точке и зная, что решение
вдоль характеристик не меняется, можно получить решение
\(u\left(x,\ t\right)\) в любой момент времени t. Мы не получили в явном
виде выражение для \(u\left(x,\ t\right)\) через переменые x и t, но
знание характеристик уравнения позволяет решить некоторые интересные
задачи.

\section{Задача дорожного
движения}\label{ux437ux430ux434ux430ux447ux430-ux434ux43eux440ux43eux436ux43dux43eux433ux43e-ux434ux432ux438ux436ux435ux43dux438ux44f}

Займемся моделью дорожного движения, в которой поток вычисляется по
формуле \(f(u)=u^2\) , а начальная плотность изо- Р И

MATHBLK0743X

С. 28.3. Начальная плотность автомобилей движущихся слева направо.

бражена на рис. 28.3. Другими словами, нам нужно решить задачу

MATHBLK0744X

Начнем с определения характеристик, выходящих из точек вида \((x_0, 0)\)
. Для \(x_0 < 0\) получаем

MATHBLK0745X

Эти прямые линии изображены на рис. 28.4. Теперь рассмотрим точки
\(0 < x_0 < 1\) . Здесь характеристики определяются следующими
соотношениями:

MATHBLK0746X

Разремая последнее соотношение относительно t, получаем

MATHBLK0747X

\section{Декартова система
координат}\label{ux434ux435ux43aux430ux440ux442ux43eux432ux430-ux441ux438ux441ux442ux435ux43cux430-ux43aux43eux43eux440ux434ux438ux43dux430ux442}

MATHBLK0748X

где \(u_R\) и \(u_L\) ---значения решения справа и слева от волнового
фронта, а \(f(u_R)\) и \(f(u_L)\) ---значения потока при этих
плотностях.

MATHBLK0749X

РИС. 28.5. Плотность дорожного движения в различные моменты времени.

В нашем примере скорость распространения разрыва оказывается равной
единице. В самом деле, если вспомнить, что \(f(u) = u^2\) , то сразу
получаем

MATHBLK0750X

\section{Циліндрична система
координат}\label{ux446ux438ux43bux456ux43dux434ux440ux438ux447ux43dux430-ux441ux438ux441ux442ux435ux43cux430-ux43aux43eux43eux440ux434ux438ux43dux430ux442}

MATHBLK0751X

изображено на рис. 28.5.

\section{ЗАМЕЧАНИЯ}\label{ux437ux430ux43cux435ux447ux430ux43dux438ux44f-9}

\begin{enumerate}
\def\labelenumi{\arabic{enumi}.}
\item
  Ударная волна в нашем примере возникает благодаря тому, что поток
  быстро растет с ростом плотности u. Если бы поток задавался формулой
  \(f(u)\)=u, то уравнение имело бы вид \(u_1+u_x=0\) , а решение этого
  уравнения представляет волну, движущуюся вправо без искажений. Если на
  мгновение задуматься над смыслом формулы \(f(u)\)=u, то становится
  очевидным, что решение будет двигаться именно таким образом.
\item
  Прямой подстановкой решения в виде
\end{enumerate}

MATHBLK0752X

\begin{verbatim}
\end{aligned}
\end{verbatim}

\section{Сферична система
координат}\label{ux441ux444ux435ux440ux438ux447ux43dux430-ux441ux438ux441ux442ux435ux43cux430-ux43aux43eux43eux440ux434ux438ux43dux430ux442}

MATHBLK0753X

{[} \textbackslash begin\{array\}\{ll\} u\_t +
g\textbackslash left(u\textbackslash right) u\_x = 0, \&
-\textbackslash infty \textless{} x \textless{} \textbackslash infty, \&
0 \textless{} t \textless{} \textbackslash infty, \textbackslash{}
u\textbackslash left(x, ~0\textbackslash right) =
\textbackslash varphi\textbackslash left(x\textbackslash right), \&
-\textbackslash infty \textless{} x \textless{} \textbackslash infty.
\textbackslash end\{array\} {]}

MATHBLK0754X

Рівняння Лапласа Рівняння Гельмгольца Рівняння Пуассона-Шредінгера

\section{Рівняння гіперболічного
типу}\label{ux440ux456ux432ux43dux44fux43dux43dux44f-ux433ux456ux43fux435ux440ux431ux43eux43bux456ux447ux43dux43eux433ux43e-ux442ux438ux43fux443}

MATHBLK0755X

В нашем частном случае можно получить решение в явном виде (во многих
других случаях это сделать не удается)

MATHBLK0756X

Проверьте, что эта функция действительно является решением нашей задачи.

\section{\texorpdfstring{\textbf{ЗАДАЧИ}}{ЗАДАЧИ}}\label{ux437ux430ux434ux430ux447ux438-9}

\begin{enumerate}
\def\labelenumi{\arabic{enumi}.}
\tightlist
\item
  Решите задачу Коши
\end{enumerate}

MATHBLK0757X

Какова здесь связь между потоком и плотностью? Совнадает ли поведение
полученного решения с тем, что вы ожидали? Сравните решения задачи 1 и
задачи 2.

\begin{enumerate}
\def\labelenumi{\arabic{enumi}.}
\setcounter{enumi}{2}
\tightlist
\item
  Предположим, что невязкая жидкость течет по трубе и при этом
  просачивается через стенки по закону \(F(u)\) = ku (г/см·с). Закон
  сохранения (скорее закон изменения) в этом случае записывается в виде
\end{enumerate}

MATHBLK0758X.

Найдите решение этого уравнения, если \(f(u)\) = u и \(\varphi(x) = u\)
(x, 0) (начальное распределение задано в общем виде). Дайте
интерпретацию полученного решения.

\begin{enumerate}
\def\labelenumi{\arabic{enumi}.}
\setcounter{enumi}{3}
\item
  Найдите решение предыдущего уравнения, если потери определяются
  функцией \(F(x, t)\) = 1/x. Проверьте полученное решение. Ясен ли вам
  его физический смысл?
\item
  Проверьте, что функция и, неявно заданная соотношением
\end{enumerate}

MATHBLK0759X

является решением нелинейной задачи

MATHBLK1240X

MATHBLK1241X

\section{СИСТЕМЫ УРАВНЕНИЙ С ЧАСТНЫМИ
ПРОИЗВОДНЫМИ}\label{ux441ux438ux441ux442ux435ux43cux44b-ux443ux440ux430ux432ux43dux435ux43dux438ux439-ux441-ux447ux430ux441ux442ux43dux44bux43cux438-ux43fux440ux43eux438ux437ux432ux43eux434ux43dux44bux43cux438}

ЦЕЛЬ ЛЕКЦИИ: Познакомить читателя с тем фактом, что многие физические
системы (особенно в динамике жидкости) нельзя описать одним уравнением.
Для описания таких систем требуются системы связанных уравнений. В эти
системы входят некоторые неизвестные функции, такие, как давление
\(p(x, y, z, t)\), плотность \(\rho(x, y, z, t)\) , температура
\(T(x, y, z, t)\) и их частные производные. Связи между этими функциями
определяются физическими законами и наша задача --- найти все эти
функции одновременно.

В лекции будет показано, как можно решить линейную систему уравнений

MATHBLK1242X

если преобразовать ее в новую систему независимых уравнений

MATHBLK1243X

а затем независимо решить каждое уравнение этой системы.

Во многих областях науки встречаются системы величин, которые связаны
между собой не одним, а сразу несколькими соотношениями. Например, в
динамике жидкостей---это четыре уравнения

MATHBLK1244X

(сохранение импульса вдоль оси \(x\) ), (29.1)
\(v_t + uv_x + vv_y + \frac{1}{\rho} p_y = 0\) (сохранение импульса
вдоль оси \(y\) ), (сохранение массы), (сохранение массы), (сохранение
массы),
\(\left(\frac{p}{\rho^{\gamma}}\right)_t + u\left(\frac{p}{\rho^{\gamma}}\right)_x + v\left(\frac{p}{\rho^{\gamma}}\right)_y = 0\)
(сохранение энергии), \(p(x, y, z, t)\) --- давление в жидкости,
\(u(x, y, z, t)\) --- \(x\) -компонента скорости, \(v(x, y, z, t)\) ---
\(y\) -компонента скорости, \(\rho(x, y, z, t)\) --- плотность жидкости.

Приведенная нелинейная система уравнений называется уравнениями Эйлера
движения жидкости. Задача состоит в одновременном отыскании всех
неизвестных функций р, и, v и р, которые удовлетворяли бы всем четырем
уравнениям (вместе с начальными и граничными условиями).

Есть и другие причины для изучения систем уравнений. В теории
обыкновенных дифференциальных уравнений (если читатель помнит) показано,
что одно уравнение второго порядка можно представить в виде системы двух
уравнений первого порядка. И хотя в теории уравнений с частными
производными все не так просто, очень часто одно уравнение с частными
производными высокого порядка можно свести к системе уравнений с
частными производными первого порядка. Например, телеграфное уравнение

MATHBLK1245X

путем введения трех новых переменных

MATHBLK1246X

MATHBLK1247X

MATHBLK1248X

можно свести к системе уравнений первого порядка

(29.2)

MATHBLK1249X

MATHBLK1250X

{[} \textbackslash begin\{cases\}
\textbackslash frac\{\textbackslash partial
u\_1\}\{\textbackslash partial x\} = u\_2, \textbackslash{}
\textbackslash frac\{\textbackslash partial
a\_1\}\{\textbackslash partial t\} = u\_3, \textbackslash{}
\textbackslash frac\{\textbackslash partial
u\_3\}\{\textbackslash partial t\} = c\^{}2
\textbackslash frac\{\textbackslash partial
u\_2\}\{\textbackslash partial x\} + au\_3 + bu\_4.
\textbackslash end\{cases\} {]}

MATHBLK1251X

MATHBLK1252X

Для того чтобы решать системы типа (29.2), необходимо хорошо знать такие
понятия, как собственные значения и собственные векторы матрицы. По эфой
причине мы предлагаем вниманию читателя краткий обзор линейной алгебры.

\section{Обзор линейной
алгебры}\label{ux43eux431ux437ux43eux440-ux43bux438ux43dux435ux439ux43dux43eux439-ux430ux43bux433ux435ux431ux440ux44b}

Общий случай

Частный случай

Определение. Если А---квадратная матрица порядка n, то собственными
значениями матрицы называются n корней характе-

Пусть

MATHBLK1253X

MATHBLK1254X

MATHBLK1255X.

ристического уравнения

MATHBLK1256X

где \(\det(A-\lambda I)\) --- определитель матрицы \(A-\lambda I\) .

Некоторые собственные значения могут совпадать.

Тогда собственные значения являются корнями уравнения

MATHBLK1257X

MATHBLK1258X

MATHBLK1259X

MATHBLK1260X

Откуда

MATHBLK1261X, \(\lambda_2 = 3\) .

Определение. Если \(\lambda\) ---собственное значение матрицы A, то
соответствующим ему собственным вектором называется такой ненулевой
вектор, который удовлетворяет соотношению

MATHBLK1262X.

Пример. Вектор X =

{[} \textbackslash begin\{bmatrix\} 1 \textbackslash{} 2
\textbackslash end\{bmatrix\} {]}

\(является собственным вектором, соответствующим собственному значению\)
\textbackslash lambda\_2 = 3 , так как

MATHBLK1263X

MATHBLK1264X

MATHBLK1265X

MATHBLK1266X

MATHBLK1267X

MATHBLK1268X

Определение. Матрица \(A^{-1}\) называется обратной к матрице A, если

MATHBLK1269X

Пример, Обратной к

MATHBLK1270X

MATHBLK1271X

MATHBLK1272X

является матрица

MATHBLK1273X

MATHBLK1274X

MATHBLK1275X

где 1-единичная матрица.

поскольку

MATHBLK1276X

MATHBLK1277X

MATHBLK1278X

\section{Диагонализация
матрицы}\label{ux434ux438ux430ux433ux43eux43dux430ux43bux438ux437ux430ux446ux438ux44f-ux43cux430ux442ux440ux438ux446ux44b}

Пусть A --- квадратная матрина порядка n и пусть все ее собственные
значения \(\lambda_1\) , \(\lambda_2\) \ldots{} \(\lambda_n\) различны.
Составим матрицу P так, чтобы у нее \emph{столбцами} были координаты
собственных векторов (в k-м столбце стоит собственный вектор \(X_k\) ,
соответствующий собственному значению \(\lambda_k\) ), т. е.

MATHBLK1279X

MATHBLK1280X

{[} \textbackslash begin\{bmatrix\} X\_1 \&
X\_\{\textbackslash tilde\{2\}\} \& \textbackslash dots \& X\_n
\textbackslash{} X\_\{1\} \& X\_\{2\} \& \textbackslash dots \& X\_n
\textbackslash end\{bmatrix\} {]}

MATHBLK1281X

MATHBLK1282X

MATHBLK1283X

{[} \textbackslash begin\{bmatrix\} x\_\{11\} \&
x\_\{1\textbackslash tilde\{2\}\} \& \textbackslash dots \& x\_\{1n\}
\textbackslash{} x\_\{21\} \& x\_\{22\} \& \textbackslash dots \&
x\_\{2n\} \textbackslash{} \textbackslash vdots \& \textbackslash vdots
\& \textbackslash ddots \& \textbackslash vdots \textbackslash{}
x\_\{n1\} \& x\_\{n\textbackslash tilde\{3\}\} \& \textbackslash dots \&
x\_\{nn\} \textbackslash end\{bmatrix\} {]}

MATHBLK1284X

MATHBLK1285X

Тогда матрица Л, определяемая соотношением

MATHBLK1286X

где \(P^{-1}\) --- обратная к матрице P, будет диагональной матрицей
вида

MATHBLK1287X

MATHBLK1288X

{[} \textbackslash begin\{bmatrix\} \textbackslash lambda\_1 \& \& \&
\textbackslash{} \& \textbackslash lambda\_2 \& 0 \& \textbackslash{} \&
\& \textbackslash ddots \& \textbackslash{} \& \& \&
\textbackslash lambda\_n \textbackslash end\{bmatrix\} {]}

MATHBLK1289X

MATHBLK1290X

Например, у матрицы

MATHBLK1291X

MATHBLK1292X

MATHBLK1293X

собственные числа равны \(\lambda_ 1 = 4\) и \(\lambda_2 = -4\) , а
собственные векторы соответственно

MATHBLK1294X

MATHBLK1295X

MATHBLK1296X

in X\_2 =

{[} \textbackslash begin\{bmatrix\} -2 \textbackslash{} 1
\textbackslash end\{bmatrix\} {]}

.

Значит, произведение

MATHBLK1297X

MATHBLK1298X

MATHBLK1299X

MATHBLK1300X

MATHBLK1301X

MATHBLK1302X

MATHBLK1303X

MATHBLK1304X

MATHBLK1305X

MATHBLK1306X

равно диагональной матрице

MATHBLK1307X

MATHBLK1308X

MATHBLK1309X

(Читатель может проверить это самостоятельно.)

Теперь можно решить простую систему двух уравнений с частными
производными (с соответствующими начальными условиями).

\section{\texorpdfstring{Решение линейной системы \(u_t + Au_x = 0\)
.}{Решение линейной системы u\_t + Au\_x = 0 .}}\label{ux440ux435ux448ux435ux43dux438ux435-ux43bux438ux43dux435ux439ux43dux43eux439-ux441ux438ux441ux442ux435ux43cux44b-u_t-au_x-0-.}

Рассмотрим задачу Коши для системы, содержащей два уравнения и два
начальных условия:

(УЧП 1)

MATHBLK1310X,\\
(29.3) (УЧП 2)
\(\frac{\partial u_2}{\partial t} + 2 \frac{\partial u_1}{\partial x} = 0\)
, \(-\infty < x < \infty\) , \(0 < t < \infty\) ,\\
(HУ 1) \(u_1(x, 0) = f(x)\) ,\\
(HУ 2) \(u_2(x, 0) = g(x)\) , \(-\infty < x < \infty\) .

Эта задача может соответствовать определению давления \(u_1(x, t)\) и
плотности \(u_2(x, t)\) как функции пространственной координаты x и
времени t по известным распределениям этих величин в начальный момент.

Запишем систему уравнений в матричной форме

MATHBLK1311X

MATHBLK1312X

{[} \textbackslash begin\{bmatrix\}
\textbackslash frac\{\textbackslash partial
u\_1\}\{\textbackslash partial t\} \textbackslash{}
\textbackslash frac\{\textbackslash partial
u\_2\}\{\textbackslash partial t\} \textbackslash end\{bmatrix\} {]}

MATHBLK1313X

\begin{itemize}
\tightlist
\item
\end{itemize}

MATHBLK1314X

MATHBLK1315X

{[} \textbackslash begin\{bmatrix\}
\textbackslash frac\{\textbackslash partial
u\_1\}\{\textbackslash partial x\} \textbackslash{}
\textbackslash frac\{\textbackslash partial
u\_2\}\{\textbackslash partial x\} \textbackslash end\{bmatrix\} {]}

MATHBLK1316X

MATHBLK1317X

MATHBLK1318X

MATHBLK1319X

или \((29 \ 4)\) MATHBLK1320X

где

MATHBLK1321X

MATHBLK1322X

MATHBLK1323X

MATHBLK1324X

{[} \textbackslash begin\{bmatrix\}
\textbackslash frac\{\textbackslash partial
u\_1\}\{\textbackslash partial t\} \textbackslash{}
\textbackslash frac\{\textbackslash partial
u\_2\}\{\textbackslash partial t\} \textbackslash end\{bmatrix\} {]}

MATHBLK1325X

MATHBLK1326X

MATHBLK1327X

{[} \textbackslash begin\{bmatrix\}
\textbackslash frac\{\textbackslash partial
u\_1\}\{\textbackslash partial x\} \textbackslash{}
\textbackslash frac\{\textbackslash partial
u\_2\}\{\textbackslash partial x\} \textbackslash end\{bmatrix\} {]}

MATHBLK1328X

MATHBLK1329X

MATHBLK1330X

MATHBLK1331X

Введем новые неизвестные функции

MATHBLK1332X

MATHBLK1333X

{[} \textbackslash begin\{bmatrix\} v\_1 \textbackslash{} v\_2
\textbackslash end\{bmatrix\} {]}

MATHBLK1334X

MATHBLK1335X

с помощью преобразования

MATHBLK1336X.

где P---матрица, по столбцам которой стоят собственные векторы матрицы A
(матрица P нам уже известна). Оказывается, что после такого
преобразования для определения v получится очень простая система (два
уравнения относительно новых неизвестных \(v_1\) и \(v_2\) оказываются
независимыми). Значит, \(v_1\) и \(v_2\) легко находятся. После того как
\(v_1\) и \(v_2\) найдены, по формуле

MATHBLK1337X

находятся искомые функции и, и и2.

Сначала, однако, выясним, как выглядит система для определения v. Чтобы
разобраться в этом, читатель должен обратиться к матричному исчислению.
Продифференцировав обе части соотношения u = Pv, получаем

(29 5)

MATHBLK1338X

MATHBLK1339X

(Чтобы лучше понять дальнейшее изложение, читателю рекомендуется
записать эти уравнен я в развернутой форме.)

Теперь подставим соотношения (29.5) в систему

MATHBLK1340X

и получим

MATHBLK1341X

Если умножить теперь обе части последнего уравнения на \(P^{-1}\) , то
получим

MATHBLK1342X

или

MATHBLK1343X

Раньше мы уже видели, что для нашей матрицы

MATHBLK1344X

MATHBLK1345X

MATHBLK1346X

Запишем (29.6) в развернутой форме

(29.7)

MATHBLK1347X

Получилась система из двух несвязанных уравнений, которые можно решать
независимо. Читателю должно быть известно, что решениями этих уравнений
являются бегущие волны

MATHBLK1348X \(v_2(x, t) = \psi(x+4t),\) где ф и ф-произвольные
дифференцируемые функции.

Таким образом, мы нашли функции \(v_1\) и \(v_2\) . Для получения общего
решения u нужно вычислить

MATHBLK1349X

MATHBLK1350X

MATHBLK1351X

{[} \textbackslash begin\{bmatrix\} v\_1 \textbackslash{} v\_2
\textbackslash end\{bmatrix\} {]}

MATHBLK1352X

MATHBLK1353X

MATHBLK1354X

MATHBLK1355X

MATHBLK1356X

MATHBLK1357X

MATHBLK1358X

MATHBLK1359X

MATHBLK1360X

Другими словами,

(29.8)

MATHBLK1361X

Следовательно, мы нашли общее решение системы двух уравнений (29.3).

Типичным решением (а их бесконечное число) могут быть, например, функции

MATHBLK1362X

(две произвольные функции)

и, следовательно,

MATHBLK1363X \(u_2(x, t) = \sin(x-4t) + (x+4t)^2.\) Подставим теперь
общее решение (29.8) в начальные условия

MATHBLK1364X \(u_2(x, 0) = g(x),\) в результате получаем

MATHBLK1365X

MATHBLK1366X

откуда функци і ф и ф легко находятся:

MATHBLK1367X

MATHBLK1368X

\textbf{Тепе}рь можно написать решение задачи (29.3) в окончательном
виде

MATHBLK1369X

MATHBLK1370X

(29.9)

MATHBLK1371X

MATHBLK1372X

\section{ЗАМЕЧАНИЯ}\label{ux437ux430ux43cux435ux447ux430ux43dux438ux44f-10}

\begin{enumerate}
\def\labelenumi{\arabic{enumi}.}
\item
  Многие численные методы ориентированы на решение систем уравнений и,
  значит, многне программы для ЭВМ написаны так, чтобы решать системы
  уравнений первого порядка. По этой причине, прежде чем воспользоваться
  программои, приходится преобразовать свое уравнение высокого порядка в
  систему уравнений первого порядка.
\item
  Линейную систему
\end{enumerate}

MATHBLK1373X

можно решать тем же методом, что и систему \(u_t + Au_x = 0\) . Разница
в том, что в этом случае матрица P из собственных векторов матрицы
\(A(x, t)\) будет функцией x и t.

\section{ЗАДАЧИ}\label{ux437ux430ux434ux430ux447ux438-10}

\begin{enumerate}
\def\labelenumi{\arabic{enumi}.}
\tightlist
\item
  Запишите систему уравнений (29.2) в матричной форме
\end{enumerate}

MATHBLK1374X

где A, B и C-матрицы третьего порядка.

\begin{enumerate}
\def\labelenumi{\arabic{enumi}.}
\setcounter{enumi}{1}
\tightlist
\item
  Найдите собственные векторы и собственные значения матрицы
\end{enumerate}

MATHBLK1375X

MATHBLK1376X.

MATHBLK1377X

MATHBLK1378X

{[} \textbackslash begin\{cases\}
\textbackslash frac\{\textbackslash partial
u\_1\}\{\textbackslash partial t\} +
\textbackslash frac\{\textbackslash partial
u\_1\}\{\textbackslash partial x\} +
\textbackslash frac\{\textbackslash partial
u\_2\}\{\textbackslash partial x\} = 0, \textbackslash{}
\textbackslash frac\{\textbackslash partial
u\_2\}\{\textbackslash partial t\} + 4
\textbackslash frac\{\textbackslash partial
u\_1\}\{\textbackslash partial x\} +
\textbackslash frac\{\textbackslash partial
u\_2\}\{\textbackslash partial x\} = 0. \textbackslash end\{cases\} {]}

MATHBLK1379X

MATHBLK1380X

УКАЗАНИЕ. Запишите сначала систему в матричной форме.

\begin{enumerate}
\def\labelenumi{\arabic{enumi}.}
\setcounter{enumi}{3}
\item
  Проверьте формулы (29.5), записав их в скалярной форме.
\item
  Проверьте, что функции \(u_1\) и \(u_2\) из (29.8) удовлетворяют обоим
  уравнениям (29.3).
\end{enumerate}

\section{КОЛЕБАНИЯ МЕМБРАНЫ (ВОЛНОВОЕ УРАВНЕНИЕ В ПОЛЯРНЫХ
КООРДИНАТАХ)}\label{ux43aux43eux43bux435ux431ux430ux43dux438ux44f-ux43cux435ux43cux431ux440ux430ux43dux44b-ux432ux43eux43bux43dux43eux432ux43eux435-ux443ux440ux430ux432ux43dux435ux43dux438ux435-ux432-ux43fux43eux43bux44fux440ux43dux44bux445-ux43aux43eux43eux440ux434ux438ux43dux430ux442ux430ux445}

ЦЕЛЬ ЛЕКЦИИ: Показать, что колебания мембраны можно описать волновым
уравнением

MATHBLK1381X.

Если искать решение этого уравнения в виде стоячих волн:

MATHBLK1382X

то оно распадается на три обыкновенных дифференциальных уравнения:

MATHBLK1383X

(уравнение гармонических колебаний),
\(r^2 R'' + rR' + (\lambda^2 r^2 - n^2) R = 0\) (уравнение Бесселя),
(уравнение гармонических колебаний).

Все произведения \(R(r) \Theta(\theta) T(t)\) решений этих трех
обыкновенных дифференциальных уравнений описывают основные колебания
мембраны, а функции \(R(r) \Theta(\theta)\) описывают форму мембраны.
Чтобы найти колебания мембраны при произвольных начальных условиях, мы
составим такую комбинацию фундаментальных решений, чтобы начальные
условия удовлетворялись.

Цель этой лекции ясна---найти колебания круглой мембраны с заданными
начальными и граничными условиями. Для простоты будем считать, что
радиус мембраны равен единице, а смещение по границе равно нулю. Пусть
\(u\left(r,\theta,t\right)\) обозначает возвышение точек мембраны под
плоскостью равновесия. Нам необходимо решить задачу (см. рис. 30.1)

MATHBLK1384X

MATHBLK1385X

MATHBLK1386X

Чтобы решить эту задачу, вспомним задачу о колебаниях скрипичной струны
из лекции 20, которая решалась суперпозицией бесконечного числа простых
колебаний. Если воспользовать-

ся этим же подходом, то решения следует искать в виде

MATHBLK1387X.

Форма этих колебаний определяется функцией \(U(r, \theta)\) , а характер
осцилляций --- множителем \(T(t)\).

Подставив это представление в волновое уравнение (см. задачу 1), получим
два уравнения

MATHBLK1388X

(уравнение Гельм-гольца),

MATHBLK1389X

(уравнение гармонических колебаний),

\pandocbounded{\includegraphics[keepaspectratio]{images/_page_221_Picture_9.jpeg}}

РИС. 30.1. Қолеблющаяся мембрана (гиперболическая смешанная задача).

где

MATHBLK1390X

Отметим, что здесь мы требуем, чтобы константа разделения была
положительна (поэтому мы обозначим ее \(\lambda^{2-1}\) )). Только в
этом случае функции \(T(t)\) будут периодическими.

На следующем шаге мы решим уравнение Гельмгольца, но сначала разберемся
с граничным условием для него. Чтобы найти его, подставим
\(u(r, \theta, t) = U(r, \theta) T(t)\) в граничное условие для мембраны

MATHBLK1391X

или

MATHBLK1392X

Следовательно, для того чтобы найти форму фундаментальных колебаний
мембраны, необходимо решить задачу

MATHBLK1393X

MATHBLK1394X

MATHBLK1395X

MATHBLK1396X

MATHBLK1397X

MATHBLK1398X

{[} \textbackslash begin\{array\}\{ll\}
(\textbackslash mathsf\{Y\}\textbackslash mathsf\{H\}\textbackslash mathsf\{\textbackslash Pi\})
\& \textbackslash Delta U + \textbackslash lambda\^{}2 U = 0,
\textbackslash{}
(\textbackslash mathsf\{\textbackslash Gamma\}\textbackslash mathsf\{Y\})
\& U(1, \textbackslash theta) = 0, \textbackslash end\{array\} {]}

MATHBLK1399X

MATHBLK1400X

подставим

MATHBLK1401X

в задачу Гельмгольца. Сделав это, получаем

MATHBLK1402X

(уравнение Бесселя), \(R(1) = 0\) , \(\Theta'' + n^2\Theta = 0\)

(проделайте все это сами). Отметим, что мы выбрали константу разделения
в виде \(n^2\) , \(n=0,. \ 1,\ 2,\ \ldots\) , так как хотим, чтобы
функции \(\Theta\left(\theta\right)\) были периодическими с периодом
\(2\pi\) . Очевидно, что форма мембраны является периодической функцией
по \(\theta\) . Итак, для того чтобы решить задачу Гельмгольца, мы
должны решить два обыкновенных дифференциальных уравнения

MATHBLK1403X, \(0 < r < 1\) , \(R(0) < \infty\) (физическое
ограничение), \(R(1) = 0\) , \(\Theta'' + n^2\Theta = 0\) .

\section{Уравнение
Бесселя}\label{ux443ux440ux430ux432ux43dux435ux43dux438ux435-ux431ux435ux441ux441ux435ux43bux44f}

Уравнение

MATHBLK1404X

жорошо известно в теории обыкновенных дифференциальных уравнений. Это
так называемое уравнение Бесселя. Оно имеет два

линейно-независимых решения

\(R_{1}(r) = A J_{n}(\lambda r)\) функция Бесселя \emph{первого} рода
\emph{п}-го порядка,

\(R_{2}(r) = BY_{n}(\lambda r)\) функция Бесселя второго рода n-го
порядка,

и, следовательно, общее решение этого уравнения представимо в виде

MATHBLK1405X

Отметим, что рещения зависят от параметров n и \(\lambda\) , входящих в
уравнение. Некоторые графики функции Бесселя приведены на рис. 30.2. Для
определения функций \(J_ n(\lambda r)\) и \(Y_n(\lambda r)\) можно

\pandocbounded{\includegraphics[keepaspectratio]{images/_page_223_Figure_8.jpeg}}

\pandocbounded{\includegraphics[keepaspectratio]{images/_page_223_Figure_9.jpeg}}

Рис. 30 2 Функции Бесселя. a --- функции Бесселя первого рода;
\(\pmb{\delta}\) --- функции Бесселя второго рода.

воспользоваться методом Фробениуса, который позволяет найти \(R(r)\) в
виде степенных рядов. Этим методом находятся два линейно-независимых
степенных ряда \(J_n(\lambda r)\) и \(Y_n(\lambda r)\) . Поскольку
функция \(Y_n(\lambda r)\) не ограничена при r=0, мы выбираем решение в
виде

MATHBLK1406X

Последний шаг в определении \(R\left( r\right)\) связан с использованием
граничного условия \(R\left( \mathbf{l}\right) =0\) для определения всех
\(\lambda\) (в данный

Таблица 30.1

{\def\LTcaptype{none} % do not increment counter
\begin{longtable}[]{@{}lll@{}}
\toprule\noalign{}
Корни & уравнений. & \(J_n(r) = 0\) \\
\midrule\noalign{}
\endhead
\bottomrule\noalign{}
\endlastfoot
------- & ------------ & -------------- \\
\end{longtable}
}

{\def\LTcaptype{none} % do not increment counter
\begin{longtable}[]{@{}lllllllll@{}}
\toprule\noalign{}
Γ & 1 & n & & & & & & \\
\midrule\noalign{}
\endhead
\bottomrule\noalign{}
\endlastfoot
& & 0 & 1 & 2 & 3 & 4 & & \\
& 1 & 2,40 & 3,83 & 5,13 & 6,38 & 7,59 & & \\
ı & 2 & 5,52 & 7,02 & 8,42 & 9,76 & 11,06 & & \\
ľ & 3 & 8,65 & 10,17 & 11,62 & 13,02 & 14,37 & & \\
& 4 & 11,79 & 13,32 & 14,80 & 16,22 & 17,62 & & \\
j & 5 & 14,93 & 16,47 & 17,96 & 19,41 & 20,83 & & \\
1 & . & • & & • & & • & & \\
& & & • & & • & • & & \\
- & • & • & • & • & • & • & & \\
\end{longtable}
}

момент нас не интересует коэффициент A). Подставляя \(R (1)\) = 0, в
\(R (r)\) = \(J_n\) ( \(\lambda r\) ), получаем \(J_n\) ( \(\lambda\) )
= 0. Другими словами, для того чтобы функция \(R (r)\) обращалась в нуль
на границе круга, мы должны выбрать константу разделения так, чтобы она
была корнем уравнения \(J_n\) (r) = 0, \(\tau\) . e.

MATHBLK1407X

где \(k_{nm}\) --- m-й корень уравнения \(J_n(r) = 0\) .

Хорошо известны таблицы этих корней, но их можно вычислять на ЭВМ.
Значения нескольких первых корней представлены в табл. 30.1. Зная эти
корни, мы сразу можем построить решение краевой задачи для уравнения
Гельмгольца. Если собственные значения \(\lambda = k_{nm}\) , то
собственные функции \(U_{ nm}(r)\) определяются формулой

MATHBLK1408X

На рис. 30.3 изображены некоторые из этих функций для различных значений
n и m. Общая форма функции \(U_{nm}(r,\theta)\) не зависит от значений
констант A и B. Эти константы влияют только на амплитуду колебаний и
расположение линии узлов относительно \(\theta=0\) .

Каждая функция представляет собой фундаментальное колебание круглой
мембраны с частотой

MATHBLK1409X

Эти частоты мы нашли из решений

MATHBLK1410X

\section{уравнения}\label{ux443ux440ux430ux432ux43dux435ux43dux438ux44f}

MATHBLK1411X

Интересно отметить, что отношение частоты колебания
\(U_{nm}(\mathbf{r}, \mathbf{\theta})\) к частоте основного колебания
U\_\{nt\}(\textbackslash ma

{[} thbf\{r\}, \textbackslash mathbf\{\textbackslash theta\}) , т.е. {]}

MATHBLK1412X

\begin{verbatim}
в отличие от одномерного случая не является целым числом. Другими словами, высшие колебания мембраны не являются чистыми обертонами основной частоты. Отметим также, что число m совпадает с числом круговых узловых линий у функции   U_{um}(r,\theta)
\end{verbatim}

\pandocbounded{\includegraphics[keepaspectratio]{images/_page_225_Figure_7.jpeg}}

РИС. 30.3. Линин узлов для основных гармоник круглой мембраны.

(амплитуда колебаний в точках узловых линий равна нулю). Радиусы этих
узловых линий определяются отношениями

\begin{verbatim}
MATHBLK1413X,   \frac{k_{02}}{k_{0m}}  ,   \frac{k_{03}}{k_{0m}}  , ..., 1.
\end{verbatim}

Теперь, когда найдены собственные функции задачи Гельмгольца, их следует
умножить на временной множитель

MATHBLK1414X

и найти такую линейную комбинацию получившихся функций, которая бы
удовлетворяла начальным условиям. Это значит, что решение исходной
задачи мы будем представлять в виде

MATHBLK1415X

\begin{verbatim}
Отметим, что, заменив величину   A\sin\left(n\theta\right)+B\cos\left(n\theta\right)   величиной   \cos\left(n\theta\right)  , мы только по-другому выбрали начало отсчета угла   \theta  . Константы   A_{nm}   и   B_{nm}   тоже можно было объединить в одну, но это уже несущественно.

Вместо того чтобы строить сложную процедуру отыскания коэффициентов   A_{nm}   и   B_{nm}   в общем случае, мы рассмотрим случай, когда u не зависит от   \theta   (тоже довольно общий). Другими словами, предположим, что начальное смещение мембраны зависит только от r; т. е.

MATHBLK1416X  u_i(r, \theta, 0) = 0.  (Так же легко можно рассмотреть случай   u_t \neq 0  .) При сделанных предположениях решение следует искать в виде
\end{verbatim}

(30.1)

MATHBLK1417X

Наша задача --- подобрать коэффициенты \(A_m\) таким образом, чтобы
выполнялось начальное условие

(30.2)

MATHBLK1418X

\begin{verbatim}
(условие   u_t=0   выполняется при этом автоматически). Для определения коэффициентов   A_m   воспользуемся ортогональностью функций Бесселя   \{J_0(k_{0m}r)\colon m=1,\ 2,\ \dots\}
\end{verbatim}

MATHBLK1419X

MATHBLK1420X

{[} \textbackslash begin\{cases\} 0, \& i \textbackslash neq j,
\textbackslash{} \textbackslash frac\{1\}\{2\}
J\_\{1\}\^{}\{2\}(k\_\{0i\}), \& i = j, \textbackslash end\{cases\} {]}

MATHBLK1421X

MATHBLK1422X

\begin{verbatim}
(метод вычисления этого интеграла можно найти в книгах по функциям Бесселя). Умножим обе части уравнения (30.2) на   rJ_0\left(k_{bf}r\right)   и проинтегрируем от 0 до 1:
\end{verbatim}

MATHBLK1423X

Из последнего соотношения определяем коэффициенты

(30.3)

MATHBLK1424X

\begin{verbatim}
Итак, не зависящие от   \theta   колебания мембраны описываются функцией вида (30.1), коэффициенты которой вычисляются по формулам (30.3). Это решение не так сложно, как могло бы показаться читателю. Его можно интерпретировать следующим образом. Разложим начальное условие f(r) в сумму
\end{verbatim}

MATHBLK1425X

\begin{verbatim}
а затем умножим каждое слагаемое на соответствующий осцилляторный множитель   \cos{(k_{0m}ct)}  , т. е.
\end{verbatim}

MATHBLK1426X

В результате получили решение исходной задачи. Например, решение задачи
о колебаниях мембраны с начальными условиями

\begin{verbatim}
MATHBLK1427X  u_t(r, \theta, 0) = 0  имеет вид
\end{verbatim}

MATHBLK1428X

\begin{verbatim}
(Поскольку   k_{01}=2,4048\ldots=2,40  ,   k_{03}=8,6537\ldots\neq 8,65  , то это решение не является точным.)
\end{verbatim}

\section{\texorpdfstring{Графики функций \(J_{0}\)(\(k_{01}\)r) ,
\(J_{0}\)(\(k_{02}\)r) , . .
.}{Графики функций J\_\{0\}(k\_\{01\}r) , J\_\{0\}(k\_\{02\}r) , . . .}}\label{ux433ux440ux430ux444ux438ux43aux438-ux444ux443ux43dux43aux446ux438ux439-j_0k_01r-j_0k_02r-.-.-.}

\begin{verbatim}
Читателю будет полезно иметь наглядное представление о графиках функций   J_{\ Mathfrak{g}}(k_{\mathfrak{g}_1}R)  ,   J_{\mathfrak{g}}(k_{\mathfrak{g}_2}r)   ..., так как решение выражается через эти функции. На рис. 30.4 изображен график   J_{\mathfrak{g}}(r)  .
\end{verbatim}

\pandocbounded{\includegraphics[keepaspectratio]{images/_page_227_Figure_15.jpeg}}

\begin{verbatim}
РИС. 30.4. Функция Бесселя нулевого порядка   J_{\phi}(r)  .
\end{verbatim}

Теперь, для того чтобы построить графики функций \(J_0\)(\(k_{01}\)r) ,
\(J_0\)(\(k_{02}\)r) , \ldots, \(J_0\)(\(k_{0m}\)r) , сожмем графики
вдоль оси x так, чтобы m-й нуль попал в точку r=1 (рис. 30.5).

\pandocbounded{\includegraphics[keepaspectratio]{images/_page_228_Figure_3.jpeg}}

\begin{verbatim}
РИС. 30 5. Графики функций   J_0\left(k_{\theta m}r\right)   (основные гармоники колебаний мембраны).
\end{verbatim}

\section{ЗАМЕЧАНИЯ}\label{ux437ux430ux43cux435ux447ux430ux43dux438ux44f-11}

\begin{enumerate}
\def\labelenumi{\arabic{enumi}.}
\tightlist
\item
  Предположим, что начальное смещение мембраны, у которой c=1, задается
  формулами
\end{enumerate}

(НУ)

MATHBLK1429X

MATHBLK1430X

MATHBLK1431X

\begin{verbatim}
В соответствии с общим методом мы должны разложить начальное смещение мембраны в сумму функций вида   J_{0}\left(k_{0m}r\right)  , найти колебания каждого слагаемого, а затем все колебания сложить. Основная особенность колебаний мембраны Р И С
\end{verbatim}

\pandocbounded{\includegraphics[keepaspectratio]{images/_page_228_Figure_9.jpeg}}

. 30 б. Начальные условия разлагаются по основным гармоникам.

\begin{verbatim}
в том, что каждое из смещений вида   J_{\mathrm{ }0}   (2,4 r),   J_{\mathrm{ }0}   (5,52 r), ...   J_{\mathrm{ }0}   (  k_{\mathrm{ }0} {\mathrm{ }0m}   r) совершает простое гармоническое движение (т. е. все точки мембраны колеблются с одной и той же частотой). В нашем случае начальное смещение является суммой двух простейших и, значит, мембрана будет участвовать в двух
\end{verbatim}

гармонических колебаннях (см. рис. 30.6)

MATHBLK1432X

В результате сложения этих простейших колебаний получается
результирующее сложное движение мембраны. Отметим еще, что частота
высокочастотного колебания не кратна частоте низкочастотного.

\begin{verbatim}
2. Если начальные условия заданы в виде u=f(r),   u_t=0  , то мы можем разложить начальное смещение f(r) в сумму элементарных «кирпичиков»   A_m J_a(k_{0m}r)  , каждый из которых совершает колебания вида   \cos(k_{0m}ct)   на своей собственной частоте. Сложив фундаментальные колебания
\end{verbatim}

MATHBLK1433X,

мы получим результирующие колебания, вызванные начальным смещением
\(f(r)\).

\section{ЗАДАЧИ}\label{ux437ux430ux434ux430ux447ux438-11}

\begin{verbatim}
1. Подставьте   u(r, \theta, t) = U(r, \theta) T(t)   в волновое уравнение
\end{verbatim}

MATHBLK1434X,

чтобы разделить его на два уравнения

MATHBLK1435X

MATHBLK1436X

\begin{verbatim}
2. Подстановкой   U\left(r,\;\theta\right)=R\left(r\right)\Theta\left(\theta\right)   приведите задачу Гельмгольца
\end{verbatim}

MATHBLK1437X

MATHBLK1438X

к виду

\begin{verbatim}
MATHBLK1439X  R(1) = 0,    \Theta'' + n^2\Theta = 0.

Константу разделения мы обозначили   n^2  . Почему n может принимать только значения   0, 1, 2, \ldots  ?
\end{verbatim}

\begin{enumerate}
\def\labelenumi{\arabic{enumi}.}
\setcounter{enumi}{2}
\tightlist
\item
  Решите задачу
\end{enumerate}

MATHBLK1440X

MATHBLK1441X

MATHBLK1442X

MATHBLK1443X

MATHBLK1444X

\begin{verbatim}
- 5. Решите задачу 4, если начальное смещение мембраны задано в виде   u(r, \theta, 0) = J_{\mathrm{ }o}(2,) 4r) 0.5\ J_{\mathrm{ }o}(8,65r) + 0.25\ J_{\mathrm{ }o}(14,93r)  . Какова наивысшая частота колебаний мембраны в этом случае?
\end{verbatim}

\begin{itemize}
\item
  \begin{enumerate}
  \def\labelenumi{\arabic{enumi}.}
  \setcounter{enumi}{5}
  \tightlist
  \item
    Постройте графики следующих функций на отрезке 0≤r≤1
  \end{enumerate}

  \begin{itemize}
  \item
    \begin{enumerate}
    \def\labelenumi{\alph{enumi})}
    \tightlist
    \item
      J\_\{o\}(5,52 r) ,
    \end{enumerate}

    \begin{itemize}
    \tightlist
    \item
      \begin{enumerate}
      \def\labelenumi{\arabic{enumi})}
      \setcounter{enumi}{5}
      \tightlist
      \item
        J\_0(14,93 \textbackslash r) .
      \end{enumerate}
    \end{itemize}
  \end{itemize}
\end{itemize}

\section{ЛАПЛАСИАН (ИНТУИТИВНОЕ
ОПИСАНИЕ)}\label{ux43bux430ux43fux43bux430ux441ux438ux430ux43d-ux438ux43dux442ux443ux438ux442ux438ux432ux43dux43eux435-ux43eux43fux438ux441ux430ux43dux438ux435}

ЦЕЛЬ ЛЕКЦИИ: Дать интуитивное описание лапласиана и привести его
выражения в различных системах координат. В наиболее важных системах
координат двух-и трехмерного пространства лапласиан записывается в виде

трехмерного пространства лапласнан записывается в виде

MATHBLK1445X

Поскольку преобразования координат вызывают затруднения у некогорых
студентов, мы обсудим правила, которыми надлежит пользоваться при
выполнении таких преобразований.

Оператор Лапласа

(31.1)

MATHBLK1446X

вероятно, самый важный оператор математической физики. Каков смысл этого
оператора и какое отношение имеет сумма трех частных производных второго
порядка к законам природы? Ответы на эти вопросы связаны с тем фактом,
что лапласиан функции позволяет оценить значение функции в точке через
значения функции в соседних точках. Кроме того, лапласиан можно считать
обобщением второй производной функции одной переменной на многомерный
случай.

Теперь мы переходим к основному свойству лапласиана, благодаря которому
он широко используется.

\section{Основное свойство двумерного оператора
А}\label{ux43eux441ux43dux43eux432ux43dux43eux435-ux441ux432ux43eux439ux441ux442ux432ux43e-ux434ux432ux443ux43cux435ux440ux43dux43eux433ux43e-ux43eux43fux435ux440ux430ux442ux43eux440ux430-ux430}

Если

\begin{verbatim}
1.   \Delta u > 0   в точке (x, y), то u(x, y) меньше «среднего значения функции в соседних точках» 1) (например, на окружности с центром в точке (x, y));

2.   \Delta u = 0   в точке (x, y), то u(x, y) равна «среднему значению
\end{verbatim}

функции в соседних точках»;

\begin{verbatim}
3.   \Delta u < 0   в точке (x, y), то u(x, y) больше «среднего значения функции в соседних точках».
\end{verbatim}

Используя этот факт 2), дадим интерпретацию основных уравнений
математической физики.

\section{Интуитивный смысл некоторых законов
физики}\label{ux438ux43dux442ux443ux438ux442ux438ux432ux43dux44bux439-ux441ux43cux44bux441ux43b-ux43dux435ux43aux43eux442ux43eux440ux44bux445-ux437ux430ux43aux43eux43dux43eux432-ux444ux438ux437ux438ux43aux438}

\begin{verbatim}
Согласно уравнению теплопроводности   u_t = \alpha^2 \Delta u   для температуры (или концентрации) и, скорость изменения температуры   u_t   пропорциональна величине   \Delta u  . Значит, если температура в точке температура на окружности, окружающей данную точку, то температура в данной точке будет возрастать.
\end{verbatim}

\pandocbounded{\includegraphics[keepaspectratio]{images/_page_232_Picture_11.jpeg}}

\pandocbounded{
\includegraphics[keepaspectratio]{images/_page_232_Picture_12.jpeg}}

\begin{verbatim}
РИС. 31.1. Интуитивное описание лапласиана   \Delta u  : a— среднее значение функции u на окружности с центром в точке p равно u (p),   \Delta u   (p) = 0;   \delta  — среднее значение функции u на окружности с центром в точке p больше, чем u (p),   \Delta u   (p) > 0;   \delta  — среднее значение функции u на окружности с центром в точке p меньше, чем u (p),   \Delta u   (p) < 0.

Согласно волновому уравнению   u_{tt} = \alpha^2 \Delta u   для смещения мембраны, ускорение точки мембраны   u_{tt}   (или сила, действующая на точку) пропорционально величине   \Delta u  . Значит, точка мембраны
\end{verbatim}

2) Этот факт (теорема о среднем) не зависит от размерности пространства.
--- Прим. ред.

1) Под «средним значением функции в соседних точках» понимается среднее
значение функции либо по окружности, либо по кругу с центром в точке (x,
y). --- Прим. ред.

ускоряется вверх (сила направлена вверх), если ее смещение (по высоте)
меньше, чем среднее смещение соседних точек.

\begin{verbatim}
Если функция удовлетворяет уравнению Лапласа   \Delta u=0  , то ее вначение всегда совпадает со средним значением. Например, натянутая неподвижная резиновая мембрана удовлетворяет уравнению Лапласа, следовательно, смещение мембраны в любой точке равно среднему смещению мембраны на окружности с центром в этой точке. Рис. 31.1 иллюстрирует это свойство лапласиана.

У равнение Пуассона   \Delta u = f  , где f — функция, зависящая от пространственных переменных (может быть и константой), описывает значительное число различных явлений.

- 1.   \Delta u = -\rho   описывает потенциал электростатического поля, если   \rho  —плотность статических зарядов. Зная основное свойство лапласиана, что бы вы могли сказать о свойствах потенциального поля?

- 2.   \Delta u = -g(x, y)   описывает стационарное распределение температуры от теплового источника g(x, y). Если g(x, y) положительна в точке, значит, в этой точке тепло выделяется, если же g(x, y) отрицательна, то тепло в этой точке поглощается.
\end{verbatim}

\begin{itemize}
\tightlist
\item
  \begin{enumerate}
  \def\labelenumi{\arabic{enumi}.}
  \setcounter{enumi}{2}
  \tightlist
  \item
    Δu + λu = 0 это уравнение Гельмгольца (приведенное волновое
    уравнение), которое описывает (помимо многого другого)
    пространственную часть собственных колебаний мембраны. Это уравнение
    возникает при разделении пространственных и временной переменных в
    волновом уравнении и уравнении теплопроводности.
  \end{enumerate}
\end{itemize}

Пока что мы говорили о физическом смысле лапласиана.

MATHBLK1447X

\begin{verbatim}
(двумерный случай),   \Delta u = u_{xx} + u_{yy} + u_{zz}   (трехмерный случай).

Однако во многих задачах возникает необходимость знать выражение для лапласнана в других координатах. Если, например, граница области — окружность, то естественно перейти к полярным координатам   (r, \theta)  , а если область трехмерна, а ее граница — сфера, то следует перейти к сферическим координатам   (r, \theta, \phi)  . Итак, нам предстоит ответить на вопрос: как выглядит лапласиан в различных системах координат?
\end{verbatim}

Прежде чем приступить к ответу на этот вопрос, напомним кратко о пяти
основных системах координат в двумерном и трехмерном пространствах:

декартова прямоугольная система в двумерном пространстве, декартова
прямоугольная система в трехмерном пространстве, полярные координаты в
двумерном пространстве,

цилиндрические координаты в трехмерном пространстве, сферические
координаты в трехмерном пространстве.

Полярные координаты определяются соотношениями

MATHBLK1448X

MATHBLK1449X

{[} r\^{}2 = x\^{}2 + y\^{}2, \& x = r \textbackslash cos
\textbackslash theta, \textbackslash{} {]}

MATHBLK1450X

{[} \textbackslash text\{tg \} \textbackslash theta = y/x, \& y = r
\textbackslash sin \textbackslash theta. {]}

MATHBLK1451X

MATHBLK1452X Ц и л и

\pandocbounded{\includegraphics[keepaspectratio]{images/_page_234_Picture_4.jpeg}}

ндрические координаты связаны с декартовыми следующими соотношениями:

\begin{verbatim}
MATHBLK1453X,   x = r \cos \theta  ,
      tg \theta = y/x  ,   tg \theta = y/x  ,   tg \theta = z  ,   tg \theta  ,   TG \theta = z
\end{verbatim}

\pandocbounded{\includegraphics[keepaspectratio]{images/_page_234_Picture_7.jpeg}}

Сферические координаты определяются следующим образом:

\begin{verbatim}
MATHBLK1454X,   x = r \sin \varphi \cos \theta  ,   \cos \varphi = z/r  , или   y = r \sin \varphi \sin \theta  ,   tg \theta = y/x  ,   z = r \cos \varphi  .
\end{verbatim}

\pandocbounded{\includegraphics[keepaspectratio]{images/_page_235_Picture_2.jpeg}}

\section{Преобразование
координат}\label{ux43fux440ux435ux43eux431ux440ux430ux437ux43eux432ux430ux43dux438ux435-ux43aux43eux43eux440ux434ux438ux43dux430ux442}

Теперь мы подготовлены к преобразованию координат. В качестве примера
посмотрим, как преобразуется лапласиан при переходе к полярным
координатам. Если мы разберемся в этом, то Р И С

\pandocbounded{\includegraphics[keepaspectratio]{images/_page_235_Picture_5.jpeg}}

\begin{verbatim}
312 Функционатьная зависимость функции u от переменных   (r, \theta)   и (x, y)
\end{verbatim}

\pandocbounded{\includegraphics[keepaspectratio]{images/_page_235_Picture_7.jpeg}}

РИС. 31 3 Функциональная зависимость производной \(u_x\) от переменных
(r, 0) и (x, y).

сможем легко переходить от одной системы координат к другой. Итак, у нас
есть лапласиан

MATHBLK1455X

\begin{verbatim}
Чтобы найти лапласиан в новой системе координат, обратимся к диаграмме, изображенной на рис. 31.2. Цель диаграммы — точно показать, что от чего зависит. Например, из этой диаграммы видно, что функция u зависит от новых переменных r и   \theta  .
\end{verbatim}

Каждая из этих переменных в свою очередь зависит от x и y (соединяющие
линии указывают на эту зависимость).

\begin{verbatim}
Мы могли бы нарисовать диаграмму, в основании которой находились бы r и   \theta  , но поскольку нам предстоит вычислять   u_{xx}   и   u_{yy}  , то больше подходит именно эта диаграмма. Как же теперь выразить   u_{xx}   и   u_{yy}   через r и   \theta  ?
\end{verbatim}

На самом деле сделать это очень легко. Сначала найдем \(u_x\) и \(u_y\)
. Мы найдем \(u_x\) , складывая пути на диаграмме, ведущие от u к x. В
данном случае существуют только два таких пути, так что

MATHBLK1456X

По правилам дифференцирования сложных функций следует перемножить
производные, написанные у ребер каждого пути, а затем сложить их.
Аналогично

MATHBLK1457X

Обратимся теперь ко вторым производным. Найдем \(u_{xx}\) (в качестве
упражнения предоставим читателю возможность найти \(u_{yy}\) ).

\begin{verbatim}
"Сначала, однако, построим другую диаграмму (см. рис. 31.3). Читателю, вероятно, ясно, что если u зависит от r и   \theta   (которые в свою очередь зависят от x и y), то   u_x   так же (вообще говоря) зависит от тех же переменных. Таким образом,
\end{verbatim}

MATHBLK1458X

MATHBLK1459X

MATHBLK1460X

MATHBLK1461X

MATHBLK1462X

Аналогично

MATHBLK1463X

Складывая \(u_{xx}\) и \(u_{yy}\) , получаем окончательный результат

MATHBLK1464X

т. е. выражение для лапласиана в полярных координатах. Его смысл такой
же, как и у лапласиана в декартовых координатах, но форма записи другая.
К сожалению, теперь это оператор с переменными коэффициентами, так что
уравнение, содержащее лапласианы в полярных координатах, решать труднее.

Аналогично можно преобразовать и трехмерный лапласиан

MATHBLK1465X

В результате получаем в цилиндрических координатах

MATHBLK1466X

в сферических координатах

MATHBLK1467X

\section{ЗАМЕЧАНИЯ}\label{ux437ux430ux43cux435ux447ux430ux43dux438ux44f-12}

\begin{enumerate}
\def\labelenumi{\arabic{enumi}.}
\item
  Только в декартовой системе координат лапласиан обладает постоянными
  коэффициентами. Этим объясняется, почему задачи в других системах
  координат решать труднее. Тем не менее к этим уравнениям с переменными
  коэффициентами все еще применим метод разделения переменных, правда,
  получающиеся при этом обыкновенные дифференциальные уравнения тоже
  будут с переменными коэффициентами. Так получаются уравнение Бесселя,
  уравнение Лежандра и другие, так называемые классические уравнения
  математической физики. Такие уравнения решаются обычно методом
  разложения в степенные ряды.
\item
  Кроме упомянутых здесь, существует много др гих координатных систем
  1). Например, задача о колебаниях эллиптической мембраны требует
  представления лапласиана в эллиптической системе координат. Чтобы
  сделать это, требуется несколько иной подход: преобразовать эллипс в
  круг, решить получившуюся задачу в новых (полярных) координатах, и
  затем полученное решение преобразовать к исходным координатам 2).
\end{enumerate}

2) См. Тихонов А. Н., Самарский А. А. Уравнения математической физики.
--- М.: Наука, 1972 и др. издания. --- Прим. ред.

2) Целесообразнее непосредственно перейти к эллиптическим координатам.
--- Прим. ред.

\section{ЗАЛАЧИ}\label{ux437ux430ux43bux430ux447ux438}

\begin{verbatim}
- 1. Как записать волновое уравнение   u_{tt} = \alpha^2 \Delta u   в сферических координатах, если вам известно, что решение зависит только OT r H t?

- 2. Как записать волновое уравнение   u_{tt} = \alpha^2 \Delta u   в полярных координатах, если известно, что решение зависит только от г - 3. Как записать уравнение Лапласа   \Delta u = 0   в полярных координатах, если известно, что u зависит только от r? Каково решение этого уравнения? Эти функции называются потенциалами с'криговой симметрией.
\end{verbatim}

\begin{itemize}
\item
  \begin{enumerate}
  \def\labelenumi{\arabic{enumi}.}
  \setcounter{enumi}{3}
  \tightlist
  \item
    Как записать уравнение Лапласа в сферических координатах, если
    решение и зависит только от г? Можете ли вы найти решения этого
    уравнения? Эти функции называются сферически-симметричными
    потенциалами,
  \end{enumerate}
\item
  \begin{enumerate}
  \def\labelenumi{\arabic{enumi}.}
  \setcounter{enumi}{4}
  \tightlist
  \item
    Что вы можете сказать о поверхности \(u(x, y)\) = xy\^{}2
  \end{enumerate}

  \begin{itemize}
  \tightlist
  \item
    \begin{enumerate}
    \def\labelenumi{\arabic{enumi}.}
    \setcounter{enumi}{5}
    \tightlist
    \item
      Преобразуйте трехмерный лапласиан \textbackslash Delta u =
      u\_\{xx\} + u\_\{yy\} + u\_\{zz\} к сферическим координатам.
    \end{enumerate}
  \end{itemize}
\end{itemize}

Лекция 32

\section{ОБЩИЕ СВОЙСТВА КРАЕВЫХ
ЗАДАЧ}\label{ux43eux431ux449ux438ux435-ux441ux432ux43eux439ux441ux442ux432ux430-ux43aux440ux430ux435ux432ux44bux445-ux437ux430ux434ux430ux447}

ЦЕЛЬ ЛЕКЦИИ: Объяснить, как возникают уравнения с частными производными,
не содержащие производных по времени. В отличие от гиперболического
волнового уравнения или параболического уравнения теплопроводности эти
уравнения не требуют начальных условий. Для них нужны только граничные
условия. По этой причине задачи для таких уравнений называются краевыми
задачами

Укажем три наиболее важных типа граничных условий.

\begin{itemize}
\tightlist
\item
  Граничное условие первого рода (условие Дирихле).
\end{itemize}

Граничное условие второго рода (условие Неймана).

\begin{itemize}
\tightlist
\item
  \begin{enumerate}
  \def\labelenumi{\arabic{enumi}.}
  \setcounter{enumi}{2}
  \tightlist
  \item
    Граничное условие третьего рода (условие Робена).
  \end{enumerate}
\item
  В лекции проводятся примеры и пояснения этих граничных условий.
\end{itemize}

До сих пор мы изучали явления, которые развивались в пространстве и во
времени. Существует, однако, много явлений, которые не изменяются с
течением времени. Эти явления в большинстве случаев описываются краевыми
эллиптическими задачами, и цель нашей лекции---познакомить читателя с
такого типа задачами.

Можно указать два случая, когда в уравнениях математической физики не
содержится время

\begin{enumerate}
\def\labelenumi{\arabic{enumi}.}
\tightlist
\item
  Стационарные задачи
\end{enumerate}

2 Задачи, в которых зависимость от времени входит в виде временного
множителя.

Сначала давайте рассмотрим стационарные задачи.

\section{Стационарные
задачи}\label{ux441ux442ux430ux446ux438ux43eux43dux430ux440ux43dux44bux435-ux437ux430ux434ux430ux447ux438}

\begin{verbatim}
Предположим, что нас интересует стационарное (при   t \to \infty  ) решение уравнения теплопроводности
\end{verbatim}

MATHBLK1468X.

Очевидно, если решение не изменяется во времени, то \(u_t = 0\)и
равнение теплопроводности переходит в уравнение Лапласа

MATHBLK1469X

Проиллюстрируем подробно идею стационарности на примере. Рассмотрим
задачу

MATHBLK1470X

MATHBLK1471X

MATHBLK1472X

MATHBLK1473X

{[} \textbackslash begin\{cases\}
\textbackslash frac\{d\textsuperscript{\{2\}u\}\{dx}\{2\}\} =
-\textbackslash sin(\textbackslash pi x), \& 0 \textless{} x \textless{}
1, \textbackslash{} u(0) = 0, \textbackslash{} u(1) = 0.
\textbackslash end\{cases\} {]}

MATHBLK1474X

MATHBLK1475X

Решение этой задачи имеет вид

MATHBLK1476X

\begin{verbatim}
Если решить задачу (32.1), то ее решение в начальный момент будет совпадать с   \sin{(3\pi x)}  , а затем все в большей и большей степен походить на
\end{verbatim}

MATHBLK1477X

\begin{verbatim}
Для некоторых задач может не существовать стационарное решение (не зависящее от времени), а для других стационарным решением может оказаться временная синусоида. Следовательно, не всегда справедливо, что производные   u_t, u_{tt} \dots   равны нулю. Поэтому мы, как правило, должны иметь некоторую информацию о физическом смысле задачи.
\end{verbatim}

\section{Выделение временной компоненты в гиперболических и
параболических
задачах}\label{ux432ux44bux434ux435ux43bux435ux43dux438ux435-ux432ux440ux435ux43cux435ux43dux43dux43eux439-ux43aux43eux43cux43fux43eux43dux435ux43dux442ux44b-ux432-ux433ux438ux43fux435ux440ux431ux43eux43bux438ux447ux435ux441ux43aux438ux445-ux438-ux43fux430ux440ux430ux431ux43eux43bux438ux447ux435ux441ux43aux438ux445-ux437ux430ux434ux430ux447ux430ux445}

В задаче о колебаниях мембраны

(УЧП)

MATHBLK1478X

\begin{verbatim}
(ГУ)   u = 0, \quad \text{on a circle,}
\end{verbatim}

(НУ)

MATHBLK1479X

MATHBLK1480X

{[} \textbackslash begin\{cases\} u(r, \textbackslash theta, 0) = f(r,
\textbackslash theta), \textbackslash{} u\_t(r, \textbackslash theta, 0)
= g(r, \textbackslash theta) \textbackslash end\{cases\} {]}

MATHBLK1481X

MATHBLK1482X

MATHBLK1483X

{[} \textbackslash begin\{aligned\}
(\textbackslash mathrm\{Y\}\textbackslash mathrm\{H\}\textbackslash mathrm\{\textbackslash Pi\})
\& \textbackslash Delta U + \textbackslash lambda\^{}2 U = 0,
\textbackslash{} (\textbackslash mathrm\{H\}\textbackslash mathrm\{Y\})
\& U\textbackslash left(1,,\textbackslash theta\textbackslash right) =
0. \textbackslash end\{aligned\} {]}

MATHBLK1484X

MATHBLK1485X

\begin{verbatim}
Такая ситуация характерна для теории уравнений с частными производными, где решение зачастую представляет собой произведение пространственной части   U\left(r,\theta\right)   на временную часть   T\left(t\right)  . Фактически мы получим то же самое уравнение Гельмгольца, если попытаемся выделить временную компоненту в уравнении теплопроводности
\end{verbatim}

MATHBLK1486X.

При решении краевых задач чаще всего используются три типа граничных
условий. На них мы сейчас остановимся.

\section{Три основных типа граничных условий в краевых
задачах}\label{ux442ux440ux438-ux43eux441ux43dux43eux432ux43dux44bux445-ux442ux438ux43fux430-ux433ux440ux430ux43dux438ux447ux43dux44bux445-ux443ux441ux43bux43eux432ux438ux439-ux432-ux43aux440ux430ux435ux432ux44bux445-ux437ux430ux434ux430ux447ux430ux445}

Краевая задача с граничными условиями первого рода (задача Дирихле)

Требуется найти решение уравнения в некоторой области пространства,
которое принимает на границе области заданные значения. В качестве
примера можно привести задачу о нахождении стационарного распределения
температуры внутри области,

если задана температура на границе этой области. Другой пример: найти
распределение потенциала внутри области, если известен потенциал на
границе. Ниже приведен еще ряд примеров.

\section{ПРИМЕРЫ ЗАДАЧ
ДИРИХЛЕ}\label{ux43fux440ux438ux43cux435ux440ux44b-ux437ux430ux434ux430ux447-ux434ux438ux440ux438ux445ux43bux435}

Рассмотрим уравнение Лапласа внутри круга с заданными значениями решения
на границе (см. рис. 32.1):

(УЧП)

MATHBLK1487X

\begin{verbatim}
(ГУ)   u(1, \theta) = \sin \theta, \quad 0 \le \theta < 2\pi.
\end{verbatim}

В качестве еще одного примера можно указать внешнюю задачу Дирихле, для
которой решение уравнения Лапласа пишется вне круга, а граничные
значения заданы по окружности (см. рис. 32.2).

\pandocbounded{\includegraphics[keepaspectratio]{images/_page_241_Picture_7.jpeg}}

РИС, 32.1. Внутренняя задача Дирихле.

\pandocbounded{\includegraphics[keepaspectratio]{images/_page_241_Picture_9.jpeg}}

РИС. 32.2. Внешняя задача Дирихле.

\section{Краевые задачи с граничным условием второго рода (задача
Неймана)}\label{ux43aux440ux430ux435ux432ux44bux435-ux437ux430ux434ux430ux447ux438-ux441-ux433ux440ux430ux43dux438ux447ux43dux44bux43c-ux443ux441ux43bux43eux432ux438ux435ux43c-ux432ux442ux43eux440ux43eux433ux43e-ux440ux43eux434ux430-ux437ux430ux434ux430ux447ux430-ux43dux435ux439ux43cux430ux43dux430}

Требуется найти решение уравнения в некоторой области пространства, на
границе которой задана внешняя нормальная производная ди/дп (которая
пропорциональна втекающему потоку). Это общая задача и для стационарной
теплопроводности, и для электростатики, если на границе задан поток
(тепла, электронов и т. д.).

Например, предположим, что тепловой поток на границе круга

изменяется по закону

MATHBLK1488X

Тогда стационарное распределение температуры внутри круга является
решением краевой задачи

MATHBLK1489X

MATHBLK1490X

{[} \textbackslash begin\{cases\} \textbackslash Delta u = 0, \& 0
\textless{} r \textless{} 1, \textbackslash{}
\textbackslash frac\{\textbackslash partial u\}\{\textbackslash partial
r\} = \textbackslash sin \textbackslash theta, \& r = 1, \& 0
\textbackslash le \textbackslash theta \textless{} 2\textbackslash pi.
\textbackslash end\{cases\} {]}

MATHBLK1491X

MATHBLK1492X

\begin{verbatim}
В этой задаче *поток* тепла (кал/см·с) через границу направлен енутрь при   0 \leqslant \theta < \pi   и наружу при   \pi \leqslant \theta < 2\pi  .
\end{verbatim}

Однако, поскольку полный поток

MATHBLK1493X

(условие, которое всегда должно выполняться в задаче Неймана), можно
утверждать, что температура каждой точки внутри круга не изменяется во
времени.

\pandocbounded{\includegraphics[keepaspectratio]{images/_page_242_Picture_9.jpeg}}

\textbf{РИС.} 32.3. Поток тепла в задаче Неймана. Температура в каждой
точке \textbf{мостоянна}, поскольку полное количество тепла, втекающего
в точку и вытежающего из него, равно нулю.

Другими словами, задача Неймана (для уравнения Лапласа) имеет смысл
только в том случае, когда полный поток тепла через границу равен нулю.
Математически это значит, что на границе должно выполняться соотношение

MATHBLK1494X

иначе задача не будет иметь решения. Например, внутренняя задача Неймана

MATHBLK1495X

MATHBLK1496X

{[} \textbackslash begin\{cases\} \textbackslash Delta u = 0, \& 0
\textless{} r \textless{} 1, \textbackslash{}
\textbackslash frac\{\textbackslash partial u\}\{\textbackslash partial
r\}(1, \textbackslash theta) = 1, \& r = 1, \& 0 \textbackslash le
\textbackslash theta \textless{} 2\textbackslash pi,
\textbackslash end\{cases\} {]}

MATHBLK1497X

MATHBLK1498X

MATHBLK1499X

{[} \textbackslash begin\{cases\} \textbackslash Delta u = 0, \& 0
\textless{} r \textless{} 1, \textbackslash{}
\textbackslash frac\{\textbackslash partial u\}\{\textbackslash partial
r\}(1, \textbackslash theta) =
\textbackslash cos(2\textbackslash theta), \& r = 1, \& 0
\textbackslash le \textbackslash theta \textless{} 2\textbackslash pi,
\textbackslash end\{cases\} {]}

MATHBLK1500X

MATHBLK1501X

\begin{verbatim}
имеет бесконечное множество решений   u(r,) \theta)   (заметим, что общий поток через границу равен нулю). Если у нас есть одно решение, то, добавляя к нему константы, можно пайти и все другие решения. Например, одно из решений нашей задачи Неймана имеет вид
\end{verbatim}

MATHBLK1502X

Если прибавить к нему произвольную константу, то получится другое
решение. По этой причине для выделения единственного решения нужно иметь
дополнительную информацию (например, знать значение решения в одной
точке).

\section{\texorpdfstring{\textbf{Краевые задачи с граничным условием}
третьего
рода}{Краевые задачи с граничным условием третьего рода}}\label{ux43aux440ux430ux435ux432ux44bux435-ux437ux430ux434ux430ux447ux438-ux441-ux433ux440ux430ux43dux438ux447ux43dux44bux43c-ux443ux441ux43bux43eux432ux438ux435ux43c-ux442ux440ux435ux442ux44cux435ux433ux43e-ux440ux43eux434ux430}

Требуется найти такое решение уравнения в некоторой области
пространства, которое удовлетворяет на границе условию вида

MATHBLK1503X

где h---заданная константа, а g---заданная функция, которая, вообще
говоря, меняется вдоль границы. Обратим внимание на то, что новое
граничное условие является как бы смесью граничных условий первых двух
типов. Более привычная форма

этого условия имеет вид

MATHBLK1504X

Согласно этому граничному условию, поток, втекающий в область через
границу, пропорционален разности между температурой u и некоторой
заданной температурой g. Это означает, что (при h\textgreater0)

\begin{enumerate}
\def\labelenumi{\arabic{enumi})}
\item
  если температура \emph{и} на границе выше температуры окружающей среды
  на границе \emph{g}, то тепло вытекает из области;
\item
  если и меньше температуры g, то тепло втекает в область. Конечно, это
  не что иное, как закон теплообмена Ньютона, а граничные условия всех
  трех типов совершенно естественны для стационарных задач
  теплопроводности.
\end{enumerate}

\section{ПРИМЕР}\label{ux43fux440ux438ux43cux435ux440-1}

\begin{verbatim}
Предположим, что вне единичного круга температура распределена по закону   g(\theta) = \sin \theta  . В этом случае поток тепла через границу должен задаваться в виде
\end{verbatim}

MATHBLK1505X

Следовательно, стационарное распределение температуры внутри круга
является решением краевой задачи

(32.2)

MATHBLK1506X

MATHBLK1507X

MATHBLK1508X

\begin{verbatim}
Константа h является физическим параметром и численно равна количеству тепла, протекающему через границу при разности температур в один градус. Если h велика, то при разности температур в один градус тепловой поток велик и решение становится похожим на решение задачи Дирихле с граничным условием f =   \sin \theta  . С другой стороны, если f = 0, то граничное условне переходит в условие f f f f f f f f f f
\end{verbatim}

MATHBLK1509X

\begin{verbatim}
Читатель может представить, как изменяются решения задачи (32.2), когда параметр h пробегает диапазон от 0 до   +\infty  . Если h=0, то граница теплоизолирована и, следовательно, решение   u(r,\theta)   является константой (решение в данном случае не единственное, так что любая константа является решением). Если h становится все больше и больше, то решение становится все всявие и больше похожим на решение задачи Дирихле с гранич-

ным условием   u=\sin\theta  . Если h>0 и стремится к нулю, то и решение также будет стремиться к тождественному нулю (поскольку нуль—это среднее значение   \sin\theta   на границе).
\end{verbatim}

\section{ЗАДАЧИ}\label{ux437ux430ux434ux430ux447ux438-12}

\begin{enumerate}
\def\labelenumi{\arabic{enumi}.}
\tightlist
\item
  Можете ли вы из интуптивных соображений решить задачу Дирихле?
\end{enumerate}

MATHBLK1510X

MATHBLK1511X

MATHBLK1512X

\begin{enumerate}
\def\labelenumi{\arabic{enumi}.}
\setcounter{enumi}{1}
\tightlist
\item
  Имеет ли задача Неймана
\end{enumerate}

MATHBLK1513X

{[} \textbackslash Delta u = 0 \textless{} , 0 r \textless{} 1 ,
(\textbackslash Gamma \textbackslash forall)
\textbackslash frac\{\textbackslash partial u\}\{\textbackslash partial
r\} = \textbackslash sin\^{}2 \textbackslash theta {]}

решение внутри круга?

\begin{enumerate}
\def\labelenumi{\arabic{enumi}.}
\setcounter{enumi}{2}
\tightlist
\item
  Попытайтесь представить, как выглядит решение задачи
\end{enumerate}

(УЧП)

\begin{verbatim}
MATHBLK1514X,   0 < r < 1  ,

(ГУ)   \frac{\partial u}{\partial r} + h(u - \sin \theta) = 0
\end{verbatim}

при различных значениях параметра h.

\begin{enumerate}
\def\labelenumi{\arabic{enumi}.}
\setcounter{enumi}{3}
\tightlist
\item
  Какую краевую задачу нужно решить, чтобы найти стационарные решения
  следующей задачи:
\end{enumerate}

MATHBLK1515X

MATHBLK1516X

\begin{verbatim}
\right. \quad 0 < t < \infty, \\ (\text{НУ}) & \left\{
\end{verbatim}

MATHBLK1517X

{[} \textbackslash begin\{aligned\} u\textbackslash left(x,
~0\textbackslash right) =
\textbackslash sin\textbackslash left(3\textbackslash pi
x\textbackslash right), \textbackslash{} u\_t\textbackslash left(x,
~0\textbackslash right) = 0? \textbackslash end\{aligned\} {]}

MATHBLK1518X

MATHBLK1519X

Исходя из физической интерпретации лапласиана, что вы можете сказать об
общих свойствах решений краевой задачи для уравнения Гельмгольца

(УЧП)

\begin{verbatim}
MATHBLK1520X,   0 < r < 1  ,   (\Gamma Y)     u(1, \theta) = 0  ,   0 \le \theta < 2\pi  ?
\end{verbatim}

\begin{enumerate}
\def\labelenumi{\arabic{enumi}.}
\setcounter{enumi}{5}
\tightlist
\item
  Дайте физическую интерпретацию краевой задачи с граничными
\end{enumerate}

условиями смешанного типа внутри квадрата

(YUII)

MATHBLK1521X

MATHBLK1522X

MATHBLK1523X

{[} \textbackslash begin\{cases\} u\_y(x, 0) - h{[}u(x, 0) - 2{]} = 0,
\& 0 \textless{} x \textless{} 1, \textbackslash{} u(x, 1) = 1, \& 0
\textless{} x \textless{} 1, \textbackslash{} u\_x(0, y) = 0, \& 0
\textless{} y \textless{} 1. \textbackslash end\{cases\} {]}

MATHBLK1524X

MATHBLK1525X

MATHBLK1526X

{[} \textbackslash begin\{cases\} u\_y(x, 0) - h{[}u(x, 0) - 2{]} = 0,
\& 0 \textless{} x \textless{} 1, \textbackslash{} u\_x(0, y) = 0, \& 0
\textless{} y \textless{} 1. \textbackslash end\{cases\} {]}

MATHBLK1527X

MATHBLK1528X

Лекция 33

\section{ВНУТРЕННЯЯ ЗАДАЧА ДИРИХЛЕ ДЛЯ
КРУГА}\label{ux432ux43dux443ux442ux440ux435ux43dux43dux44fux44f-ux437ux430ux434ux430ux447ux430-ux434ux438ux440ux438ux445ux43bux435-ux434ux43bux44f-ux43aux440ux443ux433ux430}

ЦЕЛЬ ЛЕКЦИИ: Показать, как методом разделения переменных решигь
внутреннюю задачу Дирихле для круга

(УЧП)

MATHBLK1529X

\begin{verbatim}
(ГУ)   u(1, \theta) = g(\theta), \quad 0 \le \theta < 2\pi.
\end{verbatim}

Сначала граничное условие разлагается в ряд

MATHBLK1530X

а затем ищется решение каждой из задач

\begin{verbatim}
MATHBLK1531X  \Delta u = 0,    u(1, \theta) = \sin(n\theta),     u(1, \theta) = \cos n\theta.
\end{verbatim}

Решения этих задач записываются в виде

MATHBLK1532X

откуда получается решение задачи Дирихле

MATHBLK1533X

После того как решение в виде ряда получено, путем алгебраических
преобразований решению можно придать другую, интегральную форму, которая
называется интегральной формулой Пуассона. В этой лекции предлагается
множество новых идей, как решить задачу Дирихле

(33.1)

MATHBLK1534X

MATHBLK1535X

MATHBLK1536X

{[} u(r, \textbackslash theta) =
\textbackslash sum\_\{n=0\}\^{}\{\textbackslash infty\} r\^{}n
\textbackslash left{[} a\_n \textbackslash cos(n\textbackslash theta) +
b\_n \textbackslash sin(n\textbackslash theta) \textbackslash right{]}.
{]}

MATHBLK1537X

{[} g(\textbackslash theta) =
\textbackslash sum\_\{n=0\}\^{}\{\textbackslash infty\} {[}a\_n
\textbackslash cos(n\textbackslash theta) + b\_n
\textbackslash sin(n\textbackslash theta){]}, {]}

MATHBLK1538X

{[} u(r, \textbackslash theta) =
\textbackslash sum\_\{n=0\}\^{}\{\textbackslash infty\} r\^{}n {[}a\_n
\textbackslash cos(n\textbackslash theta) + b\_n
\textbackslash sin(n\textbackslash theta){]}. {]}

MATHBLK1539X

{[} \textbackslash Delta u = 0 {]}

MATHBLK1540X

{[} u(r, \textbackslash theta) = 1 + r \textbackslash sin
\textbackslash theta + \textbackslash frac\{r\^{}3\}\{2\}
\textbackslash sin (3\textbackslash theta) + r\^{}4 \textbackslash cos
(4\textbackslash theta) {]}

MATHBLK1541X

{[} u(r, \textbackslash theta) =
\textbackslash sum\_\{n=0\}\^{}\{\textbackslash infty\} (r/R)\^{}n
{[}a\_n \textbackslash cos(n\textbackslash theta) + b\_n
\textbackslash sin(n\textbackslash theta){]}. {]}

MATHBLK1542X

{[} a\_0 = \textbackslash frac\{1\}\{2\textbackslash pi\}
\textbackslash int\_0\^{}\{2\textbackslash pi\} g(\textbackslash theta)
d\textbackslash theta. {]}

MATHBLK1543X

{[} u(r, \textbackslash theta) =
\textbackslash sum\_\{n=0\}\^{}\{\textbackslash infty\} (r/R)\^{}n
\textbackslash left{[} a\_n \textbackslash cos(n\textbackslash theta) +
b\_n \textbackslash sin(n\textbackslash theta) \textbackslash right{]},
{]}

MATHBLK1544X

{[} u(r, \textbackslash theta) =
\textbackslash frac\{1\}\{2\textbackslash pi\}
\textbackslash int\_\{0\}\^{}\{2\textbackslash pi\}
g(\textbackslash theta) d\textbackslash theta +
\textbackslash frac\{1\}\{\textbackslash pi\}
\textbackslash sum\_\{n=1\}\^{}\{\textbackslash infty\} (r/R)\^{}\{n\}
\textbackslash int\_\{0\}\^{}\{2\textbackslash pi\}
g(\textbackslash alpha) \textbackslash times \textbackslash{}
\textbackslash times
{[}\textbackslash cos(n\textbackslash alpha)\textbackslash cos(n\textbackslash theta)
+
\textbackslash sin(n\textbackslash alpha)\textbackslash sin(n\textbackslash theta){]}
d\textbackslash alpha = \textbackslash{} =
\textbackslash frac\{1\}\{2\textbackslash pi\}
\textbackslash int\_\{0\}\^{}\{2\textbackslash pi\}
\textbackslash left\{ 1 + 2
\textbackslash sum\_\{n=1\}\^{}\{\textbackslash infty\} (r/R)\^{}\{n\}
\textbackslash cos{[}n(\textbackslash theta - \textbackslash alpha){]}
\textbackslash right\} g(\textbackslash alpha) d\textbackslash alpha =
\textbackslash{} = \textbackslash frac\{1\}\{2\textbackslash pi\}
\textbackslash int\_\{0\}\^{}\{2\textbackslash pi\}
\textbackslash left\{ 1 +
\textbackslash sum\_\{n=1\}\^{}\{\textbackslash infty\} (r/R)\^{}\{n\}
\textbackslash left{[} e\^{}\{in(\textbackslash theta -
\textbackslash alpha)\} + e\^{}\{-in(\textbackslash theta -
\textbackslash alpha)\} \textbackslash right{]} \textbackslash right\}
g(\textbackslash alpha) d\textbackslash alpha = \textbackslash{} =
\textbackslash frac\{1\}\{2\textbackslash pi\}
\textbackslash int\_\{0\}\^{}\{2\textbackslash pi\}
\textbackslash left\{ 1 +
\textbackslash frac\{re\^{}\{i(\textbackslash theta -
\textbackslash alpha)\}\}\{\textbackslash left{[}R -
re\^{}\{i(\textbackslash theta -
\textbackslash alpha)\}\textbackslash right{]}\} +
\textbackslash frac\{re\^{}\{-i(\textbackslash theta -
\textbackslash alpha)\}\}\{\textbackslash left{[}R -
re\^{}\{-i(\textbackslash theta -
\textbackslash alpha)\}\textbackslash right{]}\} \textbackslash right\}
g(\textbackslash alpha) d\textbackslash alpha = \textbackslash{} =
\textbackslash frac\{1\}\{2\textbackslash pi\}
\textbackslash int\_\{0\}\^{}\{2\textbackslash pi\}
\textbackslash left{[} \textbackslash frac\{R\^{}\{2\} -
r\textsuperscript{\{2\}\}\{R}\{2\} -
2rR\textbackslash cos(\textbackslash theta - \textbackslash alpha) +
r\^{}\{2\}\} \textbackslash right{]} g(\textbackslash alpha)
d\textbackslash alpha. {]}

MATHBLK1545X

{[} u(r, \textbackslash theta) =
\textbackslash frac\{1\}\{2\textbackslash pi\}
\textbackslash int\_\{0\}\^{}\{2\textbackslash pi\}
\textbackslash left{[} \textbackslash frac\{R\^{}\{2\} -
r\textsuperscript{\{2\}\}\{R}\{2\} - 2rR\textbackslash cos((0 -
\textbackslash alpha) + r\^{}\{2\})\} \textbackslash right{]}
g(\textbackslash alpha) d\textbackslash alpha. {]}

MATHBLK1546X

{[}
\textbackslash frac\{R\textsuperscript{2-r}2\}\{R\textsuperscript{2-2rR\textbackslash cos(\textbackslash theta-\textbackslash alpha)+r}2\}.
{]}

MATHBLK1547X

{[} u(0, 0) = \textbackslash frac\{1\}\{2\textbackslash pi\}
\textbackslash int\_\{0\}\^{}\{2\textbackslash pi\}
g(\textbackslash alpha) d\textbackslash alpha. {]}

MATHBLK1548X

(УЧП)

MATHBLK1549X

{[} \textbackslash Delta W = 0 {]}

MATHBLK1550X

(УЧП)

MATHBLK1551X

{[} \textbackslash Delta W = \textbackslash Delta V {]}

MATHBLK1552X

{[} u\_\{rr\} + \textbackslash frac\{1\}\{r\} u\_r +
\textbackslash frac\{1\}\{r\^{}3\}
u\_\{\textbackslash theta\textbackslash theta\} = 0 {]}

MATHBLK1553X

{[} u(1, \textbackslash theta) = 1 + \textbackslash sin
\textbackslash theta + \textbackslash frac\{1\}\{2\} \textbackslash cos
\textbackslash theta {]}

MATHBLK1554X

{[} \textbackslash Delta u = 0 {]}

MATHBLK1555X

{[}
(\textbackslash mathsf\{Y\}\textbackslash mathsf\{Y\}\textbackslash mathsf\{\textbackslash Pi\})
\textbackslash quad \textbackslash Delta u = 0, \textbackslash quad 0
\textless{} r \textless{} 1, {]}

MATHBLK1556X

{[}
(\textbackslash mathsf\{\textbackslash Gamma\}\textbackslash mathsf\{Y\})
\textbackslash quad u(1, , \textbackslash theta) = {]}

MATHBLK1557X

MATHBLK1558X

Прикиньте, как выглядит график решения?

\begin{verbatim}
- 6. Пусть в ядре Пуассона r=3R/4,   \theta=\pi/2  . Постройте график зависимости ядра Пуассона от угла   \alpha   в диапазоне   0 \le \alpha < 2\pi  .
\end{verbatim}

\begin{itemize}
\tightlist
\item
  \begin{enumerate}
  \def\labelenumi{\arabic{enumi}.}
  \setcounter{enumi}{6}
  \tightlist
  \item
    Убедитесь в справедливости замечания 1 из этой лекции.
  \end{enumerate}
\item
  \begin{enumerate}
  \def\labelenumi{\arabic{enumi}.}
  \setcounter{enumi}{7}
  \tightlist
  \item
    Убедитесь в справедливости замечания 2 из этой лекции.
  \end{enumerate}
\end{itemize}

\section{Лекция 34}\label{ux43bux435ux43aux446ux438ux44f-34}

\section{ЗАДАЧА ДИРИХЛЕ В
КОЛЬЦЕ}\label{ux437ux430ux434ux430ux447ux430-ux434ux438ux440ux438ux445ux43bux435-ux432-ux43aux43eux43bux44cux446ux435}

ЦЕЛЬ ЛЕКЦИИ: Решить задачу Дирихле в кольце

(УЧП)

MATHBLK1559X

(ГУ)

{[} \textbackslash begin\{cases\} u(R\_1, \textbackslash theta) =
g\_1(\textbackslash theta), \textbackslash{} u(R\_2,
\textbackslash theta) = g\_2(\textbackslash theta).
\textbackslash end\{cases\} {]}

{[} 0 \textbackslash le \textbackslash theta \textbackslash le
2\textbackslash pi, {]}

Здесь мы опять воспользуемся методом разделения переменных, но в отличие
от внутренней зодачи Дирихле мы не станем отбрасывать решения вида

MATHBLK1560X

\begin{verbatim}
Следовательно, решение должно записываться в виде   u(r, \theta) = a_0 + b_0 \ln r +
\end{verbatim}

MATHBLK1561X

\begin{verbatim}
Мы также кратко обсудим решение внешней задачи Дирихле. В этом случае нам придется отбросить слагаемые, неограниченные при   r=\infty  . Следовательно, решение внешней задачи Дирихле
\end{verbatim}

MATHBLK1562X

MATHBLK1563X

MATHBLK1564X

следует искать в виде

MATHBLK1565X

Существует множество разнообразных по форме областей, для которых можно
в явном виде выписать решение задачи Дирихле. Мы перечислим только
некоторые из них:

\begin{itemize}
\item
  а) в круге (лекция 33),
\item
  б) в кольце (данная лекция),
\item
  в) вне круга (данная лекция),
\item
  г) в шаре (в последующих лекциях),
\item
  д) в сферическом слое (в последующих лекциях),
\item
  е) в полосе (двумерный случай),
\item
  ж) в слое между плоскостями (трехмерный случай).
\end{itemize}

Эгот список можно продолжать до бесконечности. Мы хотим решить простую,
но достаточно типичную задачу Дприхле, так чтобы читатель освоил общие
принципы и мог при необходимости решить свою задачу.

\pandocbounded{\includegraphics[keepaspectratio]{images/_page_254_Picture_3.jpeg}}

РИС. 34.1. Уравнение Лапласа в кольце.

Здесь наша цель---найти форму мыльной пленки, натинутой между двумя
проволочными ободками. Вероятно, в этой задаче интунция уже не играет
такой роли, как во внутренней задаче Дирихле. Математическую модель
задачи запишем следующим образом:

(34.1)

MATHBLK1566X

MATHBLK1567X

MATHBLK1568X

Общая схема этой задачи изображена на рис. 34.1.

\begin{verbatim}
Задача состоит в том, чтобы найти решение   u(r, \theta)   между окружностями   r = R_1   и   r = R_2  , совпадающее на окружностях с   g_1(\theta)   и соответственно с   g_2(\theta)  . Будем искать решения вида
\end{verbatim}

MATHBLK1569X.

\begin{verbatim}
Подставляя их в уравнение Лапласа, приходим к двум обыкновенным дифференциальным уравнениям для функций R(r) и   \Theta(\theta)
\end{verbatim}

MATHBLK1570X

\begin{verbatim}
(уравнение Эйлера),   \Theta'' + \lambda^2\Theta = 0  .

Отметим, что константа разделения в этих уравнениях должна быть больше или равна нулю (мы обозначили ее   \lambda^2  ), иначе реше-

ние для   \Theta   (  \theta  ) не будет периодическим. Далее мы решим оба этих уравнения, а затем перемножим эти решения, в результате получим все решения уравнения Лапласа, которые представимы в виде R (r)   \Theta   (  \theta  ).
\end{verbatim}

Из этих двух уравнений наибольший интерес представляет уравнение Эйлера.
К счастью, это одно из немногих обыкновенных дифференциальных уравнений
с переменными коэффициентами, которые решаются сравнительно легко.

\section{РЕШЕНИЕ УРАВНЕНИЯ
ЭЙЛЕРА}\label{ux440ux435ux448ux435ux43dux438ux435-ux443ux440ux430ux432ux43dux435ux43dux438ux44f-ux44dux439ux43bux435ux440ux430}

MATHBLK1571X

\begin{verbatim}
СЛУЧАЙ I (  \lambda = 0  ). В этом случае уравнение Эйлера принимает вид
\end{verbatim}

MATHBLK1572X

общее решение которого легко находится:

MATHBLK1573X

\begin{verbatim}
В самом деле, сделав подстановку   V\left(r\right) \Rightarrow R'\left(r\right)  , получаем для V уравнение с разделяющимися переменными
\end{verbatim}

MATHBLK1574X.

\begin{verbatim}
Решив это уравнение, находим   V\left(r\right)=c_{1}/r  , откуда
\end{verbatim}

MATHBLK1575X.

\begin{verbatim}
СЛУЧАЙ 2. (  \lambda > 0  ). В этом случае уравнение Эйлера имеет вид
\end{verbatim}

MATHBLK1576X

\begin{verbatim}
Будем искать решение этого уравнения в виде   R(r)=r^{\alpha}  . Наша цель—найти два значения   \alpha   (скажем,   \alpha_1   и   \alpha_2  ), таких, чтобы функция
\end{verbatim}

MATHBLK1577X

\begin{verbatim}
являлась общим решением уравнения. Подстановка   R(r) = r^{\alpha}   в уравнение Эйлера приводит к   \alpha = \lambda   и   \alpha = -\lambda  . Следовательно,
\end{verbatim}

MATHBLK1578X

Итак, используя решение уравнения Эйлера, мы получили

MATHBLK1579X

MATHBLK1580X

{[} \textbackslash begin\{cases\} R(r) = a + b \textbackslash ln r,
\textbackslash{} \textbackslash Theta(\textbackslash theta) = c +
d\textbackslash theta, \textbackslash{} \textbackslash lambda
\textgreater{} 0 \textbackslash end\{cases\} {]}

MATHBLK1581X

MATHBLK1582X

MATHBLK1583X

{[} \textbackslash begin\{cases\} R(r) = ar\^{}\{\textbackslash lambda\}
+ br\^{}\{-\textbackslash lambda\}, \textbackslash{}
\textbackslash Theta(\textbackslash theta) = c
\textbackslash cos(\textbackslash lambda \textbackslash theta) + d
\textbackslash sin(\textbackslash lambda \textbackslash theta).
\textbackslash end\{cases\} {]}

MATHBLK1584X

MATHBLK1585X

{[} c - \textbackslash text\{константа\}, {]}

MATHBLK1586X

{[} u(r, \textbackslash theta) = a\_0 + b\_0 \textbackslash ln r +
\textbackslash sum\_\{n=1\}\^{}\{\textbackslash infty\}
\textbackslash left{[} (a\_n r\^{}n + b\_n r\^{}\{-n\})
\textbackslash cos(n\textbackslash theta) + (c\_n r\^{}n + d\_n
r\^{}\{-n\}) \textbackslash sin(n\textbackslash theta)
\textbackslash right{]}. {]}

MATHBLK1587X

MATHBLK1588X

{[} \textbackslash begin\{cases\} u(R\_1, \textbackslash theta) =
g\_1(\textbackslash theta), \textbackslash{} u(R\_2,
\textbackslash theta) = g\_2(\textbackslash theta).
\textbackslash end\{cases\} {]}

MATHBLK1589X

MATHBLK1590X

MATHBLK1591X

{[} \textbackslash begin\{cases\} a\_0 + b\_0 \textbackslash ln R\_1 =
\textbackslash frac\{1\}\{2\textbackslash pi\}
\textbackslash int\textbackslash limits\_0\^{}\{2\textbackslash pi\}
g\_1(s) , ds \textbackslash{} a\_0 + b\_0 \textbackslash ln R\_2 =
\textbackslash frac\{1\}\{2\textbackslash pi\}
\textbackslash int\textbackslash limits\_0\^{}\{2\textbackslash pi\}
g\_2(s) , ds \textbackslash end\{cases\} {]}

MATHBLK1592X

MATHBLK1593X

MATHBLK1594X

{[} \textbackslash begin\{cases\} a\_n R\_1\^{}n + b\_n R\_1\^{}\{-n\} =
\textbackslash frac\{1\}\{\textbackslash pi\}
\textbackslash int\textbackslash limits\_0\^{}\{2\textbackslash pi\}
g\_1(s) \textbackslash cos(ns) , ds \textbackslash{} a\_n R\_2\^{}n +
b\_n R\_2\^{}\{-n\} = \textbackslash frac\{1\}\{\textbackslash pi\}
\textbackslash int\textbackslash limits\_0\^{}\{2\textbackslash pi\}
g\_2(s) \textbackslash cos(ns) , ds \textbackslash end\{cases\} {]}

MATHBLK1595X

MATHBLK1596X

MATHBLK1597X

MATHBLK1598X

Тем самым из этих уравнений находятся все неизвестные коэффициенты
\(a_0\) , \(b_0\) , \(a_n\) , \(b_n\) , \(c_n\) и \(d_n\) . Теперь
задача (34.1) полностью решена. Решение дается выражением (34.2),
коэффициенты которого определяются по формулам (34.3).

Чтобы лучше понять полученные результаты, рассмотрим

несколько примеров.

\section{Некоторые варианты задачи Дирихле в
кольце}\label{ux43dux435ux43aux43eux442ux43eux440ux44bux435-ux432ux430ux440ux438ux430ux43dux442ux44b-ux437ux430ux434ux430ux447ux438-ux434ux438ux440ux438ux445ux43bux435-ux432-ux43aux43eux43bux44cux446ux435}

ПРИМЕР 1.

\begin{verbatim}
Предположим, что потенциал на внутренней окружности равен нулю, а потенциал на внешней равен   \sin\theta  . Тогда
\end{verbatim}

(УЧП)

\begin{verbatim}
MATHBLK1599X,   1 < r < 2  ,
\end{verbatim}

(ГУ)

MATHBLK1600X

MATHBLK1601X

{[} \textbackslash begin\{cases\} u(1, \textbackslash theta) = 0,
\textbackslash{} u(2, \textbackslash theta) = \textbackslash sin
\textbackslash theta. \textbackslash end\{cases\} {]}

MATHBLK1602X

MATHBLK1603X

{[} 0 \textbackslash le \textbackslash theta \textbackslash le
2\textbackslash pi , {]}

\begin{verbatim}
Первое, что мы должны сделать—это вычислить все интегралы в (34.3) и решить соответствующие уравнения для   a_0  ,   b_0  ,   a_n  ,   b_n  ,   c_n   и   d_n  . Целесообразно сначала попытаться подобрать такие частные решения (из выписанных выше), линейная комбинация которых удовлетворяет граничным условням (здесь   u = Ar \sin \theta + Br^{-1} \sin \theta  ). Затем из граничных условий найти коэффициенты (здесь A = 2/3, B = -2/3). В результате этих простых вычислений получаем

MATHBLK1604X,   b_0=0  ,   a_n=0  ,   b_n=0     (n=1, 2, ...)  ,   c_n=

\begin{cases} 2/3, & n=1, \\0 & \text{для всіх інших } n, \\ d_n=\begin{cases} -2/3, & n=1, \\0 & \text{для всіх інших} n. \end{cases}
\end{verbatim}

Зная все коэффициенты, выпишем решение задачи

MATHBLK1605X

Легко проверить, что эта функция удовлетворяет обоим граничным условиям.
Очевидно и то, что она удовлетворяет уравнению Лапласа, так как
принадлежит классу функций вида (34.2).

\section{ПРИМЕР 2.}\label{ux43fux440ux438ux43cux435ux440-2.}

Рассмотрим задачу с \emph{постоянными потенциалами} на гранинах кольца

(УЧП)

\begin{verbatim}
MATHBLK1606X,   1 < r < 2  .
\end{verbatim}

(ГУ)

MATHBLK1607X

MATHBLK1608X

{[} \textbackslash begin\{cases\} u(1, \textbackslash theta) = 3,
\textbackslash{} u(2, \textbackslash theta) = 5.
\textbackslash end\{cases\} {]}

MATHBLK1609X

MATHBLK1610X

\begin{verbatim}
В этом случае можно сэкономить массу времени, поскольку очевидно, что решение не зависит от   \theta   (так как граничные условия не зависят от   \theta  ). Другими словами, наше решение следует искать среди функций вида   a_0 + b_0 \ln r  . Подставляя эту функцию в граничные условия, получаем
\end{verbatim}

MATHBLK1611X,

{[} a\_0 + b\_0 \textbackslash ln 2 = 5 , {]}

\begin{verbatim}
откуда   a_0 = 3  ,   b_0 = 2/\ln 2 = 2.9  . Следовательно, решением будет функция
\end{verbatim}

MATHBLK1612X.

Ее график изображен на рис. 342.

\pandocbounded{\includegraphics[keepaspectratio]{images/_page_258_Figure_10.jpeg}}

РИС. 34.2. Зависимость потенциала от радиуса внутри кольца.

\section{ПРИМЕР 3}\label{ux43fux440ux438ux43cux435ux440-3}

Решим еще одну интересную задачу

(УЧП)

\begin{verbatim}
MATHBLK1613X,   1 < r < 2  ,
\end{verbatim}

(ГУ)

MATHBLK1614X

MATHBLK1615X

{[} \textbackslash begin\{cases\} u(1, \textbackslash theta) =
\textbackslash sin \textbackslash theta, \textbackslash{} u(2,
\textbackslash theta) = \textbackslash sin \textbackslash theta.
\textbackslash end\{cases\} {]}

MATHBLK1616X

MATHBLK1617X

Можно быстро проверить, что все коэффициенты \(a_0\) , \(b_0\) , \(a_n\)
, \(b_n\) , \(c_n\) и \(d_n\) равны нулю, а коэффициенты \(c_i\) и
\(d_i\) определяются из системы уравнений

MATHBLK1618X

MATHBLK1619X

Решая эту систему, получаем \(c_1=1/3\)и \(d_1=2/3\).Следовательно,
решением задачи является функция

MATHBLK1620X

\begin{verbatim}
Форма кривых при различных значениях   \theta   показана на рис. 34.3.
\end{verbatim}

\pandocbounded{\includegraphics[keepaspectratio]{images/_page_259_Figure_7.jpeg}}

\begin{verbatim}
РИС. 34.3 Мыльная оченка между   u\left(1,\;\theta\right)=\sin\theta   и   u\left(2,\;\theta\right)=\sin\theta  .
\end{verbatim}

Наша лекция завершается кратким обсуждением внешней за-дачи Дирихле для
круга.

\section{Внешняя задача
Дирихле}\label{ux432ux43dux435ux448ux43dux44fux44f-ux437ux430ux434ux430ux447ux430-ux434ux438ux440ux438ux445ux43bux435}

Внешняя задача Дирихле

MATHBLK1621X

MATHBLK1622X

MATHBLK1623X

решается точно так же, как решалась внутренняя задача в лекции 33. За
исключением того, что теперь необходимо отбросить те решения, которые не
ограничены при стремлении г к бесконечности, т. е.

\begin{verbatim}
MATHBLK1624X,   r^n \cos(n\theta)     H     r^n \sin(n\theta)  .
\end{verbatim}

Следовательно, в качестве решения остается взять функцию

(34.4)

MATHBLK1625X

где коэффициенты \(a_n\) и \(b_n\) вычисляются по тем же формулам, что и
раньше:

MATHBLK1626X

MATHBLK1627X

\begin{verbatim}
Другими словами, мы просто разлагаем функцию   u\left(1,\ ;\theta\right)=g\left(\theta\right)   в ряд Фурье
\end{verbatim}

MATHBLK1628X

а затем каждый член этого ряда домножаем на коэффициент \(r^{-n}\) .
Чтобы теснее познакомиться с этим решением, приведем два примера.

Примеры внешней задачи Дирихле

Внешняя задача

(УЧП)

\begin{verbatim}
MATHBLK1629X,   1 < r < \infty  ,

(ГУ)   u(1, \theta) = 1 + \sin \theta + \cos (3\theta)  ,   0 \le \theta \le 2\pi  ,
\end{verbatim}

имеет решение

MATHBLK1630X

ПРИМЕР 2

Внешняя задача

(УЧП)

\begin{verbatim}
MATHBLK1631X,   1 < r < \infty  ,

(ГУ)   u(r, \theta) = \cos(4\theta)  ,   0 \le \theta \le 2\pi  ,
\end{verbatim}

имеет решение

MATHBLK1632X

Читатель должен попытаться представить, как выглядит это решение

\section{ЗАМЕЧАНИЯ}\label{ux437ux430ux43cux435ux447ux430ux43dux438ux44f-13}

Решение внешней задачи Дирихле при произвольном радиусе R

MATHBLK1633X

MATHBLK1634X

{[} \textbackslash begin\{aligned\} (\textbackslash text\{YHI\}) \&
\textbackslash Delta u = 0, \& R \textless{} r \textless{}
\textbackslash infty, \textbackslash{}
(\textbackslash text\{\textbackslash GammaY\}) \& u (R,
\textbackslash theta) = g (\textbackslash theta), \& 0
\textbackslash leqslant \textbackslash theta \textbackslash leqslant
2\textbackslash pi, \textbackslash end\{aligned\} {]}

MATHBLK1635X

MATHBLK1636X

имеет вид

MATHBLK1637X

\begin{verbatim}
2 Существует только два решения двумерного уравнения Лапласа, которые не зависят от   \theta  , а зависят только от r. Это константа c и   \ln r   Потенциал   \ln r   играет большую роль в теорин и называется логарифмическим потенциалом. Ниже мы подробнее остановимся на его свойствах
\end{verbatim}

\section{ЗАДАЧИ}\label{ux437ux430ux434ux430ux447ux438-13}

1 Решите задачу Дирихле

(УЧП)

\begin{verbatim}
MATHBLK1638X,   1 < r < 2  ,
\end{verbatim}

(ГУ)

MATHBLK1639X

MATHBLK1640X

{[} \textbackslash begin\{cases\} u(1, \textbackslash theta) =
\textbackslash cos \textbackslash theta, \textbackslash{} u(2,
\textbackslash theta) = \textbackslash sin \textbackslash theta.
\textbackslash end\{cases\} {]}

MATHBLK1641X

MATHBLK1642X

\begin{enumerate}
\def\labelenumi{\arabic{enumi}.}
\setcounter{enumi}{1}
\tightlist
\item
  Найдите решение внешней задачи Дирихле
\end{enumerate}

MATHBLK1643X

при следующих граничных условиях:

\begin{enumerate}
\def\labelenumi{\alph{enumi})}
\tightlist
\item
\end{enumerate}

MATHBLK1644X,

\begin{enumerate}
\def\labelenumi{\arabic{enumi})}
\setcounter{enumi}{5}
\tightlist
\item
\end{enumerate}

MATHBLK1645X,

\begin{enumerate}
\def\labelenumi{\Alph{enumi})}
\setcounter{enumi}{1}
\tightlist
\item
\end{enumerate}

MATHBLK1646X,

\begin{enumerate}
\def\labelenumi{\Alph{enumi})}
\setcounter{enumi}{1}
\tightlist
\item
\end{enumerate}

MATHBLK1647X,

{[} r) u(1, \textbackslash theta) = {]}

{[} \textbackslash begin\{cases\} 1, \& 0 \textbackslash le
\textbackslash theta \textless{} \textbackslash pi, \textbackslash{} 0,
\& \textbackslash pi \textbackslash le \textbackslash theta \textless{}
2\textbackslash pi. \textbackslash end\{cases\} {]}

3 Рещение внешней задачи Неймана

MATHBLK1648X

MATHBLK1649X

{[} \textbackslash begin\{aligned\} (\textbackslash forall
\textbackslash Pi \textbackslash Pi) \& \textbackslash Delta u = 0,
\textbackslash quad 1 \textless{} r \textless{} \textbackslash infty,
\textbackslash{} (\textbackslash Gamma \textbackslash forall) \&
\textbackslash frac\{\textbackslash partial u\}\{\textbackslash partial
r\} (1, \textbackslash theta) = g(\textbackslash theta),
\textbackslash quad 0 \textbackslash leqslant \textbackslash theta
\textbackslash leqslant 2\textbackslash pi,
\textbackslash end\{aligned\} {]}

MATHBLK1650X

MATHBLK1651X

имеет тот же вид, что и решение задачи Дирихле, т. е.

MATHBLK1652X

но теперь коэффициенты \(a_n\) и \(b_n\) должны удовлетворять новым
граничным условиям. Подставьте это решение в граничное условие

MATHBLK1653X

и получите решение задачи

MATHBLK1654X

MATHBLK1655X

{[} \textbackslash begin\{cases\} \textbackslash Delta u = 0, \& 1
\textless{} r \textless{} \textbackslash infty, \textbackslash{}
\textbackslash frac\{\textbackslash partial u\}\{\textbackslash partial
r\}(1, \textbackslash theta) = \textbackslash sin \textbackslash theta.
\textbackslash end\{cases\} {]}

MATHBLK1656X

MATHBLK1657X

MATHBLK1658X

{[} \textbackslash begin\{cases\} u(R\_1, \textbackslash theta) =
g\_1(\textbackslash theta), \textbackslash{} u(R\_2,
\textbackslash theta) = g\_2(\textbackslash theta)
\textbackslash end\{cases\} {]}

MATHBLK1659X

MATHBLK1660X

и получите уравнения (34.3).

Лекция 35

\section{УРАВНЕНИЕ ПАПЛАСА В СФЕРИЧЕСКИХ КООРДИНАТАХ (СФЕРИЧЕСКИЕ
ГАРМОНИКИ)}\label{ux443ux440ux430ux432ux43dux435ux43dux438ux435-ux43fux430ux43fux43bux430ux441ux430-ux432-ux441ux444ux435ux440ux438ux447ux435ux441ux43aux438ux445-ux43aux43eux43eux440ux434ux438ux43dux430ux442ux430ux445-ux441ux444ux435ux440ux438ux447ux435ux441ux43aux438ux435-ux433ux430ux440ux43cux43eux43dux438ux43aux438}

ЦЕЛЬ ЛЕКЦИИ: Найти такие частные решения уравнения Лапласа в сферических
координатах, чтобы из них можно было строить решения различных задач
(типа задач Дирихле, Неймана). Кроме того, мы решим внутреннюю задачу
Дирихле

(УЧП)

MATHBLK1661X),

(ГУ)

MATHBLK1662X

\begin{verbatim}
для одного частного случая, когда потенциал на границе зависит только от   \phi  .

<sup>1)</sup> Здесь левая часть уравнения есть   r^2 \Lambda u   в сферических координатах.—   \Pi \rho u m  .   \rho e \partial_v

Для решения задачи нам придется разложить потенциал на границе   g(\varphi)   в ряд по сферическим гармоникам
\end{verbatim}

MATHBLK1663X

\begin{verbatim}
где частные решения уравнения Лапласа, так называемые сферические гармоники   P_n   (  \cos \varphi  ), являются полиномами Лежандра от   \cos \varphi  . После того как разложение произведено, решение находится сразу:
\end{verbatim}

MATHBLK1664X

Аналогично находится решение внешней задачи Дирихле:

MATHBLK1665X

Одной из важных задач физики является задача определения потенциала
внутри или вие сферы, если задан потенциал на поверхности сферы. Для
внутренней задачи мы должны найти функцию, которая является решением
задачи

(35 1)

MATHBLK1666X

MATHBLK1667X

\begin{verbatim}
Обратим внимание на то, что сферический лапласнан записан не так, как раньше. Эта форма записи более компактна и более удобна для использования. Другая типичная интерпретация этой задачи такова: найти распределение температуры внутри сферы, если температура на границе сферы задана. Очень часто функция   g(\theta, \phi)   имеет специальный вид, так что нет необходимости решать задачу в самом общем виде.

В этой лекции мы остановимся на двух важных частных случаях. В одном случае функция   g(\theta, \phi)   будет константой, а в другом — будет зависеть только от угла   \phi   (этот угол отсчитывается от северного полюса сферы).
\end{verbatim}

\section{Частные случаи задачи
Дирихле}\label{ux447ux430ux441ux442ux43dux44bux435-ux441ux43bux443ux447ux430ux438-ux437ux430ux434ux430ux447ux438-ux434ux438ux440ux438ux445ux43bux435}

\begin{verbatim}
Частный случай   I(g(\theta, \phi) = \text{const})
\end{verbatim}

Ясно, что в этом случае решение не зависит от углов ф и 0 и уравнение
Лапласа переходит в обыкновенное дифференциальное уравнение

MATHBLK1668X

Это очень простое уравнение, и читатель может легко решить его. Общее
решение этого уравнения имеет вид

(35.2)

MATHBLK1669X.

Другими словами, среди всех потенциалов существует только два таких,
которые не зависят от углов, а зависят только от r. Эти потенциалы---
константа и c/r. Потенциал 1/r играет очень важную роль в физике и
называется ньютоновским потенциалом.

\pandocbounded{\includegraphics[keepaspectratio]{images/_page_264_Picture_9.jpeg}}

РИС. 35.1. Задача Дирихле в щаре

Теперь мы решим две задачи, в каждой из которых потенциал зависит только
от r.

ЗАДАЧА 1 (Потенциал внутри сферы)

(УЧП)

\begin{verbatim}
MATHBLK1670X,   0 < r < 1  , (ГУ)   u(1, \theta, \varphi) = 3  .

Согласно (35.2), получаем ограниченное решение   u(r, \theta, \phi) = 3  .   3A \квадрат A \ квадрат A A = 2   (Потенциал между двумя сферами, каждая из которых обладает постоянным потенциалом)
\end{verbatim}

Предположим, что мы хотим найти стационарное распределение температуры
между двумя сферами, которые поддерживаются при различных температурах:

MATHBLK1671X

MATHBLK1672X

MATHBLK1673X

{[} \textbackslash begin\{cases\} u(R\_1, \textbackslash theta,
\textbackslash varphi) = A, \textbackslash{} u(R\_2,
\textbackslash theta, \textbackslash varphi) = B.
\textbackslash end\{cases\} {]}

MATHBLK1674X

MATHBLK1675X

\begin{verbatim}
Поскольку мы знаем, что общая форма потенциала, не зависящего от   \theta   и   \varphi  ,
\end{verbatim}

MATHBLK1676X

то подставим его в граничное условие и решим получившуюся систему
относительно а и b. В результате получаем

\begin{verbatim}
MATHBLK1677X,   b = \frac{R_2B-R_1A}{R_2-R_1}  .
\end{verbatim}

Графики этого потенциала при различных значениях потенциалов A и B
изображены на рис. 35.2.

\pandocbounded{\includegraphics[keepaspectratio]{images/_page_265_Figure_9.jpeg}}

РИС 35 2. Потенциал между двумя концентрическими сферами, одна из
которых обладает потенциалом A, а другая---потенциалом B.

\begin{verbatim}
Частный случай 2 (  g(\theta, \varphi)   зависит только от   \varphi  )
\end{verbatim}

В этом случае задача Дирихле записывается в виде

(УЧП)

\begin{verbatim}
MATHBLK1678X,   0 < r < 1  ,   0 \le \varphi \le \pi  , (ГУ)   u(1) , \theta, \varphi) = g(\varphi)  ,   0 \le \varphi \le \pi  .
\end{verbatim}

Воспользуемся методом разделення переменных и найдем решения вида

MATHBLK1679X

В результате приходим к двум обыкновенным дифференциальным

уравнениям:

MATHBLK1680X

\begin{verbatim}
(уравнение Эйлера), [  \sin \phi \Phi'  ]' +   n(n+1)\sin \phi \Phi = 0   (уравнение Лежандра).
\end{verbatim}

Немного позже читатель увидит, почему константа разделения выбрана в
виде \(n(n+1)\).

\begin{verbatim}
Уравнение Эйлера мы решим, подставив в него   R(r) = r^x   и определив возможные значения   \alpha  . Сделав это, получим
\end{verbatim}

MATHBLK1681X

MATHBLK1682X

MATHBLK1683X

Следовательно, общее решение уравнения Эйлера имеет вид \(R(r)\) =
ar\^{}n + br\^{}\{-(n+1)\} .

Уравнение Лежандра решить не так легко. Общий метод решения основан на
подстановке

MATHBLK1684X.

Проделав эту замену переменной, получаем новую форму уравнения Лежандра

MATHBLK1685X

\begin{verbatim}
Теперь нужно решить это уравнение и подставить в решение   x = \cos \varphi  . Уравнение Лежандра является линейным уравнением вгорого порядка с переменными коэффициентами. Одна из особенностей этого уравнения в том, что коэффициент   (1-x^2)   при   d^2\Phi/dx^2   обращается в нуль на концах отрезка   -1 \leqslant x \leqslant 1  . Уравнения этого типа называются сингулярными дифференциальными уравнениями и очень часто решаются методом Фробеннуса Не вдаваясь в подробности этого метода, мы приведем только наиболее важный результат. Ограниченные решения уравнения Лежандра могут существовать, только если   n=0,1,2,\ldots   При этом оказывается, что они являются полиномами   P_n(x)   (полиномами Лежандра)
\end{verbatim}

MATHBLK1686X

MATHBLK1687X

MATHBLK1688X

Графики нескольких полиномов Лежандра изображены на рис. 35.3. Теперь
ясно, что ограниченные решения уравнений

MATHBLK1689X

{[} 0 \textless{} r \textless{} 1, 0 \textbackslash le
\textbackslash phi \textbackslash le \textbackslash pi, {]}

имеют вид

MATHBLK1690X,

{[} \textbackslash Phi(\textbackslash varphi) = aP\_n(\textbackslash cos
\textbackslash varphi) . {]}

\begin{verbatim}
Функция   P_n(\cos\phi)  —эго полином Лежандра n-го порядка, в котором x заменен на   \cos\phi  . И наконец, ищем решение исходной задачи в виде ряда
\end{verbatim}

(35.3)

MATHBLK1691X

\begin{verbatim}
коэффициенты которого выбираем так, чтобы удовлетворить граничному условию   u\left(1,\,\ varphi\right)=g\left(\varphi\right)  . Подстановка решения (35-3)
\end{verbatim}

\pandocbounded{\includegraphics[keepaspectratio]{images/_page_267_Figure_9.jpeg}}

\begin{verbatim}
РИС 35 3. Полиномы Лежандра   P_{n}\left(x\right)  .
\end{verbatim}

в граничные условия дает

MATHBLK1692X

\begin{verbatim}
Если умножить обе части этого соотношения на   P_m (\cos \varphi) \sin \varphi   и проинтегрировать по   \varphi   от 0 до   \pi  , то получим
\end{verbatim}

MATHBLK1693X

MATHBLK1694X

MATHBLK1695X

MATHBLK1696X

{[} \textbackslash begin\{cases\} 0, \& n \textbackslash neq m,
\textbackslash{} \textbackslash frac\{2a\_\{m\}\}\{2m+1\}, \& n = m.
\textbackslash end\{cases\} {]}

MATHBLK1697X

MATHBLK1698X

Можно убедиться в том, что полиномы Лежандра ортогональны на отрезке
{[}-1, 1{]}. Следовательно,

(35.4)

MATHBLK1699X

Итак, решение задачи Дирихле (35.1) дается выражением

(35.5)

MATHBLK1700X

где коэффициенты \(a_n\) вычисляются по формулам (35.4).

Приведем теперь пример цилиндрически-симметричного потенциала.

\begin{verbatim}
Цилиндрически-симметричный потенциал (не зависящий от   \theta  )
\end{verbatim}

Предположим, что температура поверхности сферы определяется формулой

MATHBLK1701X

\begin{verbatim}
и предположим, что мы хотим найти температуру внутри сферы. В нашей задаче температура постоянная на линиях постоянной широты (например, температура равна 2 на экваторе:   \varphi = \pi/2  , и Р И С
\end{verbatim}

\pandocbounded{\includegraphics[keepaspectratio]{images/_page_268_Figure_12.jpeg}}

. 35.4. Зависимость температуры от широты в радианах (цачало отсчета на
северном жолюсе).

\begin{verbatim}
нулю—на полюсах:   \phi = 0  ,   \phi < \pi  ). Чтобы найти u, нужно решить задачу
\end{verbatim}

(УЧП)

\begin{verbatim}
MATHBLK1702X,   0 < r < 1  ,

(ГУ)   u(1, \theta, \varphi) = 1 - \cos(2\varphi)  ,   0 \le \varphi \le \pi  .
\end{verbatim}

График температуры на поверхности изображен на рис. 35.4. Нам необходимо
найти коэффициенты \(a_n\) для решения (35.5). Чтобы сделать это, мы
должны

\begin{itemize}
\item
  а) либо воспользоваться формулой (35.4),
\item
  б) либо воспользоваться программой, существующей на вашем
  вычислительном центре и предназначенной для вычисления коэффициентов
  разложения по полиномам Лежандра,
\item
  в) либо проявить немного здравого смысла.
\end{itemize}

\pandocbounded{\includegraphics[keepaspectratio]{images/_page_269_Figure_5.jpeg}}

РИС. 35.5. Зависимость температуры от расстояния до центра шара.

Попытаемся воспользоваться последним рецептом. Рассмотрим
тригонометрическое тождество

MATHBLK1703X

которое позволяет переписать температуру на границе в виде

MATHBLK1704X

MATHBLK1705X

MATHBLK1706X

\begin{verbatim}
Мы получили разложение функции   g(\phi)   в ряд по полиномам Лежандра. Теперь решение задачи легко находится
\end{verbatim}

MATHBLK1707X

MATHBLK1708X

Графики этого решения на различных широтах приведены на рис. 35.5.

\section{ЗАМЕЧАНИЯ}\label{ux437ux430ux43cux435ux447ux430ux43dux438ux44f-14}

\begin{verbatim}
1. Обратим внимание на то, что после разложения функции   g(\phi)   в ряд
\end{verbatim}

MATHBLK1709X

нужно только умножить n-й член ряда на r\^{}n , чтобы получить

решение

MATHBLK1710X

\begin{enumerate}
\def\labelenumi{\arabic{enumi}.}
\setcounter{enumi}{1}
\tightlist
\item
  Решением внешней задачи Дирихле
\end{enumerate}

(УЧП)

\begin{verbatim}
MATHBLK1711X,   1 < r < \infty  ,   (\Gamma Y)     u(1, \theta, \varphi) = g(\varphi)
\end{verbatim}

является функция

MATHBLK1712X

где

MATHBLK1713X

\begin{verbatim}
Например, если граничное условие   g(\phi) = 3  , то решение   u(r, \phi) = 3/r  . Отметим, что в этой задаче решение стремится к нулю при   r \to \infty  , а в случае двух переменных решение внешней задачи с постоянным граничным условием само будет постоянным.
\end{verbatim}

\section{ЗАДАЧИ}\label{ux437ux430ux434ux430ux447ux438-14}

\begin{verbatim}
1. Подставьте   R(r) = r^{\alpha}   в уравнение Эйлера
\end{verbatim}

MATHBLK1714X

\begin{verbatim}
и получите   \alpha = n   и   \alpha = -(n+1)  .

2. Сделайте замену переменной   x = \cos \varphi   в уравнении Лежандра   [\sin \varphi \Phi']' + n(n+1)\sin \varphi \Phi = 0, \quad 0 \le \varphi \le \pi.
\end{verbatim}

и приведите его к виду

MATHBLK1715X

\begin{enumerate}
\def\labelenumi{\arabic{enumi}.}
\setcounter{enumi}{2}
\tightlist
\item
  Проверьте формулу Родрига
\end{enumerate}

MATHBLK1716X

для полиномов Лежандра.

\begin{enumerate}
\def\labelenumi{\arabic{enumi}.}
\setcounter{enumi}{3}
\tightlist
\item
  Решите внутреннюю задачу Дирихле в шаре
\end{enumerate}

(УЧП)

\begin{verbatim}
MATHBLK1717X,   0 < r < 1  ,   0 \le \varphi \le \pi  ,   0 \le \theta < 2\pi MATH
\end{verbatim}

INL1448X (\textbackslash Gamma Y) u(1, \textbackslash varphi) =
\textbackslash cos(3\textbackslash varphi) .

Используя тригонометрию, попытайтесь выразить \(\cos{(3\phi)}\) через
\(\cos{\phi}, \, \cos^2{\phi}, \, \cos^3{\phi}, \, a\) затем
скомбинируйте их так, чтобы по-

лучилось разложение

MATHBLK1718X

\begin{enumerate}
\def\labelenumi{\arabic{enumi}.}
\setcounter{enumi}{4}
\tightlist
\item
  Решите задачу Дирихле в шаре
\end{enumerate}

MATHBLK1719X

MATHBLK1720X

MATHBLK1721X

MATHBLK1722X

В этой задаче верхняя полусфера нагрета (+1), а нижняя охлаждена (-1).

\begin{enumerate}
\def\labelenumi{\arabic{enumi}.}
\setcounter{enumi}{5}
\tightlist
\item
  Найдите решение задачи Дирихле в шаре
\end{enumerate}

(УЧП)

MATHBLK1723X, \(1 < r < \infty\) , \(0 \le \varphi \le \pi\) ,
\(0 \le \theta < 2\pi\) , \((\Gamma Y)\)
\(u(1, \varphi) = 1 + \cos \varphi\) .

Проверьте ответ.

Лекция 36

\section{НЕОДНОРОДНАЯ ЗАДАЧА ДИРИХЛЕ (ФУНКЦИЯ
ГРИНА)}\label{ux43dux435ux43eux434ux43dux43eux440ux43eux434ux43dux430ux44f-ux437ux430ux434ux430ux447ux430-ux434ux438ux440ux438ux445ux43bux435-ux444ux443ux43dux43aux446ux438ux44f-ux433ux440ux438ux43dux430}

ЦЕЛЬ ЛЕКЦИИ: Показать, как с помощью функции Грина (функции источника)
можно решить неоднородную задачу Дирихле. В этом важном методе правая
часть уравнения рассматривается как некоторое входное воздействие и
разлагается в непрерывную совокупность дельтаобразных источников,
распределенных по некоторой заданной области. Затем находится отклик
системы на каждый такой источник (функция Грина) и все отклики
суммируются (интегрируются). В результате получаем полное решение
задачи.

Достаточно общей задачей прикладной математики является определение
потенциала в некоторой области пространства, если задано распределение
источников \(f(x, y)\) енутри этой области. В электростатике потенциал
(в вольтах) в области D обусловлен распределением зарядов с плотностью
\(f(x, y)\) по этой области. Типичная задача --- определение потенциала
в круге в двумерном случае.

В такой задаче потенциал должен удовлетворять уравнению

Пуассона (например, с нулевыми граничными условиями)

(36.1) (УЧП)

MATHBLK1724X

(ГУ) \(u(1, \theta) = 0, \quad 0 \le \theta \le 2\pi.\)

Мы выбрали нулевые условия не случайно. Если бы мы хотели решить задачу
в общем случае, с ненулевыми граничными условиями и неоднородным
уравнением, то вклад неоднородных граничных условий можно было бы найти
по интегральной формуле Пуассона из лекции 33.

Чтобы получить хоть какое-то представление о неоднородных
дифференциальных уравнениях, давайте рассмотрим графически решение
следующей задачи для уравнения Пуассона:

(УЧП)

MATHBLK1725X, \(0 < r < 1\) , \(0 \le \theta \le 2\pi\) ( \(q\)
---положительная константа),

(HY)

MATHBLK1726X

На границе потенциал (или температура, если вам это больше нравится)
поддерживается равным нулю, а лапласиан равен --- q во всех точках
области. Поскольку \(\Delta u\left(p\right)\) служит мерой разности
между функцией \(u\left(p\right)\) и ее средним значением, уравнение
Пуассона утверждает, что поверхность \(u\left(r\right)\) всегда выпукла
вниз. Другими словами, она будет похожа на тонкую мембрану с
закрепленным краем, которую сверху вниз обдувает поток воздуха.

Если правая часть уравнения \(f(x, y)\) изменяется от точки к точке, то
и выпуклость также будет переменной функцией точки области.

Перейдем теперь к главной части нашей лекции: определению функции Грина
и решению уравнения (36.1).

Сначала, однако, мы должны ввести понятие потенциала точеных источников
и стоков.

\section{Потенциалы точечных источников и
стоков}\label{ux43fux43eux442ux435ux43dux446ux438ux430ux43bux44b-ux442ux43eux447ux435ux447ux43dux44bux445-ux438ux441ux442ux43eux447ux43dux438ux43aux43eux432-ux438-ux441ux442ux43eux43aux43eux432}

Для решения неоднородного линейного уравнения достаточно решить это
уравнение с точечным источником, поскольку решение задачи в общем случае
можно найти, суммируя вклады от точечных источников. Наша задача ---
определить потенциал, который создает в заданной области точечный
источник (или сток). Этим источникам можно дать различную интерпретацию.
В теории теплопроводности мы можем рассматривать источник как точку, в
которой возникает тепло, а сток --- как точку, где тепло поглощается. В
электростатике источник может рассматриваться как отдельный
положительный заряд (протон), а сток --- как отдельный отрицательный
заряд (электрон). В любом случае независимо от интерпретации мы найдем
двумерный потенциал и,

создаваемый единичным точечным источником (поиск трехмерного потенциала
мы оставляем читателю в качестве упражнения).

Предположим, что отдельный точечный источник величины +q помещен в
начало координат Ясно, что тепло (или что-то другое)

будет вытекать из области вдоль радиальных линий, и, следовательно, для
везничны полного потока, протекающего через окружность радиуса r можно
записать

Полный поток, вытекающий через окруж-

HOCTD

MATHBLK1727X

Но полный поток должен совпадать с количеством тепла, созданного внутри
круга (закон сохранения), т. е.

MATHBLK1728X

Решая это диф реренциальное уравнение относительно функции \(u(r)\),
получаем Р И С

\pandocbounded{\includegraphics[keepaspectratio]{images/_page_273_Picture_10.jpeg}}

36 1 Радиальный поток тепла от точечного источника; a ---тепло
выделяется в начале координат.

MATHBLK1729X

(Решение изображено на рис. 362)

В электростатике разность потенциа тов \(u(B)\) - \(u(A)\) равна работе,
которую необходимо совершить, чтобы переместить единичный положительный
заряд из точки A в точку B (рис. 36.2).

\pandocbounded{\includegraphics[keepaspectratio]{images/_page_273_Figure_15.jpeg}}

РИС 36 2 Потенциал точечного заряда в двумерном случае. \(\alpha\) ---
работа, необходимая для перемещения единичного заряда из A в B

Сток можно интерпретировать как отрицательный источник, так что сток
величины -q создает потенциал

MATHBLK1730X

Этим завершается обсуждение потенциалов точечных зарядов. Теперь мы
готовы к решению неоднородного уравнения методом функции Грина.

\section{Уравнение Пуассона в
круге}\label{ux443ux440ux430ux432ux43dux435ux43dux438ux435-ux43fux443ux430ux441ux441ux43eux43dux430-ux432-ux43aux440ux443ux433ux435}

Займемся решением очень важной задачи

(36.2)

MATHBLK1731X
\(u_{rr} + \frac{1}{r} u_r + \frac{1}{r^2} u_{\theta\theta} = f(r, \theta), \quad 0 < r < 1, \quad 0 \le \theta \le 2\pi,\)
\((\Gamma Y)\) \(u(1, \theta) = 0, \quad 0 \le \theta \ le 2\pi.\)

\textbf{М}етод функции Грина (метод функции источника) состоит из двух
шагов.

\begin{enumerate}
\def\labelenumi{\arabic{enumi}.}
\item
  Помещаем в точку \((\rho, \varphi)\) точечный единичный заряд и
  находим потенциал \(G(r, \theta, \rho, \varphi)\) , который создается
  этим зарядом и обращается в нуль на границе.
\item
  Суммируем индивидуальные отклики \(G(r, \theta, \rho, \phi)\) ,
  взвешенные правой частью (плотностью зарядов) \(f(r, \theta)\) , по
  всему кругу. В результате получаем искомое решение
\end{enumerate}

MATHBLK1732X

Найдем теперь функцию источника для нашей задачи.

\section{\texorpdfstring{Определение функции источника
\(G(r, \theta, \rho, \phi)\)}{Определение функции источника G(r, \textbackslash theta, \textbackslash rho, \textbackslash phi)}}\label{ux43eux43fux440ux435ux434ux435ux43bux435ux43dux438ux435-ux444ux443ux43dux43aux446ux438ux438-ux438ux441ux442ux43eux447ux43dux438ux43aux430-gr-theta-rho-phi}

Заменим правую часть \(f(r,\theta)\) нашего уравнения точечным
источником величины 1, помещенным в некоторую произвольно выбранную
точку \((\rho, \varphi)\) . Математически точечный источник
представляется дельта-функцией \(\delta(r-\rho, \theta-\varphi)^1\) ).
Мы считаем, что дельта-функция равна нулю всюду, кроме точки
\((\rho, \varphi)\) , в которой находится единичный заряд. На языке сил
дельта-функцию можно интерпретировать как точечную единичную силу,
приложенную к точке \((\rho, \varphi)\) . Наша задача---найти потенциал
точечного заряда, если известно, что на границе он равен нулю. Функция,
удовлетворяющая этим условиям, называется функцией Грина или функцией
источника. Она равна \emph{отклику} системы в точке \((r, \theta)\) на
\emph{источник}, помещенный в точке \((\rho, \varphi)\) . Трудность
рещения этой задачи в том, что искомый потенциал должен обращаться в
нуль на границе. Если бы не это условие, то потенциал можно было бы
легко найти, так как мы уже знаем, что

См.: Владимиров В. С. Обобщенные функции в математической физике --- М.:
Наука, 1976. --- Прим ред.

потенциал точечного заряда равен

MATHBLK1733X,

где r --- расстояние от заряда до точки наблюдения.

Искомый потенциал может иметь следующий физический смысл:

\begin{enumerate}
\def\labelenumi{\arabic{enumi}.}
\tightlist
\item
  В теории теплопроводности --- это равновесная температура внутри
  круга, если тепловой источник размещается в точке (о. ф),
\end{enumerate}

а на границе поддерживается ну-

левая температура.

\begin{itemize}
\tightlist
\item
  \begin{enumerate}
  \def\labelenumi{\arabic{enumi}.}
  \setcounter{enumi}{1}
  \tightlist
  \item
    Он описывает отклонение мембраны от положения равновесия, если
    отклонения на границе равны нулю, а в точке (р. ф) мембрана отгянута
    на очень большую высоту.
  \end{enumerate}
\item
  \begin{enumerate}
  \def\labelenumi{\arabic{enumi}.}
  \setcounter{enumi}{2}
  \tightlist
  \item
    В электростатике это распределение потенциала внутри круга, если
    точечный положительный заряд помещен в точку \((\rho, \phi)\) , а
    круг заземлен и его потенциал поддерживается равным нулю.
  \end{enumerate}
\end{itemize}

\pandocbounded{\includegraphics[keepaspectratio]{images/_page_275_Picture_11.jpeg}}

Единичная окружность РИС. 36 3. Функция Гряна 6 (г. 0. р, ф) точечного
источника, расположенного в точке (р, ф).

Теперь найдем функцию Грина; она будет похожа на функцию, изображенную
на рис. 36.3.

\section{Построение
решения}\label{ux43fux43eux441ux442ux440ux43eux435ux43dux438ux435-ux440ux435ux448ux435ux43dux438ux44f}

ШАГ 1. Поскольку функция

MATHBLK1734X

является потенциалом поля в точке \(p = (r, \varphi) \theta)\) , которое
создано единичным точечным зарядом, помещенным в точку
\(Q = (\rho, \phi)\) , \((R-paccmonnue\ between\ points\ P\ and\ Q)\)
единственное, что нам осталось сделать---так изменить эту функцию, чтобы
она обратилась в нуль на границе.

ЩАГ 2. Из эксперимента известно, что эквипотенциальные линии точечных
зарядов (положительных или отрицательных) представляют собой окружности
(рис. 36.4).

Основная идея построения функции Грина состоит в следующем: нужно
снаружи круга разместить другой (отрицательный) заряд таким образом,
чтобы полный потенциал на окружности r=1 был константой. За гем можно
вычесть из потенциала эту константу и получить нулевой потенциал на
границе. Очевидно. найденный таким образом потенциал будет обладать
всеми свойствами функции \(G(r, \theta, \rho, \phi)\) . Остается
неясным, в какой точке вне круга нужно разместить отрицательный заряд,
чтобы потенциал на окружности стал постоянным. Можно сравнительно Р И С

\pandocbounded{\includegraphics[keepaspectratio]{images/_page_276_Picture_3.jpeg}}

. 36 4. Потенциал поля, образованного двумя противоположно заряженными
частицами; a, b--- линии (окружности) постоянного потенциала.

\pandocbounded{\includegraphics[keepaspectratio]{images/_page_276_Picture_5.jpeg}}

\textbf{РИС.} 36.5. Заряды Q в \(\overline{Q}\) дают постоянный
потенциал на окружности r=1.

легко убедиться в том, что если отрицательный заряд поместить в точку
\(\bar{Q} = (\bar{\rho}, \bar{\phi}) = (1/\rho, \phi)\) , то полный
потенциал

MATHBLK1735X

будет постоянным на окружности r=1 Здесь через R и \(\overline{R}\)
обозначены расстояния от зарядов до точки (г, в) (см. рис. 36.5). Можно
показать, что величина потенциала на окружности r=1определяется
выражением

MATHBLK1736X

(положительная константа).

С учетом этого результата строим функцию Грина

(36.3)

MATHBLK1737X

Потенциал положентельного заряда, расположенного в тенциал границы
постоянного в тенциал границы постоянного в тенциал границы постоянного
в тенциал границы постоянного в тенциал границы постоянного в тенциал
границы постоянного в тенциал границы постоянного в тенциал границы
постоянного в тенциал границы постоянного в тенциал границы постоянного
в тенциал границы постоянного в тенциал границы постоянного в тенциал
границы постоянного в тенциал границы постоянного в тенциал границы
постоянного в тенциал границы постоянного в тенциал границы постоянного
в тенциал границы постоянного в тенциал границы постоянного в тенциал
границы постоянного в тенциал границы постоянного в тенциал границы
постоянного в тенциал границы постоянного в тенциал границы постоянного
в тенциал границы постоянного в тенциал границы постоянного в тенциал
границы постоянного в тенциал границы постоянного в тенциал границы
постоянного в тенциал границы постоянного в тенциал границы постоянного
в тенциал границы постоянного в тенциал границы постоянного в тенциал
границы постоянного в тенциал границы постоянного в тенциал границы
постоянного в тенциал границы постоянного в тенциал границы постоянного
в тенциал границы постоянного в тенциал границы постоянного в тенциал
границы постоянного в тенциал границы постоянного в тенциал границы
постоянного в тенциал границы постоянного в тенциал границы постоянного
в тенциал границы постоянного в тенциал границы постоянного в тенциал
границы постоянного в тенциал границы постоянного в тенциал границы
постоянного в тенциал границы постоянного в тенциал границы постоянного
в тенциал границы постоянного в тенциал границы постоянного в тенциал
границы постоянного в тенциал границы постоянного в тенциал границы
постоянного в тенциал границы постоянного в тенциал границы постоянного
в тенциал границы постоянного в тенциал границы постоянного в тенциал
границы постоян границы постоян границы постоян границы постоян границы
постоян

где

MATHBLK1738X

MATHBLK1739X

формулы, выражающие расстояние между двумя точками в полярных
координатах. Чтобы найти решение исходной задачи, осталось провести
наложение всех импульсных откликов системы. Итак, мы переходим к
последнему шагу.

ШАГ 3. Суперпозиция импульсных откликов Этот шаг совсем прост. Из
соотношения

MATHBLK1740X

получаем

(36.4)

MATHBLK1741X

Формула (36.4) выражает решение задачи Дирихле для уравнения Пуассона
внутри круга через функцию Грина этой задачи. Если плотность зарядов
\(f(r, \theta)\) известна, то интеграл можно вычислить, например
численно.

\section{ЗАМЕЧАНИЯ}\label{ux437ux430ux43cux435ux447ux430ux43dux438ux44f-15}

\begin{enumerate}
\def\labelenumi{\arabic{enumi}.}
\tightlist
\item
  Задачу
\end{enumerate}

MATHBLK1742X

MATHBLK1743X

MATHBLK1744X

тоже можно решить \emph{методом функции Грина}. В этом случае решение
записывается в виде

MATHBLK1745X

Последнему соотношению можно придать форму более удобную для вычислений,
если раскрыть в явном виде выражение для \(\partial G/\partial r\) . В
результате получим интегральную формулу Пуассона, найденную нами ранее в
лекции 33.

(36.5)

MATHBLK1746X

\begin{enumerate}
\def\labelenumi{\arabic{enumi}.}
\setcounter{enumi}{1}
\tightlist
\item
  Решение общей задачи Дирихле
\end{enumerate}

MATHBLK1747X

MATHBLK1748X

MATHBLK1749X

находится как сумма решений (36.4) и (36.5).

\begin{enumerate}
\def\labelenumi{\arabic{enumi}.}
\setcounter{enumi}{2}
\item
  Метод функций Грина позволяет получить решение многил задач в областях
  различной формы. Однако для каждой области и (точнее, для каждого
  оператора, стоящего в левой части граничного условия) для каждого
  уравнения нужно находить свою функцию Грина, а это не всегда легко
  сделать.
\item
  Чтобы реально найти решение по формуле (36.4), в большинстве случаев
  необходимо вычислять интегралы на ЭВМ.
\end{enumerate}

\section{ЗАДАЧИ}\label{ux437ux430ux434ux430ux447ux438-15}

Найти потенциал точечного источника в трехмерном пространстве.

\begin{enumerate}
\def\labelenumi{\arabic{enumi}.}
\setcounter{enumi}{1}
\tightlist
\item
  Найти функцию Грина \(G(x, y, \xi, \eta)\) для уравнения Лапласа в
  верхней полуплоскости y \textgreater{} 0. Другими словами, найти
  потенциал в точке (x, y) в верхней полуплоскости, если заряд
  расположен в точке \((\xi, \eta)\) , а потенциал на границе y = 0
  равен нулю (см. рис.).
\end{enumerate}

\pandocbounded{\includegraphics[keepaspectratio]{images/_page_278_Figure_11.jpeg}}

УКАЗАНИЕ. Если отрицательный заряд поместить в точку

\(\overline{Q} = (\xi, -\eta)\) , то ясво, что полный потенциал на линии
y = 0 будет равен нулю. Следовательно, функция Грина будет совпадать с
потенциалом этих двух зарядов.

\begin{itemize}
\tightlist
\item
  \begin{enumerate}
  \def\labelenumi{\arabic{enumi}.}
  \setcounter{enumi}{2}
  \tightlist
  \item
    Используя результаты задачи 2, найдите решение уравнения Пуассона
    \(\Delta u = -k\) в верхней полуплоскости с нулевым граничным
    условием.
  \end{enumerate}
\item
  \begin{enumerate}
  \def\labelenumi{\arabic{enumi}.}
  \setcounter{enumi}{3}
  \tightlist
  \item
    Как бы вы строили функцию Грина для первого квадранта \(x>0,\ y>0\)
    ?
  \end{enumerate}
\item
  \begin{enumerate}
  \def\labelenumi{\arabic{enumi}.}
  \setcounter{enumi}{4}
  \tightlist
  \item
    Другой подход к решению уравнения Пуассона состоит в следующем.
    Предположим, что вы хотите решить задачу
  \end{enumerate}
\end{itemize}

(УЧП)

MATHBLK1750X, \(0 < r < 1\) , \(0 \le \theta < 2\pi\) , \((\Gamma Y)\)
\(u(1, \theta) = \sin \theta\) , \(0 \le \theta \le 2\pi\) .

Сначала попытайтесь найти частное решение уравнения с помощью
подстановки

MATHBLK1751X.

Затем подставьте это выражение в уравнение Пуассона и определите
константу \(A_{\cdot}\)

После того как частное решение \(u_p\) найдено, надо искать полное
решение в виде \(u=w+u_p\) . Какому граничному условию должна
удовлетворять функция w? Найдите \(w(r, \theta)\) . Что является
решением исходной задачи \(u(r, \theta)\) ? Проверьте это. Тщательно
проанализируйте ответ. Дайте интерпретацию кажлого члена.

\section{Часть 5. ЧИСЛЕННЫЕ И ПРИБЛИЖЕННЫЕ
МЕТОДЫ}\label{ux447ux430ux441ux442ux44c-5.-ux447ux438ux441ux43bux435ux43dux43dux44bux435-ux438-ux43fux440ux438ux431ux43bux438ux436ux435ux43dux43dux44bux435-ux43cux435ux442ux43eux434ux44b}

Лекция 37

\section{ЧИСЛЕННЫЕ РЕШЕНИЯ (ЭЛЛИПТИЧЕСКИЕ
ЗАДАЧИ)}\label{ux447ux438ux441ux43bux435ux43dux43dux44bux435-ux440ux435ux448ux435ux43dux438ux44f-ux44dux43bux43bux438ux43fux442ux438ux447ux435ux441ux43aux438ux435-ux437ux430ux434ux430ux447ux438}

ЦЕЛЬ ЛЕКЦИИ: Показать, как уравнения в частных производных можно свести
к системам алгебраических уравнений, если частные производные заменить
их конечно-разностными аппроксимациями. Решение этой системы
алгебраических уравнений итерационным методом позволяет получить
приближенное решение

Дается информация о пакете прикладных программ ELLPACK, предназначенном
для решения эллиптических

задач общего вида на ЭВМ.

К настоящему моменту мы познакомились с несколькими методами решения
линейных уравнений с частными производными. Однако большинство
уравнений, с которыми мы имели дело, были очень простыми. Так же просты
были и формы областей, в которых решались задачи. Тем не менее многие
задачи нельзя упростить до такой степени, чтобы они свелись к заданному
набору шаблонных задач. Такие задачи следует решать приближенно
численными методами. Прогресс в области высокопроизводительных ЭВМ
привел к созданию новых численных методов. Сегодня уже решены такие
нелинейные задачи динамики жидкости, теории упругости и теории
потенциала, о решении которых десять лет назад никто даже не
задумывался.

Под общим названием «численные методы» объединены несколько различных
подходов к решению задач. Подробное описание всех этих подходов можно
найти в книге {[}4{]} из рекомендованного списка литературы. В этой и
двух последующих лекциях будет показано, как решать эллиптические,
гиперболические и параболические уравнения методом конечных разностей.

Сначала познакомимся с понятием конечных разностей, а затем покажем, как
с их помощью решить задачу Дирихле в квадрате.

\section{Конечно-разностные
аппроксимации}\label{ux43aux43eux43dux435ux447ux43dux43e-ux440ux430ux437ux43dux43eux441ux442ux43dux44bux435-ux430ux43fux43fux440ux43eux43aux441ux438ux43cux430ux446ux438ux438}

Вспомним ряд Тейлора для функции \(f(x)\):

MATHBLK1752X

Если оборвать этот ряд на втором члене, то получим

MATHBLK1753X.

Откуда

MATHBLK1754X

Выражение, стоящее в правой части, называется правой разностной
производной. Она аппроксимирует первую производную f'(x) в точке x.

В разложении Тейлора функции \(f(\mathbf{r})\) можно заменить h на -h и
получить левую разностную производную

(37.2)

MATHBLK1755X.

Buyutan \(f(x - h) \cong f(x) - f'(x) h\) \textbackslash nif
\(f(x + h) \cong f( x) + f'(x) h\) ,

получаем центральную разностную производную

(37.3)

MATHBLK1756X

Если в ряде Тейлора оставить на одно слагаемое больше, то совершенно
аналогично можно получить центральную разностную производную для
аппроксимации f'\,'(x)

(37.4)

MATHBLK1757X

Теперь можно распространить понятие конечно-разностпой аппроксимации на
\emph{частные производные}. Если исходить из разложения Тейлора функции
двух переменных

MATHBLK1758X

MATHBLK1759X

можно получить следующие аппроксимации частных производных:

MATHBLK1760X

MATHBLK1761X

{[} \textbackslash begin\{split\} u\_x(x,~y) \& \textbackslash cong
\textbackslash frac\{u(x+h,~y)-u(x,~y)\}\{h\}, \textbackslash{}
u\_\{xx\}(x,~y) \& \textbackslash cong
\textbackslash frac\{1\}\{h\^{}2\} \textbackslash big{[}
u(x+h,~y)-2u(x,~y)+u(x-h,y) \textbackslash big{]}, \textbackslash{}
u\_u(x,~y) \& \textbackslash cong
\textbackslash frac\{u(x,~y+k)-u(x,~y)\}\{k\}, \textbackslash{}
u\_\{uy\}(x,~y) \& \textbackslash cong
\textbackslash frac\{1\}\{h\^{}2\} \textbackslash big{[}
u(x,~y+k)-2u(x,~y)+u(x,~y-k) \textbackslash big{]}.
\textbackslash end\{split\} {]}

MATHBLK1762X

MATHBLK1763X

В этих формулах частные производные аппроксимируются правыми,
центральными и левыми разностными производными, но в данной лекции мы
будем пользоваться только центральными аппроксимациями.

\pandocbounded{\includegraphics[keepaspectratio]{images/_page_282_Figure_5.jpeg}}

РИС. 37.1 Сетка для задачи Дирихле внутри квадрата.

Чтобы познакомиться с основными правилами использования
конечно-разностных аппроксимаций, рассмогрим простую задачу Дирихле.

Решение задачи Дирихле методом конечных разностей

Пусть необходимо решить задачу Дирихле

(ЗЧП)

MATHBLK1764X, \(0 < x < 1\) , \(0 < y < 1\) , \(u = 0\) , на верхней и
боковых сторонах квадрата, \(u(x, 0) = \sin(x\pi)\) ,
\(0 \le \le x \le 1\) , \(y = 0\) .

Построим в плоскости (x, y) сетку, как показано на рис. 37.1. Удобно
(особенно, если мы собираемся применять ЭВМ)

непользовать следующие обозначения:

MATHBLK1765X

MATHBLK1766X

{[} \textbackslash begin\{split\}
u\textbackslash left(x,,y\textbackslash right) \&=
u\_\{ij\},,\textbackslash{}
u\textbackslash left(x,,y+k\textbackslash right) \&=
u\_\{i+1,j\},,\textbackslash{}
u\textbackslash left(x,,y-k\textbackslash right) \&=
u\_\{i-1,j\},,\textbackslash{}
u\textbackslash left(x-h,,y\textbackslash right) \&=
u\_\{i,j-1\},,\textbackslash{}
u\textbackslash left(x+h,,y\textbackslash right) \&=
u\_\{i,j+1\},,\textbackslash{}
u\_\{x\}\textbackslash left(x,,y\textbackslash right) \&=
\textbackslash frac\{1\}\{2h\}\textbackslash left(u\_\{i,j+1\}-u\_\{i,j-1\}\textbackslash right),,\textbackslash{}
u\_\{y\}\textbackslash left(x,,y\textbackslash right) \&=
\textbackslash frac\{1\}\{2k\}\textbackslash left(u\_\{i+1,j\}-u\_\{i-1,j\}\textbackslash right),,\textbackslash{}
u\_\{xx\}\textbackslash left(x,,y\textbackslash right) \&=
\textbackslash frac\{1\}\{h\^{}\{2\}\}\textbackslash left(u\_\{i,j+1\}-2u\_\{i,j\}+u\_\{i,j-1\}\textbackslash right),,\textbackslash{}
u\_\{yy\}\textbackslash left(x,,y\textbackslash right) \&=
\textbackslash frac\{1\}\{h\^{}\{2\}\}\textbackslash left(u\_\{i+1,j\}-2u\_\{i,j\}+u\_\{i-1,j\}\textbackslash right),.
\textbackslash end\{split\} {]}

MATHBLK1767X

MATHBLK1768X

Основная идея решения задачи Дирихле основана на замене частных
производных в уравнении Лапласа

MATHBLK1769X

их конечно-разностными аппроксимациями. Проделав это и воспользовавшись
компактными обозначениями, получаем следующее разностное уравнение:

MATHBLK1770X

В случае когда h=k, уравнение Лапласа приводится к виду (37.6)
\((u_{i+1,j}+u_{i-1,j}+u_{i,j+1}+u_{i,j+1}-4u_{ij})=0.\)

Разрешив его относительно \(u_{II}\) , получим

MATHBLK1771X

Отметим, что в этом соотношении все значения \(u_{tj}\) берутся для
внутренних узлов сетки. Согласно последнему соотношению, решение
\(u_{ij}\) аппроксимируется средним значением решения по четырем
соседним точкам. Следовательно, мы можем разработать некоторый численный
метод решения этой задачи.

Алгоритм численного решения задачи Дирихле (метод Либманна)

ШАГ 1. Присвоим величинам \(u_{ij}\) во внутренних узлах численное
вначение, равное среднему всех значений граничных условий.

ШАГ 2. Будем пересчитывать значения во всех внутренних точках сетки
путем замены старого значения средним по четырем соседним точкам. Не
очень важно, как будет организован процесс счета, но обычно его проводят
по строчкам (либо по столбцам). После нескольких операций процесс
сойдется к приближенному решению задачи. Скорость сходимости этого
алгоритма невелика, но ее можно увеличить несколькими способами.
Заинтересованному читателю мы рекомендуем обратиться к книге {[}4{]} из
списка рекомендованной литературы.

\section{ЗАМЕЧАНИЯ}\label{ux437ux430ux43cux435ux447ux430ux43dux438ux44f-16}

\begin{enumerate}
\def\labelenumi{\arabic{enumi}.}
\tightlist
\item
  Если взять сетку с чегырьмя внутренними точками (m=n=4), то система
  (37.6) запишется в виде
\end{enumerate}

MATHBLK1772X

MATHBLK1773X

MATHBLK1774X

MATHBLK1775X

Ее необходимо решить относительно \(u_{22}\) , \(u_{21}\) , \(u_{32}\) и
\(u_{33}\) . Решение такой системы можно найти \emph{итерационным
методом}. Метод Либманна также является одним из итерационных методов.

\begin{itemize}
\tightlist
\item
  \begin{enumerate}
  \def\labelenumi{\arabic{enumi}.}
  \setcounter{enumi}{1}
  \tightlist
  \item
    Если уменьшить шаги сетки h и k (число узлов сетки при этом будет
    возрастать), то мы получим систему уравнений типа (37.7), но больших
    размеров В общем случае число уравнений всегда совпадает с числом
    внутренних узлов сетки.
  \end{enumerate}
\item
  \begin{enumerate}
  \def\labelenumi{\arabic{enumi}.}
  \setcounter{enumi}{2}
  \tightlist
  \item
    Систему уравнений (37.7) можно переписать в матричной форме
  \end{enumerate}
\end{itemize}

MATHBLK1776X

MATHBLK1777X

MATHBLK1778X

{[} \textbackslash begin\{bmatrix\} u\_\{22\} \textbackslash{} u\_\{23\}
\textbackslash{} u\_\{32\} \textbackslash{} u\_\{33\}
\textbackslash end\{bmatrix\} {]}

MATHBLK1779X

MATHBLK1780X

MATHBLK1781X

MATHBLK1782X

MATHBLK1783X.

MATHBLK1784X

{[} u\_\{xx\} + u\_\{yy\} = f(x, y) {]}

MATHBLK1785X

{[} xu\_\{xx\} + u\_\{yy\} + 2u = \textbackslash sin(x - y) {]}

MATHBLK1786X

{[} \textbackslash sin xu\_\{xx\} + u\_\{xy\} + 3u = 0 {]}

MATHBLK1787X

(УЧ\textbackslash Pi)

MATHBLK1788X

{[} u\_\{xx\} + u\_\{yy\} + 2u = 0 {]}

MATHBLK1789X

{[} u\_\{xx\} + u\_\{yy\} = 0 {]}

MATHBLK1790X

MATHBLK1791X

MATHBLK1792X

{[} (\textbackslash text\{УЧП\}) {]}

MATHBLK1793X

{[} u\_t = u\_\{xx\} {]}

MATHBLK1794X

{[} u\_t = u\_\{xx\}, \textbackslash quad 0 \textless{} x \textless{} 1,
\textbackslash quad 0 \textless{} t \textless{} \textbackslash infty,
{]}

MATHBLK1795X

MATHBLK1796X

{[} \textbackslash begin\{cases\} u(0, t) = 1, \textbackslash{} u\_x(1,
t) = -{[}u(1, t) - g(t){]}, \textbackslash{} u(x, 0) = 0,
\textbackslash quad 0 \textbackslash le x \textbackslash le 1.
\textbackslash end\{cases\} {]}

MATHBLK1797X

MATHBLK1798X

Для решения этой задачи методом конечных разностей построим
прямоугольную сетку, узлы которой определяются формулами (см. рис. 38.1)

MATHBLK1799X, \(j = 0, 1, 2, ..., n\) , \(y_i = ik\) ,
\(i = 0, 1, 2, ..., m\) .

Отметим, что значения \(u_{ij}\) на левой и нижней сторонах сетки на
рис. 38.1 известны из граничных и начальных условий и наша вадача
состоит в отыскании остальных значений \(u_{ij}\) . Для решения этой
задачи заменим частные производные в уравнении теплопроводности их
конечно-разностными аппроксимациями

MATHBLK1800X

MATHBLK1801X

Подставим эти выражения в уравнение \(u_t = u_{xx}\) и разрешим
получившееся уравнение относительно значений функции на верхнем
временном слое. В результате получаем

(38.2)

MATHBLK1802X

Это и есть искомая формула, поскольку она выражает решение в данный
момент времени через решение в предыдущий момент времени (индекс
\emph{i} относится к временной переменной). На рис. 38.1 выделены те
значения, которые входят в данную формулу.

Теперь можно приступить к вычислениям. Однако сначала нужно
аппроксимировать производную в граничном условии на правом конце

MATHBLK1803X

В результате аппроксимации получаем

(38.3)

MATHBLK1804X

где значения \(g_i = g(ik)\) известны. Здесь мы заменили \(u_x(1, t)\)
левой разностной производной, поскольку правая разностная производная
потребовала бы значений функции за пределами сетки. Из (38.3) находим

(38.4)

MATHBLK1805X

Формулы (38.2) и (38.4) позволяют начать вычисления.

Алгоритм вычислений по явной схеме

ШАГ 1. Находим решение на сеточном слое \(t = \Delta t\) , используя
явную формулу (см. рис. 38.2)

MATHBLK1806X

ШАГ 2. Величину \(u_{2,n}\) находим по формуле (38.4)

MATHBLK1807X

Совершив шаги 1 и 2, получаем решение для \(t=\Delta t\) . Для получения
решения при \(t=2\Delta t\) (вторая строка снизу на рис. 38.2) повторяем
шаги 1 и 2, поднявшись на одну строку вверх, т. е. увеличив t на 1 и
используя \(u_{t,j}\) с предыдущей строки. Аналогично вычисляется
решение в последующие моменты времени \(t=3\Delta t,\ 4\Delta t\dots\)

Для того чтобы помочь читателю в проведении расчетов по этому методу, на
рис. 38.3 приведена достаточно подробная блоксхема алгоритма. Те, кто не
очень хорошо знаком с блок-схемами.

\pandocbounded{\includegraphics[keepaspectratio]{images/_page_290_Figure_3.jpeg}}

РИС. 38.2. Иллюстрация явной разностной схемы: n--- число узлов сетки
вдоль оси \(x;\ m\) --- число узлов сетки вдоль оси
\(y;\ h=1/(m-1),\ k=T/(m-1)\)

могут счигать их промежуточным звеном между вычислительными алгоритмами
и подробными программами для ЭВМ. Блок-схемы дают строгое описание
процесса вычислений.

\section{ЗАМЕЧАНИЯ}\label{ux437ux430ux43cux435ux447ux430ux43dux438ux44f-17}

\begin{itemize}
\tightlist
\item
  У явной схемы есть серьезный недостаток. Если шаг по времени
  оказывается достаточно большим по сравнению с шагом по х, погрешности
  округления могут стать настолько большими, что полученное решение
  потеряет смысл. Отношение шагов по t и х зависит от уравнения и
  граничных условий, но в общем случае шаг по времени должен быть много
  меньше шага по координате. В книге {[}4{]} доказано, что для
  применимости явной схемы должно выполняться условие k/h² ≤ 0,5.
\item
  \begin{enumerate}
  \def\labelenumi{\arabic{enumi}.}
  \setcounter{enumi}{1}
  \tightlist
  \item
    Справедливо следующее эмпирическое правило: если уменьшать шаги
    \(\Delta t\) и \(\Delta x\) , то погрешность аппроксимации частных
    производных конечными разностями тоже будет уменьшаться, однако чем
    мельче сетка, тем больше вычислений необходимо совершить и,
    следовательно, тем больше будут погрешности округления. Это явление
    иллюстрируется на рис. 38.5.
  \end{enumerate}
\item
  \begin{enumerate}
  \def\labelenumi{\arabic{enumi}.}
  \setcounter{enumi}{2}
  \tightlist
  \item
    Для гиперболической задачи
  \end{enumerate}
\end{itemize}

(УЧП)

MATHBLK1808X

(ГУ)

MATHBLK1809X

MATHBLK1810X

{[} \textbackslash begin\{cases\} u(0, t) = g\_1(t), \textbackslash{}
u(1, t) = g\_2(t), \textbackslash end\{cases\} {]}

MATHBLK1811X

MATHBLK1812X

(HY)

MATHBLK1813X

MATHBLK1814X

{[} \textbackslash begin\{cases\} u(x, 0) = \textbackslash varphi(x),
\textbackslash{} u\_t(x, 0) = \textbackslash psi(x),
\textbackslash end\{cases\} {]}

MATHBLK1815X

MATHBLK1816X

также можно построить явную разностную схему Для этого аппроксимируем
\(u_{tt}\) и \(u_{xy}\) центральными разностными произ- Р И С

\pandocbounded{\includegraphics[keepaspectratio]{images/_page_291_Figure_4.jpeg}}

. 38.3. Алгоритм решения задачи по явной схеме.

водными

MATHBLK1817X

MATHBLK1818X

а начальное условие--- по схеме

MATHBLK1819X

В результате для вычисления величины \(u(x, t+k)\) получаем следующую
явную схему:

(38.5)

MATHBLK1820X

Из (38.5) видно, что для вычисления решения на данном времениом слое
необходимо знать решение на двух предыдущих слоях. Значит, для начала
счета необходимо воспользоваться Р И С

\pandocbounded{\includegraphics[keepaspectratio]{images/_page_292_Figure_6.jpeg}}

. 38.4. Использование массивов при реализации явной схемы.

начальным условнем на скорость

MATHBLK1821X

из которого получаем \(u(x, k) = \varphi(x) + k\psi(x)\) , т. е.
значение решения при \(t = \Delta t\) . Решение в последующие моменты
времени можно находить по явной формуле (38.5).

\section{ЗАДАЧИ}\label{ux437ux430ux434ux430ux447ux438-16}

\begin{enumerate}
\def\labelenumi{\arabic{enumi}.}
\tightlist
\item
  Постройте явную конечно-разностную схему для задачи
\end{enumerate}

(УЧП)

MATHBLK1822X, \(0 < x < 1\) , \(0 < t < \infty\) ,\\
(ГУ)

MATHBLK1823X

MATHBLK1824X

{[} \textbackslash begin\{cases\} u(0, t) = 0, \textbackslash{} u(1, t)
= 0, \textbackslash{} u(x, 0) = \textbackslash sin(\textbackslash pi x),
\textbackslash quad 0 \textbackslash le x \textbackslash le 1.
\textbackslash end\{cases\} {]}

MATHBLK1825X

MATHBLK1826X

MATHBLK1827X

\begin{verbatim}
\begin{array}{ccccccccccccccccccccccccccccccccccc
\end{verbatim}

MATHBLK1828X

{[} \textbackslash frac\{\textbackslash Delta t\}\{(\textbackslash Delta
x)\^{}2\} \textbackslash leqslant 0.5. {]}

MATHBLK1829X

\begin{verbatim}
\begin{array}{ccccccccccccccccccccccccccccccccccc
\end{verbatim}

MATHBLK1830X

{[} \textbackslash begin\{split\} u\_t(x,~t) \&=
\textbackslash frac\{1\}\{k\} \textbackslash big{[} u(x,~t+k) - u(x,~t)
\textbackslash big{]}, \textbackslash{} u\_\{xx\}(x,~t) \&=
\textbackslash frac\{\textbackslash lambda\}\{h\^{}2\}
\textbackslash big{[} u(x+h,~t+k) - 2u(x,~t+k) + u(x-h,~t+k)
\textbackslash big{]} + \textbackslash{} \&\textbackslash quad +
\textbackslash frac\{(1-\textbackslash lambda)\}\{h\^{}2\}
\textbackslash big{[} u(x+h,~t) - 2u(x,~t) + u(x-h,~t)
\textbackslash big{]}, \textbackslash end\{split\} {]}

MATHBLK1831X

{[} \textbackslash frac\{1\}\{k\} (u\_\{i+1, j\} - u\_\{i, j\}) =
\textbackslash frac\{\textbackslash lambda\}\{h\^{}2\} (u\_\{i+1, j+1\}
- 2u\_\{i+1, j\} + u\_\{i+i, j-1\}) +
\textbackslash frac\{(1-\textbackslash lambda)\}\{h\^{}2\} (u\_\{i,
j+1\} - 2u\_\{i, j\} + u\_\{i, j-1\}), {]}

MATHBLK1832X

(39.3)

MATHBLK1833X

{[} (\textbackslash Gamma Y) \textbackslash begin\{cases\} u\_\{i, 1\} =
0, \& i = 1, 2, \ldots, m, \textbackslash{} u\_\{i, n\} = 0, \& i = 1,
2, \ldots, m, \textbackslash{} u\_\{i, j\} = 1, \& j = 2, \ldots, n-1.
\textbackslash end\{cases\} {]}

MATHBLK1834X

{[} -\textbackslash lambda r u\_\{i+1, j+1\} +
(1+2r\textbackslash lambda) u\_\{i+1, j\} - \textbackslash lambda r
u\_\{i+1, j-1\} = = r(1-\textbackslash lambda) u\_\{i, j+1\} +
{[}1-2r(1-\textbackslash lambda){]} u\_\{i, j\} +
r(1-\textbackslash lambda) u\_\{i, j-1\}, {]}

MATHBLK1835X

{[} \textbackslash begin\{array\}\{l\}
-u\_\{21\}+3u\_\{22\}-u\_\{23\}=u\_\{11\}-u\_\{12\}+u\_\{13\}=1,\textbackslash{}
-u\_\{22\}+3u\_\{23\}-u\_\{24\}=u\_\{12\}-u\_\{13\}+u\_\{14\}=1,\textbackslash{}
-u\_\{23\}+3u\_\{24\}-u\_\{25\}=u\_\{13\}-u\_\{14\}+u\_\{15\}=1,\textbackslash{}
-u\_\{24\}+3u\_\{25\}-u\_\{26\}=u\_\{14\}-u\_\{15\}+u\_\{16\}=1.
\textbackslash end\{array\} {]}

MATHBLK1836X

MATHBLK1837X

MATHBLK1838X

MATHBLK1839X

{[} \textbackslash begin\{bmatrix\} u\_\{2\textbackslash bar\{z\}\}
\textbackslash{} u\_\{\textbackslash bar\{z\}a\} \textbackslash{}
u\_\{24\} \textbackslash{} u\_\{\textbackslash bar\{z\}x\}
\textbackslash end\{bmatrix\} {]}

MATHBLK1840X

MATHBLK1841X

MATHBLK1842X.

MATHBLK1843X

MATHBLK1844X

{[} \textbackslash begin\{bmatrix\} b\_1 \& c\_1 \& 0 \& 0 \&
\textbackslash dots \& \& \& 0 \textbackslash{} a\_1 \& b\_2 \& c\_2 \&
0 \& \textbackslash dots \& \& \& 0 \textbackslash{} 0 \&
a\_\{\textbackslash bar\{z\}\} \& b\_3 \& c\_3 \& \textbackslash dots \&
\& \& 0 \textbackslash{} \& \& \& \& \& \& \& \& \textbackslash{}
\textbackslash vdots \& \& \& \& \& \& \& \& \textbackslash{} 0 \& 0 \&
0 \& \textbackslash dots \& \& \& \& a\_\{n-1\} \& b\_n
\textbackslash end\{bmatrix\} {]}

MATHBLK1845X

MATHBLK1846X

{[} \textbackslash begin\{bmatrix\} x\_1 \textbackslash{}
x\_\{\textbackslash bar\{z\}\} \textbackslash{}
x\_\{\textbackslash bar\{s\}\} \textbackslash{} \textbackslash vdots
\textbackslash{} x\_n \textbackslash end\{bmatrix\} {]}

MATHBLK1847X

MATHBLK1848X

MATHBLK1849X

{[} \textbackslash begin\{bmatrix\} d\_t \textbackslash{}
d\_\{\textbackslash bar\{z\}\} \textbackslash{}
d\_\{\textbackslash bar\{s\}\} \textbackslash{} \textbackslash vdots
\textbackslash{} \textbackslash vdots \textbackslash{} d\_n
\textbackslash end\{bmatrix\} {]}

MATHBLK1850X,

MATHBLK1851X

MATHBLK1852X

{[} \textbackslash begin\{bmatrix\} 1 \& c\_1\^{}* \& 0 \& 0 \&
\textbackslash dots \& 0 \textbackslash{} 0 \& 1 \& c\_2\^{}* \& 0 \&
\textbackslash dots \& 0 \textbackslash{} 0 \& 0 \& 1 \& c\_2\^{}* \&
\textbackslash dots \& 0 \textbackslash{} \textbackslash vdots \&
\textbackslash vdots \& \& \& \textbackslash vdots \&
\textbackslash vdots \textbackslash{} 0 \& 0 \& \textbackslash dots \&
\& 1 \textbackslash end\{bmatrix\} {]}

MATHBLK1853X

MATHBLK1854X

{[} \textbackslash begin\{bmatrix\} x\_1 \textbackslash{} x\_2
\textbackslash{} x\_3 \textbackslash{} \textbackslash vdots
\textbackslash{} x\_n \textbackslash end\{bmatrix\} {]}

MATHBLK1855X

MATHBLK1856X

MATHBLK1857X

{[} \textbackslash begin\{bmatrix\} d\_1\^{}* \textbackslash{} d\_2\^{}*
\textbackslash{} d\_3\^{}* \textbackslash{} \textbackslash vdots
\textbackslash{} d\_n\^{}* \textbackslash end\{bmatrix\} {]}

MATHBLK1858X

MATHBLK1859X

где

MATHBLK1860X, \(c_{j+1}^* = \frac{c_{j+1}}{b_{j+1} - a_j c_j^*}\) ,
\(j = 1, 2, \ldots, n-2\) , и \(d_1^* = d_1/b_1\) ,
\(d_{j+1}^* = \frac{d_{j+1} - a_j d_j^*}{b_{j+1} - a_j c_j^*}\) ,
\(j = 1, 2, \ldots, n-1\) .

В этом преобразовании нет ничего таниственного, поскольку вторая система
полностью эквивалентна первой. Матрица новой системы устроена так, что
эта система очень легко решается. Решая уравнения последовательно снизу
вверх, получаем

MATHBLK1861X

Применяя этот метод к системе четырех уравнений (39.5), находим решение
\(u_{22}=0,60\) ; \(u_{23}=0,80\) ; \(u_{24}=0,80\) ; \(u_{23}=0,60\) .
Мы получили приближенные значения решения во внутренних точках сетки при
\(t=\Delta t\) . Теперь можно сделать следующий шаг по времени, но для
этого придется решить новую систему уравнений.

В неявной схеме объем вычислений на каждом шаге больше, чем в явной, но
хорошую точность можно получить даже при гораздо большем шаге.

\section{ЗАДАЧИ}\label{ux437ux430ux434ux430ux447ux438-17}

\begin{itemize}
\tightlist
\item
  \begin{enumerate}
  \def\labelenumi{\arabic{enumi}.}
  \tightlist
  \item
    Из дифференциального уравнения (39.3) получите уравнение (39.4).
  \end{enumerate}
\item
  \begin{enumerate}
  \def\labelenumi{\arabic{enumi}.}
  \setcounter{enumi}{1}
  \tightlist
  \item
    Постройте неявную конечно-разностную схему для задачи
  \end{enumerate}
\end{itemize}

MATHBLK1862X

MATHBLK1863X

\begin{verbatim}
\begin{array}{ll} (\mathrm{Y}\mathrm{H}\Pi) & u_t = u_{xx}, \quad 0 < x < 1, \quad 0 < t < \infty, \\ (\Gamma\mathrm{Y}) & \left\{ \begin{array}{ll} u_{\,0}(0, \ t) = 1, \\ u_{\,x}(1, \ t) + u_{\,}(1, \ t) = g_{\,}(t), \\ u_{\,x}(1, \ t) = 0, \quad 0 \leqslant \kappa \leqslant 1. \end{array}
\end{verbatim}

MATHBLK1864X

MATHBLK1865X

\begin{enumerate}
\def\labelenumi{\arabic{enumi}.}
\setcounter{enumi}{2}
\tightlist
\item
  Как с помощью неявной схемы решить задачу
\end{enumerate}

(УЧП)

MATHBLK1866X

(ГУ)

{[} \textbackslash begin\{cases\} u(0, t) = 0, \textbackslash{} u(1, t)
= 0, \textbackslash{} u(x, 0) = 1, \textbackslash quad 0
\textbackslash le x \textbackslash le 1 \textbackslash end\{cases\} {]}

\begin{itemize}
\tightlist
\item
  \begin{enumerate}
  \def\labelenumi{\arabic{enumi}.}
  \setcounter{enumi}{3}
  \tightlist
  \item
    Как выглядит шаблон для уравнения (39.4), когда \(\lambda = 1\) ?
  \end{enumerate}
\item
  \begin{enumerate}
  \def\labelenumi{\arabic{enumi}.}
  \setcounter{enumi}{4}
  \tightlist
  \item
    Постройте блок-схему решения уравнения теплопроводности (39.2). Если
    есть возможность, то напишите программу для ЭВМ Было бы очень
    полезно решить эту задачу численно при различных значениях λ с
    начальным условием и (x, 0) = sin (πx). Сравните полученные при
    различных λ численные результаты с аналитическим решением, которое в
    данном случае имеет вид
  \end{enumerate}
\end{itemize}

MATHBLK1867X

6 Воспользовавшись формулами, приведенными в тексте, решите систему
алгебраических уравнений (39.5).

Лекция 40

\section{СРАВНЕНИЕ АНАЛИТИЧЕСКИХ РЕШЕНИЙ С
ЧИСЛЕННЫМИ}\label{ux441ux440ux430ux432ux43dux435ux43dux438ux435-ux430ux43dux430ux43bux438ux442ux438ux447ux435ux441ux43aux438ux445-ux440ux435ux448ux435ux43dux438ux439-ux441-ux447ux438ux441ux43bux435ux43dux43dux44bux43cux438}

ЦЕЛЬ ЛЕКЦИИ: Сравнить достоинства и недостатки аналитических и численных
методов решения уравнений с частными производными. Приведена важная
математическая задача об идентификации физических величин
(параметрическая идентификация). В связи с этой задачей рассматривается
один важный пример из биологии.

Вероятно, пришло время обсудить достоинства и недостатки аналитических и
численных решений уравнений с частными производными. Сначала выясним,
что мы подразумеваем под этими двумя типами решений.

\section{Аналитические
решения}\label{ux430ux43dux430ux43bux438ux442ux438ux447ux435ux441ux43aux438ux435-ux440ux435ux448ux435ux43dux438ux44f}

Под аналитическими решениями мы подразумеваем такие решения, в которых
неизвестная функция и выражена через независимые переменные и параметры
системы в виде формул, бесконечных рядов и интегралов.

\section{Численные
решения}\label{ux447ux438ux441ux43bux435ux43dux43dux44bux435-ux440ux435ux448ux435ux43dux438ux44f}

Под численными мы будем понимать решения, полученные численно после
приближенной замены исходного уравнения другим, более простым
уравнением. Например, в конечно-разностном методе частные производные
выражаются через конечные разности и решение уравнения с частными
производными аппроксимируется решением разностного уравнения.
Результатом такой процедуры обычно является табличка значений решения и
при некоторых значениях независимых переменных.

Таблина 40.1

{\def\LTcaptype{none} % do not increment counter
\begin{longtable}[]{@{}
  >{\raggedright\arraybackslash}p{(\linewidth - 22\tabcolsep) * \real{0.1788}}
  >{\raggedright\arraybackslash}p{(\linewidth - 22\tabcolsep) * \real{0.0861}}
  >{\raggedright\arraybackslash}p{(\linewidth - 22\tabcolsep) * \real{0.1192}}
  >{\raggedright\arraybackslash}p{(\linewidth - 22\tabcolsep) * \real{0.0728}}
  >{\raggedright\arraybackslash}p{(\linewidth - 22\tabcolsep) * \real{0.0662}}
  >{\raggedright\arraybackslash}p{(\linewidth - 22\tabcolsep) * \real{0.0662}}
  >{\raggedright\arraybackslash}p{(\linewidth - 22\tabcolsep) * \real{0.0331}}
  >{\raggedright\arraybackslash}p{(\linewidth - 22\tabcolsep) * \real{0.0331}}
  >{\raggedright\arraybackslash}p{(\linewidth - 22\tabcolsep) * \real{0.0331}}
  >{\raggedright\arraybackslash}p{(\linewidth - 22\tabcolsep) * \real{0.0728}}
  >{\raggedright\arraybackslash}p{(\linewidth - 22\tabcolsep) * \real{0.1192}}
  >{\raggedright\arraybackslash}p{(\linewidth - 22\tabcolsep) * \real{0.1192}}@{}}
\toprule\noalign{}
\begin{minipage}[b]{\linewidth}\raggedright
\end{minipage} & \begin{minipage}[b]{\linewidth}\raggedright
\end{minipage} & \begin{minipage}[b]{\linewidth}\raggedright
\end{minipage} & \begin{minipage}[b]{\linewidth}\raggedright
\end{minipage} & \begin{minipage}[b]{\linewidth}\raggedright
\end{minipage} & \begin{minipage}[b]{\linewidth}\raggedright
\end{minipage} & \begin{minipage}[b]{\linewidth}\raggedright
х
\end{minipage} & \begin{minipage}[b]{\linewidth}\raggedright
\end{minipage} & \begin{minipage}[b]{\linewidth}\raggedright
\end{minipage} & \begin{minipage}[b]{\linewidth}\raggedright
\end{minipage} & \begin{minipage}[b]{\linewidth}\raggedright
\end{minipage} & \begin{minipage}[b]{\linewidth}\raggedright
\end{minipage} \\
\midrule\noalign{}
\endhead
\bottomrule\noalign{}
\endlastfoot
& 0 & 0,1 & 0,2 & 0,3 & 0,4 & 0,5 & 0,6 & 0,7 & 0,8 & 0,9 & 1 \\
00,010,020,03 & 100 & 10,20,15 & 10,34 & l: & : & : & ! & : & 10,34 &
10,20,15 & 1000 \\
; & & : & : & : & : & : & : & : & : & : & : \\
\end{longtable}
}

Теперь, когда мы знаем, что подразумевается под решением каждого из этих
типов, давайте сравним их между собой.

\section{\texorpdfstring{\textbf{Сравнительный анализ численного} и
аналитического
решений}{Сравнительный анализ численного и аналитического решений}}\label{ux441ux440ux430ux432ux43dux438ux442ux435ux43bux44cux43dux44bux439-ux430ux43dux430ux43bux438ux437-ux447ux438ux441ux43bux435ux43dux43dux43eux433ux43e-ux438-ux430ux43dux430ux43bux438ux442ux438ux447ux435ux441ux43aux43eux433ux43e-ux440ux435ux448ux435ux43dux438ux439}

Рассмотрим простейшую смещанную параболическую задачу

(40.1)

MATHBLK1868X

MATHBLK1869X

решение которой изображено на рис. 40.1. Что же лучше, иметь
аналитическое решение этой задачи

(40.2)

MATHBLK1870X

или численное, представленное при \(\alpha = 1\) в табл. 40.1?

Вопрос хороший, а ответ на него зависит от того, что мы хотим делать с
этим решением дальше. Однако каждому из типов решений присущи очевидные
преимущества.

\section{Преимущества аналитического
решения}\label{ux43fux440ux435ux438ux43cux443ux449ux435ux441ux442ux432ux430-ux430ux43dux430ux43bux438ux442ux438ux447ux435ux441ux43aux43eux433ux43e-ux440ux435ux448ux435ux43dux438ux44f}

\begin{enumerate}
\def\labelenumi{\arabic{enumi}.}
\tightlist
\item
  Очевидно, решение (40.2) более информативно, чем таблица чисел. Если
  мы хотим вычислить решение в любой конкретной точке (x, t), мы можем
  сделать это как угодно точно, просто увеличивая число учитываемых
  членов ряда. При этом без труда получается оценка сверху на величину
  допускаемой ошибки.
\end{enumerate}

\pandocbounded{\includegraphics[keepaspectratio]{images/_page_301_Figure_4.jpeg}}

РИС. 40.1. Решение уравнения теплопроводности (40.1) в различные моменты
времени.

\begin{itemize}
\tightlist
\item
  \begin{enumerate}
  \def\labelenumi{\arabic{enumi}.}
  \setcounter{enumi}{1}
  \tightlist
  \item
    Аналитическое решение всегда позволяет вычислить значения решения в
    одной точке (x, t), не прибегая к вычислению значений решений в
    других точках, как это бывает при решении задач явными или неявными
    разностными схемами.
  \end{enumerate}
\item
  \begin{enumerate}
  \def\labelenumi{\arabic{enumi}.}
  \setcounter{enumi}{2}
  \tightlist
  \item
    Аналитическое решение позволяет определить решение в любой точке, а
    не только в узлах сетки.
  \end{enumerate}
\item
  \begin{enumerate}
  \def\labelenumi{\arabic{enumi}.}
  \setcounter{enumi}{3}
  \tightlist
  \item
    Наиболее важным для нас преимуществом аналитического решения
    является возможность проследить влияние физических параметров,
    начальных и граничных условий на характер решения.
  \end{enumerate}
\end{itemize}

Численные методы не выявляют этих закономерностей, поскольку они
позволяют находить решение только при заданных параметрах, начальных и
граничных условиях. Иногда важно знать соотношения между параметрами
модели и решением, особенно если речь идет об оценке физических
параметров по виду решения. Например, предположим, что решение и
определено экспериментально и нам известно аналитическое решение

u = функция от параметров.

Тогда можно поставить вопрос об определении параметров нак функции
входных данных по схеме

Параметры = функция от u = функция исходных данных.

Такая задача называется параметрической идентификацией. Для
параметрической идентификации необходимо уметь решать уравнение с
частными производными. Несколько позже в этой лекции мы приведем пример
параметрической идентификации в биологии, но сначала посмотрим, чем же
хороши численные решения.

\section{Преимущества численных
решений}\label{ux43fux440ux435ux438ux43cux443ux449ux435ux441ux442ux432ux430-ux447ux438ux441ux43bux435ux43dux43dux44bux445-ux440ux435ux448ux435ux43dux438ux439}

Главное преимущество численных решений состоит в том, что их можно
получить даже в том случае, когда аналитическое решение получить
невозможно. Практически все нелинейные уравнения с частными производными
необходимо решать численными методами, а большинство реально
существующих моделей физики, химии, биологии и т. д. являются по своей
природе нелинейными. Как правило, линейные модели лишь приближенно
аппроксимируют нелинейные, если в них отбросить нелинейные члены. В
качестве примера приведем некоторые нелинейные уравнения:

\begin{enumerate}
\def\labelenumi{\arabic{enumi})}
\item
  нелинейное волновое уравнение \(u_{tt} = u_{xx} + f(u)\) ;
\item
  уравнение реакции с нелинейной диффузией \(u_t = u_{xx} + f(u)\) ;
\item
  система Ходжкина---Хаксли
\end{enumerate}

{[} \textbackslash begin\{cases\} u\_t = u\_\{xx\} + f(u, v),
\textbackslash{} v\_t = g(u, v). \textbackslash end\{cases\} {]}

Аналитические решения для этих уравнений не известны ни при каких
нелинейных функциях f и \(g^{(1)}\) . Поэтому общий подход к решению
нелинейных (а в ряде случаев и линейных) задач обычно базируется на
численных решениях.

Теперь мы рассмотрим пример, как с помощью аналитического решения можно
определять физические параметры. Более подробное изложение этих задач
можно найти в рекомендованной литературе.

\section{Параметрическая идентификация (в
биологии)}\label{ux43fux430ux440ux430ux43cux435ux442ux440ux438ux447ux435ux441ux43aux430ux44f-ux438ux434ux435ux43dux442ux438ux444ux438ux43aux430ux446ux438ux44f-ux432-ux431ux438ux43eux43bux43eux433ux438ux438}

Предположим, что биолог пытается определить, с какой скоростью ионы
поташа \((K^+)\) диффундируют в растворе экзоплазмы. Если знать этот
коэффициент диффузии, можно много сказать о том, как нервные импульсы
передаются вдоль аксонов. Проблема заключается в том, что этот
коэффициент практически невозможно найти прямым измерением. Можно,
однако, найти математическую

1) Это не так. В качестве примера можно привести уравнение
\(u_{ii} = u_{xx} - \sin u_x - \Pi pum\) , ped.

связь между концентрацией поташа \(u\left(x,\,t\right)\) и коэффициентом
диффузии \(D,\,\,a\,\,\) затем по измерениям величины
\(u\left(x,\,t\right)\) определить величину D.

Биологи Ходжкин и Кейс обнаружили, что после помещения гигантского
аксона в специальный солевой раствор в начальный момент концентрация
радиочктивного поташа (*2K) вдоль аксона приближенно описывается кривой

MATHBLK1871X.

Значит, параметры A н a следует подобрать так, чтобы кривая
соответствовала экспериментальным данным (например, в смысле метода
наименьших квадратов). Было также обнаружено, что Р И С

\pandocbounded{\includegraphics[keepaspectratio]{images/_page_303_Figure_6.jpeg}}

. 40.2. Начальная концентрация ионов ( \(^{42}K\) ): a---теоретическая
кривая; \(\delta\) ---экспериментальные точки.

с течением времени ноны поташа под действием диффузии и конвекции
растекаются вдоль аксона. Ходжкин и Кейс предположили, что концентрация
и является решением следующей задачи с начальными данными для уравнения
диффузии:

(40.3)

MATHBLK1872X

Эта задача была решена в лекции 15 путем перехода в движущуюся систему
координат. Если читатель помнит, то общая схема метода выглядит
следующим образом. Сначала отбрасываем конвективный член и получаем
решение чисто диффузионной задачи (при V = 0):

MATHBLK1873X

MATHBLK1874X

а затем в этом решении заменяем x на x-Vt и получаем решение задачи
(40.3):

(40.4)

MATHBLK1875X

В опыте Ходжкина и Кейса проводилось прямое измерение скорости, а
формула (40.4) устанавливала связь между концентрацией ионов (
\(^{42}K\) ) и коэффициентом диффузии D (величины A, a и V были
известны). Экспериментальная установка позволяла измерять Р И С

\pandocbounded{\includegraphics[keepaspectratio]{images/_page_304_Figure_3.jpeg}}

. 40.3. Кривая \(u\left(x,\ t\right)\) соответствует экспериментальным
точкам: a---кривая, наилучшим образом соответствующая экспериментальным
данным; b---экспериментально найденная концентрация в точке.

концентрацию поташа \(u(x, t)\) в одной точке \(x_0\) аксона, но в
различные моменты времени. Оставалось найти коэффициент D таким образом,
чтобы теоретическая кривая (40.4) наилучшим образом соответствовала
экспериментальным данным (рис. 40.3).

Другими словами, для различных коэффициентов D мы получим различные
кривые \(u\left(x_{0},\ t\right)\) , и поэтому величину D нужно выбрать
так, чтобы решение уравнения соответствовало опыту. В своем эксперименте
Ходжкин и Кейс установили, что коэффициент диффузии ионов \((K^{+})\) в
экзоплазме равен \(D=1.5\times10^{-5}\) см²/с.

\section{ЗАДАЧИ}\label{ux437ux430ux434ux430ux447ux438-18}

\begin{itemize}
\tightlist
\item
  \begin{enumerate}
  \def\labelenumi{\arabic{enumi}.}
  \tightlist
  \item
    Какой эксперимент могли бы вы предложить для оценки параметра α в
    задаче (40.1), если в вашем распоряжении есть аналитическое решение
    вида (40.2)?
  \end{enumerate}
\item
  \begin{enumerate}
  \def\labelenumi{\arabic{enumi}.}
  \setcounter{enumi}{1}
  \tightlist
  \item
    В методе наименьших квадратов минимизируется сумма квадратов
    уклонений кривой от экспериментальных точек. Например, согласно
    методу наименьших квадратов, прямая \(y(x)\) = a + bx минимизирует
    (см. рис. 40.4) функцию
  \end{enumerate}
\end{itemize}

MATHBLK1876X

Выразите константы a и b через начальные данные из условия минимума
функции SS.

\begin{enumerate}
\def\labelenumi{\arabic{enumi}.}
\setcounter{enumi}{2}
\tightlist
\item
  Одной из важных задач в биохимии является определение молекулярного
  веса макромолекул многлобина. В одном из способов кровь помещают в
  ультрацентрифугу и вращают в течение заданного промежутка времени.
  Концентрация жидкости в центрифуге описывается уравнением Ламма
\end{enumerate}

MATHBLK1877X

где

r-расстояние от центра центрифуги,

D---коэффициент диффузии (зависит от молек лярного веса миоглобина),

s -- коэффициент седиментации (определяется экспериментально),

\pandocbounded{\includegraphics[keepaspectratio]{images/_page_305_Figure_8.jpeg}}

РИС. 40.4. Аппрокенмация экспериментальных данных по методу наименьших
квадратов.

\(\omega\) --- угловая скорость вращения центрифуги (известна),
\(u(r, t)\) --- концентрация вещества в центрифуге. Молекулярный вес
миоглобина определяется с помощью стационарного ( \(u_t \equiv 0\) )
решения \(u(r, \infty)\) уравнения Ламма. Для определения
\(u(r, \infty)\) получается обыкновенное дифференциальное уравнение.
Вопрос ставится следующим образом: как по экспериментально найденному
стационарному состоянию оценить молекулярный вес многлобина?

\section{\texorpdfstring{КЛАССИФИКАЦИЯ УРАВНЕНИЙ (ПАРАБОЛИЧЕСКИЕ И
ЭЛЛИПТИЧЕСКИЕ
\textbf{УРАВНЕНИЯ})}{КЛАССИФИКАЦИЯ УРАВНЕНИЙ (ПАРАБОЛИЧЕСКИЕ И ЭЛЛИПТИЧЕСКИЕ УРАВНЕНИЯ)}}\label{ux43aux43bux430ux441ux441ux438ux444ux438ux43aux430ux446ux438ux44f-ux443ux440ux430ux432ux43dux435ux43dux438ux439-ux43fux430ux440ux430ux431ux43eux43bux438ux447ux435ux441ux43aux438ux435-ux438-ux44dux43bux43bux438ux43fux442ux438ux447ux435ux441ux43aux438ux435-ux443ux440ux430ux432ux43dux435ux43dux438ux44f}

ЦЕЛЬ ЛЕКЦИИ: Для линейного дифференциального уравнения второго порядка

MATHBLK1878X

найти координатную систему ( , η), у якій рівняння зводиться до однієї з
двох канонічних форм

{[} u\_\{\textbackslash eta\textbackslash eta\} = \textbackslash varphi
(\textbackslash xi, \textbackslash eta, u, u\_\{\textbackslash xi\},
u\_\{\textbackslash eta\}) {]}

(параболічний канонічний

форма),
\(u_{\xi\xi} + u_{\eta\eta} = q \; (\xi, \, \eta, \, u, \, u_{\xi}, \, u_{\eta})\)
(еліптична канонічна форма)

залежно від значення \(B^{2} - 4AC = 0\) є параболічним випадком,
\(B^2 - 4AC < 0\) є еліптичним випадком.

У лекції 23 ми розділили лінійні рівняння другого порядку двома змінними
вигляду

(41.1)

{[} Au\_\{xx\} + Bu\_\{xy\} + Cu\_\{yx\} + Du\_x + Eu\_y + Fu = G {]}

на три основні класи та знайшли дві канонічні форми рівнянь
гіперболічного типу.

У цій лекції ми покажемо, як канонічно зводити параболічні та еліптичні
рівняння.

\section{Зведення параболічних рівнянь до канонічної
форми}\label{ux437ux432ux435ux434ux435ux43dux43dux44f-ux43fux430ux440ux430ux431ux43eux43bux456ux447ux43dux438ux445-ux440ux456ux432ux43dux44fux43dux44c-ux434ux43e-ux43aux430ux43dux43eux43dux456ux447ux43dux43eux457-ux444ux43eux440ux43cux438}

Розглянемо рівняння виду (41.1) з \(B^2 - 4AC = 0\) і введемо нові
координати (E, p) таким чином, щоб звести рівняння до вигляду

{[} u\_\{\textbackslash eta\textbackslash eta\} =
\textbackslash varphi(\textbackslash xi, \textbackslash eta, u,
u\_\{\textbackslash xi\}, u\_\{\textbackslash eta\}). {]}

З формальною заміною змінних \(\xi = \xi(x, y)\) та
\(\eta = \eta(x, y)\) (41.1) зводиться до тієї ж форми

(41.2)

{[}
\textbackslash overline\{A\}u\_\{\textbackslash xi\textbackslash xi\} +
\textbackslash overline\{B\}u\_\{\textbackslash xi\textbackslash eta\} +
\textbackslash overline\{C\}u\_\{\textbackslash eta\textbackslash eta\}
+ \textbackslash overline\{D\}u\_\{\textbackslash xi\} +
\textbackslash overline\{E\}u\_\{\textbackslash eta\} +
\textbackslash overline\{F\}u = \textbackslash overline\{G\} {]}

, рей

{[} \textbackslash overline\{A\} =
A\textbackslash xi\_\{\textbackslash lambda\}\^{}\{2\} +
B\textbackslash xi\_\{\textbackslash lambda\}\textbackslash xi\_\{y\} +
C\textbackslash xi\_\{y\}\^{}\{2\}, {]}

{[} \textbackslash overline\{B\} =
2A\textbackslash xi\_\{x\}\textbackslash eta\_\{x\} +
B(\textbackslash xi\_\{x\}\textbackslash eta\_\{y\} +
\textbackslash xi\_\{y\}\textbackslash eta\_\{x\}) +
2C\textbackslash xi\_\{y\}\textbackslash eta\_\{y\}, {]}

{[} \textbackslash overline\{C\} =
A\textbackslash eta\_\{\textbackslash lambda\}\^{}\{2\} +
B\textbackslash eta\_\{x\}\textbackslash eta\_\{y\} +
C\textbackslash eta\_\{y\}\^{}\{2\}, {]}

{[} \textbackslash overline\{D\} = A\textbackslash xi\_\{xx\} +
B\textbackslash xi\_\{xy\} + C\textbackslash xi\_\{yy\} +
D\textbackslash xi\_\{x\} + E\textbackslash xi\_\{y\}, {]}

{[} \textbackslash overline\{E\} =
A\textbackslash eta\_\{\textbackslash lambda\} +
B\textbackslash eta\_\{\textbackslash lambda y\} +
C\textbackslash eta\_\{yy\} + D\textbackslash eta\_\{x\} +
E\textbackslash eta\_\{y\}, {]}

{[} \textbackslash overline\{F\} = F, \textbackslash quad
\textbackslash overline\{G\} = G. {]}

Наша мета --- змусити коефіцієнти \textbackslash overline\{A\} і
\textbackslash overline\{B\} (або \textbackslash overline\{C\} і
\textbackslash overline\{B\} ) стати нулем (ви швидко можете побачити,
що B\^{}2-4AC має дорівнювати нулю для реалізації нашого плану). Для
впевненості припустимо, що \textbackslash overline\{A\}=0 і
\textbackslash overline\{B\}=0 і розв'язуємо отримані рівняння. Спершу
розв'яжемо рівняння \textbackslash overline\{A\}=0 за значенням
{[}\textbackslash xi\_\{\textbackslash lambda\}/\textbackslash xi\_\{\textbackslash nu\}{]}
:

{[} {[}\textbackslash xi\_x/\textbackslash xi\_y{]} = -B/2A {]}

.

\begin{verbatim}
Задовольняємо це рівняння, якщо виберемо нову координату \xi = \xi \, (x, \, y) як загальний інтеграл
\end{verbatim}

{[} \textbackslash xi(x, y) = c {]}

Диференціальне рівняння

{[} \textbackslash frac\{dy\}\{dx\} =
-\textbackslash left{[}\textbackslash xi\_x/\textbackslash xi\_y\textbackslash right{]}
= B\_t 2A. {]}

Наприклад, якщо dy/dx = B/2A = 3, то y - 3x = c, а отже
\textbackslash xi(x, y) = y - 3x . Ця функція задовольняє умову
{[}\textbackslash xi\_x/\textbackslash xi\_y{]} = -3 і, відповідно,
повертає коефіцієнт \textbackslash overline\{A\} до нуля.

\begin{verbatim}
Отже, половина роботи виконана: ми знайшли нову координатну \xi = \xi(x, y) , яка забезпечила рівність \overline{A} = 0 . Залишається знайти \eta(x, y) таким чином, щоб перетворити коефіцієнт \overline{B} на нуль.

Час зупинитися тут. Виявляється, що якщо вибрати B^2-4AC=0 і \xi так, щоб \overline{A}=0 , то коефіцієнт \overline{B} автоматично стане нулем 1). Тепер ми побачимо це.

Коефіцієнт \overline{B} задається виразом
\end{verbatim}

{[} \textbackslash overline\{B\} =
2A\textbackslash xi\_x\textbackslash eta\_x +
B(\textbackslash xi\_x\textbackslash eta\_y +
\textbackslash xi\_y\textbackslash eta\_x) +
2C\textbackslash xi\_y\textbackslash eta\_y, {]}

{[} \textbackslash overline\{B\}\^{}2 -
4\textbackslash overline\{A\}\textbackslash overline\{C\} = (B\^{}2 -
4AC) (\textbackslash xi\_x \textbackslash eta\_y - \textbackslash xi\_y
\textbackslash eta\_x)\^{}2 - \textbackslash Pi \textbackslash rho u u.
~\textbackslash rho e \textbackslash partial. {]}

1) Завдяки рівності

І з \(B^2 - 4AC = 0\) році , то

{[} \textbackslash overline\{B\} = 2 \textbackslash left( V
\textbackslash overline\{A\} \textbackslash xi\_x + V
\textbackslash overline\{C\} \textbackslash xi\_y \textbackslash right)
\textbackslash left( V \textbackslash overline\{A\}
\textbackslash eta\_x + V \textbackslash overline\{C\}
\textbackslash eta\_y \textbackslash right). {]}

Оскільки
\(\xi_x/\xi_y = -B \ 2A = -2 \sqrt{A\tilde{C}}/2A = -\sqrt{C/A}\)
останній вираз для B буде записаний як

{[} \textbackslash overline\{B\} = 2
\textbackslash sqrt\{\textbackslash overline\{A\}\}
\textbackslash left{[}
\textbackslash sqrt\{\textbackslash overline\{A\}\} ,
\textbackslash eta\_x +
\textbackslash sqrt\{\textbackslash overline\{C\}\} ,
\textbackslash eta\_y \textbackslash right{]}. {]}

І оскільки \(\overline{A}\) дорівнює нулю, то \(\overline{B}\) , звісно,
також дорівнює нулю.

Оскільки при нашому виборі \(\xi\) обидва коефіцієнти \(\overline{A}\) і
\(\overline{B}\) повертаються до нуля, ми можемо взяти будь-яку функцію
як \(\eta\) (якщо вона не пропорційна координаті \(\xi\) ). У прикладі
можна вибрати \(\eta = y\) .

Залишається знайти канонічну форму рівняння в нових координатах. Для
цього потрібно просто підставити формули \(\xi\) і \(\eta\) (41.3) і
знайти всі шанси \(\overline{A}\) , \(\overline{B}\) , \(\overline{C}\)
, \(\overline{D}\) , \(\overline{E}\) , \(\overline{F}\) і
\(\overline{G}\) . На цьому вивчення параболічного випадку завершується,
але перш ніж перейти до іншої теми, розглянемо простий приклад.

\section{\texorpdfstring{Перетворення параболічного рівняння
\(u_{xx} + 2u_{xy} + u_{yy} = 0\) у канонічну
форму}{Перетворення параболічного рівняння u\_\{xx\} + 2u\_\{xy\} + u\_\{yy\} = 0 у канонічну форму}}\label{ux43fux435ux440ux435ux442ux432ux43eux440ux435ux43dux43dux44f-ux43fux430ux440ux430ux431ux43eux43bux456ux447ux43dux43eux433ux43e-ux440ux456ux432ux43dux44fux43dux43dux44f-u_xx-2u_xy-u_yy-0-ux443-ux43aux430ux43dux43eux43dux456ux447ux43dux443-ux444ux43eux440ux43cux443}

Розглянемо рівняння

{[} u\_\{xy\} + 2u\_\{xy\} + u\_\{yy\} = 0. {]}

Його коефіцієнти мають такі значення: A=1, B=2, C=1, D=E=F=G=0. Отже,
\(B^2-4AC=0\) для всіх значень x і y. Щоб знайти нові координати \(\xi\)
і \(\eta\) та канонічну форму рівняння, ми рухаємося наступним чином.

КРОК 1. Запишемо \emph{характеристичне рівняння} (тепер лише одне, а не
два)

{[} \textbackslash frac\{dy\}\{dx\} =
-{[}\textbackslash xi\_x/\textbackslash xi\_y{]} = B/2A = 1. {]}

Розв'язуючи це рівняння відносно y, отримуємо

{[} y = x + c {]}

,

Де --- характеристика \(\xi = y - x\) визначена. Використання цієї
характеристики забезпечує досягнення рівності \(\overline{A} = 0\) .
Координату \(\eta\) можна вибрати довільно (якщо вона не залежить від
\(\xi\) ). Ми оберемо його наступним чином:

{[} \textbackslash eta = y {]}

.

Нові координати

\begin{verbatim}
\begin{cases}
\end{verbatim}

{[} \textbackslash xi = y - x, \textbackslash{} {]}

{[} \textbackslash eta = y {]}

\begin{verbatim}
\end{cases}
\end{verbatim}

изображены на рис. 41 1.

\pandocbounded{\includegraphics[keepaspectratio]{images/_page_309_Figure_3.jpeg}}

Рис. 41.1. Новая система координат \(\xi = y - x\) , \(\eta = y\) .

ШАГ 2. Осталось найти каноническую форму уравнения. Подставляя \(\xi\) и
\(\eta\) в формулы для коефіцієнтів \(\overline{A}\) , \(\overline{B}\)
, \(\overline{C}\) , \(\overline{D}\) , \(\overline{E}\) ,
\(\overline{F}\) и \(\overline{G}\) , получаем \(\overline{A} = 0\) (так
и должно быть; мы находили \(\xi\) из условия равенства нулю цього
коэффициента), \(\overline{B}=0\) (мы уже раньше показали, що этог
коефіцієнт також обращается в нуль),

(41.5)

{[} \textbackslash overline\{C\} = A \textbackslash eta\_x\^{}2 + B
\textbackslash eta\_x \textbackslash eta\_y + C
\textbackslash eta\_y\^{}2 = 1, {]}

{[} \textbackslash overline\{D\} = A \textbackslash xi\_\{xx\} + B
\textbackslash xi\_\{xy\} + C \textbackslash eta\_\{yy\} + D
\textbackslash xi\_x + E \textbackslash xi\_y = 0, {]}

{[} \textbackslash overline\{E\} = A \textbackslash eta\_\{xx\} + B
\textbackslash eta\_\{xy\} + C \textbackslash eta\_\{yy\} + D
\textbackslash eta\_x + E \textbackslash eta\_y = 0, {]}

{[} \textbackslash overline\{F\} = F = 0, {]}

{[} \textbackslash overline\{G\} = G = 0. {]}

Отже, нове рівняння

{[}
\textbackslash overline\{A\}u\_\{\textbackslash xi\textbackslash xi\} +
\textbackslash overline\{B\}u\_\{\textbackslash xi\textbackslash eta\} +
\textbackslash overline\{C\}u\_\{\textbackslash eta\textbackslash eta\}
+ \textbackslash overline\{D\}u\_\{\textbackslash xi\} +
\textbackslash overline\{E\}u\_\{\textbackslash eta\} +
\textbackslash overline\{F\}u = \textbackslash overline\{G\} {]}

перетворюється на \(u_{nn} = 0\) .

На цьому наш приклад завершено, але перш ніж рухатися далі, звернемо
увагу на одну обставину у позначення канонічної форми. Це рівняння
настільки просте, що можна знайти загальне розв'язання (тобто всі
розв'язки). Інтегруючи \(u_{\eta\eta} = 0\) до \(\eta\) , ми отримуємо

{[} u\_\{\textbackslash eta\} = f(\textbackslash xi), {]}

де \(f(\xi)\) --- довільна функція. Реінтеграція веде до спільного
рішення

{[} u(\textbackslash xi, \textbackslash eta) = \textbackslash eta
f(\textbackslash xi) + g(\textbackslash xi), {]}

де \(g(\xi)\) --- інша довільна функція.

Читач легко може переконати себе, що будь-яка функція такого роду
задовольняє рівняння \(u_{nn} = 0\) .

Залишається зробити останній крок і замінити початкові координати x і y
у загальний розв'язок.

Внаслідок підстановки отримуємо, що

{[} u(x, y) = yf(y-x) + g(y-x) {]}

є загальним розв'язком рівняння

{[} u\_\{xx\} + 2u\_\{xy\} + u\_\{yy\} = 0. {]}

Наприклад, якщо вибрати \(f(x) = \sin x\) і \(g(x) = x^2\) , то

{[} u(x, y) = y \textbackslash sin(y-x) + (y-x)\^{}2 {]}

буде одним із нескінченної кількості розв'язків (звісно, розв'язок, що
моделює певну фізичну задачу, залежить від початкових і граничних умов).

\section{Зведення еліптичних рівнянь до канонічної
форми}\label{ux437ux432ux435ux434ux435ux43dux43dux44f-ux435ux43bux456ux43fux442ux438ux447ux43dux438ux445-ux440ux456ux432ux43dux44fux43dux44c-ux434ux43e-ux43aux430ux43dux43eux43dux456ux447ux43dux43eux457-ux444ux43eux440ux43cux438}

Розглянемо загальне рівняння

{[} Au\_\{xx\} + Bu\_\{xy\} + Cu\_\{yy\} + Du\_x + Eu\_y + Fu = G {]}

Але тепер у \(B^2-4AC<0\) . Переходячи до нових пояснювальних змінних,
ми хочемо перетворити її у нову форму

{[} u\_\{\textbackslash xi\textbackslash xi\} +
u\_\{\textbackslash eta\textbackslash eta\} =
\textbackslash varphi(\textbackslash xi, \textbackslash eta, u,
u\_\{\textbackslash xi\}, u\_\{\textbackslash eta\}). {]}

Щоб знайти ці нові змінні, ми зробимо ті ж обчислення, що й у двох
попередніх випадках: отримаємо трансформоване рівняння

{[}
\textbackslash overline\{A\}u\_\{\textbackslash xi\textbackslash xi\} +
\textbackslash overline\{B\}u\_\{\textbackslash xi\textbackslash eta\} +
\textbackslash overline\{C\}u\_\{\textbackslash eta\textbackslash eta\}
+ \textbackslash overline\{D\}u\_\{\textbackslash xi\} +
\textbackslash overline\{E\}u\_\{\textbackslash eta\} +
\textbackslash overline\{F\}u = \textbackslash overline\{G\} {]}

і вимагати, щоб \(\xi\) і \(\eta\) задовольняли рівняння
\(\overline{A}=\overline{C}\) і \(\overline{B}=0\) . На жаль, ці
рівняння не дозволяють одразу знайти \(\xi\) і \(\eta\) , як це було у
попередніх випадках, тому ми продовжимо трохи пізніше.

Ми будемо шукати необхідне перетворення як композицію двох перетворень.

\section{Конверсія
1}\label{ux43aux43eux43dux432ux435ux440ux441ux456ux44f-1}

За першим перетворенням у нові координати \textbf{ξ і ŋ ми} надаємо
нашому рівнянню таку форму

{[} u\_\{\textbackslash xi\textbackslash eta\} =
\textbackslash psi(\textbackslash xi, \textbackslash eta, u,
u\_\{\textbackslash xi\}, u\_\{\textbackslash eta\}). {]}

Це можливо, але лише шляхом введення складних координат. Щоб знайти ці
комплексні координати \(\xi\) і \(\eta\) , розв'яжемо характеристичні
рівняння, так само як у гіперболічному випадку:

{[} \textbackslash frac\{dy\}\{dx\} = \textbackslash frac\{B -
\textbackslash sqrt\{B\^{}2 - 4AC\}\}\{2A\} {]}

, (нагадуємо, що \(B^2 - 4AC < 0\) ),

і ми отримаємо

{[} \textbackslash xi(x, y) = \textbackslash text\{const\},
\textbackslash quad \textbackslash eta(x, y) =
\textbackslash text\{const\}. {]}

Наприклад, характеристичні рівняння

{[} \textbackslash frac\{dy\}\{dx\} = -\textbackslash sqrt\{-4x\^{}2\} =
-2ix, {]}

{[} \textbackslash frac\{dy\}\{dx\} = \textbackslash sqrt\{-4x\^{}2\} =
2ix {]}

також задаються комплексно-спряженим координатам

{[} \textbackslash xi = y + ix\^{}2, \textbackslash quad
\textbackslash eta = y - ix\^{}2. {]}

\section{Конверсія
2}\label{ux43aux43eux43dux432ux435ux440ux441ux456ux44f-2}

\begin{verbatim}
Наступне перетворення з (\xi,\eta) на (\alpha,\beta) виконується за формулами
\end{verbatim}

{[} \textbackslash alpha = \textbackslash frac\{\textbackslash xi +
\textbackslash eta\}\{2\} {]}

{[} , \textbackslash beta = \textbackslash frac\{\textbackslash xi -
\textbackslash eta\}\{2i\} . {]}

Внаслідок застосування другого перетворення, рівняння

{[} u\_\{\textbackslash xi\textbackslash eta\} =
\textbackslash psi(\textbackslash xi, \textbackslash eta, u,
u\_\{\textbackslash xi\}, u\_\{\textbackslash eta\}) {]}

набуває остаточної форми

{[} u\_\{\textbackslash alpha\textbackslash alpha\} +
u\_\{\textbackslash beta\textbackslash beta\} =
\textbackslash varphi(\textbackslash alpha, \textbackslash beta, u,
u\_\{\textbackslash alpha\}, u\_\{\textbackslash beta\}), {]}

де \textbackslash phi та \textbackslash phi використовуються для
позначення правих сторін рівнянь у загальному вигляді. Замість того, щоб
доводити, що ці дві трансформації дійсно ведуть до бажаного результату,
давайте розглянемо простий приклад.

\section{Переведення рівняння y\^{}2u\_\{xx\} + x\^{}2u\_\{yy\} = 0 до
канонічної
форми}\label{ux43fux435ux440ux435ux432ux435ux434ux435ux43dux43dux44f-ux440ux456ux432ux43dux44fux43dux43dux44f-y2u_xx-x2u_yy-0-ux434ux43e-ux43aux430ux43dux43eux43dux456ux447ux43dux43eux457-ux444ux43eux440ux43cux438}

Розглянемо рівняння

{[} y\^{}2 u\_\{xx\} + x\^{}2 u\_\{yy\} = 0, {]}

у яких A=y\^{}2 , B=0, C=x\^{}2 , D=E=F=G=0. Дискримінантний B\^{}2-4AC
дорівнює -4x\textsuperscript{2y}2 , і ми зведемо це рівняння до
канонічної форми в першому квадранті x\textgreater0, y\textgreater0.

КРОК 1 (Перше перетворення)

Запишемо характеристичні рівняння

{[} \textbackslash frac\{dy\}\{dx\} = \textbackslash frac\{B -
\textbackslash sqrt\{B\^{}2 - 4AC\}\}\{2A\} =
-\textbackslash frac\{\textbackslash sqrt\{-4x\textsuperscript{2y}2\}\}\{2y\^{}2\}
= -i\textbackslash frac\{x\}\{y\}, {]}

{[} \textbackslash frac\{dy\}\{dx\} = \textbackslash frac\{B +
\textbackslash sqrt\{B\^{}2 - 4AC\}\}\{2A\} =
\textbackslash frac\{\textbackslash sqrt\{-4x\textsuperscript{2y}2\}\}\{2y\^{}2\}
= i\textbackslash frac\{x\}\{y\} {]}

Крім того, S. S. Розв'язуючи їх, отримуємо перші інтеграли

{[} y\^{}2 + ix\^{}2 = \textbackslash text\{const\}, {]}

{[} y\^{}2 - ix\^{}2 = \textbackslash text\{const\}, {]}

Від

{[} \textbackslash xi(x, y) = y\^{}2 + ix\^{}2, {]}

\textbackslash eta(x, y) = y\^{}2 - ix\^{}2. (Неважливо, яка з двох
функцій називається \textbackslash xi , а яка --- \textbackslash eta ;
ви можете поміняти їх місцями, якщо хочете.) Після цього перетворення
початкове рівняння набуває вигляду

{[} u\_\{\textbackslash xi\textbackslash eta\} =
\textbackslash psi(\textbackslash xi, \textbackslash eta, u,
u\_\{\textbackslash xi\}, u\_\{\textbackslash eta\}). {]}

Ми не звертатимемо уваги на це рівняння (яке є комплексним гіперболічним
рівнянням), а одразу перейдемо до другого перетворення.

КРОК 2. (Друга трансформація)

При здійсненні другої трансформації ми маємо

{[} \textbackslash alpha = \textbackslash frac\{\textbackslash xi +
\textbackslash eta\}\{2\} = y\^{}2 {]}

(реальна частина \textbackslash xi і \textbackslash eta ), \$beta =
\textbackslash frac\{\textbackslash xi -
\textbackslash eta\}\{2\textbackslash iota\} = x\^{}2 (\$уявна частина
\textbackslash xi і \textbackslash eta ).

\begin{verbatim}
Щоб спростити нотацію, найкраще перепозначити (\alpha,\beta) (\xi,\eta) окремо і записати отримане \epsilon перетворення як
\end{verbatim}

{[} \textbackslash xi(x, y) = y\^{}2, \textbackslash quad
\textbackslash eta(x, y) = x\^{}2. {]}

КРОК 3 (Отримання нового рівняння)

\begin{verbatim}
Нову канонічну форму можна отримати, обчисливши всі коефіцієнти \overline{A}, \overline{B}, \overline{C}, \overline{D}, \overline{E}, \overline{F} і \overline{G} у рівнянні
\end{verbatim}

{[} \textbackslash overline\{A\}u\_\{28\} +
\textbackslash overline\{B\}u\_\{2n\} +
\textbackslash overline\{C\}u\_\{nn\} + \textbackslash overline\{D\}u +
\textbackslash overline\{E\}u\_\{n\} + \textbackslash overline\{F\}u =
\textbackslash overline\{G\} {]}

за формулами (41.5) з \textbackslash xi = y\^{}2 та \textbackslash eta =
x\^{}2 . Зробивши це, отримуємо канонічну форму еліптичного рівняння

{[} u\_\{s\textbackslash xi\} +
u\_\{\textbackslash eta\textbackslash eta\} =
\textbackslash frac\{-\textbackslash epsilon u\_\{\textbackslash eta\} -
\textbackslash eta
u\_\{\textbackslash xi\}\}\{2\textbackslash xi\textbackslash eta\}. {]}

\section{ЗАУВАЖЕННЯ}\label{ux437ux430ux443ux432ux430ux436ux435ux43dux43dux44f-5}

\begin{enumerate}
\def\labelenumi{\arabic{enumi}.}
\tightlist
\item
  Аналогічно можна класифікувати лінійні рівняння другого порядку з
  \emph{трьох або більше змінних}, але для цього потрібен матричний
  аналіз. Наприклад, рівняння другого порядку з трьома змінними
\end{enumerate}

{[} u\_t = u\_\{xx\} + u\_\{yy\}, {]}

згідно з цією класифікацією, буде параболічним, а рівняння

{[} u\_\{tt\} = u\_\{xx\} + u\_\{yy\} {]}

\begin{itemize}
\tightlist
\item
  Гіперболічно.
\end{itemize}

\begin{enumerate}
\def\labelenumi{\arabic{enumi}.}
\setcounter{enumi}{1}
\tightlist
\item
  Наш інтерес до класифікації диференціальних рівнянь у похідних
  пов'язаний із тим, що три основні типи рівнянь описують три різні типи
  фізичних явищ, і ми хотіли б отримати математичну класифікацію цих
  трьох класів фізичних явищ.
\end{enumerate}

\section{ЗАВДАННЯ}\label{ux437ux430ux432ux434ux430ux43dux43dux44f-9}

\begin{itemize}
\item
  \begin{enumerate}
  \def\labelenumi{\arabic{enumi}.}
  \tightlist
  \item
    Які з цих параболічних і еліптичних рівнянь записуються у канонічній
    формі:
  \end{enumerate}

  \begin{itemize}
  \tightlist
  \item
    \begin{enumerate}
    \def\labelenumi{\alph{enumi})}
    \tightlist
    \item
      u\_t = u\_\{xx\} hu ,

      \begin{itemize}
      \tightlist
      \item
        \begin{enumerate}
        \def\labelenumii{\arabic{enumii})}
        \setcounter{enumii}{5}
        \tightlist
        \item
          u\_\{xy\} + u\_\{xx\} + 3u = \textbackslash sin x ,
        \end{enumerate}
      \item
        \textbackslash mathbf\{B\}) ~\textbackslash hat\{u\_\{xx\}\} +
        2\textbackslash hat\{u\_\{yy\}\} = 0, -
        \textbackslash mathbf\{r\}) ~u\_\{xx\} = \textbackslash sin\^{}2
        x ?
      \end{itemize}
    \end{enumerate}
  \end{itemize}
\item
  \begin{enumerate}
  \def\labelenumi{\arabic{enumi}.}
  \setcounter{enumi}{1}
  \tightlist
  \item
    Перетворимо параболічне рівняння \(u_{xx}\) + 2u\_\{xy\} +
    \(u_{yy}\) + u = 2 у канонічну форму.
  \end{enumerate}
\item
  \begin{enumerate}
  \def\labelenumi{\arabic{enumi}.}
  \setcounter{enumi}{2}
  \tightlist
  \item
    Перетворимо еліптичне рівняння
    \(u_{xx}\)+2u\_\{yy\}+x\textsuperscript{2u\_x=\$e}\{-x\^{}2/2\}\$ у
    форму каноха.
  \end{enumerate}
\end{itemize}

\section{МЕТОД МОНТЕ-КАРЛО
(ВСТУП)}\label{ux43cux435ux442ux43eux434-ux43cux43eux43dux442ux435-ux43aux430ux440ux43bux43e-ux432ux441ux442ux443ux43f}

МЕТА ЛЕКЦІЇ: Пояснити основні ідени, що лежать в основі методу
Монте-Карло, і показати, як цей метод можна використовувати для
розв'язання різних задач. Виявляється, що за допомогою азартних ігор
можна змоделювати (загалом, на комп'ютері) наближені розв'язки реальних
задач. Як простий приклад можна вказати обчислення інтеграла

{[} I = \textbackslash int\_0\^{}1 x\^{}2 dx {]}

кидаючи дротик у одиничний квадрат \{(x, y): 0 \textless{} x \textless{}
1, 0 \textless{} y \textless{} 1\} Після приблизно 100 кидків відносна
частина дротиків, що опускається нижче кривої y = x\^{}2 , вважається
приблизним значенням інтеграла. Узагальнення такої гри (з випадковим
киданням дрогіка) завжди містить процедуру генерації послідовності
випадкових чисел.

У цій лекції ми коротко опишемо метод Монте-Карло, а в наступній лекції
покажемо, як його використовувати для розв'язання диференціальних
рівнянь у похідних

Методи Монте-Карло зазвичай називають групою методів розв'язання
детермінованих задач (тобто задач без випадковості), у яких суттєво
використовуються елементи випадковості. Загальна філософія етнічних
методів показана на рис. 42 1.

Ймовірнісна гра

\begin{verbatim}
Деторміноване завдання \rho — розв'язання задачі (наприклад, обчислення інтеграла, розв'язання UBP тощо).

Результат = \hat{p} Наближення p Рішення = p Рис. 42.1. Загальна теорія методу Монте-Карло.
\end{verbatim}

\section{Обчислення
інтегралу}\label{ux43eux431ux447ux438ux441ux43bux435ux43dux43dux44f-ux456ux43dux442ux435ux433ux440ux430ux43bux443}

Припустимо, ми хочемо обчислити інтеграл

{[} I = \textbackslash int\_\{a\}\^{}\{b\} f(x) , dx {]}

(детерміноване завдання). Щоб використати метод Монте-Карло, потрібно
винайти гру на удачу, яка дає приблизне значення цього інтеграла.

\begin{verbatim}
Звісно, можна придумати багато різних ігор. Остаточна версія залежить від точності розрахунків, простоти гри тощо. Очевидний варіант гри для обчислення інтегралу — кинути дротик у прямокутний R = \{(x,y): a \le x \le b, 0 \le y \le \max f(x)\} (див. рис. 42.2)
\end{verbatim}

\pandocbounded{\includegraphics[keepaspectratio]{images/_page_315_Figure_7.jpeg}}

Рис. 42 2. Обчислення інтегралу методом Монте-Карло.

Очевидно, якщо випадковим чином кинути близько 100 дротиків у
прямокутник R, що містить графік функції, можна отримати оцінку значення
інтеграла, помноживши площу прямокутника R на відносну кількість
дротиків під графіком функції. Отже, результат нашої гри

{[} \textbackslash widehat\{I\} = \textbackslash Big{[} {]}

\begin{verbatim}
\begin{aligned} \mbox{относительное число дротнков} \mbox{ дротнков} \mbox{ } \mbox{]} \times (Площадь области R) \mbox{ } \end{aligned}
\end{verbatim}

--- це інтегральний бал 1.

Щоб виконати ці обчислення на комп'ютері, нам потрібно якимось чином
отримати послідовність випадкових точок (подивимось, як це буде зроблено
пізніше), тобто наказати комп'ютеру «кидати дротики». Блок-схема на рис.
42.3 показує, як комп'ютер вирішить цю проблему.

\pandocbounded{\includegraphics[keepaspectratio]{images/_page_316_Figure_2.jpeg}}

\textbf{РИС.} 42.3. Блок-схема розрахунку інтегрального
\textbackslash int\_\{a\}\^{}\{b\} \(f(x)\) dx методом Монте-Карло (100
кидків); M = \textbackslash max\_\{a\} \(f(x)\) .

\section{Випадкові
числа}\label{ux432ux438ux43fux430ux434ux43aux43eux432ux456-ux447ux438ux441ux43bux430}

Перш ніж перейти до рівняння Монте-Карло з частковими похідними,
зупинимося на важливому питанні випадкових чисел. Щойно, обчислюючи
інтеграл, нам потрібно було побудувати послідовність випадкових точок
\(P_i = (x_i, y_i)\)всередині прямокутника R. Іншими словами,
x-координата точки P має бути випадковим числом з інтервалу {[}a, b{]},
а \(y_i\) від інтервалу {[}0, M{]}. Щоб знайти випадкові числа з цих
інтервалів, звертаємося до послідовності випадкових чисел, рівномірно
розподілених по сегменту {[}0, 1{]}. Тоді числа, рівномірно розподілені
на сегменті {[}a, b{]}, можна обчислити за формулою

{[} x\_i = a + (b - a) r\_i. {]}

Тепер, звичайно, виникає питання: як отримати послідовність випадкових
чисел, \{\(r_i\): i=1,2,\textbackslash ldots\} рівномірно розподілених
на {[}0,1{]}? Найпоширенішим методом сьогодні є метод дедукцій (метод
порівнянь). Використовуючи цей метод, отримуємо послідовність випадкових
цілих чисел (наприклад, 2120, 1401, 177, 3013,\ldots), потім ліворуч від
кожного числа розміщується десяткова крапка, так що всі числа
знаходяться між нулем і одним (0,212; 0,1401; 0,0177; 0,3013;
\textbackslash ldots) .

Отже, щоб отримати послідовність випадкових цілих чисел (між 0 і P), ми
використаємо метод віднімання.

\section{Отримання випадкових чисел за допомогою методу
відрахування.}\label{ux43eux442ux440ux438ux43cux430ux43dux43dux44f-ux432ux438ux43fux430ux434ux43aux43eux432ux438ux445-ux447ux438ux441ux435ux43b-ux437ux430-ux434ux43eux43fux43eux43cux43eux433ux43eux44e-ux43cux435ux442ux43eux434ux443-ux432ux456ux434ux440ux430ux445ux443ux432ux430ux43dux43dux44f.}

\begin{enumerate}
\def\labelenumi{\arabic{enumi}.}
\item
  Як перше випадкове ціле число, обирайте будь-яке число, що знаходиться
  між 0 і P (число P обирається наперед).
\item
  Помножте це випадкове число на певне фіксоване число.
\end{enumerate}

Попередньо обраний множник M.

\begin{enumerate}
\def\labelenumi{\arabic{enumi}.}
\setcounter{enumi}{2}
\item
  Додайте до добутку деяке фіксоване ціле число \emph{K}, обране
  заздалегідь.
\item
  Поділити отриману суму на P і залишок
\end{enumerate}

Виберіть нове випадкове число.

\begin{verbatim}
Метод виведення можна записати у формі формули r_{i+1} \equiv (Mr_i + K) \bmod P , що означає, що для отримання наступного r_{i+1} потрібно взяти попередній r_i , помножити його на M, додати результат на K, поділити суму на P і взяти решту поділу. Наприклад, якщо P = 100, M = 37, K = 16, r_0 = 15 , то отримується наступна послідовність випадкових чисел:
\end{verbatim}

{[} Mr\_0 + K = 571 {]}

{[} Mr\_1 + K = 2643 {]}

{[} Mr\_2 + K = 1607 {]}

{[} r = 71 {]}

{[} r\_2 = 43 {]}

Y

Оскільки P = 100, усі випадкові числа лежать між нулем і 100 (так званий
клас дедукції за модулем 100). Розміщуючи десяткову крапку перед кожним
із цих чисел, ми отримуємо послідовність випадкових чисел, що лежать між
нулем і одним \(r_0=0.15;\ r_1=0.71;\ r_2=0.43;\ r_3=0.07.\) (Ми також
позначали ці числа \(r_L\) )

\section{ЗАУВАЖЕННЯ}\label{ux437ux430ux443ux432ux430ux436ux435ux43dux43dux44f-6}

\begin{enumerate}
\def\labelenumi{\arabic{enumi}.}
\tightlist
\item
  Оскільки наш генератор випадкових чисел має P=100, дедукції можуть
  приймати лише значення \(0,1,2,\ldots,99\) Тому випадкові числа в
  нашому процесі дуже скоро почнуть повторюватися. Насправді, якщо
  продовжити наведений вище приклад, отримаємо
\end{enumerate}

15, 71, 43, 7, 43, 7, 43, 7\ldots{} (період із двох чисел).

Це означає, що наш метод не дуже хороший. Ідеально, якщо алгоритм
випадково видає весь клас відрахувань \(\{0, 1, 2, \ldots, 99\}\) до
початку повторних спроб. Математично можна довести, що якщо числа M, K і
P обираються за певними правилами, то незалежно від вибору першого
випадкового числа \(r_o\) алгоритм генеруватиме весь клас дедукцій. Тож
якщо ви оберете P дуже великим (наприклад, \(2^{40}\) ), можна бути
впевненим, що (з практичної точки зору) цей процес ніколи не
повториться.

\begin{enumerate}
\def\labelenumi{\arabic{enumi}.}
\setcounter{enumi}{1}
\tightlist
\item
  Звичайні чисельні методи (наприклад, метод Сімпсона), загалом вищі за
  якістю за метод Монте-Карло. Але якщо ми говоримо про обчислення
  багатовимірних інтегралів типу
\end{enumerate}

{[} I = \textbackslash int\_\{0\}\^{}\{1\}
\textbackslash int\_\{0\}\^{}\{1\} \textbackslash int\_\{0\}\^{}\{1\}
\textbackslash int\_\{0\}\^{}\{1\}
e\textsuperscript{\{-(x}\{2\}+y\textsuperscript{\{2\}+z}\{2\}+\textbackslash omega\^{}\{2\})\}
dx dy dz dw, {]}

тоді краще використовувати метод Монте-Карло.

\begin{enumerate}
\def\labelenumi{\arabic{enumi}.}
\setcounter{enumi}{2}
\tightlist
\item
  Можливо отримати випадкові числа з різними законами розподілу (і не
  лише рівномірними на сегменті {[}0, 1{]}). Комп'ютерні програми можуть
  генерувати випадкові числа з біноміальними, нормальними,
  гамма-розподілами та кратними
\end{enumerate}

дистрибуції. Ці розподіли зустрічаються в різних галузях природничих
наук і стають необхідними для моделювання реального світу.

\section{ЗАВДАННЯ}\label{ux437ux430ux432ux434ux430ux43dux43dux44f-10}

\begin{enumerate}
\def\labelenumi{\arabic{enumi}.}
\tightlist
\item
  Напишіть програму для оцінки інтегралу
\end{enumerate}

{[} I = \textbackslash int\_\{0\}\^{}\{1\} e\^{}\{\textbackslash sin x\}
dx. {]}

Яка точність вашої відповіді? Побудуйте графік точності та кількості
кидків. Чи збігається цей процес?

\begin{enumerate}
\def\labelenumi{\arabic{enumi}.}
\setcounter{enumi}{1}
\tightlist
\item
  Отримати послідовність випадкових чисел за допомогою алгоритму
\end{enumerate}

{[} r\_\{i+1\} \textbackslash equiv (3r\_i + 4) \textbackslash mod 7,
\textbackslash quad r\_0 = 0. {]}

\begin{enumerate}
\def\labelenumi{\arabic{enumi}.}
\setcounter{enumi}{2}
\tightlist
\item
  Напишіть програму (або принаймні блок-схему) для оцінки інтегралу
\end{enumerate}

{[} I = \textbackslash int\_\{0\}\^{}\{1\}
\textbackslash int\_\{0\}\^{}\{1\} \textbackslash int\_\{0\}\^{}\{1\}
e\textsuperscript{\{-(x}\{2\}+y\textsuperscript{\{2\}+z}\{2\})\} dx dy
dz. {]}

Що ви можете сказати про точність?

\begin{enumerate}
\def\labelenumi{\arabic{enumi}.}
\setcounter{enumi}{3}
\tightlist
\item
  Як отримати послідовність випадкових точок, що лежать всередині
  трикутника T?
\end{enumerate}

\pandocbounded{
\includegraphics[keepaspectratio]{images/_page_319_Figure_13.jpeg}}

\pandocbounded{\includegraphics[keepaspectratio]{images/_page_319_Figure_14.jpeg}}

\begin{enumerate}
\def\labelenumi{\arabic{enumi}.}
\setcounter{enumi}{4}
\tightlist
\item
  Як отримати послідовність випадкових чисел, розподіл яких показаний на
  рисунку?
\end{enumerate}

\pandocbounded{\includegraphics[keepaspectratio]{images/_page_319_Figure_16.jpeg}}

Іншими словами, як отримати послідовність чисел \{0, 1, 2\}, якщо
ймовірність виникнення 0 і 2 дорівнює 0,25, а ймовірність появи 1 ---
0,5?

Лекція 43

\section{РОЗВ'ЯЗАННЯ ДИФЕРЕНЦІАЛЬНИХ РІВНЯНЬ У ПОХІДНИХ ПОХІДНИХ ВІД
МОНТЕ-КАРЛО}\label{ux440ux43eux437ux432ux44fux437ux430ux43dux43dux44f-ux434ux438ux444ux435ux440ux435ux43dux446ux456ux430ux43bux44cux43dux438ux445-ux440ux456ux432ux43dux44fux43dux44c-ux443-ux43fux43eux445ux456ux434ux43dux438ux445-ux43fux43eux445ux456ux434ux43dux438ux445-ux432ux456ux434-ux43cux43eux43dux442ux435-ux43aux430ux440ux43bux43e}

МЕТА ЛЕКЦІЇ: Показати, як побудувати таку гру, результатом якої є
приблизне розв'язання диференціального рівняння. Одна з цих ігор (модель
випадкової ходьби) веде до скінченної різниці апроксимації задачі
Діріхле в квадраті. Цю модель можна узагальнити для отримання розв'язків
інших задач.

У попередній лекції було сказано, що можливо винайти такі ігри на удачу,
результатом яких є розв'язки (можливо, приблизні) задач для рівнянь з
частковими коефіцієнтами проникнення. У цій лекції ми покажемо, як
побудувати таку гру для наближеного розв'язання задачі Діріхле (PPP)
\(u_{xx} + u_{yy} = 0\) , 0 \textless{} x \textless{} 1, 0 \textless{} y
\textless{} 1,

(GI)

{[} u(x,y) = g(x,y) = {]}

{[} \textbackslash begin\{cases\} 1, \& \textbackslash text\{на верхней
стороне квадрата,\} \textbackslash{} 0, \& \textbackslash text\{на
боковых и пижней сторонах квадрата.\} \textbackslash end\{cases\} {]}

Після того как ця частная задача буде решена методом Монте-Карло, мы
обсудим, как решать этим методом более общие задачи.

Для иллюстрации метода Монте-Карло рассмотрим игру, яка називається
«Блуждающий пьяница». Чтобы играть в эту игру, необходим стол с
нанесенной на нем сеткой из линий, как показано на рис. 43.1.

Перейдем к описанию правил игры.

\section{Как играть в «Блуждающего
пьяницу»?}\label{ux43aux430ux43a-ux438ux433ux440ux430ux442ux44c-ux432-ux431ux43bux443ux436ux434ux430ux44eux449ux435ux433ux43e-ux43fux44cux44fux43dux438ux446ux443}

ШАГ 1. Блуждания пьяницы начинаются из произвольной точки сетки (в нашем
случае це точка A).

\pandocbounded{\includegraphics[keepaspectratio]{images/_page_321_Figure_4.jpeg}}

\begin{itemize}
\item
  А начальная точка - - внутренние точки сетки - □ -Конечная точка р;
\item
  \(g_i\) вознаграждение в m . \(p_I\) Рис. 43.1. Игровое поле для игры
  «Блуждающий пьяница».
\end{itemize}

ШАГ 2. На каждом ходе игры пьяница случайным образом перебирается в одну
из четырех соседних точек сетки (в нашем случае соседними для A являются
точки B, C, D и E и вероятность попадания в каждую из цих точек равна
0,25).

Таблица 43.1

\begin{verbatim}
Вероятность окончания случайного блуждания в граничной точке p_i и величина вознаграждения g_i | Граничная точка р <sub>і</sub> | P_{\mathcal{A}}\left(p_{i}\right) — относительная доля блужданий, закончившихся в точке p_{i} | g <sub>i</sub> — величина возна-<br>граждения при дости-<br>жении точки p <sub>i</sub> | | |
\end{verbatim}

\textbar-------------------------------------------------------\textbar---------------------------------------------------------------------------------------------------\textbar----------------------------------------------------------------------------------------\textbar--\textbar--\textbar{}

1234567891011 \textbar{}
0,040,150,030,060,170,050,060,150,030,060,160,04 \textbar{} 110000000
\textbar{} \textbar{} \textbar{}

ti

ШАГ 3. Після перехода в соседнюю точку процесс блуждания возобновляется.
Так пьяница идет от точки к точке до тех пор, пока случайным образом не
окажется в граничной точке \(p_t\) . Здесь он останавливается, а мы
записываем номер точки \(p_t\) . На цьому заканчивается одна случайная
прогулка.

ШАГ 4. Повторим щаги 1-3 достаточно много раз и найдем относительное
число прибытий пьяницы в каждую из граничных точек \(p_i\) . В табл.
43.1 приведены типичные результаты, полученные в результате 100
случайных прогулок.

\begin{verbatim}
ШАГ Б. Предположим, що пьяница получает вознаграждение g_i = g\left(p_i\right) (величина граничного условия в точке p_i , якщо он достиг точки p_i ) и предположим, що цель игры—обчислити среднее вознаграждение R\left(A\right) для всех случайных прогулок, начинающихся из точки A. Цей средний выигрыш определяется по формуле
\end{verbatim}

{[} R(A) = g\_1 P\_A(p\_1) + g\_2 P\_A(p\_2) + \textbackslash dots +
g\_\{12\} P\_A(p\_\{12\}). {]}

Игра завершается, когда величина \(R\left(A\right)\) найдена. В
соответствии с данными табл. 43.1 для \(R\left(A\right)\) получаем

{[} R(A) = 1 \textbackslash cdot (0.04) + 1 \textbackslash cdot (0.15) +
1 \textbackslash cdot (0.03) + 0 \textbackslash cdot (0.06) +
\textbackslash dot + 0 \textbackslash cdot (0.04). {]}

Теперь объясним, зачем нам понадобилось играть в такую игру.

\section{Для чего играть в «Блуждающего
пьяницу»?}\label{ux434ux43bux44f-ux447ux435ux433ux43e-ux438ux433ux440ux430ux442ux44c-ux432-ux431ux43bux443ux436ux434ux430ux44eux449ux435ux433ux43e-ux43fux44cux44fux43dux438ux446ux443}

Оказывается, що среднее вознаграждение, о котором лише що шла речь,
является приближенным решением задачи Дирихле в точке A. Це интересное
наблюдение основано на двух фактах.

\begin{enumerate}
\def\labelenumi{\arabic{enumi}.}
\item
  Предположим, що пъяница начал свою прогулку из точки, лежащей на
  еранице. Каждая такая прогулка немедленно заканчивается в той же
  точке, и пъяница получает вознаграждение \(g_i\) . Таким образом,
  среднее вознаграждение для каждой граничной точки равно \(g_i\) .
\item
  Теперь предположим, що прогулка нач нается из внутренней точки. Тогда
  ясно, що среднее вознаграждечие для точки \(R(A)\) буде средним
  арифметическим от средних вознаграждений для четырех соседних точек
\end{enumerate}

{[} R(A) = \textbackslash frac\{1\}\{4\} {[}R(B) + R(C) + R(D) +
R(E){]}. {]}

Зададимся еще раз вопросом: почему же среднее значение вознаграждения
\(R(A)\) аппроксимирует розв'язок задачи Дирихле в точке A? Мы видели,
що величина \(R(A)\) удовлетворяет двум

уравнениям

{[} R(A) = \textbackslash frac\{1\}\{4\} {[}R(B) + R(C) + R(D) + R(E){]}
{]}

(во внутренних точках), \(R(A)\) = \(g_i\) (в граничных точках).

Якщо \(g_i\) ---це значения функції \(g(x, y)\) из граничного условия в
граничных точках \(p_i\) , то два наших уравнения точно совпадают с
двумя уравнениями, які були получены при решении задачи Дирихле методом
конечных разностей. То есть величина \(R(A)\) соответствует величине
\(u_{i,f}\) в разностных уравнениях

{[} u\_\{i,j\} = \textbackslash frac\{1\}\{4\} \textbackslash left{[}
u\_\{i-1,j\} + u\_\{i+1,j\} + u\_\{i,j-1\} + u\_\{i,j+1\}
\textbackslash right{]}, {]}

(i, j)---внутренняя точка, \(u_{ij} = g_{ij}, g_{ij}\)---значение
розв'язки в граничной точке (i, j).

Отже, величина \(R(A)\) действительно аппроксимирует розв'язок уравнения
з частинними похідними в точке A.

Теперь можна считать, що игра в «Блуждающего пьяницу» состоит из трех
шагов.

\section{Розв'язок уравнения Лапласа методом
Монте-Карло}\label{ux440ux43eux437ux432ux44fux437ux43eux43a-ux443ux440ux430ux432ux43dux435ux43dux438ux44f-ux43bux430ux43fux43bux430ux441ux430-ux43cux435ux442ux43eux434ux43eux43c-ux43cux43eux43dux442ux435-ux43aux430ux440ux43bux43e}

\begin{itemize}
\item
  Приведенные ниже правила позволяют получить розв'язок в одной
  внутренней точке квадрата.
\item
  ШАГ 1. Совершим некоторое число случайных прогулок, начинающихся в
  точке А и заканчивающихся в одной из граничных точек. Проследим за
  тем, сколько раз прогулки заканчивались в каждой граничной точке.
\item
  ШАГ 2. Після того как все прогулки завершены, для каждой граничной
  точки вычисляем относительное число прогулок, завершившихся в цієї
  точке. Обозначим эти величины \(P_A\)(\(p_i\)) .

  \begin{itemize}
  \tightlist
  \item
    ШАГ 3. Вычисляем приближенное розв'язок u(A) по формуле u(A) = g\_1
    P\_A(p\_1) + g\_2 P\_A(p\_2) + \textbackslash dots + g\_N P\_A(p\_N)
    ,
  \end{itemize}
\end{itemize}

где \(g_i\) --- значение функції g в точке \(p_i\) , а N --- число
граничных точек.

Игру «Блуждающий пьяница» можна модернизировать, чтобы решить более
сложную задачу. Ниже приводится пример такой модернизации.

Розв'язок задачи Дирихле с змінними коэффициентами.

Рассмотрим эллиптическую краевую вадачу в квадрате

(PPP)

\(u_{xx}\) + (\textbackslash sin x)
\(u_{yy} = 0 , 0 < x < \pi , 0 < y < \pi , (\Gamma Y)\)u(x, y)\(=\)g(x,
y)\$\$на границе квадрата.

\begin{verbatim}
Чтобы решить эту задачу, заменим u_{xx} , u_{yy} и \sin x следующим образом:
\end{verbatim}

{[} u\_\{xx\}=u\_\{l,~j+1\}-2u\_\{l,~j\}+u\_\{l,~j-1\}, {]}

и \emph{\{yy\}=\(u_{l+1,\ j}\)-2u}\{l,~j\}+\(u_{l-1,\ j}\), (центральные
разностные производные),

и подставим их в уравнение з частинними похідними. Соответствующая сетка
показана на рис. 43.2. Разрешив получившееся уравнение относительно
\(u_{i,l}\) , получаем

{[} u\_\{i, j\} = \textbackslash frac\{u\_\{i, j+1\} + u\_\{i, j-1\} +
\textbackslash sin x\_j (u\_\{i+1, j\} + u\_\{i-1, j\})\}\{2 (1 +
\textbackslash sin x\_j)\}. {]}

Посмотрим внимательно на це уравнение. Коефіцієнти при
\(u_{i+1,\ j}\),~\(u_i\),~j+i,~\(u_l\),~j-1 н \(u_{i-1}\),~j
положительны, и их сумма равна единице. Другими словами, розв'язок
\(u_i\),~j является взвешенным средним решений в четырех соседних
точках. Отже, можна Р И С

\pandocbounded{\includegraphics[keepaspectratio]{images/_page_324_Picture_9.jpeg}}

. 43.2. Узлы сетки для случайного маршруга.

модифицировать «Блуждающего пьяницу» так, чтобы вероятности перехода в
соседние точки були не 0,25, а равнялись бы коэффициенту в
соответствующем члене. Другими словами, якщо пьяница находится в точке
(i, j), то он переходит в точку

(i, j+1)

з імовірністю \(\frac{1}{2(1+\sin x_j)}\) у точку \((i, j-1)\), з
імовірністю \(\frac{1}{2(1+\sin x_j)}\) у точку \((i+1, j)\), з
імовірністю \(\frac{\sin x_j}{2(1+\sin x_j)}\) у точку \((i-1, j)\) і з
імовірністю \(\frac{\sin x_j}{2(1+\sin x_j)}\).

В остальном игра не меняется. Модификация для других задач может
оказаться более хитроумной, но основная идея остается той же. Читатель
может заняться разработкой своих игр для розв'язки других задач.
Параболический случай рассмотрен в разделе «Задачи».

\section{ЗАМЕЧАНИЯ}\label{ux437ux430ux43cux435ux447ux430ux43dux438ux44f-18}

\begin{verbatim}
1. Отметим, що якщо относительные частоты попадания P_A(p_l) известны, то можна легко найти величину u\left(A\right) при других граничных условиях g_l . Для цього достаточно подставить g_l в формулу
\end{verbatim}

{[} u(A) = g\_1 P\_A(p\_1) + g\_2 P\_A(p\_2) + \textbackslash dots +
g\_N P\_A(p\_N), {]}

а не разыгрывать вновь случайные прогулки.

\begin{enumerate}
\def\labelenumi{\arabic{enumi}.}
\setcounter{enumi}{1}
\tightlist
\item
  Во многих случаях исследователя интересует розв'язок уравнения з
  частинними похідними в одной точке. Якщо граница области сложна и
  уравнение зависит от трех или четырех независимых змінних, метод
  Монте-Карло может оказаться удобным. Методы Монте-Карло возникли при
  решении очень трудных задач диффузии нейтронов. Эти задачи невозможно
  решать аналитически.
\end{enumerate}

\section{ЗАДАЧИ}\label{ux437ux430ux434ux430ux447ux438-19}

\begin{enumerate}
\def\labelenumi{\arabic{enumi}.}
\tightlist
\item
  Построить блок-схему игры «Блуждающий пьяница» для розв'язки задачи
\end{enumerate}

(УЧП)

{[} u\_\{xx\} + u\_\{yy\} = 0 {]}

, 0 \textless{} x \textless{} 1 , 0 \textless{} y \textless{} 1 , (ГУ)
\(u(x, y)\) = \(g(x, y)\) на границе,

\begin{itemize}
\item
  во внутренних узлах сетки, якщо число узлов сетки произвольно. 2.
  Напишите программу для розв'язки задачи по блок-схеме из задачи 1.
\item
  \begin{enumerate}
  \def\labelenumi{\arabic{enumi}.}
  \setcounter{enumi}{2}
  \tightlist
  \item
    Как модифицировать игру, чтобы решить задачу
  \end{enumerate}
\end{itemize}

(УЧП)

{[} u\_\{xx\} + x\^{}2 u\_\{yy\} = 0 {]}

, 0 \textless{} x \textless{} 1 , 0 \textless{} y \textless{} 1 , (ГУ)
\(u(x, y)\) = \(g(x, y)\) на границе?

Как выглядит в цієї игре случайная прогулка?

\begin{enumerate}
\def\labelenumi{\arabic{enumi}.}
\setcounter{enumi}{3}
\tightlist
\item
  Можете ли вы так модифицировать игру, чтобы решить задачу
\end{enumerate}

{[} (\textbackslash text\{UBP\}) \textbackslash qquad {]}

{[} \textbackslash begin\{aligned\} u\_\{xx\} + u\_\{yy\} = 0,
\textbackslash quad 0 \textless{} x \textless{} 1, \textbackslash quad 0
\textless{} y \textless{} 1, \textbackslash{} u(x, 1) = 0,
\textbackslash quad 0 \textless{} x \textless{} 1, \textbackslash quad 0
\textless{} y \textless{} 1, \textbackslash{} u(x, 0) = 0,
\textbackslash quad 0 \textless{} x \textless{} 1, \textbackslash quad 0
\textless{} y \textless{} 1, \textbackslash{} u(0, y) = 1,
\textbackslash quad \textbackslash frac\{\textbackslash partial
u\}\{\textbackslash partial x\}(1, y) = 0 \textbackslash end\{aligned\}
{]}

\begin{enumerate}
\def\labelenumi{\arabic{enumi}.}
\setcounter{enumi}{4}
\item
  Разработайте схему метода Монте-Карло для розв'язки смешанной
  параболической задачи

  \textbackslash begin\{aligned\} (\textbackslash text\{УЧП\}) \& u\_t =
  \textbackslash alpha\^{}2 u\_\{xx\} - \textbackslash beta u\_x -
  \textbackslash gamma u, \textbackslash quad 0 \textless{} x
  \textless{} 1, \textbackslash quad 0 \textless{} t \textless{}
  \textbackslash infty \textbackslash{} (\textbackslash text\{ГУ\}) \&
  \textbackslash left\{ \textbackslash begin\{aligned\}
  u\textbackslash left(0,~t\textbackslash right) =
  f\textbackslash left(t\textbackslash right), \& 0 \textless{} t
  \textless{} \textbackslash infty, \textbackslash{}
  u\textbackslash left(1,~t\textbackslash right) =
  g\textbackslash left(t\textbackslash right), \& 0 \textless{} t
  \textless{} \textbackslash infty, \textbackslash{}
  u\textbackslash left(x,~0\textbackslash right) =
  \textbackslash varphi\textbackslash left(x\textbackslash right), \& 0
  \textbackslash leqslant x \textbackslash leqslant 1.
  \textbackslash end\{aligned\} \textbackslash end\{aligned\}

  \textbackslash right. \textbackslash end\{aligned\}
\end{enumerate}

ПРИМІТКА. Замінимо це рівняння на скінченну різницю Кренка--Ніколсона,
виразимо \(u_{i+1,j}\) через п'ять значень \(u_{i+1,j-1}\) ,
\(u_{i+1,j+1}\) , \(u_{i,j+1}\) , \(u_{i,j+1}\) , \(u_{i,j+1}\) у
суміжних точках.

\section{ВАРІАЦІЙНЕ ЧИСЛЕННЯ (РІВНЯННЯ
ЕЙЛЕРА--ЛАГРАНЖА)}\label{ux432ux430ux440ux456ux430ux446ux456ux439ux43dux435-ux447ux438ux441ux43bux435ux43dux43dux44f-ux440ux456ux432ux43dux44fux43dux43dux44f-ux435ux439ux43bux435ux440ux430ux43bux430ux433ux440ux430ux43dux436ux430}

МЕТА ЛЕКЦІЇ: Ввести поняття функціоналу (функцію з функції) і пояснити,
як функціонали природно виникають у фізиці. Типовий функціонал є
інтегралом

{[} J{[}y{]} = \textbackslash int\_\{b\}\^{}\{a\} F(x, y, y') dx, {]}

що є функцією функції y (припускається, що інтегративна функція
F\textbackslash left(x,,y,,y'\textbackslash right) задана). Приклад
функціональності:

{[} J{[}y{]} = \textbackslash int\_\{0\}\^{}\{1\} {[}y\^{}\{2\}(x) +
y'\^{}\{2\}(x){]} dx. {]}

Ми покажемо, як знайти функцію \textbackslash overline\{y\}(x) , яка
мінімізує функціональність J{[}y{]}. Виявляється, що мінімізуюча функція
має \textbackslash overline\{y\} задовольняти так зване рівняння
Ейлера-Лагранжа. Це рівняння виконує ту ж роль, що й необхідна умова
мінімуму

{[} \textbackslash frac\{df(x)\}\{dx\} = 0 {]}

функції \(f(x)\) у точці x у диференціальному численні.

Варіаційне числення тісно пов'язане з диференціальними рівняннями, але,
на жаль, більшість студентів його не вивчає. Ця лекція (і наступна) є
вступом до варіаційного числення та показує, як можна розв'язувати
диференціальні рівняння в похідних на основі варіаційних принципів.

Калькулюс варіацій виник одночасно з математичним аналізом у зв'язку з
розв'язанням задач максимізації та мінімізації функцій функцій (які
називаються функціоналами). Першою задачею варіаційного числення була
задача брахістохрона, сформульована Йоганном Бернуллі у 1696 році.

Бернуллі показав, що час сходження записується у формі

{[} T = \textbackslash int\_\{0\}\^{}\{T\} dt =
\textbackslash int\_\{a\}\^{}\{L\} \textbackslash frac\{dt\}\{ds\} ds =
\textbackslash int\_\{0\}\^{}\{L\} \textbackslash frac\{ds\}\{v\} =
\textbackslash frac\{1\}\{\textbackslash sqrt\{2mg\}\}
\textbackslash int\_\{0\}\^{}\{L\}
\textbackslash frac\{ds\}\{\textbackslash sqrt\{y\}\} =
\textbackslash frac\{1\}\{\textbackslash sqrt\{2mg\}\}
\textbackslash int\_\{a\}\^{}\{b\}
\textbackslash sqrt\{\textbackslash frac\{1+y'\^{}\{2\}\}\{y\}\} dy {]}

і, отже, є функцією T\{y\} з функції. Оскільки багато функціоналів
впорядковані таким чином, ми зосередимося на вивченні загальної
функціональності

(44.1)

{[} J{[}y{]} = \textbackslash int\_\{a\}\^{}\{b\} F(x, y, y') dx. {]}

Розв'язок якої \(u_0\)(r, \textbackslash theta) = r \textbackslash cos
\textbackslash theta відомий. Як змінити \(u_0\)(r,
\textbackslash theta) , щоб отримати рішення проблеми (46.2)? Як уже
згадувалося, розглянемо клас задач \textbackslash Delta u +
\textbackslash epsilon u\^{}2 = 0 з граничними умовами і шукаємо
розв'язок кожної задачі з цього класу у вигляді

(46.4)

{[} u(r, \textbackslash theta) = u\_0(r, \textbackslash theta) +
\textbackslash varepsilon u\_1(r, \textbackslash theta) +
\textbackslash varepsilon\^{}2 u\_2(r, \textbackslash theta) +
\textbackslash dots {]}

Потім знаходимо розв'язок задачі (46.2), поклавши
\textbackslash varepsilon = 1\^{}\{\textbackslash circ\} ).

Вираз (46.4), який претендує на розв'язок, буде замінений на обурену
проблему

(PPP)

{[} \textbackslash Delta u + \textbackslash varepsilon u\^{}2 = 0 {]}

,

{[} (GI) u(1, \textbackslash theta) = \textbackslash cos
\textbackslash theta . {]}

У результаті ми отримуємо

{[} \textbackslash Delta (u\_0 + \textbackslash varepsilon u\_1 +
\textbackslash varepsilon\^{}2 u\_2 + \textbackslash ldots) +
\textbackslash varepsilon (u\_0 + \textbackslash varepsilon u\_1 +
\textbackslash varepsilon\^{}2 u\_2 + \textbackslash ldots)\^{}2 = 0,
{]}

{[} u\_0 (1, \textbackslash theta) + \textbackslash varepsilon u\_1 (1,
\textbackslash theta) + \textbackslash varepsilon\^{}2 u\_2 (1,
\textbackslash theta) + \textbackslash ldots = \textbackslash cos
\textbackslash theta. {]}

Виконуючи прості алгебраїчні перетворення та прирівнюючи коефіцієнти за
однаковими ступенями \(\varepsilon\) , ми отримуємо послідовність задач
\(P_0,\ P_1,\ P_2,\ \dots\) розв'язання, \(u_0,\ u_1,\ \dots\) невідомих
функцій (усі ці задачі лінійні).

Ця серія, очевидно, сходиться на \textbar{} є \textbar{} \textless{} 1
443.--- Ред. Ед

.

гетерогенні):

{[} P\_\{0\} {]}

{[} \textbackslash begin\{cases\} \textbackslash Delta u\_\{0\} = 0; \&
0 \textless{} r \textless{} 1, \textbackslash{} u\_\{0\}(1,
\textbackslash theta) = \textbackslash cos \textbackslash theta,
\textbackslash end\{cases\} {]}

{[} \textbackslash qquad u\_\{\textbackslash theta\}(r,
\textbackslash theta) = r \textbackslash cos \textbackslash theta, {]}

{[} P\_\{1\} {]}

{[} \textbackslash begin\{cases\} \textbackslash Delta u\_\{1\} =
-u\_\{0\}\^{}\{2\}, \textbackslash{} u\_\{1\}(1, \textbackslash theta) =
0, \textbackslash{} u\_\{2\}(1, \textbackslash theta).
\textbackslash end\{cases\} {]}

{[} P\_\{2\} {]}

{[} \textbackslash begin\{cases\} \textbackslash Delta u\_\{2\} =
-2u\_\{0\}u\_\{1\}, \textbackslash{} u\_\{2\}(1, \textbackslash theta).
\textbackslash end\{cases\} {]}

Почнемо з визначення функції \(u_1\left(r,\;\theta\right)\) з завдання
\(P_1\) (ця функція вже відома \(u_n\) ):

{[} P\_\{1\} {]}

{[} \textbackslash begin\{cases\} \textbackslash Delta u\_\{1\} =
-r\^{}\{2\} \textbackslash cos\^{}\{2\} \textbackslash theta =
-\textbackslash frac\{r\^{}\{2\}\}\{2\} \textbackslash left{[}1 +
\textbackslash cos(2\textbackslash theta)\textbackslash right{]} =
-\textbackslash frac\{r\^{}\{2\}\}\{2\} -
\textbackslash frac\{r\^{}\{2\}\}\{2\}
\textbackslash cos(2\textbackslash theta), \textbackslash{} u\_\{1\}(1,
\textbackslash theta) = 0. \textbackslash end\{cases\} {]}

Щоб розв'язати \emph{неоднорідне рівняння}, ми також звернемося до
методу, який використовувався в теорії звичайних диференціальних рівнянь
і складається з наступного:

\begin{itemize}
\item
  \begin{enumerate}
  \def\labelenumi{\arabic{enumi})}
  \tightlist
  \item
    знайти загальний розв'язок \(y_n\) однорідного рівняння;
  \end{enumerate}
\item
  \begin{enumerate}
  \def\labelenumi{\arabic{enumi})}
  \setcounter{enumi}{1}
  \tightlist
  \item
    знайти часткове розв'язок \(y_p\) гетерогенного рівняння;
  \end{enumerate}
\item
  \begin{enumerate}
  \def\labelenumi{\arabic{enumi})}
  \setcounter{enumi}{2}
  \tightlist
  \item
    Підставити \(y_h\) + \(y_p\) у початкову умову та визначити
    константи.
  \end{enumerate}

  Ми використаємо цей метод для розв'язання нашої проблеми. У нашій
  задачі P\_1 загальний розв'язок однорідного рівняння
  \textbackslash Delta u=0 має вигляд із розділеними змінними
\end{itemize}

{[} u\_h(r, \textbackslash theta) =
\textbackslash sum\_\{n=0\}\^{}\{\textbackslash infty\} r\^{}n {[}a\_n
\textbackslash cos(n\textbackslash theta) + b\_n
\textbackslash sin(n\textbackslash theta){]}. {]}

Часткове (лише одне) розв'язок неоднорідного рівняння

{[} \textbackslash Delta u = -\textbackslash frac\{r\^{}2\}\{2\} -
\textbackslash frac\{r\^{}3\}\{2\}
\textbackslash cos\{(2\textbackslash theta)\} {]}

Ми розглядаємо у формі

{[} u\_p(r, \textbackslash theta) = Ar\^{}4 + Br\^{}4
\textbackslash cos(2\textbackslash theta). {]}

Правильні частини форми r\^{}n , r\^{}n
\textbackslash cos\{(n\textbackslash theta)\} , r\^{}n
\textbackslash sin\{(n\textbackslash theta)\} ведуть до рішень
Ar\^{}\{n+2\} , Br\^{}\{n+2\}
\textbackslash cos\{(n\textbackslash theta)\} Cr\^{}\{n+2\}
\textbackslash sin\{(n\textbackslash theta)\} відповідно.

\begin{verbatim}
Підставляючи функції u_p(r, \theta) у гетерогенне рівняння, отримуємо A = -1/32, B = -1/24.
\end{verbatim}

Отже,

{[} u\_p(r, \textbackslash theta) = -\textbackslash frac\{r\^{}4\}\{32\}
- \textbackslash frac\{r\^{}4\}\{24\}
\textbackslash cos(2\textbackslash theta). {]}

Щоб розв'язати задачу, \(P_1\) має зробити останній крок --- підставити
загальне розв'язання u(r, \textbackslash theta) = \(u_h\)(r,
\textbackslash theta) + \(u_p\)(r, \textbackslash theta) на граничні
умови u(1, \textbackslash theta) = 0 і знайти коефіцієнти \(a_n\) і
\(b_n\) .

Після цього ми отримуємо

{[} \textbackslash sum\_\{n=0\}\^{}\{\textbackslash infty\} {[}a\_n
\textbackslash cos(n\textbackslash theta) + b\_n
\textbackslash sin(n\textbackslash theta){]} -
\textbackslash frac\{1\}\{32\} - \textbackslash frac\{1\}\{24\}
\textbackslash cos(2\textbackslash theta) = 0. {]}

Отже, \(a_0=1/32 , a_2=1/24\),і всі інші \(a_n\) і \(b_n\) дорівнюють
нулю. Отже, рішення проблеми \(P_1\) знайдено:

{[} u\_1(r, \textbackslash theta) = \textbackslash frac\{1\}\{32\} +
\textbackslash frac\{1\}\{24\}r\^{}2\textbackslash cos(2\textbackslash theta)
-
\textbackslash frac\{r\^{}4\}\{24\}\textbackslash cos(2\textbackslash theta)
- \textbackslash frac\{r\^{}4\}\{32\} = \textbackslash frac\{(r\^{}4 -
1)\}\{32\} - \textbackslash frac\{(r\^{}4 -
r\^{}2)\}\{24\}\textbackslash cos(2\textbackslash theta). {]}

Функція \(u_1\)(r, \textbackslash theta) називається пертурбацією
першого порядку до \(u_4\)(r, \textbackslash theta) . Додавши \(u_0\) і
\(u_1\) , отримуємо новий підхід до розв'язання задачі (46.2):

(46.5)

{[} u = u\_0 + u\_1 = r \textbackslash cos \textbackslash theta -
\textbackslash frac\{(r\^{}4 - 1)\}\{32\} - \textbackslash frac\{(r\^{}4
- r\^{}2)\}\{24\} \textbackslash cos (2\textbackslash theta). {]}

Щоб знайти збурення наступного порядку \(u_{z}\)(r,
\textbackslash theta) , потрібно розв'язати задачу

{[} P\_\{2\} {]}

{[} \textbackslash begin\{cases\} \textbackslash Delta u\_\{1\} =
-2u\_\{0\}u\_\{1\}, \textbackslash{} u\_\{2\}(1, \textbackslash theta) =
0, \textbackslash end\{cases\} {]}

\begin{verbatim}
де u_{\phi}(r,\theta) і u_{1}(r,0) вже відомі, раніше знайдені функції. Звісно, виконувати всі ці алгебраїчні перетворення без допомоги комп'ютера надзвичайно складно. На щастя, вираз (46.5) є майже точним рішенням нашої проблеми. Насправді, якщо підставити цей вираз зліва від \Delta u + u^{2} нашого нелінійного рівняння, отримаємо вираз, який майже всюди відрізняється від нуля в межах кола 0 < r < 1.
\end{verbatim}

Окрім розв'язання нелінійних задач, теорія збурень може
використовуватися для розв'язання задач у областях з неправильними
межами (звісно, якщо нерегулярність не дуже велика). Давайте розглянемо
простий приклад.

\section{Приклад проблеми з порушеною
межею}\label{ux43fux440ux438ux43aux43bux430ux434-ux43fux440ux43eux431ux43bux435ux43cux438-ux437-ux43fux43eux440ux443ux448ux435ux43dux43eux44e-ux43cux435ux436ux435ux44e}

\begin{verbatim}
Порушення можна порушити не лише рівняння з похідними похідними. Можливо знайти розв'язок рівняння Лапласа всередині деформованого кола, збуривши розв'язок рівняння Лапласа всередині кругової області. Наприклад, припустимо, ми хочемо знайти потенціал у площі, обмеженій кривою r=1+\frac{1}{4}\sin\theta (у полярних координатах). Нехай потенціал u на межі заданий. Тобто, нам потрібно вирішити проблему
\end{verbatim}

(46.6)

{[} \textbackslash Delta u = 0, \textbackslash quad 0 \textless{} r
\textless{} 1 + \textbackslash frac\{1\}\{4\} \textbackslash sin
\textbackslash theta, {]}

{[} u \textbackslash left( 1 + \textbackslash frac\{1\}\{4\}
\textbackslash sin \textbackslash theta, ~\textbackslash theta
\textbackslash right) = \textbackslash cos \textbackslash theta,
\textbackslash quad 0 \textbackslash le \textbackslash theta
\textbackslash le 2\textbackslash pi. {]}

Цю задачу можна розглядати як задачу визначення форми мильної плівки,
натягнутої на деформоване коло, висота якого вище цієї межею дорівнює
\(\cos\theta\) (див. рис. 46.2).

\pandocbounded{\includegraphics[keepaspectratio]{images/_page_349_Figure_5.jpeg}}

Рис. 46.2. Рівняння Лапласа всередині деформованого кола.

Оскільки загальна ідея теорії збурень полягає в тому, щоб звести складну
задачу до простої, припустимо, що наша межа умовна від деформованого
кола

{[}
u\textbackslash left(1+\textbackslash frac\{1\}\{4\}\textbackslash sin\textbackslash theta,~\textbackslash theta\textbackslash right)=\textbackslash cos\textbackslash theta
{]}

досягається в кінцевій точці сімейства граничних умов виду (див. рисунок
46 3)

{[} u(1+\textbackslash epsilon\textbackslash sin\textbackslash theta,
\textbackslash theta) = \textbackslash cos\textbackslash theta,
\textbackslash quad 0 \textbackslash le \textbackslash epsilon
\textbackslash le 1/4. {]}

Ця ідея приводить нас до серії про Тейлора. Використання серії Тейлора

{[} f(x+h) = f(x) + f'(x)h + f'\,'(x)\textbackslash frac\{h\^{}2\}\{2!\}
+ \textbackslash dots, {]}

Розкладемо складну граничну умову на ряд відповідно до простої межі Р

\pandocbounded{\includegraphics[keepaspectratio]{images/_page_350_Figure_2.jpeg}}

ис. 46 3. Схематичне представлення методу збурення.

умови, тобто

{[}
u(1+\textbackslash varepsilon\textbackslash sin\textbackslash theta,~\textbackslash theta)=u(1,~\textbackslash theta)+u\_r(1,~\textbackslash theta)(\textbackslash varepsilon\textbackslash sin\textbackslash theta)+u\_\{rr\}(1,~\textbackslash theta)\textbackslash frac\{(\textbackslash varepsilon\textbackslash sin\textbackslash theta)\^{}2\}\{21\}+\textbackslash ldots
{]}

Підставляючи цю декомпозицію в початкову задачу, отримуємо нову
еквівалентну задачу

(46.7)

{[} \textbackslash Delta u = 0, \textbackslash quad 0 \textless{} r
\textless{} 1 + \textbackslash varepsilon \textbackslash sin
\textbackslash theta, {]}

{[} u(1, \textbackslash theta) + u\_r(1, \textbackslash theta)
(\textbackslash varepsilon \textbackslash sin \textbackslash theta) +
u\_\{rr\}(1, \textbackslash theta)
\textbackslash frac\{(\textbackslash varepsilon \textbackslash sin
\textbackslash theta)\^{}2\}\{2!\} + \textbackslash dots =
\textbackslash cos \textbackslash theta. {]}

Коли всі розрахунки будуть виконані, ми внесемо
\textbackslash varepsilon=1/4 у кінцеві результати. Звісно, здається, що
цю задачу не можна легко розв'язати, але її можна звести до
послідовності задач, у кожній з яких визначена одна з функцій
u\_\{\textbackslash mathrm\{ \}e\} , u\_\{\textbackslash mathrm\{ \}e\}
, \ldots, що включені у розклад

{[} (46.8) u = u\_0 + \textbackslash varepsilon u\_1 +
\textbackslash varepsilon\^{}2 u\_2 + \textbackslash dots. {]}

Після підстави ряду в (467) отримуємо таку послідовність задач для
визначення функцій \(u_a\) , \(u_1\) , \ldots;

{[} P\_\{0\} {]}

{[} \textbackslash begin\{cases\} \textbackslash Delta u\_\{0\} = 0, \&
0 \textless{} r \textless{} 1 \textbackslash text\{ (внутри круга),\}
\textbackslash{} u\_\{0\}(1, \textbackslash theta) = \textbackslash cos
\textbackslash theta, \& u\_\{0\}(r, \textbackslash theta) = r
\textbackslash cos \textbackslash theta, \textbackslash end\{cases\} {]}

{[} P\_\{1\} {]}

{[} \textbackslash begin\{cases\} \textbackslash Delta u\_\{1\} = 0, \&
0 \textless{} r \textless{} 1 \textbackslash text\{ (внутри круга),\}
\textbackslash{} u\_\{1\}(1, \textbackslash theta) = -\textbackslash sin
\textbackslash theta \textbackslash frac\{\textbackslash partial
u\_\{0\}(1, \textbackslash theta)\}\{\textbackslash partial r\} =
-\textbackslash sin \textbackslash theta \textbackslash cos
\textbackslash theta \textbackslash end\{cases\} {]}

Отже, розв'язуючи кожну з задач Дірікла всередині кола для функцій
\(u_0\) , \(u_1\) , \(u_2\) , отримуємо розв'язок

(46.9)

{[} u = u\_0 + \textbackslash frac\{1\}\{4\}u\_1 +
\textbackslash left(\textbackslash frac\{1\}\{4\}\textbackslash right)\^{}2
u\_1 + \textbackslash cdots {]}

Задачі Діріхле в деформованій області (зверніть увагу, що тут

Ми вже взяли \(\epsilon = 1/4\) :)

Як вправу, читача просять знайти пертурбацію першого порядку \(u_1\) і
перевірити, наскільки добре є приблизне розв'язання

{[} u\_u + \textbackslash frac\{1\}\{4\} u\_t {]}

задовольняє рівняння (46 b)

\section{ЗАУВАЖЕННЯ}\label{ux437ux430ux443ux432ux430ux436ux435ux43dux43dux44f-7}

ŗ

Методи теорії збурень можна використовувати не лише для розв'язання
нелінійних рівнянь із частковими похідними. Наприклад, можливо
розв'язати задачу Коші.

(46 10)

{[} u\_t = (1+x) u\_\{xx\}, \textbackslash{} u(x, 0) = q(x),
\textbackslash quad -\textbackslash infty \textless{} x \textless{}
\textbackslash infty, {]}

якщо розглянути рівняння з параметром

(46.11)

{[} u\_1 = (1 + \textbackslash varepsilon\_1) u\_\{xx\} {]}

і побачити його розв'язок у вигляді

{[} u = u\_1 + \textbackslash varepsilon u\_1 +
\textbackslash varepsilon\^{}2 u\_2 + \textbackslash cdots {]}

Підставляючи цю декомпозицію на (46.11), отримуємо таку послідовність
завдань

{[} P\_\{0\} {]}

\begin{verbatim}
\begin{cases} \frac{\partial u_{0}}{\partial t} = \frac{\partial^{2} u_{1}}{\partial x^{2}}, \\ u_{0}(x, 0) = \varphi(x), \\ P_{1} \begin{cases} \frac{\partial u_{1}}{\partial t} - \frac{\partial^{2} u_{1}}{\partial x^{2}} = x \frac{\partial^{2} u_{0}}{\partial x^{2}}, \\ u_{1}(x, 0) = 0 \\ \vdots & \vdots \end{cases}
\end{verbatim}

Отметим, що в первых двух задачах коефіцієнти постоянны. Читатель
повинен понимать, що параметр в повинен быть достаточно мал, ниаче
бескопечный ряд может оказаться расходящимся

\section{ЗАДАЧИ}\label{ux437ux430ux434ux430ux447ux438-20}

\begin{verbatim}
- 1. Подставьте разложение (46.8) в задачу (46.7) и получите последовательность задач P_0,\ P_1,\ P_2,\ \dots . 2. Покажите, що нелинейную задачу
\end{verbatim}

{[} \textbackslash begin\{array\}\{ll\} \textbackslash Delta u + u\^{}2
= 0, \& 0 \textbackslash leqslant r \textless{} 1, \textbackslash quad 0
\textbackslash leqslant \textbackslash theta \textless{}
2\textbackslash pi, \textbackslash{} u(1, \textbackslash theta) =
\textbackslash cos \textbackslash theta, \& 0 \textbackslash leqslant
\textbackslash theta \textless{} 2\textbackslash pi,
\textbackslash end\{array\} {]}

можна звести до послідовності лінійних задач \(P_{o}\) , \(P_{i}\) ,
\(P_{z}\) .

\begin{itemize}
\tightlist
\item
  \begin{enumerate}
  \def\labelenumi{\arabic{enumi}.}
  \setcounter{enumi}{2}
  \tightlist
  \item
    Замініть нелінійну задачу 2 на приблизне розв'язання (46,5) і
    оцініть її точність.
  \end{enumerate}
\item
  \begin{enumerate}
  \def\labelenumi{\arabic{enumi}.}
  \setcounter{enumi}{3}
  \tightlist
  \item
    Розв'яжіть задачу \(P_1\) з прикладу з збуреною межою і перевірте,
    наскільки добре є наближене розв'язання
  \end{enumerate}
\end{itemize}

{[} u(r, \textbackslash theta) = u\_0(r, \textbackslash theta) +
\textbackslash frac\{1\}\{4\}u\_1(r, \textbackslash theta) {]}

задовольняє відношення

{[} \textbackslash Delta u = 0, {]}

(u\textbackslash left(1 +
\textbackslash frac\{1\}\{4\}\textbackslash sin\textbackslash theta,
\textbackslash theta\textbackslash right) =
\textbackslash cos\textbackslash theta).

\subsubsection{Лекція 47}\label{ux43bux435ux43aux446ux456ux44f-47}

\section{РОЗВ'ЯЗОК ДИФЕРЕНЦІАЛЬНИХ РІВНЯНЬ У ПОХІДНИХ ЗА ДОПОМОГОЮ
МЕТОДУ КОНФОРМНИХ
ВІДОБРАЖЕНЬ}\label{ux440ux43eux437ux432ux44fux437ux43eux43a-ux434ux438ux444ux435ux440ux435ux43dux446ux456ux430ux43bux44cux43dux438ux445-ux440ux456ux432ux43dux44fux43dux44c-ux443-ux43fux43eux445ux456ux434ux43dux438ux445-ux437ux430-ux434ux43eux43fux43eux43cux43eux433ux43eux44e-ux43cux435ux442ux43eux434ux443-ux43aux43eux43dux444ux43eux440ux43cux43dux438ux445-ux432ux456ux434ux43eux431ux440ux430ux436ux435ux43dux44c}

МЕТА ЛЕКЦІЇ: Показати, як деякі двовимірні граничні задачі можна
перетворити на інші, простіші задачі за допомогою конформного
відображення. Припустимо, наприклад, що ми хочемо розв'язати рівняння
Лапласа

{[} q\_\{xx\} + \textbackslash varphi\_\{uu\} = 0 {]}

з певною граничною умовою в області комплексної форми (площина змінних x
і y). Цю задачу можна перетворити на нову, в якій потрібно знайти
розв'язок рівняння Лапласа

{[} q\_\{uu\} + q\_\{vv\} = 0 {]}

у простішій області (площина змінних і і і).

\pandocbounded{\includegraphics[keepaspectratio]{images/_page_353_Figure_9.jpeg}}

\begin{verbatim}
У конформному відображенні рівняння Лапласа \varphi_{xx}+\varphi_{yy}=0 у координатній площині (x,y) повертається у рівняння Лапласа \varphi_{nn}\dashv -\varphi_{vo}=0 у координатній площині (u,v). Інші рівняння, такі як \varphi_{xx}+2\varphi_{xy}+\varphi_{yy}=0 , не мають такої властивості. Після розв'язання рівняння Лапласа \varphi(u,v) у простій області (

наприклад, у колі, напівплощині, квадраті) знайдено, достатньо підставити вирази U=U(X, Y), V=V(X, Y) у цей розв'язок, і ми отримаємо розв'язок нашої задачі \phi[u(x, y), v(x, y)] , вираженого через початкові змінні.
\end{verbatim}

У цій лекції ми розглянемо:

\begin{enumerate}
\def\labelenumi{\arabic{enumi})}
\item
  про властивості конформних відображень,
\item
  як їх будувати (для деяких районів),
\item
  як розв'язувати задачі методом конформних відображень.
\end{enumerate}

Однією з головних труднощів, що виникають при розв'язанні задач меж, є
складна форма кордонів. Навіть межі відносно простої форми часто роблять
задачу дуже складною для розв'язання. Звісно, можна спробувати
розв'язати задачу за допомогою методів теорії збурень, але теорія
збурень працює добре лише у випадках, коли межа близька до простої.

Однак існує спосіб розв'язати двовимірне рівняння Лапласа в областях з
комплексною межею. Цей метод базується на конформних відображеннях.

Перш ніж почати розв'язувати задачі цим методом, давайте трохи
ознайомимося з поняттям конформного відображення та комплексних змінних
загалом.

\section{Комплексні функції та конформні
відображення}\label{ux43aux43eux43cux43fux43bux435ux43aux441ux43dux456-ux444ux443ux43dux43aux446ux456ux457-ux442ux430-ux43aux43eux43dux444ux43eux440ux43cux43dux456-ux432ux456ux434ux43eux431ux440ux430ux436ux435ux43dux43dux44f}

У цій лекції ми ознайомимося лише з тими поняттями з теорії функцій
комплексної змінної, які нам знадобляться пізніше. Представимо
комплексне число z = x + iy як точку на площині комплексного xy (рис.
47.1).

Щоб ввести поняття конформного відображення, спершу ознайомимося з
поняттям функції комплексної змінної

{[} w = f(z) {]}

.

Тут z --- це комплексна змінна (з деякої області площини z), а w ---
нова комплексна змінна, яка отримає значення відповідно до формули w =
\(f(z)\) у новій комплексній площині w.

Наприклад, комплексна функція, визначена в \emph{першому квадранті}
комплексної z-площини, w=z\^{}2 відображає цю область у всю півплощину
v\textgreater0 у комплексній площині w.

Таблиця 47 2

\begin{verbatim}
| Площина d | z \longrightarrow z^2 | Площина S |
\end{verbatim}

\textbar--------------------------\textbar-------------------------\textbar------------------------------------------------------------\textbar{}

0 \textbar{} \textbar{} 0 \textbar{}

t1 → t \textbar{} \textbar{} -1(2) \textbar{}

Вісь x Перший квадрант \textbar{} \textbar{} Позитивна дійсна піввісь
верхньої півплощини \textbar{}

Таблиця 47.2 показує відповідність, яка встановлюється шляхом
відображення w=z\^{}2 між деякими точками на площинах z і w.

\pandocbounded{\includegraphics[keepaspectratio]{images/_page_355_Figure_5.jpeg}}

Рис. 47.1. Площина комплексної змінної \(\boldsymbol{z}\) та
некольорових формул бігунів

\pandocbounded{\includegraphics[keepaspectratio]{images/_page_355_Figure_7.jpeg}}

РИСУНОК 47.2. Відображення першого квадранта площини \{\textbackslash bf
z\} у верхню півплощину \{\textbackslash bf w\} .

Інакше кажучи, якщо ми говоримо про комплексісні функції змінної w =
\(f(z)\), то можна розглянути питання, як криві площини z перетворюються
на криві в площині

Це питання не ставиться в теорії функцій дійсної змінної, з якою читач
уже добре знайомий.

Перш ніж рухатися далі, давайте розглянемо, як w=z\^{}2 складне
відображення типів у еквівалентній дійсній формі. Дійсна форма показує,
як початкові координати (x,y) у площині z перетворюються на нові
координати (u,v) у площині w w=z\^{}2 .

{[} u + iv = (x + iy)\^{}2 = x\^{}2 - y\^{}2 + 2ixy {]}

І давайте окремо прирівняємо реальні та уявні частини в цьому
співвідношенні. У результаті ми отримуємо

\begin{verbatim}
\left\{
\end{verbatim}

{[} \textbackslash begin\{aligned\}
u=x\textsuperscript{2-y}2,\textbackslash{} v=2xy,
\textbackslash end\{aligned\} {]}

{[} \textbackslash Rightarrow. {]}

(действительная форма отображения w=z\^{}2 ).

Ця форма понадобится нам позже. В цьому примере она интересна тем, що
показывает, как гиперболы в плоскости z переходят в координатные линии u
= const и v = const в плоскости w.

Теперь мы перейдем к изучению одного частного вида функцій комплексного
переменного, які осуществляют конформные отображения.

\section{Определение конформного
отображения}\label{ux43eux43fux440ux435ux434ux435ux43bux435ux43dux438ux435-ux43aux43eux43dux444ux43eux440ux43cux43dux43eux433ux43e-ux43eux442ux43eux431ux440ux430ux436ux435ux43dux438ux44f}

\begin{verbatim}
Отображение w=f(z) комплексиой плоскости z на комплексную плоскость w називається конформным в точке z_{\bullet} плоскости z, якщо производная f'(z_{\bullet}) \neq 0 . Отображение f(z) називається конформным в области D, якщо f'(z) \neq 0 в каждой точке области D.

Наприклад, отображение f(z) = z^2 конформно всюду, за исключением точки z = 0, поскольку z'(z) = 2z \neq 0 для всех z \neq 0 . C другой стороны, отображение e^z конформно во всей плоскости z, так как всюду f'(z) = e^z \neq 0 . Какая же польза от конформных отображений? Ответ состоит в том, що, решая рівняння Лапласа \phi_{xx} + \phi_{yy} = 0 внутри какой-то области плоскости змінних x и y, мы можем рассматривать эту плоскость как плоскость комплексного переменного z. Рассмотрим комплексное отображение w = f(z) плоскости z на плоскость w, в которой введены координаты (u, v). При отображении w = f(z) рівняння Лапласа \varphi_{xx} + \varphi_{yy} = 0 преобразуется в некоторое новое уравнение з частинними похідними, зависящее от новых координат и и v. Замечательным является следующее: якщо отображение w = f(z) является конформным в области, где задано уравнение \phi_{xy} + \phi_{uy} = 0 , то новое уравнение також буде уравнением Лапласа в координатах и и v. To есть уравнение \phi_{xx} + \phi_{uu} = 0 перейдет снова в уравнение \varphi_{nn} + \varphi_{nn} = 0 . Так мы приходим к идее найти такое конформное отображение, яке переводит область со сложной границей в область с простой границей (напомним, що отображение w=z^* переводит границу первого квадранта плоскости z в действительную ось плоскости w).
\end{verbatim}

Приведем несколько примеров, в которых буде показано, как можна решать
уравнения з частинними похідними с помощью конформных отображений.

Рівняння лапласа в верхней полуплоскости.

Предположим, що нам необходимо решить следующую задачу Дирихле в верхней
полувлоскости (рис. 47.3):

(47.1)

\begin{verbatim}
\begin{aligned} (Y & \Pi) & \varphi_{xx} + \varphi_{yy} = 0, & -\infty < x < \infty, & 0 < y < \infty, \\ (47.1) & \varphi(x, 0) = \begin{cases} 0, & |x| > 1, \\ 1, & |x| \le 1. \end{cases}
\end{verbatim}

Один із способів розв'язання цієї проблеми --- застосувати перетворення
Фур'є до змінної x, але набагато краще конформно відобразити верхню
половину площини у площину (u, v). Читачеві не буде складно побачити, що
конформне відображення

{[} w =
\textbackslash ln\textbackslash left\{\textbackslash frac\{z-1\}\{z+1\}\textbackslash right\}
{]}

Вона має такі властивості (див. RNS. 47.4).

\begin{verbatim}
1. Вона відображає верхню півплощину площини z у смуги -\infty < u < \infty , 0 < v < \pi , у площині w.
\end{verbatim}

\pandocbounded{\includegraphics[keepaspectratio]{images/_page_357_Figure_11.jpeg}}

РИС 47.3 Задача граничного значення у верхній половині площини.

\begin{itemize}
\tightlist
\item
  \begin{enumerate}
  \def\labelenumi{\arabic{enumi}.}
  \setcounter{enumi}{1}
  \tightlist
  \item
    Ребро прямої y=0, -1 \textless{} x \textless{} 1, у площині z (де
    потенціал дорівнює 1), переходить у пряму \(v=\pi,\) -
    \(\infty < u < \infty\) , у площині w. - 3. Два промені
    x\textgreater1, y=0, і x\textless-1, y=0, у площині z проходять
    відповідно у позитивний і негативний
  \end{enumerate}
\end{itemize}

Читач повинен бути впевнений, що може підтвердити ці твердження.

\pandocbounded{\includegraphics[keepaspectratio]{images/_page_358_Figure_2.jpeg}}

Рис. 47.4. Перетворення складної задачі на легку за допомогою
конформного відображення.

Важливість цих властивостей полягає в тому, що вони перетворюють
початкову вадаху (47.1) на дуже просту задачу

(47.2)

{[} \textbackslash begin\{cases\} (\textbackslash text\{УЧП\}) \&
\textbackslash varphi\_\{uu\} + \textbackslash varphi\_\{vv\} = 0,
\textbackslash quad -\textbackslash infty \textless{} u \textless{}
\textbackslash infty, \textbackslash quad 0 \textless{} v \textless{}
\textbackslash pi, \textbackslash{} (\textbackslash text\{47.2\}) \&
\textbackslash begin\{cases\} \textbackslash varphi(u, 0) = 0,
\textbackslash{} \textbackslash varphi(u, \textbackslash pi) = 1.
\textbackslash end\{cases\} \textbackslash end\{cases\} {]}

Розв'язок очевидний: \(\varphi(u, v) = \frac{1}{\pi}v\) .

Тепер, щоб знайти розв'язок початкової задачі, потрібно пояснити v через
x і y і підставити цей вираз на \(q(u,v)=v'\pi\) . Після цього ми
отримуємо

{[} w = u + iv =
\textbackslash ln\textbackslash left{[}\textbackslash frac\{z-1\}\{z+1\}\textbackslash right{]}
= {]}

{[} =
\textbackslash ln\textbackslash left{[}\textbackslash frac\{z-1\}\{z+1\}\textbackslash right{]}
+ i
\textbackslash arg\textbackslash left{[}\textbackslash frac\{z-1\}\{z+1\}\textbackslash right{]}
{]}

і, відповідно,

{[} v =
\textbackslash arg\textbackslash left{[}\textbackslash frac\{z-1\}\{z+1\}\textbackslash right{]}
=
\textbackslash arg\textbackslash left{[}\textbackslash frac\{x+iy-1\}\{x+1+iy\}\textbackslash right{]}
= {]}

{[} =
\textbackslash arg\textbackslash left{[}\textbackslash frac\{x\textsuperscript{2+y}2-1+i2y\}\{(x+1)\textsuperscript{2+y}2\}\textbackslash right{]}
= {]}

{[} =
\textbackslash arg\textbackslash left{[}\textbackslash frac\{2y\}\{x\textsuperscript{2+y}2-1\}\textbackslash right{]}
{]}

(згідно з загальним правилом
\textbackslash arg\{(x+iy)\}=\textbackslash arg\{(\(y^{\dagger}\)x)\} ).
Отже, розв'язок задачі (47.1) має вигляд

{[} \textbackslash varphi(x, y) =
\textbackslash frac\{1\}\{\textbackslash pi\}
\textbackslash arctan\textbackslash left{[}\textbackslash frac\{2y\}\{x\^{}2
+ y\^{}2 - 1\}\textbackslash right{]}. {]}

Щоб візуально зрозуміти, як працює ця функція, ми радимо читачеві
побудувати її за кількома лініями y=c.

У другому прикладі ми перетворюємо область між двома неконцентричними
колами у кільце.

Задача Діріхле в області між двома некопоцентричними колами

Припустимо, ви хочете знайти потенціал між двома колами

{[} x\^{}\{2\} + y\^{}\{2\} = 1, {]}

{[} (x-1)\^{}\{2\} + y\^{}\{2\} = 9, {]}

якщо потенціал на внутрішньому колі дорівнює одинці, а на зовнішньому
колі дорівнює двом (рис. 47.5). Іншими словами, нам потрібно розв'язати
проблему

(47.3)

{[} \textbackslash begin\{aligned\} \& (\textbackslash text\{УЧП\}) \&
\textbackslash phi\_\{xx\} + \textbackslash phi\_\{yy\} = 0, \&
\textbackslash text\{внутри \} D, \textbackslash{} \&
(\textbackslash text\{ГУ\}) \& \textbackslash begin\{cases\}
\textbackslash varphi(x, y) = 1 \& \textbackslash text\{на \} x\^{}2 +
y\^{}2 = 1, \textbackslash{} \textbackslash varphi(x, y) = 2 \&
\textbackslash text\{на \} (x - 1)\^{}2 + y\^{}2 = 9.
\textbackslash end\{cases\} \textbackslash end\{aligned\} {]}

Проблема полягає в тому, щоб знайти конформне відображення, яке
перетворює цю область на щось простіше, де нашу проблему легко
розв'язати. У цьому випадку очевидно, що ця проста ділянка має бути
кільцем. Відповідна експозиція

\pandocbounded{\includegraphics[keepaspectratio]{images/_page_359_Figure_10.jpeg}}

РИСУНОК 47.5. Конформне відображення на кільце.

Це вже не так очевидно, але існують довідники, де можна знайти дисплей
із вказаними властивостями. Одним із найзагальних є «Довідник конформних
відображень»

G. Kober 1), яка містить сотні конформних відображень між різними
регіонами. Якщо читач витратить трохи часу і краще ознайомиться з
конформними таблицями відображення, він знайде відповідне конформне
відображення для задачі (47.3). У нашому конкретному випадку це буде
образ виду

(47.4)

{[} w = 2t \textbackslash left{[} \textbackslash frac\{z-s\}\{z-t\}
\textbackslash right{]}, {]}

де s=-0,146, а t=-6,85. У своїй еквівалентній дійсній формі це буде

(47.5)

{[} u = \textbackslash gamma {[}(x-s)(x-t)-\textbackslash gamma\^{}2{]},
{]}

{[} v = \textbackslash gamma {[}y(x-t){]}, {]}

де

{[} \textbackslash gamma = 2t/{[}(x-t)\^{}2 + y\^{}2{]}. {]}

Під впливом цього конформного відображення область між неконцентричними
колами переходить у кільце, межі якого визначаються рівняннями
\(u^2 + v^2 = 1\) і \(u^2 + v^2 = 6,86\) (читач може перевірити це сам)

Якщо застосувати це конформне відображення до задачі (47 3), отримаємо
задачу Діріхле всередині кільця

(47.6)

{[} \textbackslash begin\{aligned\} \textbackslash varphi\_\{uu\} +
\textbackslash varphi\_\{vv\} \&= 0, \& \textbackslash text\{в кольче,\}
\textbackslash{} (\textbackslash Gamma Y) \& \& \& \& \& \& \&
\textbackslash{} (\textbackslash Gamma Y) \& \& \& \& \& \&
\textbackslash{} (\textbackslash Gamma Y) \& \& \& \& \&
\textbackslash{} \& \& \& \& \& \textbackslash{} (\textbackslash Gamma
Y) \& \& \& \& \& \textbackslash{} \& \& \& \& \textbackslash{} \& \& \&
\& \textbackslash{} \& \& \& \& \textbackslash{} \& \& \& \&
\textbackslash{} \& \& \& \textbackslash{} \& \& \& \textbackslash{} \&
\& \& \textbackslash{} \& \& \& \textbackslash{} \& \& \&
\textbackslash{} \& \& \& \textbackslash{} \& \& \textbackslash{} \& \&
\textbackslash{} \& \& \textbackslash{} \& \& \textbackslash{} \& \&
\textbackslash{} \& \& \textbackslash{} \& \& \textbackslash{} \& \&
\textbackslash{} \& \& \textbackslash{} \& \& \textbackslash{} \& \&
\textbackslash{} \& \& \textbackslash{} \& \& \textbackslash{} \& \&
\textbackslash{} \& \& \textbackslash{} \& \& \textbackslash{} \& \&
\textbackslash{} \& \& \textbackslash{} \& \& \textbackslash{} \& \&
\textbackslash{} \& \& \textbackslash{} \& \& \textbackslash{} \& \&
\textbackslash{} \& \& \textbackslash{} \& \& \textbackslash{} \& \&
\textbackslash{} \& \& \textbackslash{} \& \& \textbackslash{} \& \&
\textbackslash{} \& \& \textbackslash{} \& \& \textbackslash{} \& \&
\textbackslash{} \& \& \textbackslash{} \& \& \textbackslash{} \& \&
\textbackslash{} \& \& \textbackslash{} \& \& \textbackslash{} \& \&
\textbackslash{} \& \& \textbackslash{} \& \& \textbackslash{} \& \&
\textbackslash{} \& \& \textbackslash{} \& \& \textbackslash{} \& \&
\textbackslash{} \& \& \textbackslash{} \& \& \textbackslash{} \& \&
\textbackslash{} \& \& \textbackslash{} \& \& \textbackslash{} \& \&
\textbackslash{} \& \& \textbackslash{} \& \& \textbackslash{} \& \&
\textbackslash{} \& \& \textbackslash{} \& \& \textbackslash{} \& \&
\textbackslash{} \& \& \textbackslash{} \& \& \textbackslash{} \& \&
\textbackslash{} \& \& \textbackslash{} \& \& \textbackslash{} \& \&
\textbackslash{} \& \& \textbackslash{} \& \& \textbackslash{} \& \&
\textbackslash{} \& \& \textbackslash{} \& \& \textbackslash{} \& \&
\textbackslash{} \& \& \textbackslash{} \& \& \textbackslash{} \& \&
\textbackslash{} \& \& \textbackslash{} \& \& \textbackslash{} \& \&
\textbackslash{} \& \& \textbackslash{} \& \& \textbackslash{} \& \&
\textbackslash{} \& \& \textbackslash{} \& \& \textbackslash{} \& \&
\textbackslash{} \& \& \textbackslash{} \& \& \textbackslash{} \& \&
\textbackslash{} \& \& \textbackslash{} \& \& \textbackslash{} \& \&
\textbackslash{} \& \& \textbackslash{} \& \& \textbackslash{} \& \&
\textbackslash{} \& \& \textbackslash{} \& \& \textbackslash{} \& \&
\textbackslash{} \& \& \textbackslash{} \& \& \textbackslash{} \& \&
\textbackslash{} \& \& \textbackslash{} \& \& \textbackslash{} \& \&
\textbackslash{} \& \& \textbackslash{} \& \& \textbackslash{} \& \&
\textbackslash{} \& \& \textbackslash{} \& \& \textbackslash{} \& \&
\textbackslash{} \& \& \textbackslash{} \& \& \textbackslash{} \& \&
\textbackslash{} \& \& \textbackslash{} \& \& \textbackslash{} \& \&
\textbackslash{} \& \& \textbackslash{} \& \& \textbackslash{} \& \&
\textbackslash{} \& \& \textbackslash{} \& \& \textbackslash{} \& \&
\textbackslash{} \& \& \textbackslash{} \& \& \textbackslash{} \& \&
\textbackslash{} \& \& \textbackslash{} \& \& \textbackslash{} \& \&
\textbackslash{} \& \& \textbackslash{} \& \& \textbackslash{} \& \&
\textbackslash{} \& \& \textbackslash{} \& \& \textbackslash{} \& \&
\textbackslash{} \& \& \textbackslash{} \& \& \textbackslash{} \& \&
\textbackslash{} \& \& \textbackslash{} \& \& \textbackslash{} \& \&
\textbackslash{} \& \& \textbackslash{} \& \& \textbackslash{} \& \&
\textbackslash{} \& \& \textbackslash{} \& \& \textbackslash{} \& \&
\textbackslash{} \& \& \textbackslash{} \& \& \textbackslash{} \& \&
\textbackslash{} \& \& \textbackslash{} \& \& \textbackslash{} \& \&
\textbackslash{} \& \& \textbackslash{} \& \& \textbackslash{} \& \&
\textbackslash{} \& \& \textbackslash{} \& \& \textbackslash{} \& \&
\textbackslash{} \& \& \textbackslash{} \& \& \textbackslash{} \& \&
\textbackslash{} \& \& \textbackslash{} \& \& \textbackslash{} \& \&
\textbackslash{} \& \& \textbackslash{} \& \& \textbackslash{} \& \&
\textbackslash{} \& \& \textbackslash{} \& \& \textbackslash{} \& \&
\textbackslash{} \& \& \textbackslash{} \& \& \textbackslash{} \& \&
\textbackslash{} \& \& \textbackslash{} \& \& \textbackslash{} \& \&
\textbackslash{} \& \& \textbackslash{} \& \& \textbackslash{} \& \&
\textbackslash{} \& \& \textbackslash{} \& \& \textbackslash{} \& \&
\textbackslash{} \& \& \textbackslash{} \& \& \textbackslash{} \& \&
\textbackslash{} \& \& \textbackslash{} \& \& \textbackslash{} \& \&
\textbackslash{} \& \& \textbackslash{} \& \& \textbackslash{} \& \&
\textbackslash{} \& \& \textbackslash{} \& \& \textbackslash{} \& \&
\textbackslash{} \& \& \textbackslash{} \& \& \textbackslash{} \& \&
\textbackslash{} \& \& \textbackslash{} \& \& \textbackslash{} \& \&
\textbackslash{} \& \& \textbackslash{} \& \& \textbackslash{} \& \&
\textbackslash{} \& \& \textbackslash{} \& \& \textbackslash{} \& \&
\textbackslash{} \& \& \textbackslash{} \& \& \textbackslash{} \& \&
\textbackslash{} \& \& \textbackslash{} \& \& \textbackslash{} \& \&
\textbackslash{} \& \& \textbackslash{} \& \& \textbackslash{} \& \&
\textbackslash{} \& \& \textbackslash{} \& \& \textbackslash{} \& \&
\textbackslash{} \& \& \textbackslash{} \& \& \textbackslash{} \& \&
\textbackslash{} \& \& \textbackslash{} \& \& \textbackslash{} \& \&
\textbackslash{} \& \& \textbackslash{} \& \& \textbackslash{} \& \&
\textbackslash{} \& \& \textbackslash{} \& \& \textbackslash{} \& \&
\textbackslash{} \& \& \textbackslash{} \& \& \textbackslash{} \& \&
\textbackslash{} \& \& \textbackslash{} \& \& \textbackslash{} \& \&
\textbackslash{} \& \& \textbackslash{} \& \& \textbackslash{} \& \&
\textbackslash{} \& \& \textbackslash{} \& \& \textbackslash{} \& \&
\textbackslash{} \& \& \textbackslash{} \& \& \textbackslash{} \& \&
\textbackslash{} \& \& \textbackslash{} \& \& \textbackslash{} \& \&
\textbackslash{} \& \& \textbackslash{} \& \& \textbackslash{} \& \&
\textbackslash{} \& \& \textbackslash{} \& \& \textbackslash{} \& \&
\textbackslash{} \& \& \textbackslash{} \& \& \textbackslash{} \& \&
\textbackslash{} \& \& \textbackslash{} \& \& \textbackslash{} \& \&
\textbackslash{} \& \& \textbackslash{} \& \& \textbackslash{} \& \&
\textbackslash{} \& \& \textbackslash{} \& \& \textbackslash{} \& \&
\textbackslash{} \& \& \textbackslash{} \& \& \textbackslash{} \& \&
\textbackslash{} \& \& \textbackslash{} \& \& \textbackslash{} \& \&
\textbackslash{} \& \& \textbackslash{} \& \& \textbackslash{} \& \&
\textbackslash{} \& \& \textbackslash{} \& \& \textbackslash{} \& \&
\textbackslash{} \& \& \textbackslash{} \& \& \textbackslash{} \& \&
\textbackslash{} \& \& \textbackslash{} \& \& \textbackslash{} \& \&
\textbackslash{} \& \& \textbackslash{} \& \& \textbackslash{} \& \&
\textbackslash{} \& \& \textbackslash{} \& \& \textbackslash{} \& \&
\textbackslash{} \& \& \textbackslash{} \& \& \textbackslash{} \& \&
\textbackslash{} \& \& \textbackslash{} \& \& \textbackslash{} \& \&
\textbackslash{} \& \& \textbackslash{} \& \& \textbackslash{} \& \&
\textbackslash{} \& \& \textbackslash{} \& \& \textbackslash{} \& \&
\textbackslash{} \& \& \textbackslash{} \& \& \textbackslash{} \& \&
\textbackslash{} \& \& \textbackslash{} \& \& \textbackslash{} \& \&
\textbackslash{} \& \& \textbackslash end\{aligned\} {]}

Розв'язати цю задачу досить просто, оскільки розв'язок має бути
радіально симетричним, а радіально симетричний розв'язок рівняння
Лапласа, як відомо, має вигляд

{[} \textbackslash varphi(r) = a \textbackslash ln r + b,
\textbackslash quad r = u\^{}2 + v\^{}2. {]}

Підставляючи це рівняння у граничну умову (47,6), отримуємо

{[} \textbackslash varphi(u, v) = 0.57 \textbackslash ln(u\^{}2 +
v\^{}2) + 1 {]}

Якщо повернутися до початкових координат x і y відповідно до
відображення (47.5), отримаємо розв'язок задачі (47.3)

{[} \textbackslash varphi(x, y) = 0.57 \textbackslash ln(u\^{}2 +
v\^{}2) + 1, {]}

де w і v визначаються за рівняннями (47 5).

Бачиш. Див. також Коппенфельс В., Штальман Ф. Практика конформних
відображень. --- Москва, Іллінойс, 1963. --- Примітка редактора.

\section{ЗАУВАЖЕННЯ}\label{ux437ux430ux443ux432ux430ux436ux435ux43dux43dux44f-8}

\begin{enumerate}
\def\labelenumi{\arabic{enumi}.}
\item
  Для побудови конформного відображення області з достатньо загальним
  виглядом у площині z на верхній площині w використовується
  відображення Крістоффеля-Шварца (коли область відображена у верхню
  півплощину, тоді цю напівплощину можна відобразити в іншу область).
  Інтеграл Крістоффеля--Шварца можна знайти у книзі {[}28{]} з
  рекомендованої літератури.
\item
  Метод конформних відображень має обмежені можливості, оскільки op
  застосовний лише до двовимірного рівняння Лапласа (і, з незначними
  змінами, до рівняння
\end{enumerate}

Пуассон).

\section{ЗАВДАННЯ}\label{ux437ux430ux432ux434ux430ux43dux43dux44f-11}

\begin{enumerate}
\def\labelenumi{\arabic{enumi}.}
\tightlist
\item
  Визначити область відповідності відображення
\end{enumerate}

{[} w = \textbackslash ln \textbackslash left{[}
\textbackslash frac\{z-1\}\{z+1\} \textbackslash right{]}. {]}

Переконайтеся, що це відображення переносить верхню половину площини
змінної z у зону --- \(\infty < u < \infty\) , \(0 < v < \pi\) ,
w-площину.

\begin{enumerate}
\def\labelenumi{\arabic{enumi}.}
\setcounter{enumi}{1}
\tightlist
\item
  Яке зображення першого квадранта при відображенні \(w=z^{8}\) ? Щоб
  розв'язати задачу, корисно перейти до експоненціальної форми
  комплексного числа
\end{enumerate}

{[} z = re\^{}\{i\textbackslash theta\} {]}

.

\begin{enumerate}
\def\labelenumi{\arabic{enumi}.}
\setcounter{enumi}{2}
\item
  За допомогою конформного відображення \(w=z^2\) розв'язати задачу
  Діріхле у першому квадранті

  \textbackslash begin\{aligned\} (\textbackslash text\{Y\}
  \textbackslash forall \textbackslash Pi) \& \textbackslash qquad
  \textbackslash phi\_\{\textbackslash text\{vx\}\} +
  \textbackslash phi\_\{yy\} = 0, \textbackslash quad 0 \textless{} x
  \textless{} \textbackslash infty, \textbackslash quad 0 \textless{} y
  \textless{} \textbackslash infty, \textbackslash{}
  \textbackslash left\{ \textbackslash begin\{aligned\}
  \textbackslash phi(x, ~0) = \textbackslash begin\{cases\} 1, \& 0
  \textless{} x \textless{} 1, \textbackslash{} 0, \& 1
  \textbackslash leqslant x \textless{} \textbackslash infty,
  \textbackslash end\{cases\} \textbackslash{} \textbackslash phi(0, ~y)
  = \textbackslash begin\{cases\} 1, \& 0 \textless{} y \textless{} 1,
  \textbackslash{} 0, \& 1 \textbackslash leqslant y \textless{}
  \textbackslash infty. \textbackslash end\{cases\}
  \textbackslash end\{aligned\} \textbackslash end\{aligned\}
\end{enumerate}

{[} \textbackslash right. {]}

Зверніть увагу на зображення.

\pandocbounded{
\includegraphics[keepaspectratio]{images/_page_362_Figure_1.jpeg}}

\pandocbounded{\includegraphics[keepaspectratio]{images/_page_362_Figure_2.jpeg}}

\begin{enumerate}
\def\labelenumi{\arabic{enumi}.}
\setcounter{enumi}{3}
\tightlist
\item
  Розв'язати змішану задачу Діріхле-Неймана з кутовим порошком за
  допомогою розчину під кутом 45°
\end{enumerate}

\pandocbounded{\includegraphics[keepaspectratio]{images/_page_362_Figure_5.jpeg}}

\begin{verbatim}
ПРИМІТКА. Комплексна функція w = \ln z = \ln |z| + i \arg(z) відображає промінь \theta = c_1 у площині z у пряму v = c_1 , а коло r = c_2 у пряму u = \ln c_2 у площині w (див. рисунки).
\end{verbatim}

\pandocbounded{\includegraphics[keepaspectratio]{images/_page_362_Figure_7.jpeg}}

\# ТАБЛИЦІ ІНТЕГРАЛЬНИХ ПЕРЕТВОРЕНЬ

Таблиця A: Експоненціальне перетворення Фур'є,

Таблиця B: Перетворення синуса Фур'є.

Таблиця St Cosine Фур'є-перетворення.

Таблиця Dr скінченне перетворення Фур'є, таблиця та скінченне косинусне
перетворення Фур'є.

Таблиця D: Трансформація Лапласа.

\section{\texorpdfstring{Прийняті позначення \(\delta\left(x\right)\)
---
дельта-функція,}{Прийняті позначення \textbackslash delta\textbackslash left(x\textbackslash right) --- дельта-функція,}}\label{ux43fux440ux438ux439ux43dux44fux442ux456-ux43fux43eux437ux43dux430ux447ux435ux43dux43dux44f-deltaleftxright-ux434ux435ux43bux44cux442ux430-ux444ux443ux43dux43aux446ux456ux44f}

{[} H(x-a) = {]}

{[} \textbackslash begin\{cases\} 0, x \textless{} a, \textbackslash{}
1, x \textbackslash geqslant a, \textbackslash end\{cases\} {]}

Функція Хевісайду,

{[} H(x-a) = {]}

{[} \textbackslash begin\{cases\} 0, \& x \textless{} a,
\textbackslash{} 1, \& x \textbackslash geqslant a,
\textbackslash end\{cases\} {]}

Функція Хевісайду,

{[} H(a-x) = \textbackslash begin\{cases\} 1, \& x
\textbackslash leqslant a, \textbackslash{} 0, \& x \textgreater{} a,
\textbackslash end\{cases\} {]}

дзеркальною функцією Хевісайду.

ТАБЛИЦЯ A. Експоненціальне перетворення Фур'є

{[} f(x) = \textbackslash mathcal\{F\}\^{}\{-1\}{[}F{]} =
\textbackslash frac\{1\}\{\textbackslash sqrt\{2\textbackslash pi\}\}
\textbackslash int\_\{-\textbackslash infty\}\^{}\{+\textbackslash infty\}
F(\textbackslash omega) e\^{}\{-i\textbackslash omega x\}
d\textbackslash omega \textbackslash qquad F(\textbackslash omega) =
\textbackslash mathcal\{F\}{[}f{]} =
\textbackslash frac\{1\}\{\textbackslash sqrt\{2\textbackslash pi\}\}
\textbackslash int\_\{-\textbackslash infty\}\^{}\{+\textbackslash infty\}
f(x) e\^{}\{-i\textbackslash omega x\} dx {]}

{\def\LTcaptype{none} % do not increment counter
\begin{longtable}[]{@{}
  >{\raggedright\arraybackslash}p{(\linewidth - 2\tabcolsep) * \real{0.3797}}
  >{\raggedright\arraybackslash}p{(\linewidth - 2\tabcolsep) * \real{0.6203}}@{}}
\toprule\noalign{}
\begin{minipage}[b]{\linewidth}\raggedright
1. \(f'(x)\)
\end{minipage} & \begin{minipage}[b]{\linewidth}\raggedright
iωF (ω)
\end{minipage} \\
\midrule\noalign{}
\endhead
\bottomrule\noalign{}
\endlastfoot
2. \(f''(x)\) & \(-\omega^2 F(\omega)\) \\
3. fn (x), (n-та похідна) & \((i\omega)^n F(\omega)\) \\
4. \(f(ax), a > 0\) & \(\frac{1}{a} F\left(\frac{\omega}{a}\right)\) \\
5. \(f(x-a)\) & \(e^{-ia\omega}F(\omega)\) \\
6. \(e^{-a^2x^2}\) & \(\frac{1}{a\sqrt{2}}e^{-\omega^2/4a^2}\) \\
7. \(e^{-a_1 x}\) & \(\sqrt{\frac{2}{\pi}} \frac{a}{a^2 + \omega^2}\) \\
\end{longtable}
}

\begin{enumerate}
\def\labelenumi{\arabic{enumi}.}
\setcounter{enumi}{19}
\item
  гріх (ax)
\item
\end{enumerate}

{[} \textbackslash begin\{cases\} 1, \textbar x\textbar{} \textless{} a
\textbackslash{} 0, \textbar x\textbar{} \textgreater{} a
\textbackslash end\{cases\} {]}

\begin{enumerate}
\def\labelenumi{\arabic{enumi}.}
\setcounter{enumi}{8}
\tightlist
\item
\end{enumerate}

{[} \textbackslash begin\{cases\} 1, \textbar x\textbar{} \textless{} 1
\textbackslash{} 0, \textbar x\textbar{} \textgreater{} 1
\textbackslash end\{cases\} {]}

\begin{enumerate}
\def\labelenumi{\arabic{enumi}.}
\setcounter{enumi}{9}
\tightlist
\item
\end{enumerate}

{[} \textbackslash delta(x-a) {]}

\begin{enumerate}
\def\labelenumi{\arabic{enumi}.}
\setcounter{enumi}{10}
\tightlist
\item
\end{enumerate}

\(f(x)\) * \(g(x)\)

\begin{enumerate}
\def\labelenumi{\arabic{enumi}.}
\setcounter{enumi}{11}
\tightlist
\item
\end{enumerate}

{[} (1+x\textsuperscript{2)}\{-1\} {]}

\begin{enumerate}
\def\labelenumi{\arabic{enumi}.}
\setcounter{enumi}{12}
\tightlist
\item
\end{enumerate}

{[} xe\textsuperscript{\{-a+x}4\}, a \textgreater{} 0 {]}

\begin{enumerate}
\def\labelenumi{\arabic{enumi}.}
\setcounter{enumi}{13}
\tightlist
\item
\end{enumerate}

\(H(x+a)\) - \(H(x-a)\)

\begin{enumerate}
\def\labelenumi{\arabic{enumi}.}
\setcounter{enumi}{14}
\tightlist
\item
\end{enumerate}

{[} \textbackslash frac\{a\}\{x\textsuperscript{2+a}2\} {]}

\begin{enumerate}
\def\labelenumi{\arabic{enumi}.}
\setcounter{enumi}{15}
\tightlist
\item
\end{enumerate}

{[} \textbackslash frac\{2ax\}\{(x\textsuperscript{2+a}2)\^{}2\} {]}

\begin{enumerate}
\def\labelenumi{\arabic{enumi}.}
\setcounter{enumi}{16}
\tightlist
\item
\end{enumerate}

{[} \textbackslash begin\{cases\} \textbackslash cos(ax),
\textbar x\textbar{} \textgreater{} \textbackslash pi/2a
\textbackslash{} 0, \textbar x\textbar{} \textgreater{}
\textbackslash pi/2a \textbackslash end\{cases\} {]}

\begin{enumerate}
\def\labelenumi{\arabic{enumi}.}
\setcounter{enumi}{17}
\tightlist
\item
\end{enumerate}

{[} \textbackslash begin\{cases\} 1-\textbar x\textbar,
\textbar x\textbar{} \textless{} 1 \textbackslash{} 0,
\textbar x\textbar{} \textgreater{} 1 \textbackslash end\{cases\} {]}

\begin{enumerate}
\def\labelenumi{\arabic{enumi}.}
\setcounter{enumi}{18}
\tightlist
\item
\end{enumerate}

{[} \textbackslash cos(ax) {]}

\begin{enumerate}
\def\labelenumi{\arabic{enumi}.}
\setcounter{enumi}{9}
\tightlist
\item
\end{enumerate}

{[} \textbackslash frac\{2\}\{\textbackslash pi\}
\textbackslash frac\{\textbackslash sin
a\textbackslash omega\}\{\textbackslash omega\} {]}

\begin{enumerate}
\def\labelenumi{\arabic{enumi}.}
\setcounter{enumi}{9}
\tightlist
\item
\end{enumerate}

{[} \textbackslash frac\{2\}\{\textbackslash pi\}
\textbackslash frac\{\textbackslash sin
a\textbackslash omega\}\{(\textbackslash omega\textsuperscript{2+a}2)\^{}2\}
{]}

\begin{enumerate}
\def\labelenumi{\arabic{enumi}.}
\setcounter{enumi}{10}
\tightlist
\item
\end{enumerate}

{[} \textbackslash frac\{\textbackslash pi\}\{2\}
e\^{}\{-a+\textbackslash omega\} {]}

\begin{enumerate}
\def\labelenumi{\arabic{enumi}.}
\setcounter{enumi}{11}
\tightlist
\item
\end{enumerate}

{[} \textbackslash frac\{\textbackslash pi\}\{2\}
e\^{}\{-a+\textbackslash omega\} {]}

\begin{enumerate}
\def\labelenumi{\arabic{enumi}.}
\setcounter{enumi}{12}
\tightlist
\item
\end{enumerate}

{[} \textbackslash frac\{\textbackslash pi\}\{2\}
e\^{}\{-a+\textbackslash omega\} {]}

\begin{enumerate}
\def\labelenumi{\arabic{enumi}.}
\setcounter{enumi}{13}
\tightlist
\item
\end{enumerate}

{[} \textbackslash frac\{\textbackslash pi\}\{2\}
e\^{}\{-a+\textbackslash omega\} {]}

\begin{enumerate}
\def\labelenumi{\arabic{enumi}.}
\setcounter{enumi}{14}
\tightlist
\item
\end{enumerate}

{[} \textbackslash frac\{\textbackslash pi\}\{2\}
e\^{}\{-a+\textbackslash omega\} {]}

\begin{enumerate}
\def\labelenumi{\arabic{enumi}.}
\setcounter{enumi}{15}
\tightlist
\item
\end{enumerate}

{[} \textbackslash frac\{2ax\}\{(x\textsuperscript{2+a}2)\^{}2\} {]}

\begin{enumerate}
\def\labelenumi{\arabic{enumi}.}
\setcounter{enumi}{16}
\tightlist
\item
\end{enumerate}

{[} \textbackslash frac\{\textbackslash cos(ax), \textbar x\textbar{}
\textgreater{} \textbackslash pi/2a\}\{0, \textbar x\textbar{}
\textgreater{} \textbackslash pi/2a\} {]}

\begin{enumerate}
\def\labelenumi{\arabic{enumi}.}
\setcounter{enumi}{17}
\tightlist
\item
\end{enumerate}

{[} \textbackslash frac\{1-\textbar x\textbar, \textbar x\textbar{}
\textless{} 1\}\{0, \textbar x\textbar{} \textgreater{} 1\} {]}

\begin{enumerate}
\def\labelenumi{\arabic{enumi}.}
\setcounter{enumi}{18}
\tightlist
\item
\end{enumerate}

{[} \textbackslash cos(ax) {]}

\begin{enumerate}
\def\labelenumi{\arabic{enumi}.}
\setcounter{enumi}{18}
\tightlist
\item
\end{enumerate}

{[} \textbackslash cos(ax) {]}

\begin{enumerate}
\def\labelenumi{\arabic{enumi}.}
\setcounter{enumi}{9}
\tightlist
\item
\end{enumerate}

{[} \textbackslash frac\{\textbackslash pi\}\{2\}
{[}\textbackslash delta(\textbackslash omega+a) +
\textbackslash delta(\textbackslash omega-a){]} {]}

\section{\texorpdfstring{ТАБЛИЦЯ B. Перетворення Фур'є
\(i \sqrt{\frac{\pi}{2}} [\delta(\omega+a) - \delta(\omega-a)]\)}{ТАБЛИЦЯ B. Перетворення Фур'є i \textbackslash sqrt\{\textbackslash frac\{\textbackslash pi\}\{2\}\} {[}\textbackslash delta(\textbackslash omega+a) - \textbackslash delta(\textbackslash omega-a){]}}}\label{ux442ux430ux431ux43bux438ux446ux44f-b.-ux43fux435ux440ux435ux442ux432ux43eux440ux435ux43dux43dux44f-ux444ux443ux440ux454-i-sqrtfracpi2-deltaomegaa---deltaomega-a}

{[} f(x) = \textbackslash int\_\{0\}\^{}\{\textbackslash infty\}
F(\textbackslash omega) \textbackslash sin(\textbackslash omega x)
d\textbackslash omega {]}

{[} F(\textbackslash omega) =
\textbackslash frac\{?\}\{\textbackslash pi\}
\textbackslash int\_\{0\}\^{}\{\textbackslash infty\} f(x)
\textbackslash sin(\textbackslash omega x) dx {]}

{[} 0 \textless{} x \textless{} \textbackslash infty {]}

{[} 0 \textless{} \textbackslash omega \textless{} x {]}

\begin{enumerate}
\def\labelenumi{\arabic{enumi}.}
\tightlist
\item
\end{enumerate}

f'\,'(0)

{[} -\textbackslash omega\^{}2 F(\textbackslash omega) +
\textbackslash frac\{2\}\{\textbackslash pi\} \textbackslash omega f(0)
{]}

\begin{enumerate}
\def\labelenumi{\arabic{enumi}.}
\setcounter{enumi}{1}
\tightlist
\item
  \(f(\alpha x)\)
\end{enumerate}

{[} \textbackslash frac\{1\}\{\textbackslash alpha\}
F\textbackslash left(\textbackslash frac\{\textbackslash omega\}\{\textbackslash alpha\}\textbackslash right)
{]}

\begin{enumerate}
\def\labelenumi{\arabic{enumi}.}
\setcounter{enumi}{2}
\tightlist
\item
  \(e^{-\alpha x}\)
\end{enumerate}

{[} \textbackslash frac\{2\textbackslash omega\}\{\textbackslash pi
(\textbackslash alpha\^{}2 + \textbackslash omega\^{}2)\} {]}

\begin{enumerate}
\def\labelenumi{\arabic{enumi}.}
\setcounter{enumi}{3}
\tightlist
\item
\end{enumerate}

{[} x\^{}\{-1/2\} {]}

{[} {[}2/\textbackslash pi\textbackslash omega{]}\^{}\{1/2\} {]}

\begin{enumerate}
\def\labelenumi{\arabic{enumi}.}
\setcounter{enumi}{4}
\tightlist
\item
  \(H(a-x)\)
\end{enumerate}

{[} \textbackslash frac\{2\}\{\textbackslash pi\textbackslash omega\}
{[}1-\textbackslash cos(\textbackslash omega a){]} {]}

\begin{verbatim}
6. x^{-1}
\end{verbatim}

1

\begin{verbatim}
7. \frac{x}{x^2+a^2}
\end{verbatim}

{[} e\^{}\{-1\}a + \textbackslash omega {]}

\begin{verbatim}
8. \frac{x}{x^4+4}
\end{verbatim}

{[}
\textbackslash frac\{1\}\{2\}e\^{}\{-\textbackslash alpha\}\textbackslash sin\textbackslash omega
{]}

\begin{verbatim}
9. \arctan \frac{a}{x}
\end{verbatim}

{[}
\textbackslash frac\{1-e\^{}\{-a\textbackslash omega\}\}\{\textbackslash omega\}
{]}

\begin{verbatim}
10. -x^2/(x)
\end{verbatim}

{[}
\textbackslash frac\{2\}\{\textbackslash pi\}F'\,'(\textbackslash omega)
{]}

\section{ТАБЛИЦЯ C. Перетворення
Фур'є}\label{ux442ux430ux431ux43bux438ux446ux44f-c.-ux43fux435ux440ux435ux442ux432ux43eux440ux435ux43dux43dux44f-ux444ux443ux440ux454}

{[} f(x) = \textbackslash int\_\{0\}\^{}\{\textbackslash infty\}
F(\textbackslash omega) \textbackslash cos(\textbackslash omega x)
d\textbackslash omega \textbackslash qquad \textbackslash qquad F(x) =
\textbackslash frac\{2\}\{\textbackslash pi\}
\textbackslash int\_\{0\}\^{}\{\textbackslash infty\} f(x)
\textbackslash cos(\textbackslash omega x) dx {]}

{[} 0 \textless{} x \textless{} \textbackslash infty
\textbackslash qquad \textbackslash qquad 0 \textless{}
\textbackslash omega \textless{} \textbackslash infty {]}

\begin{enumerate}
\def\labelenumi{\arabic{enumi}.}
\tightlist
\item
\end{enumerate}

j'\,'(x)

\section{ТАБЛИЦЯ D. Скінченне перетворення синусоїда
Фур'є}\label{ux442ux430ux431ux43bux438ux446ux44f-d.-ux441ux43aux456ux43dux447ux435ux43dux43dux435-ux43fux435ux440ux435ux442ux432ux43eux440ux435ux43dux43dux44f-ux441ux438ux43dux443ux441ux43eux457ux434ux430-ux444ux443ux440ux454}

У скінченному перетворенні синуса Фур'є функція \(f(x)\), визначена в
інтервалі \(0 \le x \le \pi\) , перетворюється у послідовність \(S_n\) ,
n = 1, 2 \ldots{} За формулою

{[} S\_n = \textbackslash frac\{2\}\{\textbackslash pi\}
\textbackslash int\_0\^{}\{\textbackslash pi\} f(x)
\textbackslash sin(nx) dx. {]}

Коли йдеться про таблиці, зручно розглядати змінну x як належну до
інтервалу {[}0, l{]}. Будь-який інший чисельний інтервал {[}a, b{]}
можна перетворити на \([0, \pi]\) за допомогою формули

{[} y = \textbackslash frac\{\textbackslash pi\}\{b - a\} (x - a). {]}

У скінченному синусоїдальному перетворенні немає нічого загадкового:
елементи послідовності \(\mathcal{S}_n\) збігаються з коефіцієнтами в
\(\sin(nx)\) у розкладі функції \(f(x)\) у ряд Фур'є, тобто

{[} f(x) = \textbackslash sum\_\{n=1\}\^{}\{\textbackslash infty\} S\_n
\textbackslash sin(nx). {]}

{[} f(x) = \textbackslash sum\_\{n=1\}\^{}\{\textbackslash infty\} S\_n
\textbackslash sin(nx) {]}

{[} S\_n = \textbackslash frac\{2\}\{\textbackslash pi\}
\textbackslash int\_\{0\}\^{}\{\textbackslash pi\} f(x)
\textbackslash sin(nx) dx {]}

{[} 0 \textbackslash le x \textbackslash le \textbackslash pi {]}

{[} n = 1, 2, \ldots{} {]}

\begin{enumerate}
\def\labelenumi{\arabic{enumi}.}
\tightlist
\item
\end{enumerate}

f'\,'(x)

\begin{enumerate}
\def\labelenumi{\arabic{enumi}.}
\setcounter{enumi}{8}
\tightlist
\item
\end{enumerate}

{[} \textbackslash frac\{\textbackslash pi\}\{2\} e\^{}\{ax\} {]}

\(\frac{n}{n^2 + a^2} [1 - (-1)^n e^{ax}]\) 10.
\(\frac{\sin a (\pi - x)}{\sin a\pi}\) \(\frac{2n}{\pi (n^2 + a^2)}\)
ТАБЛИЦЯ E. Скінченне перетворення Косніуса Фур'є

У скінченному косинусному перетворенні Фур'є функція \(f(x)\), визначена
на інтервалі \(0 \le x \le \pi\) , перетворюється у послідовність
\(C_n\) , \(n=0,\ 1,\ 2,\ \ldots\) , за формулою

{[} C\_n = \textbackslash frac\{2\}\{\textbackslash pi\}
\textbackslash int\_0\^{}\{\textbackslash pi\} f(x)
\textbackslash cos(nx) dx. {]}

Коли йдеться про таблиці, зручно розглядати змінну x як таку, що
належить до інтервалу \([0,\pi]\) . Будь-який інший числовий інтервал
{[}a,b{]} можна перетворити на \([0,\pi]\) за допомогою формули

{[} y = \textbackslash frac\{\textbackslash pi\}\{b-a\}(x-a). {]}

У скінченному косинусному перетворенні немає нічого загадкового:
елементи послідовності \(C_n\) збігаються з коефіцієнтами, коли
\(\cos{(nx)}\) у косинусному розкладі функції \(f(x)\) у ряд Фур'є,
тобто

{[} f(x) = \textbackslash frac\{C\_0\}\{2\} +
\textbackslash sum\_\{n=1\}\^{}\{\textbackslash infty\} C\_n
\textbackslash cos(nx). {]}

{[} f(x) = \textbackslash frac\{C\_0\}\{2\} +
\textbackslash sum\_\{n=1\}\^{}\{\textbackslash infty\} C\_n
\textbackslash cos(nx) \textbackslash qquad C\_n =
\textbackslash frac\{2\}\{\textbackslash pi\}
\textbackslash int\_\{0\}\^{}\{\textbackslash pi\} f(x)
\textbackslash cos nx , dx {]}

{[} 0 \textbackslash le x \textbackslash le \textbackslash pi
\textbackslash qquad n = 0, 1, 2, \textbackslash dots {]}

\begin{enumerate}
\def\labelenumi{\arabic{enumi}.}
\tightlist
\item
\end{enumerate}

f'\,'(x)

{[} -n\^{}\{2\}C\_\{n\} - \textbackslash frac\{2\}\{\textbackslash pi\}
\{f'(0) - (-1)\^{}\{n\} f'(n)\} {]}

\begin{verbatim}
2. \frac{a_{0}}{2} + \sum_{n=1}^{\infty} a_{n} \cos(nx) a_{n} 3. f(n-x) (-1)^{n} \frac{2}{\pi} C_{n} 4. 1
\end{verbatim}

{[} \textbackslash begin\{cases\} 2, \&n=0 \textbackslash{} 0, \&n=1, 2,
\textbackslash dots \textbackslash end\{cases\} {]}

\begin{verbatim}
5. \cos(mx), m=1, 2, \dots
\end{verbatim}

{[} \textbackslash begin\{cases\} 1, \& n=m \textbackslash{} 0, \&
n\textbackslash neq m \textbackslash end\{cases\} {]}

\begin{enumerate}
\def\labelenumi{\arabic{enumi}.}
\setcounter{enumi}{5}
\tightlist
\item
\end{enumerate}

x

\begin{verbatim}
\begin{cases}
\end{verbatim}

ТАБЛИЦЯ F. Перетворення Лапласа

\begin{verbatim}
| F(s) = \mathscr{L}[f(t)] |
\end{verbatim}

\textbar----------------------------------\textbar{}

1/с, с \textgreater{} 0 \textbar{}

\begin{verbatim}
| \frac{1}{s-a} , s>a |

| \frac{a}{s^2 + a^2}, s > 0 |

| \frac{s}{s^2 + a^2}, s > 0 |

| \frac{a}{s^2 - a^2}, s > a |

| \frac{s}{s^2-a^2}, s> a |

| \frac{b}{(s-a)^2+b^2}, s>a |

| \frac{s-a}{(s-a)^2+b^2}, s>a |

| \frac{n!}{8^{n+1}} <>0 |

| \frac{n!}{(s-a)^{n+1}}, s > a |

| \frac{e^{-as}}{s} , s>0 |
\end{verbatim}

e\^{}\{-as\}F(s) . \textbar{}

F(s-a) \textbar{}

F(s) G(s) \textbar{}

\_ \textbar{}

\begin{verbatim}
| 15. f^{(n)}(t) , (n-похідна) | s^{n}F(s) - s^{n-1}f(0) - \dots - f^{(n-1)}(0) |
\end{verbatim}

\textbar------------------------------------------------------------------------\textbar--------------------------------------------------\textbar{}

\begin{verbatim}
| 16. ж. (at) | \frac{1}{a}F\left(\frac{s}{a}\right),  a>0 |

| 17. \int_{0}^{t} f(\tau) d\tau | \frac{1}{s} F(s) |

| 18. ERF (T/2A) | \frac{1}{s} e^{a^2s^2} erfe (AS) |

| 19. Ерік (A/2 <b>V</b> <del>T</del> ) | \frac{1}{s}e^{-a} |

| 20. J_{\theta}(at) | (s^2+a^2)^{1/2} |

| 21. \delta(t-a) | e^{-sa} |

| 22. \frac{1}{\sqrt{\pi t}} \exp\left(-\frac{a^2}{4t}\right) | \frac{e^{-a/s}}{\sqrt{s}}, a \ge 0 |

| 23. \frac{1}{\sqrt{nt}} - ae^{a^2t} \operatorname{eric} (a \sqrt{t}) | \frac{1}{\sqrt{s}+a} |
\end{verbatim}

\section{Додаток 2}\label{ux434ux43eux434ux430ux442ux43eux43a-2}

\section{ПРЕДСТАВЛЕННЯ ЛАПЛАСІАНА В РІЗНИХ СИСТЕМАХ
КООРДИНАТ}\label{ux43fux440ux435ux434ux441ux442ux430ux432ux43bux435ux43dux43dux44f-ux43bux430ux43fux43bux430ux441ux456ux430ux43dux430-ux432-ux440ux456ux437ux43dux438ux445-ux441ux438ux441ux442ux435ux43cux430ux445-ux43aux43eux43eux440ux434ux438ux43dux430ux442}

{[} \textbackslash Delta u = u\_\{xx\} + u\_\{yy\} {]}

{[} \textbackslash Delta u = u\_\{rr\} +
\textbackslash frac\{1\}\{r\}u\_r +
\textbackslash frac\{1\}\{r\^{}2\}u\_\{\textbackslash theta\textbackslash theta\}
{]}

{[} \textbackslash Delta u = u\_\{xx\} + u\_\{yy\} + u\_\{zz\} {]}

{[} \textbackslash Delta u = u\_\{rr\} +
\textbackslash frac\{1\}\{r\}u\_r +
\textbackslash frac\{1\}\{r\^{}2\}u\_\{\textbackslash theta\textbackslash theta\}
+ u\_\{zz\} {]}

{[} \textbackslash Delta u = u\_\{rr\} +
\textbackslash frac\{2\}\{r\}u\_r +
\textbackslash frac\{1\}\{r\^{}2\}u\_\{\textbackslash theta\textbackslash theta\}
+
\textbackslash frac\{\textbackslash cot\textbackslash theta\}\{r\^{}2\}
{]}

{[} \textbackslash Delta u = u\_\{rr\} +
\textbackslash frac\{2\}\{r\}u\_r +
\textbackslash frac\{1\}\{r\^{}2\}u\_\{\textbackslash theta\textbackslash theta\}
+
\textbackslash frac\{\textbackslash cot\textbackslash theta\}\{r\^{}2\}u\_0
+
\textbackslash frac\{1\}\{r\^{}2\textbackslash sin\textbackslash theta\}u\_\{\textbackslash phi\textbackslash phi\}
{]}

У тривимірній сферичній системі координат

\pandocbounded{\includegraphics[keepaspectratio]{images/_page_370_Figure_1.jpeg}}

\pandocbounded{\includegraphics[keepaspectratio]{images/_page_370_Picture_2.jpeg}}

\pandocbounded{\includegraphics[keepaspectratio]{images/_page_370_Figure_3.jpeg}}

\section{Декартова система
координат}\label{ux434ux435ux43aux430ux440ux442ux43eux432ux430-ux441ux438ux441ux442ux435ux43cux430-ux43aux43eux43eux440ux434ux438ux43dux430ux442-1}

{[} -\textbackslash infty \textless{} x \textless{}
+\textbackslash infty {]}

{[} -\textbackslash infty \textless{} y \textless{}
+\textbackslash infty {]}

{[} -\textbackslash infty \textless{} z \textless{}
+\textbackslash infty {]}

\section{Циліндрична система
координат}\label{ux446ux438ux43bux456ux43dux434ux440ux438ux447ux43dux430-ux441ux438ux441ux442ux435ux43cux430-ux43aux43eux43eux440ux434ux438ux43dux430ux442-1}

\begin{verbatim}
\begin{aligned}
\end{verbatim}

{[} -\textbackslash infty \& \textless{} x \textless{}
+\textbackslash infty \& r \textbackslash geqslant 0, ~0
\textbackslash leqslant \textbackslash theta \textbackslash leqslant
2\textbackslash pi, \textbackslash{} {]}

{[} -\textbackslash infty \& \textless{} y \textless{}
+\textbackslash infty \& -\textbackslash alpha \textbackslash leqslant z
\textbackslash leqslant \textbackslash alpha, \textbackslash{} {]}

{[} -\textbackslash infty \& \textless{} z \textless{}
+\textbackslash infty \& x = r \textbackslash cos \textbackslash theta,
\textbackslash{} {]}

{[} y = r \textbackslash sin \textbackslash theta, \textbackslash{} {]}

{[} z = z. {]}

\begin{verbatim}
\end{aligned}
\end{verbatim}

\section{Сферична система
координат}\label{ux441ux444ux435ux440ux438ux447ux43dux430-ux441ux438ux441ux442ux435ux43cux430-ux43aux43eux43eux440ux434ux438ux43dux430ux442-1}

{[} r \textbackslash geqslant 0, 0 \textbackslash leqslant
\textbackslash theta \textbackslash leqslant 2\textbackslash pi , 0
\textbackslash leqslant \textbackslash varphi \textbackslash leqslant
\textbackslash pi , x = r \textbackslash sin \textbackslash varphi
\textbackslash cos \textbackslash theta , y = r \textbackslash sin
\textbackslash varphi \textbackslash sin \textbackslash theta , z = r
\textbackslash cos \textbackslash varphi . {]}

\section{Еліптичні типові
рівняння}\label{ux435ux43bux456ux43fux442ux438ux447ux43dux456-ux442ux438ux43fux43eux432ux456-ux440ux456ux432ux43dux44fux43dux43dux44f}

{[} \textbackslash Delta u = 0 {]}

{[} \textbackslash Delta u + \textbackslash lambda\^{}2 u = 0 {]}

{[} \textbackslash Delta u = k {]}

{[} \textbackslash Delta u + k(E - V)u = 0 {]}

Рівняння Лапласа Рівняння Гельмгольца Рівняння Пуассона-Шредінгера

\section{Рівняння гіперболічного
типу}\label{ux440ux456ux432ux43dux44fux43dux43dux44f-ux433ux456ux43fux435ux440ux431ux43eux43bux456ux447ux43dux43eux433ux43e-ux442ux438ux43fux443-1}

{[} \textbackslash begin\{aligned\} u\_\{tt\} = c\^{}2 u\_\{xx\}
\textbackslash{} u\_\{tt\} = c\^{}2 u\_\{xx\} - h u\_t \textbackslash{}
u\_\{tt\} = c\^{}2 u\_\{xx\} - h u\_t - k u \textbackslash{} u\_\{tt\} =
c\^{}2 u\_\{xx\} + f(x, t) \textbackslash end\{aligned\} {]}

{[} u\_\{tt\} = c\^{}2 \textbackslash Delta u {]}

{[} u\_\{tt\} = c\^{}2 \textbackslash Delta u - hu\_t {]}

Рівняння коливання струни Рівняння коливання з тертям

Телеграфне рівняння

Рівняння коливань під впливом зовнішньої сили

Багатовимірне хвильове рівняння з тертям

\section{Параболічні типові
рівняння}\label{ux43fux430ux440ux430ux431ux43eux43bux456ux447ux43dux456-ux442ux438ux43fux43eux432ux456-ux440ux456ux432ux43dux44fux43dux43dux44f}

{[} u\_t = \textbackslash alpha\^{}2 u\_\{xx\} {]}

{[} u\_t = \textbackslash alpha\^{}2 u\_\{xx\} - hu\_x {]}

{[} u\_t = \textbackslash alpha\^{}2 u\_\{xx\} - ku {]}

{[} u\_t = \textbackslash alpha\^{}2 u\_\{xx\} + f(x, t) {]}

Уравнение одномерной диффузии Уравнение конвективной диффузии

Уравнение теплопроводности с поглощением Уравнение теплопроводности с
источником

\section{\texorpdfstring{\textbf{КРОССВОРД}}{КРОССВОРД}}\label{ux43aux440ux43eux441ux441ux432ux43eux440ux434-1}

\pandocbounded{\includegraphics[keepaspectratio]{images/_page_371_Figure_3.jpeg}}

\begin{verbatim}
Горизонтальна 1. Характеристики векторного поля. 3. Накладання. 8. Довжина. 9. Оператор Лаплас. 12. Поширення обурення. 13. Довжина вектора. 14. Робота функцій прокладки. 15. Множина \int \Delta u = 0 множина точок, на яких будується схема різниці. 16. |u|_{\Gamma} = \varphi Завдання ______ 22. Рух рідини описується рівняннями _______ Стокс. 24 Функція f(x) = f(-x), якщо f(x) = f(-x). 25. Французький математик, автор методу ділення змінних. 26. Тонкий натягнутий дріт. 27. Будь-яке нетривіальне розв'язання задачі Штурма-Ліувіля називається ___ 28. Французький математик, на честь якого названа задача з початковою умовою. 29. Основний принцип геометричної оптики називається принципом ______ 30. Колишня назва студентської книги записів. 32. Математик, який запропонував прямий метод розв'язання варіаційних задач. 33. Функцію джерела також називають функцією ____ 34. Гармоніка. 36. Площа літака D = \{r \colon R_1 \leqslant r \leqslant R_2\} . 38. Англійський математик і механік. 40. Головна характеристика періодичної функції. 42. Хвилі в середовищі поширюються відповідно до принципу ___ 43. Спектральний параметр. 44. Узагальнення поняття функції. 45. Множина частот під час розкладу функції.
\end{verbatim}

Вертикальний

поля.

\begin{enumerate}
\def\labelenumi{\arabic{enumi}.}
\setcounter{enumi}{1}
\item
  Термофізичні характеристики матерії. 3. Радянський математик,
  академік, автор найвідомішого підручника з рівнянь математичної
  фізики. 4. Задача з граничними умовами третього типу називається
  задачею \_\_\_\_\_\_\_\_\_\_\_

  K(x,s) у інтегральному перетворенні F(s) =
  \textbackslash int\_\{a\}\^{}\{b\} K(x, s) f(s) ds 7. Значення, яке не
  змінюється під час трансформацій. 10. Характеристики векторного поля.
  11. Французький математик. 13. Стю-

  Функція піни H(t) = \textbackslash begin\{cases\} 0, \& t \textless{}
  0 \textbackslash{} 1, \& t \textbackslash ge 0
  \textbackslash end\{cases\} \textbackslash end\{cases\}
\end{enumerate}

називається функцією 19. Відображення складної площини,

Той, що зберігає кути між кривими, називається \_\_\_\_\_. 20. Лінія у
фазовому просторі, вздовж якої розв'язок не змінюється. 21. Дисплей 23.
Зміна функції. 29. Функція з функції. 31. Комп'ютер. 33. Наука про
просторові форми. 35. Французький математик, на честь якого названо один
із принципів розв'язання неоднорідних диференціальних рівнянь. 37.
Німецький астроном, на честь якого названо клас циліндричних функцій.
39. Великий англійський математик, фізик, механік і астроном. 41.
Інтегральна характеристика вектора

\section{Джерела ()}\label{ux434ux436ux435ux440ux435ux43bux430}

\section{Головне}\label{ux433ux43eux43bux43eux432ux43dux435}

\begin{itemize}
\tightlist
\item
  Бідадзе А. В. Рівняння математичної фізики.--- Москва: Наука, 1976.
\item
  Володимиров В. С. Рівняння математичної фізики.--- Москва: Наука,
  1971. Кошляков Н. С., Глінер Е. Б., Смірнов М. М. Рівняння в зокрема
  математичній фізиці.--- Москва: Вища школа, 1970.
\item
  \begin{enumerate}
  \def\labelenumi{\arabic{enumi}.}
  \setcounter{enumi}{3}
  \tightlist
  \item
    Самарський, А. А. Вступ до чисельних методів. Москва, Наукове
    видавництво, 1982.
  \end{enumerate}
\item
  \begin{enumerate}
  \def\labelenumi{\arabic{enumi}.}
  \setcounter{enumi}{4}
  \tightlist
  \item
    Смирнов В. І. Курс, виший математіки, т. 2.--- М.: Наука, 1967.
  \end{enumerate}
\item
  \begin{enumerate}
  \def\labelenumi{\arabic{enumi}.}
  \setcounter{enumi}{5}
  \tightlist
  \item
    Тихонов А. Н., Самарський А. А. Рівняння математичної фізики.
    Москва, Наукове видання, 1972.
  \end{enumerate}
\end{itemize}

\section{Додатково}\label{ux434ux43eux434ux430ux442ux43aux43eux432ux43e}

\begin{itemize}
\tightlist
\item
  \begin{enumerate}
  \def\labelenumi{\arabic{enumi}.}
  \setcounter{enumi}{6}
  \tightlist
  \item
    Задача Адамара Дж. Коші для лінійних диференціальних рівнянь у
    похідних похідних гіперболічного типу. Москва, Наукове видання,
    1978.
  \end{enumerate}
\item
  Арсенін В. Я. Методи математичної фізики та спеціальні функції.
  Москва; Наука, 1974.
\item
  \begin{enumerate}
  \def\labelenumi{\arabic{enumi}.}
  \setcounter{enumi}{8}
  \tightlist
  \item
    Арфкен Г. Математичні методи у фізиці.--- Москва: Атомиздат, 1970.
  \end{enumerate}
\item
  \begin{enumerate}
  \def\labelenumi{\arabic{enumi}.}
  \setcounter{enumi}{9}
  \tightlist
  \item
    Берс Л., Джон Ф., Шехтер Б. Диференціальні рівняння з частинними
    похідними. Москва, видавництво Мир, 1966.
  \end{enumerate}
\item
  \begin{enumerate}
  \def\labelenumi{\arabic{enumi}.}
  \setcounter{enumi}{10}
  \tightlist
  \item
    Бінадзе, А. В. Крайові задачі для еліптичних рівнянь другого
    порядку.--- Москва: Наука, 1966.
  \end{enumerate}
\item
  \begin{enumerate}
  \def\labelenumi{\arabic{enumi}.}
  \setcounter{enumi}{11}
  \tightlist
  \item
    Бідадзе А. В. Деякі класи диференціальних рівнянь у похідних
    похідних. Москва, Наукове видання, 1981.
  \end{enumerate}
\item
  \begin{enumerate}
  \def\labelenumi{\arabic{enumi}.}
  \setcounter{enumi}{12}
  \tightlist
  \item
    Біцадзе А. В., Калініченко Д. К. Збірка задач з рівнянь математичної
    фізики. Москва, Наука, 1977.
  \end{enumerate}
\item
  \begin{enumerate}
  \def\labelenumi{\arabic{enumi}.}
  \setcounter{enumi}{13}
  \tightlist
  \item
    Будак Б. М., Самарський А. А., Тихонов А. Н. Сборник задач по
    математичній фізиці {[}Збірка задач математичної фізики{]}. Москва,
    Наукове видання, 1972.
  \end{enumerate}
\item
  Вазов В., Форсайт Дж. Методи різниці для розв'язання диференціальних
  рівнянь у похідних похідних.--- Москва: Іл, 1963
\item
  Нові методи розв'язання еліптичних рівнянь {[}Нові методи розв'язання
  еліптичних рівнянь{]}. Москва, Гостехіддат, 1948. Векуа, І. Н.
  Узагальнені аналітичні функції. Москва, Фізматгиз, 1959.
\item
  Володимиров В. С. Узагальнені функції в математичній фізиці, Москва,
  Наука, 1976.
\item
  \begin{enumerate}
  \def\labelenumi{\arabic{enumi}.}
  \setcounter{enumi}{18}
  \tightlist
  \item
    Володимиров В. С. та ін. Збірка задач з рівнянь математичної фізики.
    --- Москва: Наука, 1974.
  \end{enumerate}
\item
  Годунов С. К. Рівняння математичної фізики.--- Москва: Наука, 1971.
  Годунов С. К., Рябенський В. С. Схеми різниці.--- Москва: Наука, 1977.
\item
  Гордін: задача Л. Коші для гіперболічних рівнянь.--- Москва: IL, 1961.
  Джеффріс Г., Спірлс Б. Методи математичної фізики. т. 1-3, --- Москва:
  Мир,
\item
  \begin{enumerate}
  \def\labelenumi{\arabic{enumi}.}
  \setcounter{enumi}{1969}
  \tightlist
  \item
  \end{enumerate}
\item
  \begin{enumerate}
  \def\labelenumi{\arabic{enumi}.}
  \setcounter{enumi}{23}
  \tightlist
  \item
    Зоммерфельд А. Рівняння в часткових похідних математичної фізики.---
    Москва: Іл, 1950. 25. Кализкін, Н. Н. Численільна методика.--- М.:
    Наука, 1978.
  \end{enumerate}
\item
  25 Курант Р., Гілберт Д. Методи математичної фізики. т. 1-2, - Москва:
  Гостехиздат, 1951.
\item
  \begin{enumerate}
  \def\labelenumi{\arabic{enumi}.}
  \setcounter{enumi}{26}
  \tightlist
  \item
    Курант Р. Рівняння з частковими виробництвами. Москва, Видавництво
    Мир, 1964.
  \end{enumerate}
\item
  \begin{enumerate}
  \def\labelenumi{\arabic{enumi}.}
  \setcounter{enumi}{27}
  \tightlist
  \item
    Лаврентьєв М. А., Шабат Б. В. Методи теорії функцій комплексної
    змінної. - Москва: Наука, 1965.
  \end{enumerate}
\item
  \begin{enumerate}
  \def\labelenumi{\arabic{enumi}.}
  \setcounter{enumi}{28}
  \tightlist
  \item
    Ладиженська О. А. Прикордонні задачі математичної фізики. Москва:
    Наука,
  \end{enumerate}
\item
  Ладиженська О. А. Змішана задача для гіперболічних рівнянь.--- Москва:
  Фізматгиз, 1953.
\end{itemize}

Список був складений редактором перекладу.

\begin{itemize}
\tightlist
\item
  Ладиженська О. А., Солоняков В. А., Уралцева Н. Н. Лінійні та
  кварілінні рівняння параболічного типу.--- Москва: Наука, 1967.
\item
  \begin{enumerate}
  \def\labelenumi{\arabic{enumi}.}
  \setcounter{enumi}{31}
  \tightlist
  \item
    Ладиженська О. А., Уралцева Н. Н. Лінійні та квазілінійні рівняння
    еліптичного типу. - Москва: Наука, 1964, 1973.
  \end{enumerate}
\item
  Лебедєв Н. Н., Скальська М. П., Уфляйнд Ю. С. Сборник задах по
  математичній фізиці.--- Москва: Гостехиддат, 1965.
\item
  \begin{enumerate}
  \def\labelenumi{\arabic{enumi}.}
  \setcounter{enumi}{33}
  \tightlist
  \item
    Лайонс, Дж.--- Л. Деякі методи розв'язання задач з нелінійними
    краями.--- М.З. Мір. 1972.
  \end{enumerate}
\item
  \begin{enumerate}
  \def\labelenumi{\arabic{enumi}.}
  \setcounter{enumi}{34}
  \tightlist
  \item
    Лайонс Дж.---Л., Maggenes E. Neobreed Grancinum Problems та їх
    застосування. Москва, Видавництво Мир, 1971. 36. Марчук Г. І. Методи
    обчислювальної математики. Москва, Наукове видання, 1980.
  \end{enumerate}
\item
  Мізохата С. Теорія рівнянь з приватними виробниками. --- Москва: Мир,
  1977.
\item
  \begin{enumerate}
  \def\labelenumi{\arabic{enumi}.}
  \setcounter{enumi}{37}
  \tightlist
  \item
    Міранда К. Рівняння з частковими похідними еліптичного типу.---
    Москва: Іл, 1957.
  \end{enumerate}
\item
  Михайлов, В. П. Диференціальні рівняння в похідних з похідними. ---
  Москва: Наука, 1976.
\item
  Міклін С. Г. Курс математичної фізики.--- Москва: Наука, 1968. Міхлін,
  С. Г. Лінійні диференціальні рівняння. --- Москва; Вишня школа, 1977.
\item
  Морс П. М., Фецибах X, Методи теоретичної фізики. т. 1-2.--- Москва:
  IL, 1958.
\item
  \begin{enumerate}
  \def\labelenumi{\arabic{enumi}.}
  \setcounter{enumi}{42}
  \tightlist
  \item
    Меттьюз Дж., Вокер Р. Математичні методи фізики.--- Москва:
    Атомиздат, 1972.
  \end{enumerate}
\item
  Лекції Нагумо М. з сучасної теорії рівнянь у приватних виробництвах
  Москва, видавництво Мвр, 1967. 45. Обен, Дж.-П. Приблизне розв'язання
  еліптичних граничних задач.--- Москва:
\item
  Мір, 1977. 46. Петровський І. Г. Лекції з рівнянь із частковими
  похідними. --- Москва: Фізматгиз, 1961.
\item
  \begin{enumerate}
  \def\labelenumi{\arabic{enumi}.}
  \setcounter{enumi}{46}
  \tightlist
  \item
    Положий Г. Н. Еквалізація математичної фізики.--- Москва: Вища
    школа.
  \end{enumerate}
\item
  \begin{enumerate}
  \def\labelenumi{\arabic{enumi}.}
  \setcounter{enumi}{47}
  \tightlist
  \item
    Рід М., Саймон Б. Методи сучасної математичної фізики.--- Москва:
    Мир,
  \end{enumerate}
\item
  т. 1-1977, т. 2-1978, т. 3, 4-1982. Ріхамаєр Р. Принципи сучасної
  математичної фізики.--- Москва: Мир,
\item
  \begin{enumerate}
  \def\labelenumi{\arabic{enumi}.}
  \setcounter{enumi}{49}
  \tightlist
  \item
    Ріхтмаєр Р., Мортон К. Методи різниці для розв'язання задач з
    граничними значеннями.--- Москва:
  \end{enumerate}
\item
  Мір, 1972,
\item
  Рождественський Б. Л., Япенко, Н. Н. Системи квазіполіцейських рівнянь
  та їх застосування в газовій динаміці.--- Москва: Наука, 1978.
\item
  52, Самарський А. А. Теорія різницьких схем. Москва, Наука, 1977.
\item
  \begin{enumerate}
  \def\labelenumi{\arabic{enumi}.}
  \setcounter{enumi}{52}
  \tightlist
  \item
    Самарський А. А., Андрєєв В. Б. Різниці для еліптичних рівнянь---
    Москва, Наука, 1976.
  \end{enumerate}
\item
  \begin{enumerate}
  \def\labelenumi{\arabic{enumi}.}
  \setcounter{enumi}{53}
  \tightlist
  \item
    Смірнов В. І. Курс вищої математики.--- Москва: Наука, 1974.
  \end{enumerate}
\item
  \begin{enumerate}
  \def\labelenumi{\arabic{enumi}.}
  \setcounter{enumi}{54}
  \tightlist
  \item
    Смірнов М. М. Дегенеративні еліптичні та гіперболічні рівняння.--
    Москва: Наука, 1966.
  \end{enumerate}
\item
  56, Смірнов М. М. Диференціальні рівняння у похідних другого порядку,
  -- Мінськ: Видавництво БДУ, 1974.
\item
  \begin{enumerate}
  \def\labelenumi{\arabic{enumi}.}
  \setcounter{enumi}{56}
  \tightlist
  \item
    Смірнов М. М. Проблеми рівнянь математичної фізики. Москва, Наукове
    видавництво, 1965.
  \end{enumerate}
\item
  \begin{enumerate}
  \def\labelenumi{\arabic{enumi}.}
  \setcounter{enumi}{57}
  \tightlist
  \item
    Соболєв С. Л. Пекаки primeneniya funktsional'nogo analiza v
    matematicheskoi fizike {[}Застосування функціонального аналізу в
    математичній фізиці{]}. Ленінград: Видавництво ЛГУ, 1950,
    Новосибірськ: СБ АН РСРСР, 1962.
  \end{enumerate}
\item
  Соболєв, С. Л. Рівняння математичної фізики.--- М · Наука, 1966.
\item
  Траптер К. Дж. Інтегральні перетворення у математичній фізиці.---
  Москва: Гостехиддат, 1957.
\item
  \begin{enumerate}
  \def\labelenumi{\arabic{enumi}.}
  \setcounter{enumi}{60}
  \tightlist
  \item
    Трев Ф. Лекції з теорії диференціальних рівнянь у похідних
    похідних.--- Москва: Мир, 1965.
  \end{enumerate}
\item
  Трікомі Ф. Лекції з диференціальних рівнянь у похідних похідних.
  Москва, IL Publ., 1957. Фрідман А. Рівняння з частковими похідними
  параболічного типу. --- Москва: Мир, 1968.
\item
  Германдер Л. Лінійні диференціальні оператори з частковими
  похідними.--- Москва: Мир. 1965.
\item
  \begin{enumerate}
  \def\labelenumi{\arabic{enumi}.}
  \setcounter{enumi}{64}
  \tightlist
  \item
    Шварц Л. Математичні методи фізичних наук.--- Москва: Мир, 1965.
  \end{enumerate}
\item
  \begin{enumerate}
  \def\labelenumi{\arabic{enumi}.}
  \setcounter{enumi}{65}
  \tightlist
  \item
    Шилов Г. Е. Математичний аналіз. Спеціальний курс. Москва,
    Фізмат-гіз видавництво, 1960.
  \end{enumerate}
\item
  \begin{enumerate}
  \def\labelenumi{\arabic{enumi}.}
  \setcounter{enumi}{66}
  \tightlist
  \item
    Щилов, Г. Е. Математичний аналіз. Другий спеціальний курс.---
    Москва: Наука, 1965.
  \end{enumerate}
\item
  \begin{enumerate}
  \def\labelenumi{\arabic{enumi}.}
  \setcounter{enumi}{67}
  \tightlist
  \item
    Яненко Н. Н. Метод дробових щагів для розв'язання багатовимірних
    задач математичної фізики.--- Новосибірськ: Наука, 1967.
  \end{enumerate}
\end{itemize}

\section{ІНДЕКС
ІМЕНІ}\label{ux456ux43dux434ux435ux43aux441-ux456ux43cux435ux43dux456}

Бернуллі 329, Бессель 106, Буффон 322, Ніколсон 296, Ньютон 296

Гамільтон 335, 336 Гаусс 286 Гельмгольц 223 Грін 10 Гук 147 Гюйгенс 177,
181 Пуассон 180, 235

Д'Аламбер 126--129, Діріхле 85, 240, Дюамель 107, 108, 110--112, Рітц
338, 340, Робін 240, Родюг 268

** Seidel 286

Сімсон 320 Стокс 8, 180

Кіз 305, 306 Крхгоф 124

Кірхгоф 124 Кранк 296 Крістоффель 363 Ферма 335 Фур'є 24, 25, 36, 41,
74, 84, 91

Гакслі 304, Ханкел 105, 106, Хевісайд 82, 115, 116, Ходжкін 304-306

Лагранж 329 Ламм 307 Лаплас 8, 50, 98-107 Лежандр 235 Лельбніц 180
Лібман 285 Ліувіль 51, 58, 66, 148

Шварц 363 Шредінгер 8 Штурм 51, 58, 66, 148

Ейлер 84, 215, 329

\textbf{М}аксвелл 8 Меллін 105, 106

10sr (Молодий) 124

Navier 8 Neumann 240, 286

Jacobi 286

\section{ІНДЕКС
ПРЕДМЕТІВ}\label{ux456ux43dux434ux435ux43aux441-ux43fux440ux435ux434ux43cux435ux442ux456ux432}

Матрична алгебра 215, 216 Аналітичне розв'язання 301 Наближення сплайну
10

Позачасовий час 165

Калькулюс варіацій 329 Методи варіацій 10 Зовнішня похідна 79 Хвилі, що
рухаються, 149

-- квартира 177 -- стоячи 222

сферичний 178

--- ударні 210 циліндричні 177, 182

Гармонік 84, 154 Сферичний 265 градієнів 22 Граничні умови 17, 22, 23,
26

--- лінійний 26

--- однорідні 36, 46 --- гетерогенні 26, 49, 50

Дельта-функція 276 Діагоналізація матриці 216 Задача дифузії 113, 163
Дифузійний потік 114 --- член 113 Дифузія 113

Задача ** 147 -- Діріхле 242, 243, 246

\begin{itemize}
\item
  Діріхле - Нойман 364 - Кірхгоф 124
\item
  Коуу 69, 95, 116, 126, 207
\end{itemize}

--- мінімізація 10 --- Нойман 243---245

Ньютон 52, 121

--- на брахістохроні 329 Штурму --- Ліувіль 51, 58, 59, 66,

Закон збереження енергії 31, 114, 205

--- Теплопередача Ньютона 8 246

--- Фур'є 24, 25 Звукові хвилі 124 Момент згинання 157 Інтеграл Фур'є 87
Інтегральні перетворення 9, 73.76. Рівняння 10

Класифікація UCP 308 Крутильні осциляції 123 Поздовжні 123 Конвективні
хвилі 116 Конвективний транспорт 113 --- Член 113 Конвекція 113
Скінченні різниці 282 Канонічні координати 127, 132 --- Полярні 233,
235-238 Сферичні 233, 235, 236, 239, 266 Циліндричні 233, 235, 236, 239
--- еліптичний 239 Неперервний набір 96 Косинусне перетворення 73 ---
скінченне 185, 186 --- Обернене 75, 185, 186 Коефіцієнт дифузії 119 ---
теплопередача 23 Задача граничного значення 240, 241

Лапласіан 233 Лінійність 11

--- --- Еліптичний тренажер 241

Метод дедукції 319 -- Лібман 285, 286 -- Монте-Карло 316 -- Рітц 338,
340 -- Сімпсон 320 спуск 182 -- Фур'є 36 Режим коливання 154 Модуль Юнга
124, 145

Початковий зсув 231 Початкові умови 18

Оператор Лапласа 233 Ортогональність 41, 84, 90

\begin{itemize}
\tightlist
\item
  Ходжкіна - Гакслі 304
\end{itemize}

Власні вектори 67, 215, 217

Випадкові номери 319

Змішане завдання 19

--- enachenia 54, 215, 216 Відгук 96 функція 10, 65 Конформне
відображення 253, 356 Крістоффель-Шварц 363 --- частота 162 Спектральна
функція 76, 77 --- неперервна 77 --- частота 85 Stock 274 Параметрична
ідентифікація 304 Змінна залежна 8 Стоячі хвилі 148, 149 Суперпозиція
97, 193, 223 --- незалежна 8 Схема підрахунку тікання 289 Щільність
матеріалу 30 --- явна 289 Многочлен Лежандра 265 Кранк--Ніколсон 296,
299, 328 Порядок рівняння 10 Потенціал Pewtop 266 з круговою симетрією
240. сферично-симетрична 240 циліндрично симетрична 270 Температурна
характеристика 69 Теорема Діріхле 85 Правило Лейбніца 180--Стокс 180
Перетворення залежних змінних Теорія збурень 10 Тепловий імпульс 96
інтеграл 9, 73, 76 Теплоємність 30 --- координати 9 --- Лаплас 50, 75,
98---107, 115, 116, теплопровідність 25, 30, 33 126 Точкове джерело 274
--- обернене 75 --- Меллін 105, 106 --- Ханкель 76, 105, 106 ---
обернене 76. Питома теплоємність 34 --- Фур'є 74, 75, 83, 85, 89, 91
Вирівнювання Бесселя 106, 224, 239 --- хвиля 8, 142, 169, 234 --- ---
скінченне 185 --- нелінійне 304 --- пряме 88, 116, 126 --- обернене 88,
74, 75 гармонічні коливання 223 --- Гельмгольц 223, 224, 235 Принцип
Гамільтона 335, 336 --- Гюйгенс 177, 181 --- Дюамель 107, 108, 110, 112
Рух Ньютона 8 дифузія 18, 197 - Суперлосоцитності 191, 196, 250
конвектива 19, 63, 113, 119, 120 ** - Ферма 335** - Ламм 307 - Лаплас 8,
170, 235, 258, 285, 363 - Лежандр 239, 268 Змінне розділення 9, 191
лінійне 11 - гіперболічне 11, 12, 121, 169 Розклад власних функцій 10,
66 - параболічне 11, 12, 14, 169 - еліптичне 11, 12, 169 - Максвелл 8 -
Нав'є - Стокс 8 Згортка 90, 93 неліцейних {[}{]} гетерогенних 11, 19,
65. --- скінченна 102 Гравітаційна сила 122 гомогенна P Пуассон 235. --
--- ранкового перетворення 122 синусоїдального перетворення 75 постпост
285 скінченне 185, 186, 193 зворотне 75, 185, 186 реакції з нелінійною
дифузією Система телеграфних рівнянь 124 --- з прямокутними
коефіцієнтами 60.

170

з частковими похідними 7

Телеграф 8, 123, 125, 215

Зберігання тепла 16

{\def\LTcaptype{none} % do not increment counter
\begin{longtable}[]{@{}llllllll@{}}
\toprule\noalign{}
\_ & tep & Lonre & Доступність води & 8, & 16, & 30, & 169. \\
\midrule\noalign{}
\endhead
\bottomrule\noalign{}
\endlastfoot
& & 234. & & • & • & • & , \\
\end{longtable}
}

\begin{itemize}
\tightlist
\item
  характеристики 202, 208
\end{itemize}

--- Шредінгер, 8 років

--- Ейлер 215, 256, 268

--- Ейлер-Лагранж 329, 330, \textbf{33}2

Стан Діріхле 240

\begin{verbatim}
- Neumann 240 - Robin 240
\end{verbatim}

Сингулярний диференціал 268

Гріна 10, 96, 196, 273, 276

д'Аламбер 126-132

\begin{itemize}
\tightlist
\item
  додатково до ймовірнісного інтеграла 80
\end{itemize}

Неточника 96, 276

-- неперіодичні 77 -- періодичні 77 -- Хевісайд 82, 115, 116, 365.
Фур'є-образ 88, 89, 91, 93

Формула д'Аламбера 180, 183

--- Лейбніц 184

Джерела 99, 100

\begin{itemize}
\tightlist
\item
  Пуассон 180
\end{itemize}

--- --- інтеграл 248, 249, 251

--- Родріга 268

--- Ейлер 84, 85 Фундаментальне рішення 37 Функціональність 329

Функція Бесселя 224, 225

Чисельні методи 9 Кількість персемних 11

Рівнопотенціальні лінії 277

Ядро перетворення 73 -- перетворення Пуассона 252

\section{ЗМІСТ}\label{ux437ux43cux456ux441ux442}

{\def\LTcaptype{none} % do not increment counter
\begin{longtable}[]{@{}
  >{\raggedright\arraybackslash}p{(\linewidth - 4\tabcolsep) * \real{0.1856}}
  >{\raggedright\arraybackslash}p{(\linewidth - 4\tabcolsep) * \real{0.1237}}
  >{\raggedright\arraybackslash}p{(\linewidth - 4\tabcolsep) * \real{0.6907}}@{}}
\toprule\noalign{}
\begin{minipage}[b]{\linewidth}\raggedright
З редагування
\end{minipage} & \begin{minipage}[b]{\linewidth}\raggedright
Гора
\end{minipage} & \begin{minipage}[b]{\linewidth}\raggedright
Переклад
\end{minipage} \\
\midrule\noalign{}
\endhead
\bottomrule\noalign{}
\endlastfoot
Передмова & viye & ****** \\
& & \\
Частина 1. & Введено & \\
Лекція & 1. & Вступ до теорії диференціальних рівнянь з похідними
похідними \\
Частина 2. & Діф & Завдання з синтезу \\
Лекція & 2. & Задачі типу дифузії (параболічні рівняння) \\
Лекція & 3. & Граничні умови в задачах дифузії \\
Лекція & 4. & Виведення рівняння теплопровідності , \\
Лекція & 5. & Розщеплення змінних \\
Лекція & 6. & Перетворення неоднорідних граничних умов у одно- \\
& & Родичі \\
Лекція & 7. & Розв'язання складніших задач методом розділення змін- \\
Тторене \textasciitilde{} & o & Перетворити комплексні рівняння у просту
форму \\
Лекція-лекція & o. & Перетворення комплексних рівнянь у простий
вигляд \\
Мексика & σ, & Власні функції , \\
Pection & 10 & Інтегральні перетворення (перетворення синусів і
козіпус- \\
Боулз & 10. & ) \\
Лекція & 11. & Перетворення рядів і четвірок \\
Лекіня & 12 & Перетворення Фур'є та його застосування до розв'язку \\
* 10HM 11 /1 & , & Рівняння з похідними похідними \\
Лекція & 13. & Перетворення Лапласа \\
Лекція & 14. & Принцип Дюамеля \\
Лекчич & 15. & Принцип Дюамеля \\
& & \\
Частина 3. & Guipe & Проблеми \\
Лекпея & 16. & Одномірне хвильове рівняння (гіперболічні еквалізації) \\
& & \\
Лекція & 17. & Формула Д'Аламбера \\
Лекція & 18. & Формула Д'Аламбера (продовження) \\
Лекція & 19. & Рівняння води та граничні умови \\
Лекція & 20. & Вібрації обмеженої струни (стоячі хвилі) 14 \\
Лекція & 21. & Осциляції пучка (рівняння часткового виробника \\
& & Четвертий порядок) \\
& & Перехід на безрозмірні змінні \\
Лекція & 23. & Класифікація диференціальних рівнянь у похідних \\
& & (Канонічна форма гіперболічного рівняння) 10 \\
Лекція & 24. & Хвильове рівняння у вільному просторі (двовимірне \\
\_ & \textasciitilde{} - & та тривимірні задачі) \\
Лекція & 25. & Скінченні перетворення Фур'є (синус і косинус до - \\
& 00 & Освіта) \\
& & Принцип суперпозиції --- основи теорії динамічних систем 19 \\
& & Рівняння першого порядку (характеристичний метод) 19 \\
Лекція & 28. & Нединамічні рівняння першого порядку (закони
збереження \\
t. & 00 & Ньяя) , \\
Лекція & 29. & Системи диференціальних рівнянь у похідних 2 \\
≢Tection & & Мембранні осциляції (покращення хвиль у Polar 2000
Godinates) \\
& & БУУУУММаалака аа аа аа аа аа \\
\end{longtable}
}

{\def\LTcaptype{none} % do not increment counter
\begin{longtable}[]{@{}
  >{\raggedright\arraybackslash}p{(\linewidth - 2\tabcolsep) * \real{0.7658}}
  >{\raggedright\arraybackslash}p{(\linewidth - 2\tabcolsep) * \real{0.2342}}@{}}
\toprule\noalign{}
\begin{minipage}[b]{\linewidth}\raggedright
Частина 4. Еліптичні задачі
\end{minipage} & \begin{minipage}[b]{\linewidth}\raggedright
233
\end{minipage} \\
\midrule\noalign{}
\endhead
\bottomrule\noalign{}
\endlastfoot
Лекція 31. Лапласнан (інтутивна олісанація) & 233240248255 \\
гармоніки) & 264273 \\
Частина 5. Чисельні та наближені методи & 282 \\
Лекція 37. Чисельні розв'язки (еліптичні задачі) & 282289 \\
Lsson) & 296301 \\
Рівняння) & 308316 \\
Монте-Карло & 322329 \\
Лекція 45. Варіаційні методи розв'язання диференціальних рівнянь у
похідних похідних & 338 \\
Лекція 46. Розв'язок рівнянь часткових похідних теоретиків збурень &
346 \\
Лекція 47. Розв'язання диференціальних рівнянь у похідних похідних за
методом конформних відображень & 355 \\
Додаток 1 & 365 \\
Додаток 2 & 371 \\
Spinning 3 & 373 \\
Джерела & 375 \\
Індекс назв & 378 \\
Індекс предметів & 379 \\
\end{longtable}
}

\section{Дорогий
читачу!}\label{ux434ux43eux440ux43eux433ux438ux439-ux447ux438ux442ux430ux447ux443}

Будь ласка, надсилайте свої коментарі щодо змісту приладів, його
конструкції, якості перекладу та іншого за адресою: 129820, Москва,
I-110, GSP. 1-й Рижський провулок, 2, видавництво «Мир».

\section{Фарлоу
Стейпле}\label{ux444ux430ux440ux43bux43eux443-ux441ux442ux435ux439ux43fux43bux435}

\section{ДИФЕРЕНЦІАЛЬНІ РІВНЯННЯ В ПОХІДНИХ ПОХІДНИХ ДЛЯ НАУКОВЦІВ І
ІНЖЕНЕРІВ}\label{ux434ux438ux444ux435ux440ux435ux43dux446ux456ux430ux43bux44cux43dux456-ux440ux456ux432ux43dux44fux43dux43dux44f-ux432-ux43fux43eux445ux456ux434ux43dux438ux445-ux43fux43eux445ux456ux434ux43dux438ux445-ux434ux43bux44f-ux43dux430ux443ux43aux43eux432ux446ux456ux432-ux456-ux456ux43dux436ux435ux43dux435ux440ux456ux432}

Науковий редактор М. В. Хатунцева-молодший науковий редактор Е. І.
Качулина Художник В. Е. Карпов Художній редактор В. І. Шаповалов
Технічний редактор Т. А. Максимова Коректор Т. І. Стіфсева

\section{UB No 5289}\label{ub-no-5289}

Набор 27.12.84. Подано до публікації 01.07.85. Формат \(60\times90^{1}\)
є. Папір книги. журналу. Печатка висока. Літературний Гарінтура. Том L2
--- бум. n, Usl. Печ. с. 24. Стан. Cr.-Ott. 24. Уч.-изд. л. 19.69, Ізд.
No 1/3807. Тираж 5 000 примірників. Замовлення L6 241. Ціна 1 ранд. 70
тис.

MIR 129820, GSP, Москва, I-110, 1-й Рижський провулок, 2

Надруковано в Ленінградській друкарні No 6 Ордена Червоного Прапора
Праці Ленінградської асоціації «Технічний король» імені Євга. Соколова
Союзполіграф при Державному комітеті СРСР з видавничої справи,
друкарства та книжкової торгівлі, 193144, Ленінград, вулиця Моісеєнко,
10, з матрицями Ордена Жовтневого Регулюта та Ордена Червоного Прапора
Праці МПО «Перша зразкова друкарня» імені А. А. Жданога
Союзполіграфірому при Державному комітеті СССГ з видавничої, друкарської
та книжкової торгівлі, 113054, Москам, Валовія, 28

\end{document}
